\documentclass[11pt]{article}
\usepackage[T1]{fontenc}
\usepackage{amsmath,amssymb,amsthm}
\usepackage[margin=1in]{geometry}
\usepackage{enumitem}
\usepackage[colorlinks=true,linkcolor=blue,urlcolor=blue]{hyperref}

\newtheorem{theorem}{Theorem}[section]
\newtheorem{proposition}[theorem]{Proposition}
\newtheorem{lemma}[theorem]{Lemma}
\newtheorem{corollary}[theorem]{Corollary}
\newtheorem{definition}[theorem]{Definition}
\newtheorem{warning}[theorem]{Warning}
\newtheorem{firewall}[theorem]{Claim firewall}
\newtheorem{decision}[theorem]{Next decision}
\newtheorem{assessment}[theorem]{Assessment}
\newtheorem{ledger}[theorem]{Acquisition ledger}
\theoremstyle{remark}
\newtheorem{remark}[theorem]{Remark}

\title{Renewal Standard Model--Einstein Continuum Research Notes\\
Thirteenth Cycle, V1}
\author{}
\date{9 September 2026}

\begin{document}
\maketitle
\tableofcontents

\section{Context and purpose of the thirteenth cycle}
\label{sec:v13v1}

These notes continue a long-form attempt to construct a rigorous
continuum limit of a regulated Standard Model coupled to Einstein
gravity.  The intended endpoint is a compatible local continuum
state carrying the gauge, fermion, Higgs, and gravitational sectors,
with a controlled diffeomorphism quotient, anomaly-free physical
Ward identities, Osterwalder--Schrader positivity, and the correct
large-distance massive and massless sectors.

That endpoint has not been reached.  The preceding notes prove many
conditional implications and several no-go statements, but the
physical hypotheses have not all been populated by one regulator.
The purpose of this compact restart is to state the acquired
mathematics without reproducing its proofs, expose the remaining
gates, and continue from the first unresolved local coercivity
problem.

The completed twelfth cycle is frozen as
\[
 \texttt{Renewal\_SM\_Einstein\_Continuum\_Research\_Notes\_V100\_f12.tex}.
                                                                    \tag{1.1}
\]
Its exact record is
\[
\begin{array}{c|c}
\text{lines}&27627\\
\text{bytes}&944425\\
\text{SHA--256}&
\texttt{82DDF1C141E9A72BD9DF6CAF1BE24D10D715B19EF9BA689554272A6FC01A10FB}.
\end{array}                                                     \tag{1.2}
\]
Earlier frozen cycles remain separate archives.  Future versions
append to this compact V1 and replace only the immediate active
predecessor.  At every fifth version the newest Yang--Mills and
Navier--Stokes companions are inspected.  At V100 the active file is
again frozen and the sequence restarts.

\section{Major mathematical results acquired}
\label{sec:v13v1-results}

This section is a map of proved implications.  Every item remains
subject to the hypotheses stated in the archived theorem that
establishes it.

\subsection{Finite regulators, topology, and determinants}

\begin{enumerate}[leftmargin=2.2em]
\item Finite-cutoff renewal maps, slice carriers, Schur--Feshbach
reductions, and determinant regularizations were separated into
explicit algebraic, spectral, and stability hypotheses.
\item Exact orbit--stabilizer weights and full-density coordinate
covariance prevent symmetry factors from being inferred from a
partially compressed diagram.
\item Relative trace-class or Hilbert--Schmidt increments control
Fredholm determinant moduli and phases on zero-free paths.  Local
phase one-forms can be removed by counterterms only when they are
exact; cycle holonomy and determinant-line topology remain separate.
\item Gapped Hermitian carriers give integer sector labels stable
under certified homotopies.  An action cost exceeding sector entropy
produces exponential topological-sector tails.
\item A quantitative Faddeev--Popov gap yields a local
diffeomorphism slice and relative determinant control.  A global
quotient still requires compatible overlap Jacobians, stabilizer and
Gribov multiplicities, and equivariant refinement.
\end{enumerate}

\subsection{Anomaly and Ward structure}

\begin{enumerate}[leftmargin=2.2em]
\item Finite cochain Hodge decompositions separate a regulated anomaly
into a harmonic obstruction and an exact removable part.  Controlled
chain homotopies transport cohomology through refinement.
\item Standard Model representation arithmetic cancels the universal
local anomaly polynomial under the stated charge assignments.  This
does not by itself cancel regulator remainders, torsion anomalies, or
global determinant-line anomalies.
\item A retained harmonic anomaly below a proved integral-lattice
systole vanishes exactly.  The lattice membership, systole, and
torsion checks remain physical inputs.
\item Symmetry-restoring counterterms preserve nonvanishing
conditional normalizers and shell variance bounds when their
multiplicative perturbations satisfy the recorded smallness budgets.
\item Exact or vanishing-defect Ward identities pass through the
normalized complex shell martingale and through convergent joint
states on declared local cores.
\item For closed unbounded BRST derivations, predual-norm convergence
and graph-norm operator control pass a possibly nonzero anomaly
functional to the completed graph domain.  The anomaly is a natural
projective functional under BRST-intertwining local inclusions.
\end{enumerate}

\subsection{Gravity, contours, and normalization}

\begin{enumerate}[leftmargin=2.2em]
\item The bare Euclidean Einstein--Hilbert action is unbounded below
along bounded fixed-volume, fixed-topology conformal oscillations.
It cannot serve without qualification as a positive amplitude base.
\item A coercive real Gaussian conformal reference has a
Sobolev-valued projective continuum.  A naive fourth-order regulator
cannot simultaneously retain nondegenerate fixed modes and provide
the uniform reweighting moment needed to remove the regulator.
\item Independent negative quadratic conformal modes admit an exact
imaginary-axis Gaussian contour.  Orientation normalization gives
projective consistency and an infinite Sobolev-valued amplitude
under a trace condition.
\item General complex contours require real-part coercivity,
singular-divisor avoidance, relative homology, normalized new-mode
factors, and a denominator anchor.  Antilinear contour invariance
gives reality, while positivity requires a separate Gram
factorization.
\item A nonzero normalized phase anchor survives summable
positive-measure, charge-admissibility, determinant-transport, and
holonomy defects.  This constructs a conditional nonzero complex
functional, not yet a positive physical Schwinger measure.
\item Uniform gravitational \(C^{2,\alpha}\) moments and a coercive
action floor yield compact containment and Wasserstein subsequences.
The large Hölder state space must use a separable little-Hölder or
Sobolev screen.
\end{enumerate}

\subsection{Continuum laws and observable reconstruction}

\begin{enumerate}[leftmargin=2.2em]
\item Projectively consistent Gaussian reference laws exist under
covariance consistency and spectral-tail tightness.  Summable
adjacent covariance, density, or conditional-state defects produce
cofinal limits with quantitative tails.
\item Fibre-first compact containment combines a tight Einstein
background law, a conditional field tail, and a partition floor.
Orbit quotient tightness can be proved without a global gauge
section when the metric quotient is Polish.
\item A countable Hilbert-cube observable compactification gives
automatic subsequences.  Countable distance coordinates reconstruct
a bounded Polish quotient metric, its topology, and its Borel laws.
A dense observable ledger upgrades subsequential agreement to
full-sequence convergence.
\item Variance floors or mass on disjoint compact regions prevent a
limiting state from collapsing to a deterministic point.
Superlinear moment bounds pass unbounded stress-energy and Einstein
writers through local convergence.
\item A joint classical--quantum state on
\(C(K)\otimes_{\min}\widetilde{\mathcal K}({\cal H})\) has a
trace-class conditional-density disintegration when a
compact-resolvent energy moment excludes quantum mass at infinity.
\item Compatible local restrictions define a quasi-local state.
Regionwise energy bounds suffice; no impossible uniform bound on
the total infinite-volume energy is required.
\end{enumerate}

\subsection{Local normal algebras and BRST compactness}

The final part of the twelfth cycle established the following chain.

\begin{enumerate}[leftmargin=2.2em]
\item A normal unital inclusion
\(\iota_{m\ell}:{\cal M}_m\to{\cal M}_\ell\) has a positive
contractive preadjoint restriction.  Local predual compactness and
vanishing compatibility defects yield a projectively compatible
normal family by diagonal extraction.  This replaces a
finite-dimensional tensor-shell partial trace and accommodates edge
centres.
\item Sector tightness and compactness inside retained sectors are
independent.  A proper central sector penalty controls the first;
uniform finite-rank normal coarse grainings control the second.
Finite trace or finite weight alone is insufficient, as shown by a
diffuse finite-algebra counterexample.
\item Uniform elliptic form comparison to one compact-resolvent
reference operator, a heat trace, and a partition floor give an
explicit common finite-rank quantum channel and trace-norm
compactness.
\item For compact regulated gauge groups, an invariant reference
operator makes that channel equivariant.  Peter--Weyl charge sectors
then realize both the sector tail and the internal compact core.
Continuum gauge and diffeomorphism groups are not covered by Haar
averaging.
\item A closed nilpotent BRST charge defines a Hodge Laplacian whose
spectral projections reduce the differential.  A harmonic
replacement state makes the finite-rank channel exactly BRST
equivariant; without one the graph-Ward defect is bounded by the
positive-energy tail.
\item Finite-dimensional BRST harmonic state modes can be retained
exactly while compactness is imposed only on the positive Hodge
block.  These state zero modes are distinct from the harmonic
degree-one anomaly cochains.
\item A projected G{\aa}rding estimate transfers a gauge-fixed
Hamiltonian moment to the positive BRST-Hodge moment while leaving
the harmonic kernel uncharged.  Its validity depends on a common
boundary form domain, a positive constrained principal symbol, and
uniform lower-order bounds.
\end{enumerate}

\section{Structural statements preserved by convergence}
\label{sec:v13v1-structure}

The archived arguments prove that the following properties are
closed under the indicated modes of convergence:

\begin{enumerate}[leftmargin=2.2em]
\item Osterwalder--Schrader reflection positivity is weak-star
closed when the reflection and positive-time core converge on the
OS products.
\item Ward identities are closed when the states, derivations, and
residuals converge on a common core; graph-norm convergence extends
this to a closed unbounded domain.
\item Translation invariance and named compact covariance identities
pass pointwise to a quasi-local limit.
\item Uniform cutoff clustering envelopes pass to the limit.
Exponential decay in a specified OS channel gives a spectral
threshold in that channel, not a universal mass gap.  Photon and
graviton channels remain massless.
\item Mean Einstein equations, Schwinger--Dyson identities with
contact terms, and almost-sure residual constraints are distinct
claims.  The archived notes give separate sufficient limit
conditions for each and do not infer one from another.
\item Cross-regulator universality follows when two constructions
share a local observable net, their growing weighted ledgers agree,
and their reconstruction and renormalized-parameter mismatch errors
vanish.
\end{enumerate}

\section{The live conditional theorem chain}
\label{sec:v13v1-chain}

For each fixed region \(m\), the latest construction can be written
as
\[
\begin{split}
&\text{compact Einstein-background core}\\
&\Longrightarrow
\text{uniform projected elliptic comparison}\\
&\Longrightarrow
\text{positive BRST-Hodge energy moment}\\
&\Longrightarrow
\text{finite-rank local predual compactness}\\
&\Longrightarrow
\text{projectively compatible locally normal states}\\
&\Longrightarrow
\text{graph-domain anomalous Ward identity}.
\end{split}                                                     \tag{1.3}
\]
Every arrow is a proved implication under explicit hypotheses.  The
programme remains incomplete because no single coupled regulator has
yet populated every premise.

The V100 companion audit adds a sharper proof strategy.  Yang--Mills
V63 derives an exact endpoint Loewner floor from a factorization and
a constrained alias identity, improving its inherited inverse floor
by a factor of seven.  That numerical factor is specific to the
Yang--Mills fibre.  The transferable point is to prove the strongest
reference-background factorization floor before spending any
background perturbation budget.

\section{Unresolved physical and analytic gates}
\label{sec:v13v1-open}

\begin{ledger}[Requirements still needed for the full continuum]
\label{led:v13v1-open}
A complete construction still requires at least the following
compatible inputs:

\begin{enumerate}
\item one nonlinear gravitational ultraviolet construction, or a
coupled conformal contour with uniform coercivity, divisor avoidance,
and projective normalization;
\item a global or high-probability diffeomorphism quotient atlas with
verified overlaps, stabilizers, Gribov multiplicities, and controlled
background tails;
\item a complete chiral fermion and ghost determinant line with local
and global anomaly cancellation and a nonzero phase anchor;
\item a coupled gauge-fixed SM--Einstein operator with a common
boundary form domain and uniform positive-block coercivity over the
admitted gravitational core;
\item compactness, edge-sector tightness, and local inclusion
compatibility for the physical observable algebras;
\item reflection factorization for the full interacting kernel,
including fermion, ghost, gauge, Higgs, and metric blocks;
\item removal of all Ward, BRST, refinement, and background-tail
defects on a common local core;
\item construction of the desired infinite-volume state, phase or
vacuum selection, massive-channel clustering, and the correct
massless photon and graviton sectors;
\item a cross-regulator universality comparison with matched
renormalized parameters.
\end{enumerate}
\end{ledger}

\begin{warning}[No completion claim]
\label{warn:v13v1-incomplete}
The existence of a quasi-local state under conditional compactness
hypotheses is not the construction of the physical Standard
Model--Einstein continuum.  None of the unresolved inputs above is
silently assumed by this summary.
\end{warning}

\section{Current direction}
\label{sec:v13v1-direction}

The next note will work on the first quantitative link in (1.3).
At one reference background, suppose the positive gauge-fixed
Hessian and the positive BRST-Hodge operator factor as
\[
 K_{x_0}^+={\cal C}_{x_0}^*{\cal C}_{x_0},\qquad
 \Delta_{\mathbb Q}^+={\cal D}^*{\cal D}.                       \tag{1.4}
\]
The immediate task is to turn a lower bound for
\({\cal C}_{x_0}\) on the constrained range of \({\cal D}\) into an
exact Loewner comparison, then quantify how much of that floor
survives a relative form or operator perturbation of the Einstein
background.  The harmonic kernel will remain outside the coercive
estimate.

\begin{assessment}[Status at compact restart]
\label{ass:v13v1-status}
The archive now contains a rigorous conditional route from local
elliptic and BRST data to compatible locally normal continuum
states, together with explicit warnings at every unpopulated
physical gate.  The thirteenth cycle begins at the remaining
factorization and perturbation problem rather than repeating the
earlier compactness, anomaly, or convergence proofs.
\end{assessment}

% =============================================================================
% BEGIN APPENDED THIRTEENTH-CYCLE V2 MATERIAL
% =============================================================================

\section{Factorization floors and background-stable Loewner
coercivity}
\label{sec:v13v2}

The V100 audit imported one proof strategy from Yang--Mills V63:
derive the reference floor from the exact constrained factorization
before estimating background perturbations.  This note proves the
abstract result.  The theorem retains the kernel of the BRST
differential and therefore applies to the positive Hodge comparison
without assigning energy to harmonic modes.

\subsection{Exact reference-background factorization}

Let \({\cal H},{\cal Y},{\cal Z}\) be Hilbert spaces.  Let
\[
 {\cal D}:\operatorname{Dom}{\cal D}\subset{\cal H}
 \longrightarrow{\cal Y}                                      \tag{2.1}
\]
be closed and densely defined.  Any linear constraints have already
been incorporated into its closed graph domain.  Let
\[
 {\cal Y}_{\cal D}:=\overline{\operatorname{Ran}{\cal D}}.       \tag{2.2}
\]
Suppose \(B_0:{\cal Y}_{\cal D}\to{\cal Z}\) is bounded and bounded
below:
\[
 \|B_0y\|_{\cal Z}^2\geq a_0\|y\|_{\cal Y}^2,
 \qquad y\in{\cal Y}_{\cal D},\qquad a_0>0.                     \tag{2.3}
\]
Define
\[
 {\cal C}_0=B_0{\cal D}
 \quad\text{on }\operatorname{Dom}{\cal D}.                     \tag{2.4}
\]

\begin{theorem}[Factorization-to-Loewner theorem]
\label{thm:v13v2-factor}
The operator \({\cal C}_0\) is closed,
\[
 \ker{\cal C}_0=\ker{\cal D},                                  \tag{2.5}
\]
and its closed quadratic form satisfies
\[
 \|{\cal C}_0u\|_{\cal Z}^2
 \geq a_0\|{\cal D}u\|_{\cal Y}^2,
 \qquad u\in\operatorname{Dom}{\cal D}.                         \tag{2.6}
\]
Equivalently,
\[
 {\cal C}_0^*{\cal C}_0
 \succeq a_0{\cal D}^*{\cal D}                                 \tag{2.7}
\]
in quadratic-form order.  The constant is the constrained squared
minimum modulus
\[
 a_0=
 \inf_{\substack{y\in{\cal Y}_{\cal D}\\\|y\|=1}}
 \|B_0y\|^2                                                     \tag{2.8}
\]
when the infimum in (2.8) is used as the definition and is positive.
\end{theorem}

\begin{proof}
Equation (2.6) follows by applying (2.3) to
\(y={\cal D}u\).  It also shows that
\({\cal C}_0u=0\) if and only if \({\cal D}u=0\), proving (2.5).

To prove closedness, let \(u_n\to u\) in \({\cal H}\) and
\({\cal C}_0u_n\to z\) in \({\cal Z}\).  The lower bound gives
\[
 \|{\cal D}u_n-{\cal D}u_k\|
 \leq a_0^{-1/2}
       \|{\cal C}_0u_n-{\cal C}_0u_k\|.                         \tag{2.9}
\]
Hence \({\cal D}u_n\) is Cauchy and converges to some
\(y\in{\cal Y}_{\cal D}\).  Closedness of \({\cal D}\) gives
\(u\in\operatorname{Dom}{\cal D}\) and
\({\cal D}u=y\).  Continuity of \(B_0\) then gives
\(z=B_0{\cal D}u={\cal C}_0u\).  Thus \({\cal C}_0\) is closed.

The nonnegative self-adjoint operators
\({\cal C}_0^*{\cal C}_0\) and
\({\cal D}^*{\cal D}\) are those represented by the closed forms
on the two sides of (2.6).  The form inequality is precisely (2.7).
Finally, (2.8) is the optimal constant in (2.3).
\end{proof}

Exact block constraints matter through the range
\({\cal Y}_{\cal D}\).  Estimating \(B_0\) on all of \({\cal Y}\)
can produce a much smaller floor than its minimum modulus on the
actual constrained range.  This is the abstract analogue of the
alias constraint used in Yang--Mills V63.

If
\[
 \Delta_{\mathbb Q}={\cal D}^*{\cal D},\qquad
 K_0={\cal C}_0^*{\cal C}_0,                                   \tag{2.10}
\]
then (2.5) says that the Hodge kernel is retained exactly.  No
inverse is taken on that kernel.

\subsection{Additivity of positive structural blocks}

\begin{proposition}[Loewner floors add]
\label{prop:v13v2-add}
Let \(q_j\) be nonnegative closed forms on one common form domain
and suppose
\[
 q_j[u]\geq a_j\|{\cal D}u\|^2,\qquad a_j\geq0,
 \quad j=0,\ldots,N.                                           \tag{2.11}
\]
Then
\[
 \sum_{j=0}^Nq_j[u]
 \geq\left(\sum_{j=0}^Na_j\right)\|{\cal D}u\|^2.               \tag{2.12}
\]
\end{proposition}

\begin{proof}
Sum the inequalities (2.11).  All terms are defined on the common
form domain, so no operator-domain multiplication is involved.
\end{proof}

This elementary statement prevents a bookkeeping loss.  When a
reference Hessian is a sum of positive Gram blocks, their certified
Loewner floors should be added before inversion or heat-trace
estimation.  In Yang--Mills V63 this is how the owner block and
Wilson block combine; the numerical constants are not transferred
to the present theory.

\subsection{Sharp relative-form perturbation}

Let \(k_x\) be a closed semibounded form on
\(\operatorname{Dom}{\cal D}\) of the form
\[
 k_x[u]=\|{\cal C}_0u\|^2+r_x[u].                               \tag{2.13}
\]

\begin{theorem}[Relative-form survival of the exact floor]
\label{thm:v13v2-form-perturb}
Suppose
\[
 r_x[u]\geq
 -\eta_x\|{\cal D}u\|^2-b_x\|u\|^2,
 \qquad 0\leq\eta_x<a_0,\quad b_x\geq0.                         \tag{2.14}
\]
Then
\[
 k_x[u]\geq
 (a_0-\eta_x)\|{\cal D}u\|^2-b_x\|u\|^2,                       \tag{2.15}
\]
or, in form order,
\[
 K_x+b_xI\succeq
 (a_0-\eta_x){\cal D}^*{\cal D}.                               \tag{2.16}
\]
The kernel of \({\cal D}\) receives no positive penalty from the
right side.
\end{theorem}

\begin{proof}
Insert (2.6) in (2.13) and then use (2.14):
\[
 k_x[u]\geq
 a_0\|{\cal D}u\|^2-\eta_x\|{\cal D}u\|^2-b_x\|u\|^2.
\]
This is (2.15); the representation theorem for closed forms gives
(2.16).
\end{proof}

The budget \(\eta_x\) is subtracted once from the exact structural
floor.  Replacing both terms by unrelated entrywise bounds before
this subtraction can lose the usable scale.

\subsection{An operator-level perturbation bound}

Suppose instead that the background changes the factor itself:
\[
 {\cal C}_x={\cal C}_0+E_x                                     \tag{2.17}
\]
on \(\operatorname{Dom}{\cal D}\), with
\[
 \|E_xu\|
 \leq\varepsilon_x\|{\cal D}u\|+c_x\|u\|,
 \qquad 0\leq\varepsilon_x<\sqrt{a_0}.                         \tag{2.18}
\]

\begin{lemma}[Two-scale reverse-triangle inequality]
\label{lem:v13v2-two-scale}
For \(A,B\geq0\) and \(0<\theta<1\),
\[
 \max\{A-B,0\}^2
 \geq(1-\theta)A^2-(\theta^{-1}-1)B^2.                         \tag{2.19}
\]
\end{lemma}

\begin{proof}
If \(A\geq B\), use
\(2AB\leq\theta A^2+\theta^{-1}B^2\) in
\((A-B)^2\).  If \(A<B\), the right side of (2.19) is nonpositive
because
\[
 (1-\theta)A^2<
 (1-\theta)B^2
 \leq(\theta^{-1}-1)B^2.
\]
Thus (2.19) holds in both cases.
\end{proof}

\begin{theorem}[Factor perturbation with explicit constants]
\label{thm:v13v2-operator-perturb}
Under (2.3), (2.17), and (2.18), for every
\(0<\theta<1\),
\[
 \|{\cal C}_xu\|^2
 \geq a_x(\theta)\|{\cal D}u\|^2-b_x(\theta)\|u\|^2,            \tag{2.20}
\]
where
\[
 \begin{split}
 a_x(\theta)
 &:=(1-\theta)(\sqrt{a_0}-\varepsilon_x)^2>0,\\
 b_x(\theta)
 &:=(\theta^{-1}-1)c_x^2.                                     \tag{2.21}
 \end{split}
\]
If \(c_x=0\), the sharper bound
\[
 \|{\cal C}_xu\|^2
 \geq(\sqrt{a_0}-\varepsilon_x)^2
      \|{\cal D}u\|^2                                         \tag{2.22}
\]
holds.
\end{theorem}

\begin{proof}
The reverse triangle inequality, (2.3), and (2.18) give
\[
 \|{\cal C}_xu\|
 \geq
 (\sqrt{a_0}-\varepsilon_x)\|{\cal D}u\|-c_x\|u\|.              \tag{2.23}
\]
The left side is nonnegative, so apply
Lemma~\ref{lem:v13v2-two-scale} with
\[
 A=(\sqrt{a_0}-\varepsilon_x)\|{\cal D}u\|,
 \qquad B=c_x\|u\|.
\]
This proves (2.20)--(2.21).  If \(c_x=0\), square (2.23) directly.
\end{proof}

The parameter \(\theta\) trades positive Hodge coercivity against a
zeroth-order shift.  It should be selected before the Gibbs
partition estimate, because \(b_x(\theta)\) enters the heat-trace
constant exponentially while \(a_x(\theta)\) controls the spectral
tail.

\subsection{A uniform Einstein-background core}

Let \((X,d_X)\) be a set of admitted classical backgrounds, with
reference point \(x_0\).  Suppose
\[
 \varepsilon_x\leq L_1d_X(x,x_0),\qquad
 c_x\leq L_0d_X(x,x_0).                                       \tag{2.24}
\]
For a radius
\[
 0<r<\frac{\sqrt{a_0}}{L_1}                                   \tag{2.25}
\]
and fixed \(0<\theta<1\), define
\[
 \begin{split}
 a_r&=(1-\theta)(\sqrt{a_0}-L_1r)^2,\\
 b_r&=(\theta^{-1}-1)L_0^2r^2.                                \tag{2.26}
 \end{split}
\]

\begin{corollary}[Uniform structural floor on a background ball]
\label{cor:v13v2-background}
For every \(x\) with \(d_X(x,x_0)\leq r\),
\[
 {\cal C}_x^*{\cal C}_x+b_rI
 \succeq a_r{\cal D}^*{\cal D}                                \tag{2.27}
\]
in form sense.  If local states satisfy
\[
 \sup_{x,n}
 \operatorname{Tr}(\rho_{x,n}{\cal C}_x^*{\cal C}_x)\leq C_r,
                                                                    \tag{2.28}
\]
then
\[
 \sup_{x,n}
 \operatorname{Tr}(\rho_{x,n}{\cal D}^*{\cal D})
 \leq\frac{C_r+b_r}{a_r}.                                     \tag{2.29}
\]
The zero-mode-retaining finite-rank channel of archived V98 has the
uniform trace-norm error
\[
 2\sqrt{\frac{C_r+b_r}{a_rR}}
 +2\frac{C_r+b_r}{a_rR}.                                      \tag{2.30}
\]
\end{corollary}

\begin{proof}
Equations (2.24)--(2.26) inserted in
Theorem~\ref{thm:v13v2-operator-perturb} give (2.27).
Apply the normal state to increasing spectral truncations of that
form inequality to obtain (2.29).  The zero-mode-safe channel bound
from archived V98 with
\(L=(C_r+b_r)/a_r\) gives (2.30).
\end{proof}

For a random Einstein background, (2.30) controls the conditional
quantum fibre only on the selected ball.  The complement must be
handled by the gravitational compact-containment tail; it cannot be
absorbed by taking \(r\) beyond the threshold (2.25).

\subsection{What must be computed in the coupled regulator}

The theorem turns the next physical calculation into four auditable
stocks:

\begin{enumerate}
\item an exact reference factorization
\({\cal C}_0=B_0{\cal D}\) on the constrained positive BRST-Hodge
domain;
\item the constrained minimum modulus \(a_0\), obtained on
\(\overline{\operatorname{Ran}{\cal D}}\) rather than the ambient
coefficient space;
\item background Lipschitz constants \(L_1,L_0\) after one fixed
Hilbert-space and boundary-domain trivialization; and
\item a gravitational probability tail outside a radius satisfying
(2.25).
\end{enumerate}

A failure of the third stock should send the argument upstream to
the background trivialization or slice atlas.  Weakening \(a_0\)
until it covers a discontinuous domain change would be circular.

\begin{assessment}[Status after thirteenth-cycle V2]
\label{ass:v13v2-status}
An exact constrained factorization now yields the optimal reference
Loewner floor, preserves the BRST kernel, and survives relative-form
or factor perturbations with explicit constants.  A background ball
strictly below the radius (2.25) gives the positive Hodge moment and
finite-rank compactness error required by the archived local-normal
closure theorem.
\end{assessment}

\begin{firewall}[Claim boundary after thirteenth-cycle V2]
\label{fire:v13v2-claim}
No coupled SM--Einstein factorization, constrained minimum modulus,
background Lipschitz stock, or probability-tail radius has been
populated.  The Yang--Mills factor-seven constant is not used.
The theorem does not establish anomaly cancellation, determinant
gluing, reflection positivity, clustering, or the full continuum.
\end{firewall}

\subsection{Next bottleneck}

The next step is to make the constrained minimum modulus stable under
a changing range.  In the coupled problem, the BRST-positive
subspace and boundary constraint projector may vary with the
Einstein background.  A fixed coefficient bound is insufficient if
the subspace rotates.  The next note should prove a gap-metric
estimate for the minimum modulus under simultaneous perturbation of
\(B_0\) and the range projection, with a separate failure mode when
the cohomology dimension jumps.

\subsection{Parent record}

This V2 note appends to the validated thirteenth-cycle V1 file with
record
\[
\begin{array}{c|c}
\text{lines}&354\\
\text{bytes}&15872\\
\text{SHA--256}&
\texttt{0AD73287D8A6EDDB34F50397A4A3DA2197174334B03FA4559B17EAE55865A20A}.
\end{array}
\]

% =============================================================================
% END APPENDED THIRTEENTH-CYCLE V2 MATERIAL
% =============================================================================

% =============================================================================
% BEGIN APPENDED THIRTEENTH-CYCLE V3 MATERIAL
% =============================================================================

\section{Rotating constrained ranges and stability of the minimum
modulus}
\label{sec:v13v3}

V2 keeps the constrained range fixed while the factor changes.  In a
coupled Einstein background, the positive BRST range and the
boundary constraint projector may also rotate.  This note measures
that rotation by the gap of orthogonal projections and proves a
minimum-modulus estimate after canonical transport.  It also shows
why a cohomology-dimension jump cannot be hidden in a small
perturbation constant.

\subsection{Direct rotation of nearby subspaces}

Let \({\cal Y}\) be a Hilbert space and let \(P_0,P_x\) be
orthogonal projections with ranges
\[
 {\cal Y}_0=P_0{\cal Y},\qquad {\cal Y}_x=P_x{\cal Y}.            \tag{3.1}
\]
Set
\[
 \delta_x:=\|P_x-P_0\|.                                        \tag{3.2}
\]

\begin{lemma}[Canonical direct rotation]
\label{lem:v13v3-rotation}
If \(\delta_x<1\), there is a unitary \(U_x\) on \({\cal Y}\)
such that
\[
 U_xP_0U_x^*=P_x                                               \tag{3.3}
\]
and
\[
 \|U_x-I\|
 \leq u(\delta_x):=
 \sqrt{\,2-2\sqrt{1-\delta_x^2}\,}.                             \tag{3.4}
\]
In particular, \(U_x:{\cal Y}_0\to{\cal Y}_x\) is unitary.
\end{lemma}

\begin{proof}
Define
\[
 S_x=P_xP_0+(I-P_x)(I-P_0).                                    \tag{3.5}
\]
The hypothesis \(\|P_x-P_0\|<1\) implies that \(S_x\) is
invertible.  Its polar part
\[
 U_x=S_x(S_x^*S_x)^{-1/2}                                      \tag{3.6}
\]
is unitary and intertwines the two projections, which proves
(3.3).  In the standard two-projection decomposition, every
nontrivial block is a two-dimensional principal-angle block with
angle \(0\leq\vartheta\leq\pi/2\), and
\[
 \|P_x-P_0\|_{\rm block}=\sin\vartheta,\qquad
 \|U_x-I\|_{\rm block}=2\sin(\vartheta/2).                      \tag{3.7}
\]
The largest principal angle satisfies
\(\sin\vartheta_{\max}=\delta_x\).  Hence
\[
 2\sin(\vartheta_{\max}/2)
 =\sqrt{2-2\cos\vartheta_{\max}}
 =\sqrt{2-2\sqrt{1-\delta_x^2}},
\]
which gives (3.4).  The trivial common and orthogonal blocks add no
larger norm.
\end{proof}

The threshold \(\delta_x<1\) is structural.  Below it, the
constraint ranges are in the same local Grassmann chart.  At the
threshold, the graph representation and direct rotation may fail.

\subsection{Intrinsic transported factor bound}

Let \(B_0:{\cal Y}_0\to{\cal Z}\) and
\(B_x:{\cal Y}_x\to{\cal Z}\) be bounded operators.  Define their
minimum moduli
\[
 m_0=
 \inf_{\substack{y\in{\cal Y}_0\\\|y\|=1}}\|B_0y\|,
 \qquad
 m_x=
 \inf_{\substack{y\in{\cal Y}_x\\\|y\|=1}}\|B_xy\|.             \tag{3.8}
\]
Let \(U_x\) be the direct rotation in
Lemma~\ref{lem:v13v3-rotation}, and define the intrinsic transported
defect
\[
 e_x:=\|B_xU_x|_{{\cal Y}_0}-B_0\|.                            \tag{3.9}
\]

\begin{theorem}[Stability on a rotating constrained range]
\label{thm:v13v3-minmod}
If \(\delta_x<1\), then
\[
 m_x\geq m_0-e_x.                                              \tag{3.10}
\]
In particular, if \(m_0>0\) and \(e_x<m_0\), then
\[
 B_x^*B_x\succeq(m_0-e_x)^2I_{{\cal Y}_x}.                     \tag{3.11}
\]
\end{theorem}

\begin{proof}
Every unit vector \(y_x\in{\cal Y}_x\) has the unique form
\(y_x=U_xy_0\) with \(y_0\in{\cal Y}_0\) and \(\|y_0\|=1\).
Therefore
\[
 \begin{split}
 \|B_xy_x\|
 &=\|B_xU_xy_0\|\\
 &\geq\|B_0y_0\|
       -\|(B_xU_x-B_0)y_0\|
 \geq m_0-e_x.
 \end{split}                                                    \tag{3.12}
\]
Take the infimum over \(y_x\).  Squaring the positive lower bound
gives (3.11).
\end{proof}

The transported norm (3.9) is the correct stock when the coefficient
spaces form a nontrivial bundle.  It compares operators only after
their constrained domains have been identified.

\subsection{An ambient sufficient estimate}

Sometimes \(B_0\) and \(B_x\) extend to bounded operators
\(\widetilde B_0,\widetilde B_x:{\cal Y}\to{\cal Z}\).  If
\[
 \|\widetilde B_x-\widetilde B_0\|\leq\varepsilon_x,\qquad
 \|\widetilde B_0\|\leq M_0,                                   \tag{3.13}
\]
then
\[
 \begin{split}
 e_x
 &\leq
 \|(\widetilde B_x-\widetilde B_0)U_xP_0\|
 +\|\widetilde B_0(U_x-I)P_0\|\\
 &\leq\varepsilon_x+M_0u(\delta_x).                            \tag{3.14}
 \end{split}
\]

\begin{corollary}[Gap-metric lower floor]
\label{cor:v13v3-gap}
Under (3.13),
\[
 m_x\geq
 m_0-\varepsilon_x-M_0u(\delta_x).                             \tag{3.15}
\]
A positive constrained floor is therefore guaranteed whenever
\[
 \varepsilon_x+M_0u(\delta_x)<m_0.                             \tag{3.16}
\]
\end{corollary}

\begin{proof}
Combine (3.14) with
Theorem~\ref{thm:v13v3-minmod}.
\end{proof}

There is also a direct bound that avoids choosing \(U_x\):
\[
 m_x\geq
 m_0\sqrt{1-\delta_x^2}-M_0\delta_x-\varepsilon_x.              \tag{3.17}
\]
Indeed, for a unit \(y\in{\cal Y}_x\),
\[
 \|(I-P_0)y\|\leq\delta_x,\qquad
 \|P_0y\|\geq\sqrt{1-\delta_x^2},                               \tag{3.18}
\]
and one applies the lower bound to \(B_0P_0y\), while estimating
\(B_0(I-P_0)y\) and \((B_x-B_0)y\) in norm.  The transported
estimate (3.15) is preferable when \(B_0\) is intrinsically defined
only on \({\cal Y}_0\).

\subsection{Application to moving BRST factorizations}

Let
\[
 {\cal D}_x:\operatorname{Dom}{\cal D}_x\subset{\cal H}
 \longrightarrow{\cal Y}
\]
be closed, with
\[
 {\cal Y}_x=\overline{\operatorname{Ran}{\cal D}_x},\qquad
 P_x=\text{the orthogonal projection onto }{\cal Y}_x.          \tag{3.19}
\]
Suppose the positive gauge-fixed factor is
\[
 {\cal C}_x=B_x{\cal D}_x.                                     \tag{3.20}
\]

\begin{corollary}[Moving-range Loewner comparison]
\label{cor:v13v3-loewner}
If \(\delta_x<1\) and \(e_x<m_0\), then
\[
 {\cal C}_x^*{\cal C}_x
 \succeq(m_0-e_x)^2{\cal D}_x^*{\cal D}_x                      \tag{3.21}
\]
in form sense, and
\[
 \ker{\cal C}_x=\ker{\cal D}_x.                                \tag{3.22}
\]
\end{corollary}

\begin{proof}
For \(u\in\operatorname{Dom}{\cal D}_x\), its image
\({\cal D}_xu\) lies in \({\cal Y}_x\).  Apply (3.11):
\[
 \|{\cal C}_xu\|^2
 =\|B_x{\cal D}_xu\|^2
 \geq(m_0-e_x)^2\|{\cal D}_xu\|^2.
\]
This proves (3.21).  The positive lower bound also proves (3.22).
\end{proof}

This is the moving-range replacement for (2.7).  It compares the
physical factor to the Hodge operator built from the same
background-dependent differential.  A further graph transport is
needed before different \({\cal D}_x^*{\cal D}_x\) can be compared
on one fixed Hilbert space.

\subsection{Dimension jumps lie at gap one}

\begin{proposition}[Finite-dimensional jump obstruction]
\label{prop:v13v3-jump}
If \(P_0\) and \(P_x\) have finite-dimensional ranges of different
dimensions, then
\[
 \|P_x-P_0\|=1.                                                \tag{3.23}
\]
The same conclusion holds for finite-dimensional kernels when their
orthogonal-complement projections are used.
\end{proposition}

\begin{proof}
Assume, for example,
\(\dim\operatorname{Ran}P_x>\dim\operatorname{Ran}P_0\).
There is a unit vector
\[
 y\in\operatorname{Ran}P_x\cap
       (\operatorname{Ran}P_0)^\perp.
\]
Then
\[
 (P_x-P_0)y=y,
\]
so the norm is at least one.  The difference of two orthogonal
projections always has norm at most one, proving equality.  The
reverse dimension inequality is symmetric.  Kernel projections
satisfy the same argument.
\end{proof}

Thus constant finite cohomology dimension is not a cosmetic
assumption.  A jump forces departure from every gap ball
\(\delta<1\), destroys the direct-rotation chart, and must be
recorded as the refinement obstruction already isolated in the
archived anomaly ledger.

\subsection{Uniformity on a compact background core}

Let \(K\) be a compact set of Einstein backgrounds.  Suppose
\(x\mapsto P_x\) is norm continuous and
\(x\mapsto B_x\) is continuous after direct rotation.  Around each
\(x_0\in K\), choose a gap chart with
\[
 \sup_{x\text{ in chart}} e_x< m_{x_0}.                        \tag{3.24}
\]
A finite subcover gives finitely many positive local floors.  Their
minimum is positive:
\[
 a_K:=
 \min_{\text{charts}}
 \left(\inf_{x\text{ in chart}}(m_{x_0}-e_x)^2\right)>0.        \tag{3.25}
\]
On overlaps, the direct rotations may differ by a unitary of the
reference constrained range.  The minimum modulus is invariant
under that unitary, so (3.25) is chart independent.  Determinant
phases and orientations can still have overlap holonomy; the
positive Loewner floor does not resolve that separate problem.

\begin{assessment}[Status after thirteenth-cycle V3]
\label{ass:v13v3-status}
The constrained minimum modulus is stable under simultaneous
coefficient and range perturbations after canonical direct
rotation.  Equations (3.15) and (3.21) give explicit positive floors,
while a finite cohomology-dimension jump is proved to occur at
projection gap one and therefore cannot be absorbed into the
perturbation budget.
\end{assessment}

\begin{firewall}[Claim boundary after thirteenth-cycle V3]
\label{fire:v13v3-claim}
No norm-continuous RPESM range projection, background direct
rotation, transported factor bound, or constant cohomology dimension
has been constructed.  The theorem treats positive quadratic
factors and does not glue determinant orientations, cancel
anomalies, establish reflection positivity, or construct the full
continuum.
\end{firewall}

\subsection{Next bottleneck}

The next required transport concerns the unbounded differential
itself.  A positive minimum modulus for \(B_x\) controls
\({\cal D}_x^*{\cal D}_x\), but V98 uses one fixed reference
compactness channel.  The next note should prove norm-resolvent and
spectral-projection stability for the positive Hodge operators under
a common graph-domain perturbation, retaining the finite harmonic
bundle and giving an explicit contour-gap condition.

\subsection{Parent record}

This V3 note appends to the validated thirteenth-cycle V2 file with
record
\[
\begin{array}{c|c}
\text{lines}&744\\
\text{bytes}&28405\\
\text{SHA--256}&
\texttt{D5F935B47622B411811A04FB874796EFDC815E7CFE9E5248CB87A36FC8A42B47}.
\end{array}
\]

% =============================================================================
% END APPENDED THIRTEENTH-CYCLE V3 MATERIAL
% =============================================================================

% =============================================================================
% BEGIN APPENDED THIRTEENTH-CYCLE V4 MATERIAL
% =============================================================================

\section{Norm-resolvent and harmonic-projection stability of the
Hodge family}
\label{sec:v13v4}

V3 controls a bounded factor after its constrained range is
transported.  The Hodge differential and its Laplacian are
unbounded, so their compact resolvents and harmonic projections
require a separate argument.  This note gives that argument for a
common transported form domain.

\subsection{Relative form metric}

Let \(A_0\geq0\) and \(A_x\geq0\) be self-adjoint operators on a
Hilbert space \({\cal H}\), with closed forms \(q_0,q_x\) and one
common form domain \({\cal V}\).  Fix \(\mu>0\) and assume
\[
 |q_x[u]-q_0[u]|
 \leq\varepsilon_x\{q_0[u]+\mu\|u\|^2\},
 \qquad u\in{\cal V},\qquad 0\leq\varepsilon_x<1.                \tag{4.1}
\]
Let
\[
 R_0=(A_0+\mu I)^{-1},\qquad
 R_x=(A_x+\mu I)^{-1}.                                         \tag{4.2}
\]

\begin{theorem}[Quantitative norm-resolvent stability]
\label{thm:v13v4-resolvent}
Under (4.1),
\[
 \|R_x-R_0\|
 \leq
 \eta_x:=
 \frac{\varepsilon_x}{\mu(1-\varepsilon_x)}.                   \tag{4.3}
\]
If \(A_0\) has compact resolvent, then so does \(A_x\).
\end{theorem}

\begin{proof}
Use
\[
 \langle u,v\rangle_{0,\mu}
 :=
 q_0[u,v]+\mu\langle u,v\rangle                                \tag{4.4}
\]
as the form-domain inner product.  The symmetric form
\(q_x-q_0\) is represented in this inner product by a bounded
self-adjoint operator \(E_x\) with
\[
 \|E_x\|\leq\varepsilon_x.                                     \tag{4.5}
\]
In form notation,
\[
 A_x+\mu I
 =
 (A_0+\mu I)^{1/2}(I+E_x)(A_0+\mu I)^{1/2}.                    \tag{4.6}
\]
Since \(\varepsilon_x<1\), \(I+E_x\) is invertible and
\[
 R_x
 =
 R_0^{1/2}(I+E_x)^{-1}R_0^{1/2}.                               \tag{4.7}
\]
Therefore
\[
 \begin{split}
 \|R_x-R_0\|
 &\leq
 \|R_0^{1/2}\|^2
 \|(I+E_x)^{-1}-I\|\\
 &\leq
 \frac1\mu\,
 \frac{\varepsilon_x}{1-\varepsilon_x},
 \end{split}                                                    \tag{4.8}
\]
which is (4.3).

If \(A_0\) has compact resolvent, then \(R_0^{1/2}\) is compact.
Equation (4.7) expresses \(R_x\) as a compact--bounded--compact
product, so \(R_x\) is compact.  This is equivalent to compact
resolvent of \(A_x\).
\end{proof}

Thus compact resolvent need only be proved at the reference
background, provided the transported form perturbation remains
strictly below one.

\subsection{Uniform heat-trace transfer}

\begin{proposition}[Heat trace under a relative form perturbation]
\label{prop:v13v4-heat}
Suppose \(A_0\) has compact resolvent.  For every \(t>0\),
\[
 \operatorname{Tr}(e^{-tA_x})
 \leq
 e^{t\varepsilon_x\mu}
 \operatorname{Tr}(e^{-t(1-\varepsilon_x)A_0}).                 \tag{4.9}
\]
In particular, if \(\varepsilon_x\leq\bar\varepsilon<1\) on a
background core and the right side is finite at
\(1-\bar\varepsilon\), then the heat traces are uniformly bounded
on that core.
\end{proposition}

\begin{proof}
Equation (4.1) gives the lower form inequality
\[
 A_x+\mu I\succeq
 (1-\varepsilon_x)(A_0+\mu I).                                 \tag{4.10}
\]
Let the eigenvalues be ordered increasingly with multiplicity.
The min--max principle gives
\[
 \lambda_k(A_x)
 \geq
 (1-\varepsilon_x)\lambda_k(A_0)-\varepsilon_x\mu.              \tag{4.11}
\]
Exponentiate, sum over \(k\), and obtain (4.9).
\end{proof}

This transfers the heat-trace input needed by the Gibbs compactness
argument.  It does not state operator monotonicity of the
exponential; the proof uses eigenvalue min--max.

\subsection{The harmonic spectral contour}

Assume \(A_0\) has a finite-dimensional kernel and a positive
fixed-region gap
\[
 d_0:=\dim\ker A_0<\infty,\qquad
 g_0:=\inf\{\lambda>0:\lambda\in\operatorname{spec}A_0\}>0.
                                                                    \tag{4.12}
\]
For the bounded resolvent \(R_0\), the harmonic eigenvalue is
\[
 r_0=\frac1\mu,
\]
and the rest of its spectrum lies below
\(r_1=(g_0+\mu)^{-1}\).  Set
\[
 \Gamma_\mu:=r_0-r_1
 =\frac{g_0}{\mu(g_0+\mu)}.                                    \tag{4.13}
\]

\begin{theorem}[Harmonic projection stability]
\label{thm:v13v4-harmonic}
If
\[
 \eta_x<\Gamma_\mu/2,                                          \tag{4.14}
\]
the circle centered at \(r_0\) with radius \(\Gamma_\mu/2\)
defines a Riesz projection \(\widehat P_x^0\) of \(R_x\) having
rank \(d_0\), and
\[
 \|\widehat P_x^0-P_0^0\|
 \leq
 \frac{2\eta_x}{\Gamma_\mu-2\eta_x},                            \tag{4.15}
\]
where \(P_0^0\) projects onto \(\ker A_0\).

If additionally
\[
 \dim\ker A_x=d_0,                                              \tag{4.16}
\]
then \(\widehat P_x^0=P_x^0\), the orthogonal projection onto
\(\ker A_x\).  If
\[
 \eta_x<\Gamma_\mu/4,                                          \tag{4.17}
\]
then \(\|P_x^0-P_0^0\|<1\), so the direct rotation of V3 exists.
\end{theorem}

\begin{proof}
Let \(C\) be the circle in the statement and write
\(r=\Gamma_\mu/2\).  Its distance from
\(\operatorname{spec}R_0\) is \(r\).  The spectral variation bound
for bounded self-adjoint operators and (4.3) show that its distance
from \(\operatorname{spec}R_x\) is at least \(r-\eta_x>0\).
The resolvent identity gives, for \(z\in C\),
\[
 \|(z-R_x)^{-1}-(z-R_0)^{-1}\|
 \leq
 \frac{\eta_x}{r(r-\eta_x)}.                                   \tag{4.18}
\]
Integrating around the circle,
\[
 \begin{split}
 \|\widehat P_x^0-P_0^0\|
 &\leq
 \frac1{2\pi}(2\pi r)
 \frac{\eta_x}{r(r-\eta_x)}\\
 &=\frac{\eta_x}{r-\eta_x}
 =\frac{2\eta_x}{\Gamma_\mu-2\eta_x},
 \end{split}
\]
which proves (4.15).  A continuous Riesz projection cannot change
finite rank while the contour remains in the resolvent, so its rank
is \(d_0\).

Every vector in \(\ker A_x\) is an eigenvector of \(R_x\) with
eigenvalue \(r_0\), hence lies inside the contour.  Under (4.16),
this kernel already has the full rank \(d_0\) of
\(\widehat P_x^0\), proving equality.  Finally, (4.17) makes the
right side of (4.15) strictly less than one, so V3 applies.
\end{proof}

The rank condition (4.16) is not implied by a small form
perturbation: zero eigenvalues can move to small positive values.
In the BRST setting it must come from nilpotence, cohomology
transport, and the constant-dimension certificate.  This keeps
analytic spectral stability separate from anomaly cohomology.

\subsection{A differential-level sufficient condition}

Let
\[
 {\cal D}_0,{\cal D}_x:{\cal V}\subset{\cal H}
 \longrightarrow{\cal Y}                                      \tag{4.19}
\]
be closed operators with a common graph domain, and set
\[
 A_0={\cal D}_0^*{\cal D}_0,\qquad
 A_x={\cal D}_x^*{\cal D}_x.                                   \tag{4.20}
\]
Suppose
\[
 \|({\cal D}_x-{\cal D}_0)u\|
 \leq\alpha_x\|{\cal D}_0u\|+\beta_x\|u\|.                     \tag{4.21}
\]
Define
\[
 \rho_x=
 \sqrt{\alpha_x^2+\beta_x^2/\mu},
 \qquad
 \varepsilon_x=2\rho_x+\rho_x^2.                               \tag{4.22}
\]

\begin{proposition}[Graph perturbation implies form perturbation]
\label{prop:v13v4-graph}
For every \(u\in{\cal V}\),
\[
 |q_x[u]-q_0[u]|
 \leq\varepsilon_x\{q_0[u]+\mu\|u\|^2\}.                       \tag{4.23}
\]
Thus all conclusions of V4 hold whenever
\[
 2\rho_x+\rho_x^2<1.                                           \tag{4.24}
\]
\end{proposition}

\begin{proof}
Put \(F_x={\cal D}_x-{\cal D}_0\).  By Cauchy--Schwarz in
\(\mathbb R^2\), (4.21) gives
\[
 \|F_xu\|
 \leq\rho_x
 \{\,\|{\cal D}_0u\|^2+\mu\|u\|^2\,\}^{1/2}.                   \tag{4.25}
\]
Since
\[
 q_x[u]-q_0[u]
 =
 2\operatorname{Re}\langle{\cal D}_0u,F_xu\rangle
 +\|F_xu\|^2,                                                  \tag{4.26}
\]
and
\(\|{\cal D}_0u\|\leq
 \{q_0[u]+\mu\|u\|^2\}^{1/2}\),
equations (4.25)--(4.26) give (4.23).
\end{proof}

For a BRST-Hodge operator one may take the row differential
\[
 {\cal D}_x=
 \begin{pmatrix}{\mathbb Q}_x\\{\mathbb Q}_x^*\end{pmatrix},
 \qquad
 {\cal D}_x^*{\cal D}_x
 =
 \Delta_{{\mathbb Q},x}.                                      \tag{4.27}
\]
The common-domain and adjoint-domain assertions in (4.19) are
substantial boundary conditions, not formal consequences of
coefficient convergence.

\subsection{One fixed compactness channel after transport}

Assume the constant-kernel condition and (4.17).  Let \(U_x\) be
the direct rotation carrying \(P_0^0\) to \(P_x^0\).  Norm-resolvent
stability also transports every isolated finite spectral band.
Let \(C_R\) be a positively oriented rectifiable contour of length
\(L_R\) which encloses the selected reference resolvent eigenvalues,
and let
\[
 d_R=\operatorname{dist}(C_R,\operatorname{spec}R_0)>0.
\]
If \(\eta_x<d_R\), the same Riesz argument gives a finite-rank
projection \(P_{x,R}\) and
\[
 \|P_{x,R}-P_{0,R}\|
 \leq
 \frac{L_R}{2\pi}\,
 \frac{\eta_x}{d_R(d_R-\eta_x)}.                               \tag{4.28}
\]
When the right side is less than one, a direct rotation identifies
the moving low-energy subspace with one reference finite-dimensional
space.

Equation (4.28) should be used only with an explicitly chosen
contour; a cutoff placed on an eigenvalue has \(d_R=0\) and no
projection-stability conclusion.  Smooth heat-kernel coarse
grainings avoid this contour placement problem but are not finite
rank.

\begin{assessment}[Status after thirteenth-cycle V4]
\label{ass:v13v4-status}
A relative common-domain perturbation now gives explicit
norm-resolvent control, preservation of compact resolvent, uniform
heat-trace bounds, and a contour estimate for the harmonic
projection.  A graph perturbation of the BRST row differential
supplies the relative form stock through (4.22)--(4.24).  Constant
cohomology dimension remains a separate discrete gate.
\end{assessment}

\begin{firewall}[Claim boundary after thirteenth-cycle V4]
\label{fire:v13v4-claim}
No common transported domain or graph perturbation estimate has
been proved for the coupled SM--Einstein BRST differential.  The
fixed-region gap \(g_0\) is not an infinite-volume mass gap and is
not claimed uniform in the region.  Kernel constancy, determinant
orientation, anomaly cancellation, reflection positivity, and the
full continuum remain unproved.
\end{firewall}

\subsection{Next bottleneck}

The immediate next step is the required companion audit at V5.
After that audit, the substantive direction is to replace the local
background ball by finitely many gap charts and glue the transported
finite-rank channels.  The gluing must preserve positivity and
unitality and must expose, rather than average away, any harmonic
bundle holonomy or determinant-line obstruction.

\subsection{Parent record}

This V4 note appends to the validated thirteenth-cycle V3 file with
record
\[
\begin{array}{c|c}
\text{lines}&1078\\
\text{bytes}&38927\\
\text{SHA--256}&
\texttt{C7FB0EBCDCA223288486334D7E532F9737435197611C2BCE4FFDC26931ED3E9B}.
\end{array}
\]

% =============================================================================
% END APPENDED THIRTEENTH-CYCLE V4 MATERIAL
% =============================================================================

% =============================================================================
% BEGIN APPENDED THIRTEENTH-CYCLE V5 MATERIAL
% =============================================================================

\section{V5 audit: large global improvements still require local
channel structure}
\label{sec:v13v5}

This is the first mandatory companion audit of the thirteenth cycle.
The newest files found in the shared directory were
\[
\begin{array}{c|c|c|c}
\text{programme}&\text{file version}&\text{bytes}&\text{SHA--256}\\ \hline
\mathrm{YM}&66&819983&
\texttt{3588991D7B16AF87C2064050B7520CE6A15BD47327725D114A9D46ED94559F5E}\\
\mathrm{NS}&26&278283&
\texttt{82C887F8C2C69E4F28FAD7C3DC8DE5B9099CB5A4BF431EBA1FFF3E91935978AD}.
\end{array}
\]
The file named Yang--Mills V66 is byte-for-byte identical to the V65
file, has \(24063\) lines, and contains no V66 section.  Its newest
effective mathematics is therefore V65.  The Navier--Stokes file is
unchanged.

\subsection{New effective Yang--Mills results}

Yang--Mills V64 reruns the differentiated inverse recursion using
the action-specific endpoint floor from V63.  With every coefficient
stock otherwise fixed, the aggregate global derivative majorant
improves by a factor approximately
\[
 117618.060493.                                                 \tag{5.1}
\]
The proof remains exact rational arithmetic, but even the improved
global trapezoid grid is computationally inadmissible.

V65 then groups complete Laurent coefficients in Frobenius norm
before propagating them through the same recursion.  This gives a
second exact improvement factor satisfying
\[
 67171<
 \frac{B_{\Sigma,64}}{B_{\Sigma,65}}
 <67172,                                                        \tag{5.2}
\]
yet the remaining bound is still of order \(10^{22}\), and the
smallest globally certified grid for an error below \(300\) still
has more than \(10^{41}\) points.  V65 therefore directs the next
Yang--Mills step to the actual cyclic trace contractions or local
momentum boxes.

\subsection{Transfer to the present programme}

The transferable conclusion is qualitative and rigorous.  Two
correct global repairs, each worth many orders of magnitude, may
still fail when the final quantity depends on local constrained
channels and cancellations.  The SM--Einstein counterpart is the
background-chart problem left by V4.  Taking the worst projection
gap, worst spectral contour, and worst channel error across all
charts before gluing can destroy a usable estimate.

The next construction will instead form a convex partition-of-unity
gluing of the actual chartwise normal unital completely positive
channels.  Positivity and trace preservation will then hold before
any norm is taken, and the conditional state error will be the
weighted local error.  Overlap disagreement, harmonic-bundle
holonomy, and noninvariance of the partition will remain explicit
defect stocks.

\begin{assessment}[Decision after the V5 audit]
\label{ass:v13v5-decision}
V6 will prove a chartwise channel-gluing theorem over a compact
Einstein-background core.  It will distinguish finite quantum rank
in each fibre from the false claim of finite rank on the full
classical--quantum algebra.  It will also give overlap,
partition-dependence, covariance, and BRST-commutator estimates, and
will identify a frame-independent harmonic replacement density when
the harmonic dimension is constant.
\end{assessment}

\begin{firewall}[Claim boundary after V5]
\label{fire:v13v5-claim}
Neither the Yang--Mills response sign nor its continuum mass gap is
proved by V64--V65.  No Yang--Mills numerical constant is imported
into the SM--Einstein estimates.  The audit does not construct the
Einstein background charts, local channels, invariant partition, or
harmonic bundle.  The full SM--Einstein continuum remains open.
\end{firewall}

\subsection{Parent record}

This V5 note appends to the validated thirteenth-cycle V4 file with
record
\[
\begin{array}{c|c}
\text{lines}&1441\\
\text{bytes}&50150\\
\text{SHA--256}&
\texttt{E33A8A752907B548A4DFB0B043C442334FA1DF36B2240436F83EE25D9B0DB897}.
\end{array}
\]

% =============================================================================
% END APPENDED THIRTEENTH-CYCLE V5 MATERIAL
% =============================================================================

% =============================================================================
% BEGIN APPENDED THIRTEENTH-CYCLE V6 MATERIAL
% =============================================================================

\section{Positive gluing of local compactness channels over the
Einstein-background core}
\label{sec:v13v6}

V4 transports low-energy subspaces only inside a spectral gap chart.
A compact background core generally needs several such charts.
This note glues their actual quantum channels by a partition of
unity.  Convexity preserves complete positivity and unitality
exactly, while overlap and symmetry defects remain visible.

\subsection{Chartwise channel data}

Let \(X\) be a compact metrizable space of admitted Einstein
backgrounds and let
\[
 {\cal A}=C(X;\widetilde{\cal K}({\cal H}))                     \tag{6.1}
\]
be the continuous-field algebra with fibre the unitization of the
compact operators on a separable Hilbert space.  The fixed fibre
description is obtained after the local Hilbert-space
trivializations of V2--V4.

Let \(\{U_i\}_{i=1}^N\) be a finite open cover of \(X\), and let
\(\{\chi_i\}_{i=1}^N\) be a continuous subordinate partition of
unity:
\[
 \chi_i\geq0,\qquad
 \operatorname{supp}\chi_i\subset U_i,\qquad
 \sum_{i=1}^N\chi_i=1.                                        \tag{6.2}
\]
For \(x\in U_i\), suppose
\[
 {\cal T}_{i,x}:{\cal B}({\cal H})\longrightarrow
 {\cal B}({\cal H})                                           \tag{6.3}
\]
is normal, unital, and completely positive, preserves
\(\widetilde{\cal K}({\cal H})\), and depends norm-continuously on
\(x\) after multiplication by \(\chi_i(x)\).  Define
\[
 ({\cal T}F)(x)
 :=
 \sum_{i=1}^N\chi_i(x){\cal T}_{i,x}(F(x)).
                                                                    \tag{6.4}
\]

\begin{theorem}[Partition-of-unity channel gluing]
\label{thm:v13v6-glue}
Equation (6.4) defines a unital completely positive contraction
\({\cal T}:{\cal A}\to{\cal A}\).  It is
\(C(X)\)-linear:
\[
 {\cal T}(fF)=f\,{\cal T}(F),
 \qquad f\in C(X),\quad F\in{\cal A}.                           \tag{6.5}
\]
Every fibre map
\[
 {\cal T}_x=\sum_i\chi_i(x){\cal T}_{i,x}                       \tag{6.6}
\]
is normal.
\end{theorem}

\begin{proof}
The support and continuity assumptions make each term in (6.4) a
continuous section, hence \({\cal T}F\in{\cal A}\).  At matrix
level \(k\), for a positive
\([F_{ab}]\in M_k({\cal A})\), every
\([{\cal T}_{i,x}(F_{ab}(x))]\) is positive.  Its convex combination
with the nonnegative coefficients \(\chi_i(x)\) is positive for
every \(x\).  Thus \({\cal T}\) is completely positive.  Moreover,
\[
 {\cal T}_x(I)=\sum_i\chi_i(x)I=I,
\]
so it is unital and therefore contractive.  Scalar multiplication
commutes with every fibre map, proving (6.5).  A finite convex
combination of normal fibre maps is normal.
\end{proof}

The map has finite quantum rank in each fibre when all chart maps do.
It need not have finite rank as a map on the full algebra
\({\cal A}\), because the classical variable \(x\) remains
continuous.  Classical compactness and quantum no-escape are
different parts of the continuum argument.

\subsection{Conditional-density error}

Let \(\mu\) be a probability measure on \(X\), and let
\(x\mapsto\rho_x\) be a measurable field of density matrices.  The
joint state is
\[
 \Omega(F)=
 \int_X\operatorname{Tr}(\rho_xF(x))\,d\mu(x).                  \tag{6.7}
\]
The preadjoint glued density is
\[
 \rho_x^{\cal T}
 =
 \sum_i\chi_i(x){\cal T}_{i,x*}(\rho_x).                       \tag{6.8}
\]
It is positive, measurable, and has trace one.

\begin{proposition}[Weighted local error]
\label{prop:v13v6-error}
If
\[
 \|\rho_x-{\cal T}_{i,x*}\rho_x\|_1
 \leq e_i(x)
 \quad\text{whenever }\chi_i(x)>0,                             \tag{6.9}
\]
then
\[
 \|\rho_x-\rho_x^{\cal T}\|_1
 \leq\sum_i\chi_i(x)e_i(x),                                    \tag{6.10}
\]
and
\[
 \|\Omega-\Omega\circ{\cal T}\|_{{\cal A}^*}
 \leq
 \int_X\sum_i\chi_i(x)e_i(x)\,d\mu(x).                         \tag{6.11}
\]
\end{proposition}

\begin{proof}
Since \(\sum_i\chi_i=1\),
\[
 \rho_x-\rho_x^{\cal T}
 =
 \sum_i\chi_i(x)
 \{\rho_x-{\cal T}_{i,x*}\rho_x\}.
\]
The trace-norm triangle inequality gives (6.10).  For
\(\|F\|_{\cal A}\leq1\),
\[
 |(\Omega-\Omega\circ{\cal T})(F)|
 \leq
 \int_X\|\rho_x-\rho_x^{\cal T}\|_1\,d\mu(x),
\]
and (6.11) follows.
\end{proof}

If the chart channel is the spectral channel of V98 and its local
positive Hodge moment is \(L_i(x)\), one may take
\[
 e_i(x)=
 2\sqrt{L_i(x)/R_i}+2L_i(x)/R_i.                               \tag{6.12}
\]
Equation (6.10), rather than the maximum over all charts, is the
natural conditional error.  A maximum remains available when one
needs a uniform bound.

\subsection{Overlap disagreement and partition dependence}

For \(x\in X\), let
\[
 I_x=\{i:\chi_i(x)>0\text{ or }\psi_i(x)>0\}                    \tag{6.13}
\]
for two subordinate partitions \(\chi,\psi\) applied to the same
chart maps.  Define the overlap diameter
\[
 d(x)=
 \max_{i,j\in I_x}
 \|{\cal T}_{i,x}-{\cal T}_{j,x}\|_{\rm cb}.                   \tag{6.14}
\]

\begin{proposition}[Gluing ambiguity bound]
\label{prop:v13v6-overlap}
Let \({\cal T}^{\chi}\) and \({\cal T}^{\psi}\) be the two glued
maps.  Then
\[
 \|{\cal T}_x^{\chi}-{\cal T}_x^{\psi}\|_{\rm cb}
 \leq2d(x).                                                    \tag{6.15}
\]
Consequently, for every density matrix \(\rho_x\),
\[
 \|{\cal T}_{x*}^{\chi}\rho_x
      -{\cal T}_{x*}^{\psi}\rho_x\|_1
 \leq2d(x).                                                    \tag{6.16}
\]
If the chart maps agree after the declared transition functions on
every overlap, then the glued channel is independent of the
partition.
\end{proposition}

\begin{proof}
Fix \(j\in I_x\).  Since
\(\sum_i(\chi_i-\psi_i)=0\),
\[
 {\cal T}_x^\chi-{\cal T}_x^\psi
 =
 \sum_{i\in I_x}(\chi_i-\psi_i)
       ({\cal T}_{i,x}-{\cal T}_{j,x}).                         \tag{6.17}
\]
The total variation distance between two probability vectors is at
most two:
\[
 \sum_i|\chi_i(x)-\psi_i(x)|\leq2.
\]
Equations (6.14) and (6.17) prove (6.15).  Preadjoints have the same
operator norm, which gives (6.16).  Exact overlap agreement makes
\(d(x)=0\).
\end{proof}

Thus a partition of unity always produces a legitimate positive
channel.  It produces a canonical global channel only when the
overlap defect vanishes or is proved irrelevant in the
\(R\to\infty\) limit.

\subsection{Harmonic replacement density and holonomy}

Suppose the harmonic projections in a gap chart have constant rank
\(d>0\).  Define the frame-independent density
\[
 \sigma_i^0(x)=\frac1dP_i^0(x).                                \tag{6.18}
\]
If a unitary transition \(V_{ij}(x)\) obeys
\[
 P_i^0(x)=V_{ij}(x)P_j^0(x)V_{ij}(x)^*,                        \tag{6.19}
\]
then
\[
 \sigma_i^0(x)
 =V_{ij}(x)\sigma_j^0(x)V_{ij}(x)^*.                           \tag{6.20}
\]

\begin{proposition}[Frame-independent harmonic channel]
\label{prop:v13v6-harmonic}
Assume both the retained spectral projection and the harmonic
projection transform by (6.19), and form the chart channel using
(6.18).  Then the channels transform by conjugation on overlaps.
After passing to one common fibre trivialization, they agree.
Non-Abelian frame holonomy inside the harmonic block does not change
\(\sigma^0\), and scalar determinant-line phases cancel under
conjugation.
\end{proposition}

\begin{proof}
Functional calculus transports the retained projection by
conjugation.  Equation (6.20) transports the replacement density.
Insert both covariance identities in the spectral-channel formula
\[
 P_RAP_R+\operatorname{Tr}(\sigma^0A)(I-P_R).
\]
Trace invariance under unitary conjugation gives the conjugated
channel formula.  The density \(P^0/d\) is invariant under every
unitary acting within the harmonic block.  A scalar phase acts
trivially by conjugation.
\end{proof}

This choice does not set the physical probability of harmonic or
superselection sectors.  It only chooses where the discarded
high-energy tail is placed by the auxiliary approximation channel,
and that redirected mass vanishes as the tail bound tends to zero.
If a finite-cutoff argument requires specified sector weights, those
weights must instead be supplied as a covariant density-bundle
section.

\subsection{Covariance and BRST defects}

Suppose a group acts on \(X\), permutes the chart indices, and acts
on fibres by automorphisms \(\beta_{g,x}\).  If
\[
 \chi_{gi}(gx)=\chi_i(x),\qquad
 {\cal T}_{gi,gx}\beta_{g,x}
 =
 \beta_{g,x}{\cal T}_{i,x},                                   \tag{6.21}
\]
then the glued channel is equivariant.  This follows by substituting
(6.21) in (6.4) and relabelling the finite sum.

For the BRST calculation, let \(s\) be a derivation on a common
field core.  Define each chart commutator by
\[
 c_i(F):={\cal T}_i(sF)-s({\cal T}_iF).                         \tag{6.22}
\]
The Leibniz rule gives
\[
 [\,{\cal T},s\,]F
 =
 \sum_i\chi_ic_i(F)
 -\sum_i(s\chi_i){\cal T}_iF.                                  \tag{6.23}
\]
Because \(\sum_i s\chi_i=0\), for any active reference chart \(j\),
\[
 [\,{\cal T},s\,]F
 =
 \sum_i\chi_ic_i(F)
 -\sum_i(s\chi_i)({\cal T}_i-{\cal T}_j)F.                     \tag{6.24}
\]
Hence
\[
 \|[\,{\cal T},s\,]F(x)\|
 \leq
 \sum_i\chi_i(x)\|c_i(F)(x)\|
 +
 d(x)\|F(x)\|
 \sum_i|(s\chi_i)(x)|.                                        \tag{6.25}
\]
If \(s\) acts only in the quantum fibre, then \(s\chi_i=0\).  If the
partition is BRST invariant and every chart channel is exactly
BRST equivariant, the glued channel is exactly equivariant.  For a
combined gauge--diffeomorphism differential, the second term in
(6.25) is a genuine stock; it may not be omitted.

\subsection{Use in the local-normal limit}

Let \(X\) be one compact gravitational core carrying background mass
at least \(1-\tau\).  Equations (6.11) and (6.12) control the
conditional quantum approximation on \(X\), while the complement
contributes at most \(2\tau\) to a state-norm estimate for bounded
observables.  Thus
\[
 \|\Omega-\Omega\circ{\cal T}\|
 \leq
 2\tau+
 \int_X\sum_i\chi_i(x)e_i(x)\,d\mu(x).                         \tag{6.26}
\]
As the core expands and the local spectral thresholds increase, the
right side tends to zero when the two tails do.

The background component is still treated in weak or Wasserstein
topology.  Equation (6.26) supplies quantum no-escape conditionally
over that background.  It does not make the continuum family of
Dirac background laws compact in total variation.

\begin{assessment}[Status after thirteenth-cycle V6]
\label{ass:v13v6-status}
Chartwise finite-rank compactness channels can now be glued over a
compact Einstein-background core while preserving complete
positivity, unitality, and trace exactly.  The state error is the
weighted local error (6.10).  Overlap disagreement, partition
dependence, covariance, and BRST failure have explicit bounds, and
the maximally mixed harmonic replacement density removes
frame-holonomy dependence without altering physical sector weights
in the continuum limit.
\end{assessment}

\begin{firewall}[Claim boundary after thirteenth-cycle V6]
\label{fire:v13v6-claim}
No physical gap-chart cover, continuous channel field, covariant
partition, or vanishing overlap defect has been constructed for the
RPESM background space.  The glued map is generally not finite rank
on the full classical--quantum algebra.  Diffeomorphism and
determinant-line holonomy, anomaly cancellation, reflection
positivity, clustering, and the full continuum remain unproved.
\end{firewall}

\subsection{Next bottleneck}

The next step is projective compatibility of the glued channels
between nested spacetime regions.  Even if each regional channel is
positive and BRST equivariant, its low-energy projection need not
commute with restriction.  The next note should quantify the
commutator between local predual restriction and coarse graining,
and show how a summable two-parameter defect permits a diagonal
choice of region and spectral threshold.

\subsection{Parent record}

This V6 note appends to the validated thirteenth-cycle V5 file with
record
\[
\begin{array}{c|c}
\text{lines}&1546\\
\text{bytes}&54478\\
\text{SHA--256}&
\texttt{FFE8EB3EB23537EE18017A26E544698FE61D64DBAE152FF04DD4E66C8AF3037E}.
\end{array}
\]

% =============================================================================
% END APPENDED THIRTEENTH-CYCLE V6 MATERIAL
% =============================================================================

% =============================================================================
% BEGIN APPENDED THIRTEENTH-CYCLE V7 MATERIAL
% =============================================================================

\section{Projective restriction versus local spectral coarse
graining}
\label{sec:v13v7}

The glued channel of V6 is constructed independently in each region.
Its low-energy projection need not commute with restriction to a
smaller region.  Exact commutation is unnecessary for the continuum
limit: if both channels approximate the identity uniformly on the
admitted local states, their naturality defect is automatically
small.  This note proves that observation and a summable
region--threshold diagonal theorem.

\subsection{Abstract local channels}

Let
\[
 R_{m\ell}:({\cal M}_\ell)_*\longrightarrow({\cal M}_m)_*,
 \qquad m\leq\ell,                                             \tag{7.1}
\]
be the contractive predual restrictions of a directed system of
normal unital inclusions.  Let \({\cal K}_m\) be a family of normal
states on \({\cal M}_m\), and assume
\[
 R_{m\ell}{\cal K}_\ell\subset{\cal K}_m.                       \tag{7.2}
\]
For a threshold \(r\), let
\[
 S_{m,r}:({\cal M}_m)_*\longrightarrow({\cal M}_m)_*            \tag{7.3}
\]
be the preadjoint of a normal unital completely positive channel.
Define its uniform state error
\[
 e_m(r):=
 \sup_{\omega\in{\cal K}_m}
 \|S_{m,r}\omega-\omega\|,                                     \tag{7.4}
\]
and assume
\[
 e_m(r)\longrightarrow0\quad\text{as }r\to\infty               \tag{7.5}
\]
for every fixed \(m\).

\begin{proposition}[Automatic asymptotic naturality]
\label{prop:v13v7-natural}
For \(m\leq\ell\), define
\[
 c_{m\ell}(r,s):=
 \sup_{\omega\in{\cal K}_\ell}
 \|R_{m\ell}S_{\ell,s}\omega
       -S_{m,r}R_{m\ell}\omega\|.                              \tag{7.6}
\]
Then
\[
 c_{m\ell}(r,s)\leq e_\ell(s)+e_m(r).                           \tag{7.7}
\]
Hence exact commutation of restriction and coarse graining is not
needed for their commutator to vanish on the admitted state family.
\end{proposition}

\begin{proof}
For \(\omega\in{\cal K}_\ell\), contractivity of \(R_{m\ell}\) and
(7.2) give
\[
\begin{split}
 &\|R_{m\ell}S_{\ell,s}\omega
       -S_{m,r}R_{m\ell}\omega\|\\
 &\quad\leq
 \|R_{m\ell}(S_{\ell,s}\omega-\omega)\|
 +\|R_{m\ell}\omega-S_{m,r}R_{m\ell}\omega\|\\
 &\quad\leq e_\ell(s)+e_m(r).
\end{split}                                                     \tag{7.8}
\]
Take the supremum.
\end{proof}

This estimate is deliberately state restricted.  The channels may
remain far apart in completely bounded norm on the full algebra
while agreeing asymptotically on the energy-controlled state family
that enters the continuum construction.

\subsection{Finite-regulator compatibility}

Let \(\omega_m^{(n)}\in{\cal K}_m\) be finite-regulator local states,
with projective defect
\[
 \delta_{m\ell}^{(n)}
 :=
 \|\omega_m^{(n)}-R_{m\ell}\omega_\ell^{(n)}\|.                 \tag{7.9}
\]
Define the coarse states
\[
 \eta_m^{(n,r)}:=S_{m,r}\omega_m^{(n)}.                         \tag{7.10}
\]

\begin{proposition}[Coarse projective defect]
\label{prop:v13v7-defect}
For any thresholds \(r,s\),
\[
 \|\eta_m^{(n,r)}
       -R_{m\ell}\eta_\ell^{(n,s)}\|
 \leq
 \delta_{m\ell}^{(n)}+c_{m\ell}(r,s)                           \tag{7.11}
\]
and therefore
\[
 \|\eta_m^{(n,r)}
       -R_{m\ell}\eta_\ell^{(n,s)}\|
 \leq
 \delta_{m\ell}^{(n)}+e_m(r)+e_\ell(s).                        \tag{7.12}
\]
\end{proposition}

\begin{proof}
Insert \(S_{m,r}R_{m\ell}\omega_\ell^{(n)}\).  Contractivity of
\(S_{m,r}\) bounds the first difference by
\(\delta_{m\ell}^{(n)}\), and (7.6) bounds the second by
\(c_{m\ell}(r,s)\).  Proposition
\ref{prop:v13v7-natural} gives (7.12).
\end{proof}

Thus incompatible low-energy projections do not introduce a new
independent limit obstruction when their state approximation errors
already vanish.  A separate commutator stock is useful only when one
needs a stronger finite-threshold identity.

\subsection{A summable two-index diagonal}

\begin{theorem}[Region--cutoff--threshold diagonal]
\label{thm:v13v7-diagonal}
Assume:

\begin{enumerate}
\item every \({\cal K}_m\) is relatively compact in predual norm;
\item for each fixed \(m\leq\ell\),
\[
 \delta_{m\ell}^{(n)}\longrightarrow0;                         \tag{7.13}
\]
\item the channel errors satisfy (7.5).
\end{enumerate}

Then there exist an increasing regulator subsequence \(n_j\),
normal states \(\omega_m\), and thresholds \(r_{m,j}\) defined for
\(m\leq j\), such that for every fixed \(m\),
\[
 \eta_m^{(j)}
 :=S_{m,r_{m,j}}\omega_m^{(n_j)}
 \longrightarrow\omega_m                                     \tag{7.14}
\]
in predual norm, the limit family is exactly projective, and the
following estimates hold after passing to a further subsequence if
necessary:
\[
 \|\eta_m^{(j)}-\omega_m\|\leq2^{1-j},
 \qquad m\leq j,                                               \tag{7.15}
\]
\[
 \|\eta_m^{(j)}
       -R_{m\ell}\eta_\ell^{(j)}\|
 \leq3\,2^{-j},
 \qquad m\leq\ell\leq j,                                      \tag{7.16}
\]
and
\[
 \sum_{j\geq m}
 \|\eta_m^{(j+1)}-\eta_m^{(j)}\|<\infty.                       \tag{7.17}
\]
\end{theorem}

\begin{proof}
Relative compactness and diagonal subsequence extraction give
\(n_j\) and normal states \(\omega_m\) such that
\[
 \|\omega_m^{(n_j)}-\omega_m\|\longrightarrow0
\]
for every fixed \(m\).  By (7.13) and contractivity of restriction,
the limit family satisfies
\[
 \omega_m=R_{m\ell}\omega_\ell.                                \tag{7.18}
\]
Pass to a further increasing subsequence so that, simultaneously
for the finitely many indices \(m\leq\ell\leq j\),
\[
 \|\omega_m^{(n_j)}-\omega_m\|\leq2^{-j},
 \qquad
 \delta_{m\ell}^{(n_j)}\leq2^{-j}.                             \tag{7.19}
\]
For every \(m\leq j\), choose \(r_{m,j}\) large enough that
\[
 e_m(r_{m,j})\leq2^{-j}.                                      \tag{7.20}
\]
Equations (7.19)--(7.20) prove (7.15).  Proposition
\ref{prop:v13v7-defect} proves (7.16).

For fixed \(m\) and \(j\geq m\),
\[
 \|\eta_m^{(j+1)}-\eta_m^{(j)}\|
 \leq
 \|\eta_m^{(j+1)}-\omega_m\|
 +\|\eta_m^{(j)}-\omega_m\|
 \leq3\,2^{-j}.                                                \tag{7.21}
\]
The geometric series proves (7.17).
\end{proof}

This theorem gives direct Cauchy witnesses rather than only an
unnamed compactness subsequence.  At stage \(j\), only finitely many
regions and thresholds must be coordinated.

\subsection{Adding overlap and BRST stocks}

Let \(o_m(r)\) bound the V6 overlap or partition ambiguity on the
admitted state family, and let \(w_m(r)\) bound the graph-dual Ward
defect of the corresponding channel.  If
\[
 o_m(r)\to0,\qquad w_m(r)\to0                                  \tag{7.22}
\]
for every fixed \(m\), the thresholds in (7.20) can be chosen to
satisfy simultaneously
\[
 e_m(r_{m,j})+o_m(r_{m,j})+w_m(r_{m,j})
 \leq2^{-j}.                                                   \tag{7.23}
\]
Only finitely many requirements occur at each \(j\), so the same
diagonal proof applies.  The resulting overlap and Ward defects are
summable in \(j\), and V6 together with the graph-domain theorem
summarized in compact V1 passes both structures to the local limit.

A stock that fails to tend to zero cannot be hidden by the diagonal
choice.  In particular, persistent harmonic holonomy, a nonzero
retained anomaly, or a projection-gap crossing remains visible.

\subsection{Classical background tails}

For a noncompact background space, choose compact cores
\(X_j\) with
\[
 \mu_{n_j}(X_j^c)\leq2^{-j}.                                   \tag{7.24}
\]
Perform the finite chart construction of V6 on \(X_j\).  The
bounded-observable contribution from the complement is at most
\(2^{1-j}\), so it can be added to the right sides of
(7.15)--(7.16) without spoiling summability.  Uniformity of the
channel errors is required only on the chosen core at that stage.

This ordering combines the fibre-first compact containment from the
archive with the present two-index diagonal:
\[
 \text{background core }X_j
 \ \longrightarrow\
 \text{finite regional chart cover}
 \ \longrightarrow\
 \text{regional spectral thresholds}
 \ \longrightarrow\
 \text{regulator subsequence}.                                \tag{7.25}
\]
The regulator is selected last so that all finitely many defects at
stage \(j\) are simultaneously small.

\begin{assessment}[Status after thirteenth-cycle V7]
\label{ass:v13v7-status}
Local spectral coarse grainings need not commute exactly with
restriction.  Their naturality defect on the admitted state family
is bounded by the two local approximation errors.  Relative
predual compactness, vanishing regulator compatibility, and
vanishing channel errors admit a common diagonal with summable
state increments and projective defects.  Background, overlap, and
BRST tails can be inserted into the same finite-stage selection.
\end{assessment}

\begin{firewall}[Claim boundary after thirteenth-cycle V7]
\label{fire:v13v7-claim}
No RPESM state families, compact background cores, regional
channels, or numerical defect stocks have been populated.  The
theorem does not remove a persistent anomaly, harmonic holonomy, or
cohomology jump.  It establishes a conditional diagonal mechanism,
not the physical continuum itself.
\end{firewall}

\subsection{Next bottleneck}

The next unresolved structure is reflection positivity.  Convex
channel gluing preserves ordinary positivity, but an arbitrary
channel need not preserve the Osterwalder--Schrader cone.  The next
note should characterize a sufficient reflected Kraus
factorization for each local channel and prove that partition gluing
preserves the OS Gram form only when the background partition and
chart transitions respect time reflection.

\subsection{Parent record}

This V7 note appends to the validated thirteenth-cycle V6 file with
record
\[
\begin{array}{c|c}
\text{lines}&1922\\
\text{bytes}&67042\\
\text{SHA--256}&
\texttt{0288C50B733BF81D6E9F14F637D4CDBC9D640AB98AD999CE0863FDA1D79255D4}.
\end{array}
\]

% =============================================================================
% END APPENDED THIRTEENTH-CYCLE V7 MATERIAL
% =============================================================================

% =============================================================================
% BEGIN APPENDED THIRTEENTH-CYCLE V8 MATERIAL
% =============================================================================

\section{Reflection-factorized compactness channels and OS-positive
gluing}
\label{sec:v13v8}

V6 preserves ordinary positivity because convex combinations of
unital completely positive maps remain unital and completely
positive.  Osterwalder--Schrader positivity is a different cone and
is not preserved by an arbitrary quantum channel.  This note gives
a sufficient reflected factorization, proves its stability under a
reflection-invariant partition of unity, and records the residual
when the factorization or partition is only approximate.

\subsection{Reflection algebra and the channel criterion}

Let \({\cal A}\) be a unital C-star algebra, let
\({\cal A}_+\subset{\cal A}\) be the positive-time algebra, and let
\[
 \Theta:{\cal A}\longrightarrow{\cal A}                        \tag{8.1}
\]
be the norm-isometric conjugate-linear reflection used in the OS
form.  A state \(\omega\) is OS positive when
\[
 \omega(\Theta(F)F)\geq0,\qquad F\in{\cal A}_+.                 \tag{8.2}
\]
By polarization, this is equivalent to positive semidefiniteness of
every finite Gram matrix
\[
 [\,\omega(\Theta(F_a)F_b)\,]_{a,b}.                            \tag{8.3}
\]

\begin{definition}[Reflected channel factorization]
\label{def:v13v8-factor}
A linear map \({\cal T}:{\cal A}\to{\cal A}\) has a finite
reflected factorization on \({\cal A}_+\) if there are linear maps
\[
 L_\alpha:{\cal A}_+\longrightarrow{\cal A}_+,
 \qquad \alpha=1,\ldots,r,                                     \tag{8.4}
\]
such that
\[
 {\cal T}(\Theta(F)G)
 =
 \sum_{\alpha=1}^r
 \Theta(L_\alpha F)L_\alpha G,
 \qquad F,G\in{\cal A}_+.                                      \tag{8.5}
\]
\end{definition}

\begin{theorem}[Reflected factorization preserves the OS cone]
\label{thm:v13v8-preserve}
Let \({\cal T}\) be unital and completely positive and satisfy
(8.5).  If \(\omega\) is OS positive, then
\[
 \omega_{\cal T}:=\omega\circ{\cal T}                           \tag{8.6}
\]
is a normalized positive state and is OS positive.
\end{theorem}

\begin{proof}
Unital complete positivity makes (8.6) a state.  For
\(F\in{\cal A}_+\), (8.5) gives
\[
 \omega_{\cal T}(\Theta(F)F)
 =
 \sum_{\alpha=1}^r
 \omega(\Theta(L_\alpha F)L_\alpha F)\geq0.                    \tag{8.7}
\]
Each term is nonnegative by (8.2).  More generally, for coefficients
\(c_a\), equation (8.5) gives
\[
 \sum_{a,b}\overline{c_a}c_b
 \omega_{\cal T}(\Theta(F_a)F_b)
 =
 \sum_\alpha
 \omega\!\left(
 \Theta\!\left(\sum_ac_aL_\alpha F_a\right)
 \left(\sum_bc_bL_\alpha F_b\right)\right)\geq0,                \tag{8.8}
\]
which proves the full Gram condition.
\end{proof}

Complete positivity alone does not imply (8.5).  The latter is the
extra reflection structure that must be checked for the cutoff
channel.

\subsection{A product construction from one half-space channel}

Suppose the algebra has a commuting time-split realization generated
by \({\cal A}_+\) and
\({\cal A}_-=\Theta({\cal A}_+)\).  Let
\[
 \Phi_+:{\cal A}_+\to{\cal A}_+                                \tag{8.9}
\]
be a unital completely positive map.  Define its reflected
negative-time partner by
\[
 \Phi_-(\Theta(F)):=\Theta(\Phi_+(F)).                          \tag{8.10}
\]
On the split tensor product, define
\[
 {\cal T}^{\rm split}:=\Phi_-\otimes\Phi_+.                     \tag{8.11}
\]

\begin{proposition}[Paired half-space channel]
\label{prop:v13v8-split}
The map (8.11) is unital and completely positive and satisfies
\[
 {\cal T}^{\rm split}(\Theta(F)G)
 =
 \Theta(\Phi_+F)\Phi_+G,
 \qquad F,G\in{\cal A}_+.                                      \tag{8.12}
\]
It therefore preserves every OS-positive state.  If \(\Phi_+\) has
finite rank, then the split channel has finite rank on each split
quantum fibre.
\end{proposition}

\begin{proof}
The tensor product of two unital completely positive maps is unital
and completely positive on the selected C-star tensor product.
Equations (8.10)--(8.11) give (8.12), which is (8.5) with one map
\(L=\Phi_+\).  Apply
Theorem~\ref{thm:v13v8-preserve}.  The tensor product of two
finite-dimensional ranges is finite-dimensional.
\end{proof}

For a spectral compactness channel, this construction requires a
positive-time reference operator and low-energy channel whose
reflected copy gives the negative-time channel.  A spectral
projection of a full operator crossing the reflection plane need not
factor in this way.

\subsection{Reflection-invariant partition gluing}

Let \({\cal T}_i\), \(i=1,\ldots,N\), be unital completely positive
maps satisfying reflected factorizations with maps \(L_{i\alpha}\).
Let \(h_i\) be central positive elements of \({\cal A}_+\) such that
\[
 \Theta(h_i)=h_i,\qquad
 \sum_{i=1}^Nh_i^2=I.                                         \tag{8.13}
\]
Define
\[
 {\cal T}(A)=
 \sum_{i=1}^Nh_i{\cal T}_i(A)h_i.                              \tag{8.14}
\]

\begin{theorem}[OS-positive channel gluing]
\label{thm:v13v8-glue}
The map (8.14) is unital and completely positive and has the
reflected factorization
\[
 {\cal T}(\Theta(F)G)
 =
 \sum_{i,\alpha}
 \Theta(h_iL_{i\alpha}F)\,
 h_iL_{i\alpha}G.                                              \tag{8.15}
\]
Consequently it preserves the OS cone.
\end{theorem}

\begin{proof}
Each map \(A\mapsto h_i{\cal T}_i(A)h_i\) is completely positive.
Their sum is completely positive, and (8.13) gives
\[
 {\cal T}(I)=\sum_i h_i^2=I.
\]
For \(F,G\in{\cal A}_+\), centrality, reflection invariance, and the
chart factorizations give
\[
\begin{split}
 {\cal T}(\Theta(F)G)
 &=
 \sum_{i,\alpha}
 h_i\Theta(L_{i\alpha}F)L_{i\alpha}G h_i\\
 &=
 \sum_{i,\alpha}
 \Theta(h_iL_{i\alpha}F)h_iL_{i\alpha}G,
\end{split}                                                     \tag{8.16}
\]
which is (8.15).  Theorem~\ref{thm:v13v8-preserve} applies.
\end{proof}

For the background algebra \(C(X)\), write
\[
 h_i(x)=\sqrt{\chi_i(x)}.                                      \tag{8.17}
\]
Condition (8.13) means that the partition is invariant under the
background time reflection \(\vartheta\):
\[
 \chi_i(\vartheta x)=\chi_i(x).                                \tag{8.18}
\]
Thus an ordinary partition subordinate to arbitrary gap charts is
not enough.  One needs reflection-stable charts or an independently
constructed invariant refinement.

\subsection{Approximate reflected factorization}

Suppose the chart identity has a remainder
\[
 {\cal T}_i(\Theta(F)F)
 =
 \sum_\alpha\Theta(L_{i\alpha}F)L_{i\alpha}F
 +D_i(F),                                                      \tag{8.19}
\]
with
\[
 \|D_i(F)\|\leq r_i(F).                                        \tag{8.20}
\]

\begin{proposition}[Quantitative OS residual]
\label{prop:v13v8-residual}
For the invariant gluing (8.14) and every OS-positive state
\(\omega\),
\[
 (\omega\circ{\cal T})(\Theta(F)F)
 \geq
 -\sum_i\|h_i\|^2r_i(F).                                      \tag{8.21}
\]
For a classical partition evaluated pointwise, the sharper
conditional residual is
\[
 -\sum_i\chi_i(x)r_i(F;x).                                    \tag{8.22}
\]
\end{proposition}

\begin{proof}
The factorized terms in (8.19) become OS squares exactly as in
(8.16).  The remaining operator is
\(\sum_i h_iD_i(F)h_i\), whose norm is at most the right side of
(8.21).  A state has norm one.  Pointwise central scalar weights
give (8.22).
\end{proof}

If \(r_i(F)\to0\) along the regional threshold diagonal of V7, then
the limiting state is OS positive by weak-star closedness of the OS
cone.

\subsection{Partition and overlap defects}

Let \(\overline{\cal T}\) be a channel obtained from an invariant
partition and exact reflected chart maps, and let \({\cal T}\) be
the implemented glued channel.  If
\[
 \|{\cal T}-\overline{\cal T}\|_{\rm cb}\leq d_{\rm OS},         \tag{8.23}
\]
then for every \(\|F\|\leq1\),
\[
 (\omega\circ{\cal T})(\Theta(F)F)\geq-d_{\rm OS}.              \tag{8.24}
\]
Indeed, the reference term is nonnegative and
\[
 |\,\omega(({\cal T}-\overline{\cal T})(\Theta(F)F))\,|
 \leq d_{\rm OS}\|\Theta(F)F\|
 \leq d_{\rm OS}.                                              \tag{8.25}
\]
The V6 overlap diameter and the distance from an invariant
partition supply explicit contributions to \(d_{\rm OS}\).

Reflection symmetrization cannot always be performed inside one
chart.  If \(\vartheta U_i\) crosses a spectral-gap failure or a
cohomology jump, the reflected channel does not lie in the same
local trivialization.  That event is a geometric obstruction, not a
small OS residual.

\subsection{Compatibility with BRST and sector data}

Suppose each half-space channel \(\Phi_{i,+}\) commutes with the
positive-time BRST differential, its reflected partner commutes on
negative time, and the invariant square roots \(h_i\) are BRST
closed.  Then the glued channel commutes with the full split
differential.  This follows either from (6.23) or directly from the
Leibniz rule applied to (8.14).

The harmonic replacement density should also factor across the
reflection plane or admit a reflected Gram decomposition.  The
frame-independent density \(P^0/d\) from V6 removes internal frame
holonomy but does not imply a time-split factorization of \(P^0\).
Thus BRST-equivariant compactness and OS-preserving compactness are
two compatible but distinct conditions on the reference channel.

\begin{assessment}[Status after thirteenth-cycle V8]
\label{ass:v13v8-status}
A reflected factorization on positive-time products is sufficient
for a unital completely positive channel to preserve the full OS
Gram cone.  Paired half-space channels provide an explicit
construction, and reflection-invariant square-root partitions glue
them without loss.  Approximate factorization, overlap, and
partition errors give explicit OS residuals that can enter the V7
summable diagonal.
\end{assessment}

\begin{firewall}[Claim boundary after thirteenth-cycle V8]
\label{fire:v13v8-claim}
No RPESM half-space spectral operator, reflected finite-rank
channel, invariant gap-chart cover, or split harmonic replacement
density has been constructed.  A generic UCP channel is not asserted
to preserve OS positivity.  The theorem does not prove anomaly
cancellation, determinant-line reflection gluing, clustering, or the
full SM--Einstein continuum.
\end{firewall}

\subsection{Next bottleneck}

The next step is to construct the half-space spectral channel from a
local operator without assuming an exact Hilbert tensor split across
the reflection plane.  A viable route is an OS Gram-kernel
factorization: approximate the cross-plane kernel by finite-rank
positive Gram matrices and control the trace-norm tail.  The next
note should prove that such positive finite-rank kernel
approximations induce reflected channel factorizations and remain
stable under the background chart gluing.

\subsection{Parent record}

This V8 note appends to the validated thirteenth-cycle V7 file with
record
\[
\begin{array}{c|c}
\text{lines}&2231\\
\text{bytes}&77336\\
\text{SHA--256}&
\texttt{A4AEF3BB9E571BFF0C49003C2D118872DC05E643C101DFD4909871040A259FE7}.
\end{array}
\]

% =============================================================================
% END APPENDED THIRTEENTH-CYCLE V8 MATERIAL
% =============================================================================

% =============================================================================
% BEGIN APPENDED THIRTEENTH-CYCLE V9 MATERIAL
% =============================================================================

\section{Positive finite-rank approximation of the OS Gram kernel}
\label{sec:v13v9}

V8 gives a universal channel criterion.  In many constructive
arguments the first available object is weaker: the OS sesquilinear
form of one regulated state is represented by a positive
cross-plane kernel.  This note proves a uniform positive
finite-rank approximation for that kernel and a safe normalization
step.  It also identifies the additional extension theorem needed
before the approximation can be called a channel.

\subsection{A dominated family of positive trace-class kernels}

Let \({\cal E}\) be a separable Hilbert space and let \(H\geq0\) be
trace class.  Choose increasing finite-rank spectral projections
\(P_N\) of \(H\) such that
\[
 P_N\longrightarrow I\quad\text{strongly},\qquad
 \tau_N:=\operatorname{Tr}((I-P_N)H)\longrightarrow0.           \tag{9.1}
\]
Let \(G_x\geq0\), \(x\in X\), be a family of positive trace-class
operators satisfying the uniform Loewner domination
\[
 0\leq G_x\leq C H.                                            \tag{9.2}
\]
Define
\[
 G_{x,N}=P_NG_xP_N.                                             \tag{9.3}
\]

\begin{theorem}[Uniform positive Gram-kernel truncation]
\label{thm:v13v9-kernel}
Every \(G_{x,N}\) is positive and finite rank, and
\[
 \sup_{x\in X}\|G_x-G_{x,N}\|_1
 \leq
 2C\sqrt{\operatorname{Tr}(H)\tau_N}+C\tau_N
 \longrightarrow0.                                            \tag{9.4}
\]
In particular, positivity is preserved before the truncation error
is estimated.
\end{theorem}

\begin{proof}
Write \(P=P_N\), \(Q=I-P\), and
\[
 \delta_{x,N}:=\operatorname{Tr}(QG_xQ)
 =\operatorname{Tr}(QG_x).
\]
Order domination gives
\[
 \delta_{x,N}\leq C\operatorname{Tr}(QH)=C\tau_N,\qquad
 \operatorname{Tr}G_x\leq C\operatorname{Tr}H.                 \tag{9.5}
\]
The trace-norm Cauchy--Schwarz inequality gives
\[
 \|PG_xQ\|_1
 \leq
 \{\operatorname{Tr}(PG_xP)
    \operatorname{Tr}(QG_xQ)\}^{1/2}
 \leq
 C\sqrt{\operatorname{Tr}(H)\tau_N}.                           \tag{9.6}
\]
The same bound holds for \(QG_xP\), while
\(\|QG_xQ\|_1=\delta_{x,N}\leq C\tau_N\).  Since
\[
 G_x-PG_xP=PG_xQ+QG_xP+QG_xQ,
\]
the triangle inequality proves (9.4).  Compression by \(P_N\)
preserves positivity and has finite-dimensional range.
\end{proof}

A uniform trace bound without the order domination (9.2) does not
imply (9.4): the mass can move through an orthonormal sequence.
The common positive majorant is the kernel analogue of the fixed
reference energy used in the compactness argument.

\subsection{OS forms and their finite feature maps}

Let \({\cal C}_+\) be a normed positive-time test space.  Suppose
the regulated OS form at background \(x\) has a representation
\[
 {\mathfrak q}_x(F,G)
 =
 \langle J_xF,G_xJ_xG\rangle_{\cal E},                         \tag{9.7}
\]
where \(J_x:{\cal C}_+\to{\cal E}\) is linear and
\[
 \sup_x\|J_x\|\leq M.                                          \tag{9.8}
\]
Set
\[
 {\mathfrak q}_{x,N}(F,G)
 =
 \langle J_xF,G_{x,N}J_xG\rangle.                              \tag{9.9}
\]

\begin{corollary}[Uniform OS-form error]
\label{cor:v13v9-form}
The form \({\mathfrak q}_{x,N}\) is positive semidefinite and
\[
 |{\mathfrak q}_x(F,G)-{\mathfrak q}_{x,N}(F,G)|
 \leq
 M^2\|G_x-G_{x,N}\|_1\,
 \|F\|\,\|G\|.                                                 \tag{9.10}
\]
Thus the right side tends to zero uniformly in \(x\) on the product
of the two test-space unit balls.
\end{corollary}

\begin{proof}
Positivity follows from
\[
 {\mathfrak q}_{x,N}(F,F)
 =
 \|G_{x,N}^{1/2}J_xF\|^2\geq0.                                \tag{9.11}
\]
For the error, use
\[
 |\langle u,(G_x-G_{x,N})v\rangle|
 \leq\|G_x-G_{x,N}\|\,\|u\|\,\|v\|
 \leq\|G_x-G_{x,N}\|_1\|u\|\|v\|
\]
with \(u=J_xF\), \(v=J_xG\), and (9.8).
\end{proof}

Choose an orthonormal basis \(e_1,\ldots,e_{d_N}\) of
\(P_N{\cal E}\), and diagonalize the positive matrix
\(G_{x,N}\) there:
\[
 G_{x,N}
 =
 \sum_{\alpha=1}^{d_N}
 \lambda_{\alpha,x,N}
 |v_{\alpha,x,N}\rangle\langle v_{\alpha,x,N}|.                \tag{9.12}
\]
Then
\[
 {\mathfrak q}_{x,N}(F,G)
 =
 \sum_{\alpha=1}^{d_N}
 \overline{\ell_{\alpha,x,N}(F)}
 \ell_{\alpha,x,N}(G),                                        \tag{9.13}
\]
where
\[
 \ell_{\alpha,x,N}(F)
 =
 \sqrt{\lambda_{\alpha,x,N}}\,
 \langle v_{\alpha,x,N},J_xF\rangle.                           \tag{9.14}
\]
Equation (9.13) is the finite scalar Gram factorization.  Degenerate
eigenvectors need not be selected continuously: the compressed
positive operator (9.3), rather than its individual eigenvectors,
is the invariant object.

\subsection{Safe normalization}

Suppose
\[
 {\mathfrak q}_x({\bf1},{\bf1})=1                              \tag{9.15}
\]
and define
\[
 z_{x,N}:={\mathfrak q}_{x,N}({\bf1},{\bf1}).                  \tag{9.16}
\]
Let
\[
 \epsilon_N:=
 M^2\sup_x\|G_x-G_{x,N}\|_1\,\|{\bf1}\|^2.                    \tag{9.17}
\]

\begin{proposition}[Positive denominator anchor]
\label{prop:v13v9-normalize}
If \(\epsilon_N<1\), then
\[
 z_{x,N}\geq1-\epsilon_N>0                                    \tag{9.18}
\]
uniformly in \(x\), and
\[
 \widehat{\mathfrak q}_{x,N}
 :=
 z_{x,N}^{-1}{\mathfrak q}_{x,N}                              \tag{9.19}
\]
is a normalized positive semidefinite OS form.  Moreover,
\[
\begin{split}
 |\widehat{\mathfrak q}_{x,N}(F,G)
       -{\mathfrak q}_x(F,G)|
 \leq{}&
 \frac{|{\mathfrak q}_{x,N}(F,G)
             -{\mathfrak q}_x(F,G)|}{1-\epsilon_N}\\
 &+
 \frac{\epsilon_N}{1-\epsilon_N}
 |{\mathfrak q}_x(F,G)|.
                                                                    \tag{9.20}
\end{split}
\]
\end{proposition}

\begin{proof}
Equations (9.10), (9.15), and (9.17) give
\[
 |z_{x,N}-1|\leq\epsilon_N,
\]
which proves (9.18).  Division by a positive scalar preserves
positive semidefiniteness and gives value one at
\(({\bf1},{\bf1})\).  Finally,
\[
 \frac{{\mathfrak q}_{x,N}}{z_{x,N}}-{\mathfrak q}_x
 =
 \frac{{\mathfrak q}_{x,N}-{\mathfrak q}_x}{z_{x,N}}
 +
 \left(\frac1{z_{x,N}}-1\right){\mathfrak q}_x,
\]
and
\[
 z_{x,N}^{-1}\leq(1-\epsilon_N)^{-1},\qquad
 |z_{x,N}^{-1}-1|
 \leq\epsilon_N(1-\epsilon_N)^{-1}.
\]
This proves (9.20).
\end{proof}

The denominator is real and positive because it is a diagonal value
of a positive Gram form.  This is stronger than a nonzero complex
phase anchor, but it applies only after the OS kernel representation
has been established.

\subsection{Background gluing}

Let \(\chi_i\) be a reflection-invariant partition of unity on a
compact background core, and suppose each chart has a positive
kernel \(G_{i,x,N}\).  Define
\[
 G_{x,N}^{\rm gl}
 =
 \sum_i\chi_i(x)G_{i,x,N}.                                    \tag{9.21}
\]
It is positive and finite rank in each fibre.  If all chart errors
are bounded by \(r_i(N,x)\), then
\[
 \|G_x^{\rm gl}-G_{x,N}^{\rm gl}\|_1
 \leq\sum_i\chi_i(x)r_i(N,x).                                 \tag{9.22}
\]
The proof is the same convexity and triangle argument as in V6.

On an overlap with unitary transition \(V_{ij}(x)\), require
\[
 G_{i,x}=V_{ij}(x)G_{j,x}V_{ij}(x)^*,\qquad
 H_i=V_{ij}(x)H_jV_{ij}(x)^*.                                 \tag{9.23}
\]
Functional calculus then transports \(P_N\), and the compressed
kernels transform covariantly.  Scalar phases and changes of basis
inside a degenerate eigenspace disappear under conjugation.  A
failure of (9.23) is an overlap kernel defect and enters (9.22);
it is not cured by diagonalizing each chart separately.

\subsection{State kernel versus universal channel}

The scalar features (9.14) prove positivity of the particular
approximating OS form.  They do not by themselves produce elements
\(L_\alpha F\in{\cal A}_+\) satisfying the operator identity (8.5),
and they do not define a unital completely positive map on the full
observable algebra.

To obtain the universal channel of V8 one needs an operator-valued
Kolmogorov factorization
\[
 {\cal T}(\Theta(F)G)
 =
 \langle{\cal L}F,{\cal L}G\rangle_{\cal E_{\cal A}}            \tag{9.24}
\]
in a Hilbert module whose inner product is a finite sum of elements
\(\Theta(L_\alpha F)L_\alpha G\), together with a completely
positive unital extension from the OS product system to
\({\cal A}\).  Neither extension follows from scalar positivity of
one state.  This distinction prevents a state-dependent Gram
truncation from being advertised as a regulator-independent
physical channel.

\subsection{Insertion into the diagonal limit}

If
\[
 \tau_N\longrightarrow0,\qquad
 \sup_x\|J_x\|\leq M,\qquad
 \epsilon_N<1
\]
on every selected compact background core, then (9.4), (9.10), and
(9.20) give a vanishing OS residual.  Choose \(N=N_j\) in the V7
diagonal so that the residual is at most \(2^{-j}\).  The normalized
finite-rank Gram forms then converge on the positive-time core, and
weak-star closedness of the OS cone gives an OS-positive limiting
state, provided the same limiting state is defined on the full
observable algebra.

\begin{assessment}[Status after thirteenth-cycle V9]
\label{ass:v13v9-status}
A common positive trace-class majorant now yields uniform positive
finite-rank approximations of the state-side OS Gram kernel with the
explicit trace-tail bound (9.4).  The normalization remains positive
under (9.18), chart gluing preserves the Gram cone, and the error can
be made summable in the V7 diagonal.  The distinction between a
state-specific Gram approximation and a universal UCP channel is
explicit.
\end{assessment}

\begin{firewall}[Claim boundary after thirteenth-cycle V9]
\label{fire:v13v9-claim}
No RPESM cross-plane Gram operator, common trace-class majorant,
bounded feature map, or completely positive extension has been
constructed.  Scalar OS-kernel positivity does not establish a
universal reflected channel, determinant-line reflection gluing,
anomaly cancellation, clustering, or the full SM--Einstein
continuum.
\end{firewall}

\subsection{Next bottleneck}

The next required step is the V10 companion audit.  After that audit,
the substantive direction is the completely positive extension
problem in (9.24).  A suitable target is a finite-dimensional OS
operator system generated by the truncated positive-time features:
one should prove when its reflected Gram map is completely positive,
when Arveson extension applies, and what extra condition makes the
extension unital and compatible across background charts.

\subsection{Parent record}

This V9 note appends to the validated thirteenth-cycle V8 file with
record
\[
\begin{array}{c|c}
\text{lines}&2565\\
\text{bytes}&88777\\
\text{SHA--256}&
\texttt{68F2CE3CB5260C740B8B366C023CDD4F2E2617C8CD403534669CA20AB7BBF060}.
\end{array}
\]

% =============================================================================
% END APPENDED THIRTEENTH-CYCLE V9 MATERIAL
% =============================================================================

% =============================================================================
% BEGIN APPENDED THIRTEENTH-CYCLE V10 MATERIAL
% =============================================================================

\section{V10 audit: complete positivity should be certified in the
relative Gram metric}
\label{sec:v13v10}

This is the mandatory companion audit following V5.  The newest
files examined were
\[
\begin{array}{c|c|c|c}
\text{programme}&\text{version}&\text{bytes}&\text{SHA--256}\\ \hline
\mathrm{YM}&68&844092&
\texttt{B0E2441D94D57E3B6BB81042321ED729B5545388A919EB76E4A00A5AE78B415C}\\
\mathrm{NS}&26&278283&
\texttt{82C887F8C2C69E4F28FAD7C3DC8DE5B9099CB5A4BF431EBA1FFF3E91935978AD}.
\end{array}
\]
The Navier--Stokes companion is unchanged, so no prior result is
repeated.

\subsection{Effective Yang--Mills advances V66--V68}

Yang--Mills V66 applies the trace-ideal inequalities
\[
 |\operatorname{Tr}(GQ)|
 \leq\sqrt{45}\|G\|_{\rm op}\|Q\|_F,\qquad
 |\operatorname{Tr}(GP\widetilde GN)|
 \leq\|G\|_{\rm op}\|\widetilde G\|_{\rm op}
      \|P\|_F\|N\|_F,                                         \tag{10.1}
\]
removing essentially the full artificial factor \(45\) from its
bubble-dominated global majorant.  The resulting global grid is
still impossible.

V67 proves a residual-certified center inverse and a rational
Laurent radius for local boxes.  Its uniform absolute-radius pilot
would require more than \(7.3\,10^{17}\) endpoint centers merely to
retain half of the existing floor.  This is a no-go result for that
literal uniform criterion, not for every adaptive method.

V68 then goes upstream.  For a positive Hessian \(H\), its
derivatives \(H_{r,m}\), and inverse derivatives \(G_{r,m}\), it
defines
\[
 T_{r,m}=H^{-1/2}H_{r,m}H^{-1/2},\qquad
 U_{r,m}=H^{1/2}G_{r,m}H^{1/2},                               \tag{10.2}
\]
and proves the exact recursion
\[
 U_{r,m}=
 -\!\!\sum_{\substack{0\leq a\leq r,\ 0\leq b\leq m\\
                     (a,b)\ne(0,0)}}
 {r\choose a}{m\choose b}T_{a,b}U_{r-a,m-b}.                  \tag{10.3}
\]
Consequently the absolute inverse floor occurs only once.  Its
reported mesh reductions are diagnostics; the relative recursion
itself is rigorous.

\subsection{Transfer to the OS extension problem}

V9 identifies complete positivity of an operator-valued extension
as the missing step.  An absolute entrywise perturbation of its Choi
or Gram matrix would repeat the loss exposed by Yang--Mills
V66--V67.  The correct finite-dimensional certificate is relative:
if \(C_0\geq0\) is the reference Choi matrix and
\[
 -\varepsilon C_0\preceq C-C_0\preceq\varepsilon C_0,
 \qquad \varepsilon<1,                                        \tag{10.4}
\]
then
\[
 C\succeq(1-\varepsilon)C_0\geq0.                              \tag{10.5}
\]
This preserves the kernel automatically and does not divide by a
small global eigenvalue.  Unitality must be checked by a separate
linear constraint on the Choi blocks.

\begin{assessment}[Decision after the V10 audit]
\label{ass:v13v10-decision}
V11 will formulate the truncated OS products as a finite-dimensional
operator system.  It will give a Choi-matrix feasibility criterion
for a unital completely positive extension, prove the
relative-Loewner stability rule (10.4)--(10.5), and exhibit the
transpose map as a counterexample showing why scalar Gram
positivity is insufficient.  Arveson extension will be invoked only
after complete positivity and relation compatibility have been
proved on the operator system.
\end{assessment}

\begin{firewall}[Claim boundary after V10]
\label{fire:v13v10-claim}
No Yang--Mills response sign or mass gap is established by V66--V68,
and their diagnostic mesh values are not imported.  No RPESM Choi
matrix, operator system, relative perturbation certificate, or UCP
extension has yet been populated.  The full SM--Einstein continuum
remains open.
\end{firewall}

\subsection{Parent record}

This V10 note appends to the validated thirteenth-cycle V9 file with
record
\[
\begin{array}{c|c}
\text{lines}&2914\\
\text{bytes}&99842\\
\text{SHA--256}&
\texttt{E0B165D68C281FB6DF38A04BE492AB1DE92440684F6060EBC9CC07BB034FF2E6}.
\end{array}
\]

% =============================================================================
% END APPENDED THIRTEENTH-CYCLE V10 MATERIAL
% =============================================================================

% =============================================================================
% BEGIN APPENDED THIRTEENTH-CYCLE V11 MATERIAL
% =============================================================================

\section{The finite OS operator system and its UCP extension}
\label{sec:v13v11}

V9 constructs positive scalar Gram forms for one state.  A universal
compactness channel requires a map on observables which is unital and
completely positive at every matrix level.  This note reduces the
finite-cutoff extension problem to a positive Choi matrix with linear
constraints, and proves a relative-Loewner stability criterion which
retains singular Gram directions.

\subsection{The truncated reflected-product system}

Let \({\cal A}_N\) be a finite-cutoff unital C-star algebra and choose
positive-time features
\[
 F_1,\ldots,F_q\in({\cal A}_N)_+.                              \tag{11.1}
\]
Let \({\cal S}_N\subset{\cal A}_N\) be the finite-dimensional
operator system generated by
\[
 I,\qquad X_{ab}:=\Theta(F_a)F_b,\qquad X_{ab}^*,
 \quad 1\leq a,b\leq q.                                       \tag{11.2}
\]
An intended reflected factorization prescribes
\[
 \Gamma_0(X_{ab})
 =
 \sum_{\alpha=1}^r
 \Theta(L_\alpha F_a)L_\alpha F_b.                            \tag{11.3}
\]

The first gate is well-definedness.  Whenever
\[
 c_0I+\sum_{a,b}c_{ab}X_{ab}
       +\sum_{a,b}d_{ab}X_{ab}^*=0
 \quad\text{in }{\cal A}_N,                                   \tag{11.4}
\]
the same linear combination of the prescribed outputs, with
\(I\) sent to \(I\), must vanish.  This relation-compatibility
condition is necessary before positivity can be discussed.

\begin{proposition}[Finite-feature reflected factorization]
\label{prop:v13v11-feature}
Suppose (11.3) and the assignment on adjoints satisfy all relations
(11.4), so they define a unital linear map
\(\Gamma_0:{\cal S}_N\to{\cal A}_N\).  If \(\Gamma_0\) is completely
positive and has a unital completely positive extension
\(\Gamma:{\cal A}_N\to{\cal A}_N\), then for every
\(F,G\in\operatorname{span}\{F_1,\ldots,F_q\}\),
\[
 \Gamma(\Theta(F)G)
 =
 \sum_{\alpha=1}^r\Theta(L_\alpha F)L_\alpha G.                \tag{11.5}
\]
Hence \(\omega\circ\Gamma\) is OS positive on the truncated feature
span for every OS-positive state \(\omega\).
\end{proposition}

\begin{proof}
Expand \(F=\sum_af_aF_a\) and \(G=\sum_bg_bF_b\).
Conjugate linearity of \(\Theta\), linearity of \(\Gamma\), and
(11.3) give
\[
\begin{split}
 \Gamma(\Theta(F)G)
 &=
 \sum_{a,b}\overline{f_a}g_b\Gamma_0(X_{ab})\\
 &=
 \sum_\alpha
 \Theta\!\left(\sum_af_aL_\alpha F_a\right)
 \left(\sum_bg_bL_\alpha F_b\right),
\end{split}                                                     \tag{11.6}
\]
which is (11.5).  The OS conclusion is
Theorem~\ref{thm:v13v8-preserve} restricted to this span.
\end{proof}

A numerical table of values for the \(X_{ab}\) is insufficient if
it violates a relation in (11.4).  Such a violation is a compiler
or overlap defect, not an anomaly and not a positivity residual.

\subsection{The finite matrix Choi certificate}

Represent the finite-cutoff algebra as a subalgebra of a matrix
algebra.  The basic full-matrix result is stated first.  For a linear
map
\[
 \Phi:M_d(\mathbb C)\longrightarrow M_s(\mathbb C),             \tag{11.7}
\]
define its Choi matrix
\[
 C_\Phi
 =
 \sum_{a,b=1}^dE_{ab}\otimes\Phi(E_{ab})
 \in M_d\otimes M_s.                                          \tag{11.8}
\]

\begin{theorem}[Choi feasibility for a prescribed operator system]
\label{thm:v13v11-choi}
Let \({\cal S}\subset M_d\) be an operator system, let
\(A_1,\ldots,A_p\) be a linear basis containing \(I_d\), and let
\(B_1,\ldots,B_p\in M_s\) be prescribed with the same star
relations and \(I_d\mapsto I_s\).  A unital completely positive
map
\[
 \Phi:M_d\longrightarrow M_s,\qquad \Phi(A_j)=B_j              \tag{11.9}
\]
exists if and only if there is \(C\in M_d\otimes M_s\) satisfying
\[
 C\succeq0,                                                     \tag{11.10}
\]
\[
 \operatorname{Tr}_1 C=I_s,                                    \tag{11.11}
\]
and
\[
 \operatorname{Tr}_1[(A_j^T\otimes I_s)C]=B_j,
 \qquad j=1,\ldots,p.                                         \tag{11.12}
\]
Here \(\operatorname{Tr}_1\) is the partial trace over \(M_d\).
\end{theorem}

\begin{proof}
If \(\Phi\) is completely positive, Choi's theorem gives
\(C_\Phi\succeq0\).  Writing \(C_\Phi=[C_{ab}]_{a,b=1}^d\), one has
\(C_{ab}=\Phi(E_{ab})\), and direct matrix expansion gives
\[
 \Phi(A)=
 \operatorname{Tr}_1[(A^T\otimes I_s)C_\Phi].                  \tag{11.13}
\]
Taking \(A=I_d\) proves (11.11), and taking \(A=A_j\) proves
(11.12).

Conversely, let \(C\succeq0\) satisfy (11.11)--(11.12), and define
\(\Phi_C\) by the right side of (11.13).  Choi's theorem makes
\(\Phi_C\) completely positive.  Equation (11.11) gives
\(\Phi_C(I_d)=I_s\), and (11.12) gives every prescribed value.
\end{proof}

Equations (11.10)--(11.12) are a finite semidefinite feasibility
problem with exact linear constraints.  A floating positive
eigenvalue is not a certificate unless accompanied by a rational or
interval residual and a proved eigenvalue floor, as in the
center--radius principle audited at V10.

For a general operator system
\({\cal S}\subset{\cal A}\subset{\cal B}({\cal H})\), once the
prescribed map on \({\cal S}\) is proved unital and completely
positive, the Arveson extension theorem gives a completely positive
extension to \({\cal A}\); because the unit already belongs to
\({\cal S}\) and is fixed, the extension is unital.  This invocation
does not repair failure of relation compatibility or complete
positivity on \({\cal S}\).

\subsection{Kraus recovery from the certificate}

If \(C\succeq0\), choose a factorization
\[
 C=\sum_{\alpha=1}^r
 |w_\alpha\rangle\langle w_\alpha|.                            \tag{11.14}
\]
Reshape \(w_\alpha\in\mathbb C^d\otimes\mathbb C^s\) into matrices
\(V_\alpha:\mathbb C^d\to\mathbb C^s\).  Then
\[
 \Phi_C(A)=
 \sum_{\alpha=1}^rV_\alpha A V_\alpha^*,                       \tag{11.15}
\]
up to the fixed transpose convention in (11.13), and
\[
 \sum_\alpha V_\alpha V_\alpha^*=I_s                           \tag{11.16}
\]
is the unitality condition.  Thus the positive Choi certificate
provides an explicit finite Kraus representation.  The reflected
factorization (11.5) remains an additional set of linear equations
on its values at \(X_{ab}\).

\subsection{Relative Choi positivity}

Let \(C_0\succeq0\) be a reference feasible Choi matrix and let
\(\Delta C=\Delta C^*\).  Suppose
\[
 -\varepsilon C_0\preceq\Delta C\preceq\varepsilon C_0,
 \qquad 0\leq\varepsilon<1.                                   \tag{11.17}
\]

\begin{theorem}[Whitened stability of complete positivity]
\label{thm:v13v11-relative}
The perturbed matrix
\[
 C=C_0+\Delta C                                                \tag{11.18}
\]
satisfies
\[
 C\succeq(1-\varepsilon)C_0\succeq0.                           \tag{11.19}
\]
If additionally
\[
 \operatorname{Tr}_1\Delta C=0                                \tag{11.20}
\]
and the required value constraints are perturbed by the linear
amount prescribed through (11.12), then \(C\) is the Choi matrix of
a UCP map with those values.

In finite dimensions, (11.17) is equivalent to
\[
 \ker C_0\subset\ker\Delta C,\qquad
 \|C_0^{+1/2}\Delta C\,C_0^{+1/2}\|_{\rm op}
 \leq\varepsilon,                                             \tag{11.21}
\]
where \(C_0^+\) is the Moore--Penrose inverse.
\end{theorem}

\begin{proof}
The lower inequality in (11.17) gives
\[
 C_0+\Delta C\ \succeq\ (1-\varepsilon)C_0,
\]
which proves (11.19).  Equation (11.20) preserves (11.11), and
linearity of (11.12) gives the stated prescribed values.
Theorem~\ref{thm:v13v11-choi} then supplies the UCP map.

For the equivalence, decompose the space as
\(\operatorname{supp}C_0\oplus\ker C_0\).  The two-sided order in
(11.17), tested on the kernel and then on kernel--support linear
combinations, forces every kernel row and column of \(\Delta C\)
to vanish.  On the support, conjugate (11.17) by
\(C_0^{+1/2}\) to obtain
\[
 -\varepsilon I\preceq
 C_0^{+1/2}\Delta C\,C_0^{+1/2}
 \preceq\varepsilon I,
\]
which is equivalent to the norm bound in (11.21).  Reversing the
conjugation proves the converse.
\end{proof}

An absolute estimate
\(\|\Delta C\|<\lambda_{\min}(C_0)\) works only when \(C_0\) is
strictly positive and spends the smallest global eigenvalue.  The
relative condition (11.17) remains meaningful for a singular Choi
matrix and preserves its forbidden kernel directions exactly.

\subsection{Why scalar positivity is insufficient}

\begin{proposition}[Transpose counterexample]
\label{prop:v13v11-transpose}
The transpose map
\[
 \mathsf t:M_2\longrightarrow M_2,\qquad
 \mathsf t(A)=A^T                                             \tag{11.22}
\]
is positive and unital but is not completely positive.
\end{proposition}

\begin{proof}
Transpose preserves the eigenvalues of a Hermitian matrix, so it is
positive, and it fixes the identity.  Its Choi matrix is
\[
 C_{\mathsf t}
 =
 \sum_{a,b=1}^2E_{ab}\otimes E_{ba},                           \tag{11.23}
\]
the flip operator on \(\mathbb C^2\otimes\mathbb C^2\).
It acts by \(-1\) on the antisymmetric vector
\[
 e_1\otimes e_2-e_2\otimes e_1,
\]
so \(C_{\mathsf t}\) is not positive.  Choi's theorem then shows
that \(\mathsf t\) is not completely positive.
\end{proof}

This example separates one-level or state-specific Gram positivity
from the matrix-level condition needed for a physical channel.

\subsection{Background chart gluing of Choi certificates}

Suppose chart \(i\) has a feasible Choi matrix \(C_i(x)\) with
\[
 C_i(x)\succeq0,\qquad
 \operatorname{Tr}_1C_i(x)=I.                                 \tag{11.24}
\]
For a partition of unity \(\chi_i\), set
\[
 C(x)=\sum_i\chi_i(x)C_i(x).                                  \tag{11.25}
\]
Then (11.24) holds for \(C(x)\), so the glued map is UCP.  This is
the Choi form of the V6 convex gluing theorem.

If
\[
 -\varepsilon_iC_{0,i}\preceq
 C_i-C_{0,i}
 \preceq\varepsilon_iC_{0,i},                                 \tag{11.26}
\]
then with
\[
 C_0=\sum_i\chi_iC_{0,i},\qquad
 \bar\varepsilon=\max_i\varepsilon_i,
\]
summing the Loewner inequalities gives
\[
 -\bar\varepsilon C_0
 \preceq C-C_0
 \preceq\bar\varepsilon C_0.                                  \tag{11.27}
\]
Thus relative complete positivity survives the chart gluing without
an entrywise global maximum.

Unitary chart transitions act on Choi matrices by a fixed unitary
congruence determined by the input and output representations.
Positivity and the relative inequalities are invariant under that
congruence.  Exact overlap equality still requires the prescribed
operator-system values to transform compatibly.

\subsection{Continuum use and remaining extension stocks}

For an increasing family of positive-time features, let
\({\cal S}_N\) exhaust a norm-dense OS product system.  At each
finite stage one must provide:

\begin{enumerate}
\item exact relation compatibility (11.4);
\item a feasible positive Choi matrix satisfying (11.11)--(11.12);
\item the reflected value identities (11.3);
\item compatibility under background-chart and regional
restriction; and
\item a summable approximation error on the growing feature ledger.
\end{enumerate}

The V7 diagonal then selects regulator, region, feature number, and
spectral threshold together.  Exact UCP maps preserve ordinary
positivity; the reflected identities preserve the OS Gram cone on
the finite ledger; vanishing ledger error and OS closedness pass the
property to the dense local core.

\begin{assessment}[Status after thirteenth-cycle V11]
\label{ass:v13v11-status}
The finite-cutoff universal extension problem is now an exact Choi
semidefinite feasibility problem with separate relation, unitality,
and reflected-value constraints.  Relative Loewner control of the
Choi perturbation preserves complete positivity and singular kernel
directions without spending a global minimum eigenvalue.  Scalar
Gram positivity is proved insufficient by the transpose example.
\end{assessment}

\begin{firewall}[Claim boundary after thirteenth-cycle V11]
\label{fire:v13v11-claim}
No RPESM operator system, Choi certificate, reflected value map, or
normal compatible Arveson extension has been constructed.  Arveson
extension does not automatically preserve compact range,
background continuity, regional compatibility, BRST equivariance,
or the compact-operator core.  Anomaly cancellation, determinant
gluing, clustering, and the full SM--Einstein continuum remain
unproved.
\end{firewall}

\subsection{Next bottleneck}

The next task is simultaneous enforcement of UCP, reflected, BRST,
and restriction constraints.  Each is linear except Choi
positivity.  The next note should formulate their intersection as a
finite semidefinite feasibility set, prove compactness after a trace
normalization, and give a quantitative stability theorem under
linear residuals using a strictly feasible reference point on the
allowed support.  If strict feasibility fails because of protected
kernel directions, the relative support condition (11.21) must be
retained.

\subsection{Parent record}

This V11 note appends to the validated thirteenth-cycle V10 file with
record
\[
\begin{array}{c|c}
\text{lines}&3031\\
\text{bytes}&104216\\
\text{SHA--256}&
\texttt{0BB6F257695E9DFDFC319515072FCC058D8AB71EAC9A9915D7362CF1E7E6365A}.
\end{array}
\]

% =============================================================================
% END APPENDED THIRTEENTH-CYCLE V11 MATERIAL
% =============================================================================

% =============================================================================
% BEGIN APPENDED THIRTEENTH-CYCLE V12 MATERIAL
% =============================================================================

\section{The joint Choi spectrahedron and quantitative residual
repair}
\label{sec:v13v12}

V11 treats complete positivity and each prescribed value through one
Choi matrix.  The physical finite-cutoff map must satisfy several
families of linear identities simultaneously: unitality, reflected
values, BRST intertwining, background overlap, and regional
restriction.  Their intersection with the positive cone is a
finite-dimensional spectrahedron.  This note proves its compactness
and a support-sensitive residual-repair theorem.

\subsection{The aggregate linear constraint map}

Let
\[
 {\cal V}_P=
 \{C=C^*\in M_d\otimes M_s:C=PCP\}                             \tag{12.1}
\]
be the real vector space of Hermitian Choi matrices supported by a
fixed projection \(P\).  The projection records protected kernel
directions from the relative condition (11.21).  Let \({\cal W}\)
be a finite-dimensional real Hilbert space and let
\[
 {\mathcal A}:{\cal V}_P\longrightarrow{\cal W}                \tag{12.2}
\]
collect all required linear equations.  Its components may include:

\begin{enumerate}
\item unitality:
\[
 \operatorname{Tr}_1C=I_s;
                                                                    \tag{12.3}
\]
\item the reflected product values (11.12);
\item BRST intertwining
\[
 \Phi_Cs_{\rm in}=s_{\rm out}\Phi_C
\]
on a finite graph core;
\item equality after a declared background transition; and
\item regional channel naturality on a finite local observable
ledger.
\end{enumerate}

After moving fixed inhomogeneous terms to the right, write the full
system as
\[
 {\mathcal A}C=b.                                               \tag{12.4}
\]
All maps are understood on the real Hermitian spaces, so adjoints
and norms below are ordinary finite-dimensional ones.

Define the supported feasibility set
\[
 {\mathfrak F}_{P,b}
 =
 \{C\in{\cal V}_P:C\succeq0,\ {\mathcal A}C=b\}.                \tag{12.5}
\]

\begin{theorem}[Compact joint feasibility set]
\label{thm:v13v12-compact}
Assume the unitality equations (12.3) are among the components of
(12.4).  Then \({\mathfrak F}_{P,b}\) is convex and compact.
Every point of it defines a UCP map satisfying all encoded finite
constraints.
\end{theorem}

\begin{proof}
The positive semidefinite cone, the support space \({\cal V}_P\),
and the affine set (12.4) are closed and convex, so their
intersection is closed and convex.  For a feasible Choi matrix,
\[
 \operatorname{Tr}C
 =
 \operatorname{Tr}(\operatorname{Tr}_1C)
 =
 \operatorname{Tr}I_s=s.                                      \tag{12.6}
\]
For \(C\succeq0\), the trace bounds every Schatten norm and in
particular \(\|C\|_{\rm op}\leq s\).  Thus the feasible set is
bounded.  Closed and bounded subsets of the finite-dimensional
space \({\cal V}_P\) are compact.  The Choi theorem from V11 gives
the final assertion.
\end{proof}

The theorem gives no existence.  It shows that once nonemptiness is
proved, optimizing any continuous error stock over the admissible
channels attains a minimum.

\subsection{A support-sensitive repair theorem}

Let
\[
 {\cal W}_0:=\operatorname{Ran}{\mathcal A}|_{{\cal V}_P}.       \tag{12.7}
\]
Assume a linear right inverse
\[
 {\mathcal Q}:{\cal W}_0\longrightarrow{\cal V}_P,\qquad
 {\mathcal A}{\mathcal Q}=I_{{\cal W}_0}                        \tag{12.8}
\]
has been certified, with
\[
 \|{\mathcal Q}w\|_{\rm op}\leq\gamma\|w\|_{\cal W}.            \tag{12.9}
\]

\begin{theorem}[Residual repair inside the allowed support]
\label{thm:v13v12-repair}
Let \(\widetilde C\in{\cal V}_P\) satisfy
\[
 \widetilde C\succeq\lambda P,\qquad \lambda>0,                \tag{12.10}
\]
and let its linear residual be
\[
 r={\mathcal A}\widetilde C-b\in{\cal W}_0.                    \tag{12.11}
\]
Define
\[
 C=\widetilde C-{\mathcal Q}r.                                 \tag{12.12}
\]
Then
\[
 {\mathcal A}C=b                                               \tag{12.13}
\]
and
\[
 C\succeq(\lambda-\gamma\|r\|_{\cal W})P.                      \tag{12.14}
\]
Consequently \(C\in{\mathfrak F}_{P,b}\) whenever
\[
 \gamma\|r\|_{\cal W}\leq\lambda.                              \tag{12.15}
\]
Strict inequality leaves a certified positive floor on \(P\).
\end{theorem}

\begin{proof}
Equations (12.8), (12.11), and (12.12) give
\[
 {\mathcal A}C
 =
 {\mathcal A}\widetilde C-r=b.
\]
The correction is Hermitian and supported in \(P\).  By (12.9),
\[
 -\gamma\|r\|P
 \preceq-{\mathcal Q}r
 \preceq\gamma\|r\|P                                         \tag{12.16}
\]
as forms on \(P\).  Add the lower inequality to (12.10) and obtain
(12.14).
\end{proof}

This theorem turns floating or interval approximations into exact
feasible Choi matrices when three quantities are certified:
the support floor \(\lambda\), the residual norm, and the
right-inverse norm \(\gamma\).

\subsection{Perturbed target constraints}

Let \(C_0\in{\mathfrak F}_{P,b_0}\) satisfy
\[
 C_0\succeq\lambda_0P.                                        \tag{12.17}
\]
For a perturbed right side \(b_x=b_0+\Delta b_x\), assume
\(\Delta b_x\in{\cal W}_0\) and define
\[
 C_x=C_0+{\mathcal Q}\Delta b_x.                               \tag{12.18}
\]

\begin{corollary}[Strict-feasibility radius]
\label{cor:v13v12-radius}
The matrix \(C_x\) satisfies
\[
 {\mathcal A}C_x=b_x,\qquad
 C_x\succeq
 (\lambda_0-\gamma\|\Delta b_x\|)P.                            \tag{12.19}
\]
Thus every target with
\[
 \|\Delta b_x\|<\lambda_0/\gamma                              \tag{12.20}
\]
has a UCP solution on the same support.
\end{corollary}

\begin{proof}
Apply (12.8)--(12.9) to (12.18), or use
Theorem~\ref{thm:v13v12-repair} with
\(\widetilde C=C_0\) and target \(b_x\).
\end{proof}

If unitality is fixed rather than perturbed, \(\Delta b_x\) has zero
component in (12.3), and the right inverse must preserve that
homogeneous equation.  The same applies to an anomaly-free BRST
right side: the repair proves anomaly freedom only when the zero
target is part of the exact certified system.

\subsection{Stability of the right inverse}

After the support projection and coordinate frames are transported
as in V3--V4, the linear constraint map itself may vary:
\[
 {\mathcal A}_x={\mathcal A}_0+\Delta{\mathcal A}_x.            \tag{12.21}
\]
Let \({\mathcal Q}_0\) be a right inverse of
\({\mathcal A}_0\) on the relevant target space with norm at most
\(\gamma_0\).

\begin{proposition}[Neumann repair of the constraint solver]
\label{prop:v13v12-neumann}
If
\[
 \|\Delta{\mathcal A}_x\|\,\gamma_0<1,                         \tag{12.22}
\]
then
\[
 {\mathcal Q}_x
 =
 {\mathcal Q}_0
 (I+\Delta{\mathcal A}_x{\mathcal Q}_0)^{-1}                  \tag{12.23}
\]
is a right inverse of \({\mathcal A}_x\), and
\[
 \|{\mathcal Q}_x\|
 \leq
 \frac{\gamma_0}
 {1-\|\Delta{\mathcal A}_x\|\gamma_0}.                         \tag{12.24}
\]
\end{proposition}

\begin{proof}
Equation (12.22) makes the inverse in (12.23) converge by the
Neumann series.  Then
\[
 {\mathcal A}_x{\mathcal Q}_0
 =
 {\mathcal A}_0{\mathcal Q}_0
 +\Delta{\mathcal A}_x{\mathcal Q}_0
 =
 I+\Delta{\mathcal A}_x{\mathcal Q}_0,
\]
so multiplication by its inverse proves
\({\mathcal A}_x{\mathcal Q}_x=I\).  The geometric-series norm
bound gives (12.24).
\end{proof}

This is the linear-constraint counterpart of the projection-gap and
norm-resolvent radii in V3--V4.  The three radii must hold on the
same background chart.

\subsection{Protected kernels and failure of ordinary Slater
interiority}

If \(P\neq I\), no feasible matrix supported in \(P\) is strictly
positive in the full Choi space.  The correct Slater condition is
relative:
\[
 C_0\succeq\lambda_0P
 \quad\text{inside }{\cal V}_P.                               \tag{12.25}
\]
Corrections must remain in \({\cal V}_P\).  A residual demanding a
component outside
\(\operatorname{Ran}{\mathcal A}|_{{\cal V}_P}\) cannot be repaired
without enlarging the support.  Such a residual may signal an
incompatible reflected relation, a changed cohomology class, or an
anomalous constraint.

An exact dual obstruction is always sufficient to prove
infeasibility.  If \(y\in{\cal W}\) satisfies
\[
 {\mathcal A}^*y\succeq0\text{ on }{\cal V}_P,\qquad
 \langle y,b\rangle<0,                                        \tag{12.26}
\]
then no \(C\in{\mathfrak F}_{P,b}\) exists.  Indeed, feasibility
would give
\[
 \langle y,b\rangle
 =\langle{\mathcal A}^*y,C\rangle\geq0,                         \tag{12.27}
\]
a contradiction.  No converse is claimed without an appropriate
closed-image or constraint qualification.

\subsection{Several regions and charts}

For finitely many regions and background charts at diagonal stage
\(j\), collect their Choi matrices in the direct sum
\[
 {\bf C}^{(j)}=\bigoplus_{\alpha=1}^{N_j}C_\alpha.              \tag{12.28}
\]
Positivity is blockwise, and all overlap and restriction equations
are linear couplings among the blocks.  Unitality fixes
\[
 \operatorname{Tr}{\bf C}^{(j)}
 =\sum_{\alpha=1}^{N_j}s_\alpha,                               \tag{12.29}
\]
so the joint feasible set is compact by the same proof as
Theorem~\ref{thm:v13v12-compact}.

A nested sequence of nonempty compact feasible sets has a compatible
point when the bonding maps are continuous and every finite set of
constraints is jointly satisfiable.  At the present stage this is a
finite-intersection principle, not a proof that the RPESM constraint
sets are nonempty.  The next construction must provide uniform
finite-stage feasibility and quantify the defect of their bonding
maps.

\begin{assessment}[Status after thirteenth-cycle V12]
\label{ass:v13v12-status}
All finite UCP, reflected, BRST, overlap, and regional equations can
now be treated as one compact Choi spectrahedron.  A strictly
positive point on the protected support and a certified right
inverse repair linear residuals without losing complete positivity.
The right inverse itself is stable on an explicit Neumann radius,
and a dual vector can certify genuine infeasibility.
\end{assessment}

\begin{firewall}[Claim boundary after thirteenth-cycle V12]
\label{fire:v13v12-claim}
No RPESM finite-stage spectrahedron has been shown nonempty, and no
strict supported Choi point or right-inverse norm has been
calculated.  Enforcing a zero BRST right side is not a proof of
physical anomaly cancellation unless feasibility is independently
certified.  Normality, compact range, infinite-ledger consistency,
determinant gluing, clustering, and the full continuum remain
unproved.
\end{firewall}

\subsection{Next bottleneck}

The next step is a projective compactness theorem for these finite
spectrahedra.  One needs approximate bonding maps between successive
feature, region, and regulator levels, a summable correction radius,
and a proof that nonempty finite-stage feasible sets do not escape
the protected support.  This should yield a compatible family of
UCP reflected BRST channels, or else an explicit limiting dual
obstruction.

\subsection{Parent record}

This V12 note appends to the validated thirteenth-cycle V11 file with
record
\[
\begin{array}{c|c}
\text{lines}&3420\\
\text{bytes}&117893\\
\text{SHA--256}&
\texttt{C5BB59CBB8CD6B828F8F2D2E7FB83005D1986989CE5B2467B79F1AC9F763E05A}.
\end{array}
\]

% =============================================================================
% END APPENDED THIRTEENTH-CYCLE V12 MATERIAL
% =============================================================================

% =============================================================================
% BEGIN APPENDED THIRTEENTH-CYCLE V13 MATERIAL
% =============================================================================

\section{Compact inverse limits of the joint channel
spectrahedra}
\label{sec:v13v13}

V12 gives compactness and residual repair at one finite stage.
Repairing each stage independently does not ensure that the repaired
maps agree on older observables.  The correct construction puts all
constraints up to stage \(j\) into one joint spectrahedron and uses
exact truncation as the bonding map.  Compactness then gives a
compatible infinite family.

\subsection{Nested finite ledgers}

Let
\[
 {\cal S}_1\subset{\cal S}_2\subset\cdots                      \tag{13.1}
\]
be finite-dimensional operator systems whose union is norm dense in
the chosen local observable algebra \({\cal A}_{\rm loc}\).
The index may simultaneously encode finitely many regions,
background charts, positive-time features, and BRST graph-core
elements.

At stage \(j\), collect every Choi block needed for the first \(j\)
items into
\[
 {\bf C}^{(j)}
 =
 (C_1^{(j)},\ldots,C_{N_j}^{(j)}).                             \tag{13.2}
\]
Let \({\mathfrak F}_j\) be the set satisfying:

\begin{enumerate}
\item blockwise Choi positivity and unitality;
\item the protected support constraints;
\item every reflected-product and BRST equation listed by stage
\(j\);
\item every background overlap equation listed by stage \(j\); and
\item exact regional and feature-restriction compatibility among
all blocks already present.
\end{enumerate}

All equalities are finite-dimensional and linear.  By the V12
trace-normalization argument,
\[
 {\mathfrak F}_j\text{ is compact and convex}.                 \tag{13.3}
\]
Because the ledger is nested, dropping the blocks and coordinates
introduced after stage \(j\) defines a continuous affine map
\[
 \pi_j:{\mathfrak F}_{j+1}\longrightarrow{\mathfrak F}_j.      \tag{13.4}
\]
The exact compatibility equations ensure that \(\pi_j\) really
takes values in \({\mathfrak F}_j\).

\begin{theorem}[Compact inverse-limit channel theorem]
\label{thm:v13v13-inverse}
If
\[
 {\mathfrak F}_j\neq\varnothing
 \quad\text{for every }j\in\mathbb N,                          \tag{13.5}
\]
then the inverse limit
\[
 \varprojlim({\mathfrak F}_j,\pi_j)
 =
 \{({\bf C}^{(j)})_j:
   {\bf C}^{(j)}=\pi_j({\bf C}^{(j+1)})\text{ for all }j\}
                                                                    \tag{13.6}
\]
is nonempty and compact.
\end{theorem}

\begin{proof}
Consider the compact product
\[
 {\mathfrak P}:=\prod_{j=1}^\infty{\mathfrak F}_j.             \tag{13.7}
\]
For \(n\geq1\), let
\[
 K_n=
 \{({\bf C}^{(j)})_j\in{\mathfrak P}:
 {\bf C}^{(j)}=\pi_j({\bf C}^{(j+1)})
 \text{ for }1\leq j\leq n\}.                                 \tag{13.8}
\]
Continuity of the bonding maps makes \(K_n\) closed.  It is
nonempty: choose any element of
\({\mathfrak F}_{n+1}\), propagate it backward with
\(\pi_n,\ldots,\pi_1\), and choose arbitrary elements in the later
coordinates.  The sets \(K_n\) are nested compact nonempty sets, so
their intersection is nonempty.  That intersection is exactly
(13.6).  It is closed in \({\mathfrak P}\), hence compact.
\end{proof}

Surjectivity of the individual bonding maps is not required for this
countable chain.  What is required is nonemptiness of every joint
finite stage and exact nesting of its constraints.

\subsection{From compatible Choi blocks to a quasi-local channel}

Choose an element of the inverse limit.  It defines compatible UCP
maps
\[
 \Phi_j:{\cal S}_j\longrightarrow{\cal A}_{\rm loc}            \tag{13.9}
\]
on every finite operator system.  For
\(A\in\bigcup_j{\cal S}_j\), set
\[
 \Phi(A)=\Phi_j(A)
 \quad\text{when }A\in{\cal S}_j.                              \tag{13.10}
\]
Compatibility makes this definition independent of \(j\).

\begin{proposition}[Extension from the dense ledger]
\label{prop:v13v13-extension}
The map (13.10) is unital, completely positive, and contractive on
the algebraic union.  It extends uniquely to a UCP map
\[
 \Phi:{\cal A}_{\rm loc}\longrightarrow{\cal A}_{\rm loc}.      \tag{13.11}
\]
Every reflected-product, BRST, overlap, and regional identity which
appears at some finite stage holds on the corresponding dense core.
\end{proposition}

\begin{proof}
Every finite matrix over the algebraic union is contained in
\(M_k({\cal S}_j)\) for some sufficiently large \(j\).  Complete
positivity of \(\Phi_j\) therefore proves complete positivity of
(13.10) at every matrix level.  Unitality gives norm one, hence
contractivity.  The map extends uniquely by continuity to the norm
closure.  A named finite identity is present at all sufficiently
large stages and is preserved by compatibility, so it holds for the
limit map on its declared core.
\end{proof}

If the reflected identities hold on a norm-dense positive-time
core, continuity extends them to its closure.  If the BRST
derivation is unbounded, graph-norm density and the estimates of the
archived graph-domain theorem are still required; norm continuity of
\(\Phi\) alone does not extend an unbounded intertwining identity.

\subsection{Certified nonemptiness from approximate finite data}

At stage \(j\), write the aggregate supported Choi space as
\({\cal V}_{P_j}\), the aggregate constraint map as
\[
 {\mathcal A}_j:{\cal V}_{P_j}\to{\cal W}_j,
 \qquad {\mathcal A}_j{\bf C}=b_j,                             \tag{13.12}
\]
and let \({\mathcal Q}_j\) be a supported right inverse with norm
at most \(\gamma_j\).

\begin{corollary}[Finite-stage residual certificate]
\label{cor:v13v13-finite}
Suppose there is an approximate joint block
\(\widetilde{\bf C}^{(j)}\) satisfying
\[
 \widetilde{\bf C}^{(j)}\succeq\lambda_jP_j,\qquad
 r_j:={\mathcal A}_j\widetilde{\bf C}^{(j)}-b_j
       \in\operatorname{Ran}{\mathcal A}_j,                    \tag{13.13}
\]
and
\[
 \gamma_j\|r_j\|\leq\lambda_j.                                \tag{13.14}
\]
Then \({\mathfrak F}_j\neq\varnothing\).  If
(13.13)--(13.14) are certified for every \(j\), the inverse limit in
Theorem~\ref{thm:v13v13-inverse} is nonempty.
\end{corollary}

\begin{proof}
Apply the V12 residual-repair theorem to the entire joint block:
\[
 {\bf C}^{(j)}
 =
 \widetilde{\bf C}^{(j)}-{\mathcal Q}_jr_j.
\]
It is positive on the protected support and satisfies every
aggregate stage-\(j\) equation exactly, so it belongs to
\({\mathfrak F}_j\).  The inverse-limit theorem applies to the
resulting nonempty compact sets.
\end{proof}

The approximate blocks at different stages need not be mutually
compatible.  Exact compatibility is recovered by compact selection
from the joint feasible sets, not by declaring the independently
repaired matrices equal.

\subsection{Exact lifting as an alternative constructive route}

Suppose one wants to extend a chosen
\({\bf C}^{(j)}\in{\mathfrak F}_j\).  At stage \(j+1\), augment the
constraint map by the truncation equation
\[
 \pi_j{\bf C}^{(j+1)}={\bf C}^{(j)}.                            \tag{13.15}
\]
If an approximate lift has a supported positive floor and its
residual is smaller than the augmented right-inverse radius, V12
repairs it to an exact lift.  Iterating gives a constructive thread
through the inverse system.

This lifting condition is stronger than (13.5).  It proves
surjectivity over the selected thread and is useful for algorithms,
but the compact existence proof needs only finite-stage
nonemptiness.

\subsection{Finite detection of an obstruction}

The contrapositive of
Theorem~\ref{thm:v13v13-inverse} is useful:
if no compatible infinite family exists while all bonding maps are
well defined as in (13.4), then
\[
 {\mathfrak F}_{j_0}=\varnothing
 \quad\text{for some finite }j_0.                              \tag{13.16}
\]
Thus compact trace-normalized Choi systems have no obstruction that
appears only after every finite ledger has passed.

A sufficient exact certificate of (13.16) is a V12 dual vector
\(y_{j_0}\) satisfying
\[
 {\mathcal A}_{j_0}^*y_{j_0}\succeq0
 \text{ on the supported cone},\qquad
 \langle y_{j_0},b_{j_0}\rangle<0.                             \tag{13.17}
\]
Such a witness locates the first incompatible combination of
unital, reflected, BRST, overlap, or restriction equations.  It
must not be renamed a physical anomaly until its component in the
BRST cohomology ledger has been identified.

\subsection{Local normality and infinite volume}

The extension (13.11) is a UCP map on a quasi-local C-star algebra.
At every finite matrix stage it is normal automatically.  Local
normality on an infinite-dimensional regional von Neumann algebra
requires the predual compactness and energy no-escape argument
summarized in compact V1.  The inverse-limit theorem does not give a
globally normal map or state in a preferred infinite-volume
representation.

For the state construction, compose any OS-positive, anomaly-free
finite-stage state with the compatible channel and use the V7
diagonal.  The channel identities then survive on the growing
ledger.  Existence of the underlying state and of the nonempty
spectrahedra remain separate inputs.

\begin{assessment}[Status after thirteenth-cycle V13]
\label{ass:v13v13-status}
Nonempty compact joint spectrahedra at every finite ledger stage
produce a compatible infinite UCP channel by an exact inverse-limit
argument.  The channel inherits every finite reflected, BRST,
overlap, and regional equation on the dense union.  Approximate
finite data certify stage nonemptiness through the supported
residual radius, and any failure of the compact inverse limit is
already visible at a finite stage.
\end{assessment}

\begin{firewall}[Claim boundary after thirteenth-cycle V13]
\label{fire:v13v13-claim}
No RPESM joint spectrahedron has been certified nonempty at even one
complete physical stage, and no family of residual and
right-inverse bounds has been produced.  The quasi-local extension
is not automatically normal, finite rank, or compatible with an
unbounded BRST domain beyond the certified graph core.  The theorem
does not construct the physical state, cancel anomalies, prove
clustering, or complete the SM--Einstein continuum.
\end{firewall}

\subsection{Next bottleneck}

Pointwise nonempty spectrahedra over the Einstein-background core do
not automatically admit continuous choices.  The next note should
treat the feasible sets as a lower-semicontinuous convex-valued
bundle, apply a continuous selection theorem under its exact
hypotheses, and give an explicit local selection from the V12
strict-feasibility radius.  Holonomy should remain in the transition
maps rather than be removed by choosing discontinuous Choi frames.

\subsection{Parent record}

This V13 note appends to the validated thirteenth-cycle V12 file with
record
\[
\begin{array}{c|c}
\text{lines}&3777\\
\text{bytes}&129531\\
\text{SHA--256}&
\texttt{E81E1A1FC2C1D6F3FB36F16489F6F0A90BA7A14050F32B8FFD1DD64CC7ED47C5}.
\end{array}
\]

% =============================================================================
% END APPENDED THIRTEENTH-CYCLE V13 MATERIAL
% =============================================================================

% =============================================================================
% BEGIN APPENDED THIRTEENTH-CYCLE V14 MATERIAL
% =============================================================================

\section{Continuous background selection from the feasible Choi
bundle}
\label{sec:v13v14}

V13 gives a compatible channel once every finite joint
spectrahedron is nonempty.  A classical--quantum continuum also
requires the Choi blocks to depend measurably, preferably
continuously, on the Einstein background.  Pointwise nonemptiness
alone does not select such a field.  This note derives local
continuous selections from supported strict feasibility and glues
them by convexity.

\subsection{The transported feasible bundle}

Let \(X\) be a compact metric background core.  In one local
trivialization, let \(P_x\) be a norm-continuous family of support
projections in a fixed finite-dimensional Hermitian matrix space
\({\cal V}\).  The direct rotations of V3 identify \(P_x{\cal V}P_x\)
with a fixed real vector space
\[
 {\cal V}_{P_0}=\{C=C^*:C=P_0CP_0\}.                            \tag{14.1}
\]
After that transport, write the finite-stage constraints as
\[
 {\mathcal A}_xC=b_x,                                         \tag{14.2}
\]
where \(x\mapsto{\mathcal A}_x\) and \(x\mapsto b_x\) are
continuous.  The fibre is
\[
 {\mathfrak F}_x=
 \{C\in{\cal V}_{P_0}:C\succeq0,\ {\mathcal A}_xC=b_x\}.        \tag{14.3}
\]
All unital, reflected, BRST, overlap, and regional equations are
included in (14.2).

\subsection{Explicit local selection}

Fix \(x_0\in X\).  Suppose there is a supported strictly feasible
point
\[
 C_0\succeq\lambda_0P_0,\qquad
 {\mathcal A}_{x_0}C_0=b_{x_0},\qquad \lambda_0>0.              \tag{14.4}
\]
Let \({\mathcal Q}_0\) be a right inverse of
\({\mathcal A}_{x_0}\) on the relevant target space, with
\[
 \|{\mathcal Q}_0\|\leq\gamma_0.                               \tag{14.5}
\]
When
\[
 \|({\mathcal A}_x-{\mathcal A}_{x_0})
          {\mathcal Q}_0\|<1,                                 \tag{14.6}
\]
define the V12 Neumann right inverse
\[
 {\mathcal Q}_x=
 {\mathcal Q}_0
 \{I+({\mathcal A}_x-{\mathcal A}_{x_0})
                 {\mathcal Q}_0\}^{-1},                        \tag{14.7}
\]
and set
\[
 C_{x_0}(x)
 =
 C_0+
 {\mathcal Q}_x\{b_x-{\mathcal A}_xC_0\}.                      \tag{14.8}
\]

\begin{theorem}[Local continuous feasible section]
\label{thm:v13v14-local}
On a sufficiently small neighborhood \(U_{x_0}\), the field
\(x\mapsto C_{x_0}(x)\) is continuous,
\[
 {\mathcal A}_xC_{x_0}(x)=b_x,                                 \tag{14.9}
\]
and
\[
 C_{x_0}(x)\succeq
 \{\lambda_0-\kappa_{x_0}(x)\}P_0,                             \tag{14.10}
\]
where
\[
 \kappa_{x_0}(x):=
 \|{\mathcal Q}_x\|\,
 \|b_x-{\mathcal A}_xC_0\|.                                   \tag{14.11}
\]
In particular, it is a feasible continuous section wherever
\(\kappa_{x_0}(x)<\lambda_0\).
\end{theorem}

\begin{proof}
Equation (14.6) makes (14.7) well defined and continuous by the
Neumann series.  As in V12,
\[
 {\mathcal A}_x{\mathcal Q}_x=I.                               \tag{14.12}
\]
Apply \({\mathcal A}_x\) to (14.8) to obtain (14.9).
The correction in (14.8) is Hermitian and supported in \(P_0\), and
its operator norm is at most (14.11).  Adding it to the floor
(14.4) gives (14.10).  Continuity and vanishing of the correction
at \(x_0\) provide a neighborhood on which
\(\kappa_{x_0}<\lambda_0\).
\end{proof}

After undoing the direct rotation, (14.8) is a section in the
original support \(P_x\).  The construction changes no protected
kernel direction.

\subsection{A quantitative chart radius}

Suppose on a metric ball of radius \(r\) about \(x_0\),
\[
 \|{\mathcal A}_x-{\mathcal A}_{x_0}\|
 \leq L_A d(x,x_0),\qquad
 \|b_x-b_{x_0}\|
 \leq L_b d(x,x_0).                                           \tag{14.13}
\]
If
\[
 L_A\gamma_0r<1,                                               \tag{14.14}
\]
then
\[
 \|{\mathcal Q}_x\|
 \leq\frac{\gamma_0}{1-L_A\gamma_0r}.                          \tag{14.15}
\]
Moreover,
\[
 \|b_x-{\mathcal A}_xC_0\|
 \leq
 (L_b+L_A\|C_0\|)d(x,x_0).                                    \tag{14.16}
\]

\begin{corollary}[Certified feasible radius]
\label{cor:v13v14-radius}
Every point in the radius-\(r\) ball has the continuous feasible
selection (14.8) when
\[
 \frac{\gamma_0(L_b+L_A\|C_0\|)r}
      {1-L_A\gamma_0r}
 <\lambda_0.                                                   \tag{14.17}
\]
The resulting supported floor is at least the difference between
the two sides of (14.17).
\end{corollary}

\begin{proof}
The Neumann estimate gives (14.15).  Subtract
\(b_{x_0}={\mathcal A}_{x_0}C_0\) to obtain (14.16).
Equations (14.11), (14.15), and (14.16) then give a correction norm
strictly below \(\lambda_0\), so
Theorem~\ref{thm:v13v14-local} applies.
\end{proof}

This radius must be intersected with the projection-gap and
norm-resolvent chart radii from V3--V4.  All constraints need one
common background neighborhood.

\subsection{Global gluing inside the feasible fibres}

Use compactness of \(X\) to choose finitely many local sections
\[
 C_i(x)\in{\mathfrak F}_x,\qquad x\in U_i,
 \quad i=1,\ldots,N,                                          \tag{14.18}
\]
and a continuous partition of unity \(\{\chi_i\}\) subordinate to
that cover.  Define
\[
 C(x)=\sum_{i=1}^N\chi_i(x)C_i(x).                             \tag{14.19}
\]

\begin{theorem}[Global continuous Choi selection]
\label{thm:v13v14-global}
The field (14.19) is continuous and belongs to
\({\mathfrak F}_x\) for every \(x\in X\).  Hence it defines a
continuous background family of UCP maps satisfying every finite
constraint in (14.2).
\end{theorem}

\begin{proof}
Subordination makes each product \(\chi_iC_i\) extend continuously
by zero outside \(U_i\).  A finite sum is continuous.  Positivity is
preserved by the convex combination, and linearity gives
\[
 {\mathcal A}_xC(x)
 =
 \sum_i\chi_i(x){\mathcal A}_xC_i(x)
 =
 \sum_i\chi_i(x)b_x=b_x.                                      \tag{14.20}
\]
Thus \(C(x)\in{\mathfrak F}_x\).  The Choi theorem gives the UCP
family.
\end{proof}

This is also a constructive verification of the lower
semicontinuity needed by the Michael convex selection theorem.
The explicit proof has the advantage of retaining the quantitative
floor and chart radius.

If the local floors are
\(C_i(x)\succeq\lambda_i(x)P_0\), then
\[
 C(x)\succeq
 \left(\sum_i\chi_i(x)\lambda_i(x)\right)P_0.                  \tag{14.21}
\]
A positive uniform floor follows from a finite cover whose local
floors have positive lower bounds.

\subsection{Compact-group and reflection covariance}

Let a compact group \(G\) act continuously on \(X\) and by
unitary congruences
\[
 \rho_{g,x}:{\cal V}_{P_x}\longrightarrow
 {\cal V}_{P_{gx}}.                                           \tag{14.22}
\]
Assume the feasible bundle is equivariant:
\[
 \rho_{g,x}({\mathfrak F}_x)={\mathfrak F}_{gx}.               \tag{14.23}
\]
For any continuous selection \(C\), define
\[
 C^G(x)=
 \int_G\rho_{g^{-1},gx}C(gx)\,dg.                              \tag{14.24}
\]

\begin{proposition}[Equivariant selection by averaging]
\label{prop:v13v14-equivariant}
The field \(C^G\) is a continuous feasible section and satisfies
\[
 C^G(gx)=\rho_{g,x}C^G(x).                                    \tag{14.25}
\]
For a reflection involution, the Haar integral is the average of
the two reflected sections.
\end{proposition}

\begin{proof}
Every integrand in (14.24) lies in the convex set
\({\mathfrak F}_x\) by (14.23).  Its Bochner integral therefore
lies in the same closed convex set.  Compactness of \(G\) and
continuity of the action give continuity in \(x\).  Haar invariance
and a change of integration variable give (14.25).
\end{proof}

There is no corresponding normalized average over the full
noncompact diffeomorphism group.  Diffeomorphism covariance must be
built into the quotient charts and linear constraints or controlled
as a separate defect.

\subsection{Holonomy is not a selection obstruction for Choi
operators}

The local support bundle may have nontrivial unitary holonomy.
Because the feasible fibres are convex sets of positive operators
and transitions act by conjugation, Theorem
\ref{thm:v13v14-global} can still produce a global Choi section.
This does not choose a phase of a determinant-line vector and does
not trivialize that line.  Scalar phase holonomy disappears in Choi
conjugation, while determinant phases remain part of the amplitude
construction summarized in compact V1.

Likewise, averaging over a harmonic-frame unitary can produce the
maximally mixed support density without selecting a global harmonic
basis.  It does not prove that a retained anomaly class vanishes.

\subsection{Compatibility with the stage inverse limit}

At each ledger stage \(j\), apply
Theorem~\ref{thm:v13v14-global} to obtain continuous feasible
fields.  If the joint spectrahedra include exact truncation
compatibility, the inverse-limit construction of V13 may be applied
in the compact product of continuous-section spaces, provided those
sets are compact in the chosen topology.

Uniform equicontinuity is needed for compactness in the supremum
norm.  Without it, work pointwise in the product topology and use
the resulting measurable or pointwise continuous finite-stage
restrictions.  The next note must decide which topology is supplied
by the quantitative chart constants rather than assuming
Arzel{\`a}--Ascoli.

\begin{assessment}[Status after thirteenth-cycle V14]
\label{ass:v13v14-status}
Supported strict feasibility and a stable constraint right inverse
produce an explicit local continuous Choi section with the certified
radius (14.17).  Finitely many such sections glue by convexity to a
global continuous background family satisfying all finite UCP,
reflected, BRST, overlap, and regional equations.  Compact-group and
reflection covariance can be restored by averaging an equivariant
feasible bundle.
\end{assessment}

\begin{firewall}[Claim boundary after thirteenth-cycle V14]
\label{fire:v13v14-claim}
No RPESM feasible Choi bundle, strict supported point, continuous
support projection, or uniform chart radius has been constructed.
Compact-group averaging does not handle the diffeomorphism group.
A continuous Choi section does not trivialize the chiral determinant
line, prove local normality at infinite dimension, cancel anomalies,
establish clustering, or complete the SM--Einstein continuum.
\end{firewall}

\subsection{Next bottleneck}

The next required note is the V15 companion audit.  After that audit,
the substantive step should establish an equicontinuity estimate for
the selected Choi fields from the Lipschitz data in (14.13), include
the variation of the right inverse and support transport, and apply
Arzel{\`a}--Ascoli to the finite-stage section spaces.  This would
permit the background-continuous selection and the ledger inverse
limit to be taken in one topology.

\subsection{Parent record}

This V14 note appends to the validated thirteenth-cycle V13 file
with record
\[
\begin{array}{c|c}
\text{lines}&4077\\
\text{bytes}&140741\\
\text{SHA--256}&
\texttt{1065C23E352EAAEA529453C170B5599CF6433FB7A3AAA55BB77C7532F879A53A}.
\end{array}
\]

% =============================================================================
% END APPENDED THIRTEENTH-CYCLE V14 MATERIAL
% =============================================================================

% =============================================================================
% BEGIN APPENDED THIRTEENTH-CYCLE V15 MATERIAL
% =============================================================================

\section{V15 audit: equicontinuity should follow the actual
constraint channel}
\label{sec:v13v15}

This is the mandatory companion audit following V10.  The newest
files examined were
\[
\begin{array}{c|c|c|c}
\text{programme}&\text{version}&\text{bytes}&\text{SHA--256}\\ \hline
\mathrm{YM}&70&859150&
\texttt{854AC97F1FEBB46C032D14DD1452D50826DBD26E41D9E5482FBC52CF626DBB33}\\
\mathrm{NS}&26&278283&
\texttt{82C887F8C2C69E4F28FAD7C3DC8DE5B9099CB5A4BF431EBA1FFF3E91935978AD}.
\end{array}
\]
The Navier--Stokes note is unchanged.

\subsection{New Yang--Mills results V69--V70}

Yang--Mills V69 gives a rational global certificate for the
energy-whitened covariance recursion introduced at V68.  It improves
the preceding rigorous response majorant by a factor exceeding
\(1.5\,10^4\), while the remaining bound is still about fourteen
orders above the observed response scale.  Its identified loss is
now the exponentwise triangle inequality inside the differentiated
Gram factor.

V70 then isolates the actual first-derivative channel.  With
\[
 H=A^*A,\qquad
 X=AH^{-1/2},\qquad
 Y_\gamma=A_\gamma H^{-1/2},\qquad
 C_\gamma=X^*Y_\gamma,                                       \tag{15.1}
\]
it proves
\[
 H^{-1/2}H_\gamma H^{-1/2}
 =C_\gamma+C_\gamma^*.                                       \tag{15.2}
\]
The component of \(Y_\gamma\) orthogonal to
\(\operatorname{Ran}X\) makes no first-order contribution.  V70
also gives square-root-free block and Loewner gates for certifying
the compressed norm.  The perpendicular channel reappears at second
order, so it is not discarded universally.

\subsection{Transfer to the Choi selection problem}

For the local selection in V14, the exact quantity corrected by the
right inverse is
\[
 r_x=b_x-{\mathcal A}_xC_0.                                   \tag{15.3}
\]
Bounding \(b_x-b_{x_0}\) and
\(({\mathcal A}_x-{\mathcal A}_{x_0})C_0\) separately is valid but
can lose a cancellation enforced by the actual constraint channel.
The first equicontinuity stock should therefore be
\[
 L_r^{\rm act}:=
 \sup_{x\ne y}
 \frac{\|
 (b_x-b_y)-({\mathcal A}_x-{\mathcal A}_y)C_0
 \|}{d(x,y)},                                                 \tag{15.4}
\]
with the separated estimate used only as a fallback.

Variation of the right inverse, support rotation, and partition of
unity remain separate terms.  At second order, derivatives of
\({\mathcal A}_x\), the right inverse, and the support transport all
re-enter, exactly as the perpendicular Gram channel re-enters the
second Yang--Mills variation.

\begin{assessment}[Decision after the V15 audit]
\label{ass:v13v15-decision}
V16 will derive a Lipschitz bound for the explicit V14 local
selection using the actual residual modulus (15.4), the Neumann
right-inverse identity, and the support-transport modulus.  It will
then glue finitely many local selections with a Lipschitz partition
and apply Arzel{\`a}--Ascoli to compact section spaces.  A compact
inverse limit will produce one background-continuous compatible
channel family under exact finite-stage feasibility.
\end{assessment}

\begin{firewall}[Claim boundary after V15]
\label{fire:v13v15-claim}
The YM V69 remainder remains nonclosing, and V70 supplies
certificate gates rather than their global positivity proof.  No
YM or NS numerical constant is transferred.  No RPESM actual
constraint residual, right-inverse modulus, support transport, or
equicontinuity bound has yet been populated.  The full
SM--Einstein continuum remains open.
\end{firewall}

\subsection{Parent record}

This V15 note appends to the validated thirteenth-cycle V14 file with
record
\[
\begin{array}{c|c}
\text{lines}&4410\\
\text{bytes}&152211\\
\text{SHA--256}&
\texttt{A045582F6E3D8AEDD7234476EF3B5AA49C2718D69684DBBF71849627B32E2AA7}.
\end{array}
\]

% =============================================================================
% END APPENDED THIRTEENTH-CYCLE V15 MATERIAL
% =============================================================================

% =============================================================================
% BEGIN APPENDED THIRTEENTH-CYCLE V16 MATERIAL
% =============================================================================

\section{Equicontinuous feasible Choi fields and the section inverse
limit}
\label{sec:v13v16}

V14 constructs continuous feasible Choi fields but does not yet give
a compact space of such fields in the supremum norm.  This note
derives an explicit Lipschitz bound from the actual constraint
residual, adds the support-transport cost, and applies
Arzel{\`a}--Ascoli at every finite ledger stage.

\subsection{Lipschitz variation of the right inverse}

Work first in one fixed supported Choi space \({\cal V}_{P_0}\).
Let
\[
 {\mathcal A}_x={\mathcal A}_0+\Delta{\mathcal A}_x,\qquad
 {\mathcal Q}_x=
 {\mathcal Q}_0(I+\Delta{\mathcal A}_x{\mathcal Q}_0)^{-1}.     \tag{16.1}
\]
Assume on a chart \(U\)
\[
 \|{\mathcal Q}_0\|\leq\gamma_0,\qquad
 \|\Delta{\mathcal A}_x{\mathcal Q}_0\|\leq\kappa<1,            \tag{16.2}
\]
and
\[
 \|{\mathcal A}_x-{\mathcal A}_y\|
 \leq L_A d(x,y).                                              \tag{16.3}
\]

\begin{lemma}[Right-inverse Lipschitz bound]
\label{lem:v13v16-right}
For \(x,y\in U\),
\[
 \|{\mathcal Q}_x\|
 \leq\frac{\gamma_0}{1-\kappa},                                \tag{16.4}
\]
and
\[
 \|{\mathcal Q}_x-{\mathcal Q}_y\|
 \leq
 \frac{\gamma_0^2L_A}{(1-\kappa)^2}d(x,y).                     \tag{16.5}
\]
\end{lemma}

\begin{proof}
Set
\[
 M_x=I+\Delta{\mathcal A}_x{\mathcal Q}_0.
\]
Equation (16.2) gives
\(\|M_x^{-1}\|\leq(1-\kappa)^{-1}\), proving (16.4).
The inverse identity
\[
 M_x^{-1}-M_y^{-1}
 =
 M_x^{-1}(M_y-M_x)M_y^{-1}                                    \tag{16.6}
\]
and
\[
 \|M_y-M_x\|
 \leq L_A\gamma_0d(x,y)
\]
give (16.5) after left multiplication by
\({\mathcal Q}_0\).
\end{proof}

\subsection{The actual residual channel}

Fix a supported strictly feasible anchor \(C_0\) at the chart
center and define
\[
 r_x=b_x-{\mathcal A}_xC_0,\qquad
 C_x=C_0+{\mathcal Q}_xr_x.                                   \tag{16.7}
\]
Assume the actual residual has the bounds
\[
 \|r_x\|\leq R,\qquad
 \|r_x-r_y\|\leq L_r^{\rm act}d(x,y).                          \tag{16.8}
\]
The separated fallback is
\[
 L_r^{\rm act}
 \leq L_b+L_A\|C_0\|,                                         \tag{16.9}
\]
but cancellation in
\[
 r_x-r_y=
 (b_x-b_y)-({\mathcal A}_x-{\mathcal A}_y)C_0                 \tag{16.10}
\]
may make the actual constant smaller.

\begin{theorem}[Local feasible-section Lipschitz estimate]
\label{thm:v13v16-local}
The selection (16.7) satisfies
\[
 \|C_x-C_y\|
 \leq L_Cd(x,y),                                               \tag{16.11}
\]
where
\[
 L_C=
 \frac{\gamma_0^2L_AR}{(1-\kappa)^2}
 +
 \frac{\gamma_0L_r^{\rm act}}{1-\kappa}.                       \tag{16.12}
\]
If \(r_{x_0}=0\) on a radius-\(r\) chart, one may take
\[
 R\leq L_r^{\rm act}r.                                        \tag{16.13}
\]
\end{theorem}

\begin{proof}
Use (16.7) and insert
\({\mathcal Q}_yr_x\):
\[
 \begin{split}
 \|C_x-C_y\|
 &\leq
 \|{\mathcal Q}_x-{\mathcal Q}_y\|\,\|r_x\|
 +\|{\mathcal Q}_y\|\,\|r_x-r_y\|.
 \end{split}                                                    \tag{16.14}
\]
Apply Lemma~\ref{lem:v13v16-right} and (16.8) to obtain
(16.12).  Equation (16.13) follows from (16.8) with \(y=x_0\).
\end{proof}

Along a differentiable background curve, (16.10) becomes
\[
 \dot r_x=\dot b_x-\dot{\mathcal A}_xC_0.                      \tag{16.15}
\]
Only this compressed constraint channel enters the second term of
(16.12).  Components which cancel in (16.15) are not reintroduced
by separately taking their norms.

\subsection{Support transport}

Let \(W_x\) be the unitary congruence transporting the fixed support
to the physical Choi support at \(x\), and assume
\[
 \|W_x-W_y\|\leq L_Wd(x,y).                                   \tag{16.16}
\]
Set
\[
 \widehat C_x=W_xC_xW_x^*.                                    \tag{16.17}
\]
Unitality gives a fixed trace for every positive Choi matrix, so
\[
 \|C_x\|_{\rm op}\leq M_C                                     \tag{16.18}
\]
with \(M_C\) no larger than the output matrix dimension.

\begin{corollary}[Transported Lipschitz bound]
\label{cor:v13v16-transport}
The physical section satisfies
\[
 \|\widehat C_x-\widehat C_y\|
 \leq
 (L_C+2M_CL_W)d(x,y).                                         \tag{16.19}
\]
\end{corollary}

\begin{proof}
Insert \(W_xC_yW_x^*\):
\[
\begin{split}
 \|\widehat C_x-\widehat C_y\|
 &\leq\|C_x-C_y\|\\
 &\quad+\|(W_x-W_y)C_yW_x^*\|
       +\|W_yC_y(W_x^*-W_y^*)\|.
\end{split}                                                     \tag{16.20}
\]
Use (16.16), (16.18), and
Theorem~\ref{thm:v13v16-local}.
\end{proof}

The support-rotation term is added after the actual residual has been
estimated.  This mirrors the actual-channel distinction in the V15
audit.

\subsection{Lipschitz partition gluing}

A finite open cover of a compact metric space admits a finite
shrinking and a Lipschitz partition of unity
\[
 \chi_i\geq0,\qquad
 \sum_i\chi_i=1,\qquad
 \operatorname{supp}\chi_i\subset U_i.                         \tag{16.21}
\]
One construction uses distance-to-complement functions on the
shrunk cover and divides by their positive sum.

Suppose \(\widehat C_i\) is a feasible local section on a
neighborhood of \(\operatorname{supp}\chi_i\), with
\[
 \|\widehat C_i(x)\|\leq M_i,\qquad
 \operatorname{Lip}(\widehat C_i)\leq L_i,\qquad
 \operatorname{Lip}(\chi_i)\leq\ell_i.                         \tag{16.22}
\]
Extend \(\chi_i\widehat C_i\) by zero and define
\[
 C(x)=\sum_i\chi_i(x)\widehat C_i(x).                           \tag{16.23}
\]

\begin{proposition}[Global equicontinuity]
\label{prop:v13v16-glue}
The field (16.23) is feasible and Lipschitz, with
\[
 \operatorname{Lip}(C)
 \leq
 \sum_i(M_i\ell_i+L_i).                                       \tag{16.24}
\]
\end{proposition}

\begin{proof}
Feasibility is the convexity argument of V14.  If \(x,y\) both lie
in the local domain,
\[
\begin{split}
 \|\chi_i(x)\widehat C_i(x)
       -\chi_i(y)\widehat C_i(y)\|
 &\leq
 |\chi_i(x)-\chi_i(y)|M_i\\
 &\quad+\chi_i(y)
 \|\widehat C_i(x)-\widehat C_i(y)\|\\
 &\leq(M_i\ell_i+L_i)d(x,y).                                  \tag{16.25}
\end{split}
\]
If one point lies outside the local domain, the support condition
and the Lipschitz bound for \(\chi_i\) give the same zero-extension
estimate.  Sum (16.25) over \(i\).
\end{proof}

The factor \(\sum_iL_i\) is a safe bound.  Overlap-compatible local
sections can yield a smaller weighted estimate, but no improvement
is used without an explicit relation.

\subsection{Compact spaces of feasible sections}

At finite ledger stage \(j\), let \({\mathfrak F}_j(x)\) be the
compact Choi spectrahedron over \(x\).  Fix \(L_j<\infty\) and
define
\[
 \Gamma_j=
 \{C\in C(X;{\cal V}_j):
 C(x)\in{\mathfrak F}_j(x)\ \forall x,\quad
 \operatorname{Lip}(C)\leq L_j\}.                              \tag{16.26}
\]

\begin{theorem}[Arzel{\`a}--Ascoli compactness]
\label{thm:v13v16-AA}
If \(\Gamma_j\neq\varnothing\), then \(\Gamma_j\) is compact in the
supremum norm.
\end{theorem}

\begin{proof}
Unital Choi normalization bounds every value uniformly in the
finite-dimensional space \({\cal V}_j\).  The common Lipschitz bound
makes \(\Gamma_j\) equicontinuous.  Arzel{\`a}--Ascoli therefore
makes its closure compact in \(C(X;{\cal V}_j)\).

It remains to show closedness.  A uniform limit of
\(L_j\)-Lipschitz fields is \(L_j\)-Lipschitz.  Positivity, support,
and every linear constraint are closed and depend continuously on
\(x\).  Hence a uniform limit still belongs pointwise to
\({\mathfrak F}_j(x)\).  Thus \(\Gamma_j\) is closed and compact.
\end{proof}

\subsection{Inverse limit in the continuous-section topology}

Assume the exact truncation maps act continuously on sections,
\[
 \Pi_j:\Gamma_{j+1}\longrightarrow\Gamma_j,                    \tag{16.27}
\]
and that the chosen bounds \(L_j\) make these maps well defined.

\begin{corollary}[Background-continuous compatible channel family]
\label{cor:v13v16-inverse}
If every \(\Gamma_j\) is nonempty, then
\[
 \varprojlim(\Gamma_j,\Pi_j)\neq\varnothing.                   \tag{16.28}
\]
Every element defines a background-continuous family of compatible
UCP maps satisfying all finite reflected, BRST, overlap, and
regional identities.
\end{corollary}

\begin{proof}
Each \(\Gamma_j\) is nonempty compact by
Theorem~\ref{thm:v13v16-AA}.  Apply the compact inverse-limit proof
of V13 to the continuous bonding maps (16.27).  Compatibility on
the dense operator-system union gives the UCP family as in V13.
\end{proof}

This resolves the topology left open in V14.  It does not require a
uniform Lipschitz constant across all ledger stages; each fixed
\(\Gamma_j\) may have its own finite \(L_j\).  Exact truncation must
send the stage-\(j+1\) bound into the declared stage-\(j\) bound.

\subsection{Insertion into the state limit}

For each fixed local observable \(F\), background continuity of the
Choi field gives continuity of
\[
 x\longmapsto\Phi_x(F).                                       \tag{16.29}
\]
This places the channel-transformed observable in the
classical--quantum algebra used by the compact background law.
The V7 state diagonal can then pass expectations through the
background and regional limits.  Uniform convergence on finite
operator-system unit balls gives the corresponding finite-ledger
error.

If only measurable coefficient fields are available, the present
Arzel{\`a}--Ascoli argument does not apply.  One may still construct
measurable selections under different hypotheses, but that weaker
route must not be labelled background-continuous.

\begin{assessment}[Status after thirteenth-cycle V16]
\label{ass:v13v16-status}
The explicit feasible Choi selection now has a quantitative
Lipschitz bound based on the actual constraint residual, the
right-inverse variation, and support rotation.  Lipschitz partitions
glue local sections with a controlled modulus.  At each finite
ledger stage the resulting feasible section space is compact in
supremum norm, so a compact inverse limit yields one
background-continuous compatible UCP channel family.
\end{assessment}

\begin{firewall}[Claim boundary after thirteenth-cycle V16]
\label{fire:v13v16-claim}
No RPESM Lipschitz constraint data, actual residual cancellation,
support transport, feasible local anchor, or compatible
stage-by-stage Lipschitz bound has been populated.  The result does
not supply diffeomorphism quotient charts, determinant-line
trivialization, local normality, anomaly cancellation, clustering,
or the full continuum.
\end{firewall}

\subsection{Next bottleneck}

The remaining continuity problem is the infinite-dimensional local
normal extension.  Finite Choi sections give UCP maps on growing
operator systems, but normality on a regional von Neumann algebra
can be lost under pointwise limits.  The next note should impose a
uniform predual compactness or energy-tightness condition on the
channel outputs and prove that the background-continuous UCP limit
has locally normal preadjoints.

\subsection{Parent record}

This V16 note appends to the validated thirteenth-cycle V15 file with
record
\[
\begin{array}{c|c}
\text{lines}&4523\\
\text{bytes}&156412\\
\text{SHA--256}&
\texttt{1FFEA091E91E0717AAFA22013B24B0A5FC90EE222889DC9F868C525D10FDA61F}.
\end{array}
\]

% =============================================================================
% END APPENDED THIRTEENTH-CYCLE V16 MATERIAL
% =============================================================================

% =============================================================================
% APPENDED THIRTEENTH-CYCLE V17 MATERIAL
% =============================================================================

\section{Thirteenth-cycle V17: energy-tight preduals and local normality}
\label{sec:v13v17}

\subsection{Purpose and relation to V16}

V16 constructed background-continuous compatible UCP maps on the
norm-dense union of the finite ledger operator systems.  A bounded
pointwise limit of normal maps need not be normal.  Consequently that
construction alone does not yet give a channel on a regional von
Neumann algebra with a trace-class preadjoint.

This note isolates a sufficient compactness certificate.  Uniform
transfer of a compact-resolvent input energy into a compact-resolvent
output energy makes the Schr\"odinger orbits trace-norm precompact.
A diagonal argument then produces a trace-preserving completely
positive map on the full trace class.  Its Banach adjoint is the
required normal UCP channel.

Throughout this note \({\cal H}_{\rm in}\) and
\({\cal H}_{\rm out}\) are separable Hilbert spaces,
\({\mathfrak T}_{\rm in}={\mathfrak S}_1({\cal H}_{\rm in})\), and
\({\mathfrak T}_{\rm out}={\mathfrak S}_1({\cal H}_{\rm out})\).
The pairing is
\[
 \langle \rho,A\rangle=\operatorname{Tr}(\rho A).                \tag{17.1}
\]
A normal Heisenberg channel
\[
 \Phi_n:{\cal B}({\cal H}_{\rm out})
       \longrightarrow{\cal B}({\cal H}_{\rm in})                \tag{17.2}
\]
has a completely positive trace-preserving preadjoint
\[
 S_n=(\Phi_n)_*:{\mathfrak T}_{\rm in}
       \longrightarrow{\mathfrak T}_{\rm out}.                   \tag{17.3}
\]

\subsection{The output-energy tightness lemma}

Let \(K_{\rm out}\geq0\) be self-adjoint with compact resolvent.  Write
\[
 P_R={\bf 1}_{[0,R]}(K_{\rm out}),\qquad Q_R=I-P_R.               \tag{17.4}
\]
Then \(P_R\) has finite rank and \(P_R\uparrow I\) strongly.

\begin{lemma}[Energy tails are trace-norm tight]
\label{lem:v13v17-tight}
For \(t,E<\infty\), define
\[
 {\cal K}(t,E)=
 \{\sigma\in{\mathfrak T}_{\rm out}:
   \sigma\geq0,\ \operatorname{Tr}\sigma\leq t,\
   \operatorname{Tr}(K_{\rm out}\sigma)\leq E\}.                 \tag{17.5}
\]
The set \({\cal K}(t,E)\) is relatively compact in trace norm.
Moreover, for every \(\sigma\in{\cal K}(t,E)\),
\[
 \operatorname{Tr}(Q_R\sigma)\leq {E\over R},\qquad
 \|\sigma-P_R\sigma P_R\|_1
 \leq2\sqrt{{tE\over R}}.                                       \tag{17.6}
\]
\end{lemma}

\begin{proof}
Functional calculus gives \(K_{\rm out}\geq RQ_R\), hence the first
bound.  The gentle compression estimate for a positive trace-class
operator gives
\[
 \|\sigma-P_R\sigma P_R\|_1
 \leq2\sqrt{\operatorname{Tr}(\sigma)
             \operatorname{Tr}(Q_R\sigma)},                     \tag{17.7}
\]
which proves (17.6).

For fixed \(R\), the compressed family
\(P_R{\cal K}(t,E)P_R\) is bounded in the finite-dimensional space
\(P_R{\mathfrak T}_{\rm out}P_R\), and therefore is relatively
compact.  Equation (17.6) approximates the full family uniformly by
that finite-dimensional family as \(R\to\infty\).  Total boundedness,
and hence relative compactness in the complete trace class, follows.
\end{proof}

The compact-resolvent hypothesis is doing real work.  A bounded
energy functional with infinite-dimensional low-energy spectral
subspaces need not make (17.5) precompact.

\subsection{A dense finite-energy test space}

Let \(K_{\rm in}\geq0\) be self-adjoint with compact resolvent.  Choose
an orthonormal eigenbasis \((e_p)_{p\geq1}\), repeated according to
finite multiplicity, and set
\[
 E_{pq}=|e_p\rangle\langle e_q|.                                \tag{17.8}
\]
The rational complex span
\[
 {\cal D}_{\mathbb Q}
 =
 \operatorname{span}_{\mathbb Q+i\mathbb Q}
 \{E_{pq}:p,q\geq1\}                                            \tag{17.9}
\]
is countable and trace-norm dense in
\({\mathfrak T}_{\rm in}\).

For later compactness, off-diagonal matrix units can be recovered by
polarization from positive rank-one operators.  If
\(u=e_p\) and \(v=e_q\), then
\[
 |u\rangle\langle v|
 =
 {1\over4}\sum_{\ell=0}^{3}i^\ell
 |u+i^\ell v\rangle\langle u+i^\ell v|.                         \tag{17.10}
\]
Every vector in (17.10) has finite \(K_{\rm in}\)-energy.

\subsection{The normal compactness theorem}

Assume constants \(a,b\geq0\) satisfy the uniform energy-transfer
inequality
\[
 \operatorname{Tr}\!\left(K_{\rm out}S_n(\rho)\right)
 \leq
 a\,\operatorname{Tr}(K_{\rm in}\rho)
 +b\,\operatorname{Tr}\rho                                     \tag{17.11}
\]
for every \(n\) and every positive finite-\(K_{\rm in}\)-energy
\(\rho\).

\begin{theorem}[Energy-tight channel compactness]
\label{thm:v13v17-normal}
Under (17.11), the sequence \((S_n)\) has a subsequence
\((S_{n_k})\) and a completely positive trace-preserving contraction
\[
 S:{\mathfrak T}_{\rm in}\longrightarrow{\mathfrak T}_{\rm out} \tag{17.12}
\]
such that
\[
 \|S_{n_k}(\rho)-S(\rho)\|_1\longrightarrow0
 \quad\text{for every }\rho\in{\mathfrak T}_{\rm in}.            \tag{17.13}
\]
Consequently
\[
 \Phi=S^*:{\cal B}({\cal H}_{\rm out})
       \longrightarrow{\cal B}({\cal H}_{\rm in})                \tag{17.14}
\]
is normal and UCP, and
\[
 \operatorname{Tr}\!\left(\rho\,\Phi_{n_k}(A)\right)
 \longrightarrow
 \operatorname{Tr}\!\left(\rho\,\Phi(A)\right)                  \tag{17.15}
\]
for every \(\rho\in{\mathfrak T}_{\rm in}\) and
\(A\in{\cal B}({\cal H}_{\rm out})\).  The convergence in (17.15)
is uniform over \(\|A\|\leq1\) for each fixed \(\rho\).
\end{theorem}

\begin{proof}
First fix a positive rank-one operator
\(\rho=|w\rangle\langle w|\), where \(w\) lies in a finite spectral
subspace of \(K_{\rm in}\).  Trace preservation gives
\[
 \operatorname{Tr}S_n(\rho)=\|w\|^2,                            \tag{17.16}
\]
and (17.11) gives a uniform output-energy bound.  By
Lemma~\ref{lem:v13v17-tight}, the orbit
\(\{S_n(\rho):n\geq1\}\) is relatively compact in trace norm.
Polarization (17.10) shows that the same is true for the orbit of
each \(E_{pq}\), and finite linear combinations give relative
compactness for every element of \({\cal D}_{\mathbb Q}\).

Enumerate \({\cal D}_{\mathbb Q}=\{d_1,d_2,\ldots\}\).
Successive subsequence extraction, followed by the Cantor diagonal
choice, gives a subsequence \((S_{n_k})\) for which
\(S_{n_k}(d_j)\) converges in trace norm for every \(j\).

Each \(S_n\) has Banach norm one on the trace class.  Indeed, its
adjoint is unital and positive, hence has operator norm one.
For arbitrary \(\rho\in{\mathfrak T}_{\rm in}\), choose
\(d_j\in{\cal D}_{\mathbb Q}\) with
\(\|\rho-d_j\|_1\) small.  Then
\[
 \|S_{n_k}\rho-S_{n_\ell}\rho\|_1
 \leq2\|\rho-d_j\|_1
     +\|S_{n_k}d_j-S_{n_\ell}d_j\|_1.                           \tag{17.17}
\]
Thus \(S_{n_k}\rho\) is Cauchy for every \(\rho\).  Its limit defines
the linear contraction \(S\), and (17.13) holds.

Trace preservation passes to the limit because the trace is
trace-norm continuous.  Positivity passes because the positive cone
of the trace class is trace-norm closed.  To prove complete
positivity, fix \(m<\infty\) and a positive
\(Z=[Z_{rs}]\in{\mathfrak S}_1({\cal H}_{\rm in}\otimes\mathbb C^m)\).
Entrywise convergence in (17.13) implies
\[
 (S_{n_k}\otimes{\rm id}_m)(Z)
 \longrightarrow(S\otimes{\rm id}_m)(Z)                         \tag{17.18}
\]
in trace norm: there are only \(m^2\) entries and the trace norm of a
finite block matrix is bounded by the sum of the entry trace norms.
Every term on the left is positive, so the closedness of the
positive cone proves positivity of every amplification of \(S\).

A bounded map on the trace class has a Banach adjoint on
\({\cal B}\), and a map on \({\cal B}\) is normal precisely when it
has a preadjoint.  Complete positivity of \(S\) gives complete
positivity of \(S^*\), while trace preservation gives
\(S^*(I)=I\).  Hence (17.14) is normal UCP.  Finally,
\[
 \left|\operatorname{Tr}\rho
       \bigl(\Phi_{n_k}(A)-\Phi(A)\bigr)\right|
 =
 \left|\operatorname{Tr}
       \bigl(S_{n_k}\rho-S\rho\bigr)A\right|
 \leq\|S_{n_k}\rho-S\rho\|_1\|A\|,                              \tag{17.19}
\]
which proves both assertions about (17.15).
\end{proof}

\subsection{Passage of the energy inequality}

The limit retains the same physical energy control.

\begin{proposition}[Lower-semicontinuous energy transfer]
\label{prop:v13v17-energy}
The limiting map in Theorem~\ref{thm:v13v17-normal} satisfies
\[
 \operatorname{Tr}\!\left(K_{\rm out}S(\rho)\right)
 \leq
 a\,\operatorname{Tr}(K_{\rm in}\rho)
 +b\,\operatorname{Tr}\rho                                     \tag{17.20}
\]
for every positive finite-\(K_{\rm in}\)-energy \(\rho\).
\end{proposition}

\begin{proof}
Let \(K_{\rm out}^{(R)}=\min(K_{\rm out},R)\).  This is bounded and
positive.  Trace-norm convergence gives
\[
 \operatorname{Tr}\!\left(K_{\rm out}^{(R)}S(\rho)\right)
 =
 \lim_{k\to\infty}
 \operatorname{Tr}\!\left(K_{\rm out}^{(R)}S_{n_k}(\rho)\right)
 \leq a\operatorname{Tr}(K_{\rm in}\rho)+b\operatorname{Tr}\rho.
                                                                    \tag{17.21}
\]
Monotone convergence as \(R\to\infty\) proves (17.20).
\end{proof}

\subsection{Compatibility with the finite Choi inverse limit}

Let \({\cal E}_1\subset{\cal E}_2\subset\cdots\) be the increasing
finite-dimensional output operator systems from V13--V16, and
suppose their union \({\cal E}_\infty\) is ultraweakly dense in the
regional von Neumann algebra \({\cal M}_{\rm out}\).  Suppose the
stage-\(n\) normal extensions \(\Phi_n\) agree with a compatible
finite-ledger family on every \({\cal E}_j\) once \(n\geq j\).

\begin{corollary}[Normal realization of a compatible ledger]
\label{cor:v13v17-ledger}
If the preadjoints of the stage extensions obey (17.11), then a
subsequence has a normal UCP limit \(\Phi\).  For every
\(A\in{\cal E}_j\),
\[
 \Phi(A)=\Phi_j(A),                                             \tag{17.22}
\]
where the right side denotes the stable compatible ledger value.
Thus the abstract UCP inverse-limit channel of V13--V16 has a normal
realization on \({\cal M}_{\rm out}\).
\end{corollary}

\begin{proof}
Theorem~\ref{thm:v13v17-normal} produces a normal UCP limit.  If
\(A\in{\cal E}_j\), then \(\Phi_n(A)\) equals the stable ledger value
for all \(n\geq j\).  Equation (17.15), tested against every
trace-class \(\rho\), gives (17.22).  Normal maps agreeing on an
ultraweakly dense operator system agree on its ultraweak closure.
\end{proof}

Ultraweak density is the appropriate hypothesis in the last
sentence.  Norm density in a proper \(C^*\)-subalgebra alone does not
determine a normal map on the ambient von Neumann algebra.

\subsection{Preservation of finite-ledger identities}

Every reflected, BRST, overlap, regional, or covariance constraint
in V11--V16 is a finite family of linear equalities on a fixed
\({\cal E}_j\).  Since those values are eventually constant along
the compatible ledger sequence, they pass exactly to \(\Phi\).
Likewise, complete positivity and unitality pass through the
predual construction rather than through a merely scalar
positive-definiteness test.

If a constraint is available only with residual
\(\varepsilon_j\), it passes to the limit only when its tested
observable lies in stage \(j\) and the relevant residual tends to
zero along the chosen cofinal sequence.  The compactness theorem
does not turn a fixed nonzero residual into an exact identity.

\subsection{A quantitative finite-energy convergence metric}

Choose a trace-norm dense sequence \((\rho_\ell)_{\ell\geq1}\) in
the state space consisting of finite-energy states, and define
\[
 d_{\rm pre}(S,T)
 =
 \sum_{\ell=1}^{\infty}2^{-\ell}
 \min\{1,\|S(\rho_\ell)-T(\rho_\ell)\|_1\}.                      \tag{17.23}
\]
On trace-preserving positive contractions, convergence in
\(d_{\rm pre}\) is equivalent to pointwise trace-norm convergence
on the whole trace class.  Indeed, convergence on the dense states
extends to all states by contractivity, and the complex linear span
of states is the trace class.  Thus the subsequence in
Theorem~\ref{thm:v13v17-normal} converges in the explicit metric
(17.23).  This metric will be used for background-uniform
compactness in the next note.

\begin{assessment}[Status after thirteenth-cycle V17]
\label{ass:v13v17-status}
A uniform compact-resolvent energy-transfer inequality now upgrades
the compatible finite Choi inverse limit to a normal UCP regional
channel.  The proof constructs a CPTP trace-class preadjoint by
energy tightness, diagonal extraction, trace-norm extension, and
matrix-level closedness.  The energy estimate and every stable
finite-ledger identity survive the limit.
\end{assessment}

\begin{firewall}[Claim boundary after thirteenth-cycle V17]
\label{fire:v13v17-claim}
No RPESM Hamiltonians or stress-energy generators have yet been
shown to satisfy (17.11), and the theorem does not manufacture
normal stage extensions.  The result is conditional on a common
Hilbert realization, compact-resolvent energies, ultraweak density
of the ledger union, and one uniform energy-transfer estimate.  It
does not prove background-uniform normality, diffeomorphism descent,
anomaly cancellation, clustering, or the full SM--Einstein
continuum.
\end{firewall}

\subsection{Next bottleneck}

V16 supplied background equicontinuity at each finite Choi stage,
whereas V17 supplied local normal compactness at one fixed
background.  The next step is to combine them.  One needs a uniform
version of (17.11) and equicontinuity of the preadjoints in
\(d_{\rm pre}\), then an Arzel\`a--Ascoli argument with values in the
energy-tight channel space.  This must produce one subsequence valid
for every background, rather than a background-dependent
subsequence.

\subsection{Parent record}

This V17 note appends to the validated thirteenth-cycle V16 file with
record
\[
\begin{array}{c|c}
\text{lines}&4893\\
\text{bytes}&168065\\
\text{SHA--256}&
\texttt{306553B9011AAA643AACA89FC77CC70588D7CAF65572433F9C1E265E5BECFBE8}.
\end{array}
\]

% =============================================================================
% END APPENDED THIRTEENTH-CYCLE V17 MATERIAL
% =============================================================================

% =============================================================================
% APPENDED THIRTEENTH-CYCLE V18 MATERIAL
% =============================================================================

\section{Thirteenth-cycle V18: one normal-channel limit over the background core}
\label{sec:v13v18}

\subsection{The uniformity problem}

V17 gives a normal subsequential limit at each fixed background.
Choosing that subsequence separately at every point does not produce
a channel field: the choices need not agree, need not be measurable,
and need not preserve the finite-ledger compatibility uniformly.
This note proves a background-uniform compactness theorem.

Let \((X,d_X)\) be the compact metric Einstein-background core used
in V14--V16.  Work in one fixed Hilbert trivialization; support
transport between trivializations is accounted for separately as in
V16.  For every \(n\) and \(x\in X\), let
\[
 S_{n,x}:{\mathfrak T}_{\rm in}
       \longrightarrow{\mathfrak T}_{\rm out}                   \tag{18.1}
\]
be CPTP, with normal UCP adjoint \(\Phi_{n,x}=S_{n,x}^*\).

Fix the dense sequence of finite-energy input states
\((\rho_\ell)_{\ell\geq1}\) from V17 and put
\[
 d_{\rm pre}(S,T)
 =
 \sum_{\ell=1}^{\infty}2^{-\ell}
 \min\{1,\|S(\rho_\ell)-T(\rho_\ell)\|_1\}.                      \tag{18.2}
\]

\subsection{Completeness of the predual channel metric}

Let \({\mathfrak C}\) be the set of CPTP contractions from
\({\mathfrak T}_{\rm in}\) to \({\mathfrak T}_{\rm out}\).

\begin{lemma}[Complete metrization of pointwise predual convergence]
\label{lem:v13v18-complete}
The function \(d_{\rm pre}\) is a metric on \({\mathfrak C}\).
A sequence converges in this metric precisely when it converges in
trace norm on every input trace-class operator.  The metric space
\(({\mathfrak C},d_{\rm pre})\) is complete.
\end{lemma}

\begin{proof}
If \(d_{\rm pre}(S,T)=0\), the maps agree on the dense sequence of
states.  Contractivity extends equality to all states, and complex
linearity extends it to the trace class.

Suppose \(d_{\rm pre}(S_j,S)\to0\).  For every fixed \(\ell\), the
\(\ell\)-th summand forces
\(\|S_j\rho_\ell-S\rho_\ell\|_1\to0\).  If \(\rho\) is a state and
\(\|\rho-\rho_\ell\|_1<\varepsilon\), then
\[
 \|S_j\rho-S\rho\|_1
 \leq2\varepsilon+\|S_j\rho_\ell-S\rho_\ell\|_1.                 \tag{18.3}
\]
The decomposition of a trace-class operator into a complex linear
combination of four positive trace-class operators proves pointwise
convergence on all of \({\mathfrak T}_{\rm in}\).  The converse
follows from dominated convergence in the series (18.2).

Now let \((S_j)\) be \(d_{\rm pre}\)-Cauchy.  Each
\((S_j\rho_\ell)\) is trace-norm Cauchy.  The estimate (18.3), with
two sequence indices, defines a pointwise trace-norm limit \(S\) on
the entire trace class.  The proof of
Theorem~\ref{thm:v13v17-normal} shows that trace preservation,
contractivity, and complete positivity pass to this limit.
Therefore \(S\in{\mathfrak C}\), and dominated convergence in
(18.2) gives \(d_{\rm pre}(S_j,S)\to0\).
\end{proof}

For \(a,b\geq0\), denote by \({\mathfrak C}_{a,b}\) the channels
satisfying
\[
 \operatorname{Tr}(K_{\rm out}S\rho)
 \leq a\operatorname{Tr}(K_{\rm in}\rho)+b\operatorname{Tr}\rho \tag{18.4}
\]
on positive finite-energy inputs.  Proposition
\ref{prop:v13v17-energy} shows that
\({\mathfrak C}_{a,b}\) is closed in \(d_{\rm pre}\), hence complete.

\subsection{Pointwise compactness and a common modulus}

Assume, uniformly in \(n\) and \(x\),
\[
 S_{n,x}\in{\mathfrak C}_{a,b}.                                 \tag{18.5}
\]
For each \(\ell\), assume a modulus
\(\omega_\ell:[0,\infty)\to[0,\infty)\), with
\(\omega_\ell(t)\to0\) as \(t\downarrow0\), such that
\[
 \|S_{n,x}(\rho_\ell)-S_{n,y}(\rho_\ell)\|_1
 \leq\omega_\ell(d_X(x,y)).                                    \tag{18.6}
\]
Define
\[
 \Omega(t)=\sum_{\ell=1}^{\infty}2^{-\ell}
              \min\{1,\omega_\ell(t)\}.                         \tag{18.7}
\]
Dominated convergence gives \(\Omega(t)\to0\) as \(t\downarrow0\),
and (18.6) gives
\[
 d_{\rm pre}(S_{n,x},S_{n,y})
 \leq\Omega(d_X(x,y)).                                         \tag{18.8}
\]

\begin{lemma}[Pointwise relative compactness]
\label{lem:v13v18-pointcompact}
For every \(x\in X\), the set
\(\{S_{n,x}:n\geq1\}\) is relatively compact in
\(({\mathfrak C}_{a,b},d_{\rm pre})\).
\end{lemma}

\begin{proof}
Apply Theorem~\ref{thm:v13v17-normal} to an arbitrary sequence from
this set.  It has a subsequence converging in trace norm on every
input.  By Lemma~\ref{lem:v13v18-complete}, this is convergence in
\(d_{\rm pre}\), and Proposition~\ref{prop:v13v17-energy} keeps the
limit in \({\mathfrak C}_{a,b}\).
\end{proof}

\subsection{Background-uniform compactness}

\begin{theorem}[Energy-tight Arzel\`a--Ascoli theorem for channels]
\label{thm:v13v18-AA}
Under (18.5) and (18.6), there are a subsequence
\((n_k)\) independent of \(x\) and a map
\[
 S_\bullet:X\longrightarrow{\mathfrak C}_{a,b}                  \tag{18.9}
\]
such that
\[
 \sup_{x\in X}d_{\rm pre}(S_{n_k,x},S_x)\longrightarrow0.       \tag{18.10}
\]
The limit has the same modulus,
\[
 d_{\rm pre}(S_x,S_y)\leq\Omega(d_X(x,y)).                      \tag{18.11}
\]
For every \(\rho\in{\mathfrak T}_{\rm in}\), the map
\[
 x\longmapsto S_x(\rho)                                        \tag{18.12}
\]
is trace-norm continuous.  Consequently
\(\Phi_x=S_x^*\) is a background-continuous field of normal UCP maps
in the sense that
\[
 x\longmapsto
 \operatorname{Tr}\bigl(\rho\,\Phi_x(A)\bigr)                   \tag{18.13}
\]
is continuous for every trace-class \(\rho\) and every bounded
observable \(A\).
\end{theorem}

\begin{proof}
The maps \(x\mapsto S_{n,x}\), valued in the complete metric space
\({\mathfrak C}_{a,b}\), are equicontinuous by (18.8).  Their values
are pointwise relatively compact by
Lemma~\ref{lem:v13v18-pointcompact}.  The metric-valued
Arzel\`a--Ascoli theorem therefore supplies a uniformly
\(d_{\rm pre}\)-convergent subsequence and a continuous limit
\(S_\bullet\).  Passing to the limit in (18.8) proves (18.11).

We verify the stronger assertion (18.12).  First take
\(\rho=\rho_\ell\).  If \(d_{\rm pre}(S_x,S_y)<2^{-\ell}\), then
the \(\ell\)-th summand in (18.2) is below \(2^{-\ell}\), so
\[
 \|S_x\rho_\ell-S_y\rho_\ell\|_1
 \leq2^\ell d_{\rm pre}(S_x,S_y).                               \tag{18.14}
\]
Thus the orbit of every dense test state is continuous.  For an
arbitrary state \(\rho\), choose \(\rho_\ell\) close in trace norm
and use
\[
 \|S_x\rho-S_y\rho\|_1
 \leq2\|\rho-\rho_\ell\|_1
    +\|S_x\rho_\ell-S_y\rho_\ell\|_1.                           \tag{18.15}
\]
Linearity gives the result for all trace-class inputs.  Finally,
\[
 \left|\operatorname{Tr}\rho
       \bigl(\Phi_x(A)-\Phi_y(A)\bigr)\right|
 =
 \left|\operatorname{Tr}
       \bigl(S_x\rho-S_y\rho\bigr)A\right|
 \leq\|S_x\rho-S_y\rho\|_1\|A\|,                               \tag{18.16}
\]
which proves (18.13).
\end{proof}

\subsection{Uniform pointwise trace-norm convergence}

Uniform convergence in (18.10) also controls any one fixed input
uniformly over the background.

\begin{proposition}[Evaluation of the uniform channel limit]
\label{prop:v13v18-eval}
For every \(\rho\in{\mathfrak T}_{\rm in}\),
\[
 \sup_{x\in X}
 \|S_{n_k,x}(\rho)-S_x(\rho)\|_1\longrightarrow0.               \tag{18.17}
\]
Equivalently, for fixed \(\rho\),
\[
 \sup_{\substack{x\in X\\ \|A\|\leq1}}
 \left|\operatorname{Tr}\rho
       \bigl(\Phi_{n_k,x}(A)-\Phi_x(A)\bigr)\right|
 \longrightarrow0.                                             \tag{18.18}
\]
\end{proposition}

\begin{proof}
For a dense test state \(\rho_\ell\), (18.10) and the
\(\ell\)-th summand of (18.2) give uniform trace-norm convergence.
For an arbitrary state, approximate it by \(\rho_\ell\) and use
contractivity uniformly in \(n\) and \(x\).  Extend by linearity.
Trace duality gives (18.18).
\end{proof}

\subsection{Compatibility with the continuous Choi sections}

Suppose that at finite ledger stage \(j\), the Choi section from V16
prescribes the value of \(\Phi_{n,x}\) on \({\cal E}_j\) for every
\(n\geq j\), and that these prescribed values are continuous in
\(x\).  Then, for \(A\in{\cal E}_j\) and every \(\rho\),
\[
 \operatorname{Tr}\bigl(\rho\,\Phi_x(A)\bigr)
 =
 \lim_{k\to\infty}
 \operatorname{Tr}\bigl(\rho\,\Phi_{n_k,x}(A)\bigr),             \tag{18.19}
\]
uniformly in \(x\) by (18.18).  Hence the normal field recovers the
same finite Choi sections.  All exact reflected, BRST, overlap,
restriction, and covariance equations carried by a fixed stage
pass unchanged to the normal limit.

There is no conflict between the two compactness arguments.  V16
takes an inverse limit inside finite-dimensional spaces of
continuous Choi sections.  The present theorem takes a uniform
limit in the complete predual channel space after normal stage
extensions and the energy certificate have been supplied.

\subsection{How V16 estimates feed (18.6)}

An operator-norm estimate for a finite-dimensional Choi matrix gives
a trace-norm estimate for its Schr\"odinger map on the corresponding
finite input system, with a dimension-dependent norm-equivalence
constant.  Therefore the V16 Lipschitz bound
\[
 L_C+2M_CL_W                                                   \tag{18.20}
\]
produces a candidate \(\omega_\ell(t)=L_\ell t\) whenever the test
state \(\rho_\ell\) is already supported in that finite system.
For a fixed \(\ell\) this is sufficient; no single
dimension-dependent constant must remain bounded as
\(\ell\to\infty\), because (18.7) uses dominated convergence.

What remains indispensable is coherence between the chosen normal
stage extensions.  A Choi estimate on the compressed system does
not control an arbitrary extension outside that compression.  The
normal extensions must be chosen so that (18.6) holds on every
fixed finite-energy test state.

\begin{assessment}[Status after thirteenth-cycle V18]
\label{ass:v13v18-status}
Uniform output-energy transfer and equicontinuity on a countable
dense finite-energy state family produce one subsequence valid over
the entire compact background core.  Its limit is a
background-continuous family of normal UCP maps, converges uniformly
on each fixed trace-class input, preserves the energy inequality,
and realizes all stable continuous Choi-ledger identities.
\end{assessment}

\begin{firewall}[Claim boundary after thirteenth-cycle V18]
\label{fire:v13v18-claim}
The RPESM programme has not yet supplied the uniform energy constants
\(a,b\), the moduli \(\omega_\ell\), or coherent normal stage
extensions.  A fixed Hilbert trivialization and compact background
core were assumed.  The theorem does not yet prove compatibility
over a directed regional net, diffeomorphism descent,
determinant-line triviality, anomaly cancellation, clustering, or
the full SM--Einstein continuum.
\end{firewall}

\subsection{Next bottleneck}

Local normality must now be made compatible with the directed net
of spacetime regions.  Separate subsequences for different regions
would again be inadequate.  The next note should use a countable
cofinal regional exhaustion, impose one compatible energy-tight
metric ledger on all regions, and extract a single
background-continuous locally normal channel whose preadjoints
intertwine every restriction map.

\subsection{Parent record}

This V18 note appends to the validated thirteenth-cycle V17 file with
record
\[
\begin{array}{c|c}
\text{lines}&5265\\
\text{bytes}&182347\\
\text{SHA--256}&
\texttt{6C5339C4A0908302977DD172B087997D45E577ADD9BF6750E50666988290B779}.
\end{array}
\]

% =============================================================================
% END APPENDED THIRTEENTH-CYCLE V18 MATERIAL
% =============================================================================

% =============================================================================
% APPENDED THIRTEENTH-CYCLE V19 MATERIAL
% =============================================================================

\section{Thirteenth-cycle V19: a compatible locally normal regional net}
\label{sec:v13v19}

\subsection{Regional setup}

Let
\[
 {\cal O}_1\Subset{\cal O}_2\Subset\cdots                         \tag{19.1}
\]
be a countable cofinal exhaustion of the spacetime region class
under consideration.  For \(r\leq s\), let
\[
 \alpha^{\rm in}_{rs}:{\cal M}^{\rm in}_r\longrightarrow
                         {\cal M}^{\rm in}_s,\qquad
 \alpha^{\rm out}_{rs}:{\cal M}^{\rm out}_r\longrightarrow
                          {\cal M}^{\rm out}_s                    \tag{19.2}
\]
be normal unital injective \(*\)-homomorphisms satisfying the
direct-system identities.  Their preadjoints are contractions
\[
 (\alpha^{\rm in}_{rs})_*:({\cal M}^{\rm in}_s)_*
       \longrightarrow({\cal M}^{\rm in}_r)_*,\qquad
 (\alpha^{\rm out}_{rs})_*:({\cal M}^{\rm out}_s)_*
       \longrightarrow({\cal M}^{\rm out}_r)_*.                  \tag{19.3}
\]

For every ledger cutoff \(n\), background \(x\in X\), and
\(r\leq n\), suppose there is a normal UCP local channel
\[
 \Phi^{(r)}_{n,x}:{\cal M}^{\rm out}_r
       \longrightarrow{\cal M}^{\rm in}_r                       \tag{19.4}
\]
with preadjoint \(S^{(r)}_{n,x}\).  Assume exact regional
naturality at every finite cutoff:
\[
 \alpha^{\rm in}_{rs}\Phi^{(r)}_{n,x}
 =
 \Phi^{(s)}_{n,x}\alpha^{\rm out}_{rs},
 \qquad r\leq s\leq n.                                         \tag{19.5}
\]
Taking preadjoints gives the equivalent identity
\[
 S^{(r)}_{n,x}(\alpha^{\rm in}_{rs})_*
 =
 (\alpha^{\rm out}_{rs})_*S^{(s)}_{n,x}.                        \tag{19.6}
\]

For clarity, the compactness statements below are written in
standard trace-class representations of the local preduals.  The
same proof works intrinsically for separable von Neumann preduals
once their tight subsets and dense test sequences have been
specified.

\subsection{Regionwise hypotheses}

For every fixed \(r\), assume:

\begin{enumerate}
\item positive compact-resolvent operators
\(K_{{\rm in},r}\) and \(K_{{\rm out},r}\);
\item constants \(a_r,b_r<\infty\), independent of \(n\geq r\) and
\(x\in X\), such that
\[
 \operatorname{Tr}\!\left(
 K_{{\rm out},r}S^{(r)}_{n,x}\rho\right)
 \leq
 a_r\operatorname{Tr}(K_{{\rm in},r}\rho)
 +b_r\operatorname{Tr}\rho;                                    \tag{19.7}
\]
\item a dense sequence of finite-energy states
\((\rho^{(r)}_\ell)_{\ell\geq1}\) and moduli
\(\omega_{r\ell}(t)\to0\) as \(t\downarrow0\), uniformly satisfying
\[
 \|S^{(r)}_{n,x}\rho^{(r)}_\ell
     -S^{(r)}_{n,y}\rho^{(r)}_\ell\|_1
 \leq\omega_{r\ell}(d_X(x,y)).                                  \tag{19.8}
\]
\end{enumerate}

Define the regionwise metric
\[
 d_r(S,T)=
 \sum_{\ell=1}^{\infty}2^{-\ell}
 \min\{1,\|S\rho^{(r)}_\ell-T\rho^{(r)}_\ell\|_1\}.              \tag{19.9}
\]
No uniform bound in \(r\) is assumed.  Each region may have its own
energy constants and equicontinuity moduli.

\subsection{One cofinal subsequence for all regions}

\begin{theorem}[Countable regional diagonal compactness]
\label{thm:v13v19-diagonal}
Under (19.5)--(19.8), there is one increasing sequence
\(n_k\to\infty\) and, for every \(r\), a background-continuous
normal UCP channel
\[
 \Phi_x^{(r)}:{\cal M}^{\rm out}_r
       \longrightarrow{\cal M}^{\rm in}_r                       \tag{19.10}
\]
with preadjoint \(S_x^{(r)}\), such that
\[
 \sup_{x\in X}
 d_r(S^{(r)}_{n_k,x},S_x^{(r)})\longrightarrow0                 \tag{19.11}
\]
for every fixed \(r\).  More strongly, for every fixed
\(\rho\in({\cal M}^{\rm in}_r)_*\),
\[
 \sup_{x\in X}
 \|S^{(r)}_{n_k,x}\rho-S_x^{(r)}\rho\|\longrightarrow0.          \tag{19.12}
\]
The limiting channels satisfy exact regional naturality:
\[
 \alpha^{\rm in}_{rs}\Phi_x^{(r)}
 =
 \Phi_x^{(s)}\alpha^{\rm out}_{rs},
 \qquad r\leq s.                                                \tag{19.13}
\]
\end{theorem}

\begin{proof}
Discard the finitely many terms \(n<r\) whenever region \(r\) is
considered.  Apply Theorem~\ref{thm:v13v18-AA} to region \(1\).
It gives a subsequence with uniform convergence over \(X\).  Apply
the same theorem to region \(2\) inside that subsequence, then to
region \(3\), and so on.  Denote the resulting nested subsequences
by
\[
 N_1\supset N_2\supset\cdots.                                   \tag{19.14}
\]
Choose \(n_k\) to be the \(k\)-th retained index in \(N_k\), enlarged
if necessary so that \(n_k\geq k\).  For fixed \(r\), the tail
\((n_k)_{k\geq r}\) lies in \(N_r\), so it has the uniform
region-\(r\) limit supplied by V18.  This proves (19.10)--(19.12).

It remains to pass naturality.  Fix \(r\leq s\),
\(\eta\in({\cal M}^{\rm in}_s)_*\), and \(x\in X\).  Equation
(19.6) says
\[
 S^{(r)}_{n_k,x}(\alpha^{\rm in}_{rs})_*\eta
 =
 (\alpha^{\rm out}_{rs})_*S^{(s)}_{n_k,x}\eta                  \tag{19.15}
\]
for all sufficiently large \(k\).  Both sides converge in predual
norm by (19.12), and
\((\alpha^{\rm out}_{rs})_*\) is bounded.  Hence
\[
 S_x^{(r)}(\alpha^{\rm in}_{rs})_*\eta
 =
 (\alpha^{\rm out}_{rs})_*S_x^{(s)}\eta.                        \tag{19.16}
\]
Taking adjoints proves (19.13).
\end{proof}

The theorem requires countable cofinality.  An uncountable directed
set without a countable cofinal subsystem cannot be handled by this
sequential diagonal argument; one would need compactness of nets or
an additional separability reduction.

\subsection{The locally normal quasi-local channel}

Let \({\cal A}^{\rm out}_{\rm ql}\) and
\({\cal A}^{\rm in}_{\rm ql}\) be the norm completions of the
algebraic direct limits of the two local nets.

\begin{corollary}[Quasi-local extension]
\label{cor:v13v19-quasilocal}
For every \(x\in X\), the compatible family
\((\Phi_x^{(r)})_r\) defines a unique UCP map
\[
 \Phi_x^{\rm ql}:{\cal A}^{\rm out}_{\rm ql}
       \longrightarrow{\cal A}^{\rm in}_{\rm ql}                \tag{19.17}
\]
whose restriction to every local algebra is \(\Phi_x^{(r)}\) under
the canonical embeddings.  The map is locally normal: every local
restriction has the preadjoint \(S_x^{(r)}\).
\end{corollary}

\begin{proof}
For an algebraic direct-limit element represented by
\(A\in{\cal M}^{\rm out}_r\), set
\[
 \Phi_x^{\rm alg}([A]_r)=[\Phi_x^{(r)}(A)]_r.                   \tag{19.18}
\]
Equation (19.13) makes this independent of the representative.
For any matrix \([A_{ij}]\) contained in one local stage,
complete positivity of \(\Phi_x^{(r)}\) gives positivity of
\([\Phi_x^{\rm alg}(A_{ij})]\).  Unitality is also local.
A unital positive map between \(C^*\)-algebras has norm one, so
\(\Phi_x^{\rm alg}\) extends uniquely by norm continuity to
(19.17).  Complete positivity passes to the norm closure at every
finite matrix level.  Local normality is exactly (19.10).
\end{proof}

This conclusion is local normality on the quasi-local algebra.  It
does not assert that \(\Phi_x^{\rm ql}\) extends normally to every
chosen global von Neumann completion; such an extension depends on
the global representation and requires its own tightness or
folium-preservation statement.

\subsection{Continuity of local expectations}

Let \(A\in{\cal M}^{\rm out}_r\) and
\(\rho\in({\cal M}^{\rm in}_r)_*\).  Theorem
\ref{thm:v13v19-diagonal} and V18 give continuity of
\[
 x\longmapsto
 \langle\rho,\Phi_x^{(r)}(A)\rangle.                            \tag{19.19}
\]
If a quasi-local observable is approximated in norm by a local
observable and the background state field has uniformly bounded
norm, the same conclusion extends by a three-term approximation
argument.  Explicitly, if
\(\|A-A_r\|\leq\varepsilon\), then the two channel approximation
terms are bounded by \(\varepsilon\|\rho\|\), uniformly in \(x\).
Thus the classical--quantum expectation functions needed in the
V7 state diagonal remain continuous on \(X\).

\subsection{Approximate regional compatibility}

Exact identity (19.5) is the clean ledger requirement.  If instead
\[
 \sup_{x\in X}
 \|\alpha^{\rm in}_{rs}\Phi^{(r)}_{n,x}(A)
   -\Phi^{(s)}_{n,x}\alpha^{\rm out}_{rs}(A)\|
 \leq\varepsilon_{n,r,s}\|A\|                                  \tag{19.20}
\]
on a fixed finite operator system and
\(\varepsilon_{n_k,r,s}\to0\), then testing against normal
functionals and using (19.12) passes the corresponding equality to
the limit.  A residual bounded away from zero does not yield
naturality.

For infinitely many pairs \(r\leq s\), the cofinal diagonal must be
chosen so that every fixed residual tends to zero.  Countability
allows this: enumerate the pairs and incorporate their tail
conditions into the same nested extraction.  No uniform convergence
over all pairs is required for the local conclusion.

\subsection{Interaction with the V7 diagonal}

The V7 construction already uses a diagonal over region, spectral
threshold, and ledger depth.  The present extraction can be folded
into that diagonal by requiring, at its \(k\)-th step,

\begin{enumerate}
\item \(n_k\geq k\);
\item the first \(k\) regional predual metrics are within \(2^{-k}\)
of their limiting fields;
\item the first \(k\) restriction residuals are at most \(2^{-k}\);
\item the previously declared overlap, BRST, reflection, background,
and state-approximation errors meet their summable tail bounds.
\end{enumerate}

Every fixed regional observable then eventually lies behind all
relevant cutoffs.  The finite number of its constraints converge,
while summability retains the state-limit Cauchy property.

\begin{assessment}[Status after thirteenth-cycle V19]
\label{ass:v13v19-status}
A countable cofinal regional exhaustion, regionwise energy tightness,
and background equicontinuity now yield one compatible family of
normal local channels.  The family defines a UCP channel on the
quasi-local inductive-limit algebra, is normal on every local
von Neumann algebra, is continuous in background expectations, and
preserves exact or vanishing-residual regional restriction
identities.
\end{assessment}

\begin{firewall}[Claim boundary after thirteenth-cycle V19]
\label{fire:v13v19-claim}
No RPESM regional net, compact-resolvent energy family, energy
transfer estimate, or coherent normal extension has been
constructed.  Countable cofinality and fixed local Hilbert
trivializations were assumed.  Local normality does not imply
normality in an arbitrary global representation.  The argument
does not establish causal factorization, the time-slice axiom,
diffeomorphism descent, anomaly cancellation, clustering, or the
full SM--Einstein continuum.
\end{firewall}

\subsection{Next bottleneck}

The channel is now compatible as a locally normal regional object.
The next mandatory cross-program audit occurs at V20.  It should
check whether the newest Yang--Mills and Navier--Stokes notes supply
a sharper energy-transfer, residual, or gluing mechanism.  After
that audit, the programme should attack equivariant descent of the
locally normal channel through the residual gauge and
diffeomorphism groupoids, keeping determinant-line holonomy as a
separate obstruction.

\subsection{Parent record}

This V19 note appends to the validated thirteenth-cycle V18 file with
record
\[
\begin{array}{c|c}
\text{lines}&5579\\
\text{bytes}&194177\\
\text{SHA--256}&
\texttt{9775868B4D2ECDDA6A6A0978BDB5399B1A3D3D0C35BC1FE0514AE22B18562300}.
\end{array}
\]

% =============================================================================
% END APPENDED THIRTEENTH-CYCLE V19 MATERIAL
% =============================================================================

% =============================================================================
% APPENDED THIRTEENTH-CYCLE V20 MATERIAL
% =============================================================================

\section{Thirteenth-cycle V20: compulsory YM--NS audit}
\label{sec:v13v20}

\subsection{Files inspected}

This is the compulsory five-version cross-program audit.  The newest
companion files present in the shared \texttt{latex\_notes} folder
were inspected read-only:

\[
\begin{array}{c|c|c|l}
\text{programme}&\text{version}&\text{lines}&\text{SHA--256}\\ \hline
\text{Yang--Mills}&73&25930&
\texttt{D99B0C2F4E9B1DA9019253FA89461357}\\
&&&
\texttt{E731047D3BFD4A174922955E443B29C8}\\
\text{Navier--Stokes}&26&5319&
\texttt{82C887F8C2C69E4F28FAD7C3DC8DE5B}\\
&&&
\texttt{9099CB5A4BF431EBA1FFF3E91935978AD}
\end{array}                                                     \tag{20.1}
\]
The Navier--Stokes file and digest are unchanged from the V15 audit.
The Yang--Mills file now contains the substantive V71--V73
continuation beyond the V70 material inspected at V15.

\subsection{Yang--Mills V71: eliminate the wrong sufficient gate}

YM V71 gives exact integer corner witnesses showing that
\[
 A_\mu^*A_\mu\preceq A^*A                                    \tag{20.2}
\]
fails for both the Wilson and endpoint factors.  This does not
disprove the desired direct first-variation inequalities
\[
 -2H\preceq H_\mu\preceq2H.                                  \tag{20.3}
\]
The distinction is structural: (20.2) measures derivative-row
components and anti-Hermitian pieces that cancel or disappear from
the actual Hermitian first variation.

This agrees with the policy used in V16--V19.  Compactness and
continuity must be estimated on the actual predual output or actual
constraint residual.  A stronger factorwise surrogate may fail
although the physical compressed quantity remains controlled.
No numerical YM constant is imported into the SM--Einstein ledger.

\subsection{Yang--Mills V72: exact corner floors}

YM V72 supplies rational congruence certificates for all dyadic
corners, axes, and signs:
\[
 2H_{\rm W}\pm\partial_\mu H_{\rm W}\succeq10^{-3}I,\qquad
 2H_*\pm\partial_\mu H_*\succeq10^{-2}I.                       \tag{20.4}
\]
The actual certified common floors are larger, approximately
\(0.02067257\) and \(0.28513039\).  The proof uses rounded rational
bases, Gershgorin lower bounds, and independent inverse-residual
checks.

These are finite-corner results for the Yang--Mills symbols.  They
do not establish an all-torus inequality and do not provide an
energy-transfer estimate of the form (17.11) for the present
programme.  The transferable lesson is certificate architecture:
first certify a finite anchor by rational congruence, then transport
it using an actual residual bound.  This architecture is already
encoded in V12, V14, and V16.

\subsection{Yang--Mills V73: the moving-frame connection}

YM V73 writes the normalized parity frame as \(V(k)\) with
\[
 \partial_\mu V=V\Gamma_\mu,\qquad \Gamma_\mu^*=-\Gamma_\mu,     \tag{20.5}
\]
and defines
\[
 {\boldsymbol\nabla}_\mu X
 =\partial_\mu X-[X,\Gamma_\mu].                               \tag{20.6}
\]
The connection is constant and flat.  For every closed trace word,
the inserted commutators telescope:
\[
 \partial_\mu\operatorname{Tr}(X_1\cdots X_m)
 =
 \sum_{j=1}^m
 \operatorname{Tr}\!\left(
 X_1\cdots({\boldsymbol\nabla}_\mu X_j)\cdots X_m\right).        \tag{20.7}
\]
Thus frame motion must be removed before estimating the physical
variation.

Equation (20.7) is relevant to the next SM--Einstein step.  A
background-dependent unitary trivialization of a channel produces
the same commutator connection.  Gauge-invariant closed trace
observables do not pay its norm as an analytic defect.  V21 will
prove this identity for the source and target channel frames and
combine it with equivariant averaging.

\subsection{Navier--Stokes V26}

NS V26 computes the pointwise rank-one norms of the first two Oseen
kernel derivatives.  The first-derivative rank-one norm equals the
full symmetric-traceless norm; the second derivative improves it on
an intermediate radial interval, with the largest reported
pointwise reduction near eleven percent.

Those statements concern rank-one inputs to explicit Oseen tensors.
They provide no regional quantum energy, predual tightness,
groupoid action, or determinant-line datum for this programme.
Nothing is imported.  Their methodological content is consistent
with the actual-input principle already adopted: optimize over the
physical input class before applying a larger ambient operator norm.

\begin{decision}[V20 cross-program decision]
\label{dec:v13v20}
The next SM--Einstein note will use the YM V73 moving-frame
calculation only at the structural level.  It will:
\begin{enumerate}
\item formulate residual gauge and diffeomorphism covariance as a
groupoid equivariance equation for the locally normal regional
channel;
\item prove compact-group and proper-groupoid averaging criteria
that preserve normality, complete positivity, regional naturality,
and energy control;
\item show that infinitesimal unitary frame motion contributes
source and target commutators which cancel in invariant closed
trace observables;
\item keep determinant-line holonomy and nonproper
diffeomorphism orbits as explicit unresolved obstructions.
\end{enumerate}
No Yang--Mills or Navier--Stokes physical constant is transferred.
\end{decision}

\begin{assessment}[Status after thirteenth-cycle V20]
\label{ass:v13v20-status}
The required companion audit is complete.  It confirms the use of
actual compressed residuals and identifies one precise tool for the
next descent step: flat moving-frame commutators cancel inside
closed invariant trace words.  The locally normal regional
construction of V17--V19 is unchanged.
\end{assessment}

\begin{firewall}[Claim boundary after thirteenth-cycle V20]
\label{fire:v13v20-claim}
The audit imports no theorem about the RPESM operators and supplies
none of the missing energy constants, background moduli, anomaly
coefficients, or groupoid properness data.  The YM corner floors and
NS rank-one norms belong to their own models.  No continuum
SM--Einstein construction is claimed.
\end{firewall}

\subsection{Parent record}

This V20 note appends to the validated thirteenth-cycle V19 file with
record
\[
\begin{array}{c|c}
\text{lines}&5891\\
\text{bytes}&205807\\
\text{SHA--256}&
\texttt{B347106DAF6E6EC413D2B33E7293654F2A990BA071CC2E8CE39253D9C640E0FF}.
\end{array}
\]

% =============================================================================
% END APPENDED THIRTEENTH-CYCLE V20 MATERIAL
% =============================================================================

% =============================================================================
% APPENDED THIRTEENTH-CYCLE V21 MATERIAL
% =============================================================================

\section{Thirteenth-cycle V21: equivariant descent of the locally normal channel}
\label{sec:v13v21}

\subsection{Covariance as a groupoid equation}

Let \({\cal G}\rightrightarrows X\) be a topological groupoid acting
on the compact background core.  An arrow
\(\gamma:x\to y\) induces normal unital \(*\)-isomorphisms
\[
 \beta^{\rm in}_\gamma:{\cal M}^{\rm in}_x
       \longrightarrow{\cal M}^{\rm in}_y,\qquad
 \beta^{\rm out}_\gamma:{\cal M}^{\rm out}_x
       \longrightarrow{\cal M}^{\rm out}_y.                    \tag{21.1}
\]
They obey
\(\beta_{\eta\gamma}=\beta_\eta\beta_\gamma\) whenever
\(\gamma:x\to y\) and \(\eta:y\to z\).

For a channel field
\(\Phi_x:{\cal M}^{\rm out}_x\to{\cal M}^{\rm in}_x\), define its
transport back along \(\gamma\) by
\[
 (T_\gamma\Phi)_x
 :=
 \beta^{\rm in}_{\gamma^{-1}}\,
 \Phi_y\,\beta^{\rm out}_\gamma.                               \tag{21.2}
\]
Equivariance is the fixed-point equation
\[
 T_\gamma\Phi=\Phi_x,                                          \tag{21.3}
\]
equivalently
\[
 \beta^{\rm in}_\gamma\Phi_x
 =
 \Phi_y\beta^{\rm out}_\gamma.                                \tag{21.4}
\]

If \(S_x=(\Phi_x)_*\), then
\[
 (T_\gamma\Phi)_{x,*}
 =
 (\beta^{\rm out}_\gamma)_*\,
 S_y\,
 (\beta^{\rm in}_{\gamma^{-1}})_*.                             \tag{21.5}
\]
Thus normality is most safely preserved by averaging (21.5) in the
predual.

\subsection{Compact-group averaging}

First suppose a compact group \(G\) acts on \(X\), and write
\(\beta_{g,x}\) for the arrow \(x\to gx\).  Let \(dg\) be normalized
Haar measure.  Assume for every normal functional \(\rho\) that the
predual integrand in (21.5) is trace-norm continuous in \(g\).
Define
\[
 \overline S_x\rho
 =
 \int_G
 (\beta^{\rm out}_{g,x})_*
 S_{gx}
 (\beta^{\rm in}_{g^{-1},gx})_*\rho\,dg,                       \tag{21.6}
\]
as a Bochner integral, and let
\(\overline\Phi_x=\overline S_x^*\).

\begin{theorem}[Compact symmetry descent]
\label{thm:v13v21-compact}
The field \(\overline\Phi\) is normal, UCP, and exactly
\(G\)-equivariant.  If the original field is regionally natural and
the \(G\)-action commutes with the regional embeddings, then
\(\overline\Phi\) is regionally natural.  Every invariant linear
finite-ledger equality is preserved.  More generally, an invariant
closed convex constraint is preserved whenever all transported
channels lie in its feasible set; this covers the Choi positive cone
and convex reflected feasibility conditions.
\end{theorem}

\begin{proof}
Each integrand in (21.6) is a CPTP contraction.  A Bochner integral
of positive trace-preserving maps is positive and trace preserving.
At matrix level the same integral preserves the closed positive
cone, so \(\overline S_x\) is completely positive.  Its adjoint is
therefore normal and UCP.

The transformations \(T_g\) form a right action on channel fields:
\[
 T_gT_h=T_{hg}.                                                \tag{21.7}
\]
This follows directly from the cocycle identity for the two
\(\beta\)'s.  Right invariance of normalized Haar measure gives
\[
 T_h\overline\Phi
 =
 \int_G T_hT_g\Phi\,dg
 =
 \int_G T_{gh}\Phi\,dg
 =
 \overline\Phi.                                                \tag{21.8}
\]
This is (21.3).

If regional embeddings commute with every \(\beta_g\), each
\(T_g\Phi\) satisfies the same regional naturality square as
\(\Phi\).  Integrating that linear equality proves naturality of the
average.  The same argument applies to BRST, overlap, restriction,
and other invariant affine equalities.  A Bochner integral is a
limit of convex combinations, so invariant closed convex feasible
sets are also preserved.
\end{proof}

\begin{proposition}[Averaging displacement]
\label{prop:v13v21-displacement}
For every input normal functional \(\rho\),
\[
 \|\overline S_x\rho-S_x\rho\|
 \leq
 \int_G
 \|(T_g\Phi)_{x,*}\rho-S_x\rho\|\,dg.                           \tag{21.9}
\]
Consequently an already approximately equivariant field is changed
only by its mean actual predual defect.
\end{proposition}

\begin{proof}
Subtract \(S_x\rho\) inside (21.6) and use the triangle inequality
for the Bochner integral.
\end{proof}

No operator norm of the separate source and target transports
appears in (21.9); the quantity is the actual transported channel
defect, in accord with V16 and the V20 audit.

\subsection{Energy transfer under averaging}

Assume the local energy operators transform covariantly under the
unitary implementations of \(\beta\):
\[
 K_{{\rm in},gx}=U^{\rm in}_{g,x}K_{{\rm in},x}
                    (U^{\rm in}_{g,x})^*,\qquad
 K_{{\rm out},gx}=U^{\rm out}_{g,x}K_{{\rm out},x}
                    (U^{\rm out}_{g,x})^*.                     \tag{21.10}
\]
Suppose every \(S_x\) obeys the uniform transfer inequality with
constants \(a,b\).

\begin{proposition}[Energy-tight equivariant average]
\label{prop:v13v21-energy}
The averaged preadjoint obeys the same inequality:
\[
 \operatorname{Tr}\!\left(
 K_{{\rm out},x}\overline S_x\rho\right)
 \leq
 a\operatorname{Tr}(K_{{\rm in},x}\rho)+b\operatorname{Tr}\rho \tag{21.11}
\]
for every positive finite-energy \(\rho\).
\end{proposition}

\begin{proof}
For each integrand of (21.6), unitary covariance (21.10) transports
the output energy at \(gx\) back to the output energy at \(x\), and
the input energy at \(x\) forward to that at \(gx\).  Apply the
original energy-transfer inequality before transporting back.  Its
right side is independent of \(g\).  Integrate the resulting scalar
inequality.  For unbounded energies, first use the bounded spectral
truncations and then monotone convergence, as in
Proposition~\ref{prop:v13v17-energy}.
\end{proof}

Thus the compact-group projection can be applied before the
energy-tight extraction of V17--V19 or after it.  In either order,
the same constants remain valid.

\subsection{Proper-groupoid averaging criterion}

Compactness of the entire gauge-diffeomorphism group is usually
unavailable.  The precise ingredient used by the preceding proof is
an invariant family of probability measures, not compactness by
itself.

For \(x\in X\), write
\[
 {\cal G}_x=\{\gamma\in{\cal G}:s(\gamma)=x\}.                  \tag{21.12}
\]
Assume there are probability measures \(\mu_x\) on
\({\cal G}_x\), continuously depending on \(x\), with the right
invariance property
\[
 (R_\eta)_*\mu_y=\mu_x,\qquad
 R_\eta(\zeta)=\zeta\eta                                      \tag{21.13}
\]
for every arrow \(\eta:x\to y\).  Assume also that the predual
integrands are Bochner integrable.  Such a probability system is
obtained for a proper locally compact groupoid with a Haar system
and a normalized cutoff function, provided the relevant continuity
and support conditions hold.

Define
\[
 \overline S_x\rho
 =
 \int_{{\cal G}_x}(T_\gamma\Phi)_{x,*}\rho\,d\mu_x(\gamma).
                                                                    \tag{21.14}
\]

\begin{theorem}[Invariant-probability groupoid descent]
\label{thm:v13v21-groupoid}
Under (21.13), the adjoint field
\(\overline\Phi_x=\overline S_x^*\) is normal UCP and satisfies the
exact groupoid equivariance equation (21.4).  Regional naturality,
invariant finite-ledger identities, and a covariant uniform energy
bound pass to the average exactly as in
Theorem~\ref{thm:v13v21-compact} and
Proposition~\ref{prop:v13v21-energy}.
\end{theorem}

\begin{proof}
Complete positivity, trace preservation, and normality follow from
the predual integral exactly as before.  Let \(\eta:x\to y\).
Using the right-action identity,
\[
 T_\eta\overline\Phi
 =
 \int_{{\cal G}_y}T_\eta T_\zeta\Phi\,d\mu_y(\zeta)
 =
 \int_{{\cal G}_y}T_{\zeta\eta}\Phi\,d\mu_y(\zeta).             \tag{21.15}
\]
The substitution \(\gamma=\zeta\eta\) and (21.13) turn the last
integral into (21.14) at \(x\).  Hence
\(T_\eta\overline\Phi=\overline\Phi_x\), which is equivariance.
The remaining claims follow by integrating their transported
linear equalities or energy inequalities.
\end{proof}

The existence of (21.13) is a genuine gate.  It must not be presumed
for a nonproper residual diffeomorphism groupoid.  Gauge fixing,
slices, or amenability data may supply a replacement, but the
present theorem does not do so.

\subsection{Moving-frame covariant differentiation}

The V20 audit identified a separate source of false derivative
growth.  Let \(U(t)\) be a differentiable unitary frame and
\[
 \Gamma(t)=U(t)^*\dot U(t),\qquad \Gamma(t)^*=-\Gamma(t).        \tag{21.16}
\]
If \(X(t)=U(t)^*\widehat X(t)U(t)\), set
\[
 \nabla_tX=\dot X-[X,\Gamma].                                  \tag{21.17}
\]

\begin{lemma}[Closed trace words do not see frame commutators]
\label{lem:v13v21-trace}
For differentiable square trace-class products,
\[
 {d\over dt}\operatorname{Tr}(X_1\cdots X_m)
 =
 \sum_{j=1}^m
 \operatorname{Tr}\!\left(
 X_1\cdots(\nabla_tX_j)\cdots X_m\right).                      \tag{21.18}
\]
\end{lemma}

\begin{proof}
Since \(\dot X_j=\nabla_tX_j+[X_j,\Gamma]\), the sum of all
commutator insertions is
\[
 [X_1\cdots X_m,\Gamma].                                      \tag{21.19}
\]
Its trace vanishes by cyclicity.
\end{proof}

For a channel with independent input and output frames, let
\[
 \widetilde\Phi_t(A)
 =
 (U^{\rm in}_t)^*
 \Phi_t(U^{\rm out}_tA(U^{\rm out}_t)^*)U^{\rm in}_t.           \tag{21.20}
\]
Writing
\(\Gamma_{\rm in}=(U^{\rm in})^*\dot U^{\rm in}\) and similarly for
the output, direct differentiation gives
\[
\begin{split}
 {d\over dt}\widetilde\Phi_t(A)
 ={}&[\widetilde\Phi_t(A),\Gamma_{\rm in}]
    +\widetilde\Phi_t([\Gamma_{\rm out},A])\\
 &+
 (U^{\rm in}_t)^*
 \dot\Phi_t(U^{\rm out}_tA(U^{\rm out}_t)^*)U^{\rm in}_t.
                                                                    \tag{21.21}
\end{split}
\]
The first term cancels against the input-state frame commutator in
a trace pairing by Lemma~\ref{lem:v13v21-trace}.  The second cancels
against the contravariant frame derivative of a fixed physical
output observable.  Therefore derivative moduli for invariant
expectations must be computed from the last, physical term in
(21.21), not from the norms of the two frame connections.

\subsection{Projective implementations and the anomaly firewall}

Fermionic symmetry implementers may be only projective:
\[
 U_\eta U_\gamma
 =
 \omega(\eta,\gamma)U_{\eta\gamma},\qquad
 |\omega(\eta,\gamma)|=1.                                     \tag{21.22}
\]

\begin{proposition}[Observable descent does not trivialize the
determinant line]
\label{prop:v13v21-projective}
The conjugation maps
\(\operatorname{Ad}U_\gamma\) satisfy the exact groupoid cocycle law
even when (21.22) has nontrivial multiplier.  Hence an
observable-level channel may satisfy (21.4) although the
determinant line is nontrivial.
\end{proposition}

\begin{proof}
The scalar multiplier cancels between \(U\) and \(U^*\):
\[
 \operatorname{Ad}(U_\eta U_\gamma)
 =
 \operatorname{Ad}(\omega(\eta,\gamma)U_{\eta\gamma})
 =
 \operatorname{Ad}U_{\eta\gamma}.                             \tag{21.23}
\]
Thus conjugation is an honest action.  Trivializing the determinant
line would require the multiplier to be a coboundary, which does not
follow from (21.23).
\end{proof}

This separates two gates that must remain distinct.  V21 can descend
the locally normal observable channel when an averaging system
exists.  Gauge and gravitational anomaly cancellation is still
needed for fermionic amplitudes, partition functions, and a fully
defined quantum theory.

\begin{assessment}[Status after thirteenth-cycle V21]
\label{ass:v13v21-status}
A compact residual symmetry, or more generally a groupoid equipped
with an invariant probability system, admits an explicit predual
average of the V19 locally normal channel.  The average is normal,
UCP, regionally natural, energy tight, and exactly equivariant.
Moving unitary frames contribute commutators that cancel in closed
invariant trace observables.  Observable covariance remains
logically separate from determinant-line triviality.
\end{assessment}

\begin{firewall}[Claim boundary after thirteenth-cycle V21]
\label{fire:v13v21-claim}
Properness, a Haar system, a normalized cutoff, coherent unitary
implementers, and invariant RPESM energy operators have not been
constructed.  In particular, no invariant probability system has
been proved for the relevant diffeomorphism groupoid.  Projective
observable covariance does not prove anomaly cancellation.  The
result does not establish causal factorization, the time-slice
axiom, clustering, or the full SM--Einstein continuum.
\end{firewall}

\subsection{Next bottleneck}

The remaining descent obstruction is no longer algebraic
equivariance but anomaly and orbit geometry.  The next note should
formulate a finite-ledger determinant-line cocycle, prove a
quantitative stability and gluing theorem for local trivializations,
and identify the exact cohomology class whose vanishing is required.
This should be kept separate from the proper-groupoid averaging
gate established here.

\subsection{Parent record}

This V21 note appends to the validated thirteenth-cycle V20 file with
record
\[
\begin{array}{c|c}
\text{lines}&6063\\
\text{bytes}&212553\\
\text{SHA--256}&
\texttt{CF4170391E2FF42988B560ED4CB090ECB069C0108D415F02B9D3379F6C18BA28}.
\end{array}
\]

% =============================================================================
% END APPENDED THIRTEENTH-CYCLE V21 MATERIAL
% =============================================================================

% =============================================================================
% APPENDED THIRTEENTH-CYCLE V22 MATERIAL
% =============================================================================

\section{Thirteenth-cycle V22: determinant-line descent to the orbit quotient}
\label{sec:v13v22}

\subsection{What is new here}

Twelfth-cycle V68 already proves compact-exhausted convergence of
determinant-line transition functions, reflected real structures,
connections, sections, independent dual sections, and their paired
complex amplitudes.  Those results are not repeated.

V21 descends observable channels by conjugation, which is insensitive
to scalar projective phases.  This note supplies the missing
distinction at the line level.  Three gates occur in order:

\begin{enumerate}
\item a projective multiplier must be removed to obtain an honest
equivariant determinant line;
\item isotropy must act trivially on its fibres before that line can
descend to the coarse orbit space;
\item the descended line must be topologically trivial, and its
connection must have trivial holonomy if a parallel determinant
phase is required.
\end{enumerate}

These are independent of the normal-channel descent proved in V21.

\subsection{Projective lifts and the multiplier class}

Let \({\cal G}\rightrightarrows X\) act on the limiting Hermitian
determinant line \(\mathscr L\to X\).  A projective lift consists of
unitary maps
\[
 \tau_\gamma:\mathscr L_{s(\gamma)}
       \longrightarrow\mathscr L_{r(\gamma)}                   \tag{22.1}
\]
and a continuous multiplier
\[
 \omega:{\cal G}^{(2)}\longrightarrow{\rm U}(1)                \tag{22.2}
\]
such that
\[
 \tau_\eta\tau_\gamma
 =
 \omega(\eta,\gamma)\tau_{\eta\gamma}.                         \tag{22.3}
\]
Associativity forces the two-cocycle identity
\[
 \omega(\zeta,\eta)\omega(\zeta\eta,\gamma)
 =
 \omega(\eta,\gamma)\omega(\zeta,\eta\gamma).                  \tag{22.4}
\]

A continuous one-cochain \(b:{\cal G}\to{\rm U}(1)\) changes the
lift to
\[
 \tau_\gamma^b=b(\gamma)\tau_\gamma                            \tag{22.5}
\]
and changes the multiplier to
\[
 \omega^b(\eta,\gamma)
 =
 {b(\eta)b(\gamma)\over b(\eta\gamma)}
 \omega(\eta,\gamma).                                         \tag{22.6}
\]

\begin{proposition}[Exact multiplier-removal criterion]
\label{prop:v13v22-multiplier}
The projective lift can be changed by phases into an honest
groupoid action if and only if
\[
 [\omega]=0\quad\text{in}\quad
 H^2_{\rm cont}({\cal G};{\rm U}(1)).                           \tag{22.7}
\]
\end{proposition}

\begin{proof}
An honest action is precisely the equation
\(\omega^b=1\).  By (22.6), this is equivalent to
\[
 \omega(\eta,\gamma)
 =
 {b(\eta\gamma)\over b(\eta)b(\gamma)},                        \tag{22.8}
\]
which says exactly that \(\omega\) is the inverse coboundary of the
one-cochain \(b\).  Conversely, any \(b\) satisfying (22.8) makes
(22.5) associative.
\end{proof}

The finite-cochain Hodge and integral-systole mechanisms summarized
in V1 are applicable to a finite ledger approximation of (22.7).
They must, however, be applied to the actual groupoid multiplier
(22.2).  Cancellation of the local Standard Model anomaly
polynomial alone does not establish (22.7).

\subsection{The stabilizer character}

Assume from now on that the multiplier has been removed.  For
\(x\in X\), let
\[
 {\cal G}_x^x=\{\gamma:s(\gamma)=r(\gamma)=x\}                  \tag{22.9}
\]
be the isotropy group.  Since \(\mathscr L_x\) is one-dimensional,
its action is a unitary character
\[
 \chi_x:{\cal G}_x^x\longrightarrow{\rm U}(1),\qquad
 \tau_\gamma v=\chi_x(\gamma)v.                               \tag{22.10}
\]

\begin{proposition}[Necessary isotropy condition for coarse descent]
\label{prop:v13v22-isotropy-necessary}
If there is a line bundle
\(\overline{\mathscr L}\to X/{\cal G}\) and an equivariant line
isomorphism
\[
 q^*\overline{\mathscr L}\simeq\mathscr L,                     \tag{22.11}
\]
where \(q:X\to X/{\cal G}\), then
\[
 \chi_x\equiv1\qquad\text{for every }x\in X.                   \tag{22.12}
\]
\end{proposition}

\begin{proof}
The pullback fibre at \(x\) is
\(\overline{\mathscr L}_{q(x)}\).  Every isotropy arrow fixes
\(q(x)\), and the canonical groupoid action on a pullback from the
coarse quotient is the identity on that fibre.  Transporting this
action through (22.11) proves (22.12).
\end{proof}

Thus the vanishing of the two-cocycle (22.7) is insufficient.
It produces an honest action, but that action can still carry a
nontrivial stabilizer character.

\subsection{A quantitative stabilizer-character gap}

The stabilizer gate has a useful universal separation when the
whole compact isotropy group is tested.

\begin{lemma}[Universal compact-character gap]
\label{lem:v13v22-character-gap}
Let \(H\) be a compact topological group and
\(\chi:H\to{\rm U}(1)\) a continuous character.  If \(\chi\) is
nontrivial, then
\[
 \sup_{h\in H}|\chi(h)-1|\geq\sqrt3.                            \tag{22.13}
\]
Consequently
\[
 \sup_{h\in H}|\chi(h)-1|<\sqrt3
 \quad\Longrightarrow\quad\chi\equiv1.                         \tag{22.14}
\]
\end{lemma}

\begin{proof}
The image \(\chi(H)\) is a compact subgroup of the circle.  A closed
subgroup of \({\rm U}(1)\) is either the whole circle or the cyclic
group of \(m\)-th roots of unity.  In the first case the supremum in
(22.13) is \(2\).  In the finite nontrivial case \(m\geq2\).  If
\(m\) is even, \(-1\) belongs to the image and the supremum is \(2\).
If \(m\) is odd, the farthest root has distance
\[
 2\cos{\pi\over2m}\geq2\cos{\pi\over6}=\sqrt3.                 \tag{22.15}
\]
This proves the claim.
\end{proof}

The strict threshold in (22.14) can turn a uniform numerical
stabilizer estimate into an exact conclusion.  Sampling finitely
many isotropy elements does not suffice unless a separate
interpolation estimate controls the entire group.

\subsection{Free proper descent}

We first give a complete sufficient theorem without stabilizer
singularities.

\begin{theorem}[Free proper determinant-line descent]
\label{thm:v13v22-free-descent}
Let a locally compact group \(G\) act freely and properly on the
paracompact space \(X\), and assume
\(q:X\to Y=X/G\) is a principal \(G\)-bundle with local continuous
sections.  Let \(\mathscr L\to X\) be a Hermitian line bundle with
an honest continuous unitary \(G\)-action covering the action on
\(X\).  Then
\[
 \overline{\mathscr L}:=\mathscr L/G\longrightarrow Y          \tag{22.16}
\]
is a Hermitian line bundle and the canonical map
\[
 q^*\overline{\mathscr L}\longrightarrow\mathscr L             \tag{22.17}
\]
is an equivariant unitary isomorphism.
\end{theorem}

\begin{proof}
Let \(V\subset Y\) admit a section \(\sigma:V\to X\).  Every
\(x\in q^{-1}(V)\) has a unique representation
\(x=g\sigma(q(x))\).  The \(G\)-action identifies
\(\mathscr L_x\) unitarily with
\(\mathscr L_{\sigma(q(x))}\).  Hence the quotient of
\(\mathscr L|_{q^{-1}(V)}\) by \(G\) is canonically the line bundle
\(\sigma^*\mathscr L\to V\).

On the overlap of two sections, their unique transition
\(g_{ij}(y)\) acts unitarily between the two pulled-back line fibres.
The action cocycle gives the ordinary line-bundle cocycle on triple
overlaps.  These local quotients therefore glue to the Hermitian
line (22.16).  The local action maps define (22.17), and their
compatibility on overlaps makes it a global equivariant unitary
isomorphism.
\end{proof}

For a nonfree proper action, Proposition
\ref{prop:v13v22-isotropy-necessary} remains mandatory.  Trivial
isotropy action plus a valid slice-descent theorem can give the same
conclusion, but orbit-type singularities must be checked.  V22 does
not silently replace the quotient stack by its coarse space.

\subsection{An explicit groupoid descent-atlas criterion}

The following formulation covers nonfree cases when the needed
local orbit data have actually been provided.  Suppose the quotient
\(Y=X/{\cal G}\) has an open cover \((V_i)\), continuous lifts
\[
 \sigma_i:V_i\longrightarrow X,\qquad q\sigma_i={\rm id},       \tag{22.18}
\]
and continuous arrows
\[
 \gamma_{ij}(y):\sigma_j(y)\longrightarrow\sigma_i(y)           \tag{22.19}
\]
on \(V_i\cap V_j\), satisfying
\[
 \gamma_{ii}=1,\qquad
 \gamma_{ij}\gamma_{jk}=\gamma_{ik}.                           \tag{22.20}
\]
Assume every point over \(V_i\) is connected to \(\sigma_i(q(x))\)
by the declared groupoid action and that different choices differ
only by isotropy acting trivially on \(\mathscr L\).

\begin{proposition}[Descent from a populated orbit atlas]
\label{prop:v13v22-atlas}
Under (22.18)--(22.20) and trivial isotropy action, the transition
maps
\[
 \overline g_{ij}(y):=\tau_{\gamma_{ij}(y)}
 :
 \mathscr L_{\sigma_j(y)}\longrightarrow
 \mathscr L_{\sigma_i(y)}                                     \tag{22.21}
\]
define a Hermitian line bundle
\(\overline{\mathscr L}\to Y\) whose pullback is equivariantly
isomorphic to \(\mathscr L\).
\end{proposition}

\begin{proof}
Continuity and unitarity follow from the honest line action.
Equation (22.20) and associativity of \(\tau\) give
\[
 \overline g_{ij}\overline g_{jk}=\overline g_{ik}.             \tag{22.22}
\]
Thus the pulled-back local lines \(\sigma_i^*\mathscr L\) glue over
\(Y\).  To identify the pullback, transport from an arbitrary
\(x\) to the selected representative over \(q(x)\).  If two arrows
perform this transport, their ratio is an isotropy arrow; its fibre
action is the identity.  Hence the identification is independent
of the choice and is equivariant.
\end{proof}

This proposition turns quotient descent into a finite or countable
ledger: quotient charts, representative lifts, overlap arrows,
their cocycle, and the stabilizer-character row must all be
populated.

\subsection{Trivial line, parallel phase, and holonomy}

Descent produces a line on the quotient; it does not produce a
nowhere-zero determinant amplitude.

\begin{proposition}[Three distinct post-descent gates]
\label{prop:v13v22-three-gates}
Let \(Y\) be paracompact and have the homotopy type of a CW complex.
For the descended complex line \(\overline{\mathscr L}\):
\begin{enumerate}
\item it admits a nowhere-zero continuous section if and only if
\[
 c_1(\overline{\mathscr L})=0\in H^2(Y;\mathbb Z);              \tag{22.23}
\]
\item if a unitary connection \(\nabla\) admits a nowhere-zero
parallel section, then its curvature and all loop holonomies are
trivial;
\item conversely, on each connected component, a flat unitary
connection with trivial holonomy admits a nowhere-zero parallel
unit section.
\end{enumerate}
\end{proposition}

\begin{proof}
Complex line bundles over such a space are classified by their
first Chern class, proving the first assertion.  A parallel section
is preserved around every loop and is annihilated by
\(\nabla^2\), proving the second.  For the converse, choose a unit
vector at one base point and extend it by parallel transport.
Flatness makes transport invariant under endpoint-fixed homotopy,
and trivial holonomy makes it independent of the chosen path.
The resulting unit section is parallel.
\end{proof}

Curvature cancellation is therefore not the last condition:
a flat line can retain nontrivial torsion or continuous holonomy.
Likewise, a topologically trivial line can carry a nontrivial
connection.  These distinctions agree with the inherited anomaly
firewalls.

\subsection{Combined determinant-channel descent gate}

\begin{theorem}[Simultaneous observable and determinant descent]
\label{thm:v13v22-combined}
Assume:
\begin{enumerate}
\item the V19 locally normal regional channel satisfies the V21
proper-groupoid averaging hypotheses;
\item the V68 limiting determinant line and independent dual line
carry projective lifts with multipliers whose classes vanish as in
(22.7);
\item after multiplier removal, the line and dual satisfy the free
proper hypotheses of Theorem~\ref{thm:v13v22-free-descent}, or the
populated descent-atlas hypotheses of
Proposition~\ref{prop:v13v22-atlas};
\item the descended paired line has the triviality and, when a
parallel phase is required, holonomy properties in
Proposition~\ref{prop:v13v22-three-gates}.
\end{enumerate}
Then the observable channel descends as a locally normal UCP
regional channel, the determinant line and its dual descend to the
coarse quotient, and their invariant pairing defines a scalar
quotient amplitude.  If the V68 visibility floor and convergence
budgets also hold, its normalized complex functional converges with
the V68 bound.
\end{theorem}

\begin{proof}
Apply Theorem~\ref{thm:v13v21-groupoid} to the channel.  Remove the
two line multipliers using
Proposition~\ref{prop:v13v22-multiplier}.  Descend the resulting
honest line actions by the stated free or atlas theorem.  The
evaluation pairing
\(\mathscr L^*\otimes\mathscr L\to\mathbb C\) is invariant under the
dual actions and therefore factors through the quotient.  The
triviality or holonomy condition supplies the declared scalar phase
choice.  Finally apply the already-proved V68 paired-amplitude
estimate; none of its norm or visibility bounds is changed by
unitary descent.
\end{proof}

\begin{assessment}[Status after thirteenth-cycle V22]
\label{ass:v13v22-status}
The determinant obstruction has been separated into exact,
checkable layers.  The projective multiplier is an
\(H^2_{\rm cont}({\cal G};{\rm U}(1))\) obstruction; an honest line
still has a stabilizer-character obstruction to coarse descent; the
descended line has a Chern obstruction to a nowhere-zero section;
and a flat connection still has a holonomy obstruction to a
parallel phase.  A universal \(\sqrt3\) gap can certify triviality of
a compact stabilizer character from a uniform full-group estimate.
\end{assessment}

\begin{firewall}[Claim boundary after thirteenth-cycle V22]
\label{fire:v13v22-claim}
No RPESM projective multiplier, counterterm cochain, stabilizer
character, quotient atlas, Chern class, or global holonomy has been
computed.  The free proper or explicit descent-atlas hypotheses are
not established for the physical diffeomorphism groupoid.
Observable covariance from V21 remains insufficient.  The theorem
does not establish reflection-positive scalar measure,
clustering, causal factorization, or the full SM--Einstein
continuum.
\end{firewall}

\subsection{Next bottleneck}

The inherited notes already contain local Standard Model anomaly
arithmetic and finite-cochain obstruction machinery.  The next
useful step is to connect those data to the new quotient ledger:
construct a cofinal anomaly packet containing the multiplier,
stabilizer character, curvature, and loop holonomy, then prove that
summable adjacent defects and discrete systoles preserve exact
vanishing through the continuum limit.  This is a stability theorem
for all four obstruction layers, rather than another local anomaly
calculation.

\subsection{Parent record}

This V22 note appends to the validated thirteenth-cycle V21 file with
record
\[
\begin{array}{c|c}
\text{lines}&6457\\
\text{bytes}&226211\\
\text{SHA--256}&
\texttt{E5E641D42839A472C1B1A291A97CB007D2800C04E445AE09473361F4B6E2CFCA}.
\end{array}
\]

% =============================================================================
% END APPENDED THIRTEENTH-CYCLE V22 MATERIAL
% =============================================================================

% =============================================================================
% APPENDED THIRTEENTH-CYCLE V23 MATERIAL
% =============================================================================

\section{Thirteenth-cycle V23: cofinal stability of the full anomaly packet}
\label{sec:v13v23}

\subsection{Why four obstruction layers need two kinds of estimates}

V22 separated multiplier, stabilizer, Chern, curvature, and holonomy
gates.  They do not all obey the same stability principle.
Quantized cohomology classes and full-group characters live in
uniformly separated sets, so a sufficiently small tail forces exact
stabilization.  Curvature is a Banach-space-valued quantity and can
be arbitrarily small without vanishing.  Its continuum value is zero
only when an actual curvature residual tends to zero.

This note gives one cofinal theorem that makes that distinction
explicit.  It does not redo the transition-function convergence of
twelfth-cycle V68.

\subsection{A general discrete-tail lemma}

Let \((A,+)\) be an abelian group with a translation-invariant
complete metric \(d_A\).  Define its zero systole
\[
 {\mathfrak s}_A
 :=
 \inf\{d_A(a,0):a\in A,\ a\ne0\}.                              \tag{23.1}
\]

\begin{lemma}[Summable transport in a discrete obstruction group]
\label{lem:v13v23-discrete}
Assume \({\mathfrak s}_A>0\) and
\[
 \sum_{n=1}^{\infty}d_A(a_{n+1},a_n)<\infty.                   \tag{23.2}
\]
Then \((a_n)\) is eventually constant.  If
\[
 d_A(a_N,0)+
 \sum_{n=N}^{\infty}d_A(a_{n+1},a_n)
 <{\mathfrak s}_A,                                             \tag{23.3}
\]
then its eventual value is zero.
\end{lemma}

\begin{proof}
The summability in (23.2) makes \((a_n)\) Cauchy.  In particular,
\(d_A(a_{n+1},a_n)<{\mathfrak s}_A\) for all sufficiently large
\(n\).  Translation invariance gives
\[
 d_A(a_{n+1}-a_n,0)<{\mathfrak s}_A,                            \tag{23.4}
\]
so the definition of the systole forces \(a_{n+1}=a_n\).  Hence the
sequence is eventually constant, say with value \(a_\infty\).
The triangle inequality and (23.3) give
\(d_A(a_\infty,0)<{\mathfrak s}_A\), which forces
\(a_\infty=0\).
\end{proof}

A finitely generated integral cohomology group can be given such a
metric after separating its free lattice and finite torsion group.
The torsion coordinates must be retained explicitly; a real
harmonic representative cannot detect them.

\subsection{Uniform metrics on character fields}

Let \(H\) be a group and let
\(\operatorname{Char}(H)=\operatorname{Hom}(H,{\rm U}(1))\).
Define
\[
 d_{\rm char}(\chi,\psi)
 :=
 \sup_{h\in H}|\chi(h)-\psi(h)|.                               \tag{23.5}
\]
This is the same as the distance of the quotient character
\(\chi\psi^{-1}\) from the trivial character.

\begin{corollary}[Character sequences are uniformly discrete]
\label{cor:v13v23-character}
For continuous characters of a compact group, or for characters of
an arbitrary group when the supremum in (23.5) is over the entire
group,
\[
 \chi\ne\psi\quad\Longrightarrow\quad
 d_{\rm char}(\chi,\psi)\geq\sqrt3.                            \tag{23.6}
\]
Therefore a sequence of characters with summable adjacent
\(d_{\rm char}\)-distances is eventually constant.  The one-shot
condition
\[
 d_{\rm char}(\chi_N,1)+
 \sum_{n=N}^{\infty}
 d_{\rm char}(\chi_{n+1},\chi_n)<\sqrt3                        \tag{23.7}
\]
forces the limiting character to be trivial.
\end{corollary}

\begin{proof}
The quotient \(\chi\psi^{-1}\) is a nontrivial character when
\(\chi\ne\psi\).  Its image has compact closure in the circle.
The proof of Lemma~\ref{lem:v13v22-character-gap}, applied to that
closure, gives (23.6).  Now use
Lemma~\ref{lem:v13v23-discrete} with systole \(\sqrt3\).
\end{proof}

For stabilizers varying with \(x\), assume they have been transported
to a common group bundle over the compact core and take the supremum
over both \(x\) and the full stabilizer fibre.  If two character
fields differ at one \(x\), (23.6) at that fibre gives the same
global separation.

The full-group supremum is essential.  On \(H=\mathbb Z\), the
characters \(\chi_\theta(m)=e^{im\theta}\) can be arbitrarily close
on any fixed finite set of generators as \(\theta\to0\), even though
they are nontrivial.  A finite loop ledger needs an interpolation,
quantization, or torsion-order certificate before (23.6) can be
invoked.

\subsection{The cofinal packet}

After transporting all stages to one fixed compact ledger screen,
write the obstruction packet at cutoff \(n\) as
\[
 {\mathfrak A}_n
 =
 \bigl(
   a_n^{\rm mult},
   \chi_n^{\rm iso},
   c_n^{\rm det},
   F_n,
   h_n^{\rm hol}
 \bigr).                                                       \tag{23.8}
\]
Its entries have the following declared types.

\begin{enumerate}
\item \(a_n^{\rm mult}\in A_{\rm mult}\) is the quantized
cohomology class of the actual groupoid multiplier.  The group
\(A_{\rm mult}\) has a declared positive systole
\({\mathfrak s}_{\rm mult}\).  This assumption is the finite
cochain integral-lattice gate inherited in V1; it is not automatic
for arbitrary continuous \({\rm U}(1)\) cohomology.
\item \(\chi_n^{\rm iso}\) is the full stabilizer-character field,
measured in the uniform version of (23.5).
\item \(c_n^{\rm det}\in A_{\rm Chern}\) is the integral Chern class,
including its torsion coordinate, with positive systole
\({\mathfrak s}_{\rm Chern}\).
\item \(F_n\) is the determinant-connection curvature in a Banach
space \({\cal B}_F\) of continuous two-forms on the screen.
\item \(h_n^{\rm hol}\) is the holonomy character on the full
declared loop group \(\Pi\), measured by
\[
 d_{\rm hol}(h,k)=\sup_{\ell\in\Pi}|h(\ell)-k(\ell)|.           \tag{23.9}
\]
\end{enumerate}

Suppose adjacent cutoff transports give bounds
\[
\begin{split}
 d_{\rm mult}(a_{n+1}^{\rm mult},a_n^{\rm mult})
     &\leq\varepsilon_n^{\rm mult},\\
 d_{\rm char}(\chi_{n+1}^{\rm iso},\chi_n^{\rm iso})
     &\leq\varepsilon_n^{\rm iso},\\
 d_{\rm Chern}(c_{n+1}^{\rm det},c_n^{\rm det})
     &\leq\varepsilon_n^{\rm Chern},\\
 \|F_{n+1}-F_n\|_{{\cal B}_F}
     &\leq\varepsilon_n^F,\\
 d_{\rm hol}(h_{n+1}^{\rm hol},h_n^{\rm hol})
     &\leq\varepsilon_n^{\rm hol}.
                                                                    \tag{23.10}
\end{split}
\]
All comparisons in (23.10) must be made after the physical pullback
and the declared cochain, line, and loop transports.  Rephasing each
cutoff independently can conceal the actual adjacent defect.

\subsection{Simultaneous stabilization theorem}

\begin{theorem}[Cofinal anomaly-packet stability]
\label{thm:v13v23-packet}
Assume
\[
 \sum_{n=1}^{\infty}
 \left(
 \varepsilon_n^{\rm mult}+\varepsilon_n^{\rm iso}
 +\varepsilon_n^{\rm Chern}+\varepsilon_n^F
 +\varepsilon_n^{\rm hol}
 \right)<\infty.                                               \tag{23.11}
\]
Then:
\begin{enumerate}
\item \(a_n^{\rm mult}\), \(\chi_n^{\rm iso}\),
\(c_n^{\rm det}\), and \(h_n^{\rm hol}\) are eventually constant;
\item \(F_n\) converges in \({\cal B}_F\) to a curvature
\(F_\infty\);
\item if for some \(N\),
\[
\begin{split}
 d_{\rm mult}(a_N^{\rm mult},0)
 +\sum_{n=N}^\infty\varepsilon_n^{\rm mult}
 &<{\mathfrak s}_{\rm mult},\\
 d_{\rm char}(\chi_N^{\rm iso},1)
 +\sum_{n=N}^\infty\varepsilon_n^{\rm iso}
 &<\sqrt3,\\
 d_{\rm Chern}(c_N^{\rm det},0)
 +\sum_{n=N}^\infty\varepsilon_n^{\rm Chern}
 &<{\mathfrak s}_{\rm Chern},\\
 d_{\rm hol}(h_N^{\rm hol},1)
 +\sum_{n=N}^\infty\varepsilon_n^{\rm hol}
 &<\sqrt3,
                                                                    \tag{23.12}
\end{split}
\]
then all four discrete or character obstructions vanish exactly in
the continuum packet;
\item the curvature vanishes exactly if, in addition,
\[
 \|F_n\|_{{\cal B}_F}\longrightarrow0.                          \tag{23.13}
\]
\end{enumerate}
\end{theorem}

\begin{proof}
Apply Lemma~\ref{lem:v13v23-discrete} to the multiplier and Chern
classes.  Apply Corollary~\ref{cor:v13v23-character} to the
stabilizer and full-loop holonomy characters.  This proves eventual
constancy and the four conclusions from (23.12).

The fourth adjacent bound in (23.10) and (23.11) make \(F_n\)
Cauchy in the Banach space \({\cal B}_F\), so it has a limit.
Condition (23.13) identifies that limit as zero.  Without (23.13),
summable increments give convergence but do not determine the
limit.
\end{proof}

\subsection{Convergence of the counterterm trivializations}

Vanishing of the multiplier class at every sufficiently large stage
gives one-cochains \(b_n\) satisfying
\[
 \omega_n(\eta,\gamma)
 =
 {b_n(\eta\gamma)\over b_n(\eta)b_n(\gamma)}.                  \tag{23.14}
\]
These cochains are not unique: they can differ by a one-cocycle.
Fix a common normalization, for example values on a declared
spanning tree of the finite groupoid ledger.

\begin{proposition}[Cofinal multiplier repair]
\label{prop:v13v23-repair}
Suppose the normalized repairs satisfy
\[
 \sum_{n=N}^{\infty}
 \|b_{n+1}-b_n\|_{C^0({\cal G}_{\rm screen})}<\infty,           \tag{23.15}
\]
and \(\omega_n\to\omega_\infty\) uniformly on composable pairs.
Then \(b_n\) converges uniformly to a continuous \(b_\infty\), and
\[
 \omega_\infty(\eta,\gamma)
 =
 {b_\infty(\eta\gamma)\over
  b_\infty(\eta)b_\infty(\gamma)}.                             \tag{23.16}
\]
Thus the continuum projective lift has an actual continuous
counterterm repair, not merely a vanishing abstract class.
\end{proposition}

\begin{proof}
The summable estimate makes \(b_n\) uniformly Cauchy.  The circle is
complete and closed, so its uniform limit is continuous and
circle-valued.  Multiplication and inversion on the circle are
continuous.  Pass to the uniform limit in (23.14) to obtain
(23.16).
\end{proof}

Class vanishing without (23.15) is insufficient to select coherent
cutoff counterterms.  This is the line-level analogue of requiring
one subsequence over all backgrounds and regions in V18--V19.

\subsection{Descent after the packet vanishes}

\begin{corollary}[Stable input for V22 descent]
\label{cor:v13v23-descent}
Assume Theorem~\ref{thm:v13v23-packet}, the repair convergence in
Proposition~\ref{prop:v13v23-repair}, and the V22 free proper or
populated orbit-atlas hypotheses.  If (23.12)--(23.13) hold, then:
\begin{enumerate}
\item the limiting determinant line has an honest groupoid action;
\item its stabilizers act trivially, so the line descends to the
coarse quotient under the V22 descent hypotheses;
\item the descended line has zero first Chern class;
\item its limiting connection is flat with trivial full-loop
holonomy and hence has a parallel unit section on every connected
component.
\end{enumerate}
The V21 locally normal equivariant channel and this parallel
determinant phase may therefore be placed on the same quotient
ledger.
\end{corollary}

\begin{proof}
Items one through three follow respectively from (23.16), the
isotropy conclusion in (23.12), and the Chern conclusion in
(23.12), followed by V22 descent.  Condition (23.13) gives flatness,
and the holonomy conclusion in (23.12) gives path independence.
Proposition~\ref{prop:v13v22-three-gates} supplies the parallel unit
section.  V21 places the observable channel on the same quotient.
\end{proof}

This corollary concerns a determinant phase and a complex paired
functional.  Reflection positivity of the resulting scalar
Schwinger functional remains a separate cone condition.

\subsection{What finite sampling cannot certify}

There are two sharp limitations.

First, a small multiplier on finitely many composable pairs does not
prove that its continuous cohomology class lies in the discrete
group \(A_{\rm mult}\).  The integral-lattice membership row must be
checked before the systole is applied.

Second, small holonomy on finitely many loops does not imply trivial
holonomy on the full loop group.  If the loop group is finitely
presented and the connection is flat, values on a generating set do
determine the character, but arbitrarily small nonzero generator
phases still exist unless a quantization or bounded-order torsion
condition supplies a discrete gap.

\begin{assessment}[Status after thirteenth-cycle V23]
\label{ass:v13v23-status}
All determinant-anomaly obstruction layers now have a common
cofinal stability theorem.  Quantized multiplier and Chern classes,
full stabilizer characters, and full-loop holonomy characters
stabilize exactly under summable adjacent transport.  Explicit
systole inequalities force their continuum values to vanish.
Curvature converges under the same tail mechanism but vanishes only
under an independent zero residual.  Summable normalized
counterterms produce an actual continuum multiplier repair.
\end{assessment}

\begin{firewall}[Claim boundary after thirteenth-cycle V23]
\label{fire:v13v23-claim}
No common RPESM anomaly screen, discrete multiplier group, integral
lattice membership, systole, full stabilizer estimate, Chern packet,
curvature residual, full-loop holonomy estimate, or normalized
counterterm tail has been populated.  Finite sampling is explicitly
insufficient for the character and holonomy conclusions.  The
result does not prove reflection positivity, causal factorization,
clustering, or the full SM--Einstein continuum.
\end{firewall}

\subsection{Next bottleneck}

With local normality, regional compatibility, quotient covariance,
and determinant-anomaly descent reduced to explicit certificates,
the next structural gate is causal compatibility.  The regional UCP
channel must respect spacelike commutation and causal factorization,
and its continuum state must satisfy a time-slice property.  The
next note should formulate these as closed finite-ledger relations
and prove their passage through the V19 cofinal normal limit without
confusing local algebraic causality with clustering.

\subsection{Parent record}

This V23 note appends to the validated thirteenth-cycle V22 file with
record
\[
\begin{array}{c|c}
\text{lines}&6869\\
\text{bytes}&241578\\
\text{SHA--256}&
\texttt{B717B282B60B0B7C2CA87EC787C9089766CDE359AED672BBDD684B9FD9FF8685}.
\end{array}
\]

% =============================================================================
% END APPENDED THIRTEENTH-CYCLE V23 MATERIAL
% =============================================================================

% =============================================================================
% APPENDED THIRTEENTH-CYCLE V24 MATERIAL
% =============================================================================

\section{Thirteenth-cycle V24: net causality, split factorization, and an averaging obstruction}
\label{sec:v13v24}

\subsection{Relation to the inherited locality results}

The frozen notes contain causal response estimates, exponential and
Coulomb locality bounds, and a time-slice-local gravitational
constraint substitution.  This note does not repeat them.  Its
purpose is different: it asks which algebraic causality properties
survive the locally normal regional limit of V19 and the equivariant
descent of V21.

Three notions must be separated:

\begin{enumerate}
\item Einstein causality is commutation of spacelike local algebras;
\item split product factorization identifies a spacelike joint
channel with the tensor product of its two marginal channels;
\item the time-slice axiom says that a neighborhood of a Cauchy
surface generates the algebra of its causal completion.
\end{enumerate}

The first follows from regional naturality and locality of the input
net.  The second needs a separate tensor certificate.  The third is
a property of the two nets and then determines the natural channel
on a causal completion.

\subsection{Einstein causality from regional naturality}

Let \({\cal O}_1\) and \({\cal O}_2\) be spacelike separated and
contained in \({\cal O}\).  Write the inclusions as
\[
 \alpha^{\rm in}_{i{\cal O}}:{\cal M}^{\rm in}_i
       \longrightarrow{\cal M}^{\rm in}_{\cal O},\qquad
 \alpha^{\rm out}_{i{\cal O}}:{\cal M}^{\rm out}_i
       \longrightarrow{\cal M}^{\rm out}_{\cal O}.             \tag{24.1}
\]
Assume the input net satisfies Einstein causality:
\[
 [\alpha^{\rm in}_{1{\cal O}}(C),
   \alpha^{\rm in}_{2{\cal O}}(D)]=0                           \tag{24.2}
\]
for all local \(C,D\).

\begin{theorem}[Causal commutation of channel images]
\label{thm:v13v24-commutation}
Let \((\Phi^{(r)})_r\) be the regionally natural locally normal
channel of V19.  Then, for
\(A\in{\cal M}^{\rm out}_1\) and
\(B\in{\cal M}^{\rm out}_2\),
\[
 \left[
 \Phi^{({\cal O})}
   (\alpha^{\rm out}_{1{\cal O}}A),
 \Phi^{({\cal O})}
   (\alpha^{\rm out}_{2{\cal O}}B)
 \right]=0.                                                    \tag{24.3}
\]
The same statement holds for the V21 equivariant average when the
symmetry transports commute with the regional inclusions.
\end{theorem}

\begin{proof}
Regional naturality gives
\[
\begin{split}
 \Phi^{({\cal O})}\alpha^{\rm out}_{1{\cal O}}(A)
 &=\alpha^{\rm in}_{1{\cal O}}\Phi^{(1)}(A),\\
 \Phi^{({\cal O})}\alpha^{\rm out}_{2{\cal O}}(B)
 &=\alpha^{\rm in}_{2{\cal O}}\Phi^{(2)}(B).                   \tag{24.4}
\end{split}
\]
Their commutator vanishes by (24.2).  The V21 average remains
regionally natural, so the same proof applies after averaging.
\end{proof}

This proof does not take a limit of operator products.  That is
important because multiplication is not jointly continuous in the
ultraweak topology.  Naturality places each limiting image directly
inside a commuting local input algebra.

The output net must independently satisfy Einstein causality if it
is to be a physical observable net.  Equation (24.3) only states
that the channel does not spoil commutation after mapping into the
input net.

\subsection{Split product channels}

Assume the spacelike pair has declared normal spatial split
isomorphisms
\[
 \sigma_{\rm out}:
 {\cal M}^{\rm out}_1\ \overline\otimes\
 {\cal M}^{\rm out}_2
 \longrightarrow{\cal M}^{\rm out}_{12},\qquad
 \sigma_{\rm in}:
 {\cal M}^{\rm in}_1\ \overline\otimes\
 {\cal M}^{\rm in}_2
 \longrightarrow{\cal M}^{\rm in}_{12}.                       \tag{24.5}
\]
A split-factorizing channel obeys
\[
 \Phi^{(12)}\sigma_{\rm out}(A\otimes B)
 =
 \sigma_{\rm in}
 \bigl(\Phi^{(1)}(A)\otimes\Phi^{(2)}(B)\bigr).                 \tag{24.6}
\]

\begin{lemma}[Ultraweak convergence of bounded simple tensors]
\label{lem:v13v24-tensor-limit}
If \(C_n\to C\) ultraweakly in \({\cal M}_1\),
\(D_n\to D\) ultraweakly in \({\cal M}_2\), and both sequences are
norm bounded, then
\[
 C_n\otimes D_n\longrightarrow C\otimes D                      \tag{24.7}
\]
ultraweakly in
\({\cal M}_1\overline\otimes{\cal M}_2\).
\end{lemma}

\begin{proof}
For a product normal functional \(f\otimes g\),
\[
 (f\otimes g)(C_n\otimes D_n)=f(C_n)g(D_n)
 \longrightarrow f(C)g(D).                                   \tag{24.8}
\]
Finite linear combinations of product normal functionals are norm
dense in the predual of the spatial tensor product.  The tensors in
(24.7) have a common norm bound.  Approximate an arbitrary normal
functional in predual norm by such a finite sum, use (24.8), and
then remove the approximation.
\end{proof}

\begin{theorem}[Passage of split factorization]
\label{thm:v13v24-factor-limit}
Suppose the cofinal stage channels in V19 satisfy (24.6) for a fixed
spacelike pair and fixed split isomorphisms.  Then the locally
normal limiting channels satisfy (24.6).  The conclusion also holds
when the two sides have an ultraweak residual tending to zero
uniformly on each fixed finite-ledger set of \(A,B\) and normal test
functionals.
\end{theorem}

\begin{proof}
V19 gives point-ultraweak convergence of every local Heisenberg
channel along one common subsequence.  For fixed \(A,B\), unital
positivity gives the uniform bounds
\[
 \|\Phi_n^{(1)}(A)\|\leq\|A\|,\qquad
 \|\Phi_n^{(2)}(B)\|\leq\|B\|.                                \tag{24.9}
\]
Lemma~\ref{lem:v13v24-tensor-limit} therefore passes the right side
of the stage identity (24.6) to the tensor product of the limiting
marginals.  Point-ultraweak convergence of the joint channel passes
the left side.  Normality of the split isomorphisms completes the
proof.  A residual tending to zero disappears after testing against
an arbitrary fixed normal functional.
\end{proof}

At a finite matrix ledger, (24.6) is equivalent, after fixing one
Choi convention and its tensor-factor ordering, to
\[
 C_{12}=C_1\otimes C_2.                                       \tag{24.10}
\]
This is a closed algebraic condition but not a convex one.

\subsection{Equivariant averaging can destroy factorization}

The convex V21 group average need not preserve (24.10).

\begin{proposition}[Correlated-average obstruction]
\label{prop:v13v24-average-obstruction}
There exist product channels
\(\Phi_g^{(1)}\otimes\Phi_g^{(2)}\), indexed by a compact group,
whose diagonal group average is not the tensor product of the two
averaged marginal channels.
\end{proposition}

\begin{proof}
Take \(G=\mathbb Z_2=\{0,1\}\), two qubit algebras, and
\[
 \Phi_0^{(i)}={\rm id},\qquad
 \Phi_1^{(i)}=\operatorname{Ad}X,                              \tag{24.11}
\]
where \(X\) is the Pauli flip.  The diagonal average of the product
channels is
\[
 \Psi_{12}
 =
 {1\over2}\bigl(
 {\rm id}\otimes{\rm id}
 +\operatorname{Ad}X\otimes\operatorname{Ad}X
 \bigr).                                                       \tag{24.12}
\]
Each marginal average is
\[
 \Psi_i={1\over2}({\rm id}+\operatorname{Ad}X).                 \tag{24.13}
\]
For the Pauli matrix \(Z\), \(XZX=-Z\).  Hence
\[
 \Psi_{12}(Z\otimes Z)=Z\otimes Z,\qquad
 (\Psi_1\otimes\Psi_2)(Z\otimes Z)=0.                          \tag{24.14}
\]
Thus \(\Psi_{12}\ne\Psi_1\otimes\Psi_2\).
\end{proof}

The group average introduces shared classical randomness: the same
group element is applied to both regions.  It preserves complete
positivity, normality, regional naturality, and the commutation
statement (24.3), but can break statistical product factorization.

\begin{corollary}[Factorization-compatible order of construction]
\label{cor:v13v24-order}
Suppose the two marginal channels are separately made equivariant
and their tensor product is compatible with the diagonal symmetry
and the split isomorphisms.  Defining the joint channel by
\[
 \Phi^{(12)}
 =
 \sigma_{\rm in}
 (\Phi^{(1)}\overline\otimes\Phi^{(2)})
 \sigma_{\rm out}^{-1}                                        \tag{24.15}
\]
gives an equivariant split-factorizing normal UCP channel.
\end{corollary}

\begin{proof}
The spatial tensor product of normal UCP maps is normal UCP.
Equation (24.15) gives (24.6) by construction.  Equivariance of each
factor and equivariance of the two split isomorphisms give
equivariance of their tensor product.
\end{proof}

This construction is available only if the desired physical joint
channel is a product channel on the declared split pair.  Vacuum
correlations and interacting causal factorization generally require
a different structure.  They must not be erased merely to make
(24.10) true.

\subsection{Time-slice determination}

Let \({\cal N}\subset{\cal O}\) contain a Cauchy surface for
\({\cal O}\).  Suppose both nets satisfy the time-slice axiom, so
the inclusion morphisms
\[
 \alpha^{\rm in}_{{\cal N}{\cal O}},\qquad
 \alpha^{\rm out}_{{\cal N}{\cal O}}                           \tag{24.16}
\]
are normal \(*\)-isomorphisms.

\begin{theorem}[Time-slice determination of a natural channel]
\label{thm:v13v24-timeslice}
A regionally natural channel satisfies
\[
 \boxed{
 \Phi^{({\cal O})}
 =
 \alpha^{\rm in}_{{\cal N}{\cal O}}\,
 \Phi^{({\cal N})}\,
 \bigl(\alpha^{\rm out}_{{\cal N}{\cal O}}\bigr)^{-1}.}         \tag{24.17}
\]
Hence its value on the causal completion is uniquely determined by
its value on the time-slice neighborhood.  Normality and complete
positivity are preserved by (24.17).
\end{theorem}

\begin{proof}
Regional naturality is
\[
 \alpha^{\rm in}_{{\cal N}{\cal O}}\Phi^{({\cal N})}
 =
 \Phi^{({\cal O})}\alpha^{\rm out}_{{\cal N}{\cal O}}.         \tag{24.18}
\]
Compose on the right with the inverse output inclusion to obtain
(24.17).  Normal \(*\)-isomorphisms and their inverses are normal
and completely positive, proving the last statement.
\end{proof}

The theorem does not prove the time-slice axiom for either limiting
net.  It says that once the two time-slice isomorphisms have survived
the regulator limit, the natural channel automatically intertwines
them.

\subsection{Finite-ledger causal packet}

For every declared spacelike pair and Cauchy inclusion, the
finite-ledger causal packet must contain:

\begin{enumerate}
\item input and output locality commutators;
\item regional naturality squares;
\item split isomorphisms when product factorization is asserted;
\item joint and marginal Choi matrices with the tensor equation
(24.10), or an explicitly vanishing tested residual;
\item time-slice inclusion maps and certified inverses;
\item compatibility of symmetry transport with all these maps.
\end{enumerate}

The first two rows pass to V19 by Theorem
\ref{thm:v13v24-commutation}.  The product row passes by Theorem
\ref{thm:v13v24-factor-limit}.  The time-slice row passes only if
the limiting inclusions remain isomorphisms; pointwise convergence
of finite-dimensional inverses without a uniform inverse bound does
not suffice.

\begin{assessment}[Status after thirteenth-cycle V24]
\label{ass:v13v24-status}
Regional naturality now gives exact causal commutation of the
limiting channel images without multiplying ultraweak limits.
Under a declared split property, tensor-product channel
factorization passes through the V19 normal limit.  Natural channels
intertwine any surviving time-slice isomorphisms.  A concrete
two-qubit example proves that diagonal equivariant averaging can
destroy product factorization by introducing shared randomness.
\end{assessment}

\begin{firewall}[Claim boundary after thirteenth-cycle V24]
\label{fire:v13v24-claim}
No RPESM local von Neumann net, Einstein-causality relation, split
inclusion, product Choi identity, time-slice isomorphism, or
factorization-compatible equivariant selection has been populated.
Split product factorization is stronger than algebraic locality and
is not asserted for an interacting vacuum.  The result does not
establish perturbative causal factorization of time-ordered
products, clustering, or the full SM--Einstein continuum.
\end{firewall}

\subsection{Next bottleneck}

V25 is the next compulsory companion audit.  After it, the programme
must choose the correct interacting replacement for the overly
strong product equation (24.10).  A useful target is a causal
conditional-expectation or Markov-square identity on nested
diamonds, formulated as a closed predual relation and compatible
with the V21 symmetry descent.  Any such choice must retain vacuum
correlations while preventing superluminal signalling.

\subsection{Parent record}

This V24 note appends to the validated thirteenth-cycle V23 file with
record
\[
\begin{array}{c|c}
\text{lines}&7250\\
\text{bytes}&255858\\
\text{SHA--256}&
\texttt{22B078EFDC19D93FD063A9F38E557AEEE0F6504704D17B73304789CBFB548BBB}.
\end{array}
\]

% =============================================================================
% END APPENDED THIRTEENTH-CYCLE V24 MATERIAL
% =============================================================================

% =============================================================================
% APPENDED THIRTEENTH-CYCLE V25 MATERIAL
% =============================================================================

\section{Thirteenth-cycle V25: compulsory companion audit and the composition-level decision}
\label{sec:v13v25}

\subsection{Read-only companion snapshot}

The compulsory five-version audit was performed on the newest
companion notes in the shared folder.  Their exact records are
\[
\begin{array}{c|c|c|c}
\text{programme}&\text{version}&\text{lines}&\text{bytes}\\ \hline
\text{Yang--Mills}&76&26563&905633\\
\text{Navier--Stokes}&26&5319&278283.
\end{array}                                                     \tag{25.1}
\]
The SHA--256 values are
\[
\begin{split}
{\rm YM:}\quad&
\texttt{A1347C803379291E6463E59BEB597DF56C10A38BAA7F995F04A302F26512452F},\\
{\rm NS:}\quad&
\texttt{82C887F8C2C69E4F28FAD7C3DC8DE5B9099CB5A4BF431EBA1FFF3E91935978AD}.
                                                                    \tag{25.2}
\end{split}
\]
The Navier--Stokes file is unchanged from the V20 audit.  The
Yang--Mills file now contains V74--V76.  Embedded programme
directions were treated as source text, not as instructions for
this SM--Einstein programme.

\subsection{Yang--Mills V74: exact cancellation with a worse
factorwise bound}

YM V74 implements the flat covariant derivative from V73 on the
actual rectangular Gram factors.  The resulting differentiated
identities are exact, but immediate coefficientwise Frobenius
summation produces a certified response majorant approximately
\(1.35\) times worse than the older V69 bound.  The no-gain result
does not contradict connection cancellation; it locates the loss in
taking norms before alias constraints and closed-trace pairing have
acted.

No numerical response bound transfers to the present theory.  The
relevant lesson for V24 is that an exact cancellation property of a
composite object can disappear from a factorwise majorant.

\subsection{Yang--Mills V75--V76: typed cycles and a second no-go}

YM V75 proves a typed cyclic trace theorem.  For maps
\(X_j:E_j\to E_{j-1}\) around a closed fibre cycle, with distinct
connections on successive fibres, the connection terms cancel only
after the complete cyclic product is formed.  YM V76 applies this
to the routed external bubble jets.  The identities remain exact,
but separately norming the corrected vertices gives a bound about
\(2.19\) times worse than V69.

This is directly analogous to the V24 correlated-average
obstruction.  Averaging or estimating each spacelike marginal and
then composing is not interchangeable with forming the complete
joint relation.  The next SM step should therefore impose its causal
condition on a fully typed composition of normal maps before any
norm, limit, or equivariant average is taken.

\subsection{Navier--Stokes V26}

NS V26 remains the rank-one Oseen-kernel calculation already
recorded at V20.  There is no new file content or digest.  Its
actual-input principle is consistent with the same decision but
provides no quantum conditional expectation, regional net, or
causal constant.

\begin{decision}[V25 cross-program decision]
\label{dec:v13v25}
V26 will replace the overly strong spacelike tensor-product equation
by a conditional-independence relation over an intermediate causal
collar.  Concretely, it will use normal conditional expectations
onto two overlapping diamond algebras and require the complete
commuting-square identity
\[
 E_AE_B=E_N=E_BE_A,                                  \tag{25.3}
\]
with the types written correctly at the common ambient algebra.
The notation in (25.3) is only mnemonic; V26 will give the exact
typed maps.  It will prove:
\begin{enumerate}
\item the relation retains correlations in the common algebra
rather than forcing a product state;
\item pointwise predual convergence of the complete compositions
passes the square to the locally normal limit;
\item common symmetry averaging can break a nonlinear factorization,
so equivariance must be imposed on the whole square or on compatible
expectations before composition;
\item the resulting Markov property remains distinct from Einstein
causality, the time-slice axiom, and clustering.
\end{enumerate}
\end{decision}


\begin{assessment}[Status after thirteenth-cycle V25]
\label{ass:v13v25-status}
The required audit is complete.  YM V74 and V76 give rigorous
negative evidence against premature factorwise norming, while V75
provides the typed composition principle.  The SM direction is
therefore adjusted from split product channels to a full
conditional-expectation commuting square that can preserve
interacting correlations.
\end{assessment}

\begin{firewall}[Claim boundary after thirteenth-cycle V25]
\label{fire:v13v25-claim}
No YM numerical bound or NS kernel constant is imported.  Neither
companion proves a causal Markov property for the RPESM net.  The
audit does not populate the conditional expectations, their
symmetry covariance, or their continuum convergence, and it does
not establish the full SM--Einstein continuum.
\end{firewall}

\subsection{Parent record}

This V25 note appends to the validated thirteenth-cycle V24 file with
record
\[
\begin{array}{c|c}
\text{lines}&7608\\
\text{bytes}&268912\\
\text{SHA--256}&
\texttt{457494ED6C1996CE3255F90F7A49DB37AE27C2DD2182484ECD112C17022C2BB3}.
\end{array}
\]

% =============================================================================
% END APPENDED THIRTEENTH-CYCLE V25 MATERIAL
% =============================================================================

% =============================================================================
% APPENDED THIRTEENTH-CYCLE V26 MATERIAL
% =============================================================================

\section{Thirteenth-cycle V26: causal commuting squares and their normal limit}
\label{sec:v13v26}

\subsection{The interacting replacement for product factorization}

Let \({\cal M}\) be a von Neumann algebra and
\[
 {\cal N}\subset{\cal A}\cap{\cal B}\subset{\cal M}.            \tag{26.1}
\]
The intended geometry is that \({\cal A}\) is generated by a past
diamond and an intermediate collar, \({\cal B}\) by that collar and
a future diamond, and \({\cal N}\) by the collar.  Let
\[
 E_{\cal A}:{\cal M}\to{\cal A},\qquad
 E_{\cal B}:{\cal M}\to{\cal B},\qquad
 E_{\cal N}:{\cal M}\to{\cal N}                                \tag{26.2}
\]
be normal conditional expectations, regarded as normal UCP
idempotents from \({\cal M}\) into itself.

\begin{definition}[Causal conditional-expectation square]
\label{def:v13v26-square}
The quadruple
\(({\cal N},{\cal A},{\cal B},{\cal M})\), with (26.2), is a
commuting square when
\[
 \boxed{
 E_{\cal A}E_{\cal B}=E_{\cal N}
 =E_{\cal B}E_{\cal A}.}                                      \tag{26.3}
\]
\end{definition}

Unlike a tensor-product state, (26.3) permits arbitrary correlations
carried by \({\cal N}\).  It is a noncommutative conditional
independence relation over the collar.

\begin{proposition}[The common algebra is the intersection]
\label{prop:v13v26-intersection}
If (26.3) holds, then
\[
 {\cal A}\cap{\cal B}={\cal N}.                                \tag{26.4}
\]
\end{proposition}

\begin{proof}
The inclusion \({\cal N}\subset{\cal A}\cap{\cal B}\) is assumed.
If \(X\in{\cal A}\cap{\cal B}\), both \(E_{\cal A}\) and
\(E_{\cal B}\) fix \(X\).  Equation (26.3) gives
\[
 E_{\cal N}X=E_{\cal A}E_{\cal B}X=X,                          \tag{26.5}
\]
so \(X\in{\cal N}\).
\end{proof}

\subsection{Classical conditional independence is included}

Let
\({\cal M}=L^\infty(\Omega,\mathcal F,\mathbb P)\),
\({\cal A}=L^\infty(\sigma(X,Z))\),
\({\cal B}=L^\infty(\sigma(Y,Z))\), and
\({\cal N}=L^\infty(\sigma(Z))\).  Take the ordinary conditional
expectations.

\begin{proposition}[Classical meaning of the square]
\label{prop:v13v26-classical}
The square identity (26.3) holds if \(X\) and \(Y\) are conditionally
independent given \(Z\).  In particular, the square can retain
nonzero unconditional correlations mediated by \(Z\).
\end{proposition}

\begin{proof}
For a bounded \(f(Y,Z)\),
conditional independence gives
\[
 \mathbb E[f(Y,Z)\mid X,Z]
 =
 \mathbb E[f(Y,Z)\mid Z].                                     \tag{26.6}
\]
Thus \(E_{\cal A}E_{\cal B}=E_{\cal N}\) first on bounded functions
and then on all of \(L^\infty\).  The reverse identity follows
symmetrically.  If, for example, both \(X\) and \(Y\) depend on a
nonconstant \(Z\), their unconditional covariance can be nonzero.
\end{proof}

Hence (26.3) avoids the loss of interacting or background-mediated
correlations caused by the split product equation in V24.

\subsection{State-preserving squares}

Let \(\varphi\) be a faithful normal state on \({\cal M}\), and
assume the three expectations preserve \(\varphi\).  They act as
orthogonal projections on \(L^2({\cal M},\varphi)\).

\begin{lemma}[One order determines the other]
\label{lem:v13v26-one-order}
For \(\varphi\)-preserving conditional expectations, the identity
\[
 E_{\cal A}E_{\cal B}=E_{\cal N}                               \tag{26.7}
\]
implies
\[
 E_{\cal B}E_{\cal A}=E_{\cal N}.                              \tag{26.8}
\]
\end{lemma}

\begin{proof}
On the GNS Hilbert space, a \(\varphi\)-preserving conditional
expectation is the orthogonal projection onto the closed
\(L^2\)-subspace generated by its range.  Taking Hilbert-space
adjoints of
\(P_{\cal A}P_{\cal B}=P_{\cal N}\) gives
\(P_{\cal B}P_{\cal A}=P_{\cal N}\).  The GNS image of
\({\cal M}\) is dense and \(\varphi\) is faithful, so the operator
identity gives (26.8) on \({\cal M}\).
\end{proof}

The state-preserving formulation also gives a symmetry mechanism.
If a symmetry preserves \(\varphi\) and the three subalgebras, and
the \(\varphi\)-preserving expectation onto each subalgebra is
unique, then transport by the symmetry produces the same
expectation.  The complete square is already equivariant and no
convex average of its factors is needed.

\subsection{Predual convergence of typed compositions}

Let
\[
 E_{R,n}:{\cal M}\to{\cal M},
 \qquad R\in\{{\cal A},{\cal B},{\cal N}\},                    \tag{26.9}
\]
be normal conditional expectations with preadjoints
\(S_{R,n}\).  Assume
\[
 \|S_{R,n}\rho-S_R\rho\|\longrightarrow0
 \quad\text{for every }\rho\in{\cal M}_*.                      \tag{26.10}
\]
The limits are CPTP contractions by V17.

\begin{lemma}[Complete-composition convergence]
\label{lem:v13v26-composition}
Under (26.10),
\[
 S_{{\cal B},n}S_{{\cal A},n}\rho
 \longrightarrow
 S_{\cal B}S_{\cal A}\rho                                     \tag{26.11}
\]
in predual norm for every \(\rho\).  The analogous statement holds
with \({\cal A}\) and \({\cal B}\) reversed.
\end{lemma}

\begin{proof}
Use contractivity and insert
\(S_{{\cal B},n}S_{\cal A}\rho\):
\[
\begin{split}
 \|S_{{\cal B},n}S_{{\cal A},n}\rho
      -S_{\cal B}S_{\cal A}\rho\|
 \leq{}&
 \|S_{{\cal A},n}\rho-S_{\cal A}\rho\|\\
 &+
 \|(S_{{\cal B},n}-S_{\cal B})S_{\cal A}\rho\|.
                                                                    \tag{26.12}
\end{split}
\]
The first term tends to zero by (26.10), and the second does so
because its input \(S_{\cal A}\rho\) is fixed.
\end{proof}

The reversal of order in (26.11) is the preadjoint typing:
\[
 (E_{\cal A}E_{\cal B})_*=S_{\cal B}S_{\cal A}.                \tag{26.13}
\]

\subsection{The normal commuting-square limit}

Assume the ranges \({\cal A},{\cal B},{\cal N}\) are fixed after the
physical transport to a common regional representation, and every
\(E_{R,n}\) fixes its declared range pointwise.

\begin{theorem}[Closure of normal commuting squares]
\label{thm:v13v26-limit}
If every cutoff satisfies
\[
 E_{{\cal A},n}E_{{\cal B},n}
 =
 E_{{\cal N},n}
 =
 E_{{\cal B},n}E_{{\cal A},n},                                \tag{26.14}
\]
and (26.10) holds, then the adjoint limits
\(E_R=S_R^*\) are normal conditional expectations onto their
declared ranges and satisfy the exact limiting square
\[
 E_{\cal A}E_{\cal B}=E_{\cal N}
 =E_{\cal B}E_{\cal A}.                                       \tag{26.15}
\]
\end{theorem}

\begin{proof}
V17 makes each \(E_R\) normal UCP.  If \(X\in R\), then
\(E_{R,n}X=X\) for every \(n\), so point-ultraweak convergence gives
\(E_RX=X\).  For arbitrary \(X\), every \(E_{R,n}X\) lies in the
ultraweakly closed algebra \(R\), hence so does \(E_RX\).
Therefore \(E_R\) is a norm-one projection onto \(R\).  By the
conditional-expectation theorem for norm-one projections onto a
\(C^*\)-subalgebra, it is \(R\)-bimodular and is a conditional
expectation.

Take preadjoints in (26.14).  Lemma
\ref{lem:v13v26-composition} passes both complete compositions to
the limit, while (26.10) passes the middle expectation.  Equality
of the limiting preadjoints gives (26.15).
\end{proof}

The theorem is not an application of joint ultraweak continuity of
multiplication.  It uses norm convergence of the predual
compositions, which is exactly the topology constructed in
V17--V19.

\subsection{Vanishing-residual version}

Let \((\rho_\ell)\) be trace-norm dense in the normal state space.
Suppose, along the common cofinal sequence,
\[
\begin{split}
 \|S_{{\cal B},n}S_{{\cal A},n}\rho_\ell
      -S_{{\cal N},n}\rho_\ell\|
 &\leq\varepsilon_{n,\ell}^{AB},\\
 \|S_{{\cal A},n}S_{{\cal B},n}\rho_\ell
      -S_{{\cal N},n}\rho_\ell\|
 &\leq\varepsilon_{n,\ell}^{BA},                              \tag{26.16}
\end{split}
\]
where both residuals tend to zero for every fixed \(\ell\).

\begin{corollary}[Dense-ledger closure]
\label{cor:v13v26-residual}
Under (26.10) and (26.16), the limiting expectations satisfy
(26.15).
\end{corollary}

\begin{proof}
Lemma~\ref{lem:v13v26-composition} and (26.10) pass (26.16) to zero
on every \(\rho_\ell\).  All limiting preadjoint maps are
contractions.  Approximation of an arbitrary normal state by the
dense sequence extends the equality to all states, and linearity
extends it to the entire predual.  Taking adjoints proves (26.15).
\end{proof}

A fixed nonzero residual gives only an approximate square.  The
compactness theorem does not remove it.

\subsection{Intertwining the channel with causal conditioning}

Use an output square
\(E_R^{\rm out}\) and an input square \(E_R^{\rm in}\), with
\(R\in\{{\cal A},{\cal B},{\cal N}\}\), on corresponding ambient
regional algebras.  A normal UCP channel
\(\Phi:{\cal M}^{\rm out}\to{\cal M}^{\rm in}\) respects the
conditional structure when
\[
 E_R^{\rm in}\Phi=\Phi E_R^{\rm out}
 \qquad(R={\cal A},{\cal B},{\cal N}).                          \tag{26.17}
\]

\begin{theorem}[Closure of the conditioned-channel square]
\label{thm:v13v26-channel}
Suppose the cutoff channels and all six expectations have
preadjoints converging pointwise in norm, and suppose (26.17) holds
at every cutoff.  Then (26.17) holds for the limiting locally normal
channel and limiting expectations.  If both cutoff triples obey
commuting squares, both limiting triples do as well.
\end{theorem}

\begin{proof}
Taking preadjoints in the cutoff form of (26.17) gives equality of
two typed compositions of CPTP contractions.  The two-factor
estimate (26.12), applied with the channel preadjoint as either
factor, passes both sides in predual norm.  Taking adjoints yields
(26.17).  The last claim is Theorem~\ref{thm:v13v26-limit}.
\end{proof}

Equation (26.17) says that conditioning before or after the channel
has the same result.  It is a no-extra-information condition adapted
to the declared causal collar.  It does not assert that the global
state is a tensor product.

\subsection{Why separate averaging remains unsafe}

Even if each cutoff triple satisfies (26.14), separately averaging
\(E_{\cal A}\) and \(E_{\cal B}\) and then composing them need not
equal the average of \(E_{\cal A}E_{\cal B}\).  The reason is the
same bilinearity exposed by the V24 qubit example:
\[
 \left(\int T_gE_{\cal A}\,dg\right)
 \left(\int T_hE_{\cal B}\,dh\right)
 \ne
 \int T_g(E_{\cal A}E_{\cal B})\,dg                            \tag{26.18}
\]
in general.

There are two safe routes:

\begin{enumerate}
\item construct state-preserving expectations that are individually
equivariant by uniqueness, then verify their square;
\item include the three expectations and both complete composition
equalities in the joint finite Choi spectrahedron, and select an
equivariant feasible packet without averaging the factors
independently.
\end{enumerate}

The second route is not convex if idempotence or composition
equalities are imposed directly.  A convex Choi selection theorem
such as V12--V16 cannot be invoked until a convex parametrization or
a fixed-expectation formulation has been supplied.

\subsection{Energy-tight acquisition route}

Each normal conditional expectation is a normal UCP map, so its
preadjoint is CPTP.  If the six regional expectation preadjoints and
the channel preadjoint satisfy V17 energy-transfer bounds and V18
background equicontinuity, the V19 countable regional diagonal can
be applied to their finite tuple.  A finite tuple is handled by the
sum of its predual metrics.  Theorems
\ref{thm:v13v26-limit} and \ref{thm:v13v26-channel} then pass the
complete causal packet to one background-continuous locally normal
limit.

\begin{assessment}[Status after thirteenth-cycle V26]
\label{ass:v13v26-status}
A state-preserving conditional-expectation commuting square gives a
correlation-retaining replacement for spacelike product
factorization.  Exact or vanishing-residual squares are closed under
the energy-tight predual convergence of V17--V19.  The same typed
composition argument passes channel-conditioning intertwiners.
Unique state-preserving expectations provide a symmetry-compatible
route that avoids the separate-averaging obstruction.
\end{assessment}

\begin{firewall}[Claim boundary after thirteenth-cycle V26]
\label{fire:v13v26-claim}
No RPESM collar algebras, faithful normal state, modularly admissible
conditional expectations, commuting square, channel intertwiner, or
uniform energy bound for those expectations has been constructed.
Conditional expectations onto type-III local subalgebras need not
exist without additional modular or split hypotheses.  The square
is not the perturbative causal-factorization axiom for time-ordered
products and does not prove clustering or the full SM--Einstein
continuum.
\end{firewall}

\subsection{Next bottleneck}

The immediate gate is existence, not closure.  The next note should
use a faithful normal candidate state and modular invariance to
state a precise conditional-expectation existence criterion, then
identify a finite relative-entropy or recovery-map residual whose
vanishing implies the commuting square.  This goes upstream of
trying to select idempotent Choi matrices inside a nonconvex set.

\subsection{Parent record}

This V26 note appends to the validated thirteenth-cycle V25 file with
record
\[
\begin{array}{c|c}
\text{lines}&7745\\
\text{bytes}&274520\\
\text{SHA--256}&
\texttt{FDFC3770BB35006782285793A0956EAA2F4DD8F74EB8E6AF9A58347B0FB57F7A}.
\end{array}
\]

% =============================================================================
% END APPENDED THIRTEENTH-CYCLE V26 MATERIAL
% =============================================================================

% =============================================================================
% APPENDED THIRTEENTH-CYCLE V27 MATERIAL
% =============================================================================

\section{Thirteenth-cycle V27: modular conditional expectations and an exact square residual}
\label{sec:v13v27}

\subsection{The existence gate left by V26}

V26 proved that a causal conditional-expectation square is closed
under the predual convergence developed in V17--V19.  It did not
construct the expectations.  This note supplies two acquisition
tools:

\begin{enumerate}
\item modular invariance is the exact existence criterion for a
faithful-state-preserving normal conditional expectation;
\item once three such expectations exist, their commuting-square
defect is the norm of one cross-Gram matrix on the centered
\(L^2\)-subspaces.
\end{enumerate}

The second item is an actual-composition certificate.  It measures
the full collar-conditioned relation and does not separately norm
the two expectations.

\subsection{The modular existence theorem}

Let \({\cal M}\) be a von Neumann algebra, let \(\varphi\) be a
faithful normal state, and let \({\cal R}\subset{\cal M}\) be a
von Neumann subalgebra.  Denote the modular automorphism group by
\(\sigma_t^\varphi\).

\begin{theorem}[Faithful-state conditional-expectation criterion]
\label{thm:v13v27-takesaki}
There exists a normal conditional expectation
\[
 E_{\cal R}:{\cal M}\longrightarrow{\cal R},
 \qquad \varphi E_{\cal R}=\varphi,                            \tag{27.1}
\]
if and only if
\[
 \sigma_t^\varphi({\cal R})={\cal R}
 \qquad(t\in\mathbb R).                                       \tag{27.2}
\]
When it exists, the \(\varphi\)-preserving expectation is unique.
\end{theorem}

\begin{proof}
We recall the modular projection argument.  In the GNS standard
form \(({\cal H}_\varphi,\Omega_\varphi)\), let \(e_{\cal R}\) be
the orthogonal projection onto
\(\overline{{\cal R}\Omega_\varphi}\).

If \(E_{\cal R}\) exists, its \(L^2\)-implementation is
\(e_{\cal R}\):
\[
 e_{\cal R}(X\Omega_\varphi)=E_{\cal R}(X)\Omega_\varphi.
                                                                    \tag{27.3}
\]
State preservation and bimodularity make this an orthogonal
projection.  The modular operator preserves
\(\overline{{\cal R}\Omega_\varphi}\), so its modular unitaries
commute with \(e_{\cal R}\).  Therefore
\(\Delta_\varphi^{it}{\cal R}\Delta_\varphi^{-it}={\cal R}\),
which is (27.2).

Conversely, (27.2) makes
\(\overline{{\cal R}\Omega_\varphi}\) invariant under both the
modular unitaries and modular conjugation.  The modular projection
construction defines \(E_{\cal R}(X)\) by
\[
 E_{\cal R}(X)\Omega_\varphi
 =
 e_{\cal R}X\Omega_\varphi.                                   \tag{27.4}
\]
Commutation with the standard-form right action shows that the
right side is represented by a bounded element of \({\cal R}\) and
that
\(\|E_{\cal R}(X)\|\leq\|X\|\).  Applying (27.4) at every matrix
level gives complete positivity; \(e_{\cal R}\Omega_\varphi
=\Omega_\varphi\) gives unitality.  The construction is
\({\cal R}\)-bimodular, normal by monotone convergence in standard
form, fixes \({\cal R}\), and satisfies
\[
 \varphi(E_{\cal R}X)
 =
 \langle\Omega_\varphi,e_{\cal R}X\Omega_\varphi\rangle
 =
 \varphi(X).                                                   \tag{27.5}
\]
For uniqueness, if \(F\) is another such expectation, then for
\(R\in{\cal R}\),
\[
 \varphi(R^*F(X))
 =\varphi(F(R^*X))
 =\varphi(R^*X).                                               \tag{27.6}
\]
The same identity holds for \(E_{\cal R}\).  Density of
\({\cal R}\Omega_\varphi\) in its GNS subspace and faithfulness give
\(F(X)=E_{\cal R}(X)\).
\end{proof}

Thus the existence ledger is a modular invariance ledger.  A generic
UCP projection selected by a Choi feasibility problem need not
preserve the candidate physical state.

\subsection{Explicit finite-dimensional formula}

Let \({\cal M}={\cal B}({\cal H})\) and
\(\varphi(X)=\operatorname{Tr}(\rho X)\), where \(\rho>0\) and
\(\operatorname{Tr}\rho=1\).  After a unitary change of basis, any
finite-dimensional von Neumann subalgebra has the form
\[
 {\cal R}
 =
 \bigoplus_{\alpha}
 \bigl({\cal B}({\cal H}_\alpha)\otimes I_{{\cal K}_\alpha}\bigr),
 \qquad
 {\cal H}=\bigoplus_\alpha
 {\cal H}_\alpha\otimes{\cal K}_\alpha.                        \tag{27.7}
\]

\begin{proposition}[Block formula for the modular expectation]
\label{prop:v13v27-block}
The modular invariance condition (27.2) holds if and only if
\[
 \rho=
 \bigoplus_\alpha
 p_\alpha\,\rho_\alpha\otimes\kappa_\alpha,                    \tag{27.8}
\]
where \(p_\alpha>0\), \(\sum_\alpha p_\alpha=1\), and
\(\rho_\alpha,\kappa_\alpha\) are faithful density matrices.
In that case the unique state-preserving expectation is
\[
 E_{\cal R}(X)
 =
 \bigoplus_\alpha
 \left[
 \operatorname{Tr}_{{\cal K}_\alpha}
 \bigl(\kappa_\alpha X_{\alpha\alpha}\bigr)
 \otimes I_{{\cal K}_\alpha}
 \right].                                                      \tag{27.9}
\]
\end{proposition}

\begin{proof}
Modular invariance means that \(\rho^{it}\) normalizes \({\cal R}\).
The central block projections in (27.7) form a finite discrete set,
so the continuous one-parameter group fixes each of them.  On one
block, every one-parameter automorphism of
\({\cal B}({\cal H}_\alpha)\otimes I\) is inner.  Hence the generator
\(\log\rho\) on that block differs from an element of
\({\cal B}({\cal H}_\alpha)\otimes I\) by an element of the
commutant \(I\otimes{\cal B}({\cal K}_\alpha)\).  Exponentiation and
normalization give (27.8).  Conversely, (27.8) visibly normalizes
each algebra block.

The map (27.9) is unital, completely positive, fixes \({\cal R}\),
and is \({\cal R}\)-bimodular.  Direct calculation gives
\[
\begin{split}
 \operatorname{Tr}\!\left[
 (p_\alpha\rho_\alpha\otimes\kappa_\alpha)
 E_{\cal R}(X)_{\alpha\alpha}\right]
 &=
 p_\alpha\operatorname{Tr}\!\left[
 \rho_\alpha
 \operatorname{Tr}_{{\cal K}_\alpha}
 (\kappa_\alpha X_{\alpha\alpha})\right]\\
 &=
 \operatorname{Tr}\!\left[
 (p_\alpha\rho_\alpha\otimes\kappa_\alpha)
 X_{\alpha\alpha}\right].                                     \tag{27.10}
\end{split}
\]
Summing over \(\alpha\) proves state preservation.  Uniqueness is
Theorem~\ref{thm:v13v27-takesaki}.
\end{proof}

\subsection{Background continuity on a fixed block type}

Suppose the decomposition (27.7) is fixed over a background chart
and \(x\mapsto\kappa_{\alpha,x}\) is trace-norm continuous.  Formula
(27.9) defines \(E_{{\cal R},x}\).

\begin{proposition}[Quantitative expectation continuity]
\label{prop:v13v27-continuity}
For two backgrounds in the chart,
\[
 \|E_{{\cal R},x}-E_{{\cal R},y}\|
 \leq
 \max_\alpha
 \|\kappa_{\alpha,x}-\kappa_{\alpha,y}\|_1.                    \tag{27.11}
\]
Consequently the preadjoints are pointwise trace-norm continuous,
uniformly on the unit ball.
\end{proposition}

\begin{proof}
For \(\|X\|\leq1\), trace duality for the weighted partial trace gives
\[
 \left\|
 \operatorname{Tr}_{\cal K}
 ((\kappa_x-\kappa_y)X)
 \right\|
 \leq\|\kappa_x-\kappa_y\|_1.                                 \tag{27.12}
\]
Take the maximum over the direct-sum blocks.  Equality of the norm
of a map and its preadjoint gives the last assertion.
\end{proof}

A change of block multiplicity is a modular orbit-type crossing and
cannot be hidden inside (27.11).  It requires the projection and
cohomology-jump treatment already isolated in V3.

\subsection{The exact \(L^2\) square defect}

Return to
\({\cal N}\subset{\cal A}\cap{\cal B}\subset{\cal M}\) and a faithful
normal state \(\varphi\).  Assume all three subalgebras are modularly
invariant, so Theorem~\ref{thm:v13v27-takesaki} supplies their unique
\(\varphi\)-preserving expectations.

Let \(P_{\cal R}\) denote the orthogonal projection onto
\(L^2({\cal R},\varphi)\).  Since
\(P_{\cal N}\leq P_{\cal A},P_{\cal B}\), define
\[
 Q_{\cal A}=P_{\cal A}-P_{\cal N},\qquad
 Q_{\cal B}=P_{\cal B}-P_{\cal N}.                             \tag{27.13}
\]
These are the projections onto the centered subspaces
\[
 {\cal A}_0=L^2({\cal A},\varphi)\ominus L^2({\cal N},\varphi),
 \qquad
 {\cal B}_0=L^2({\cal B},\varphi)\ominus L^2({\cal N},\varphi).
                                                                    \tag{27.14}
\]

\begin{theorem}[Cross-Gram certificate for the commuting square]
\label{thm:v13v27-gram}
The following are equivalent:
\begin{enumerate}
\item \(E_{\cal A}E_{\cal B}=E_{\cal N}\);
\item \(E_{\cal B}E_{\cal A}=E_{\cal N}\);
\item \(Q_{\cal A}Q_{\cal B}=0\);
\item \({\cal A}_0\perp{\cal B}_0\) in \(L^2({\cal M},\varphi)\).
\end{enumerate}
Moreover,
\[
 \|P_{\cal A}P_{\cal B}-P_{\cal N}\|
 =
 \|Q_{\cal A}Q_{\cal B}\|.                                   \tag{27.15}
\]
If the centered spaces are finite dimensional with
\(\varphi\)-orthonormal bases \((a_i)\) and \((b_j)\), and
\[
 G_{ij}=\varphi(a_i^*b_j),                                    \tag{27.16}
\]
then
\[
 \boxed{
 \|P_{\cal A}P_{\cal B}-P_{\cal N}\|=\|G\|_{\rm op}.}          \tag{27.17}
\]
\end{theorem}

\begin{proof}
The expectations are implemented by \(P_{\cal A},P_{\cal B}\), and
\(P_{\cal N}\).  Using (27.13) and the orthogonality of the centered
spaces to \(L^2({\cal N},\varphi)\),
\[
 P_{\cal A}P_{\cal B}
 =
 (P_{\cal N}+Q_{\cal A})(P_{\cal N}+Q_{\cal B})
 =
 P_{\cal N}+Q_{\cal A}Q_{\cal B}.                             \tag{27.18}
\]
Thus the first and third conditions are equivalent and (27.15)
holds.  Taking adjoints exchanges \({\cal A}\) and \({\cal B}\), so
the second condition is equivalent as well.  Two closed subspaces
are orthogonal precisely when the product of their orthogonal
projections is zero, proving equivalence with the fourth condition.

In the declared orthonormal bases, the restriction
\(Q_{\cal A}|_{{\cal B}_0}:{\cal B}_0\to{\cal A}_0\) has matrix
\(G\).  The operator \(Q_{\cal A}Q_{\cal B}\) vanishes on
\({\cal B}_0^\perp\), so its norm equals the matrix operator norm
of \(G\).  This proves (27.17).
\end{proof}

The correct finite residual is therefore the cross-Gram norm after
subtracting the collar algebra.  Separate upper bounds on
\(Q_{\cal A}\) and \(Q_{\cal B}\) are both equal to one and contain
no useful conditional-independence information.

\subsection{A finite Choi existence spectrahedron}

At a finite ledger stage with fixed target algebra \({\cal R}\),
the conditions
\[
 E\ \hbox{is CP and unital},\qquad
 E({\cal M})\subset{\cal R},\qquad
 E|_{\cal R}={\rm id},\qquad
 \varphi E=\varphi                                             \tag{27.19}
\]
are Choi positivity plus affine linear equations.  They define a
compact spectrahedron.  Idempotence need not be imposed as a
quadratic equation: the range and fixed-point equations imply
\[
 E^2=E.                                                        \tag{27.20}
\]
When the modular condition holds and \(\varphi\) is faithful, the
spectrahedron is the singleton containing (27.9).

The square equation remains bilinear between two expectations, but
Theorem~\ref{thm:v13v27-gram} replaces it by the vanishing of an
explicit cross-Gram block once the unique modular expectations have
been constructed.

\subsection{Energy transfer as a Heisenberg Loewner gate}

Let \(K\geq0\) be the local compact-resolvent energy and let
\(S_{\cal R}=(E_{\cal R})_*\).  For a positive normal functional
\(\rho\),
\[
 \operatorname{Tr}(K S_{\cal R}\rho)
 =
 \rho(E_{\cal R}(K)),                                         \tag{27.21}
\]
with both sides defined by increasing bounded spectral
truncations.

\begin{proposition}[Energy certificate for a modular expectation]
\label{prop:v13v27-energy}
If the extended-positive operator inequality
\[
 E_{\cal R}(K)\leq aK+bI                                     \tag{27.22}
\]
holds in quadratic-form sense, then
\[
 \operatorname{Tr}(K S_{\cal R}\rho)
 \leq
 a\operatorname{Tr}(K\rho)+b\operatorname{Tr}\rho             \tag{27.23}
\]
for every positive finite-energy \(\rho\).
\end{proposition}

\begin{proof}
Apply \(\rho\) to bounded spectral truncations of (27.22), use
(27.21), and pass to the increasing limit by monotone convergence.
\end{proof}

Thus the conditional expectations can enter the V17--V19 compactness
theorems once the direct Loewner gate (27.22) is certified.  This is
again an actual-channel estimate: it bounds the expected energy
after conditioning, not unrelated blocks of a Stinespring factor.

\begin{assessment}[Status after thirteenth-cycle V27]
\label{ass:v13v27-status}
The existence problem for the V26 expectations is reduced exactly
to modular invariance.  On a finite block ledger the unique
state-preserving expectation has the explicit weighted partial-trace
formula (27.9), varies continuously with its complementary density,
and belongs to an affine Choi spectrahedron.  The full
commuting-square defect is the operator norm of the centered
cross-Gram matrix.  A direct Heisenberg Loewner inequality supplies
the energy-transfer input needed for local normal compactness.
\end{assessment}

\begin{firewall}[Claim boundary after thirteenth-cycle V27]
\label{fire:v13v27-claim}
No faithful RPESM local state, modular automorphism group, invariant
past-collar or collar-future algebra, finite cross-Gram matrix, or
energy inequality (27.22) has been populated.  Modular invariance
can fail for interacting local subalgebras, and block type can jump
over the gravitational background.  The theorem does not prove a
quantum Markov property, recovery equality, clustering, or the full
SM--Einstein continuum.
\end{firewall}

\subsection{Next bottleneck}

The cross-Gram residual is exact but expressed in
\(L^2(\varphi)\).  The next note should relate it to a state-level
recovery statement.  On a finite causal ledger, construct the Petz
recovery map for restriction to the collar, prove exact recovery
from vanishing conditional mutual information, and state a
quantitative trace-norm recovery bound.  This must be kept distinct
from merely small pair correlations.

\subsection{Parent record}

This V27 note appends to the validated thirteenth-cycle V26 file with
record
\[
\begin{array}{c|c}
\text{lines}&8127\\
\text{bytes}&288229\\
\text{SHA--256}&
\texttt{4C4BF2D7A80239EBC997A5ECA816CC3B57AED3CA3339F2C28C35EBEDDCB88164}.
\end{array}
\]

% =============================================================================
% END APPENDED THIRTEENTH-CYCLE V27 MATERIAL
% =============================================================================

% =============================================================================
% APPENDED THIRTEENTH-CYCLE V28 MATERIAL
% =============================================================================

\section{Thirteenth-cycle V28: Petz recovery and a cofinal causal-state residual}
\label{sec:v13v28}

\subsection{Finite causal ledger}

At one finite ledger stage, let
\[
 {\cal H}={\cal H}_A\otimes{\cal H}_B\otimes{\cal H}_C          \tag{28.1}
\]
represent a past region \(A\), an intermediate causal collar \(B\),
and a future region \(C\).  Let \(\rho_{ABC}\) be a faithful density
matrix and write its marginals in the usual way.  With natural
logarithms, define
\[
 I(A:C\mid B)_\rho
 =
 S(\rho_{AB})+S(\rho_{BC})-S(\rho_B)-S(\rho_{ABC}).             \tag{28.2}
\]
This is a state-level Markov defect.  It is distinct from the
algebra-level cross-Gram defect in V27 unless an explicit
identification between the two ledgers has been proved.

\subsection{Two standard finite-dimensional recovery inputs}

For a channel \({\cal N}\) and faithful density \(\sigma\), the Petz
map is
\[
 {\cal R}_{\sigma,{\cal N}}(X)
 =
 \sigma^{1/2}{\cal N}^*\!\left(
 ({\cal N}\sigma)^{-1/2}X({\cal N}\sigma)^{-1/2}
 \right)\sigma^{1/2}.                                         \tag{28.3}
\]
We use the following standard equality and strengthened-monotonicity
theorems.

\begin{theorem}[Finite-dimensional Petz equality theorem]
\label{thm:v13v28-petz-input}
For faithful densities \(\rho,\sigma\) and a finite-dimensional
CPTP map \({\cal N}\),
\[
 D(\rho\Vert\sigma)=D({\cal N}\rho\Vert{\cal N}\sigma)          \tag{28.4}
\]
if and only if
\[
 {\cal R}_{\sigma,{\cal N}}({\cal N}\rho)=\rho.                \tag{28.5}
\]
The Petz map recovers \(\sigma\) from \({\cal N}\sigma\) as well.
\end{theorem}

\begin{proof}
Complete positivity of (28.3) is immediate from its three CP
factors.  The adjoint identity
\[
 {\cal R}_{\sigma,{\cal N}}^*(I)
 =
 ({\cal N}\sigma)^{-1/2}
 {\cal N}(\sigma)
 ({\cal N}\sigma)^{-1/2}=I                                    \tag{28.6}
\]
on the faithful support gives trace preservation.  Substitution of
\({\cal N}\sigma\) gives recovery of \(\sigma\).

The remaining equivalence is the equality case in the
finite-dimensional monotonicity proof for relative entropy.  In
that proof, the contraction between the two relative modular
operators becomes an isometry on the cyclic vector exactly when
\[
 \log\rho-\log\sigma
 =
 {\cal N}^*(\log{\cal N}\rho-\log{\cal N}\sigma)               \tag{28.7}
\]
on the generated support.  Exponentiating the corresponding
relative modular relation yields (28.5).  Conversely, apply data
processing first to \({\cal N}\) and then to the recovery map:
\[
 D(\rho\Vert\sigma)
 \geq D({\cal N}\rho\Vert{\cal N}\sigma)
 \geq D(\rho\Vert\sigma).                                     \tag{28.8}
\]
Both inequalities are equalities.
\end{proof}

For density matrices \(\rho,\tau\), use the root fidelity
\[
 F(\rho,\tau)=\|\rho^{1/2}\tau^{1/2}\|_1.                      \tag{28.9}
\]

\begin{theorem}[Finite-dimensional quantitative recovery input]
\label{thm:v13v28-recovery-input}
For every tripartite state \(\rho_{ABC}\), there is a CPTP recovery
map \({\cal R}_{B\to BC}\), depending only on \(\rho_{BC}\), such
that
\[
 I(A:C\mid B)_\rho
 \geq
 -2\log F\!\left(
 \rho_{ABC},
 ({\rm id}_A\otimes{\cal R}_{B\to BC})(\rho_{AB})
 \right).                                                      \tag{28.10}
\]
Consequently
\[
 \left\|
 \rho_{ABC}
 -({\rm id}_A\otimes{\cal R}_{B\to BC})(\rho_{AB})
 \right\|_1
 \leq2\sqrt{1-e^{-I(A:C\mid B)_\rho}}
 \leq2\sqrt{I(A:C\mid B)_\rho}.                               \tag{28.11}
\]
\end{theorem}

\begin{proof}
The first inequality is the finite-dimensional strengthened
monotonicity theorem, whose recovery channel can be chosen as an
average of rotated Petz maps.  It gives
\(F\geq e^{-I/2}\).  The upper Fuchs--van de Graaf inequality gives
\[
 {1\over2}\|\rho-\tau\|_1\leq\sqrt{1-F(\rho,\tau)^2}.           \tag{28.12}
\]
Substitute the fidelity lower bound and use
\(1-e^{-u}\leq u\) for \(u\geq0\).
\end{proof}

The deep input in Theorem~\ref{thm:v13v28-recovery-input} is stated
rather than reproved.  All uses below retain its finite-dimensional
scope and its exact constants.

\subsection{The explicit collar Petz map}

Apply Theorem~\ref{thm:v13v28-petz-input} with
\[
 {\cal N}=\operatorname{Tr}_C,\qquad
 \sigma=\rho_A\otimes\rho_{BC}.                               \tag{28.13}
\]
Then
\[
 {\cal N}\rho_{ABC}=\rho_{AB},\qquad
 {\cal N}\sigma=\rho_A\otimes\rho_B,                           \tag{28.14}
\]
and the relative-entropy drop is exactly
\[
 D(\rho_{ABC}\Vert\rho_A\otimes\rho_{BC})
 -
 D(\rho_{AB}\Vert\rho_A\otimes\rho_B)
 =
 I(A:C\mid B)_\rho.                                           \tag{28.15}
\]

Define
\[
 {\cal R}^{\rm P}_{B\to BC}(X)
 =
 \rho_{BC}^{1/2}
 \left(
 \rho_B^{-1/2}X\rho_B^{-1/2}\otimes I_C
 \right)
 \rho_{BC}^{1/2}.                                              \tag{28.16}
\]

\begin{theorem}[Exact causal recovery criterion]
\label{thm:v13v28-exact}
For faithful \(\rho_{ABC}\), the following are equivalent:
\begin{enumerate}
\item \(I(A:C\mid B)_\rho=0\);
\item
\[
 \rho_{ABC}
 =
 ({\rm id}_A\otimes{\cal R}^{\rm P}_{B\to BC})(\rho_{AB});
                                                                    \tag{28.17}
\]
\item there exists a CPTP map
\({\cal R}_{B\to BC}\) such that
\[
 \rho_{ABC}
 =
 ({\rm id}_A\otimes{\cal R}_{B\to BC})(\rho_{AB}).             \tag{28.18}
\]
\end{enumerate}
\end{theorem}

\begin{proof}
Equivalence of the first two statements follows from
(28.13)--(28.16) and
Theorem~\ref{thm:v13v28-petz-input}.  The second implies the third.

Assume the third.  Taking the \(A\)-marginal of (28.18) gives
\[
 {\cal R}_{B\to BC}(\rho_B)=\rho_{BC}.                         \tag{28.19}
\]
Thus \({\rm id}_A\otimes{\cal R}\) recovers both members of the
relative-entropy pair in (28.15):
\[
 \rho_{AB}\mapsto\rho_{ABC},\qquad
 \rho_A\otimes\rho_B\mapsto\rho_A\otimes\rho_{BC}.              \tag{28.20}
\]
Data processing under \(\operatorname{Tr}_C\), followed by data
processing under the recovery channel, forces equality in (28.15).
Hence the conditional mutual information vanishes.
\end{proof}

Exact recovery does not force a product state.  The recovered
future may remain correlated with the past through the collar
degrees of freedom.

\subsection{Continuity of the finite Petz channel}

Uniform spectral floors are needed before (28.16) is used over a
background family.  Suppose
\[
 \rho_{BC},\widetilde\rho_{BC}\geq m_{BC}I,\qquad
 \rho_B,\widetilde\rho_B\geq m_BI                              \tag{28.21}
\]
with \(m_{BC},m_B>0\), and let \(d_C=\dim{\cal H}_C\).

\begin{proposition}[A finite Petz Lipschitz bound]
\label{prop:v13v28-petz-lipschitz}
The induced trace-norm operator bound satisfies
\[
 \|{\cal R}^{\rm P}_\rho-{\cal R}^{\rm P}_{\widetilde\rho}\|_{1\to1}
 \leq
 {d_C\over m_B\sqrt{m_{BC}}}
 \|\rho_{BC}-\widetilde\rho_{BC}\|_1
 +
 {d_C\over m_B^2}
 \|\rho_B-\widetilde\rho_B\|_1.                               \tag{28.22}
\]
\end{proposition}

\begin{proof}
On positive matrices bounded below by \(mI\), functional calculus
gives
\[
 \|X^{1/2}-Y^{1/2}\|
 \leq{1\over2\sqrt m}\|X-Y\|,\qquad
 \|X^{-1/2}-Y^{-1/2}\|
 \leq{1\over2m^{3/2}}\|X-Y\|.                                 \tag{28.23}
\]
Write (28.16) as \(L_{\rho_{BC}}M_{\rho_B}\), where
\[
 L_A(Z)=A^{1/2}ZA^{1/2},\qquad
 M_B(X)=(B^{-1/2}XB^{-1/2})\otimes I_C.                        \tag{28.24}
\]
Since density matrices have operator norm at most one,
\[
 \|L_A-L_{\widetilde A}\|_{1\to1}
 \leq {1\over\sqrt{m_{BC}}}\|A-\widetilde A\|_1,               \tag{28.25}
\]
while
\[
 \|M_B\|_{1\to1}\leq {d_C\over m_B},\qquad
 \|M_B-M_{\widetilde B}\|_{1\to1}
 \leq {d_C\over m_B^2}\|B-\widetilde B\|_1.                    \tag{28.26}
\]
Insert and subtract \(L_{\widetilde A}M_B\), use
\(\|L_{\widetilde A}\|_{1\to1}\leq1\), and combine
(28.25)--(28.26).
\end{proof}

The dimension and inverse-floor factors in (28.22) can diverge with
the cutoff.  The estimate is a finite-ledger continuity theorem,
not a uniform continuum theorem.

\subsection{A recovery residual that passes to the limit}

Entropy is not generally continuous in an infinite-dimensional or
growing-cutoff limit.  The trace-norm recovery residual is more
robust.  Let
\[
 r_n
 =
 \left\|
 \rho_{ABC,n}
 -({\rm id}_A\otimes{\cal R}_n)(\rho_{AB,n})
 \right\|_1.                                                   \tag{28.27}
\]

\begin{theorem}[Cofinal closure of exact recovery]
\label{thm:v13v28-cofinal}
Assume, on every fixed finite causal ledger,
\[
 \rho_{ABC,n}\to\rho_{ABC},\qquad
 \rho_{AB,n}\to\rho_{AB}                                      \tag{28.28}
\]
in trace norm,
\[
 {\cal R}_n(X)\to{\cal R}(X)
 \quad\text{in trace norm for every trace-class }X,             \tag{28.29}
\]
and \(r_n\to0\).  Then
\[
 \rho_{ABC}
 =
 ({\rm id}_A\otimes{\cal R})(\rho_{AB}).                       \tag{28.30}
\]
If the limiting ledger is finite dimensional, then
\[
 I(A:C\mid B)_\rho=0.                                         \tag{28.31}
\]
\end{theorem}

\begin{proof}
CPTP maps are trace-norm contractions.  Therefore
\[
\begin{split}
 &\left\|
 \rho_{ABC}
 -({\rm id}_A\otimes{\cal R})(\rho_{AB})
 \right\|_1\\
 &\quad\leq
 \|\rho_{ABC}-\rho_{ABC,n}\|_1+r_n
 +\|({\rm id}_A\otimes{\cal R}_n)
       (\rho_{AB,n}-\rho_{AB})\|_1\\
 &\qquad+
 \|({\rm id}_A\otimes({\cal R}_n-{\cal R}))
       (\rho_{AB})\|_1.                                       \tag{28.32}
\end{split}
\]
Every term tends to zero by the hypotheses.  This proves (28.30).
The finite-dimensional implication (28.31) is
Theorem~\ref{thm:v13v28-exact}.
\end{proof}

For a background family, uniform versions of (28.28)--(28.29) and
\(r_n\to0\) give uniform recovery.  For a regional exhaustion, the
countable diagonal of V19 applies if the recovery maps possess
uniform energy transfer and background equicontinuity.

\subsection{Energy gate for recovery maps}

Let \(K_B\) and \(K_{BC}\) be the input and output local energies of
a recovery channel.  Its Heisenberg adjoint
\(({\cal R}_{B\to BC})^*\) is normal UCP.  The direct Loewner
condition
\[
 ({\cal R}_{B\to BC})^*(K_{BC})
 \leq aK_B+bI                                                  \tag{28.33}
\]
implies
\[
 \operatorname{Tr}\!\left(
 K_{BC}{\cal R}_{B\to BC}(\eta)\right)
 \leq
 a\operatorname{Tr}(K_B\eta)+b\operatorname{Tr}\eta.           \tag{28.34}
\]
The proof is the same trace-duality and monotone-convergence
argument as Proposition~\ref{prop:v13v27-energy}.  Equation (28.33),
not finite-dimensional compactness by itself, is the required
cutoff-uniform input for V17.

\subsection{Relation to the V27 square}

On the finite tensor ledger, the V27 cross-Gram condition and the
Petz condition both express conditional independence over the
collar, but in different categories:

\begin{enumerate}
\item the cross-Gram condition is an equality between
state-preserving conditional expectations on an ambient algebra;
\item the Petz condition is an equality for one density matrix under
a recovery channel;
\item identifying them requires the modular expectations to exist
and their collar restriction/recovery maps to coincide with the
declared tensor partial trace and Petz channel.
\end{enumerate}

Without that identification, small pair correlations do not imply
small conditional mutual information, and small conditional mutual
information does not construct the V27 expectations on a type-III
continuum algebra.

\begin{assessment}[Status after thirteenth-cycle V28]
\label{ass:v13v28-status}
At every faithful finite causal ledger, zero conditional mutual
information is equivalent to exact recovery by the explicit collar
Petz map.  Strengthened monotonicity supplies a trace-distance
recovery error bounded by \(2\sqrt{I(A:C\mid B)}\).  Under spectral
floors the Petz channel has an explicit background Lipschitz bound.
Most importantly for cutoff removal, a vanishing trace-norm recovery
residual passes through pointwise trace-norm convergence without
requiring entropy continuity.
\end{assessment}

\begin{firewall}[Claim boundary after thirteenth-cycle V28]
\label{fire:v13v28-claim}
No RPESM tripartite causal factorization, faithful density, spectral
floor, conditional-mutual-information estimate, Petz channel,
recovery residual, or recovery energy inequality has been
populated.  The quantitative recovery theorem is a
finite-dimensional standard input.  Its recovery maps can lose
uniformity as dimensions grow.  Exact Markov recovery does not by
itself prove Einstein causality, the time-slice axiom, reflection
positivity, clustering, or the full SM--Einstein continuum.
\end{firewall}

\subsection{Next bottleneck}

The remaining causal-state problem is to replace finite density
matrices by local normal states on type-III algebras.  The next note
should formulate recovery through relative modular operators or
standard-form channels, and identify an energy-constrained compact
class in which the finite recovery maps admit locally normal
subsequential limits.  If that route becomes circular, the correct
upstream fallback is the trace-norm residual theorem (28.30) on an
increasing split inclusion.

\subsection{Parent record}

This V28 note appends to the validated thirteenth-cycle V27 file with
record
\[
\begin{array}{c|c}
\text{lines}&8534\\
\text{bytes}&302752\\
\text{SHA--256}&
\texttt{2DD1C636FB7309AACF91E442C88FA1DD565C083C1E2310F2EC1296D0B70EBC74}.
\end{array}
\]

% =============================================================================
% END APPENDED THIRTEENTH-CYCLE V28 MATERIAL
% =============================================================================

% =============================================================================
% APPENDED THIRTEENTH-CYCLE V29 MATERIAL
% =============================================================================

\section{Thirteenth-cycle V29: locally normal recovery on split type-III screens}
\label{sec:v13v29}

\subsection{Why the argument moves upstream}

V28 gives exact and quantitative recovery on finite tensor factors.
Local algebras in continuum quantum field theory are generally
type III, so they have no density matrix relative to a canonical
trace and their individual von Neumann entropies need not exist.
It would be circular to assume finite-dimensional entropy continuity
and call the resulting map a continuum recovery channel.

The stable datum is instead the predual-norm recovery residual.
This note extends the V17 compactness proof to arbitrary von
Neumann algebras with separable preduals, then uses a declared split
isomorphism to pass that residual.  A finite-rank phase-space
compression is the abstract replacement for the compact-resolvent
energy projection.

\subsection{Normal-channel compactness for separable preduals}

Let \({\cal M}\) and \({\cal N}\) be von Neumann algebras with
separable preduals.  Let
\[
 \Psi_n:{\cal N}\longrightarrow{\cal M}                        \tag{29.1}
\]
be normal UCP maps and
\[
 R_n=(\Psi_n)_*:{\cal M}_*\longrightarrow{\cal N}_*            \tag{29.2}
\]
their CPTP preadjoints.

\begin{theorem}[Abstract predual compactness]
\label{thm:v13v29-compact}
Let \((d_\ell)_{\ell\geq1}\) be a norm-dense sequence in the unit
ball of \({\cal M}_*\).  Assume that, for every \(\ell\),
\[
 \{R_n d_\ell:n\geq1\}
 \quad\text{is relatively compact in }{\cal N}_*.              \tag{29.3}
\]
Then a subsequence \((n_k)\) and a CPTP contraction
\[
 R:{\cal M}_*\longrightarrow{\cal N}_*                         \tag{29.4}
\]
exist such that
\[
 \|R_{n_k}\rho-R\rho\|\longrightarrow0
 \quad(\rho\in{\cal M}_*).                                     \tag{29.5}
\]
Its adjoint
\[
 \Psi=R^*:{\cal N}\longrightarrow{\cal M}                      \tag{29.6}
\]
is normal UCP, and \(\Psi_{n_k}(A)\to\Psi(A)\) ultraweakly for
every \(A\in{\cal N}\).
\end{theorem}

\begin{proof}
Enumerate the dense tests and successively extract convergent
subsequences using (29.3).  The Cantor diagonal subsequence makes
\(R_{n_k}d_\ell\) norm convergent for every \(\ell\).

Every \(R_n\) is a contraction.  If \(\rho\in{\cal M}_*\), choose
\(d_\ell\) close to \(\rho\) and use
\[
 \|R_{n_k}\rho-R_{n_j}\rho\|
 \leq2\|\rho-d_\ell\|
   +\|R_{n_k}d_\ell-R_{n_j}d_\ell\|.                           \tag{29.7}
\]
Thus \(R_{n_k}\rho\) is Cauchy.  Its limit defines the linear
contraction \(R\).  Positivity and preservation of the normal
functional value at the identity pass by norm convergence.

The Banach adjoint \(R^*\) is normal because \(R\) is its
preadjoint, and it is unital because \(R\) preserves the value at
the identity.  To prove complete positivity, let
\([A_{ij}]\geq0\) in \(M_m({\cal N})\).  Entrywise ultraweak
convergence gives
\[
 [\Psi_{n_k}(A_{ij})]\longrightarrow[\Psi(A_{ij})]             \tag{29.8}
\]
ultraweakly in \(M_m({\cal M})\).  The positive cone is ultraweakly
closed, so the limit is positive.  This holds for every \(m\).
Finally,
\[
 \rho(\Psi_{n_k}(A)-\Psi(A))
 =
 (R_{n_k}\rho-R\rho)(A)\longrightarrow0,                       \tag{29.9}
\]
which proves the asserted convergence.
\end{proof}

This theorem isolates exactly what the trace-class proof of V17
used: pointwise predual relative compactness, not a global trace on
the von Neumann algebra.

\subsection{A finite-rank phase-space sufficient condition}

The compactness condition (29.3) can be certified by finite-rank
predual compressions.

\begin{lemma}[Uniform finite-rank compression implies tightness]
\label{lem:v13v29-compression}
Fix \(\rho\in{\cal M}_*\).  Suppose that for every
\(\varepsilon>0\) there is a bounded finite-rank map
\[
 Q_{\rho,\varepsilon}:{\cal N}_*\longrightarrow{\cal N}_*      \tag{29.10}
\]
such that
\[
 \sup_n
 \|R_n\rho-Q_{\rho,\varepsilon}R_n\rho\|
 <\varepsilon                                                  \tag{29.11}
\]
and
\[
 \sup_{\rho,\varepsilon}\|Q_{\rho,\varepsilon}\|<\infty
\quad\text{on each fixed test orbit}.                          \tag{29.12}
\]
Then \(\{R_n\rho\}_n\) is relatively compact.
\end{lemma}

\begin{proof}
The compressed orbit is bounded in the finite-dimensional range of
\(Q_{\rho,\varepsilon}\), hence totally bounded.  Equation (29.11)
places the original orbit in its \(\varepsilon\)-neighborhood.
Since this holds for every \(\varepsilon\), the original orbit is
totally bounded.  Its closure in the complete predual is compact.
\end{proof}

For a type-I representation, the spectral compression in
Lemma~\ref{lem:v13v17-tight} supplies (29.10)--(29.11).  For a
type-III local algebra, a phase-space or nuclearity argument must
supply such finite-rank approximants on the actual recovery-output
orbits.  Naming nuclearity without producing (29.11) is not a
compactness proof.

\subsection{Tensor amplification of predual convergence}

Let \({\cal P}\) be another von Neumann algebra with separable
predual.  The spatial tensor predual is the norm completion of
finite sums of product normal functionals.

\begin{lemma}[Stable amplification]
\label{lem:v13v29-amplification}
If \(R_n\rho\to R\rho\) in norm for every
\(\rho\in{\cal M}_*\), then
\[
 ({\rm id}_{{\cal P}_*}\otimes R_n)(\eta)
 \longrightarrow
 ({\rm id}_{{\cal P}_*}\otimes R)(\eta)                        \tag{29.13}
\]
in predual norm for every
\(\eta\in({\cal P}\overline\otimes{\cal M})_*\).
\end{lemma}

\begin{proof}
For an elementary product functional \(\alpha\otimes\rho\),
\[
 \|(\alpha\otimes R_n\rho)-(\alpha\otimes R\rho)\|
 =
 \|\alpha\|\,\|R_n\rho-R\rho\|\longrightarrow0.                \tag{29.14}
\]
The same holds for finite sums.  Such sums are norm dense in the
spatial tensor predual.  All amplified maps are contractions, so a
three-term approximation extends convergence to arbitrary
\(\eta\).
\end{proof}

\subsection{Split causal screens}

Let \({\cal M}_A,{\cal M}_B,{\cal M}_{BC}\) be local von Neumann
algebras with separable preduals.  Assume a split causal screen has
provided normal spatial identifications
\[
 {\cal M}_{AB}\simeq
 {\cal M}_A\overline\otimes{\cal M}_B,\qquad
 {\cal M}_{ABC}\simeq
 {\cal M}_A\overline\otimes{\cal M}_{BC}.                      \tag{29.15}
\]
No claim is made that arbitrary adjacent type-III regions admit
(29.15); it is a split/nuclearity input.

Let
\[
 {\cal R}_n:({\cal M}_B)_*\longrightarrow({\cal M}_{BC})_*     \tag{29.16}
\]
be CPTP recovery maps with normal UCP adjoints.  Let
\(\varphi_n\in({\cal M}_{ABC})_*\) be normal states and
\(\varphi_{AB,n}\) their restrictions.  Define the recovery residual
\[
 r_n=
 \left\|
 \varphi_n-
 ({\rm id}_{({\cal M}_A)_*}\otimes{\cal R}_n)
 \varphi_{AB,n}
 \right\|.                                                     \tag{29.17}
\]

\begin{theorem}[Locally normal recovery on a split screen]
\label{thm:v13v29-recovery}
Assume:
\begin{enumerate}
\item \(\varphi_n\to\varphi\) in the predual norm of
\({\cal M}_{ABC}\);
\item the recovery preadjoints satisfy the compactness hypothesis
(29.3), hence along a subsequence
\({\cal R}_{n_k}\to{\cal R}\) pointwise in predual norm;
\item \(r_{n_k}\to0\).
\end{enumerate}
Then \({\cal R}\) is CPTP with a normal UCP adjoint and
\[
 \boxed{
 \varphi=
 ({\rm id}_{({\cal M}_A)_*}\otimes{\cal R})\varphi_{AB}.}      \tag{29.18}
\]
\end{theorem}

\begin{proof}
Restriction of normal functionals is contractive, so
\[
 \|\varphi_{AB,n_k}-\varphi_{AB}\|
 \leq\|\varphi_{n_k}-\varphi\|\longrightarrow0.                \tag{29.19}
\]
Theorem~\ref{thm:v13v29-compact} gives the normal limiting recovery
channel.  Lemma~\ref{lem:v13v29-amplification} gives convergence of
its amplification on the fixed input \(\varphi_{AB}\).  Therefore
\[
\begin{split}
 &\left\|
 \varphi-
 ({\rm id}\otimes{\cal R})\varphi_{AB}
 \right\|\\
 &\quad\leq
 \|\varphi-\varphi_{n_k}\|+r_{n_k}
 +\|({\rm id}\otimes{\cal R}_{n_k})
      (\varphi_{AB,n_k}-\varphi_{AB})\|\\
 &\qquad+
 \|({\rm id}\otimes({\cal R}_{n_k}-{\cal R}))
      \varphi_{AB}\|.                                         \tag{29.20}
\end{split}
\]
All four terms tend to zero.  This proves (29.18).
\end{proof}

Equation (29.18) is a type-III-compatible local Markov statement:
it uses only normal states, normal channels, split tensor products,
and predual norms.  No individual von Neumann entropy appears.

\subsection{Background-uniform recovery}

Let \(X\) again be the compact Einstein-background core.  For each
dense input functional \(d_\ell\), assume the recovery-output family
\[
 \{{\cal R}_{n,x}d_\ell:n\geq1,\ x\in X\}                      \tag{29.21}
\]
is relatively compact in the output predual and has a common
background modulus
\[
 \|{\cal R}_{n,x}d_\ell-{\cal R}_{n,y}d_\ell\|
 \leq\omega_\ell(d_X(x,y)).                                   \tag{29.22}
\]
Use the metric
\[
 d_{\rm rec}(R,T)
 =
 \sum_{\ell\geq1}2^{-\ell}
 \min\{1,\|Rd_\ell-Td_\ell\|\}.                               \tag{29.23}
\]

\begin{corollary}[One recovery field over all backgrounds]
\label{cor:v13v29-background}
Under (29.21)--(29.22), a subsequence independent of \(x\) converges
uniformly in \(d_{\rm rec}\) to a background-continuous field of
locally normal recovery channels.  If the states and residuals in
Theorem~\ref{thm:v13v29-recovery} converge uniformly in \(x\), then
(29.18) holds for every background, uniformly in predual norm.
\end{corollary}

\begin{proof}
The proof of Lemma~\ref{lem:v13v18-complete} works for separable
preduals verbatim.  Pointwise compactness follows from
Theorem~\ref{thm:v13v29-compact}, and the metric-valued
Arzel\`a--Ascoli proof of V18 supplies one uniform subsequence.
Apply estimate (29.20) uniformly in \(x\).
\end{proof}

A countable cofinal regional exhaustion can be included by the
nested diagonal of V19.  Each region may have its own phase-space
compression card; no global finite-volume energy bound is required.

\subsection{Intertwining the locally normal channel}

Let the V19 observable channel have preadjoint
\(S_x:({\cal M}^{\rm in})_*\to({\cal M}^{\rm out})_*\), and suppose
input and output causal screens have recovery maps
\({\cal R}_x^{\rm in}\) and \({\cal R}_x^{\rm out}\) with compatible
types.  The recovery intertwiner is the complete predual relation
\[
 S_x({\rm id}\otimes{\cal R}_x^{\rm in})
 =
 ({\rm id}\otimes{\cal R}_x^{\rm out})S_x^{AB}.                \tag{29.24}
\]
If (29.24) holds at every cutoff and all occurring maps converge
pointwise in predual norm, repeated use of the two-factor estimate
(26.12) passes it to the limit.  This is the recovery analogue of
the conditioned-channel square in V26.

\subsection{What nuclearity must actually deliver}

For each fixed local screen and dense input \(d_\ell\), a usable
phase-space certificate consists of finite-rank maps
\(Q_{\ell,m}\) and explicit tails
\[
 \sup_{n,x}
 \|{\cal R}_{n,x}d_\ell-Q_{\ell,m}{\cal R}_{n,x}d_\ell\|
 \leq\tau_{\ell,m},
 \qquad \tau_{\ell,m}\downarrow0.                              \tag{29.25}
\]
Together with (29.22), this supplies background-uniform relative
compactness.  A nuclearity index or split inclusion is relevant only
after it has been converted into (29.25) for the actual recovery
outputs.

\begin{assessment}[Status after thirteenth-cycle V29]
\label{ass:v13v29-status}
Normal-channel compactness and recovery closure now apply to
arbitrary von Neumann algebras with separable preduals.  On a
declared split causal screen, finite-cutoff recovery maps with
pointwise predual compactness and vanishing state residual produce
a locally normal type-III recovery channel.  Finite-rank
phase-space compression is an explicit sufficient compactness
certificate, and background and regional diagonals preserve the
construction.
\end{assessment}

\begin{firewall}[Claim boundary after thirteenth-cycle V29]
\label{fire:v13v29-claim}
No RPESM split causal screen, local nuclearity bound, finite-rank
phase-space compression, type-III recovery channel, recovery
residual, or intertwiner (29.24) has been populated.  The split
property cannot be inferred from separability, locality, or an
energy gap alone.  Equation (29.18) is a state-specific Markov
property and does not prove perturbative causal factorization,
reflection positivity, clustering, or the full SM--Einstein
continuum.
\end{firewall}

\subsection{Next bottleneck}

V30 is the next compulsory cross-program audit.  After it, the
programme should attack the phase-space compression card (29.25).
The most direct route is an energy-damped local map with a summable
singular-value or nuclear decomposition.  Its tail must control the
actual recovery outputs uniformly over the compact Einstein
background core.

\subsection{Parent record}

This V29 note appends to the validated thirteenth-cycle V28 file with
record
\[
\begin{array}{c|c}
\text{lines}&8962\\
\text{bytes}&316510\\
\text{SHA--256}&
\texttt{87CE8CFC4A8AEAD917964A11BE88CF57A2CFF57EB0CFD4AF8EFB896CC85A2CFF}.
\end{array}
\]

% =============================================================================
% END APPENDED THIRTEENTH-CYCLE V29 MATERIAL
% =============================================================================

% =============================================================================
% APPENDED THIRTEENTH-CYCLE V30 MATERIAL
% =============================================================================

\section{Thirteenth-cycle V30: compulsory companion audit and the energy-dressed compactness decision}
\label{sec:v13v30}

\subsection{Read-only companion snapshot}

The compulsory five-version audit inspected the newest companion
files.  The folder contains Yang--Mills V78 and V79 with identical
content:
\[
\begin{array}{c|c|c|c}
\text{file}&\text{lines}&\text{bytes}&\text{SHA--256}\\ \hline
{\rm YM\ V78}&27139&924440&
\texttt{5186A400C54B788054377499FBA689FFF}\\
&&&\texttt{01F682E111FFE0F939CED3E0D1ACAF7}\\
{\rm YM\ V79}&27139&924440&
\texttt{5186A400C54B788054377499FBA689FFF}\\
&&&\texttt{01F682E111FFE0F939CED3E0D1ACAF7}\\
{\rm NS\ V26}&5319&278283&
\texttt{82C887F8C2C69E4F28FAD7C3DC8DE5B}\\
&&&\texttt{9099CB5A4BF431EBA1FFF3E91935978AD}.
\end{array}                                                     \tag{30.1}
\]
Thus V79 contains no mathematical update beyond V78.  The effective
new Yang--Mills material since the V25 audit is V77--V78.  The
Navier--Stokes file is unchanged.  Companion acquisition directions
were read as source text and were not treated as instructions.

\subsection{Yang--Mills V77: energy-dressed complete words}

YM V77 rewrites covariance and vertex factors in the common Hessian
energy metric before applying Schatten inequalities.  In a complete
trace word, this absorbs inverse factors into dimensionless
normalized covariance jets and measures vertices on their actual
input and output fibres.  The resulting identities are exact and
avoid an additional absolute inverse-floor factor.

The reported numerical values in V77 are diagnostics rather than
global certificates.  No YM number is transferred.  The structural
point applies to V29: compactness should be proved for the complete
map carrying the actual recovery-output orbit, not by separately
bounding a nuclear analysis map and an unrelated synthesis factor
before their cancellations or range restrictions are used.

\subsection{Yang--Mills V78: certified anchors and relative transport}

YM V78 gives rational residual certificates for 288
energy-dressed vertex norms at sixteen dyadic corners.  It then
proves the noncommutative relative transport inequality
\[
 \|H^{-1/2}XH^{-1/2}\|_F
 \leq {1\over1-\rho}
 \left(
 \|H_0^{-1/2}X_0H_0^{-1/2}\|_F+
 \|H_0^{-1/2}\Delta XH_0^{-1/2}\|_F
 \right),                                                      \tag{30.2}
\]
when the relative Hessian perturbation has norm \(\rho<1\).
The global covering stocks remain open in the YM programme.

For the SM recovery problem, the corresponding useful architecture
is:
\[
 \text{nuclear or compact anchor for the complete damped map}
 \ \Longrightarrow\
 \text{relative operator-norm transport over a background chart}
 \ \Longrightarrow\
 \text{finite-rank orbit tails}.                              \tag{30.3}
\]
The actual RPESM map and perturbation radius still have to be
constructed.

\subsection{Navier--Stokes V26}

There is no new NS content or digest.  Its rank-one actual-input
calculation remains consistent with using the actual recovery orbit
in V29, but supplies no phase-space map, nuclear norm, singular-value
tail, or background perturbation estimate.

\begin{decision}[V30 cross-program decision]
\label{dec:v13v30}
V31 will prove a compactness card for recovery outputs which:
\begin{enumerate}
\item factors each actual output orbit through one complete nuclear
or compact operator;
\item obtains explicit finite-rank tails from a nuclear series or
approximation numbers;
\item proves collective compactness for a norm-continuous compact
operator field over the compact Einstein-background core;
\item permits finitely many anchor maps and relative transport
bounds, following the architecture of (30.2);
\item states precisely what a local phase-space or nuclearity
estimate must provide before V29 can be applied.
\end{enumerate}
No companion physical constant is imported.
\end{decision}

\begin{assessment}[Status after thirteenth-cycle V30]
\label{ass:v13v30-status}
The mandatory audit is complete.  V77--V78 support an
energy-dressed, complete-map compactness route and warn against
premature factorwise estimates.  The duplicate V79 and unchanged
NS V26 add no result.  V31 will target the finite-rank
phase-space-compression hypothesis left open in V29.
\end{assessment}

\begin{firewall}[Claim boundary after thirteenth-cycle V30]
\label{fire:v13v30-claim}
No RPESM energy-damped local map, nuclear decomposition,
approximation-number tail, or relative background transport bound
has been populated.  The YM dressed anchors concern a different
finite lattice Hessian and do not prove recovery compactness,
reflection positivity, clustering, or the full SM--Einstein
continuum.
\end{firewall}

\subsection{Parent record}

This V30 note appends to the validated thirteenth-cycle V29 file with
record
\[
\begin{array}{c|c}
\text{lines}&9340\\
\text{bytes}&329888\\
\text{SHA--256}&
\texttt{AAB33D8695CA1D1E036A6532E414B57266E8EAA73CDA1FAC9DBFECB621C22D47}.
\end{array}
\]

% =============================================================================
% END APPENDED THIRTEENTH-CYCLE V30 MATERIAL
% =============================================================================

% =============================================================================
% APPENDED THIRTEENTH-CYCLE V31 MATERIAL
% =============================================================================

\section{Thirteenth-cycle V31: nuclear recovery-output compression over the background core}
\label{sec:v13v31}

\subsection{The precise target}

V29 reduces locally normal type-III recovery to relative compactness
of the actual predual output orbits.  This note gives three
sufficient certificates for that hypothesis:

\begin{enumerate}
\item factorization through one compact operator;
\item a nuclear or singular-value expansion with an explicit
finite-rank tail;
\item finitely many compact anchor maps with operator-norm transport
over the background core.
\end{enumerate}

All statements concern the complete map carrying the recovery
output.  A compact factor occurring before an uncontrolled synthesis
map is not enough.

\subsection{Nuclear tails}

Let \(E,F\) be complex Banach spaces.  Suppose
\(T:E\to F\) has a nuclear representation
\[
 Tz=\sum_{j=1}^{\infty}\lambda_j f_j(z)y_j,\qquad
 \sum_{j=1}^{\infty}
 |\lambda_j|\,\|f_j\|\,\|y_j\|<\infty.                         \tag{31.1}
\]
Define
\[
 T_mz=\sum_{j=1}^{m}\lambda_j f_j(z)y_j,\qquad
 \nu_m(T)=\sum_{j>m}|\lambda_j|\,\|f_j\|\,\|y_j\|.              \tag{31.2}
\]

\begin{lemma}[Explicit nuclear compression]
\label{lem:v13v31-nuclear}
The map \(T_m\) has rank at most \(m\) and
\[
 \|T-T_m\|_{E\to F}\leq\nu_m(T)\longrightarrow0.               \tag{31.3}
\]
Consequently, for every \(C<\infty\),
\[
 \sup_{\|z\|_E\leq C}
 \|Tz-T_mz\|_F\leq C\nu_m(T).                                 \tag{31.4}
\]
The image \(T(CB_E)\) is relatively compact in \(F\).
\end{lemma}

\begin{proof}
For \(\|z\|\leq1\), the tail of (31.1) is bounded by
\[
 \sum_{j>m}|\lambda_j|\,|f_j(z)|\,\|y_j\|
 \leq\nu_m(T).                                                 \tag{31.5}
\]
This proves (31.3)--(31.4).  The bounded image
\(T_m(CB_E)\) lies in the finite-dimensional span of
\(y_1,\ldots,y_m\) and is therefore totally bounded.  Equation
(31.4) approximates \(T(CB_E)\) uniformly by such sets, proving
total boundedness and relative compactness.
\end{proof}

A quantitative covering bound is also available.  If
\(C\nu_m(T)\leq\varepsilon/2\), then
\(T(CB_E)\) has an \(\varepsilon\)-net with at most
\[
 \left(
 1+{4C\|T_m\|\over\varepsilon}
 \right)^{2m}                                                  \tag{31.6}
\]
points.  Indeed, the radius-\(C\|T_m\|\) ball in a complex
\(m\)-dimensional normed space has an \(\varepsilon/2\)-net with
the displayed volumetric bound, and (31.4) supplies the remaining
\(\varepsilon/2\).

\subsection{Hilbert singular-value form}

When \(E,F\) are Hilbert spaces and \(T\) is compact, let
\(s_1(T)\geq s_2(T)\geq\cdots\downarrow0\) be its singular values.

\begin{corollary}[Optimal Hilbert truncation]
\label{cor:v13v31-singular}
The rank-\(m\) singular truncation \(T_m\) obeys
\[
 \|T-T_m\|=s_{m+1}(T).                                        \tag{31.7}
\]
If \(T\) is Hilbert--Schmidt or trace class, then also
\[
 \|T-T_m\|_{\rm HS}
 =
 \left(\sum_{j>m}s_j(T)^2\right)^{1/2},\qquad
 \|T-T_m\|_1=\sum_{j>m}s_j(T).                                \tag{31.8}
\]
\end{corollary}

\begin{proof}
Use the singular-value decomposition.  The operator norm of the
tail is its largest remaining singular value, and orthogonality of
the singular vectors gives the two ideal-norm identities.
\end{proof}

The operator-norm tail in (31.7), rather than the full trace-class
tail, is sufficient for compactness of a bounded input orbit.

\subsection{Collective compactness over a background field}

Let \(X\) be compact and let
\[
 x\longmapsto T_x\in{\cal K}(E,F)                              \tag{31.9}
\]
be continuous in operator norm, where \({\cal K}(E,F)\) denotes the
compact operators.

\begin{theorem}[Collective compactness of a norm-continuous field]
\label{thm:v13v31-collective}
For every \(C<\infty\),
\[
 {\cal K}_{X,C}
 :=
 \{T_xz:x\in X,\ \|z\|_E\leq C\}                              \tag{31.10}
\]
is relatively compact in \(F\).
\end{theorem}

\begin{proof}
Fix \(\varepsilon>0\).  Operator-norm continuity and compactness of
\(X\) give points \(x_1,\ldots,x_q\) such that every \(x\) has an
\(i\) with
\[
 \|T_x-T_{x_i}\|<{\varepsilon\over2C},                         \tag{31.11}
\]
with the case \(C=0\) trivial.  Since each \(T_{x_i}\) is compact,
the set \(T_{x_i}(CB_E)\) has a finite
\(\varepsilon/2\)-net.  The union of these finitely many nets is an
\(\varepsilon\)-net for (31.10), because
\[
 \|T_xz-T_{x_i}z\|<\varepsilon/2.                              \tag{31.12}
\]
Thus (31.10) is totally bounded.
\end{proof}

This theorem is stronger than pointwise compactness of every \(T_x\).
The operator-norm continuity over compact \(X\) makes the family
collectively compact.

\subsection{Affine factorization of the actual recovery orbit}

Let \(R_{n,x}:{\cal M}_*\to{\cal N}_*\) be the recovery
preadjoints of V29, and let \((d_\ell)\) be its dense input tests.
Suppose, for every \(\ell\), there are Banach spaces \(E_\ell\),
continuous anchor fields
\[
 \sigma_\ell:X\to{\cal N}_*,                                  \tag{31.13}
\]
norm-continuous compact-operator fields
\[
 T_{\ell,x}:E_\ell\to{\cal N}_*,                              \tag{31.14}
\]
and coefficients satisfying
\[
 R_{n,x}d_\ell
 =
 \sigma_\ell(x)+T_{\ell,x}z_{n,x,\ell},\qquad
 \|z_{n,x,\ell}\|\leq C_\ell.                                 \tag{31.15}
\]

\begin{theorem}[Recovery-orbit compactness card]
\label{thm:v13v31-orbit}
Under (31.13)--(31.15),
\[
 \{R_{n,x}d_\ell:n\geq1,\ x\in X\}                            \tag{31.16}
\]
is relatively compact in \({\cal N}_*\) for every \(\ell\).
Consequently the compactness hypothesis (29.21) is satisfied.
If the left side of (31.15) also obeys the background moduli
(29.22), V29 yields one background-continuous locally normal
recovery field.
\end{theorem}

\begin{proof}
The image \(\sigma_\ell(X)\) is compact by continuity and compactness
of \(X\).  Theorem~\ref{thm:v13v31-collective} makes
\[
 \{T_{\ell,x}z:x\in X,\ \|z\|\leq C_\ell\}
\]
relatively compact.  The sum of a compact set and a relatively
compact set is relatively compact.  Equation (31.15) places the
entire actual orbit inside that sum.  Apply V29.
\end{proof}

The factorization in (31.15) is allowed to depend on the test
functional \(\ell\), but its bound must be uniform in cutoff and
background.

\subsection{Energy-damped analysis and bounded synthesis}

A common phase-space architecture is
\[
 T_{\ell,x}=J_{\ell,x}\Theta_{\beta,\ell,x},                   \tag{31.17}
\]
where
\(\Theta_{\beta,\ell,x}:E_\ell\to{\cal H}_\ell\) is an
energy-damped nuclear analysis map and
\(J_{\ell,x}:{\cal H}_\ell\to{\cal N}_*\) is a bounded synthesis
map.

\begin{proposition}[Complete damped-factor tail]
\label{prop:v13v31-damped}
Assume:
\[
 \sup_{x\in X}\|J_{\ell,x}\|\leq J_\ell,                       \tag{31.18}
\]
both operator fields in (31.17) are norm continuous, and
\(\Theta_{\beta,\ell,x}\) has nuclear expansions whose truncations
\(\Theta_{\beta,\ell,x}^{(m)}\) satisfy the uniform tail
\[
 \sup_{x\in X}
 \|\Theta_{\beta,\ell,x}
      -\Theta_{\beta,\ell,x}^{(m)}\|
 \leq\tau_{\ell,m}\downarrow0.                                \tag{31.19}
\]
Then
\[
 T_{\ell,x}^{(m)}
 :=
 J_{\ell,x}\Theta_{\beta,\ell,x}^{(m)}                         \tag{31.20}
\]
has finite rank and
\[
 \sup_{\substack{x\in X\\\|z\|\leq C_\ell}}
 \|T_{\ell,x}z-T_{\ell,x}^{(m)}z\|
 \leq C_\ell J_\ell\tau_{\ell,m}.                             \tag{31.21}
\]
Hence (31.14) is collectively compact.
\end{proposition}

\begin{proof}
Composition with \(J_{\ell,x}\) does not increase rank and gives
\[
 \|J_{\ell,x}
   (\Theta_{\beta,\ell,x}-\Theta_{\beta,\ell,x}^{(m)})z\|
 \leq J_\ell\tau_{\ell,m}\|z\|.                               \tag{31.22}
\]
This proves (31.21).  Norm continuity of the product field follows
from the two norm-continuous factors and their uniform boundedness
on compact \(X\).  Each product is compact, so Theorem
\ref{thm:v13v31-collective} applies.
\end{proof}

The factorization must be proved for the complete map in (31.15).
A conventional nuclearity map such as
\(A\mapsto e^{-\beta H}A\Omega\) has the correct phase-space flavor
but the wrong conclusion until a bounded synthesis \(J\) and the
actual recovery-output identity have been established.

\subsection{Finite anchor and relative transport certificate}

A full nuclear expansion at every background is unnecessary.
Suppose \(X\) is covered by finitely many sets \(U_i\), with compact
anchor operators \(T_i:E\to F\), finite-rank approximants
\(T_i^{(m_i)}\), and bounds
\[
 \sup_{x\in U_i}\|T_x-T_i\|\leq\delta_i,\qquad
 \|T_i-T_i^{(m_i)}\|\leq\tau_i.                               \tag{31.23}
\]

\begin{proposition}[Anchor-relative finite-rank tail]
\label{prop:v13v31-anchor}
For \(x\in U_i\) and \(\|z\|\leq C\),
\[
 \|T_xz-T_i^{(m_i)}z\|
 \leq C(\delta_i+\tau_i).                                     \tag{31.24}
\]
If, for every \(\varepsilon>0\), the cover and ranks can be chosen
so that
\[
 \max_i C(\delta_i+\tau_i)<\varepsilon,                        \tag{31.25}
\]
then the family \(\{T_xz:x\in X,\|z\|\leq C\}\) is relatively
compact.
\end{proposition}

\begin{proof}
Equation (31.24) is the triangle inequality:
\[
 \|(T_x-T_i)z+(T_i-T_i^{(m_i)})z\|
 \leq C(\delta_i+\tau_i).                                     \tag{31.26}
\]
The finitely many approximating images lie in a finite union of
finite-dimensional bounded sets and are totally bounded.
Condition (31.25) supplies arbitrary uniform accuracy.
\end{proof}

This is the abstract analogue of the YM V78 anchor-relative
transport architecture.  No Yang--Mills Hessian or numerical radius
is used.

\subsection{Relative perturbation of a dressed compact map}

Let \(D_0,D_x\) be boundedly invertible positive dressing operators
and suppose
\[
 D_x=D_0^{1/2}(I+A_x)D_0^{1/2},\qquad
 \|A_x\|\leq\rho<1.                                           \tag{31.27}
\]
If a complete physical map has the dressed form
\[
 T_x=L_xD_x^{-1/2}X_xD_x^{-1/2}R_x,                           \tag{31.28}
\]
then
\[
 \|D_x^{-1/2}X_xD_x^{-1/2}\|
 \leq{1\over1-\rho}
 \|D_0^{-1/2}X_xD_0^{-1/2}\|                                 \tag{31.29}
\]
in any unitarily invariant Hilbert ideal norm for which the products
are defined.  Splitting \(X_x=X_0+\Delta X_x\) yields the anchor plus
actual-residual bound used in (30.2).  If the outer fields \(L_x,R_x\)
are uniformly bounded and norm continuous, (31.29) supplies an
operator-norm transport estimate \(\delta_i\) for (31.23).

The proof is the same congruence calculation as the relative energy
transport theorem audited at V30: the positive comparison implies
\[
 \|D_0^{1/2}D_x^{-1}D_0^{1/2}\|\leq(1-\rho)^{-1},              \tag{31.30}
\]
and ideal-norm multiplication gives (31.29).  The estimate must be
applied to the complete dressed map before separately bounding
\(L_x,X_x,R_x\).

\subsection{Insertion into the V29 recovery theorem}

For each region and dense normal input, a sufficient RPESM
phase-space card is now:

\begin{enumerate}
\item the exact affine orbit identity (31.15);
\item a nuclear tail (31.19), a singular-value tail (31.7), or
finite anchors (31.23)--(31.25);
\item a background modulus for the complete recovery output;
\item the recovery residual \(r_n\to0\);
\item the regional and symmetry intertwiners from V29;
\item when trace-class energy is available, the direct recovery
Loewner bound (28.33).
\end{enumerate}

These rows imply the locally normal background-continuous recovery
conclusion of V29.  None of them is a consequence of nuclearity
terminology alone.

\begin{assessment}[Status after thirteenth-cycle V31]
\label{ass:v13v31-status}
The abstract phase-space hypothesis of V29 now has explicit
sufficient certificates.  Nuclear and singular-value expansions
give finite-rank tails and covering bounds; norm-continuous compact
operator fields are collectively compact over the Einstein core;
and finitely many anchor maps plus relative transport give a
checkable uniform alternative.  An exact factorization of the
actual recovery-output orbit through the complete damped map then
produces the desired locally normal recovery subsequence.
\end{assessment}

\begin{firewall}[Claim boundary after thirteenth-cycle V31]
\label{fire:v13v31-claim}
No RPESM energy-damped analysis map, bounded synthesis, affine orbit
identity, nuclear series, singular-value tail, compact anchor, or
background transport radius has been populated.  Standard AQFT
nuclearity terminology does not by itself prove (31.15).
The result does not establish the split property, recovery residual,
reflection positivity, clustering, or the full SM--Einstein
continuum.
\end{firewall}

\subsection{Next bottleneck}

The compactness implication is now complete; the open problem is the
actual factorization (31.15).  The next note should construct it from
a local Stinespring representation of the recovery channel.  The
candidate compact map is the energy-damped environment coefficient
operator.  One must prove a uniform Schmidt or nuclear tail for the
environment vectors while retaining the collar algebra and gauge
constraints.

\subsection{Parent record}

This V31 note appends to the validated thirteenth-cycle V30 file with
record
\[
\begin{array}{c|c}
\text{lines}&9476\\
\text{bytes}&335290\\
\text{SHA--256}&
\texttt{367AABC32A3A67E808135F158F4DDC9E113B5341DC7D96DB498A9EAD15D40EDA}.
\end{array}
\]

% =============================================================================
% END APPENDED THIRTEENTH-CYCLE V31 MATERIAL
% =============================================================================

% =============================================================================
% APPENDED THIRTEENTH-CYCLE V32 MATERIAL
% =============================================================================

\section{Thirteenth-cycle V32: compact Stinespring ellipsoids and system-side recovery tails}
\label{sec:v13v32}

\subsection{Why a quadratic route is needed}

The affine nuclear factorization (31.15) is a strong sufficient
condition, but a Schr\"odinger channel output is naturally quadratic
in its Stinespring vector.  Truncating only the environment is not
enough: a unitary channel has a one-dimensional environment and can
still move a state through infinitely many orthogonal system modes.

The correct Stinespring compactness condition acts on the output
system.  If every purified recovery output lies in the image of
\(D\otimes I_{\rm env}\), with \(D\) compact on the system and a
uniform bound on the preimage, then the reduced system states have
uniform finite-rank tails.  The environment dimension never enters
the estimate.

\subsection{A system-side ellipsoid lemma}

Let \({\cal H}\) and \({\cal E}\) be separable Hilbert spaces, let
\(D:{\cal H}\to{\cal H}\) be compact, and let \(P_m\) be finite-rank
orthogonal projections such that
\[
 \|(I-P_m)D\|\longrightarrow0.                                \tag{32.1}
\]
For \(w\in{\cal H}\otimes{\cal E}\), put
\[
 \xi=(D\otimes I_{\cal E})w,\qquad
 \sigma=\operatorname{Tr}_{\cal E}|\xi\rangle\langle\xi|.      \tag{32.2}
\]

\begin{lemma}[Environment-independent system tail]
\label{lem:v13v32-system-tail}
If \(\|w\|\leq C\) and \(\operatorname{Tr}\sigma\leq t\), then
\[
 \operatorname{Tr}((I-P_m)\sigma)
 \leq C^2\|(I-P_m)D\|^2                                      \tag{32.3}
\]
and
\[
 \|\sigma-P_m\sigma P_m\|_1
 \leq
 2C\sqrt t\,\|(I-P_m)D\|.                                    \tag{32.4}
\]
\end{lemma}

\begin{proof}
Writing \(Q_m=I-P_m\),
\[
\begin{split}
 \operatorname{Tr}(Q_m\sigma)
 &=
 \|(Q_m\otimes I_{\cal E})\xi\|^2\\
 &=
 \|(Q_mD\otimes I_{\cal E})w\|^2
 \leq C^2\|Q_mD\|^2.                                         \tag{32.5}
\end{split}
\]
Apply the gentle compression estimate
\[
 \|\sigma-P_m\sigma P_m\|_1
 \leq
 2\sqrt{\operatorname{Tr}\sigma\,
          \operatorname{Tr}(Q_m\sigma)}                       \tag{32.6}
\]
and use (32.5).
\end{proof}

Only the system projection \(P_m\) appears in the compressed state.
The environment remains untruncated and may vary with the cutoff.

\subsection{Compactness of the reduced-state ellipsoid}

Let
\[
 {\cal S}_{D,C,t}
 =
 \left\{
 \operatorname{Tr}_{\cal E}
 |(D\otimes I_{\cal E})w\rangle
 \langle(D\otimes I_{\cal E})w|:
 {\cal E}\ {\rm arbitrary},\
 \|w\|\leq C,\
 \operatorname{Tr}\sigma\leq t
 \right\}.                                                     \tag{32.7}
\]

\begin{theorem}[Compact reduced Stinespring ellipsoid]
\label{thm:v13v32-ellipsoid}
The set \({\cal S}_{D,C,t}\) is relatively compact in the trace
class \({\mathfrak S}_1({\cal H})\).
\end{theorem}

\begin{proof}
For fixed \(m\), the compressed family
\(P_m{\cal S}_{D,C,t}P_m\) is bounded in the finite-dimensional
space \(P_m{\mathfrak S}_1({\cal H})P_m\), hence relatively compact.
Lemma~\ref{lem:v13v32-system-tail} approximates the full family
uniformly by this compressed family as \(m\to\infty\).  Therefore
the full family is totally bounded.
\end{proof}

This theorem is nonlinear in \(w\), as a state construction should
be.  It supplies the V29 orbit compactness directly and does not
force the linear affine form (31.15).

\subsection{Energy form of the ellipsoid}

Assume now that \(D>0\) is injective with dense range, so \(D^{-1}\)
may be unbounded.  For a positive trace-class \(\sigma\), define
\[
 {\cal E}_D(\sigma)
 :=
 \operatorname{Tr}(D^{-2}\sigma)                              \tag{32.8}
\]
by monotone spectral truncation.

\begin{proposition}[Energy moment equals the purification preimage]
\label{prop:v13v32-energy}
A positive trace-class \(\sigma\) with
\({\cal E}_D(\sigma)<\infty\) has a purification
\[
 \xi=(D\otimes I_{\overline{\cal H}})w                         \tag{32.9}
\]
such that
\[
 \|w\|^2={\cal E}_D(\sigma).                                  \tag{32.10}
\]
Consequently the family
\[
 \{\sigma\geq0:\operatorname{Tr}\sigma\leq t,\
 {\cal E}_D(\sigma)\leq C^2\}                                \tag{32.11}
\]
is relatively compact whenever \(D\) is compact.
\end{proposition}

\begin{proof}
Use the Hilbert--Schmidt vectorization of
\(\sigma^{1/2}:{\overline{\cal H}}\to{\cal H}\) as a purification
\(\xi\in{\cal H}\otimes\overline{\cal H}\).  The energy condition
says that \(D^{-1}\sigma^{1/2}\) is Hilbert--Schmidt.  Let \(w\) be
its vectorization.  Then (32.9) holds and
\[
 \|w\|^2
 =
 \|D^{-1}\sigma^{1/2}\|_{\rm HS}^2
 =
 \operatorname{Tr}(D^{-2}\sigma).                             \tag{32.12}
\]
Apply Theorem~\ref{thm:v13v32-ellipsoid}.
\end{proof}

Taking \(D=(I+K)^{-1/2}\) for a positive compact-resolvent
Hamiltonian \(K\) makes (32.8) the moment
\[
 {\cal E}_D(\sigma)=\operatorname{Tr}((I+K)\sigma).            \tag{32.13}
\]
Thus V17 is the type-I energy instance of the Stinespring ellipsoid.

\subsection{Singular-value error}

Let \(D\) have singular values
\(d_1\geq d_2\geq\cdots\downarrow0\), and let \(P_m\) project onto
its first \(m\) left singular vectors.

\begin{corollary}[Explicit system compression error]
\label{cor:v13v32-singular}
For the family (32.7),
\[
 \sup_{\sigma\in{\cal S}_{D,C,t}}
 \|\sigma-P_m\sigma P_m\|_1
 \leq2C\sqrt t\,d_{m+1}.                                     \tag{32.14}
\]
\end{corollary}

\begin{proof}
The singular-value truncation gives
\(\|(I-P_m)D\|=d_{m+1}\).  Apply (32.4).
\end{proof}

This is the precise Schmidt/energy tail requested at V31.  It
controls the reduced output even though the environment Schmidt
rank is unbounded.

\subsection{Application to Stinespring recovery channels}

Let
\[
 {\cal R}_{n,x}:{\mathfrak S}_1({\cal H}_{B,x})
 \longrightarrow{\mathfrak S}_1({\cal H}_{BC,x})              \tag{32.15}
\]
be CPTP recovery maps with Stinespring isometries
\[
 V_{n,x}:{\cal H}_{B,x}
 \longrightarrow{\cal H}_{BC,x}\otimes{\cal E}_{n,x}.         \tag{32.16}
\]
Fix a dense family of input states \(d_{\ell,x}\) and choose
purifications \(\zeta_{\ell,x}\).  The purified output is
\[
 \xi_{n,x,\ell}
 =
 (V_{n,x}\otimes I)\zeta_{\ell,x},                            \tag{32.17}
\]
where the input reference space is absorbed into the environment.

Assume, after local Hilbert trivialization, there are compact
operators \(D_{\ell,x}\) on the output system and vectors
\(w_{n,x,\ell}\) such that
\[
 \xi_{n,x,\ell}
 =
 (D_{\ell,x}\otimes I)w_{n,x,\ell},\qquad
 \|w_{n,x,\ell}\|\leq C_\ell.                                 \tag{32.18}
\]

\begin{theorem}[Stinespring recovery compactness card]
\label{thm:v13v32-recovery}
If, for every \(\ell\),
\[
 x\longmapsto D_{\ell,x}
\quad\text{is operator-norm continuous and compact-valued}     \tag{32.19}
\]
over compact \(X\), then
\[
 \{
 {\cal R}_{n,x}(d_{\ell,x}):
 n\geq1,\ x\in X
 \}                                                           \tag{32.20}
\]
is relatively compact in trace norm, provided the input
trivializations and \(x\mapsto d_{\ell,x}\) are continuous.
Hence the V29 background-uniform recovery compactness hypothesis
holds.
\end{theorem}

\begin{proof}
By Theorem~\ref{thm:v13v31-collective}, the union
\[
 \{D_{\ell,x}u:x\in X,\ \|u\|\leq1\}                          \tag{32.21}
\]
is relatively compact in the output Hilbert space.  Therefore, for
every \(\varepsilon>0\), a common finite-dimensional subspace with
projection \(P\) exists such that
\[
 \sup_{x\in X}\|(I-P)D_{\ell,x}\|<\varepsilon.                 \tag{32.22}
\]
Apply Lemma~\ref{lem:v13v32-system-tail} to (32.18), with
\(t=1\).  It uniformly approximates every state in (32.20) by a
density matrix in the fixed finite-dimensional range of \(P\).
That compressed family is bounded and relatively compact, proving
the claim.
\end{proof}

The statement is insensitive to unitary changes of Stinespring
environment: such changes alter \(w\) by an environment unitary but
leave its norm and the reduced output unchanged.

\subsection{A direct output-moment version}

The Stinespring vectors need not be exhibited if an actual output
moment is available.

\begin{corollary}[Actual recovery-energy criterion]
\label{cor:v13v32-output-energy}
Suppose \(D_{\ell,x}>0\) are compact and satisfy (32.19), and
\[
 \sup_{n,x}
 \operatorname{Tr}\!\left(
 D_{\ell,x}^{-2}{\cal R}_{n,x}(d_{\ell,x})
 \right)
 \leq C_\ell^2.                                               \tag{32.23}
\]
Then the orbit (32.20) is relatively compact.
\end{corollary}

\begin{proof}
Proposition~\ref{prop:v13v32-energy} supplies (32.18) for the
canonical purification of each output.  Apply
Theorem~\ref{thm:v13v32-recovery}.
\end{proof}

Equation (32.23) is often simpler to certify from the Heisenberg
Loewner bound (28.33).  It is still an actual-channel output
estimate, not an environment or Kraus norm taken in isolation.

\subsection{Restriction to a represented type-III algebra}

Let \({\cal N}\subset{\cal B}({\cal H})\) be a faithfully and
normally represented type-III algebra.  Suppose each recovery output
normal functional on \({\cal N}\) has a positive trace-class lift
\(\sigma_{n,x,\ell}\) on \({\cal B}({\cal H})\) satisfying the
uniform ellipsoid condition above.  Restriction is contractive:
\[
 \|\sigma|_{\cal N}-\tau|_{\cal N}\|_{{\cal N}_*}
 \leq\|\sigma-\tau\|_1.                                       \tag{32.24}
\]

\begin{proposition}[Represented type-III compactness]
\label{prop:v13v32-typeIII}
A relatively compact family of such trace-class lifts restricts to
a relatively compact family in \({\cal N}_*\).  Therefore the
ellipsoid criterion is a sufficient, representation-dependent
phase-space card for Theorem~\ref{thm:v13v29-compact}.
\end{proposition}

\begin{proof}
The restriction map from the trace class onto the represented
predual is continuous by (32.24).  A continuous image of a
relatively compact set is relatively compact.
\end{proof}

Existence of uniformly controlled positive lifts is additional data.
A normal functional always has representing trace-class operators
in a normal representation, but a chosen lift need not obey the
physical energy moment.

\subsection{Symmetry covariance of the ellipsoid}

Suppose a symmetry arrow \(\gamma:x\to y\) is implemented on output
systems by \(U_\gamma\) and
\[
 D_{\ell,y}=U_\gamma D_{\ell,x}U_\gamma^*.                     \tag{32.25}
\]
Then
\[
 \operatorname{Tr}\!\left(
 D_{\ell,y}^{-2}U_\gamma\sigma U_\gamma^*
 \right)
 =
 \operatorname{Tr}(D_{\ell,x}^{-2}\sigma).                    \tag{32.26}
\]
Thus the ellipsoid is transported exactly by the symmetry.  Convex
groupoid averaging of output states preserves the moment bound when
the dressing is covariant.  The locally normal equivariant
extraction of V21 and V29 can therefore use the same compactness
card.

\subsection{Relation to the V31 linear factorization}

The two sufficient routes are complementary:

\begin{enumerate}
\item V31 uses a linear compact map on coefficient variables and is
well suited to nuclear expansions of an entire predual map;
\item V32 uses a quadratic partial trace of compactly dressed
purification vectors and is well suited to Stinespring or energy
moments;
\item either route supplies the orbit compactness required by V29;
neither follows from the other without an additional factorization.
\end{enumerate}

\begin{assessment}[Status after thirteenth-cycle V32]
\label{ass:v13v32-status}
A compact system-side Stinespring ellipsoid gives a
dimension-independent environment estimate and an explicit
trace-norm compression tail.  Equivalently, a uniform moment of
\(D^{-2}\) on the actual recovery outputs produces trace-norm
compactness.  Norm-continuous compact dressing fields yield one
finite-dimensional approximation uniformly over the compact
Einstein-background core, and controlled trace-class lifts transfer
the result to represented type-III preduals.
\end{assessment}

\begin{firewall}[Claim boundary after thirteenth-cycle V32]
\label{fire:v13v32-claim}
No RPESM Stinespring isometry, compact system dressing, purified
factorization, output moment, positive trace-class lift, singular
tail, or symmetry-covariant background field has been populated.
A one-dimensional Stinespring environment alone gives no
compactness.  The result does not prove a recovery residual,
reflection positivity, clustering, or the full SM--Einstein
continuum.
\end{firewall}

\subsection{Next bottleneck}

The remaining compactness row is now a direct physical inequality:
construct a covariant compact dressing \(D_x\) and prove (32.23)
from the local recovery adjoint.  The next note should combine the
Heisenberg Loewner estimate (28.33) with a background-uniform
compact-resolvent comparison, then propagate its singular tail
through the regional and quotient transports.

\subsection{Parent record}

This V32 note appends to the validated thirteenth-cycle V31 file with
record
\[
\begin{array}{c|c}
\text{lines}&9881\\
\text{bytes}&349086\\
\text{SHA--256}&
\texttt{C0B3FC37D7E211287B463ECC065BD1A16E6842D4CE17B7B2D4901DE87EBD2F6D}.
\end{array}
\]

% =============================================================================
% END APPENDED THIRTEENTH-CYCLE V32 MATERIAL
% =============================================================================

% =============================================================================
% APPENDED THIRTEENTH-CYCLE V33 MATERIAL
% =============================================================================

\section{Thirteenth-cycle V33: from the recovery Loewner gate to a uniform compactness tail}
\label{sec:v13v33}

\subsection{The missing implication}

V28 identified the direct recovery-energy inequality
\[
 \Psi_{n,x}(K_{{\rm out},x})
 \leq aK_{{\rm in},x}+bI,                                    \tag{33.1}
\]
where \(\Psi_{n,x}={\cal R}_{n,x}^*\).  V32 showed that a uniform
moment of a compact dressing produces compact recovery-output
orbits.  This note joins the two statements and transports the
result over the Einstein-background core with one reference
spectral projection.

\subsection{Uniform form comparison}

Let \(X\) be compact.  On a fixed output Hilbert space
\({\cal H}_{\rm out}\), let \(K_x\geq0\) be self-adjoint with compact
resolvent and common form domain.  Fix \(x_0\in X\) and assume
\[
 c_-\bigl(I+K_{x_0}\bigr)
 \preceq I+K_x
 \preceq c_+\bigl(I+K_{x_0}\bigr)                             \tag{33.2}
\]
in quadratic-form sense, with \(0<c_-\leq c_+<\infty\).

Let
\[
 0\leq\lambda_1\leq\lambda_2\leq\cdots\uparrow\infty          \tag{33.3}
\]
be the eigenvalues of \(K_{x_0}\), repeated with multiplicity, and
let \(P_m\) be the projection onto the first \(m\) eigenvectors.

\begin{lemma}[Common reference spectral tail]
\label{lem:v13v33-reference-tail}
For every positive trace-class \(\sigma\) with
\(\operatorname{Tr}\sigma\leq t\),
\[
 \operatorname{Tr}((I-P_m)\sigma)
 \leq
 {c_-^{-1}\operatorname{Tr}((I+K_x)\sigma)
  \over1+\lambda_{m+1}}.                                      \tag{33.4}
\]
Consequently
\[
 \|\sigma-P_m\sigma P_m\|_1
 \leq
 2\sqrt{
 {t\,c_-^{-1}\operatorname{Tr}((I+K_x)\sigma)
  \over1+\lambda_{m+1}}
 }.                                                           \tag{33.5}
\]
\end{lemma}

\begin{proof}
The spectral theorem gives
\[
 I-P_m
 \preceq{I+K_{x_0}\over1+\lambda_{m+1}}.                       \tag{33.6}
\]
The lower inequality in (33.2) gives
\[
 I+K_{x_0}\preceq c_-^{-1}(I+K_x).                            \tag{33.7}
\]
Apply the positive normal functional represented by \(\sigma\) to
bounded spectral truncations of (33.6)--(33.7), then use monotone
convergence.  This proves (33.4).  The gentle compression estimate
(32.6) gives (33.5).
\end{proof}

The projection \(P_m\) is independent of \(x\).  This is stronger
than choosing unrelated spectral truncations at each background.

\subsection{How relative form perturbations supply (33.2)}

Let \(q_x\) be the closed quadratic form of \(K_x\).  If
\[
 |q_x[u]-q_{x_0}[u]|
 \leq\eta\bigl(q_{x_0}[u]+\|u\|^2\bigr),
 \qquad 0\leq\eta<1,                                          \tag{33.8}
\]
uniformly in \(x\), then
\[
 (1-\eta)(I+K_{x_0})
 \preceq I+K_x
 \preceq(1+\eta)(I+K_{x_0}).                                  \tag{33.9}
\]
Thus \(c_-=1-\eta\) and \(c_+=1+\eta\).  This is the same
common-domain relative-form mechanism proved in V4; the present
use is the output recovery energy.

If \(x\mapsto(I+K_x)^{-1}\) is norm continuous, then
\[
 D_x=(I+K_x)^{-1/2}                                           \tag{33.10}
\]
is a norm-continuous compact-operator field.  Indeed, every
resolvent is compact and
\[
 \|D_x-D_y\|
 \leq
 \|(I+K_x)^{-1}-(I+K_y)^{-1}\|^{1/2}                          \tag{33.11}
\]
by the operator square-root estimate for positive contractions.
Hence the V32 collective Stinespring-ellipsoid theorem applies.

\subsection{Energy transfer for the actual recovery output}

Let
\[
 {\cal R}_{n,x}:{\mathfrak S}_1({\cal H}_{{\rm in},x})
 \longrightarrow{\mathfrak S}_1({\cal H}_{\rm out})           \tag{33.12}
\]
be CPTP recovery maps and let
\(\Psi_{n,x}={\cal R}_{n,x}^*\).  Let
\(K_{{\rm in},x}\geq0\) be the input energy.  Assume the
extended-positive Loewner inequality
\[
 \Psi_{n,x}(K_x)
 \preceq aK_{{\rm in},x}+bI                                  \tag{33.13}
\]
for constants \(a,b\geq0\), uniformly in \(n,x\).

\begin{proposition}[Uniform actual-output moment]
\label{prop:v13v33-output-moment}
If \(d_{\ell,x}\) are input states satisfying
\[
 \sup_{x\in X}
 \operatorname{Tr}(K_{{\rm in},x}d_{\ell,x})
 \leq E_\ell,                                                  \tag{33.14}
\]
then
\[
 \sup_{n,x}
 \operatorname{Tr}\!\left(
 (I+K_x){\cal R}_{n,x}(d_{\ell,x})
 \right)
 \leq C_\ell^2,\qquad
 C_\ell^2:=1+aE_\ell+b.                                       \tag{33.15}
\]
\end{proposition}

\begin{proof}
Trace preservation gives the contribution one from the identity.
For the energy term, trace duality and (33.13) give
\[
\begin{split}
 \operatorname{Tr}\!\left(
 K_x{\cal R}_{n,x}(d_{\ell,x})\right)
 &=
 \operatorname{Tr}\!\left(
 d_{\ell,x}\Psi_{n,x}(K_x)\right)\\
 &\leq
 a\operatorname{Tr}(K_{{\rm in},x}d_{\ell,x})+b
 \leq aE_\ell+b.                                              \tag{33.16}
\end{split}
\]
Use bounded spectral truncations and monotone convergence for the
unbounded operators.
\end{proof}

This estimates the complete recovery output.  No norm of a Kraus
operator or Stinespring environment is substituted.

\subsection{The explicit uniform recovery tail}

\begin{theorem}[Loewner-to-compactness theorem]
\label{thm:v13v33-main}
Assume (33.2), (33.13), and (33.14).  Then
\[
\boxed{
 \sup_{n,x}
 \left\|
 {\cal R}_{n,x}(d_{\ell,x})
 -P_m{\cal R}_{n,x}(d_{\ell,x})P_m
 \right\|_1
 \leq
 {2C_\ell\over\sqrt{c_-(1+\lambda_{m+1})}}.}                  \tag{33.17}
\]
In particular, for every \(\ell\), the full background-cutoff orbit
\[
 \{
 {\cal R}_{n,x}(d_{\ell,x}):n\geq1,\ x\in X
 \}                                                           \tag{33.18}
\]
is relatively compact in trace norm.
\end{theorem}

\begin{proof}
Apply Lemma~\ref{lem:v13v33-reference-tail} with \(t=1\) and the
moment bound (33.15).  This gives (33.17).  For fixed \(m\), the
compressed states form a bounded subset of the finite-dimensional
trace class \(P_m{\mathfrak S}_1P_m\).  The right side of (33.17)
tends to zero as \(m\to\infty\), proving total boundedness of
(33.18).
\end{proof}

Equation (33.17) directly supplies the finite-rank compression
card (29.25).  It also supplies the Stinespring ellipsoid of V32
with dressing \(D_x=(I+K_x)^{-1/2}\).

\subsection{A counting-function rate}

Suppose the reference energy has the certified counting bound
\[
 N_{x_0}(\Lambda)
 :=
 \#\{j:\lambda_j\leq\Lambda\}
 \leq C_{\rm Weyl}(1+\Lambda)^{d/2}.                           \tag{33.19}
\]
Then
\[
 1+\lambda_{m+1}
 \geq
 \left({m+1\over C_{\rm Weyl}}\right)^{2/d}.                   \tag{33.20}
\]

\begin{corollary}[Polynomial finite-rank recovery rate]
\label{cor:v13v33-rate}
Under (33.19),
\[
 \sup_{n,x}
 \|\sigma_{n,x,\ell}-P_m\sigma_{n,x,\ell}P_m\|_1
 \leq
 {2C_\ell C_{\rm Weyl}^{1/d}\over\sqrt{c_-}}\,
 (m+1)^{-1/d},                                                 \tag{33.21}
\]
where
\(\sigma_{n,x,\ell}={\cal R}_{n,x}(d_{\ell,x})\).
\end{corollary}

\begin{proof}
If \(1+\lambda_{m+1}\) were smaller than the right side of (33.20),
the first \(m+1\) eigenvalues would contradict (33.19).  Substitute
(33.20) into (33.17).
\end{proof}

The exponent in (33.21) reflects only a first energy moment.  A
faster rate requires a certified higher moment and its own
background comparison; it is not obtained by iterating (33.13).

\subsection{Background-continuous and regional recovery limit}

Assume additionally that, for every dense test state, the actual
outputs obey a common background modulus
\[
 \|{\cal R}_{n,x}(d_{\ell,x})
     -{\cal R}_{n,y}(d_{\ell,y})\|_1
 \leq\omega_\ell(d_X(x,y)).                                   \tag{33.22}
\]
Then Theorem~\ref{thm:v13v33-main} supplies the collective
compactness and (33.22) supplies equicontinuity.

\begin{corollary}[Insertion into the normal Markov limit]
\label{cor:v13v33-insertion}
Under (33.22), one subsequence converges to a
background-continuous normal recovery field as in V29.  If the
split-screen states converge in predual norm and their recovery
residuals tend to zero, then the limiting states satisfy the exact
locally normal recovery identity (29.18).  On a countable regional
exhaustion, regionwise versions of (33.2), (33.13), and (33.22)
suffice via the V19 diagonal.
\end{corollary}

\begin{proof}
Apply V29 with the relative compactness from
Theorem~\ref{thm:v13v33-main} and the modulus (33.22).  Then apply
Theorem~\ref{thm:v13v29-recovery}.  The regional assertion is the
nested diagonal argument of V19.
\end{proof}

No constant must be uniform over all infinite-volume regions.
Uniformity is required only in cutoff and background at each fixed
local region.

\subsection{Symmetry transport}

Suppose an arrow \(\gamma:x\to y\) has unitary energy transport
\[
 K_y=U_\gamma K_xU_\gamma^*,\qquad
 K_{{\rm in},y}=V_\gamma K_{{\rm in},x}V_\gamma^*.             \tag{33.23}
\]
If the recovery channel is transported covariantly, then (33.13)
and the moment bound are invariant under \(\gamma\).  The V21
groupoid average also preserves (33.13), by
Proposition~\ref{prop:v13v21-energy}.  Therefore the compactness
tail can be established before or after equivariant descent.

The fixed reference projections \(P_m\) need not themselves be
group invariant.  They are a compactness device; the limiting
channel covariance is supplied by V21.

\subsection{Represented type-III version}

For a type-III output algebra normally represented on
\({\cal H}_{\rm out}\), assume the actual normal recovery outputs
have positive trace-class lifts satisfying (33.15).  Theorem
\ref{thm:v13v33-main} makes those lifts relatively compact.
Proposition~\ref{prop:v13v32-typeIII} then makes their restrictions
relatively compact in the type-III predual.  This supplies the
abstract hypothesis of V29.

The existence of such energy-controlled positive lifts is not
automatic.  It is a separate part of the split-screen construction.

\begin{assessment}[Status after thirteenth-cycle V33]
\label{ass:v13v33-status}
The recovery compactness chain is now explicit.  A Heisenberg
Loewner inequality transfers a bounded input energy moment to the
actual recovery output.  Uniform form comparison with one
compact-resolvent reference energy then gives a common finite-rank
trace-norm tail (33.17), with rate (33.21) under a counting bound.
Together with a background modulus and vanishing recovery residual,
this yields the locally normal background-continuous Markov limit
of V29.
\end{assessment}

\begin{firewall}[Claim boundary after thirteenth-cycle V33]
\label{fire:v13v33-claim}
No RPESM recovery adjoint, input moment, compact-resolvent output
energy, reference form comparison, counting bound, positive
type-III lift, background modulus, or recovery residual has been
populated.  Compact resolvent can fail in an infinite-volume
representation and must then be replaced by a local phase-space
certificate.  The theorem does not establish reflection positivity,
clustering, or the full SM--Einstein continuum.
\end{firewall}

\subsection{Next bottleneck}

The abstract recovery sector is now closed under explicit
certificates.  The next unresolved continuum gate is the state
itself: combine the locally normal Markov recovery field with the
Osterwalder--Schrader cone constructed in V8--V9 and determine a
closed condition under which recovery preserves reflection
positivity.  A generic CPTP recovery channel need not preserve the
positive-time reflected factorization.

\subsection{Parent record}

This V33 note appends to the validated thirteenth-cycle V32 file with
record
\[
\begin{array}{c|c}
\text{lines}&10286\\
\text{bytes}&362574\\
\text{SHA--256}&
\texttt{CB7132AE83A295A666473084D5938A41ABF226ADEEC50DBF82569316710E057E}.
\end{array}
\]

% =============================================================================
% END APPENDED THIRTEENTH-CYCLE V33 MATERIAL
% =============================================================================

% =============================================================================
% APPENDED THIRTEENTH-CYCLE V34 MATERIAL
% =============================================================================

\section{Thirteenth-cycle V34: reflection-positive Markov recovery}
\label{sec:v13v34}

\subsection{Why the recovery theorem needs another hypothesis}

The recovery adjoint constructed in V28--V33 is a unital completely
positive map
\[
 \Psi:{\cal A}_{\rm out}\longrightarrow{\cal A}_{\rm in}.       \tag{34.1}
\]
It transfers an input state to the recovered output state by
\[
 \omega_{\rm rec}(A)=\omega_{\rm in}(\Psi(A)).                  \tag{34.2}
\]
V8 treated reflection-factorized endomorphisms of one algebra.  The
present note gives the two-algebra version needed by recovery,
establishes its stability under collar composition and symmetry
averaging, and joins it to the locally normal compactness limit of
V29 and V33.

Let
\[
 ({\cal A}_{\rm in},{\cal A}_{{\rm in},+},\Theta_{\rm in}),
 \qquad
 ({\cal A}_{\rm out},{\cal A}_{{\rm out},+},\Theta_{\rm out})    \tag{34.3}
\]
be reflected unital \(C^*\)-algebras.  Each reflection is
conjugate-linear and isometric.  The state \(\omega_{\rm in}\) is
OS positive when
\[
 \omega_{\rm in}\bigl(\Theta_{\rm in}(H)H\bigr)\geq0,
 \qquad H\in{\cal A}_{{\rm in},+}.                              \tag{34.4}
\]

\subsection{Cross-algebra reflected factorization}

\begin{definition}[Reflection-factorized recovery adjoint]
\label{def:v13v34-cross-factor}
A UCP map \(\Psi:{\cal A}_{\rm out}\to{\cal A}_{\rm in}\) has a
finite cross-algebra reflected factorization if there are linear maps
\[
 L_\alpha:{\cal A}_{{\rm out},+}
 \longrightarrow{\cal A}_{{\rm in},+},
 \qquad \alpha=1,\ldots,r,                                     \tag{34.5}
\]
such that
\[
 \Psi\bigl(\Theta_{\rm out}(F)G\bigr)
 =
 \sum_{\alpha=1}^r
 \Theta_{\rm in}(L_\alpha F)L_\alpha G
 \quad
 (F,G\in{\cal A}_{{\rm out},+}).                               \tag{34.6}
\]
\end{definition}

\begin{theorem}[Recovery of an OS-positive state]
\label{thm:v13v34-recovery-os}
If \(\omega_{\rm in}\) is OS positive and \(\Psi\) satisfies
Definition~\ref{def:v13v34-cross-factor}, then
\(\omega_{\rm rec}=\omega_{\rm in}\circ\Psi\) is an OS-positive
state on \({\cal A}_{\rm out}\).
\end{theorem}

\begin{proof}
Unital complete positivity of \(\Psi\) makes
\(\omega_{\rm rec}\) a state.  For
\(F_1,\ldots,F_q\in{\cal A}_{{\rm out},+}\) and
\(c=(c_1,\ldots,c_q)\in{\mathbb C}^q\), (34.6) gives
\[
\begin{split}
 &\sum_{i,j=1}^q\overline{c_i}c_j
 \omega_{\rm rec}\bigl(\Theta_{\rm out}(F_i)F_j\bigr)\\
 &\quad =
 \sum_{\alpha=1}^r
 \omega_{\rm in}\left[
 \Theta_{\rm in}\left(\sum_i c_iL_\alpha F_i\right)
 \left(\sum_jc_jL_\alpha F_j\right)\right]\geq0.                \tag{34.7}
\end{split}
\]
Thus every recovered OS Gram matrix is positive semidefinite.
\end{proof}

This theorem is recovery-specific in two ways.  The positive-time
algebras may be different, and the direction of the Heisenberg map
is opposite to the direction in which states are recovered.

\subsection{Reflection covariance is not sufficient}

It is useful to isolate a false shortcut.  Let
\[
 {\cal A}=C(\{0,1\}\times\{0,1\}),\qquad
 (\Theta f)(x,y)=\overline{f(y,x)},                              \tag{34.8}
\]
and let \({\cal A}_+\) consist of functions of the second
coordinate.  A probability matrix \(P=(p_{xy})\) defines a state,
and its OS form on \(f\in{\cal A}_+\) is
\[
 \sum_{x,y=0}^1p_{xy}\overline{f(x)}f(y).                       \tag{34.9}
\]
Hence a reflection-invariant state is OS positive precisely when
the symmetric matrix \(P\) is positive semidefinite.

Let \(q_{01}=q_{10}=1/2\) and \(q_{00}=q_{11}=0\), and define the
replacement channel in the Heisenberg picture by
\[
 \Psi_q(f)=\left(\sum_{x,y}q_{xy}f(x,y)\right)I.                \tag{34.10}
\]
This map is UCP, commutes with \(\Theta\), and maps
\({\cal A}_+\) into the constants.  Its Schr\"odinger adjoint sends
every input state to \(q\).  But the OS matrix
\[
 Q=\begin{pmatrix}0&1/2\\1/2&0\end{pmatrix}                \tag{34.11}
\]
has eigenvalues \(1/2\) and \(-1/2\).  Therefore \(q\) is not OS
positive.

\begin{proposition}[Covariance gap]
\label{prop:v13v34-covariance-gap}
UCP, reflection covariance, and preservation of the positive-time
subalgebra do not imply preservation of OS positivity.
\end{proposition}

\begin{proof}
The channel (34.10) has all three properties, while the image state
has the negative OS eigenvalue displayed in (34.11).
\end{proof}

Thus the extra identity (34.6), or an equivalent direct OS-cone
certificate, cannot be removed from the recovery ledger.

\subsection{A matrix-valued residual}

For \(F_1,\ldots,F_q\in{\cal A}_{{\rm out},+}\), candidate maps
\(L_1,\ldots,L_r\), and a UCP map \(\Psi\), define
\[
 D^{(q)}_\Psi(F)_{ij}
 :=
 \Psi\bigl(\Theta_{\rm out}(F_i)F_j\bigr)
 -
 \sum_{\alpha=1}^r
 \Theta_{\rm in}(L_\alpha F_i)L_\alpha F_j.                    \tag{34.12}
\]
The matrix belongs to \(M_q({\cal A}_{\rm in})\).

\begin{proposition}[Complete OS residual bound]
\label{prop:v13v34-matrix-residual}
If \(D^{(q)}_\Psi(F)\) is self-adjoint and
\[
 \|D^{(q)}_\Psi(F)\|_{M_q({\cal A}_{\rm in})}\leq\delta,         \tag{34.13}
\]
then the recovered OS Gram matrix \(G_{\rm rec}\) satisfies
\[
 G_{\rm rec}\succeq-\delta I_q.                                 \tag{34.14}
\]
If only the entrywise estimates
\(\|D^{(q)}_\Psi(F)_{ij}\|\leq\varepsilon\) are known, then
\[
 G_{\rm rec}\succeq-q\varepsilon I_q.                           \tag{34.15}
\]
\end{proposition}

\begin{proof}
Apply the completely positive contraction
\[
 {\rm id}_{M_q}\otimes\omega_{\rm in}:
 M_q({\cal A}_{\rm in})\longrightarrow M_q({\mathbb C})         \tag{34.16}
\]
to (34.12).  The factorized part is a sum of OS Gram matrices and
is positive semidefinite by (34.4).  The image of the residual has
operator norm at most \(\delta\), which proves (34.14).  An
entrywise \(q\times q\) matrix bounded by \(\varepsilon\) has
operator norm at most \(q\varepsilon\), proving (34.15).
\end{proof}

The matrix norm in (34.13) is the appropriate finite-ledger
quantity.  A diagonal bound only on
\(\Psi(\Theta(F)F)\) does not by itself control all polarized Gram
matrices uniformly in \(q\).

\subsection{Composition through nested collars}

Suppose
\[
 \Psi_1:{\cal A}_2\to{\cal A}_1,\qquad
 \Psi_2:{\cal A}_1\to{\cal A}_0                               \tag{34.17}
\]
have reflected factorizations \(L_\alpha\) and \(M_\beta\),
respectively.  Thus the \(L_\alpha\) map the positive-time algebra
of \({\cal A}_2\) into that of \({\cal A}_1\), while the
\(M_\beta\) map the latter into that of \({\cal A}_0\).

\begin{proposition}[Collar-composition law]
\label{prop:v13v34-composition}
The composite \(\Psi_2\circ\Psi_1\) has reflected factorization
\[
 (\Psi_2\circ\Psi_1)
 \bigl(\Theta_2(F)G\bigr)
 =
 \sum_{\alpha,\beta}
 \Theta_0(M_\beta L_\alpha F)M_\beta L_\alpha G.                \tag{34.18}
\]
Consequently every finite composition of
reflection-factorized collar recoveries preserves OS positivity.
\end{proposition}

\begin{proof}
Insert the factorization of \(\Psi_1\), use linearity of
\(\Psi_2\), and apply its factorization to each pair
\(L_\alpha F,L_\alpha G\).  The result is (34.18).  Theorem
\ref{thm:v13v34-recovery-os} then applies.
\end{proof}

This is the compatibility required by the nested regional
construction of V19.  Pairwise causal consistency alone would not
supply (34.18).

\subsection{Symmetry averaging}

Let a compact group \(G\) preserve the two reflected algebras and
their positive-time subalgebras, and commute with both reflections.
For each \(g\in G\), let \(\Psi_g\) be a transported recovery
adjoint.  Assume \(g\mapsto\Psi_g(A)\) is norm continuous and every
\(\Psi_g\) preserves OS positivity in the sense of
Theorem~\ref{thm:v13v34-recovery-os}.  Define
\[
 \overline\Psi(A)=\int_G\Psi_g(A)\,d\mu(g).                     \tag{34.19}
\]

\begin{proposition}[OS-safe compact averaging]
\label{prop:v13v34-average}
The map \(\overline\Psi\) is UCP and sends every OS-positive input
state to an OS-positive recovered state.
\end{proposition}

\begin{proof}
The Bochner integral of UCP maps is UCP.  For a fixed finite family
\(F_1,\ldots,F_q\), the OS Gram matrix recovered through
\(\overline\Psi\) is
\[
 \int_G
 [\,\omega_{\rm in}(\Psi_g(\Theta_{\rm out}(F_i)F_j))\,]_{ij}
 \,d\mu(g).                                                     \tag{34.20}
\]
Every integrand is positive semidefinite, and the positive
semidefinite cone is closed and convex.  Its integral is therefore
positive semidefinite.
\end{proof}

A common reflection is essential.  If an arrow transports the
reflection plane to a different positive-time algebra, the
integrands in (34.20) do not represent one fixed OS cone.  They
must first be pulled back to a common reflected chart.  This is the
reflection analogue of the quotient-chart compatibility in V21.

For a finite group, if \(\Psi_g\) has factor maps \(L_{g,\alpha}\),
then \(\overline\Psi\) has factor maps
\(|G|^{-1/2}L_{g,\alpha}\).  For a compact group, (34.20) supplies
the needed positivity without asserting a finite Kraus family.

\subsection{Passage to the locally normal Markov limit}

Let \({\cal D}_+\) be a norm-dense unital subalgebra of
\({\cal A}_{{\rm out},+}\).  At cutoff \(n\), let
\(\omega_{{\rm in},n}\) be OS positive and let
\(\Psi_n\) be a normal UCP recovery adjoint.  Put
\[
 \omega_{{\rm rec},n}=\omega_{{\rm in},n}\circ\Psi_n.           \tag{34.21}
\]
Assume the V29--V33 compactness hypotheses give local predual
convergence
\[
 \|\omega_{{\rm rec},n}-\omega_{\rm rec}\|_{
 {\cal A}_{\rm out}(O)_*}\longrightarrow0                      \tag{34.22}
\]
on every fixed bounded region \(O\).

Choose a countable family of finite ledgers
\[
 {\cal F}_s=\{F_{s,1},\ldots,F_{s,q_s}\}\subset{\cal D}_+        \tag{34.23}
\]
whose union is norm dense in \({\cal D}_+\).  Suppose the
matrix residuals satisfy
\[
 \delta_{n,s}:=
 \|D^{(q_s)}_{\Psi_n}({\cal F}_s)\|
 \longrightarrow0
 \quad\hbox{for every fixed }s.                                \tag{34.24}
\]

\begin{theorem}[OS-positive locally normal recovery limit]
\label{thm:v13v34-limit}
Under (34.22)--(34.24), the limiting recovered state is OS positive.
If, in addition, the recovery residuals of V29 tend to zero, the
same limiting state satisfies the exact locally normal Markov
recovery identity.  If the energy and background hypotheses of
V33 hold, the convergence can be chosen simultaneously over the
compact background core.
\end{theorem}

\begin{proof}
Fix \(s\).  Proposition~\ref{prop:v13v34-matrix-residual} gives
\[
 [\,\omega_{{\rm rec},n}
   (\Theta_{\rm out}(F_{s,i})F_{s,j}))\,]_{ij}
 \succeq-\delta_{n,s}I.                                        \tag{34.25}
\]
All entries are local for a sufficiently large fixed region.
Predual convergence in (34.22), followed by (34.24), makes the
limiting Gram matrix positive semidefinite.

Now let \(F\in{\cal A}_{{\rm out},+}\) and choose
\(F_k\) from the union of the ledgers with
\(\|F_k-F\|\to0\).  Since \(\Theta_{\rm out}\) is isometric,
\[
 \|\Theta_{\rm out}(F_k)F_k-\Theta_{\rm out}(F)F\|
 \leq
 \|F_k-F\|(\|F_k\|+\|F\|)\longrightarrow0.                     \tag{34.26}
\]
Continuity of the state extends positivity from the dense ledger to
\(F\).  The exact recovery identity follows independently from
Theorem~\ref{thm:v13v29-recovery}; both closed conditions hold on
the same subsequential limit.  The background-uniform assertion is
the diagonal supplied by Corollary~\ref{cor:v13v33-insertion}.
\end{proof}

The proof deliberately passes the finite Gram inequalities rather
than the factor maps \(L_{\alpha,n}\).  Their ranks and
representations may vary with \(n\); no convergence of individual
Kraus data is required.

\subsection{A combined cutoff certificate}

For a fixed region \(O\), background \(x\), and ledger
\({\cal F}_s\), the recovery row now has four independent numbers:
\[
 \begin{array}{c|l}
 r^{\rm Markov}_{n,O,x}
   &\hbox{trace-norm recovery defect from V28--V29},\\
 \delta^{\rm OS}_{n,O,x,s}
   &\hbox{complete reflected Gram residual (34.24)},\\
 \tau^{\rm en}_{n,O,x,\ell}(m)
   &\hbox{energy tail bounded by (33.17)},\\
 \eta^{\rm bg}_{n,O,\ell}(x,y)
   &\hbox{background modulus from (33.22)}.
 \end{array}                                                    \tag{34.27}
\]
A cofinal schedule must force the first two entries to zero, the
third uniformly to zero as \(m\to\infty\), and the fourth to a
common modulus.  None of these implications follows from the other
three.  Under these four certified rows, Theorem
\ref{thm:v13v34-limit} and V33 yield a background-continuous,
locally normal, OS-positive Markov recovery state.

\subsection{Interaction with the energy inequality}

If the group action in Proposition~\ref{prop:v13v34-average}
transports \(K_{\rm in}\) and \(K_{\rm out}\) covariantly, then the
Loewner inequality (33.13) survives averaging by V21.  Hence the
same averaged recovery is both OS safe and energy tight.  For
nested collars, repeated use of
\[
 \Psi_j(K_j)\preceq a_jK_{j-1}+b_jI                            \tag{34.28}
\]
gives
\[
 (\Psi_1\circ\cdots\circ\Psi_m)(K_m)
 \preceq
 \left(\prod_{j=1}^ma_j\right)K_0+
 \sum_{k=1}^m b_k\prod_{j=1}^{k-1}a_j\,I,                      \tag{34.29}
\]
with the empty product equal to one.  This follows by positivity
and induction.  The constants must be controlled on each fixed
regional chain; reflection factorization does not control them.

\begin{assessment}[Status after thirteenth-cycle V34]
\label{ass:v13v34-status}
The locally normal recovery programme now has a closed
reflection-positivity gate.  A cross-algebra reflected
factorization transfers an OS-positive collar state to the
recovered state, composes through nested collars, and survives
compact symmetry averaging when all arrows preserve one reflected
chart.  Matrix-valued approximate factorizations pass to the
predual limit without convergence of Kraus operators.  Combined
with V33, this yields an exact OS-positive Markov recovery limit
under four explicit cutoff certificates.
\end{assessment}

\begin{firewall}[Claim boundary after thirteenth-cycle V34]
\label{fire:v13v34-claim}
No RPESM positive-time recovery algebra, reflected recovery
factorization, OS ledger residual, common reflection chart, or
reflection-compatible symmetry average has been constructed.
The finite classical example proves that covariance and complete
positivity cannot replace this missing input.  The note does not
establish clustering, analytic continuation, anomaly cancellation,
Einstein compactness, or the full SM--Einstein continuum.
\end{firewall}

\subsection{Next bottleneck}

The remaining constructive question is finite and testable: express
(34.6) on each compressed operator system as a semidefinite
feasibility condition compatible with the ordinary Choi
spectrahedra of V11--V16, the causal expectation relations of V26,
and the Petz recovery constraints of V28.  This will determine
whether OS-positive recovery can be selected continuously over the
background rather than merely assumed at the limit.

\subsection{Parent record}

This V34 note appends to the validated thirteenth-cycle V33 file with
record
\[
\begin{array}{c|c}
\text{lines}&10646\\
\text{bytes}&374564\\
\text{SHA--256}&
\texttt{71F55AC26BB8A6D9D529F8F4B36C51F3191229B00ED3AE6506F99033EB781FE8}.
\end{array}
\]

% =============================================================================
% END APPENDED THIRTEENTH-CYCLE V34 MATERIAL
% =============================================================================


% =============================================================================
% APPENDED THIRTEENTH-CYCLE V35 MATERIAL
% =============================================================================

\section{Thirteenth-cycle V35: mandatory YM--NS companion audit}
\label{sec:v13v35}

\subsection{Files inspected}

After completing V34, the current companion programmes in the
shared \texttt{latex\_notes} directory were inspected read-only.
Their exact records are
\[
\begin{array}{c|r|r|c}
\text{file}&\text{lines}&\text{bytes}&\text{SHA--256}\\ \hline
\texttt{Renewal\_Yang\_Mills\_Mass\_Gap\_Research\_Notes\_V81.tex}
&27855&949968&
\texttt{D59C5389D96AFFC4A1A69F2B87CF75849F06748AE0B375286C712B613D7600E1}\\
\texttt{Renewal\_NS\_Stability\_Research\_Notes\_V26.tex}
&5319&278283&
\texttt{82C887F8C2C69E4F28FAD7C3DC8DE5B9099CB5A4BF431EBA1FFF3E91935978AD}.
\end{array}                                                    \tag{35.1}
\]
The recently modified auxiliary file
\[
 \texttt{Renewal\_YM\_Parallel\_Research\_Notes\_V5.tex}        \tag{35.2}
\]
was also inspected.  It has \(916\) lines, \(35778\) bytes, and
SHA--256
\[
 \texttt{9924164DA923C19A2FFA590FB14912917049256642D5BEDD2833DFA147BDD16D}.
                                                                    \tag{35.3}
\]
It is a completed side programme conditional on a main-response
margin and does not supersede the active YM V81 programme.

\subsection{Yang--Mills information relevant to the SM programme}

YM V80 replaces coefficientwise phase transport by a differential
energy-metric inequality.  Along a path it estimates the complete
contracted matrices
\[
 \dot H=\sum_\mu\xi_\mu\partial_\mu H,\qquad
 \dot X=\sum_\mu\xi_\mu\partial_\mu X                          \tag{35.4}
\]
before applying an operator or Frobenius norm.  Its Gronwall
transport has no small-displacement hypothesis.  YM V81 then finds
an exact signed-diagonal frequency collapse: the four-variable
Laurent data reduce along each selected diagonal to a degree-four
univariate trigonometric matrix bank.  The frequency collapse and
endpoint reconstruction are rigorous, while the reported path
lengths remain floating diagnostics pending rational interval
enclosures.

The transferable lesson concerns proof order.  The V34 OS condition
is a positivity statement for a complete Gram block.  Splitting that
block into entrywise absolute bounds before testing its least
eigenvalue discards the same kind of correlated cancellation that
YM V80 preserves in (35.4).  The next SM step should therefore keep
the entire reflected Gram kernel as one semidefinite variable and
transport its whole-matrix margin along background paths.

YM V81 does not provide an SM recovery map, OS-positive state, Choi
matrix, or Einstein-background compactness estimate.  Its floating
path data cannot be imported as an SM bound.

\subsection{Navier--Stokes information relevant to the SM programme}

NS V26 contracts the Oseen derivative tensors against the prescribed
rank-one input before maximizing the output norm.  This produces no
gain for the first derivative and only a localized modest gain for
the second.  The result is instructive in both directions:
structured contraction can expose a real improvement, but it can
also prove that a hoped-for improvement is absent.

For the SM recovery problem, the prescribed structured input is the
positive-time family \(F_1,\ldots,F_q\).  Its natural contraction is
the full OS Gram matrix
\[
 [\,\omega_{\rm in}\Psi(\Theta(F_i)F_j)\,]_{i,j=1}^q.           \tag{35.5}
\]
The V34 matrix residual is therefore retained.  No scalar
coefficient surrogate will replace (35.5) unless a certified
comparison proves that it loses no needed margin.

NS V26 contains no Standard-Model Choi constraint, recovery
channel, reflection factorization, determinant-line descent, or
Einstein tightness theorem.  Its numerical constants do not enter
the SM ledger.

\subsection{The auxiliary YM side programme}

The auxiliary YM V5 proves, under its declared main-response margin,
a regulator-uniform nonlinear remainder of strictly higher order
and a summable accumulation bound.  Its useful logical pattern is a
margin-protection theorem:
\[
 \|\hbox{remainder}\|<\hbox{certified interior margin}.         \tag{35.6}
\]
The SM analogue is to compare the complete perturbation of every
Choi, OS, energy, and causal LMI block with the smallest eigenvalue
margin of that block.  The side programme does not supply those
margins; it only confirms that the correct endpoint is a
quantitative interior certificate rather than a formal expansion.

\begin{proposition}[Audit decision]
\label{prop:v13v35-decision}
No companion result closes an SM--Einstein hypothesis or alters the
claim boundary of V34.  The next SM construction should:
\begin{enumerate}
\item introduce a positive Gram lift for the finite reflected
factorization, making it a genuine semidefinite feasibility
condition;
\item place that lift in the same spectrahedron as UCP, causal,
energy, symmetry, and recovery constraints;
\item transport the complete LMI blocks along background paths and
compare their operator-norm variation with certified interior
margins.
\end{enumerate}
\end{proposition}

\begin{proof}
The YM and NS results control structured finite-dimensional
contractions in their own programmes but provide none of the
objects required to instantiate an SM recovery channel.  Both show
that the prescribed contraction must precede the ambient norm.
For (35.5), a positive Gram lift retains precisely that contraction.
All listed recovery-side constraints are affine in a Choi variable
once the background and test data are fixed; hence their conjunction
can be treated at the full-block semidefinite level.  This is the
only stated consequence of the audit.
\end{proof}

\begin{assessment}[Status after thirteenth-cycle V35]
\label{ass:v13v35-status}
The required five-version audit is complete.  YM V80--V81 and NS
V26 reinforce a common rule: preserve the complete prescribed
matrix contraction until after its correlations have combined.
Accordingly, the SM programme will replace coefficientwise OS
tests by a lifted Gram spectrahedron and will measure path variation
in the operator norm of each full LMI block.
\end{assessment}

\begin{firewall}[Claim boundary after thirteenth-cycle V35]
\label{fire:v13v35-claim}
The audit imports no YM or NS numerical estimate.  The YM
signed-diagonal path bounds are not yet certified, and the NS
rank-one constants concern the Oseen kernel.  No SM reflected Gram
lift, strictly feasible recovery point, background path enclosure,
continuum measure, clustering theorem, or full SM--Einstein
continuum follows from this comparison.
\end{firewall}

\subsection{Next step}

V36 will construct the finite reflected Gram lift explicitly and
prove that its conjunction with the previously isolated recovery
conditions is a spectrahedron.  It will then formulate a
whole-block pathwise margin theorem suitable for certified
background transport.

\subsection{Parent record}

This V35 note appends to the validated thirteenth-cycle V34 file with
record
\[
\begin{array}{c|c}
\text{lines}&11077\\
\text{bytes}&390241\\
\text{SHA--256}&
\texttt{C41D413859D5A173BAD748E37F2EB7AD0EA9BD1D83AB3ABD1B41C14938B90EB0}.
\end{array}
\]

% =============================================================================
% END APPENDED THIRTEENTH-CYCLE V35 MATERIAL
% =============================================================================

% =============================================================================
% APPENDED THIRTEENTH-CYCLE V36 MATERIAL
% =============================================================================

\section{Thirteenth-cycle V36: the lifted reflected-recovery
spectrahedron}
\label{sec:v13v36}

\subsection{Finite reflected systems}

Let
\[
 E_+=\operatorname{span}\{F_1,\ldots,F_s\}
 \subset{\cal A}_{{\rm out},+},\qquad
 H_+=\operatorname{span}\{H_1,\ldots,H_t\}
 \subset{\cal A}_{{\rm in},+},                                \tag{36.1}
\]
where both displayed families are linearly independent.  Define the
fixed reflected products
\[
 A_{pq}:=\Theta_{\rm out}(F_p)F_q,\qquad
 B_{uv}:=\Theta_{\rm in}(H_u)H_v.                              \tag{36.2}
\]
The recovery adjoint is required to factor on
\(E_+\), with all factor-map values lying in \(H_+\).

The apparent obstacle is that (34.6) is quadratic in the
coefficients of the factor maps.  A Gram lift removes that
quadratic dependence without splitting the reflected matrix kernel.

\subsection{The exact positive Gram lift}

Index \({\mathbb C}^{ts}\) by pairs \((u,p)\), with
\(1\leq u\leq t\) and \(1\leq p\leq s\).

\begin{theorem}[Finite reflected-factorization lift]
\label{thm:v13v36-gram-lift}
For a linear map
\(\Psi:{\cal A}_{\rm out}\to{\cal A}_{\rm in}\), the following are
equivalent.
\begin{enumerate}
\item There are finitely many linear maps
\(L_\alpha:E_+\to H_+\) such that
\[
 \Psi(A_{pq})
 =
 \sum_\alpha\Theta_{\rm in}(L_\alpha F_p)L_\alpha F_q
 \quad(1\leq p,q\leq s).                                      \tag{36.3}
\]
\item There is a matrix \(W\in{\rm Herm}_{ts}\), \(W\succeq0\),
such that
\[
 \boxed{\quad
 \Psi(A_{pq})
 =
 \sum_{u,v=1}^t
 W_{(u,p),(v,q)}B_{uv}
 \quad(1\leq p,q\leq s).\quad}                                 \tag{36.4}
\]
\end{enumerate}
Moreover, the maps in the first statement may be chosen with their
number equal to \({\rm rank}\,W\leq ts\).
\end{theorem}

\begin{proof}
Write
\[
 L_\alpha F_p=\sum_{u=1}^t\ell_{\alpha,u p}H_u.                 \tag{36.5}
\]
Because \(\Theta_{\rm in}\) is conjugate-linear,
\[
 \Theta_{\rm in}(L_\alpha F_p)L_\alpha F_q
 =
 \sum_{u,v}
 \overline{\ell_{\alpha,u p}}\ell_{\alpha,v q}B_{uv}.           \tag{36.6}
\]
Thus (36.3) implies (36.4) with
\[
 W_{(u,p),(v,q)}
 =
 \sum_\alpha
 \overline{\ell_{\alpha,u p}}\ell_{\alpha,v q}.                \tag{36.7}
\]
This is a Gram matrix and is positive semidefinite.

Conversely, factor a positive semidefinite \(W\) as
\[
 W_{ab}=\sum_{\alpha=1}^r
 \overline{\ell_{\alpha,a}}\ell_{\alpha,b},
 \qquad r={\rm rank}\,W.                                       \tag{36.8}
\]
Define \(L_\alpha\) on the basis \(F_p\) by (36.5).
Equations (36.6)--(36.8) give (36.3).  Linear independence in
(36.1) makes the definitions unambiguous.
\end{proof}

The variable \(W\) retains all correlations among the input
positive-time basis, the output factor range, and the factor label.
Entrywise bounds on (36.4) would discard this structure.

\subsection{Choi coordinates}

For clarity, take finite matrix algebras
\[
 {\cal A}_{\rm in}=M_{d_{\rm in}}({\mathbb C}),\qquad
 {\cal A}_{\rm out}=M_{d_{\rm out}}({\mathbb C}).               \tag{36.9}
\]
For a Schr\"odinger recovery channel
\({\cal R}:M_{d_{\rm in}}\to M_{d_{\rm out}}\), use the Choi
convention
\[
 J_{\cal R}
 =
 \sum_{a,b=1}^{d_{\rm in}}
 {\cal R}(E_{ab})\otimes E_{ab}
 \in M_{d_{\rm out}}\otimes M_{d_{\rm in}}.                    \tag{36.10}
\]
Then
\[
 J_{\cal R}\succeq0,\qquad
 {\rm Tr}_{\rm out}J_{\cal R}=I_{\rm in}                        \tag{36.11}
\]
are equivalent to complete positivity and trace preservation.  The
adjoint is the linear function of \(J=J_{\cal R}\)
\[
 \Psi_J(A)
 =
 \left\{
 {\rm Tr}_{\rm out}\bigl[(A\otimes I_{\rm in})J\bigr]
 \right\}^{\mathsf T}.                                         \tag{36.12}
\]
Indeed, its \((b,a)\) entry is
\({\rm Tr}(A{\cal R}(E_{ab}))\), which is the defining trace-duality
identity.

\begin{corollary}[Universal finite OS recovery is an SDP lift]
\label{cor:v13v36-os-sdp}
The pairs \((J,W)\) for which \({\cal R}_J\) is CPTP and its
adjoint has a reflected factorization from \(E_+\) through \(H_+\)
are cut out by
\[
\begin{split}
 &J\succeq0,\qquad W\succeq0,\qquad
 {\rm Tr}_{\rm out}J=I_{\rm in},\\
 &\Psi_J(A_{pq})
 =
 \sum_{u,v}W_{(u,p),(v,q)}B_{uv}
 \quad(1\leq p,q\leq s).                                      \tag{36.13}
\end{split}
\]
Thus (36.13) is a spectrahedron in the lifted variables
\((J,W)\).  Its projection onto \(J\) is a spectrahedral shadow.
Every feasible \(J\) sends every OS-positive input state to an
OS-positive recovered state on the ledger \(E_+\).
\end{corollary}

\begin{proof}
Equations (36.11)--(36.12) and (36.4) are affine equalities in
\((J,W)\), while the two positivity conditions are LMIs.  The last
assertion follows from Theorem~\ref{thm:v13v34-recovery-os},
restricted to \(E_+\).
\end{proof}

It is important to call the projection a spectrahedral shadow:
eliminating \(W\) need not leave a set described by one LMI in
\(J\) alone.

\subsection{State-specific OS positivity as a smaller LMI}

For a fixed OS-positive input state \(\omega_{\rm in}\), define
\[
 G^{\rm OS}_{pq}(J)
 :=
 \omega_{\rm in}\bigl(\Psi_J(A_{pq})\bigr).                    \tag{36.14}
\]
The state-specific condition
\[
 G^{\rm OS}(J)\succeq0                                         \tag{36.15}
\]
is itself an LMI in \(J\).  It is weaker than (36.13): it certifies
only the selected state and ledger, while the lift (36.13)
certifies a factorization that works for every OS-positive input
state.  The distinction must be retained in any numerical ledger.

\subsection{Adding the other recovery constraints}

All finite constraints isolated earlier can be placed in the same
lifted feasibility problem.

\begin{proposition}[Joint finite recovery spectrahedron]
\label{prop:v13v36-joint}
At a fixed background and finite ledger, the conjunction of the
following requirements is a spectrahedron after adjoining the
indicated auxiliary variables:
\begin{enumerate}
\item CPTP, by (36.11);
\item universal reflected factorization, by (36.13), or
state-specific OS positivity, by (36.15);
\item the causal expectation relation
\[
 E_{\rm in}\Psi_J(C)=\Psi_J(E_{\rm out}C)                       \tag{36.16}
\]
on a finite test basis;
\item finite symmetry covariance of \({\cal R}_J\);
\item the energy-transfer inequality
\[
 aK_{\rm in}+bI-\Psi_J(K_{\rm out})\succeq0;                   \tag{36.17}
\]
\item exact recovery of any fixed finite list of input states;
\item a trace-norm recovery residual
\[
 \|{\cal R}_J(\rho)-\sigma\|_1\leq\varepsilon.                 \tag{36.18}
\]
\end{enumerate}
For (36.18), one may introduce \(T\succeq0\) and impose
\[
 -T\preceq{\cal R}_J(\rho)-\sigma\preceq T,\qquad
 {\rm Tr}\,T\leq\varepsilon.                                  \tag{36.19}
\]
\end{proposition}

\begin{proof}
Items 3, 4, and 6 are affine equations in \(J\) on the selected
finite bases.  Item 5 is an LMI because (36.12) is linear in \(J\).
For the Hermitian difference
\(\Delta={\cal R}_J(\rho)-\sigma\), (36.19) is equivalent to
\(\|\Delta\|_1\leq\varepsilon\): choosing \(T=|\Delta|\) proves one
direction, while the dual formula
\[
 \|\Delta\|_1=\max_{-I\preceq S\preceq I}{\rm Tr}(S\Delta)      \tag{36.20}
\]
and \(-T\preceq\Delta\preceq T\) give
\(\|\Delta\|_1\leq{\rm Tr}\,T\) for the other.  Every condition is
therefore an affine equality or an LMI in the lifted variables.
\end{proof}

For a fixed Petz map, its Choi matrix can either be inserted as a
candidate point or imposed by the affine equality \(J=J_{\rm
Petz}\).  The dependence of \(J_{\rm Petz}\) on the background
densities is nonlinear, but after those densities are fixed the
feasibility test above remains semidefinite.

\subsection{Exact equalities along a background path}

Let \(z\) collect \(J,W,T\), and any other finite auxiliary
variables.  Write all exact affine constraints as
\[
 {\cal C}(t)z(t)=d(t),\qquad 0\leq t\leq T.                    \tag{36.21}
\]
Here the background path has already been chosen.  Suppose
\({\cal C}(t)\) is \(C^1\), surjective, and has a continuous right
inverse \(R(t)\):
\[
 {\cal C}(t)R(t)=I.                                            \tag{36.22}
\]

\begin{lemma}[Equality-preserving transport equation]
\label{lem:v13v36-equality-ode}
If \(z(0)\) satisfies (36.21), then every solution of
\[
 \dot z
 =
 R(t)\bigl(\dot d(t)-\dot{\cal C}(t)z(t)\bigr)
 +
 \bigl(I-R(t){\cal C}(t)\bigr)v(t)                             \tag{36.23}
\]
satisfies (36.21) for all times on which it exists.  The tangent
field \(v(t)\) is arbitrary.
\end{lemma}

\begin{proof}
Differentiate the residual and use (36.22):
\[
\begin{split}
 {d\over dt}({\cal C}z-d)
 &=
 \dot{\cal C}z+
 {\cal C}R(\dot d-\dot{\cal C}z)
 +{\cal C}(I-R{\cal C})v-\dot d=0.                             \tag{36.24}
\end{split}
\]
The residual vanishes initially, hence vanishes identically.
\end{proof}

Surjectivity is a local regularity hypothesis, not an automatic
property of the recovery constraints.  Redundant equalities should
first be removed.  At a rank change, one must change charts or use a
different selection argument.

\subsection{Whole-block margin transport}

Let the LMI blocks of Proposition~\ref{prop:v13v36-joint} be denoted
\[
 {\cal B}_r(t,z(t))\in{\rm Herm}_{m_r},\qquad r=1,\ldots,N.      \tag{36.25}
\]
This includes \(J\), \(W\), the OS Gram block when used, the energy
slack, and the two trace-norm blocks \(T\pm\Delta\).

\begin{theorem}[Pathwise strict-margin protection]
\label{thm:v13v36-path-margin}
Assume \(z(t)\) obeys the exact equalities and every block in
(36.25) is absolutely continuous.  If
\[
 {\cal B}_r(0,z(0))\succeq\gamma_r I,\qquad \gamma_r>0,          \tag{36.26}
\]
and
\[
 \int_0^T
 \left\|
 {d\over dt}{\cal B}_r(t,z(t))
 \right\|_{\rm op}\,dt
 <\gamma_r
 \quad(r=1,\ldots,N),                                          \tag{36.27}
\]
then every LMI remains strictly feasible on the whole path.  More
precisely,
\[
 {\cal B}_r(t,z(t))
 \succeq
 \left(
 \gamma_r-
 \int_0^t\|\dot{\cal B}_r(s,z(s))\|_{\rm op}\,ds
 \right)I.                                                     \tag{36.28}
\]
\end{theorem}

\begin{proof}
The fundamental theorem for absolutely continuous matrix-valued
functions gives
\[
 \|{\cal B}_r(t,z(t))-{\cal B}_r(0,z(0))\|_{\rm op}
 \leq
 \int_0^t\|\dot{\cal B}_r(s,z(s))\|_{\rm op}\,ds.               \tag{36.29}
\]
For self-adjoint matrices, \(X\succeq
-\|X\|_{\rm op}I\).  Add this inequality to (36.26) to obtain
(36.28); (36.27) makes the coefficient positive.
\end{proof}

The derivative in (36.27) is taken after forming each complete LMI
block.  This is the precise SM counterpart of the complete
pathwise matrix contraction selected by the V35 companion audit.

\begin{corollary}[Certified finite path cover]
\label{cor:v13v36-path-cover}
Suppose a compact background core is covered by finitely many
piecewise \(C^1\) paths issuing from strictly feasible anchors.
If (36.21)--(36.23) and the block integrals (36.27) are certified
on every path segment, then every point reached by the path network
has a feasible lifted recovery.  Concatenation is valid when the
terminal lifted variable of one segment is the initial variable of
the next.
\end{corollary}

\begin{proof}
Apply Theorem~\ref{thm:v13v36-path-margin} successively to the
segments.  The equality residual remains zero by
Lemma~\ref{lem:v13v36-equality-ode}, and the positive margin at the
end of one segment supplies the initial margin for the next.
\end{proof}

This corollary proves feasibility along the certified network.  It
does not by itself prove path independence.  Different paths may
arrive at different feasible lifts; convexity can reconcile them
locally, but compatibility around loops is a separate descent
condition.

\subsection{Compactness of lifts and the Gram gauge}

The Choi set satisfying (36.11) is compact.  The Gram lift need not
be: linear dependencies among the products \(B_{uv}\) can leave
positive directions of \(W\) invisible in (36.4).  For continuous
selection and inverse limits one must add a gauge such as
\[
 {\rm Tr}\,W\leq M                                             \tag{36.30}
\]
or select a trace-minimizing lift.

\begin{proposition}[Bounded Gram gauge]
\label{prop:v13v36-gram-gauge}
For a fixed feasible \(J\), the set of feasible \(W\) contains a
trace minimizer.  If the feasible set is nonempty, the minimum is
finite.  Adding
\({\rm Tr}\,W\leq M\) with \(M\) at least that minimum preserves
feasibility and makes the \(W\)-fiber compact.
\end{proposition}

\begin{proof}
Choose one feasible \(W_0\), so the infimum is at most
\({\rm Tr}\,W_0<\infty\).  Intersect the closed feasible set with
\(\{W\succeq0:{\rm Tr}\,W\leq{\rm Tr}\,W_0\}\).  For positive
semidefinite matrices, the trace bounds every eigenvalue, so this
intersection is compact.  The continuous trace function attains
its minimum there.  The same trace bound makes every sublevel
fiber compact.
\end{proof}

Uniformity of \(M\) over a background core still requires proof.
A strictly feasible continuous lift on a compact core supplies one
by taking the maximum of its trace.

\subsection{Relation to finite ledgers and cofinal limits}

If \(E_+\) and \(H_+\) increase with the cutoff, restriction of a
large Gram matrix to the indices of a smaller basis is positive.
The factorization equations also restrict correctly when the
smaller reflected-product systems are literal subsystems.  Thus
principal compression gives the natural projective map
\[
 (J_{n+1},W_{n+1})\longmapsto(J_n,W_n).                         \tag{36.31}
\]
The Choi component may additionally require the coarse-graining
map of V7 and V13.  Existence of feasible points at every level is
not enough: one needs nonempty fibers compatible with the full
map (36.31), or a summable compatibility residual.  This is the
next infinite-ledger gate.

\begin{assessment}[Status after thirteenth-cycle V36]
\label{ass:v13v36-status}
Finite reflected factorization is now an exact semidefinite
condition.  The positive matrix \(W\) is a Gram lift of all
half-space factor maps, and its conjunction with Choi positivity,
causal intertwiners, symmetry, energy transfer, and recovery
residuals is a spectrahedron in lifted variables.  Exact affine
constraints can be transported by (36.23), while the whole-block
margin theorem (36.28) certifies positivity along background paths
without coefficientwise splitting.
\end{assessment}

\begin{firewall}[Claim boundary after thirteenth-cycle V36]
\label{fire:v13v36-claim}
No RPESM Choi matrix, Gram lift, strict spectral margin, equality
right inverse, path integral, uniform Gram gauge, or compatible
cofinal family has been populated.  State-specific OS positivity
must not be confused with universal reflected factorization.
Finite path feasibility does not prove loop descent, continuum
reflection positivity, clustering, analytic continuation, anomaly
cancellation, or the full SM--Einstein continuum.
\end{firewall}

\subsection{Next bottleneck}

The finite convex gate is closed abstractly.  The next task is to
construct a projectively compatible family of bounded Gram lifts
across increasing positive-time ledgers.  The main obstruction is
that choosing a trace-minimizing \(W_n\) independently at each
level need not commute with principal compression.  The next note
should formulate the exact extension fiber, prove its compactness,
and identify a surjectivity or residual condition sufficient for
an inverse-limit reflected recovery.

\subsection{Parent record}

This V36 note appends to the validated thirteenth-cycle V35 file with
record
\[
\begin{array}{c|c}
\text{lines}&11255\\
\text{bytes}&397752\\
\text{SHA--256}&
\texttt{105C9226C634FC9E2D27F7646DA551223DBF7C512EEE2DBF01112DFD5FDEB2EF}.
\end{array}
\]

% =============================================================================
% END APPENDED THIRTEENTH-CYCLE V36 MATERIAL
% =============================================================================

% =============================================================================
% APPENDED THIRTEENTH-CYCLE V37 MATERIAL
% =============================================================================

\section{Thirteenth-cycle V37: cofinal reflected Gram lifts and the
inverse-limit recovery}
\label{sec:v13v37}

\subsection{The extension problem}

V36 gives a lifted spectrahedron at every finite positive-time
ledger.  Independent nonemptiness of those spectrahedra does not
show that a particular low-level recovery can be extended.  The
correct object is the extension fiber, and the correct continuum
statement is an inverse-limit theorem.

Let \(S_n\) be the set of lifted variables
\[
 z_n=(J_n,W_n,T_n,\ldots)                                      \tag{37.1}
\]
satisfying the finite constraints through level \(n\).  Impose a
finite Gram gauge
\[
 {\rm Tr}\,W_n\leq M_n                                         \tag{37.2}
\]
and analogous finite bounds on auxiliary epigraph variables.  By
V36, each \(S_n\) is then a compact convex spectrahedron.  Let
\[
 \pi_n:S_{n+1}\longrightarrow S_n                              \tag{37.3}
\]
be the combined restriction: Choi coarse graining on the operator
system, principal compression of \(W_{n+1}\), and restriction of
the auxiliary variables.

The statement that (37.3) is a map into \(S_n\) is already a
substantive compatibility requirement.  It says that every
higher-level feasible recovery remains feasible after restriction.

\subsection{Extension fibers}

For \(z\in S_n\), define
\[
 {\rm Ext}_{n+1}(z)
 :=
 \{y\in S_{n+1}:\pi_n(y)=z\}.                                  \tag{37.4}
\]

\begin{proposition}[Finite extension is an SDP]
\label{prop:v13v37-extension-sdp}
For fixed \(z\in S_n\), the extension fiber (37.4) is compact and
convex.  It is itself a spectrahedron, obtained by adjoining the
affine equality \(\pi_n(y)=z\) to the level-\(n+1\) lifted
constraints.  Hence extendability of a certified finite recovery is
a finite semidefinite feasibility question.
\end{proposition}

\begin{proof}
The fiber is the intersection of the compact convex set \(S_{n+1}\)
with an affine subspace.  It is therefore compact and convex.
All constraints defining \(S_{n+1}\), including the Gram lift, are
affine equalities or LMIs by V36; the additional restriction
equality is affine.
\end{proof}

The bonding map is surjective precisely when
\({\rm Ext}_{n+1}(z)\neq\varnothing\) for every \(z\in S_n\).
Surjectivity is sufficient to continue every chosen anchor.  It is
stronger than what is needed merely to construct some continuum
recovery.

\subsection{Exact inverse limits}

\begin{theorem}[Nonempty compact inverse limit]
\label{thm:v13v37-inverse-limit}
Suppose every \(S_n\) is nonempty and compact and every
\(\pi_n:S_{n+1}\to S_n\) is continuous.  Then
\[
 S_\infty
 :=
 \{(z_1,z_2,\ldots)\in\prod_{n\geq1}S_n:
   \pi_n(z_{n+1})=z_n\ \hbox{for all }n\}                       \tag{37.5}
\]
is nonempty and compact.
\end{theorem}

\begin{proof}
The product \(\prod_nS_n\) is compact by Tychonoff's theorem.
For each \(n\), the compatibility set
\[
 C_n=\{(z_k)_k:\pi_n(z_{n+1})=z_n\}                             \tag{37.6}
\]
is closed.  The family \(\{C_n\}\) has the finite-intersection
property.  Indeed, for constraints through \(C_N\), choose any
\(z_{N+1}\in S_{N+1}\), define
\[
 z_N=\pi_N(z_{N+1}),\quad
 z_{N-1}=\pi_{N-1}(z_N),\quad\ldots,\quad
 z_1=\pi_1(z_2),                                                \tag{37.7}
\]
and choose arbitrary coordinates above \(N+1\).  Thus every finite
intersection is nonempty.  Compactness gives a point in
\(\bigcap_nC_n\), proving nonemptiness.  The same intersection is
closed in a compact product, hence compact.
\end{proof}

Surjectivity of the bonding maps is not required for nonemptiness
of (37.5).  It is required if one wants the coordinate projection
\(S_\infty\to S_n\) to cover every point of \(S_n\).  This
distinction prevents an unnecessary continuum hypothesis.

\begin{corollary}[Continuation of a prescribed anchor]
\label{cor:v13v37-anchor}
Fix \(z_m^\circ\in S_m\).  If every iterated extension fiber above
\(z_m^\circ\) is nonempty, equivalently if for every \(N>m\) there
is \(z_N\in S_N\) whose restriction to level \(m\) is
\(z_m^\circ\), then \(z_m^\circ\) belongs to a compatible family in
\(S_\infty\).
\end{corollary}

\begin{proof}
Intersect the compact product in the proof of
Theorem~\ref{thm:v13v37-inverse-limit} with the additional closed
condition \(z_m=z_m^\circ\).  The assumed finite-level extensions
give the finite-intersection property.
\end{proof}

\subsection{Asymptotically compatible arrays}

Exact restriction can be too rigid at a regulator cutoff.  The
closed-limit argument only needs compatibility residuals that
vanish at each fixed level.

Let \(K_n\) be compact metric supersets of \(S_n\), with continuous
maps \(\pi_n:K_{n+1}\to K_n\).  For every terminal cutoff \(N\),
suppose an array
\[
 z_n^{(N)}\in K_n,\qquad 1\leq n\leq N                         \tag{37.8}
\]
is available.

\begin{theorem}[Cofinal approximate inverse limit]
\label{thm:v13v37-approx-limit}
Assume that, for every fixed \(n\),
\[
 {\rm dist}(z_n^{(N)},S_n)\longrightarrow0,\qquad
 d_n\bigl(z_n^{(N)},\pi_n(z_{n+1}^{(N)})\bigr)
 \longrightarrow0                                             \tag{37.9}
\]
as \(N\to\infty\).  Then there is a cofinal subsequence
\(N_k\to\infty\) and points \(z_n\in S_n\) such that
\[
 z_n^{(N_k)}\longrightarrow z_n,\qquad
 \pi_n(z_{n+1})=z_n                                            \tag{37.10}
\]
for every fixed \(n\).
\end{theorem}

\begin{proof}
Compactness of \(K_1\) gives a convergent subsequence for the first
coordinate.  From it, compactness of \(K_2\) gives a further
subsequence for the second coordinate.  Continue and take the
Cantor diagonal.  Along the diagonal, every fixed coordinate
converges.  The first limit in (37.9) and closedness of \(S_n\)
place the limit \(z_n\) in \(S_n\).  Continuity of \(\pi_n\) and
the second limit in (37.9) give
\[
 d_n(z_n,\pi_n(z_{n+1}))=0.
\]
This proves (37.10).
\end{proof}

No summability over all \(n\) is needed for this existence theorem.
Pointwise vanishing at every fixed ledger suffices.  Summability is
useful only when one wants a quantitative comparison to a chosen
finite anchor without extracting a new subsequence.

\subsection{The principal Gram restriction}

Choose nested bases
\[
 E_{n,+}=\operatorname{span}\{F_1,\ldots,F_{s_n}\},\qquad
 H_{n,+}=\operatorname{span}\{H_1,\ldots,H_{t_n}\},             \tag{37.11}
\]
with \(s_n,t_n\) increasing.  Let \(P_n\) select the pair indices
\((u,p)\) present at level \(n\).  The Gram restriction is
\[
 W_n=P_nW_{n+1}P_n^*.                                          \tag{37.12}
\]
It preserves positivity and decreases trace:
\[
 W_{n+1}\succeq0
 \ \Longrightarrow\
 W_n\succeq0,\qquad
 {\rm Tr}\,W_n\leq{\rm Tr}\,W_{n+1}.                            \tag{37.13}
\]
The factorization equality also restricts from level \(n+1\) to
level \(n\), provided the products \(A_{pq}\), \(B_{uv}\), and
Choi coarse graining agree literally on the nested systems.

\begin{proposition}[Projective closure of the reflected lift]
\label{prop:v13v37-gram-projective}
Under the nesting hypothesis above, principal Gram compression
preserves the V36 reflected-factorization equations.  If every
other LMI and affine constraint is also stable under the declared
Choi and auxiliary restrictions, then \(\pi_n(S_{n+1})\subset
S_n\).
\end{proposition}

\begin{proof}
For \(p,q\leq s_n\), the level-\(n+1\) equality reads
\[
 \Psi_{J_{n+1}}(A_{pq})
 =
 \sum_{u,v\leq t_{n+1}}
 (W_{n+1})_{(u,p),(v,q)}B_{uv}.                                \tag{37.14}
\]
The projective hypothesis on the product systems requires the
terms outside \(H_{n,+}\) either to be absent under restriction or
to be absorbed by the declared coarse-graining map.  With that
hypothesis, (37.14) becomes
\[
 \Psi_{J_n}(A_{pq})
 =
 \sum_{u,v\leq t_n}
 (W_n)_{(u,p),(v,q)}B_{uv},                                    \tag{37.15}
\]
which is the level-\(n\) equation.  Positivity follows from
(37.13).  The assumed stability handles the remaining blocks.
\end{proof}

The qualification in this proposition is essential.  Principal
compression of \(W\) alone does not delete algebra-valued
contributions involving newly added \(H_u\).  A valid nested scheme
must make the product-system restriction commute with the Gram
restriction.

\subsection{Construction of the quasi-local recovery adjoint}

Let
\[
 {\cal E}_1\subset{\cal E}_2\subset\cdots,\qquad
 {\cal E}_{\rm alg}:=\bigcup_n{\cal E}_n                       \tag{37.16}
\]
be unital self-adjoint finite-dimensional operator systems whose
norm closure is \({\cal A}_{\rm out}\).  Suppose a point of
\(S_\infty\) supplies UCP maps
\[
 \Psi_n:{\cal E}_n\longrightarrow{\cal A}_{\rm in}             \tag{37.17}
\]
with exact restriction compatibility
\[
 \Psi_{n+1}|_{{\cal E}_n}=\Psi_n.                              \tag{37.18}
\]

\begin{theorem}[Quasi-local OS-preserving recovery]
\label{thm:v13v37-quasilocal}
Equations (37.17)--(37.18) define a unique UCP contraction
\[
 \Psi:{\cal A}_{\rm out}\longrightarrow{\cal A}_{\rm in}       \tag{37.19}
\]
provided the codomain is norm closed.  If the reflected
factorization constraint holds on every finite positive-time
ledger, then, for every OS-positive input state
\(\omega_{\rm in}\),
\[
 \omega_{\rm rec}:=\omega_{\rm in}\circ\Psi                   \tag{37.20}
\]
is OS positive on the norm closure of the positive-time algebra.
\end{theorem}

\begin{proof}
Compatibility makes \(\Psi_0\) well defined on
\({\cal E}_{\rm alg}\).  Every UCP map is contractive, so
\(\Psi_0\) extends uniquely by continuity to (37.19).
For each matrix level and each positive element over
\({\cal E}_{\rm alg}\), positivity is witnessed in some
\({\cal E}_n\); it passes to the norm closure.  Thus the extension
is UCP.

Every \(F\) in the algebraic positive-time union belongs to one
finite ledger.  The finite reflected factorization and
Theorem~\ref{thm:v13v34-recovery-os} give
\[
 \omega_{\rm rec}(\Theta_{\rm out}(F)F)\geq0.                  \tag{37.21}
\]
For a norm limit \(F_k\to F\), the product convergence estimate
(34.26) and continuity of the state pass (37.21) to \(F\).
\end{proof}

The theorem needs no globally convergent list of factor maps
\(L_\alpha\).  Compatible finite Gram certificates suffice because
OS positivity is tested on finite families.

\subsection{Normality and Markov recovery}

The norm extension in Theorem~\ref{thm:v13v37-quasilocal} need not
be normal in a chosen von Neumann representation.  Normality is
supplied by the independent predual compactness theorem of V29.
At each fixed local region, impose the V33 energy tail and the V29
predual convergence on the Choi/recovery component of the same
cofinal array.  The diagonal can be synchronized with
Theorem~\ref{thm:v13v37-approx-limit} because the ledgers and
regions are countable.

\begin{corollary}[Locally normal OS-Markov inverse limit]
\label{cor:v13v37-normal}
Suppose the cofinal arrays satisfy:
\begin{enumerate}
\item the feasibility and compatibility residuals (37.9);
\item the V33 local energy-tail and background-modulus bounds;
\item vanishing V29 recovery residuals;
\item the V34 complete OS residuals on every fixed ledger.
\end{enumerate}
Then a single diagonal subsequence yields a locally normal
quasi-local recovery whose recovered state is both OS positive and
an exact local Markov recovery state.
\end{corollary}

\begin{proof}
Use Theorem~\ref{thm:v13v37-approx-limit} for compatibility, V29
and V33 for locally normal predual convergence and exact recovery,
and Theorem~\ref{thm:v13v34-limit} for OS positivity.  A diagonal
over the countable product of ledgers and regions enforces all
fixed-index conditions simultaneously.
\end{proof}

\subsection{Background and loop compatibility}

When \(S_n=S_n(x)\) varies with the Einstein background, V36
constructs local path transports inside strict feasibility
margins.  To obtain a section \(x\mapsto(z_n(x))_n\), the
background transport and the ledger restriction must commute:
\[
 \pi_n\bigl({\cal P}_{n+1,\gamma}z\bigr)
 =
 {\cal P}_{n,\gamma}\bigl(\pi_nz\bigr)                         \tag{37.22}
\]
for every certified path \(\gamma\).  Around a loop, equality of
the returned lift is stronger than necessary; equality after the
declared groupoid action suffices.  Failure of (37.22) is a
projective holonomy obstruction and cannot be repaired by taking
independent trace-minimizing lifts.

\begin{assessment}[Status after thirteenth-cycle V37]
\label{ass:v13v37-status}
The finite reflected spectrahedra now have an exact cofinal
interface.  Extension of a chosen recovery is an SDP fiber;
nonempty compact projective systems possess compatible families;
and asymptotically compatible cutoff arrays converge after a
diagonal.  Compatible UCP maps extend to a quasi-local recovery,
while finite Gram certificates suffice to prove OS positivity on
the completed positive-time algebra.  V29--V34 then add local
normality and exact Markov recovery.
\end{assessment}

\begin{firewall}[Claim boundary after thirteenth-cycle V37]
\label{fire:v13v37-claim}
No RPESM finite feasible set, extension fiber, nested product-system
restriction, Gram gauge, cofinal array, background transport, or
loop descent has been populated.  Independent finite-level
feasibility does not preserve a prescribed anchor.  The norm
extension is not automatically normal.  No clustering, analytic
continuation, determinant cancellation, Einstein measure, or full
SM--Einstein continuum is proved.
\end{firewall}

\subsection{Next bottleneck}

The remaining obstruction is simultaneous background and
projective descent.  The next note should identify the holonomy of
the convex extension fibers, prove when barycentric or
minimum-norm selection commutes with restriction, and isolate a
computable loop residual whose vanishing yields one
background-continuous inverse-limit section.

\subsection{Parent record}

This V37 note appends to the validated thirteenth-cycle V36 file with
record
\[
\begin{array}{c|c}
\text{lines}&11723\\
\text{bytes}&413853\\
\text{SHA--256}&
\texttt{BF738F50A246330E07A0D12A6CABF01E014F4E58A4A47CDDFE496B6799014EBC}.
\end{array}
\]

% =============================================================================
% END APPENDED THIRTEENTH-CYCLE V37 MATERIAL
% =============================================================================

% =============================================================================
% APPENDED THIRTEENTH-CYCLE V38 MATERIAL
% =============================================================================

\section{Thirteenth-cycle V38: compact holonomy removal for the
background-continuous inverse limit}
\label{sec:v13v38}

\subsection{The remaining descent obstruction}

V36 transports a strictly feasible lifted recovery along a
background path, and V37 constructs compatible recovery families
over increasing ledgers.  A path from a base background to \(x\)
need not give the same lift as another path.  Their discrepancy is
the holonomy of the feasible convex fiber.

This note proves a fixed-point theorem that removes this discrepancy
when the holonomy has compact closure and acts affinely.  It also
shows why independent minimum-norm choices are not automatically
projective.

\subsection{The compatible inverse-limit fiber}

Fix a base background \(x_0\).  Let
\[
 {\cal S}_\infty(x_0)
 =
 \varprojlim\{S_n(x_0),\pi_n\}                                 \tag{38.1}
\]
be the nonempty compact convex inverse limit from V37.  Convexity
is coordinatewise.  Let \({\cal H}\) be a compact group acting
continuously and affinely on every \(S_n(x_0)\), and assume
\[
 \pi_n(hz)=h\pi_n(z)
 \qquad(h\in{\cal H},\ z\in S_{n+1}(x_0)).                     \tag{38.2}
\]
Then \({\cal H}\) acts continuously and affinely on
\({\cal S}_\infty(x_0)\).

\begin{theorem}[Compact affine holonomy has a compatible fixed point]
\label{thm:v13v38-fixed-point}
Under the hypotheses above, the fixed-point set
\[
 {\cal S}_\infty(x_0)^{\cal H}
 =
 \{z\in{\cal S}_\infty(x_0):hz=z\ \hbox{for all }h\in{\cal H}\}
                                                                    \tag{38.3}
\]
is nonempty, compact, and convex.
\end{theorem}

\begin{proof}
Choose \(z=(z_n)_n\in{\cal S}_\infty(x_0)\).  For each \(n\), form
the finite-dimensional Bochner integral
\[
 \overline z_n=\int_{\cal H}h z_n\,d\mu(h),                    \tag{38.4}
\]
where \(\mu\) is normalized Haar measure.  Since \(S_n(x_0)\) is
compact and convex, \(\overline z_n\in S_n(x_0)\).  Equivariance
of the restrictions and linearity of the integral give
\[
 \pi_n(\overline z_{n+1})
 =
 \int_{\cal H}h\pi_n(z_{n+1})\,d\mu(h)
 =
 \int_{\cal H}hz_n\,d\mu(h)
 =\overline z_n.                                               \tag{38.5}
\]
Thus \(\overline z=(\overline z_n)_n\) lies in the inverse limit.
Haar invariance gives \(g\overline z_n=\overline z_n\) for every
\(g\in{\cal H}\).  Hence the fixed-point set is nonempty.  It is an
intersection of closed affine equalizers inside a compact convex
set, so it is compact and convex.
\end{proof}

The averaging is performed on the complete lifted variable,
including the Choi and Gram coordinates.  Averaging independently
chosen factor maps would not be well defined when their ranks vary.

\subsection{Flat feasible transport}

Let \(X\) be a connected, locally path-connected background core.
Assume that every piecewise \(C^1\) path
\(\gamma:x\to y\) has an affine homeomorphism
\[
 {\cal P}_\gamma:{\cal S}_\infty(x)
 \longrightarrow{\cal S}_\infty(y)                            \tag{38.6}
\]
with:
\[
 {\cal P}_{\gamma_2\gamma_1}
 ={\cal P}_{\gamma_2}{\cal P}_{\gamma_1},\qquad
 {\cal P}_{\gamma^{-1}}={\cal P}_\gamma^{-1},                  \tag{38.7}
\]
and dependence only on the endpoint-fixed homotopy class.  Assume
also that transport is continuous on sufficiently small
contractible path charts.  Equations (38.6)--(38.7) are the flatness
hypothesis for the feasible bundle.

Loops at \(x_0\) give a representation
\[
 \rho:\pi_1(X,x_0)\longrightarrow
 {\rm AffHomeo}({\cal S}_\infty(x_0)).                          \tag{38.8}
\]
Let \({\cal H}\) be the closure of its image in the topology of
coordinatewise uniform convergence.

\begin{theorem}[Path-independent background recovery]
\label{thm:v13v38-descent}
Suppose \({\cal H}\) is compact, its action is continuous and
affine, and it commutes with every ledger restriction.  Choose
\(z_0\in{\cal S}_\infty(x_0)^{\cal H}\).  For \(x\in X\), set
\[
 z(x)={\cal P}_\gamma z_0                                     \tag{38.9}
\]
for any path \(\gamma:x_0\to x\).  Then (38.9) is independent of
the path, is continuous, and gives a projectively compatible
feasible recovery at every background.
\end{theorem}

\begin{proof}
If \(\gamma_1,\gamma_2:x_0\to x\), then
\(\gamma_2^{-1}\gamma_1\) is a loop at \(x_0\).  By (38.7),
\[
 {\cal P}_{\gamma_1}z_0
 =
 {\cal P}_{\gamma_2}
 {\cal P}_{\gamma_2^{-1}\gamma_1}z_0
 =
 {\cal P}_{\gamma_2}z_0,                                      \tag{38.10}
\]
because \(z_0\) is fixed by the holonomy image.  This proves path
independence.  On a contractible path chart, choose the paths
continuously; the assumed local continuity of transport proves
continuity of \(z(x)\).  Finally, each transport acts on the
inverse-limit fiber and commutes with restrictions, so every
\(z(x)\) is feasible and projectively compatible.
\end{proof}

Combined with V37, the Choi coordinate of \(z(x)\) defines a
background-continuous quasi-local OS-preserving recovery.  Combined
with the V29 predual hypotheses, it is locally normal.

\subsection{Equivariance under the physical groupoid}

Let a symmetry arrow \(g:x\to y\) act affinely on feasible fibers
by \({\cal U}_g\).  If
\[
 {\cal U}_g{\cal P}_\gamma
 =
 {\cal P}_{g\gamma}{\cal U}_g                                 \tag{38.11}
\]
and the base fixed point is invariant under the isotropy action,
then
\[
 z(y)={\cal U}_g z(x).                                         \tag{38.12}
\]

\begin{proposition}[Simultaneous path and symmetry descent]
\label{prop:v13v38-groupoid}
If the compact group generated by background holonomy and base
isotropy acts continuously and affinely on
\({\cal S}_\infty(x_0)\), preserving all finite constraints, then
one Haar average produces a base point for which both (38.9) and
(38.12) hold.
\end{proposition}

\begin{proof}
Apply Theorem~\ref{thm:v13v38-fixed-point} to the compact generated
group.  Path independence follows from
Theorem~\ref{thm:v13v38-descent}.  For an arrow \(g:x\to y\),
transport both sides of (38.12) back to the base.  Relation (38.11)
identifies their discrepancy with an element of the generated
base action, which fixes the averaged point.
\end{proof}

The determinant-line phase is not an affine action on the recovery
spectrahedron unless it has already been canceled or lifted to an
honest action.  The projective anomaly isolated in V22--V23
therefore remains a separate gate.

\subsection{Why independent minimum-norm lifts can fail}

Consider
\[
 S_2=\{(x,y)\in{\mathbb R}^2:-1\leq x\leq1,\ y=x-1\},\qquad
 \pi(x,y)=x.                                                    \tag{38.13}
\]
Then \(S_1=\pi(S_2)=[-1,1]\).  The unique minimum-norm point of
\(S_1\) is \(0\), while the squared norm on \(S_2\) is
\[
 x^2+(x-1)^2,                                                   \tag{38.14}
\]
whose unique minimizer is \(x=1/2\).  Hence
\[
 \pi\bigl(\operatorname*{argmin}_{S_2}\|\cdot\|\bigr)
 \neq
 \operatorname*{argmin}_{S_1}\|\cdot\|.                         \tag{38.15}
\]
Both feasible sets are compact and convex.  Strict convexity and
uniqueness at each level do not force projective compatibility.

A correct canonical selection needs a dynamic-programming
compatibility.  Let \(\Phi_n:S_n\to{\mathbb R}\) be strictly convex
continuous costs and put
\[
 m_n(z)
 :=
 \min\{\Phi_{n+1}(y):y\in S_{n+1},\ \pi_n(y)=z\}.               \tag{38.16}
\]

\begin{proposition}[Cost criterion for compatible minimizers]
\label{prop:v13v38-cost}
Suppose \(\pi_n\) is surjective and
\[
 m_n(z)=a_n\Phi_n(z)+b_n,\qquad a_n>0,                          \tag{38.17}
\]
for constants independent of \(z\).  If \(z_{n+1}^\star\)
minimizes \(\Phi_{n+1}\) on \(S_{n+1}\), then
\[
 \pi_n(z_{n+1}^\star)=z_n^\star,                               \tag{38.18}
\]
where \(z_n^\star\) is the unique minimizer of \(\Phi_n\).
\end{proposition}

\begin{proof}
For \(z=\pi_n(z_{n+1}^\star)\),
\[
 m_n(z)\leq\Phi_{n+1}(z_{n+1}^\star).
\]
Conversely the global minimum of \(\Phi_{n+1}\) equals
\(\min_{w\in S_n}m_n(w)\), by surjectivity.  Thus \(z\) minimizes
\(m_n\).  Equation (38.17) says that it uniquely minimizes
\(\Phi_n\), proving (38.18).
\end{proof}

Condition (38.17) is directly testable only after the extension
fibers are understood.  It must not be silently assumed for
trace-minimizing Gram lifts.

\subsection{A finite-holonomy residual}

Suppose a finite holonomy group \(H\) is generated by a symmetric
set \({\cal G}\).  Let \(L_H\) be the maximum word length of an
element of \(H\) in these generators.  For \(z\) in a normed affine
realization of a convex feasible fiber, assume
\[
 \max_{g\in{\cal G}}\|gz-z\|\leq\varepsilon.                   \tag{38.19}
\]
Let
\[
 \overline z={1\over|H|}\sum_{h\in H}hz.                       \tag{38.20}
\]

\begin{proposition}[Quantitative finite-loop correction]
\label{prop:v13v38-loop-residual}
The point \(\overline z\) is feasible and \(H\)-fixed, and
\[
 \|\overline z-z\|\leq L_H\varepsilon.                          \tag{38.21}
\]
If an LMI block at \(z\) has margin \(\gamma\) and is affine with
operator norm at most \(L_{\cal B}\) as a function of \(z\), then
the block at \(\overline z\) retains margin at least
\[
 \gamma-L_{\cal B}L_H\varepsilon.                              \tag{38.22}
\]
\end{proposition}

\begin{proof}
Feasibility and invariance follow from convexity and finite-group
averaging.  Write any \(h\) as a word of length at most \(L_H\).
Because the action is isometric in the chosen norm, a telescoping
sum gives
\(\|hz-z\|\leq L_H\varepsilon\).  Average this inequality to obtain
(38.21).  The affine block changes in operator norm by at most
\(L_{\cal B}\|\overline z-z\|\); Weyl's inequality gives (38.22).
\end{proof}

For a non-isometric action, each telescoping term acquires the
operator norm of the preceding word.  Those amplification constants
must be included explicitly.

\subsection{Proper groupoids and compact cutoffs}

If holonomy is not compact, Haar averaging is unavailable.  A
proper groupoid with a normalized cutoff can still be averaged as
in V21, but the cutoff must act on the complete inverse-limit
fiber, commute with the restrictions, and preserve the reflection
chart.  These are additional hypotheses.  An amenable mean alone
may produce a bidual-valued fixed point and does not automatically
preserve the locally normal predual class.

\subsection{Combined background-continuum theorem}

\begin{theorem}[Conditional descended OS-Markov recovery]
\label{thm:v13v38-combined}
Assume:
\begin{enumerate}
\item the compact feasible inverse-limit fibers of V37 are
nonempty over the background core;
\item V36 path transport is flat, continuous, affine, and
restriction-compatible;
\item the combined holonomy and isotropy action has compact closure
and preserves the full lifted spectrahedra;
\item the V29 and V33 local predual compactness, energy, and
recovery-residual hypotheses hold.
\end{enumerate}
Then there is a background-continuous, symmetry-equivariant,
quasi-local recovery field whose recovered state is locally normal,
OS positive, and satisfies the exact local Markov recovery
identities.
\end{theorem}

\begin{proof}
Theorems~\ref{thm:v13v38-fixed-point} and
\ref{thm:v13v38-descent} give a path-independent compatible
section.  Proposition~\ref{prop:v13v38-groupoid} gives symmetry
equivariance.  Theorem~\ref{thm:v13v37-quasilocal} gives the
quasi-local UCP map and OS positivity.  Corollary
\ref{cor:v13v37-normal} gives local normality and the exact Markov
identities.
\end{proof}

\begin{assessment}[Status after thirteenth-cycle V38]
\label{ass:v13v38-status}
The background-loop obstruction has an exact convex resolution.
Compact affine holonomy acting on the full compatible
Choi--Gram inverse limit has a fixed point, and parallel transport
of that point gives a path-independent continuous recovery field.
The same average can enforce compact isotropy equivariance.
Independent minimum-norm lifts are proved insufficient without a
dynamic-programming identity, and finite loop residuals have an
explicit margin cost.
\end{assessment}

\begin{firewall}[Claim boundary after thirteenth-cycle V38]
\label{fire:v13v38-claim}
No RPESM flat feasible transport, compact holonomy closure,
restriction-compatible action, invariant strict margin, or proper
groupoid cutoff has been constructed.  Projective determinant-line
phases must first satisfy the V22--V23 anomaly gate.  The conditional
theorem does not establish clustering, OS analytic continuation,
an Einstein measure, or the full SM--Einstein continuum.
\end{firewall}

\subsection{Next bottleneck}

With conditional OS-positive recovery and background descent in
place, the next state-level gate is quantitative time separation.
One must connect a certified transfer operator or Euclidean
time-translation contraction to clustering, while keeping the
recovery and reflection structures.  Before adding a new theorem,
the next note should inspect the existing manuscript for any
already-proved OS reconstruction or clustering result and attack
only the missing quantitative implication.

\subsection{Parent record}

This V38 note appends to the validated thirteenth-cycle V37 file with
record
\[
\begin{array}{c|c}
\text{lines}&12115\\
\text{bytes}&428399\\
\text{SHA--256}&
\texttt{2ABCA6F5A297BB8605387A0D5ED09D56E1E3AAA7C6AF30928588B5BB8DDDA66C}.
\end{array}
\]

% =============================================================================
% END APPENDED THIRTEENTH-CYCLE V38 MATERIAL
% =============================================================================

% =============================================================================
% APPENDED THIRTEENTH-CYCLE V39 MATERIAL
% =============================================================================

\section{Thirteenth-cycle V39: the OS recovery isometry and the
vacuum-multiplicity gate}
\label{sec:v13v39}

\subsection{Interface with the existing reconstruction}

The Grand Tensor manuscript already reconstructs source-cyclic OS
Hamiltonians and compressed joint spectral measures from complete
translated kernels.  It also records that a source-cyclic gap,
vacuum uniqueness, and full clustering are separate claims.  This
note does not repeat that reconstruction.  It supplies the missing
map induced by a reflection-factorized recovery and identifies the
additional zero-energy multiplicity created by the factor index.

Let the hypotheses of Theorem~\ref{thm:v13v34-recovery-os} hold.
Write
\[
 {\cal H}_{\rm in}^{\rm OS}
 =
 \overline{{\cal A}_{{\rm in},+}/{\cal N}_{\rm in}},
 \qquad
 {\cal H}_{\rm out}^{\rm OS}
 =
 \overline{{\cal A}_{{\rm out},+}/{\cal N}_{\rm out}},          \tag{39.1}
\]
where
\[
 \langle[H],[K]\rangle_{\rm in}
 =
 \omega_{\rm in}(\Theta_{\rm in}(H)K),\qquad
 \langle[F],[G]\rangle_{\rm out}
 =
 \omega_{\rm rec}(\Theta_{\rm out}(F)G).                       \tag{39.2}
\]
The null spaces are the zero seminorm vectors.

\subsection{The recovery isometry}

Suppose the recovery adjoint has \(r\) factor maps \(L_\alpha\) as
in (34.6).  On the algebraic OS quotient define
\[
 {\cal V}[F]
 :=
 \bigl([L_1F]_{\rm in},\ldots,[L_rF]_{\rm in}\bigr)
 \in({\cal H}_{\rm in}^{\rm OS})^{\oplus r}.                   \tag{39.3}
\]

\begin{theorem}[OS recovery isometric embedding]
\label{thm:v13v39-isometry}
The map (39.3) is well defined and extends uniquely to an isometry
\[
 {\cal V}:{\cal H}_{\rm out}^{\rm OS}
 \longrightarrow({\cal H}_{\rm in}^{\rm OS})^{\oplus r}.       \tag{39.4}
\]
For all positive-time \(F,G\),
\[
 \langle{\cal V}[F],{\cal V}[G]\rangle
 =
 \langle[F],[G]\rangle_{\rm out}.                              \tag{39.5}
\]
\end{theorem}

\begin{proof}
The reflected factorization gives
\[
\begin{split}
 \langle[F],[G]\rangle_{\rm out}
 &=
 \omega_{\rm in}\left(
 \Psi(\Theta_{\rm out}(F)G)\right)\\
 &=
 \sum_{\alpha=1}^r
 \omega_{\rm in}\left(
 \Theta_{\rm in}(L_\alpha F)L_\alpha G\right),
                                                                    \tag{39.6}
\end{split}
\]
which is (39.5).  In particular,
\[
 \|[F]\|_{\rm out}^2
 =
 \sum_{\alpha=1}^r\|[L_\alpha F]\|_{\rm in}^2.                 \tag{39.7}
\]
If \([F]=0\), every nonnegative term on the right vanishes, so
(39.3) is independent of the representative.  Equation (39.7)
makes the map isometric, hence it extends to the completions.
\end{proof}

This is stronger than state-level OS positivity.  It identifies the
entire recovered OS carrier as a closed subspace of a finite
multiplicity amplification of the input carrier.

At increasing finite ledgers, the Gram lift \(W_n\) of V36 gives
the finite version of (39.3) without choosing a canonical set of
factor maps.  A choice of square root of \(W_n\) changes
\({\cal V}_n\) only by a partial isometry on the factor index, so
all compressed inner products remain unchanged.

\subsection{Time-translation intertwining}

Let \(\tau_{{\rm in},t}\) and \(\tau_{{\rm out},t}\), \(t\geq0\),
be positive-time translations preserving the respective OS null
spaces.  Let
\[
 T_{\rm in}(t)[H]=[\tau_{{\rm in},t}H],\qquad
 T_{\rm out}(t)[F]=[\tau_{{\rm out},t}F]                       \tag{39.8}
\]
be the reconstructed self-adjoint contraction semigroups.

Assume the factor maps intertwine time translation modulo the input
OS null space:
\[
 [L_\alpha\tau_{{\rm out},t}F]_{\rm in}
 =
 T_{\rm in}(t)[L_\alpha F]_{\rm in}
 \quad(\alpha=1,\ldots,r).                                     \tag{39.9}
\]

\begin{theorem}[Recovery-semigroup intertwiner]
\label{thm:v13v39-semigroup}
Under (39.9),
\[
 {\cal V}T_{\rm out}(t)
 =
 T_{\rm in}(t)^{\oplus r}{\cal V}
 \qquad(t\geq0).                                                \tag{39.10}
\]
The range of \({\cal V}\) reduces
\(T_{\rm in}(t)^{\oplus r}\), and
\[
 H_{\rm out}
 =
 {\cal V}^*H_{\rm in}^{\oplus r}{\cal V}                       \tag{39.11}
\]
in the sense of the restricted self-adjoint generators.
Consequently
\[
 \sigma(H_{\rm out})
 \subset\sigma(H_{\rm in}).                                    \tag{39.12}
\]
\end{theorem}

\begin{proof}
Equation (39.9) in every component is exactly (39.10) on the dense
positive-time quotient, hence on the completion.  The range of an
isometry is closed.  It is invariant under the self-adjoint
operator \(T_{\rm in}(t)^{\oplus r}\); for a bounded self-adjoint
operator, an invariant closed subspace has invariant orthogonal
complement, so the range reduces every semigroup element.  The
restricted semigroup is unitarily equivalent through \({\cal V}\)
to \(T_{\rm out}(t)\).  Uniqueness of nonnegative self-adjoint
semigroup generators gives (39.11), and spectral restriction gives
(39.12).
\end{proof}

The spectrum inclusion concerns the full input spectrum, repeated
\(r\) times.  Multiplicity does not change spectral values, but it
does enlarge the zero-energy subspace.

\subsection{The zero-energy multiplicity obstruction}

Let
\[
 P_{{\rm in},0}=\mathbf 1_{\{0\}}(H_{\rm in}),\qquad
 Q_0=P_{{\rm in},0}^{\oplus r}.                                \tag{39.13}
\]
The output zero-energy projection satisfies
\[
 {\cal V}P_{{\rm out},0}{\cal V}^*
 =
 Q_0P_{\operatorname{Ran}{\cal V}}
 =
 P_{\operatorname{Ran}{\cal V}}Q_0.                            \tag{39.14}
\]
In particular,
\[
 \dim\ker H_{\rm out}
 =
 \dim\bigl(\operatorname{Ran}{\cal V}
            \cap(\ker H_{\rm in})^{\oplus r}\bigr).             \tag{39.15}
\]

\begin{proposition}[Vacuum-multiplicity gate]
\label{prop:v13v39-vacuum-gate}
Assume the input has a unique vacuum
\(\ker H_{\rm in}={\mathbb C}\Omega_{\rm in}\).  The recovered
theory has a unique vacuum if and only if
\[
 \operatorname{Ran}{\cal V}
 \cap({\mathbb C}\Omega_{\rm in})^{\oplus r}
 =
 {\mathbb C}{\cal V}\Omega_{\rm out}.                          \tag{39.16}
\]
Reflection factorization and time-translation intertwining alone do
not imply (39.16).
\end{proposition}

\begin{proof}
The first assertion is (39.15), noting that
\({\cal V}\Omega_{\rm out}\) is a nonzero zero-energy vector by
(39.10).  For the last assertion, take any input theory with unique
vacuum and let the output OS space be two copies of its vacuum
line, with identity semigroup.  Its inclusion into the first two
vacuum copies satisfies the isometric and intertwining relations
but has two-dimensional zero-energy space.
\end{proof}

This is the recovery counterpart of the source-complement
distinction already present in the Grand Tensor.  A recovered
source carrier can inherit every positive energy bound and still
possess an extra invariant vector in the factor multiplicity.

\subsection{Gap inheritance}

Assume
\[
 H_{\rm in}\succeq m(I-P_{{\rm in},0}),\qquad m>0.              \tag{39.17}
\]

\begin{theorem}[Recovered gap above the full invariant space]
\label{thm:v13v39-gap}
Under Theorem~\ref{thm:v13v39-semigroup} and (39.17),
\[
 H_{\rm out}\succeq m(I-P_{{\rm out},0}),                      \tag{39.18}
\]
or equivalently
\[
 \|T_{\rm out}(t)(I-P_{{\rm out},0})\|
 \leq e^{-mt}.                                                  \tag{39.19}
\]
If the vacuum gate (39.16) holds, then
\(P_{{\rm out},0}=|\Omega_{\rm out}\rangle
\langle\Omega_{\rm out}|\), so the recovered theory has the same
gap lower bound \(m\) above its unique vacuum.
\end{theorem}

\begin{proof}
Functional calculus applied to (39.17) gives
\[
 \|T_{\rm in}(t)(I-P_{{\rm in},0})\|\leq e^{-mt}.               \tag{39.20}
\]
The direct sum obeys the same estimate.  Since
\(\operatorname{Ran}{\cal V}\) reduces the semigroup and its
zero-energy projection is (39.14), restriction through the
isometry gives (39.19).  Spectral calculus makes (39.18)
equivalent.  The final statement uses Proposition
\ref{prop:v13v39-vacuum-gate}.
\end{proof}

\subsection{Recovered connected time correlations}

For positive-time \(F,G\), let
\[
 f=[F]_{\rm out},\qquad g=[G]_{\rm out}.                        \tag{39.21}
\]
The complete invariant-sector subtraction is
\[
 {\cal C}_{F,G}(t)
 :=
 \langle f,T_{\rm out}(t)g\rangle
 -
 \langle f,P_{{\rm out},0}g\rangle.                            \tag{39.22}
\]

\begin{corollary}[Exponential time clustering after the vacuum gate]
\label{cor:v13v39-clustering}
Under (39.17),
\[
 |{\cal C}_{F,G}(t)|
 \leq
 e^{-mt}
 \|(I-P_{{\rm out},0})f\|\,
 \|(I-P_{{\rm out},0})g\|.                                    \tag{39.23}
\]
If (39.16) holds, then the subtraction in (39.22) is the usual
vacuum product:
\[
 \langle f,P_{{\rm out},0}g\rangle
 =
 \langle f,\Omega_{\rm out}\rangle
 \langle\Omega_{\rm out},g\rangle.                             \tag{39.24}
\]
Thus the recovered OS state clusters exponentially in Euclidean
time on this source carrier.
\end{corollary}

\begin{proof}
Insert \(I-P_{{\rm out},0}\) on both sides of the semigroup in
(39.22), use (39.19), and apply Cauchy--Schwarz.  Under the unique
vacuum condition, the zero projection has rank one, giving
(39.24).
\end{proof}

This is source-carrier time clustering.  Spatial clustering, a
full local-net vacuum statement, and Lorentzian particle
interpretation remain stronger claims exactly as in the Grand
Tensor.

\subsection{Quantitative intertwining and vacuum leakage}

For an output vector \(f\), define its amplified input vacuum
leakage
\[
 \ell(f):=\|Q_0{\cal V}f\|.                                    \tag{39.25}
\]
Suppose that at one time \(t\) an approximate intertwining bound
holds:
\[
 \|{\cal V}T_{\rm out}(t)f
   -T_{\rm in}(t)^{\oplus r}{\cal V}f\|
 \leq\varepsilon_t(f).                                        \tag{39.26}
\]

\begin{proposition}[Approximate recovered decay]
\label{prop:v13v39-approx}
Under (39.17),
\[
 \|T_{\rm out}(t)f\|
 \leq
 \ell(f)
 +e^{-mt}\sqrt{\|f\|^2-\ell(f)^2}
 +\varepsilon_t(f).                                            \tag{39.27}
\]
In particular, a uniform exponential estimate to zero on a
centered output subspace requires both vanishing vacuum leakage and
a vanishing intertwining residual.
\end{proposition}

\begin{proof}
Because \({\cal V}\) is an isometry, the left side equals
\(\|{\cal V}T_{\rm out}(t)f\|\).  Apply (39.26) and split
\({\cal V}f\) into its \(Q_0\) and \(I-Q_0\) components.
The semigroup is the identity on the first and has norm at most
\(e^{-mt}\) on the second.  Orthogonality gives the square root in
(39.27), and the triangle inequality adds the residual.
\end{proof}

At finite ledger level, \(\ell(f)^2\) is a Gram quantity and can be
included in the semidefinite certificate once the input vacuum
projection is represented.  Equation (39.26) is a complete
mixed-time Gram residual; it should be bounded before separating
factor indices.

\subsection{Cofinal passage}

Suppose the V37 inverse-limit recovery is available, the
time-translation ledgers are nested, and the mixed-time Gram
residuals vanish at every fixed rational \(t\) and fixed source
ledger.  Strong continuity extends the intertwining from rational
times to all \(t\geq0\).  If the input gap \(m\) is uniform on the
selected local source carrier and the vacuum-leakage Grams vanish
cofinally, Theorem~\ref{thm:v13v39-gap} and Corollary
\ref{cor:v13v39-clustering} pass to the locally normal recovered
carrier.

A gap at each cutoff with \(m_n\downarrow0\) gives no positive
continuum value.  Likewise, a leakage number that is merely small
at each cutoff but has a nonzero limit leaves an invariant
plateau in (39.27).

\begin{assessment}[Status after thirteenth-cycle V39]
\label{ass:v13v39-status}
A reflection-factorized recovery now has an exact OS-Hilbert-space
meaning: it embeds the recovered carrier isometrically into a
multiplicity amplification of the input carrier.  Translation
intertwining makes the recovered Hamiltonian a reducing
restriction of the amplified input Hamiltonian, so the gap above
the complete invariant space is inherited.  Ordinary vacuum
clustering additionally requires the explicit rank-one
vacuum-multiplicity gate (39.16).  Approximate intertwining and
vacuum leakage have the quantitative bound (39.27).
\end{assessment}

\begin{firewall}[Claim boundary after thirteenth-cycle V39]
\label{fire:v13v39-claim}
No RPESM recovery factor maps, mixed-time Gram residual,
time-translation intertwiner, uniform input gap, or vacuum-leakage
certificate has been populated.  The Grand Tensor's source
reconstruction is used as an interface and is not reproved.
The result concerns the recovered source carrier and does not prove
spatial clustering, a full-net unique vacuum, Lorentzian
continuation, a particle spectrum, scattering, confinement, or the
full SM--Einstein continuum.
\end{firewall}

\subsection{Next bottleneck}

The new finite acquisition is the mixed-time Choi--Gram block:
encode (39.9), the vacuum-leakage Gram, and a finite semigroup
contraction margin in the V36 lifted spectrahedron.  The next
mandatory companion audit is V40; after that audit the programme
should continue with this mixed-time feasibility problem unless a
new companion result supplies a more direct route.

\subsection{Parent record}

This V39 note appends to the validated thirteenth-cycle V38 file with
record
\[
\begin{array}{c|c}
\text{lines}&12488\\
\text{bytes}&442145\\
\text{SHA--256}&
\texttt{E1D480FE21117FFB649865CB2AFE1FFB1303A4A40A7764AF497AE62D2B372A18}.
\end{array}
\]

% =============================================================================
% END APPENDED THIRTEENTH-CYCLE V39 MATERIAL
% =============================================================================

% =============================================================================
% APPENDED THIRTEENTH-CYCLE V40 MATERIAL
% =============================================================================

\section{Thirteenth-cycle V40: mandatory companion audit after the
OS recovery isometry}
\label{sec:v13v40}

\subsection{Current companion records}

The required checkpoint inspected the newest main YM and NS notes:
\[
\begin{array}{c|r|r|c}
\text{file}&\text{lines}&\text{bytes}&\text{SHA--256}\\ \hline
\texttt{Renewal\_Yang\_Mills\_Mass\_Gap\_Research\_Notes\_V82.tex}
&28114&958424&
\texttt{D4E3B41C26CF2E1CB05559F7449AE5A00A22955E7A50FB59BC69B80EA9858E0A}\\
\texttt{Renewal\_NS\_Stability\_Research\_Notes\_V26.tex}
&5319&278283&
\texttt{82C887F8C2C69E4F28FAD7C3DC8DE5B9099CB5A4BF431EBA1FFF3E91935978AD}.
\end{array}                                                    \tag{40.1}
\]
The updated auxiliary file
\[
 \texttt{Renewal\_YM\_Parallel\_Research\_Notes\_V5.tex}        \tag{40.2}
\]
was also inspected.  Its current record is \(1349\) lines,
\(54028\) bytes, and SHA--256
\[
 \texttt{F18D56D24CC61D3988082E014AA848DA054FEB2E05561F920A00A53D7C2A069D}.
                                                                    \tag{40.3}
\]
It now includes a terminal failure result for volume-uniform
tightness of an extensive global charge.

\subsection{New YM V82 input}

YM V82 proves a rigorous a posteriori certificate for a complete
matrix path cell.  It uses a rational trial inverse \(R\), a
certified Hessian floor, complete matrix remainders, and rational
exponential ceilings.  It thereby avoids both interval matrix
square roots and trusted floating trigonometric values.  The
theorem is still unpopulated for the YM path bank.

The direct SM use is methodological and exact.  In V36 the
background transport is controlled by complete Hermitian LMI
blocks \({\cal B}_r(t,z(t))\).  For any block whose inverse or
relative metric is needed, a rational trial inverse and the exact
identity
\[
 B(t)^{-1}-R=(I-RB(t))B(t)^{-1}                                \tag{40.4}
\]
give an a posteriori inverse error once a lower floor for \(B(t)\)
is known.  More simply, positivity itself can be certified without
an inverse from
\[
 \widetilde B\succeq\gamma I,\qquad
 \sup_{t\in J}\|B(t)-\widetilde B\|_{\rm op}<\gamma.            \tag{40.5}
\]
V41 should therefore express the mixed-time recovery conditions as
complete affine blocks and give rational residual-to-margin tests
for those blocks.  It should not diagonalize approximate floating
matrices independently.

\subsection{Auxiliary YM failure result}

The auxiliary YM note proves that a nontrivial additive global
integer charge on disjoint components cannot have a
volume-uniform exponential tail.  Linear carrier cost loses to
placement entropy, and total charge is non-tight under
tensorization.  This does not address the local recovery or OS
blocks of V39.  It reinforces the regional formulation already
used in V19, V29, and V33:
\[
 \text{constants may depend on the fixed bounded region }O,
 \quad\text{but must be uniform in cutoff and background there}. \tag{40.6}
\]
V41 will therefore impose vacuum-leakage and mixed-time margins
region by region.  It will not require a uniform bound on a total
infinite-volume energy or charge.

\subsection{NS status}

NS V26 is unchanged.  Its rank-one contraction result continues to
support complete prescribed-input contraction before norm
relaxation, but supplies no recovery semigroup, OS vacuum
projection, Choi matrix, or Standard-Model constant.

\begin{proposition}[Audit decision]
\label{prop:v13v40-decision}
No companion result closes any hypothesis of V39.  The next
substantive SM step is the finite mixed-time spectrahedron:
\begin{enumerate}
\item encode factor-map time intertwining through a joint Gram lift,
without selecting Kraus coordinates;
\item encode vacuum leakage as a positive Gram block;
\item impose rational whole-block residual-to-margin certificates
on each fixed region and background path cell;
\item keep all estimates local in region, uniform only in cutoff
and compact background.
\end{enumerate}
\end{proposition}

\begin{proof}
YM V82 provides a certificate architecture but no SM matrix.
The auxiliary YM result rules out an unrelated global extensive
tail and supports the already chosen regional quantifiers.  NS V26
has not changed.  The four listed items are exactly the unresolved
finite acquisition stated after V39, amended only by the rational
whole-block certification method.
\end{proof}

\begin{assessment}[Status after thirteenth-cycle V40]
\label{ass:v13v40-status}
The compulsory audit is complete.  YM V82 supplies a rigorous
a posteriori rational path-cell pattern, while the auxiliary YM
failure theorem rules out volume-uniform control of a nontrivial
extensive charge.  The SM direction is adjusted accordingly:
certify the full mixed-time recovery blocks with rational margins
on each fixed region and retain regional, rather than total-volume,
energy quantifiers.
\end{assessment}

\begin{firewall}[Claim boundary after thirteenth-cycle V40]
\label{fire:v13v40-claim}
No YM numerical path bound or NS kernel constant is imported.
The rational YM V82 theorem is unpopulated, and the auxiliary
global-charge failure does not prove or disprove a local SM
recovery.  No mixed-time Gram lift, vacuum leakage floor,
background cell card, continuum gap, clustering theorem, or full
SM--Einstein continuum is obtained by this audit.
\end{firewall}

\subsection{Next step}

V41 will append the joint-time Gram-lift theorem and the rational
whole-block certificate.  This continues beyond the audit as
required.

\subsection{Parent record}

This V40 note appends to the validated thirteenth-cycle V39 file with
record
\[
\begin{array}{c|c}
\text{lines}&12903\\
\text{bytes}&455984\\
\text{SHA--256}&
\texttt{9B0F17930C29371F8535765218C3836A388A4EA9E051549DB87F5CA15EF146C9}.
\end{array}
\]

% =============================================================================
% END APPENDED THIRTEENTH-CYCLE V40 MATERIAL
% =============================================================================

% =============================================================================
% APPENDED THIRTEENTH-CYCLE V41 MATERIAL
% =============================================================================

\section{Thirteenth-cycle V41: mixed-time Gram support and rational
spectral cards}
\label{sec:v13v41}

\subsection{Why factor coordinates should be eliminated}

V39 states time intertwining in terms of factor maps
\(L_\alpha\).  A Cholesky factor of the V36 Gram matrix is not
canonical and can rotate discontinuously when its rank changes.
The intertwining condition can instead be imposed on the support of
the positive Gram matrix itself.  This keeps the constraint convex
and projectively meaningful.

Retain the bases \(F_1,\ldots,F_s\) and \(H_1,\ldots,H_t\) of
V36.  Suppose a finite time ledger
\[
 {\cal T}=\{\tau_1,\ldots,\tau_m\}\subset[0,\infty)             \tag{41.1}
\]
leaves their finite positive-time spans invariant modulo the
declared null/restriction relations.  Write
\[
 \tau_{{\rm out},\tau_j}F_p
 =
 \sum_{q=1}^s(A_{{\rm out},j})_{qp}F_q,\qquad
 \tau_{{\rm in},\tau_j}H_u
 =
 \sum_{v=1}^t(A_{{\rm in},j})_{vu}H_v.                         \tag{41.2}
\]

For one factor map, let
\[
 L_\alpha F_p=\sum_{u=1}^t\ell_{\alpha,up}H_u                 \tag{41.3}
\]
and vectorize its coefficient matrix as
\(\ell_\alpha\in{\mathbb C}^{ts}\).  The coefficient form of
time intertwining is
\[
 D_j\ell_\alpha=0,\qquad
 D_j:=
 A_{{\rm out},j}^{\mathsf T}\otimes I_t
 -
 I_s\otimes A_{{\rm in},j}.                                  \tag{41.4}
\]
The order in (41.4) corresponds to column vectorization.

\subsection{The transposed Gram variable}

It is convenient to transpose the V36 Gram convention and put
\[
 \widehat W_{(u,p),(v,q)}
 :=
 \sum_\alpha
 \ell_{\alpha,up}\overline{\ell_{\alpha,vq}}
 =W_{(v,q),(u,p)}.                                             \tag{41.5}
\]
Then \(\widehat W=W^{\mathsf T}\succeq0\), and the reflected
factorization equations become
\[
 \Psi_J(A_{pq})
 =
 \sum_{u,v=1}^t
 \widehat W_{(v,q),(u,p)}B_{uv}.                               \tag{41.6}
\]

\begin{theorem}[Coordinate-free mixed-time Gram support]
\label{thm:v13v41-support}
Let \(\widehat W\succeq0\) satisfy (41.6).  The following are
equivalent.
\begin{enumerate}
\item There is a factorization of \(\widehat W\) by coefficient
vectors \(\ell_\alpha\) such that every corresponding \(L_\alpha\)
obeys all time-intertwining equations (41.4).
\item
\[
 \boxed{\quad D_j\widehat W=0
 \quad(j=1,\ldots,m).\quad}                                    \tag{41.7}
\]
\item
\[
 D_j\widehat W D_j^*=0
 \quad(j=1,\ldots,m).                                          \tag{41.8}
\]
\end{enumerate}
Consequently the finite reflection-factorization and mixed-time
intertwining conditions form one lifted spectrahedron in
\((J,\widehat W)\).
\end{theorem}

\begin{proof}
If
\(\widehat W=\sum_\alpha\ell_\alpha\ell_\alpha^*\) and
\(D_j\ell_\alpha=0\), then (41.7) follows immediately, hence so
does (41.8).

Conversely, positivity gives a square root.  From (41.8),
\[
 0=D_j\widehat W D_j^*
   =(D_j\widehat W^{1/2})(D_j\widehat W^{1/2})^*,               \tag{41.9}
\]
so \(D_j\widehat W^{1/2}=0\).  Every column in any square-root
factorization of \(\widehat W\) therefore lies in
\(\ker D_j\), simultaneously for all \(j\).  Taking those columns
as the \(\ell_\alpha\) proves the first statement.  The same
argument gives \(D_j\widehat W=0\), proving equivalence.  Finally,
(41.6)--(41.7) are affine equalities and
\(\widehat W\succeq0\) is an LMI; adjoining the Choi constraints
from V36 gives a spectrahedron.
\end{proof}

Condition (41.7) involves the complete coefficient support.  It is
unchanged by unitary rotation or rank-changing padding of the
factor index.

\subsection{Approximate support residual}

Let
\[
 {\cal K}:=\bigcap_{j=1}^m\ker D_j,\qquad
 P_{\cal K}\ \hbox{the orthogonal projection onto }{\cal K}.   \tag{41.10}
\]
For \(\widehat W\succeq0\), the exact supported part is
\[
 \widehat W_{\cal K}:=P_{\cal K}\widehat W P_{\cal K}.          \tag{41.11}
\]

\begin{proposition}[Distance to the time-intertwining face]
\label{prop:v13v41-support-distance}
One has
\[
 \|\widehat W-\widehat W_{\cal K}\|_1
 \leq
 2\sqrt{
 {\rm Tr}\widehat W\,
 {\rm Tr}((I-P_{\cal K})\widehat W)}
 +
 {\rm Tr}((I-P_{\cal K})\widehat W).                           \tag{41.12}
\]
If the stacked operator
\[
 D=\begin{pmatrix}D_1\\ \vdots\\D_m\end{pmatrix}                \tag{41.13}
\]
has smallest nonzero singular value \(\sigma_D>0\), then
\[
 {\rm Tr}((I-P_{\cal K})\widehat W)
 \leq
 {\sigma_D^{-2}}\,{\rm Tr}(D\widehat W D^*).                   \tag{41.14}
\]
\end{proposition}

\begin{proof}
Decompose relative to
\({\cal K}\oplus{\cal K}^\perp\):
\[
 \widehat W=
 \begin{pmatrix}A&B\\B^*&C\end{pmatrix},\qquad C\succeq0.       \tag{41.15}
\]
For a positive block matrix,
\(\|B\|_1\leq\sqrt{{\rm Tr}A\,{\rm Tr}C}\).  This follows by
writing \(\widehat W=XX^*\), splitting \(X=(X_0,X_1)^{\mathsf T}\),
and applying the Hilbert--Schmidt ideal inequality to
\(B=X_0X_1^*\).  Since
\(\widehat W-\widehat W_{\cal K}\) has blocks
\(\left(\begin{smallmatrix}0&B\\B^*&C\end{smallmatrix}\right)\),
the triangle inequality gives
\[
 \|\widehat W-\widehat W_{\cal K}\|_1
 \leq2\|B\|_1+{\rm Tr}C,
\]
which implies (41.12).

On \({\cal K}^\perp\), \(D^*D\succeq\sigma_D^2I\).  Hence
\[
 {\rm Tr}(D\widehat W D^*)
 =
 {\rm Tr}(D^*D\widehat W)
 \geq
 \sigma_D^2{\rm Tr}((I-P_{\cal K})\widehat W),                 \tag{41.16}
\]
proving (41.14).
\end{proof}

Thus a scalar positive quantity
\({\rm Tr}(D\widehat W D^*)\), together with a certified singular
floor for the stacked relation matrix, controls trace-norm distance
to the exact mixed-time face.

\subsection{Vacuum leakage in Gram coordinates}

Assume the input OS Hamiltonian has unique normalized vacuum
\(\Omega_{\rm in}\), and define
\[
 a_u:=\langle\Omega_{\rm in},[H_u]_{\rm in}\rangle.             \tag{41.17}
\]
For the factorization encoded by \(\widehat W\), define the
\(s\times s\) vacuum-leakage Gram
\[
 K^{\rm leak}_{pq}(\widehat W)
 :=
 \sum_{u,v=1}^t
 \overline{a_u}a_v\,
 \widehat W_{(v,q),(u,p)}.                                     \tag{41.18}
\]
This is an affine function of \(\widehat W\).

\begin{lemma}[Leakage Gram identity]
\label{lem:v13v41-leakage}
For \(f=\sum_pc_p[F_p]_{\rm out}\),
\[
 c^*K^{\rm leak}(\widehat W)c
 =
 \|Q_0{\cal V}f\|^2.                                          \tag{41.19}
\]
In particular \(K^{\rm leak}\succeq0\).
\end{lemma}

\begin{proof}
The \(\alpha\)-th factor component has vacuum amplitude
\[
 \langle\Omega_{\rm in},
 [L_\alpha\sum_pc_pF_p]_{\rm in}\rangle
 =
 \sum_{p,u}c_p\ell_{\alpha,up}a_u.                             \tag{41.20}
\]
Sum the squared moduli over \(\alpha\) and insert (41.5).  The
result is (41.19), and positivity follows.
\end{proof}

Let \(C\) be a matrix whose columns span a predeclared centered
coefficient subspace of \({\mathbb C}^s\).  The finite
vacuum-multiplicity gate is
\[
 C^*K^{\rm leak}(\widehat W)C=0.                               \tag{41.21}
\]
An approximate gate with leakage at most \(\varepsilon\) is the
LMI
\[
 C^*K^{\rm leak}(\widehat W)C
 \preceq\varepsilon^2 C^*C.                                   \tag{41.22}
\]
When the target recovered state is fixed, \(C\) can be chosen as
the orthogonal complement of its known vacuum-amplitude vector.
If the target state is a decision variable, forming that complement
is nonlinear and must not be hidden inside the SDP claim.

\subsection{A direct complete delayed-Gram certificate}

For a fixed recovered target state and source family, define
\[
 G^0_{pq}=\langle[F_p],[F_q]\rangle_{\rm out},\qquad
 G^\tau_{pq}
 =
 \langle[F_p],T_{\rm out}(\tau)[F_q]\rangle_{\rm out}.          \tag{41.23}
\]
Both are affine functions of the Choi variable when written as
recovered OS expectations of fixed observables.  Let \(G^{\rm
vac}\) be the known rank-one vacuum Gram.

\begin{proposition}[Finite-source contraction LMI]
\label{prop:v13v41-delayed-lmi}
For \(0<q<1\), the complete inequalities
\[
 0\preceq G^\tau-G^{\rm vac}
 \preceq q(G^0-G^{\rm vac})                                   \tag{41.24}
\]
are LMIs.  On the represented source span modulo the entrance
null space, they are equivalent to
\[
 0\preceq T_{\rm out}(\tau)\preceq qI                         \tag{41.25}
\]
on the centered finite source subspace.  The corresponding finite
energy floor is
\[
 H_{\rm out}\succeq-{\log q\over\tau}I                         \tag{41.26}
\]
on that invariant subspace.
\end{proposition}

\begin{proof}
Affineness makes (41.24) an LMI.  For a coefficient vector \(c\),
the two quadratic forms in (41.24) are respectively
\[
 \langle f_c,(T_{\rm out}(\tau)-P_0)f_c\rangle,\qquad
 \langle f_c,(I-P_0)f_c\rangle.                               \tag{41.27}
\]
After quotienting the null space of the second Gram, (41.24) is
exactly (41.25).  Spectral calculus for
\(T_{\rm out}(\tau)=e^{-\tau H_{\rm out}}\) gives (41.26).
\end{proof}

The upper inequality in (41.24) is a source-bank contraction.  A
full vacuum-sector gap still requires the dense-bank supremum
described in the Grand Tensor; no fixed finite bank can replace it.

\subsection{Rational whole-block cards}

The V40 audit calls for a certificate that never trusts a floating
least eigenvalue.  Let \(B(t)=B(t)^*\) be any Choi, Gram, energy,
leakage, or delayed-correlation LMI block on a path cell
\(J\).  Choose a rational Hermitian center matrix
\(\widetilde B\) and rational numbers \(\gamma,r>0\) such that
\[
 \widetilde B-\gamma I\succeq0,\qquad
 \sup_{t\in J}\|B(t)-\widetilde B\|_F\leq r<\gamma.             \tag{41.28}
\]

\begin{theorem}[Rational residual-to-margin certificate]
\label{thm:v13v41-rational-card}
Under (41.28),
\[
 \boxed{\quad B(t)\succeq(\gamma-r)I
 \quad(t\in J).\quad}                                          \tag{41.29}
\]
The first inequality in (41.28) can be certified by an exact
rational \(LDL^*\) factorization with nonnegative pivots.  If all
pivots are positive, it also supplies a strict rational center
margin.
\end{theorem}

\begin{proof}
The Frobenius norm dominates the operator norm, so
\[
 B(t)-\widetilde B\succeq-rI.                                 \tag{41.30}
\]
Add this to
\(\widetilde B\succeq\gamma I\) to obtain (41.29).
An exact \(LDL^*\) factorization is a congruence decomposition;
nonnegative diagonal pivots prove positive semidefiniteness over
the rationals.
\end{proof}

For a complex Gaussian-rational matrix, split it into an exact real
symmetric representation before applying rational \(LDL^{\mathsf
T}\).  Complete matrix remainders should be bounded before using
the Frobenius relaxation.  A factorwise or entrywise triangle
bound is valid but may destroy the strict margin.

\subsection{A posteriori inverse control}

When a relative metric is needed, let \(B(t)\succeq\lambda I\) on
the cell and let \(R\) be a rational trial inverse.  If
\[
 \sup_{t\in J}\|I-RB(t)\|_{\rm op}\leq\eta,                    \tag{41.31}
\]
then the exact identity
\[
 B(t)^{-1}-R=(I-RB(t))B(t)^{-1}                               \tag{41.32}
\]
gives
\[
 \|B(t)^{-1}-R\|_{\rm op}\leq{\eta\over\lambda}.               \tag{41.33}
\]
This is the YM V82 a posteriori mechanism specialized to the SM
LMI blocks.  The floor \(\lambda\) must itself be certified, for
example by Theorem~\ref{thm:v13v41-rational-card}; it cannot be
read from a floating eigenvalue.

\subsection{Regional cofinal certificate}

For each fixed bounded region \(O\), finite source ledger \(s\),
and rational time ledger \({\cal T}_m\), the augmented V36
spectrahedron now includes:
\[
\begin{array}{c|l}
D_j\widehat W=0
 &\text{factor-map time intertwining},\\
C^*K^{\rm leak}C\preceq\varepsilon_{\rm vac}^2C^*C
 &\text{vacuum-multiplicity leakage},\\
0\preceq G^\tau-G^{\rm vac}\preceq
q(G^0-G^{\rm vac})
 &\text{finite-source contraction},\\
B_r\succeq\gamma_r I
 &\text{strict whole-block path margins}.
\end{array}                                                    \tag{41.34}
\]
All are affine equalities or LMIs in the lifted variables once the
background, target state, time actions, and centered source space
are fixed.

\begin{corollary}[Cofinal recovered time clustering]
\label{cor:v13v41-cofinal}
Assume, on every fixed region and source ledger:
\begin{enumerate}
\item the V37 projective compatibility residuals vanish;
\item the support residual in (41.14) and the leakage bound in
(41.22) tend to zero;
\item the rational cards certify a common \(q(\tau)<1\) on a
dense nested source family;
\item the V29 local normality hypotheses hold.
\end{enumerate}
Then the recovered locally normal OS carrier has unique vacuum on
the completed source-cyclic space and
\[
 H_{\rm out}\succeq
 -{\log q(\tau)\over\tau}(I-P_{\Omega_{\rm out}}).             \tag{41.35}
\]
Its connected Euclidean-time correlations obey (39.23) with this
gap.
\end{corollary}

\begin{proof}
The support and leakage limits give the V39 exact intertwiner and
vacuum gate.  Nested dense source contraction with common \(q\)
extends by continuity from the union to the completed centered
source carrier.  Proposition~\ref{prop:v13v41-delayed-lmi} gives
(41.35), and Corollary~\ref{cor:v13v39-clustering} gives the
correlation estimate.  V29 supplies local normality.
\end{proof}

The quantifiers are local in \(O\).  No bound on total
infinite-volume energy or global topological charge is used.

\begin{assessment}[Status after thirteenth-cycle V41]
\label{ass:v13v41-status}
Time-translation compatibility and the vacuum-multiplicity gate
now have exact finite semidefinite forms.  Intertwining is the
support condition \(D_j\widehat W=0\), invariant under all choices
of factor coordinates.  Its trace residual controls distance to
the exact face.  Vacuum leakage is an affine positive Gram block,
and a complete delayed-source LMI gives a finite spectral
contraction.  Rational whole-block cards certify all strict path
margins without trusted floating eigenvalues.
\end{assessment}

\begin{firewall}[Claim boundary after thirteenth-cycle V41]
\label{fire:v13v41-claim}
No RPESM time-action matrices, supported Gram lift, singular floor,
vacuum amplitudes, centered source basis, delayed Gram, rational
cell card, common contraction \(q<1\), or dense source family has
been populated.  A finite source contraction is not a full-net mass
gap.  Spatial clustering, Lorentzian continuation, particle
structure, scattering, confinement, the Einstein measure, and the
full SM--Einstein continuum remain open.
\end{firewall}

\subsection{Next bottleneck}

The recovery and time-translation constraints are now finite convex
objects.  The next upstream obstruction is population on the
physical RPESM branch.  One must extract the finite Choi,
positive-time basis, time-action matrices, target vacuum
amplitudes, and reflected product tables from the existing finite
regulator, then run the first exact feasibility card.  If that
population is unavailable, the correct fallback is to identify the
smallest missing source row rather than add another abstract
compactness theorem.

\subsection{Parent record}

This V41 note appends to the validated thirteenth-cycle V40 file with
record
\[
\begin{array}{c|c}
\text{lines}&13059\\
\text{bytes}&462154\\
\text{SHA--256}&
\texttt{FA9076E530DD71521F35577E879A338A606EB9C1F6D325468F8ED6B9642D6ADD}.
\end{array}
\]

% =============================================================================
% END APPENDED THIRTEENTH-CYCLE V41 MATERIAL
% =============================================================================

% =============================================================================
% APPENDED THIRTEENTH-CYCLE V42 MATERIAL
% =============================================================================

\section{Thirteenth-cycle V42: population from the RPESM
conditional-detail refinement}
\label{sec:v13v42}

\subsection{What the attached regulator already supplies}

The Grand Tensor population theorem provides, at adjacent cutoffs,
a coarse stationary law \(\mu_n\), a conditional detail kernel
\(K_{n+1/n}(dy\mid x)\), and the fine law
\[
 \mu_{n+1}(dx,dy)=\mu_n(dx)K_{n+1/n}(dy\mid x).                 \tag{42.1}
\]
It also provides a reversible coarse transition \(P_n\) and the
fine transition
\[
 \widetilde P((x,y),d(x',y'))
 =
 P_n(x,dx')K_{n+1/n}(dy'\mid x').                             \tag{42.2}
\]
The projection to \(x\) intertwines the transitions exactly.  The
same finite regulator has a stationary reflected path law and the
OS square formula of the reflected-slab theorem.

This is enough to populate an exact commutative recovery channel
and its time action.  The purpose of this note is to state that
interface in operator form and separate it from the unpopulated
noncommutative fermion germ.

\subsection{The conditional-detail recovery channel}

Let \((X,\mu)\) be a probability space, let \(K(dy\mid x)\) be a
Markov kernel into a standard Borel detail space \(Y\), and put
\[
 \nu(dx,dy)=\mu(dx)K(dy\mid x).                                \tag{42.3}
\]
Define
\[
 ({\cal J}f)(x,y)=f(x),\qquad
 ({\cal C}g)(x)=\int_Yg(x,y)K(dy\mid x).                       \tag{42.4}
\]

\begin{proposition}[Exact Bayesian refinement recovery]
\label{prop:v13v42-bayesian}
The map
\[
 {\cal R}:L^1(\mu)\longrightarrow L^1(\nu),\qquad
 ({\cal R}h)(x,y)=h(x)                                        \tag{42.5}
\]
in density coordinates is positive and trace preserving.  Its
Heisenberg adjoint is the unital positive, hence completely
positive, conditional map \({\cal C}\).  Moreover
\[
 {\cal C}{\cal J}=I,\qquad
 {\cal J}^*={\cal C}\quad\hbox{on }L^2,\qquad
 {\cal R}\mu=\nu.                                              \tag{42.6}
\]
Thus the adjacent RPESM detail kernel is an exact state-recovery
channel, not merely an approximate Petz candidate.
\end{proposition}

\begin{proof}
For \(h\in L^1(\mu)\), its recovered signed measure is
\(h(x)\mu(dx)K(dy\mid x)\), so positivity and equality of total
integrals are immediate.  For bounded \(g\),
\[
 \int_{X\times Y}g(x,y)h(x)\,d\nu
 =
 \int_X({\cal C}g)(x)h(x)\,d\mu,                              \tag{42.7}
\]
which identifies the adjoint.  A positive map between commutative
\(C^*\)-algebras is completely positive, and
\({\cal C}1=1\).  Equations (42.6) follow from conditional
integration and (42.3).
\end{proof}

The formula uses densities relative to the declared stationary
laws.  It is the Schr\"odinger map from a coarse state perturbation
to the corresponding conditionally refined perturbation.

\subsection{Reversible fine dynamics as a compression dilation}

Assume \(P\) is a self-adjoint Markov contraction on \(L^2(\mu)\).
Define \(\widetilde P\) from (42.2).  Then
\[
 (\widetilde Pg)(x,y)
 =
 \int_XP(x,dx')\int_Yg(x',y')K(dy'\mid x').                   \tag{42.8}
\]

\begin{theorem}[Exact refined transfer factorization]
\label{thm:v13v42-transfer}
On \(L^2(\nu)\),
\[
 \boxed{\quad\widetilde P={\cal J}P{\cal C}.\quad}              \tag{42.9}
\]
The operator \({\cal J}\) is an isometry,
\({\cal C}={\cal J}^*\), and
\[
 L^2(\nu)=\operatorname{Ran}{\cal J}\oplus\ker{\cal C}.         \tag{42.10}
\]
Relative to this decomposition,
\[
 \widetilde P=
 \begin{pmatrix} {\cal J}P{\cal J}^*&0\\0&0\end{pmatrix}.       \tag{42.11}
\]
In particular \(\widetilde P\) is self-adjoint, its nonzero
spectrum is the nonzero spectrum of \(P\), and every pure-detail
mode in \(\ker{\cal C}\) is killed in one step.
\end{theorem}

\begin{proof}
Equation (42.8) is (42.9).  From (42.3),
\[
 \|{\cal J}f\|_{L^2(\nu)}^2
 =
 \int_X|f(x)|^2\mu(dx),
\]
so \({\cal J}\) is an isometry.  Equation (42.7) with \(L^2\)
functions gives \({\cal J}^*={\cal C}\), and the standard
isometry decomposition gives (42.10).  Now
\({\cal C}{\cal J}=I\), while \({\cal C}\) vanishes on its kernel.
This proves (42.11).  The spectral assertions follow from the
block form.
\end{proof}

\begin{corollary}[Transfer-gap preservation]
\label{cor:v13v42-gap}
If constants are the unique invariant vectors of \(P\) and
\[
 \|P|_{1^\perp}\|\leq q<1,                                    \tag{42.12}
\]
then constants are the unique invariant vectors of
\(\widetilde P\) and
\[
 \|\widetilde P|_{1^\perp}\|\leq q.                            \tag{42.13}
\]
\end{corollary}

\begin{proof}
The constant on \(X\times Y\) is \({\cal J}1\).  The block
decomposition has the coarse centered spectrum on
\({\cal J}(1^\perp)\) and eigenvalue zero on \(\ker{\cal C}\).
Thus the centered norm is at most \(\max\{q,0\}=q\), and no
additional eigenvalue one appears.
\end{proof}

This exact conclusion is stronger than an uncontrolled perturbation
argument.  It uses the special refreshed-detail form (42.2).

\subsection{OS factorization of the recovered path law}

Let \((X_k)_{k\in{\mathbb Z}}\) be the stationary reversible coarse
chain.  Conditional on the entire coarse path, sample the details
\(Y_k\) independently with law \(K(dy_k\mid X_k)\).  Its joint path
law is the stationary path law of (42.2).

For a bounded cylinder function \(F\) of times
\(0,\ldots,m\), define, in the finite-detail case,
\[
 (L_yF)(x_0,\ldots,x_m)
 :=
 K(y\mid x_0)^{1/2}
 \int F((x_0,y),(x_1,y_1),\ldots,(x_m,y_m))
 \prod_{k=1}^mK(dy_k\mid x_k).                                \tag{42.14}
\]

\begin{theorem}[Conditional-detail reflected factorization]
\label{thm:v13v42-path-os}
If the detail set is finite, then for positive-time cylinder
functions \(F,G\),
\[
 \langle[F],[G]\rangle_{\rm fine}^{\rm OS}
 =
 \sum_{y\in Y}
 \langle[L_yF],[L_yG]\rangle_{\rm coarse}^{\rm OS}.             \tag{42.15}
\]
Consequently the fine recovered path law is OS positive and its OS
space embeds isometrically into
\(({\cal H}_{\rm coarse}^{\rm OS})^{\oplus|Y|}\).

For a dominated kernel
\(K(dy\mid x)=k(x,y)\kappa(dy)\), the same statement holds with the
direct integral
\[
 \int_Y
 \langle[L_yF],[L_yG]\rangle_{\rm coarse}^{\rm OS}\,\kappa(dy),
                                                                    \tag{42.16}
\]
where \(K(y\mid x_0)^{1/2}\) in (42.14) is replaced by
\(k(x_0,y)^{1/2}\).
\end{theorem}

\begin{proof}
Condition on the coarse two-sided path and on the shared
time-zero detail \(Y_0=y\).  Given these data, the strictly
positive-time and strictly negative-time detail variables are
independent.  Reversibility and reflection covariance identify the
negative conditional expectation with the complex conjugate of
the reflected positive one.  Multiplication by
\(K(y\mid X_0)^{1/2}\) assigns the time-zero detail weight equally
to the two halves.  Summing over \(y\) gives (42.15).  The
dominated case is the same calculation followed by Tonelli's
theorem.  Each right side is a sum or integral of coarse OS Gram
forms, hence is positive.  The isometry is the construction of
Theorem~\ref{thm:v13v39-isometry}, with a direct integral in the
dominated case.
\end{proof}

The Grand Tensor reflected-slab identity proves fine OS positivity
even without a globally fixed dominating measure.  Theorem
\ref{thm:v13v42-path-os} identifies the additional conditional
structure needed to place that positivity in the recovery
factorization framework.

\subsection{Finite Gram lifts from a continuous detail fiber}

For a finite positive-time family \(F_1,\ldots,F_s\), define the
coarse half-functions \(L_yF_p\) in the dominated case.  Their
complete kernel is
\[
 {\mathsf W}_{pq}
 :=
 \int_Y |L_yF_p\rangle\langle L_yF_q|\,\kappa(dy).             \tag{42.17}
\]
It is positive as an operator-valued Gram kernel.  If a finite
coarse factor range \(H_+\) with projection \(P_H\) is chosen, then
\[
 {\mathsf W}^{H}_{pq}
 :=
 \int_Y |P_HL_yF_p\rangle\langle P_HL_yF_q|\,\kappa(dy)        \tag{42.18}
\]
has an ordinary finite positive Gram matrix of the V36 form.

The discarded reflected-Gram error obeys
\[
\begin{split}
 &\left|
 \langle L_yF_p,L_yF_q\rangle
 -
 \langle P_HL_yF_p,P_HL_yF_q\rangle
 \right|\\
 &\quad\leq
 \|(I-P_H)L_yF_p\|\,\|L_yF_q\|
 +
 \|P_HL_yF_p\|\,\|(I-P_H)L_yF_q\|.                            \tag{42.19}
\end{split}
\]
After integration, Cauchy--Schwarz gives a finite V34 matrix
residual from the individual factor-range tails.

\begin{proposition}[Physical population versus finite SDP input]
\label{prop:v13v42-compression}
The RPESM conditional kernel and reflected path law populate the
uncompressed positive kernel (42.17).  They populate a finite V36
Gram lift with certified residual once a finite range \(H_+\) and
bounds on the tails in (42.19) are supplied.  Existence of the
path law alone does not supply those compression bounds.
\end{proposition}

\begin{proof}
Positivity of (42.17) follows by integrating rank-one positive
operators.  Expanding (42.18) in a basis of \(H_+\) gives the
positive matrix of Theorem~\ref{thm:v13v36-gram-lift}.  Inequality
(42.19) and the V34 complete matrix residual theorem give the
certified finite approximation.  No estimate for the displayed
tails follows solely from normalization of \(K\).
\end{proof}

This is the precise population gap between the existential
finite regulator and a rational V41 SDP card.

\subsection{Mapping the Grand Tensor rows into the present ledger}

The attached RPESM construction now populates the following entries:
\[
\begin{array}{c|c|l}
\text{entry}&\text{status}&\text{source}\\ \hline
\mu_{n+1}=\mu_nK_{n+1/n}
 &\text{exact}&\text{projective shell identity}\\
{\cal R}_{n+1/n},\ {\cal C}_{n+1/n}
 &\text{exact}&\text{Proposition \ref{prop:v13v42-bayesian}}\\
\widetilde P={\cal J}P{\cal C}
 &\text{exact}&\text{Theorem \ref{thm:v13v42-transfer}}\\
\text{finite-cutoff OS positivity}
 &\text{exact}&\text{reflected-slab square formula}\\
\text{path reflection factorization}
 &\text{exact under finite/dominated detail}
 &\text{Theorem \ref{thm:v13v42-path-os}}\\
\text{finite Choi--Gram card}
 &\text{open}&\text{needs (42.19) and rational moments}\\
\text{cofinal local normality}
 &\text{open}&\text{needs V29--V33 uniform tails}\\
\text{physical local spectral gap}
 &\text{open}&\text{refresh gap is not this claim}.
\end{array}                                                    \tag{42.20}
\]

At each finite cutoff, the global refresh in the Grand Tensor gives
a centered \(L^2\) gap \(\rho\).  This is a valid regulator
statement.  Because the refresh is nonlocal and was introduced as
an auxiliary Renewal mechanism, it cannot be relabeled as a local
Standard-Model mass gap, a Yang--Mills coercivity estimate, or a
continuum clustering rate.  A physical time calibration and
local-dynamics comparison are separate rows.

\subsection{The first missing noncommutative row}

The Grand Tensor finite Euclidean/OS landing theorem supplies a
source-minimal finite OS carrier, the represented algebra, grading,
real structure, reflected covariance, and determinant/Pfaffian
marks on one history.  It also proves that the symmetric reflected
field generator does not determine an antisymmetric inner
derivation generated by the landed Dirac--Yukawa coefficient.

In the notation of that theorem, the first missing target is
\[
 {\cal L}^{\rm germ}_X(\pi_X(a_*))
 =
 -i[D_X^{\rm coeff},\pi_X(a_*)],                               \tag{42.21}
\]
together with its mixed incidence Gram.  Its source-minimal
residual is the positive rank-one short
\[
 r_{\delta,X}
 =
 b_{\delta,X}^*(I-P_{S_X})b_{\delta,X}>0.                      \tag{42.22}
\]

\begin{proposition}[Upstream stopping row for the active packet]
\label{prop:v13v42-first-short}
The commutative conditional-detail channel populates neither the
same-carrier CPTP germ (42.21) nor its incidence with the static
Dirac--Yukawa coefficient.  Therefore the V41 mixed-time SDP cannot
yet be claimed for the full represented SM operation algebra.
The smallest currently identified missing source direction is
(42.22).
\end{proposition}

\begin{proof}
The channel of Proposition~\ref{prop:v13v42-bayesian} acts on
commutative field observables by conditional integration, and the
transfer of Theorem~\ref{thm:v13v42-transfer} is self-adjoint.
An antisymmetric inner derivation on the represented fermion
operation algebra is not a function of those data.  The Grand
Tensor terminal-short theorem proves that its normalized witness
has nonzero orthogonal component (42.22).  Hence adding the target
to the source bank is necessary before the full mixed incidence can
vanish.
\end{proof}

This is an upstream obstruction, not a failure of compactness.
The next acquisition must be a differentiated same-history channel
panel, as specified by the Grand Tensor germ handoff.

\subsection{Cofinal consequences already available}

The exact projective identities imply that old-supported
commutative observables and their Markov transitions are compatible
at adjacent cutoffs.  Therefore their finite-dimensional
restrictions fit the V37 projective maps without a recovery
residual.  If the V33 regional energy-tail bounds and the
compression tails (42.19) are proved, the V37 diagonal yields a
locally normal OS-positive commutative recovery field.

This conditional conclusion does not require the fermionic germ.
The full SM operation algebra and its Dirac--Yukawa dynamics do
require it.

\begin{assessment}[Status after thirteenth-cycle V42]
\label{ass:v13v42-status}
The abstract recovery programme has reached the named RPESM
regulator.  Its conditional-detail kernel is an exact CPTP recovery,
and its fine reversible transfer has the exact block form
\(\widetilde P={\cal J}P{\cal C}\), preserving the coarse transfer
gap while killing pure-detail modes.  The recovered path law has a
finite or direct-integral reflected factorization.  Passing this
to a finite rational Choi--Gram card requires factor-range tail
bounds.  For the full noncommutative SM algebra, the first missing
row is the already-identified rank-one Dirac--Yukawa channel-germ
short.
\end{assessment}

\begin{firewall}[Claim boundary after thirteenth-cycle V42]
\label{fire:v13v42-claim}
The finite RPESM regulator and its commutative recovery are not
asserted to be the physical continuum.  No rational compression
tail, V41 feasible matrix, cofinal local normality estimate,
physical local gap, differentiated fermion CPTP germ, anomaly-free
cofinal line, spatial clustering, Lorentzian reconstruction, or
full SM--Einstein continuum has been proved.  The auxiliary global
refresh gap is not used as a Standard-Model or Yang--Mills mass
gap.
\end{firewall}

\subsection{Next bottleneck}

Go upstream to the rank-one short (42.22).  The next note should
derive the minimal same-carrier two-sided channel panel needed to
identify the antisymmetric derivative with
\(-i[D_X^{\rm coeff},\cdot]\), and formulate a finite conditional
complete-positivity criterion for adjoining that germ without
spoiling the already populated reflected field transfer.

\subsection{Parent record}

This V42 note appends to the validated thirteenth-cycle V41 file with
record
\[
\begin{array}{c|c}
\text{lines}&13526\\
\text{bytes}&477770\\
\text{SHA--256}&
\texttt{3179F8BDB6E32BF00F39AA0A5F21DE44FED0050FD32A7932802F017CDEB7972B}.
\end{array}
\]

% =============================================================================
% END APPENDED THIRTEENTH-CYCLE V42 MATERIAL
% =============================================================================

% =============================================================================
% APPENDED THIRTEENTH-CYCLE V43 MATERIAL
% =============================================================================

\section{Thirteenth-cycle V43: a same-history paired unitary panel
for the Dirac--Yukawa germ}
\label{sec:v13v43}

\subsection{Finite algebraic closure of the rank-one short}

At a fixed RPESM cutoff \(X\), the landed fermion operation algebra
is finite dimensional and the complete coefficient
\(D_X^{\rm coeff}=D_X^{\rm coeff\,*}\) is bounded on its represented
carrier.  Put
\[
 \alpha_{X,+}(t)(A)=e^{-itD_X^{\rm coeff}}A
 e^{itD_X^{\rm coeff}},\qquad
 \alpha_{X,-}(t)(A)=e^{itD_X^{\rm coeff}}A
 e^{-itD_X^{\rm coeff}}.                                      \tag{43.1}
\]

\begin{theorem}[Exact paired CPTP germ]
\label{thm:v13v43-paired-germ}
For every \(t\in{\mathbb R}\), both maps in (43.1) are unital
complete order isomorphisms, and their Schr\"odinger adjoints are
CPTP.  On the full finite operation algebra,
\[
 \boxed{\quad
 \lim_{t\downarrow0}
 {\alpha_{X,+}(t)-\alpha_{X,-}(t)\over2t}(A)
 =
 -i[D_X^{\rm coeff},A].\quad}                                 \tag{43.2}
\]
The convergence is in operator norm, uniformly for \(\|A\|\leq1\).
\end{theorem}

\begin{proof}
Each map is conjugation by a unitary, hence is a unital complete
order isomorphism and has a CPTP preadjoint.  The norm-convergent
finite-dimensional exponential expansion gives
\[
 \alpha_{X,\pm}(t)(A)
 =
 A\mp it[D_X^{\rm coeff},A]+O(t^2\|A\|),                      \tag{43.3}
\]
uniformly on the unit ball.  Subtracting the two expansions and
dividing by \(2t\) proves (43.2).
\end{proof}

This avoids linearizing the Choi cone at its rank-one identity
point.  Every nonzero-time sample is an actual channel.

\subsection{Coupling to the populated reflected field panel}

Let \(M_X(t)\) be any common CPTP field channel with
\(M_X(0)=I\), including a finite RPESM Markov transition.  Define
on the field--operation product
\[
 \Phi_{X,\pm}(t)
 :=
 M_X(t)\otimes\alpha_{X,\pm}(t).                              \tag{43.4}
\]

\begin{proposition}[Common field dynamics cancels from the
antisymmetric germ]
\label{prop:v13v43-common-cancel}
The maps (43.4) are CPTP.  If \(M_X(t)=I+t{\cal L}_{X}^{\rm
field}+o(t)\) strongly on the selected finite field ledger, then
\[
 \lim_{t\downarrow0}
 {\Phi_{X,+}(t)-\Phi_{X,-}(t)\over2t}
 =
 I_{\rm field}\otimes
 \bigl(-i[D_X^{\rm coeff},\,\cdot\,]\bigr).                    \tag{43.5}
\]
Thus adjoining the paired operation panel does not alter the
symmetric reflected field generator at first order.
\end{proposition}

\begin{proof}
Tensor products of CPTP maps are CPTP.  Expanding (43.4), the
common first-order field term occurs with the same sign in both
channels and cancels in their difference.  The antisymmetric
operation terms from (43.3) add, giving (43.5).
\end{proof}

Equation (43.5) is an incidence separation.  It does not identify
the antisymmetric germ with the OS Hamiltonian; it places both on
one finite marked panel without conflating them.

\subsection{Same-history occurrence}

Let \(\lambda\) denote the complete RPESM finite event and let
\(D_X^{\rm coeff}(\lambda)\) be its already populated measurable
coefficient.  Extend the event by an independent sign mark
\(s\in\{+,-\}\) and a predeclared finite set of rational durations
\(t\).  Conditional on \((\lambda,s,t)\), attach the channel
\(\Phi_{X,s}^{\lambda}(t)\) from (43.4).

\begin{proposition}[Conservative marked extension]
\label{prop:v13v43-marked-extension}
The marked construction is a probability-kernel extension of the
same RPESM history.  Forgetting \((s,t)\) returns the original
event law exactly.  Every finite matrix element of
\(\Phi_{X,s}^{\lambda}(t)\) is a direct writer on that extended
history, and its symmetric difference converges to the target
commutator (43.5).
\end{proposition}

\begin{proof}
A product with a normalized finite mark law preserves the old
marginal.  The coefficient and field channel are measurable
functions of the retained event, so finite channel matrix elements
are measurable marked writers.  Theorem
\ref{thm:v13v43-paired-germ} and Proposition
\ref{prop:v13v43-common-cancel} give the derivative.
\end{proof}

For the single witness \(a_*\) used in the Grand Tensor terminal
short, the target source now occurs.  If the source synthesis is
augmented by that normalized writer, its one-dimensional residual
becomes zero.  This closes the finite algebraic occurrence short;
it does not yet close physical calibration or the entire
differentiated channel bank.

\subsection{Symmetry, grading, and reality tests}

Let \(U_g\) act on the represented carrier.  If
\[
 D_X^{\rm coeff}(g\lambda)
 =
 U_gD_X^{\rm coeff}(\lambda)U_g^*,                            \tag{43.6}
\]
then
\[
 \alpha_{g\lambda,\pm}(t)\circ{\rm Ad}_{U_g}
 =
 {\rm Ad}_{U_g}\circ\alpha_{\lambda,\pm}(t).                   \tag{43.7}
\]
If \(D_X^{\rm coeff}\) is even for the grading, both channels
preserve the grading.  If an antiunitary real structure \(J_X\)
satisfies
\[
 J_XD_X^{\rm coeff}J_X^{-1}=\epsilon_DD_X^{\rm coeff},
 \qquad\epsilon_D\in\{+1,-1\},                                \tag{43.8}
\]
then \(J_X\) either preserves the two channel labels or exchanges
them, according to the sign and the antiunitary reversal of \(i\).

\begin{proof}
Equations (43.7)--(43.8) follow by applying covariance to the
exponentials in (43.1).  Evenness makes the exponentials commute
with the grading.
\end{proof}

These are finite exact tests.  If the complete coefficient has only
projective rather than honest gauge transport, the anomaly phase
must be resolved before (43.6) can define a descended channel
family.

\subsection{A quantitative central-difference error}

Let \(d_X=\|D_X^{\rm coeff}\|\).  The derivation
\(\delta_X(A)=-i[D_X^{\rm coeff},A]\) obeys
\[
 \|\delta_X\|\leq2d_X.                                        \tag{43.9}
\]
Since \(\alpha_{X,+}(t)=e^{t\delta_X}\) and
\(\alpha_{X,-}(t)=e^{-t\delta_X}\),
\[
 {\alpha_{X,+}(t)-\alpha_{X,-}(t)\over2t}
 =
 {\sinh(t\delta_X)\over t}.                                   \tag{43.10}
\]

\begin{proposition}[Uniform finite germ residual]
\label{prop:v13v43-germ-error}
For \(t\geq0\),
\[
 \left\|
 {\alpha_{X,+}(t)-\alpha_{X,-}(t)\over2t}
 -\delta_X
 \right\|
 \leq
 {t^2\|\delta_X\|^3\over6}e^{t\|\delta_X\|}
 \leq
 {4\over3}t^2d_X^3e^{2td_X}.                                  \tag{43.11}
\]
\end{proposition}

\begin{proof}
Use the odd exponential series:
\[
 {\sinh(t\delta_X)\over t}-\delta_X
 =
 \sum_{k\geq1}{t^{2k}\delta_X^{2k+1}\over(2k+1)!}.             \tag{43.12}
\]
Its norm is at most the first omitted power times the exponential
majorant, giving the first inequality.  Insert (43.9) for the
second.
\end{proof}

Thus a rational duration and a rational upper bound for \(d_X\)
give an explicit same-history germ residual.  A cofinal germ
requires
\[
 t_n^2d_{X_n}^3e^{2t_nd_{X_n}}\longrightarrow0.                \tag{43.13}
\]
If \(d_{X_n}\) grows, merely taking a fixed small duration is not
enough.

\subsection{Physical clock calibration}

Let \(c_X>0\) convert the marked parameter \(t\) to physical
duration \(t_{\rm phys}=c_Xt\).  Then the physical derivative is
\[
 {1\over c_X}\delta_X.                                        \tag{43.14}
\]
Identifying it with the intended coefficient germ requires either
\(c_X=1\) in the chosen units or a corresponding calibrated
rescaling of \(D_X^{\rm coeff}\).  The finite algebra does not
determine \(c_X\).

\begin{theorem}[Finite closure and cofinal requirements]
\label{thm:v13v43-closure}
At every finite RPESM cutoff, the paired marked extension closes
the rank-one algebraic occurrence short (42.22) with actual CPTP
channels and exact coefficient incidence.  A cofinal physical germ
still requires:
\begin{enumerate}
\item a common physical clock calibration \(c_X\);
\item the uniform or weighted bound (43.13);
\item compatibility of the marked channels under adjacent
cutoff restriction;
\item anomaly-free symmetry descent and the declared
grading/reality identities;
\item mixed Gram convergence with the already populated reflected
field and fermion word packets.
\end{enumerate}
\end{theorem}

\begin{proof}
Finite closure is Theorem~\ref{thm:v13v43-paired-germ} and
Proposition~\ref{prop:v13v43-marked-extension}.  Equation (43.14)
shows the calibration requirement, (43.11) gives the cofinal error,
and V23, V37, and V41 give respectively the descent, projective,
and mixed-Gram requirements.
\end{proof}

\begin{assessment}[Status after thirteenth-cycle V43]
\label{ass:v13v43-status}
The first noncommutative terminal short is closed algebraically at
every finite cutoff.  Paired unitary channels generated by the
landed complete Dirac--Yukawa coefficient are actual CPTP maps,
their symmetric difference is exactly the desired commutator
germ, and tensoring with the common reflected field channel cancels
the field generator from that antisymmetric derivative.  A
conservative same-history mark makes the witness occur without
changing the RPESM marginal.  The remaining obstruction is
cofinal physical calibration and mixed-incidence transport.
\end{assessment}

\begin{firewall}[Claim boundary after thirteenth-cycle V43]
\label{fire:v13v43-claim}
The marked unitary panel is a finite intervention construction.
It is not identified with physical Euclidean time, the OS
Hamiltonian, or observed fermion dynamics.  No uniform bound on
\(\|D_X^{\rm coeff}\|\), clock calibration, adjacent compatibility,
cofinal anomaly cancellation, mixed Gram convergence, continuum
gap, Lorentzian theory, or full SM--Einstein continuum has been
proved.
\end{firewall}

\subsection{Next bottleneck}

The next step is no longer finite complete positivity.  It is the
cofinal estimate (43.13) together with adjacent restriction of the
unitary panels.  One should compare the fine and coarse complete
coefficients by Duhamel's formula, derive a channel transport bound,
and test whether the existing graph-resolvent and determinant-line
controls supply a summable compatible duration schedule.

\subsection{Parent record}

This V43 note appends to the validated thirteenth-cycle V42 file with
record
\[
\begin{array}{c|c}
\text{lines}&13950\\
\text{bytes}&493473\\
\text{SHA--256}&
\texttt{F651B12B1CB28B3B361FD768AA0FD9F2BF394642BC1FAE11A5275BC4A1DA5925}.
\end{array}
\]

% =============================================================================
% END APPENDED THIRTEENTH-CYCLE V43 MATERIAL
% =============================================================================

% =============================================================================
% APPENDED THIRTEENTH-CYCLE V44 MATERIAL
% =============================================================================

\section{Thirteenth-cycle V44: adjacent-cutoff Duhamel transport of
the paired channel germ}
\label{sec:v13v44}

\subsection{The nonduplicating question}

The Grand Tensor already contains the finite channel-tangent,
Hamiltonian-projection, quadratic-remainder, and noisy-refocusing
theorems.  V43 supplies an exact finite paired unitary panel.  The
remaining issue is different: whether those panels represent one
germ after passage from cutoff \(n\) to cutoff \(n+1\).

Let \({\cal A}_n\) be finite matrix algebras and let
\[
 \iota_n:{\cal A}_n\longrightarrow{\cal A}_{n+1}               \tag{44.1}
\]
be unital injective star homomorphisms.  Let \(D_n=D_n^*\) and
define the bounded star derivations and automorphism groups
\[
 \delta_n(A)=-i[D_n,A],\qquad
 \alpha_n(t)=e^{t\delta_n}
 ={\rm Ad}_{e^{-itD_n}}.                                       \tag{44.2}
\]
The complete adjacent germ defect is
\[
 {\mathfrak d}_n
 :=
 \delta_{n+1}\iota_n-\iota_n\delta_n,\qquad
 d_n:=\|{\mathfrak d}_n\|_{\rm cb}.                            \tag{44.3}
\]

\subsection{Exact Duhamel identity}

\begin{theorem}[Adjacent unitary-channel Duhamel formula]
\label{thm:v13v44-duhamel}
For every \(t\in{\mathbb R}\),
\[
 \boxed{\quad
 \alpha_{n+1}(t)\iota_n-\iota_n\alpha_n(t)
 =
 \int_0^t
 \alpha_{n+1}(t-s){\mathfrak d}_n\alpha_n(s)\,ds.\quad}         \tag{44.4}
\]
The oriented integral is used when \(t<0\).  Consequently
\[
 \|\alpha_{n+1}(t)\iota_n-\iota_n\alpha_n(t)\|_{\rm cb}
 \leq |t|d_n.                                                   \tag{44.5}
\]
The same bound holds for each member of the paired panel
\(\alpha_n(\pm t)\).
\end{theorem}

\begin{proof}
For \(t\geq0\), differentiate
\[
 F(s)=\alpha_{n+1}(t-s)\iota_n\alpha_n(s).
\]
Using commutation of each generator with its own group gives
\[
 F'(s)=
 -\alpha_{n+1}(t-s)
 {\mathfrak d}_n
 \alpha_n(s).                                                   \tag{44.6}
\]
Integrating from \(0\) to \(t\) and rearranging proves (44.4).
The automorphisms are complete isometries and \(\iota_n\) is a
complete isometry, so the cb norm of the integrand is \(d_n\).
This proves (44.5).  Negative \(t\) follows by reversing the
integral.
\end{proof}

Unlike a comparison of individual unitaries, (44.4) automatically
quotients scalar or commutant shifts of the Hamiltonian.

\subsection{Coefficient mismatch modulo the commutant}

Suppose \(D_n\in{\cal A}_n\) and write
\[
 D_{n+1}=\iota_n(D_n)+C_n+R_n,\qquad
 C_n=C_n^*\in\iota_n({\cal A}_n)',\quad R_n=R_n^*.             \tag{44.7}
\]
Then
\[
 {\mathfrak d}_n(A)=-i[R_n,\iota_n(A)].                         \tag{44.8}
\]

\begin{proposition}[Commutant-quotient bound]
\label{prop:v13v44-commutant}
Under (44.7),
\[
 d_n\leq2\|R_n\|.                                               \tag{44.9}
\]
More intrinsically, with
\[
 r_n:=
 \inf_{C=C^*\in\iota_n({\cal A}_n)'}
 \|D_{n+1}-\iota_n(D_n)-C\|,                                  \tag{44.10}
\]
one has
\[
 d_n\leq2r_n.                                                   \tag{44.11}
\]
Moreover \(d_n=0\) if and only if
\[
 D_{n+1}-\iota_n(D_n)\in\iota_n({\cal A}_n)'.                  \tag{44.12}
\]
\end{proposition}

\begin{proof}
At every matrix level,
\[
 \|[R_n,X]\|\leq2\|R_n\|\,\|X\|,
\]
which proves (44.9).  Take the infimum over commutant
representatives to get (44.11).  Equation (44.3) vanishes exactly
when the Hamiltonian difference commutes with every element in the
embedded algebra, proving (44.12).
\end{proof}

Thus the complete coefficient itself need not converge.  Only its
class modulo the embedded commutant is visible to the operation
germ.

\subsection{Physical clock calibration enters before comparison}

Let \(c_n>0\) be the conversion from marked duration to physical
duration and put
\[
 \widehat D_n=c_n^{-1}D_n,\qquad
 \widehat\delta_n=-i[\widehat D_n,\,\cdot\,].                   \tag{44.13}
\]
The physical adjacent defect is
\[
 \widehat{\mathfrak d}_n
 =
 \widehat\delta_{n+1}\iota_n-\iota_n\widehat\delta_n,\qquad
 \widehat d_n=\|\widehat{\mathfrak d}_n\|_{\rm cb}.             \tag{44.14}
\]
It is controlled by the calibrated commutant quotient
\[
 \widehat r_n
 =
 \inf_{C=C^*\in\iota_n({\cal A}_n)'}
 \|\widehat D_{n+1}-\iota_n(\widehat D_n)-C\|,\qquad
 \widehat d_n\leq2\widehat r_n.                                \tag{44.15}
\]
Consequently a norm comparison of the uncalibrated \(D_n\) does
not prove physical germ transport when \(c_n\) varies.

\subsection{Why shrinking duration cannot repair the germ}

For \(t>0\), define the central-difference channel germ
\[
 g_n(t):={\alpha_n(t)-\alpha_n(-t)\over2t}.                    \tag{44.16}
\]

\begin{theorem}[Channel-value versus germ compatibility]
\label{thm:v13v44-shrinking-no-repair}
For every \(t>0\),
\[
 \|g_{n+1}(t)\iota_n-\iota_ng_n(t)\|_{\rm cb}\leq d_n.          \tag{44.17}
\]
Let
\[
 e_n(t):=\|g_n(t)-\delta_n\|_{\rm cb}.                          \tag{44.18}
\]
Then
\[
 \boxed{\quad
 d_n-e_{n+1}(t)-e_n(t)
 \leq
 \|g_{n+1}(t)\iota_n-\iota_ng_n(t)\|_{\rm cb}
 \leq
 d_n+e_{n+1}(t)+e_n(t).\quad}                                  \tag{44.19}
\]
In particular, if \(e_n(t_n)+e_{n+1}(t_n)\to0\), compatibility of
the sampled germs is equivalent to \(d_n\to0\).  The weaker
condition
\[
 t_nd_n\longrightarrow0                                       \tag{44.20}
\]
only makes the unnormalized channel values compatible; it does not
make their normalized germs compatible.
\end{theorem}

\begin{proof}
Apply (44.5) at \(t\) and at \(-t\), add the two bounds, and divide
by \(2t\) to obtain (44.17).  Insert and subtract
\(\delta_{n+1}\iota_n\) and \(\iota_n\delta_n\); the triangle
inequality gives the upper bound in (44.19), while the reverse
triangle inequality gives the lower bound.  The equivalence and
the distinction from (44.20) follow immediately.
\end{proof}

This corrects a possible misuse of the duration schedule in
(43.13).  Small duration controls the finite-difference error and
the raw channel mismatch, but the adjacent derivation defect is an
independent first-order row.

For completeness, if \(M_n=\|\delta_n\|_{\rm cb}\), the odd
exponential series gives
\[
 e_n(t)
 \leq {t^2M_n^3\over6}e^{tM_n}.                               \tag{44.21}
\]
Thus a sufficient sampled-germ schedule is
\[
 d_n\longrightarrow0,\qquad
 t_n^2M_n^3e^{t_nM_n}\longrightarrow0.                         \tag{44.22}
\]

\subsection{A first-order recovery of the defect from channel data}

The Duhamel average in (44.4) also shows that the adjacent
derivation defect is observable from sufficiently short channel
comparisons.

\begin{proposition}[Short-time identification]
\label{prop:v13v44-identification}
For \(t>0\),
\[
 \left\|
 {\,\alpha_{n+1}(t)\iota_n-\iota_n\alpha_n(t)\over t}
 -{\mathfrak d}_n
 \right\|_{\rm cb}
 \leq
 {t\over2}(M_{n+1}+M_n)d_n.                                   \tag{44.23}
\]
\end{proposition}

\begin{proof}
Divide (44.4) by \(t\).  For a bounded derivation group,
\[
 \|\alpha_k(u)-I\|_{\rm cb}\leq |u|M_k,                        \tag{44.24}
\]
by integrating its derivative and using complete isometry.
Subtract \({\mathfrak d}_n\) from the averaged integrand in
(44.4).  The two resulting terms are bounded by
\((t-s)M_{n+1}d_n\) and \(sM_nd_n\).  Integration and division by
\(t\) give (44.23).
\end{proof}

A family of channel values at times tending to zero therefore
determines the adjacent germ defect provided
\(t(M_n+M_{n+1})\to0\).  Static nonzero-time values without this
scale control do not.

\subsection{Fixed-time inductive-limit channel}

For \(m>n\), write
\[
 \iota_{m,n}=\iota_{m-1}\cdots\iota_n.                         \tag{44.25}
\]
Repeated use of (44.5) gives
\[
 \|\alpha_m(t)\iota_{m,n}
 -\iota_{m,n}\alpha_n(t)\|_{\rm cb}
 \leq |t|\sum_{k=n}^{m-1}d_k.                                  \tag{44.26}
\]

\begin{theorem}[Summable-defect automorphism limit]
\label{thm:v13v44-inductive-limit}
If
\[
 \sum_{n=1}^\infty d_n<\infty,                                 \tag{44.27}
\]
then for every fixed \(t\) the transported automorphisms converge
in norm on each finite stage to a unital star automorphism
\(\alpha_\infty(t)\) of the \(C^*\)-inductive limit.  The family
\((\alpha_\infty(t))_{t\in{\mathbb R}}\) is a group.  It is
strongly continuous on the algebraic inductive-limit core.
\end{theorem}

\begin{proof}
For an element born at stage \(n\), (44.26) and (44.27) make its
transported images Cauchy in the completed inductive limit.
The limit is linear, unital, multiplicative, star preserving, and
isometric because each approximating map has those properties.
The limits at \(t\) and \(-t\) are inverse, so the map is an
automorphism.  Passing the finite-stage group identity through the
norm limit gives the group law.  Strong continuity on each born
element follows from continuity at its finite stage plus the
uniform tail bound (44.26).
\end{proof}

Summability is stronger than the pointwise condition \(d_n\to0\)
needed for adjacent germ consistency.  It buys a route-independent
fixed-time automorphism on the full inductive limit.

\subsection{Transport of the complete V41 Gram blocks}

Let \(\omega_{n+1}\iota_n=\omega_n\) and let
\[
 {\bf A}=[A_{pq}]\in M_q({\cal A}_n)                           \tag{44.28}
\]
be a complete reflected or mixed-time observable block.  Define
the channel-transport Gram discrepancy at time \(t\) by applying
the states entrywise to
\[
 [\,\alpha_{n+1}(t)\iota_n(A_{pq})
     -\iota_n\alpha_n(t)(A_{pq})\,]_{p,q}.                     \tag{44.29}
\]

\begin{proposition}[Whole-block channel-to-Gram transport]
\label{prop:v13v44-gram}
The operator norm of the scalar matrix in (44.29) is at most
\[
 |t|d_n\,\|{\bf A}\|_{M_q({\cal A}_n)}.                        \tag{44.30}
\]
For central-difference germs, the corresponding bound is
\[
 \bigl(d_n+e_{n+1}(t)+e_n(t)\bigr)
 \|{\bf A}\|_{M_q({\cal A}_n)}.                                \tag{44.31}
\]
\end{proposition}

\begin{proof}
The amplification of a state is completely contractive.  Apply it
to the \(q\)-fold amplification of (44.5) to obtain (44.30), and
to the upper bound in (44.19) to obtain (44.31).
\end{proof}

This estimate preserves the complete Gram block until the final
operator norm, as required by the V35 and V40 audits.

\subsection{What the existing RPESM transports do and do not imply}

The exact adjacent shell relation transports the old Markov path,
clock, reflection, and source rows.  The determinant and Pfaffian
block formulas transport kernels, divisor orders, orientations,
and line products.  None of those identities bounds
\[
 \widehat r_n=
 {\rm dist}\bigl(
 \widehat D_{n+1}-\iota_n(\widehat D_n),
 \iota_n({\cal A}_n)'\bigr),                                  \tag{44.32}
\]
because determinant-line data do not determine commutators on the
embedded operation algebra.  Likewise, a graph-resolvent gap
controls functional calculus at one cutoff but does not compare
the calibrated coefficients at adjacent cutoffs without a
transport residual.

Equation (44.32) is therefore the next directly computable
physical row.  At finite cutoff it is a matrix distance to a linear
subspace and can be expressed as a convex norm-minimization
problem.  A rational upper certificate consists of a rational
commutant candidate \(C_n\) and a rational operator-norm upper
bound for the remainder.

\begin{assessment}[Status after thirteenth-cycle V44]
\label{ass:v13v44-status}
Adjacent transport of the paired unitary panel is now exact at the
abstract level.  The complete derivation defect integrates to the
channel mismatch by (44.4), commutant shifts are correctly
quotiented, and summable defects construct a fixed-time
inductive-limit automorphism.  Most importantly, shrinking the
duration cannot hide a nonvanishing normalized germ defect:
cofinal germ compatibility requires \(d_n\to0\) independently of
the central-difference schedule.  The same bound transports the
complete V41 Gram blocks.
\end{assessment}

\begin{firewall}[Claim boundary after thirteenth-cycle V44]
\label{fire:v13v44-claim}
No RPESM embedding \(\iota_n\), physical clock calibration,
calibrated commutant quotient (44.32), rational remainder,
summability bound, or mixed-Gram transport card has been populated.
Exact determinant/Pfaffian transport is not used as a substitute.
The note does not prove a continuum fermion dynamics, physical
gap, Lorentzian theory, or the full SM--Einstein continuum.
\end{firewall}

\subsection{Next bottleneck}

The next acquisition is the calibrated commutant distance (44.32)
on the named RPESM adjacent coefficient pair.  Before estimating
it, one must choose the actual operation-algebra embedding induced
by the projective fermion carrier.  The next note should prove the
finite conditional-expectation formula for that distance, separate
the visible commutator component from new-detail leakage, and
return an explicit singular vector if the distance cannot vanish.

\subsection{Parent record}

This V44 note appends to the validated thirteenth-cycle V43 file with
record
\[
\begin{array}{c|c}
\text{lines}&14249\\
\text{bytes}&504099\\
\text{SHA--256}&
\texttt{ECDB9BE0EB74E544BB4448CFC0E05DBA467708BC8268C89898B28D7B1EABDE0B}.
\end{array}
\]

% =============================================================================
% END APPENDED THIRTEENTH-CYCLE V44 MATERIAL
% =============================================================================

% =============================================================================
% APPENDED THIRTEENTH-CYCLE V45 MATERIAL
% =============================================================================

\section{Thirteenth-cycle V45: mandatory companion audit before the
commutant-distance card}
\label{sec:v13v45}

\subsection{Files inspected}

The required companion checkpoint inspected
\[
\begin{array}{c|r|r|c}
\text{file}&\text{lines}&\text{bytes}&\text{SHA--256}\\ \hline
\texttt{Renewal\_Yang\_Mills\_Mass\_Gap\_Research\_Notes\_V84.tex}
&28525&972768&
\texttt{5248AFF4B520E512C787654EA4B576A9A3A1A5065641FDDD3E7DDCFA23823525}\\
\texttt{Renewal\_NS\_Stability\_Research\_Notes\_V26.tex}
&5319&278283&
\texttt{82C887F8C2C69E4F28FAD7C3DC8DE5B9099CB5A4BF431EBA1FFF3E91935978AD}.
\end{array}                                                    \tag{45.1}
\]
The auxiliary file
\(\texttt{Renewal\_YM\_Parallel\_Research\_Notes\_V5.tex}\)
was also inspected.  Its current record is \(1351\) lines,
\(54174\) bytes, and SHA--256
\[
 \texttt{306F1BB84CFA8A8D47EA569EC17E9545F9867C8D32B548EC9D9C444B4F3C0512}.
                                                                    \tag{45.2}
\]

\subsection{New YM evidence}

YM V83 populated the first fully rational signed-diagonal path
card.  It rigorously located the first avoidable loss at separate
Frobenius bounds for a trial inverse and an interval remainder.
YM V84 then multiplied the complete interval matrices before
norming.  This improved the certified integrated energy lengths by
factors greater than \(2.85\) and \(2.95\), but the bounds remain
orders of magnitude above the floating diagnostic scale and are
not operative.

The relevant lesson for (44.32) is precise.  The fine coefficient
must first be decomposed by the exact conditional expectations onto
the embedded coarse algebra and its commutant.  The visible
coarse drift and the genuinely mixed new-detail component should
be formed as complete matrices.  Only then should Hilbert--Schmidt
or operator norms be taken.  Bounding individual blocks or
coefficient entries beforehand would repeat the V83 loss.

YM V84 supplies no Standard-Model coefficient pair, fermion-carrier
embedding, clock calibration, or commutant-distance value.

\subsection{NS and auxiliary-YM status}

NS V26 is unchanged and supplies no new input.  Its structured
rank-one contraction remains consistent with the complete-product
order selected above.

The auxiliary YM file retains its exact failure stop for a
volume-uniform tail of nontrivial extensive global charge.
Nothing in V44 uses such a tail: the adjacent germ defect is a
finite local matrix quantity.  Its cofinal estimates must remain
regionwise, as already required in V40.

\begin{proposition}[Audit decision]
\label{prop:v13v45-decision}
No companion result closes the V44 calibrated commutant-distance
row.  The next SM note should, for a factor embedding
\[
 M_d\otimes I_k\subset M_d\otimes M_k,                         \tag{45.3}
\]
construct:
\begin{enumerate}
\item the trace-preserving conditional expectation onto the
commutant \(I_d\otimes M_k\);
\item the exact Hilbert--Schmidt distance and commutator-energy
identity;
\item an orthogonal decomposition into scalar, coarse-visible,
detail-only, and mixed coarse--detail components;
\item a finite matrix-unit or singular-vector witness on every
nonzero branch;
\item operator-norm upper cards sufficient for V44.
\end{enumerate}
\end{proposition}

\begin{proof}
YM V84 shows that complete matrix contraction must precede norm
relaxation.  Conditional expectations give exactly that contraction
for the tensor-factor inclusion (45.3).  The companion notes
contain none of the required SM matrices, so only the proof
architecture, not a numerical bound, is imported.
\end{proof}

\begin{assessment}[Status after thirteenth-cycle V45]
\label{ass:v13v45-status}
The mandatory audit is complete.  YM V84 gives rigorous evidence
that changing contraction order can remove a substantial false
loss while still leaving a genuine obstruction.  V46 will apply
that order exactly: project the full adjacent coefficient mismatch
onto its commutant and tensor-interaction components before taking
any norm.
\end{assessment}

\begin{firewall}[Claim boundary after thirteenth-cycle V45]
\label{fire:v13v45-claim}
No YM path number, NS kernel constant, or auxiliary global-charge
estimate enters the SM proof.  The actual RPESM tensor-factor
embedding and coefficient matrices have not been extracted.
No commutant residual, cofinal fermion germ, continuum dynamics, or
full SM--Einstein continuum is proved by this audit.
\end{firewall}

\subsection{Next step}

V46 will carry out the finite conditional-expectation and
coarse/detail decomposition and then continue beyond this audit.

\subsection{Parent record}

This V45 note appends to the validated thirteenth-cycle V44 file with
record
\[
\begin{array}{c|c}
\text{lines}&14653\\
\text{bytes}&517706\\
\text{SHA--256}&
\texttt{CA471587E1118C2FDF9DBBD494CD40951D88797827B8197657AB60C40D266CCF}.
\end{array}
\]

% =============================================================================
% END APPENDED THIRTEENTH-CYCLE V45 MATERIAL
% =============================================================================

% =============================================================================
% APPENDED THIRTEENTH-CYCLE V46 MATERIAL
% =============================================================================

\section{Thirteenth-cycle V46: conditional-expectation
decomposition of the adjacent fermion germ}
\label{sec:v13v46}

\subsection{The factor embedding}

Fix one adjacent cutoff pair and suppose the transported old
operation algebra has the factor form
\[
 {\cal A}=M_d({\mathbb C})\otimes I_k
 \subset M_d({\mathbb C})\otimes M_k({\mathbb C})={\cal B}.    \tag{46.1}
\]
The commutant is
\[
 {\cal A}'=I_d\otimes M_k({\mathbb C}).                        \tag{46.2}
\]
Let
\[
 X:=\widehat D_{\rm fine}
   -\widehat D_{\rm coarse}\otimes I_k=X^*                    \tag{46.3}
\]
be the calibrated adjacent coefficient mismatch from V44.

All traces and Hilbert--Schmidt norms in this note are
unnormalized.  Define
\[
\begin{split}
 {\mathbb E}_{A}(Z)
 &={1\over k}{\rm Tr}_k(Z)\otimes I_k,\\
 {\mathbb E}_{C}(Z)
 &={I_d\over d}\otimes{\rm Tr}_d(Z),\\
 {\mathbb E}_{0}(Z)
 &={{\rm Tr}Z\over dk}I_{dk}.                                  \tag{46.4}
\end{split}
\]
The first map projects onto the embedded coarse factor, the second
onto its commutant, and the third onto their scalar intersection.

\begin{lemma}[Commuting conditional expectations]
\label{lem:v13v46-expectations}
The maps in (46.4) are unital trace-preserving completely positive
idempotents.  On Hilbert--Schmidt space they are orthogonal
projections and
\[
 {\mathbb E}_A{\mathbb E}_C
 =
 {\mathbb E}_C{\mathbb E}_A
 =
 {\mathbb E}_0.                                                \tag{46.5}
\]
Moreover
\[
 {\mathbb E}_C(Z)
 =
 \int_{{\cal U}(d)}
 (U\otimes I_k)Z(U^*\otimes I_k)\,dU.                          \tag{46.6}
\]
\end{lemma}

\begin{proof}
The partial-trace formulas directly give unitality,
trace preservation, complete positivity, and idempotence.
Trace duality shows that each is self-adjoint for the
Hilbert--Schmidt inner product, hence an orthogonal projection.
Applying the two partial traces in either order gives the normalized
scalar trace, proving (46.5).  Haar twirling kills every traceless
first-factor matrix and fixes \(I_d\otimes M_k\), which proves
(46.6).
\end{proof}

\subsection{Four orthogonal components}

Define
\[
\begin{split}
 X_0&={\mathbb E}_0X,\\
 X_A&=({\mathbb E}_A-{\mathbb E}_0)X,\\
 X_C&=({\mathbb E}_C-{\mathbb E}_0)X,\\
 X_{\rm mix}
 &=(I-{\mathbb E}_A-{\mathbb E}_C+{\mathbb E}_0)X.             \tag{46.7}
\end{split}
\]

\begin{theorem}[Coarse--detail Pythagoras]
\label{thm:v13v46-pythagoras}
The four components in (46.7) are pairwise Hilbert--Schmidt
orthogonal and
\[
 X=X_0+X_A+X_C+X_{\rm mix}.                                   \tag{46.8}
\]
The commutant-invisible part and the visible germ mismatch are
\[
 {\mathbb E}_CX=X_0+X_C,\qquad
 R:=(I-{\mathbb E}_C)X=X_A+X_{\rm mix},                        \tag{46.9}
\]
with the exact identity
\[
 \boxed{\quad
 \|R\|_{\rm HS}^2
 =
 \|X_A\|_{\rm HS}^2+\|X_{\rm mix}\|_{\rm HS}^2.\quad}          \tag{46.10}
\]
\end{theorem}

\begin{proof}
The commuting orthogonal projections
\({\mathbb E}_A\) and \({\mathbb E}_C\) have intersection
\({\mathbb E}_0\).  The four products
\[
 {\mathbb E}_0,\quad
 {\mathbb E}_A-{\mathbb E}_0,\quad
 {\mathbb E}_C-{\mathbb E}_0,\quad
 I-{\mathbb E}_A-{\mathbb E}_C+{\mathbb E}_0                  \tag{46.11}
\]
are mutually orthogonal projections summing to the identity.
Equations (46.8)--(46.10) follow.
\end{proof}

The meanings are now separated exactly:
\(X_A\) is the coarse-visible coefficient drift,
\(X_C\) is a detail-only Hamiltonian invisible to old
observables, and \(X_{\rm mix}\) is genuine coarse--detail
interaction leakage.  The scalar \(X_0\) is invisible to every
commutator.

\subsection{Exact commutator energy}

Let \(E_{ij}\), \(1\leq i,j\leq d\), be the standard matrix units
and define
\[
 \Delta_{\rm com}(X)
 :=
 \sum_{i,j=1}^d
 \|[X,E_{ij}\otimes I_k]\|_{\rm HS}^2.                         \tag{46.12}
\]

\begin{theorem}[Complete matrix-unit commutator identity]
\label{thm:v13v46-commutator-identity}
One has
\[
 \boxed{\quad
 \Delta_{\rm com}(X)
 =
 2d\,\|X-{\mathbb E}_CX\|_{\rm HS}^2
 =
 2d\bigl(\|X_A\|_{\rm HS}^2+
          \|X_{\rm mix}\|_{\rm HS}^2\bigr).\quad}              \tag{46.13}
\]
Consequently the adjacent germ is exactly compatible on the
embedded factor if and only if
\[
 X_A=0,\qquad X_{\rm mix}=0.                                  \tag{46.14}
\]
\end{theorem}

\begin{proof}
Expanding the commutator squares and using the matrix-unit
identities gives
\[
\begin{split}
 \Delta_{\rm com}(X)
 &=
 2d\,{\rm Tr}(X^*X)
 -2\,{\rm Tr}\bigl(({\rm Tr}_dX)^*{\rm Tr}_dX\bigr).           \tag{46.15}
\end{split}
\]
On the other hand,
\[
 \|{\mathbb E}_CX\|_{\rm HS}^2
 ={1\over d}\|{\rm Tr}_dX\|_{\rm HS}^2.                        \tag{46.16}
\]
Because \({\mathbb E}_C\) is an orthogonal projection,
(46.15)--(46.16) give the first equality in (46.13).
The second is (46.10).  Vanishing is equivalent to
\(X={\mathbb E}_CX\in{\cal A}'\), which is precisely the V44
zero-defect condition; orthogonality gives (46.14).
\end{proof}

This identity performs every coarse and detail contraction before
taking the norm.  It is the exact conditional-expectation analogue
of the contraction order selected by the V45 audit.

\subsection{Finite failure witnesses}

If \(\Delta_{\rm com}(X)>0\), at least one matrix unit satisfies
\[
 [X,E_{ij}\otimes I_k]\neq0.                                  \tag{46.17}
\]
More quantitatively,
\[
 \max_{i,j}
 \|[X,E_{ij}\otimes I_k]\|_{\rm HS}
 \geq
 \sqrt{2\over d}\,\|R\|_{\rm HS}.                              \tag{46.18}
\]

\begin{proposition}[Explicit adjacent-germ witness]
\label{prop:v13v46-witness}
Choose \((i_*,j_*)\) maximizing the left side of (46.18) and put
\[
 Q_*:=-i[X,E_{i_*j_*}\otimes I_k].                             \tag{46.19}
\]
Then \(Q_*\) is a direct nonzero source witness for failure of
adjacent germ compatibility.  The positive rank-one coefficient
Gram
\[
 {\mathfrak R}_*=
 |Q_*\rangle_{\rm HS}\langle Q_*|                              \tag{46.20}
\]
has range equal to the source-minimal missing direction for this
fixed query.  Its squared norm obeys
\[
 \|Q_*\|_{\rm HS}^2
 \geq {2\over d}\|R\|_{\rm HS}^2.                              \tag{46.21}
\]
\end{proposition}

\begin{proof}
There are \(d^2\) nonnegative summands in (46.12).  Their maximum
is at least their average.  Insert (46.13) to obtain (46.18) and
(46.21).  A nonzero vector has a positive rank-one Gram with the
stated range.
\end{proof}

If \(X_A\neq0\) and \(X_{\rm mix}=0\), the witness is a pure
coarse-coefficient drift.  If \(X_A=0\) and
\(X_{\rm mix}\neq0\), it is a new-detail coupling leakage.  When
both are nonzero, (46.10) records their costs without cross-term
cancellation.

\subsection{Singular-vector witness bank}

Define the Hilbert--Schmidt operator
\[
 {\cal K}_X:M_d({\mathbb C})\longrightarrow M_{dk}({\mathbb C}),
 \qquad
 {\cal K}_X(A)=-i[X,A\otimes I_k].                             \tag{46.22}
\]
Let \(s_{\max}({\cal K}_X)\) be its largest singular value for
Hilbert--Schmidt norms.

\begin{proposition}[Optimal Hilbert--Schmidt query]
\label{prop:v13v46-singular-query}
There is \(A_*\in M_d\) with
\(\|A_*\|_{\rm HS}=1\) such that
\[
 \|{\cal K}_X(A_*)\|_{\rm HS}
 =s_{\max}({\cal K}_X).                                       \tag{46.23}
\]
This query also has \(\|A_*\|_{\rm op}\leq1\), and therefore the
complete derivation defect \(d_{\rm adj}\) satisfies
\[
 {s_{\max}({\cal K}_X)\over\sqrt{dk}}
 \leq d_{\rm adj}.                                             \tag{46.24}
\]
The singular value vanishes exactly on the compatible branch.
\end{proposition}

\begin{proof}
Finite-dimensional singular-value decomposition gives (46.23).
The operator norm is bounded by the Hilbert--Schmidt norm, so
\(A_*\) is an admissible unit-ball query.  The output operator norm
is at least its Hilbert--Schmidt norm divided by
\(\sqrt{dk}\), yielding (46.24).  Vanishing means every commutator
in (46.22) is zero, equivalently \(X\in{\cal A}'\).
\end{proof}

This singular query is sharper than selecting a matrix unit, while
the matrix-unit bank is simpler to certify exactly.

\subsection{Operator-norm distance and an a posteriori card}

Since \({\mathbb E}_C\) is a contractive conditional expectation,
\[
 {\rm dist}_{\rm op}(X,{\cal A}')
 \leq\|R\|_{\rm op}
 \leq2\,{\rm dist}_{\rm op}(X,{\cal A}').                      \tag{46.25}
\]
Indeed, the first inequality uses the candidate
\({\mathbb E}_CX\).  For \(C\in{\cal A}'\),
\[
 R=(I-{\mathbb E}_C)(X-C),
\]
so contractivity and the triangle inequality give
\(\|R\|_{\rm op}\leq2\|X-C\|_{\rm op}\); take the infimum.

\begin{corollary}[V44 transport card]
\label{cor:v13v46-v44-card}
The adjacent cb derivation defect satisfies
\[
 d_{\rm adj}
 \leq2\|R\|_{\rm op}
 \leq2\|R\|_{\rm HS}.                                         \tag{46.26}
\]
Thus either of the following is sufficient:
\[
 \sum_n\|R_n\|_{\rm op}<\infty
 \quad\Longrightarrow\quad
 \hbox{the fixed-time V44 automorphism limit},                 \tag{46.27}
\]
and
\[
 \|R_n\|_{\rm op}\longrightarrow0
 \quad\Longrightarrow\quad
 \hbox{adjacent physical germ compatibility}.                 \tag{46.28}
\]
\end{corollary}

\begin{proof}
The first inequality is the commutator bound of V44 with the
commutant candidate \({\mathbb E}_CX\); the second is standard.
Insert the result into Theorems
\ref{thm:v13v44-inductive-limit} and
\ref{thm:v13v44-shrinking-no-repair}.
\end{proof}

The Hilbert--Schmidt criterion can become dimensionally too strong.
A normalized Hilbert--Schmidt residual tending to zero does not
imply (46.28) when dimensions grow.  The operator norm, or a
weighted local operator norm on each fixed region, is the cofinal
quantity.

\subsection{Rational evaluation}

For Gaussian-rational coefficient matrices, every component in
(46.7) is Gaussian-rational.  The squared Hilbert--Schmidt values
in (46.10) and (46.13) are exact rationals.  A rational
operator-norm upper bound follows from
\[
 \|R\|_{\rm op}
 \leq
 \sqrt{\|R\|_1\|R\|_\infty},                                  \tag{46.29}
\]
with an outward rational square-root ceiling.  Alternatively,
exact rational \(LDL^*\) certification of
\[
 \beta^2I-R^*R\succeq0                                        \tag{46.30}
\]
proves \(\|R\|_{\rm op}\leq\beta\).

If the physical coefficients are enclosed by interval matrices,
first apply the linear expectations in (46.7) to the complete
interval matrix, then multiply \(R^*R\), and only afterward take
entrywise radii or a norm.  This preserves the V45 contraction
order.

\subsection{Beyond a literal tensor factor}

If the old algebra is a finite direct sum
\[
 {\cal A}=\bigoplus_\lambda
 \bigl(M_{d_\lambda}\otimes I_{k_\lambda}\bigr),               \tag{46.31}
\]
its commutant and trace-preserving expectation are the direct sums
of the block formulas above, together with removal of
off-central-block entries.  The commutator identity becomes
\[
 \Delta_{\rm com}(X)
 =
 \sum_\lambda2d_\lambda
 \|R_\lambda\|_{\rm HS}^2
 +\Delta_{\rm off}(X),                                        \tag{46.32}
\]
where \(\Delta_{\rm off}\geq0\) is the exact matrix-unit sum from
off-central-block couplings.  Hence the same zero criterion and
witness construction apply blockwise.  The actual RPESM
representation type must be read before choosing between (46.1)
and (46.31).

\begin{assessment}[Status after thirteenth-cycle V46]
\label{ass:v13v46-status}
The calibrated adjacent coefficient mismatch now has a complete
finite decomposition.  Conditional expectations split it into
scalar, coarse-visible, detail-only, and mixed coarse--detail
parts.  Only the coarse-visible and mixed parts contribute to the
old-algebra germ, and their squared Hilbert--Schmidt costs add
exactly.  The full matrix-unit commutator identity gives an exact
zero test and an explicit failure witness, while the operator-norm
residual supplies the V44 cofinal transport card.
\end{assessment}

\begin{firewall}[Claim boundary after thirteenth-cycle V46]
\label{fire:v13v46-claim}
The tensor-factor representation (46.1) has not been identified
with the actual adjacent RPESM fermion carrier.  No calibrated
physical coefficient matrix, conditional-expectation component,
operator-norm residual, rational card, or cofinal decay has been
populated.  Hilbert--Schmidt smallness is not claimed to be
dimension-uniform.  No continuum fermion dynamics or full
SM--Einstein continuum is proved.
\end{firewall}

\subsection{Next bottleneck}

The next step is representation extraction.  Determine the actual
finite decomposition of the old operation algebra inside the new
RPESM carrier, including multiplicities and central blocks, and
express the complete graph-regulated coefficient in that basis.
Only then can (46.7) be evaluated.  If the embedding is available
only up to a source-fixing unitary, the next note should prove
unitary invariance of the four-component card and formulate a
basis-free computation from the commutant projection.

\subsection{Parent record}

This V46 note appends to the validated thirteenth-cycle V45 file with
record
\[
\begin{array}{c|c}
\text{lines}&14788\\
\text{bytes}&522990\\
\text{SHA--256}&
\texttt{8084C3BB2B53B6C018E0F9B169D18530E667A356F3E8A7516CCDCAD39EB77BEA}.
\end{array}
\]

% =============================================================================
% END APPENDED THIRTEENTH-CYCLE V46 MATERIAL
% =============================================================================

% =============================================================================
% APPENDED THIRTEENTH-CYCLE V47 MATERIAL
% =============================================================================

\section{Thirteenth-cycle V47: basis-free commutant projection and
source-fixing invariance}
\label{sec:v13v47}

\subsection{Why the total residual must be intrinsic}

The Grand Tensor finite landing is unique up to a unitary fixing
the common source vectors.  Therefore a cofinal obstruction may not
depend on a chosen matrix-unit realization.  The refined
coarse/detail split of V46 is useful after a tensor factorization is
fixed, but the total visible commutator component has a
basis-free definition for every finite represented algebra.

Let
\[
 {\cal A}\subset B({\cal H})                                   \tag{47.1}
\]
be a finite-dimensional unital star subalgebra, with the ordinary
matrix trace on \(B({\cal H})\).  Let
\[
 {\mathbb E}_{{\cal A}'}:B({\cal H})\longrightarrow{\cal A}'   \tag{47.2}
\]
be the trace-preserving conditional expectation onto the
commutant.  For \(X=X^*\), define
\[
 R_{\cal A}(X):=(I-{\mathbb E}_{{\cal A}'})X,\qquad
 \rho_{\cal A}(X):=\|R_{\cal A}(X)\|_{\rm HS}.                 \tag{47.3}
\]

\subsection{Wedderburn formula}

Choose a unitary realization
\[
 {\cal H}\simeq
 \bigoplus_{\lambda=1}^r
 {\mathbb C}^{d_\lambda}\otimes{\mathbb C}^{k_\lambda},
 \qquad
 {\cal A}\simeq
 \bigoplus_{\lambda=1}^r
 \bigl(M_{d_\lambda}\otimes I_{k_\lambda}\bigr).               \tag{47.4}
\]
Let \(P_\lambda\) be the central block projections.

\begin{theorem}[Finite commutant expectation]
\label{thm:v13v47-expectation}
In the realization (47.4),
\[
 \boxed{\quad
 {\mathbb E}_{{\cal A}'}(X)
 =
 \bigoplus_{\lambda=1}^r
 \left[
 {I_{d_\lambda}\over d_\lambda}
 \otimes
 {\rm Tr}_{d_\lambda}(P_\lambda XP_\lambda)
 \right].\quad}                                                \tag{47.5}
\]
All off-central-block entries \(P_\lambda XP_\mu\),
\(\lambda\neq\mu\), are removed.  The map is the unique
trace-preserving conditional expectation onto \({\cal A}'\) and
is the Hilbert--Schmidt orthogonal projection.
\end{theorem}

\begin{proof}
Pinching by the central projections is a trace-preserving
conditional expectation onto the block-diagonal algebra.  On a
single block, the normalized partial trace in (47.5) is the
trace-preserving conditional expectation onto
\(I_{d_\lambda}\otimes M_{k_\lambda}\), as proved in V46.
Their direct sum is unital, completely positive, trace preserving,
idempotent, and has range \({\cal A}'\).  Trace duality makes it
the Hilbert--Schmidt orthogonal projection.  A trace-preserving
conditional expectation onto a fixed finite-dimensional
star subalgebra is unique because its Hilbert--Schmidt projection
is unique.
\end{proof}

Thus every part of (47.3) is computable from the central block
projections and partial traces.  No choice of matrix units inside a
simple block affects the answer.

\subsection{Haar commutator variance}

Let \({\cal U}({\cal A})\) be the compact unitary group of
\({\cal A}\), with normalized Haar measure.

\begin{theorem}[Basis-free commutator-energy identity]
\label{thm:v13v47-haar}
For every \(X\in B({\cal H})\),
\[
 \boxed{\quad
 \int_{{\cal U}({\cal A})}
 \|[X,U]\|_{\rm HS}^2\,dU
 =
 2\|X-{\mathbb E}_{{\cal A}'}X\|_{\rm HS}^2.\quad}             \tag{47.6}
\]
Consequently
\[
 \rho_{\cal A}(X)=0
 \quad\Longleftrightarrow\quad
 [X,A]=0\ \hbox{for every }A\in{\cal A}.                       \tag{47.7}
\]
\end{theorem}

\begin{proof}
Multiplication by \(U^*\) preserves the Hilbert--Schmidt norm, so
\[
 \|[X,U]\|_{\rm HS}
 =
 \|X-UXU^*\|_{\rm HS}.                                        \tag{47.8}
\]
The Haar twirl
\[
 {\cal T}(X)=\int_{{\cal U}({\cal A})}UXU^*\,dU                \tag{47.9}
\]
is the trace-preserving conditional expectation onto the fixed
space of the conjugation action, which is \({\cal A}'\).
Expanding (47.8), integrating, and using that
\({\mathbb E}_{{\cal A}'}\) is an orthogonal projection gives
\[
\begin{split}
 \int\|X-UXU^*\|_{\rm HS}^2\,dU
 &=2\|X\|_{\rm HS}^2
   -2\langle X,{\mathbb E}_{{\cal A}'}X\rangle_{\rm HS}\\
 &=2\bigl(\|X\|_{\rm HS}^2
          -\|{\mathbb E}_{{\cal A}'}X\|_{\rm HS}^2\bigr),
\end{split}                                                     \tag{47.10}
\]
which is (47.6).  Its zero branch is exactly the fixed space.
\end{proof}

Equation (47.6) generalizes (46.13) without dimension-dependent
weights or a matrix-unit choice.

\subsection{A finite rational twirl}

For \(d\geq2\), let \({\cal G}_d\) be the finite group of signed
permutation matrices on \({\mathbb C}^d\).  Its entries lie in
\(\{0,\pm1\}\), and its commutant consists only of scalars.
For \(d=1\), use the trivial group.  On the decomposition (47.4),
let \({\cal G}_{\cal A}\) be generated by:
\begin{enumerate}
\item the blockwise actions
\(G_\lambda\otimes I_{k_\lambda}\), \(G_\lambda\in{\cal G}_{d_\lambda}\);
\item independent central signs
\(\sum_\lambda\epsilon_\lambda P_\lambda\),
\(\epsilon_\lambda\in\{\pm1\}\).
\end{enumerate}

\begin{proposition}[Exact finite twirl]
\label{prop:v13v47-finite-twirl}
One has
\[
 {\mathbb E}_{{\cal A}'}(X)
 =
 {1\over|{\cal G}_{\cal A}|}
 \sum_{G\in{\cal G}_{\cal A}}GXG^*.                            \tag{47.11}
\]
Consequently
\[
 {1\over|{\cal G}_{\cal A}|}
 \sum_{G\in{\cal G}_{\cal A}}\|[X,G]\|_{\rm HS}^2
 =
 2\rho_{\cal A}(X)^2.                                         \tag{47.12}
\]
For Gaussian-rational \(X\), both identities are exact rational
certificates.
\end{proposition}

\begin{proof}
Averaging over independent central signs kills every
off-central-block entry.  Within block \(\lambda\), averaging over
signed diagonal matrices kills off-diagonal first-factor entries,
and averaging over permutations makes the surviving diagonal
entries equal.  The result is
\(I_{d_\lambda}/d_\lambda\otimes{\rm Tr}_{d_\lambda}\), which is
(47.5).  This proves (47.11).  The proof of (47.6), with a finite
average in place of Haar measure, gives (47.12).
\end{proof}

The full group need not be enumerated in an implementation:
formula (47.5) computes the same twirl using partial traces.
Equation (47.11) is the exact finite witness that this computation
is representation intrinsic.

\subsection{Unitary transport}

Let \(V:{\cal H}\to\widetilde{\cal H}\) be unitary and put
\[
 \widetilde{\cal A}=V{\cal A}V^*,\qquad
 \widetilde X=VXV^*.                                          \tag{47.13}
\]

\begin{theorem}[Source-fixing unitary invariance]
\label{thm:v13v47-unitary}
The commutant expectations obey
\[
 {\mathbb E}_{\widetilde{\cal A}'}
 (VZV^*)
 =
 V{\mathbb E}_{{\cal A}'}(Z)V^*.                              \tag{47.14}
\]
Hence
\[
 R_{\widetilde{\cal A}}(\widetilde X)
 =
 VR_{\cal A}(X)V^*,                                           \tag{47.15}
\]
and every unitarily invariant norm, the Haar variance (47.6), the
zero verdict, and the spectrum of
\(R_{\cal A}(X)^*R_{\cal A}(X)\) are unchanged.
\end{theorem}

\begin{proof}
The right side of (47.14) is a trace-preserving conditional
expectation onto
\(V{\cal A}'V^*=(V{\cal A}V^*)'\).  Uniqueness in
Theorem~\ref{thm:v13v47-expectation} proves (47.14).  Subtraction
gives (47.15), and unitary invariance gives the remaining claims.
\end{proof}

\begin{corollary}[Well-defined card on the landed source class]
\label{cor:v13v47-landed}
Suppose two finite Grand Tensor landings are related by the
source-fixing unitary guaranteed by the finite Euclidean/OS
landing theorem, and the calibrated coefficient and operation
algebra are transported by that unitary.  Then they produce the
same V46--V47 adjacent-germ verdict and identical residual singular
values.
\end{corollary}

\begin{proof}
Apply Theorem~\ref{thm:v13v47-unitary}.
\end{proof}

Thus the source-minimal unitary ambiguity is not an obstruction to
evaluating the card.

\subsection{Central phase and commutant gauge invariance}

For \(C=C^*\in{\cal A}'\) and \(c\in{\mathbb R}\),
\[
 R_{\cal A}(X+C+cI)=R_{\cal A}(X).                             \tag{47.16}
\]
Therefore scalar phase choices and detail-only Hamiltonians do not
change the old-algebra germ residual.  This local statement does
not orient the determinant line globally; it only shows that an
already selected phase representative cannot affect the
commutator card through a scalar shift.

\subsection{Basis-free failure witness}

If \(R=R_{\cal A}(X)\neq0\), Theorem
\ref{thm:v13v47-haar} implies the existence of
\(U_*\in{\cal U}({\cal A})\) such that
\[
 \|[X,U_*]\|_{\rm HS}
 \geq\sqrt2\,\|R\|_{\rm HS}.                                  \tag{47.17}
\]
The same conclusion holds for some member of
\({\cal G}_{\cal A}\) by (47.12).

\begin{proposition}[Finite intrinsic obstruction writer]
\label{prop:v13v47-obstruction}
Choose \(G_*\in{\cal G}_{\cal A}\) maximizing the commutator norm
and set
\[
 Q_*=-i[X,G_*].                                                \tag{47.18}
\]
Then \(Q_*\neq0\) on every failure branch, its positive Gram
\[
 {\mathfrak R}_*=|Q_*\rangle_{\rm HS}\langle Q_*|              \tag{47.19}
\]
is a source-minimal rank-one witness for the fixed query, and
\[
 \|Q_*\|_{\rm HS}^2\geq2\|R\|_{\rm HS}^2.                     \tag{47.20}
\]
\end{proposition}

\begin{proof}
The maximum of the nonnegative terms in the finite average
(47.12) is at least the average.  This gives (47.20), and all
other claims follow as in Proposition
\ref{prop:v13v46-witness}.
\end{proof}

This witness uses an actual unitary observable in the old operation
algebra rather than a basis-dependent matrix unit.

\subsection{Stable approximate representations}

Suppose approximate data
\((\widetilde{\cal A},\widetilde X)\) are aligned to
\(({\cal A},X)\) by a unitary \(V\) and satisfy
\[
 \|\widetilde X-VXV^*\|_{\rm op}\leq\varepsilon_X.             \tag{47.21}
\]
If the represented algebra is transported exactly, then
contractivity of the two expectations gives
\[
 \|R_{\widetilde{\cal A}}(\widetilde X)
   -VR_{\cal A}(X)V^*\|_{\rm op}
 \leq2\varepsilon_X.                                          \tag{47.22}
\]
Hence a strict residual upper or lower margin exceeding
\(2\varepsilon_X\) is stable.

If the algebra itself is only approximately unitarily aligned,
(47.22) does not follow.  One then needs a Kadison--Kastler or
source-Gram perturbation bound controlling the two conditional
expectations.  That is a separate representation-stability row.

\subsection{Cofinal use}

At adjacent cutoff \(n\), let
\[
 X_n=\widehat D_{n+1}-\iota_n(\widehat D_n)                    \tag{47.23}
\]
in any source-fixed landed representation, and define
\[
 R_n=R_{\iota_n({\cal A}_n)}(X_n).                             \tag{47.24}
\]
By Corollary~\ref{cor:v13v47-landed}, the norm of \(R_n\) is
independent of the landed representative.  V44 gives:
\[
 \|R_n\|_{\rm op}\to0
 \quad\Longrightarrow\quad
 \hbox{adjacent physical germ compatibility},                 \tag{47.25}
\]
and
\[
 \sum_n\|R_n\|_{\rm op}<\infty
 \quad\Longrightarrow\quad
 \hbox{a fixed-time inductive-limit automorphism}.             \tag{47.26}
\]
The finite twirl provides a direct exact residual and a witness at
each level.  Population still requires the actual inclusions and
calibrated coefficients.

\begin{assessment}[Status after thirteenth-cycle V47]
\label{ass:v13v47-status}
The adjacent-germ card is now basis free.  Every finite semisimple
operation algebra has a unique trace-preserving commutant
expectation, computable by central pinching and partial trace.
Its visible coefficient residual equals one half of an exact Haar,
or finite signed-permutation, commutator variance.  The card and
all residual singular values are invariant under the source-fixing
unitary ambiguity of the Grand Tensor landing, and every failure
has a finite intrinsic unitary witness.
\end{assessment}

\begin{firewall}[Claim boundary after thirteenth-cycle V47]
\label{fire:v13v47-claim}
The actual RPESM adjacent representation, central multiplicities,
physical clock calibration, and coefficient matrices remain
unextracted.  Approximate alignment of the represented algebras is
not controlled by (47.22).  No decay or summability of the
intrinsic residual, continuum fermion automorphism, or full
SM--Einstein continuum is proved.
\end{firewall}

\subsection{Next bottleneck}

The algebraic ambiguity has been removed.  The next physical task
is to extract the finite Wedderburn multiplicities and the
source-fixing inclusion from the populated RPESM word algebra.
If only mixed source Grams are directly available, the next note
should give a Gram-only reconstruction of the inclusion and a
perturbation bound for its commutant expectation, with an explicit
rank-loss alternative.

\subsection{Parent record}

This V47 note appends to the validated thirteenth-cycle V46 file with
record
\[
\begin{array}{c|c}
\text{lines}&15212\\
\text{bytes}&536627\\
\text{SHA--256}&
\texttt{1EE2C52C93290E8775975146C95E24DDCD9653B1BB200E9153E84DC823552D60}.
\end{array}
\]

% =============================================================================
% END APPENDED THIRTEENTH-CYCLE V47 MATERIAL
% =============================================================================

% =============================================================================
% APPENDED THIRTEENTH-CYCLE V48 MATERIAL
% =============================================================================

\section{Thirteenth-cycle V48: Gram reconstruction of the supported
operation embedding}
\label{sec:v13v48}

\subsection{Source Grams and the canonical isometry}

Let \(E\) be a finite coefficient space and let
\[
 S:E\to{\cal H}_0,\qquad T:E\to{\cal H}_1                     \tag{48.1}
\]
be source synthesis maps.  Their entrance Grams are
\[
 G_S=S^*S,\qquad G_T=T^*T.                                    \tag{48.2}
\]

\begin{theorem}[Exact source-fixing isometry]
\label{thm:v13v48-source-isometry}
If \(G_S=G_T\), there is a unique isometry
\[
 U:\overline{\operatorname{Ran}S}
 \longrightarrow\overline{\operatorname{Ran}T}                \tag{48.3}
\]
such that
\[
 USc=Tc\qquad(c\in E).                                         \tag{48.4}
\]
It is unitary between the two source ranges.  If
\(G_S=G_T=:G\) is invertible, then
\[
 U=T G^{-1}S^*
 \quad\hbox{on }\operatorname{Ran}S.                           \tag{48.5}
\]
For singular \(G\), replace \(G^{-1}\) by \(G^\dagger\).
\end{theorem}

\begin{proof}
Equality of Grams gives
\[
 \|Sc\|^2=c^*G_Sc=c^*G_Tc=\|Tc\|^2.                           \tag{48.6}
\]
It also gives \(\ker S=\ker T\), so (48.4) is well defined on
\(\operatorname{Ran}S\), preserves inner products by polarization,
and has range \(\operatorname{Ran}T\).  It therefore extends to
the stated unitary.  Formula (48.5) follows from
\(TG^{-1}S^*Sc=Tc\); the Moore--Penrose formula gives the same
identity on the supported coefficient quotient.
\end{proof}

This theorem reconstructs the source-fixing transport without
choosing bases in either physical carrier.

\subsection{Action Grams reconstruct a supported representation}

Let \({\mathfrak A}\) be an abstract finite-dimensional unital
star algebra and let
\[
 \pi_S:{\mathfrak A}\to B({\cal H}_0)                          \tag{48.7}
\]
be a representation.  For \(a\in{\mathfrak A}\), define its
source action Gram
\[
 M_S(a):=S^*\pi_S(a)S.                                        \tag{48.8}
\]
Assume first that \(G_S\succ0\) and that \(S\) is source complete:
\(\operatorname{Ran}S={\cal H}_0\).  Put
\[
 \widehat S=SG_S^{-1/2},\qquad
 \Pi_S(a)=G_S^{-1/2}M_S(a)G_S^{-1/2}.                          \tag{48.9}
\]

\begin{theorem}[Gram-only representation reconstruction]
\label{thm:v13v48-representation}
Under the source-completeness hypothesis, \(\widehat S\) is
unitary and
\[
 \Pi_S(a)=\widehat S^*\pi_S(a)\widehat S.                      \tag{48.10}
\]
Consequently the finite tables \(G_S\) and \(M_S(a)\) on an
algebra basis determine the represented operation algebra up to
the unique source-fixing unitary.

If \(T\) is source complete, \(G_T=G_S\), and
\[
 M_T(a)=M_S(a)\qquad(a\ \hbox{in an algebra basis}),           \tag{48.11}
\]
then the isometry of Theorem
\ref{thm:v13v48-source-isometry} is unitary on the full carriers
and
\[
 \pi_T(a)=U\pi_S(a)U^*
 \qquad(a\in{\mathfrak A}).                                    \tag{48.12}
\]
\end{theorem}

\begin{proof}
Equation \(\widehat S^*\widehat S=I\) and source completeness make
\(\widehat S\) unitary.  Substitution gives (48.10).  Linearity
reconstructs every algebra element from a basis.  Under (48.11),
the two normalized coordinate representations are equal.
Conjugating by
\(U=\widehat T\widehat S^*\) gives (48.12).
\end{proof}

The multiplication, star, unit, grading, and reality identities can
all be checked on the reconstructed matrices \(\Pi_S(a)\).  A
table of expectation values alone is weaker than the complete
action Grams (48.8).

\subsection{A supported corner is not a unital fine representation}

Suppose \(S\) is source complete in \({\cal H}_0\), but
\(\operatorname{Ran}T={\cal K}\) is a proper subspace of
\({\cal H}_1\).  Let \(P_T\) be its projection.  Equal entrance
and action Grams determine
\[
 P_T\pi_T(a)P_T
 =
 U\pi_S(a)U^*                                                  \tag{48.13}
\]
provided \({\cal K}\) reduces \(\pi_T({\mathfrak A})\).

The reducing property has its own positive leakage:
\[
 {\cal L}_T(a)
 :=
 (I-P_T)\pi_T(a)T,\qquad
 \Delta_T(a):={\cal L}_T(a)^*{\cal L}_T(a)\succeq0.            \tag{48.14}
\]

\begin{proposition}[Supported-subrepresentation gate]
\label{prop:v13v48-supported-gate}
For a star-generating family \({\cal G}\subset{\mathfrak A}\),
the source range \({\cal K}\) reduces the represented algebra if
and only if
\[
 \Delta_T(a)=0\qquad(a\in{\cal G}).                            \tag{48.15}
\]
On the zero branch, the full fine representation has the form
\[
 \pi_T(a)
 =
 U\pi_S(a)U^*\oplus\sigma(a)                                  \tag{48.16}
\]
on \({\cal K}\oplus{\cal K}^\perp\), where
\(\sigma\) is a separate representation on the new-detail
complement.
\end{proposition}

\begin{proof}
Vanishing of (48.14) is equivalent to
\((I-P_T)\pi_T(a)P_T=0\) because \(T\) spans \({\cal K}\).
Applying the same condition to \(a^*\) gives the opposite
off-diagonal block.  A star-generating family then makes
\({\cal K}\) reducing for the entire algebra.  Block diagonalizing
gives (48.16); the lower block is a unital star representation.
\end{proof}

Thus source Grams reconstruct the old supported corner.  They do
not determine the representation \(\sigma\) carried only by new
detail.  Its Wedderburn multiplicities are independent population
data needed for the unital inclusion in V44--V47.

\subsection{Multiplicity extraction}

Write
\[
 {\mathfrak A}\simeq\bigoplus_{\lambda=1}^rM_{d_\lambda},
 \qquad z_\lambda\ \hbox{its minimal central units}.           \tag{48.17}
\]
For a representation \(\sigma\) on the detail complement,
\[
 \sigma\simeq
 \bigoplus_{\lambda=1}^r
 \bigl({\rm id}_{d_\lambda}\otimes I_{m_\lambda}\bigr).         \tag{48.18}
\]

\begin{proposition}[Exact multiplicity card]
\label{prop:v13v48-multiplicity}
The multiplicities are
\[
 \boxed{\quad
 m_\lambda={1\over d_\lambda}
 {\rm rank}\,\sigma(z_\lambda).\quad}                          \tag{48.19}
\]
They are nonnegative integers.  Conversely the list
\((m_\lambda)_\lambda\) determines the complement representation
up to unitary equivalence.
\end{proposition}

\begin{proof}
The central projection \(\sigma(z_\lambda)\) has range
\({\mathbb C}^{d_\lambda}\otimes{\mathbb C}^{m_\lambda}\), so its
rank is \(d_\lambda m_\lambda\).  The classification of
finite-dimensional star representations gives the converse.
\end{proof}

If a claimed central projection has rank not divisible by
\(d_\lambda\), it is an immediate finite obstruction to the
declared representation type.

\subsection{Rank-loss alternative}

If the source Gram varies with the cutoff, a source-fixing isometry
can fail before any algebra comparison.

\begin{proposition}[Entrance-rank alternative]
\label{prop:v13v48-rank}
Let \(G_S,G_T\succeq0\).
\begin{enumerate}
\item If \(\ker G_S=\ker G_T\) and the two Grams agree on their
common support, Theorem~\ref{thm:v13v48-source-isometry} applies
on the quotient.
\item If the kernels differ, any
\[
 c\in\ker G_S\setminus\ker G_T
 \quad\hbox{or}\quad
 c\in\ker G_T\setminus\ker G_S                                \tag{48.20}
\]
is a direct source-collapse or source-birth witness.
\item If \(G_S\succeq\kappa I\) and
\(\|G_T-G_S\|_{\rm op}<\kappa\), then both Grams are invertible
and
\[
 G_T\succeq
 (\kappa-\|G_T-G_S\|_{\rm op})I.                               \tag{48.21}
\]
\end{enumerate}
\end{proposition}

\begin{proof}
The first statement repeats the quotient part of Theorem
\ref{thm:v13v48-source-isometry}.  In the second, (48.6) gives
zero norm on one carrier and positive norm on the other.
The last statement is Weyl's inequality.
\end{proof}

A cutoff-dependent positive eigenvalue with no uniform lower floor
does not give stable normalized source frames.

\subsection{Perturbation of normalized action tables}

Suppose \(G,\widetilde G\succeq\kappa I\).  Functional calculus
gives
\[
 \|G^{-1/2}-\widetilde G^{-1/2}\|_{\rm op}
 \leq
 {1\over2\kappa^{3/2}}\|G-\widetilde G\|_{\rm op}.             \tag{48.22}
\]
For action Grams \(M,\widetilde M\), the normalized tables obey
\[
\begin{split}
 &\|G^{-1/2}MG^{-1/2}
 -\widetilde G^{-1/2}\widetilde M\widetilde G^{-1/2}\|\\
 &\quad\leq
 \kappa^{-1}\|M-\widetilde M\|
 +
 {\|M\|+\|\widetilde M\|\over2\kappa^2}
 \|G-\widetilde G\|.                                          \tag{48.23}
\end{split}
\]

\begin{proof}
The integral representation of \(x^{-1/2}\) on
\([\kappa,\infty)\) gives (48.22).  Insert and subtract the two
mixed products in (48.23), use
\(\|G^{-1/2}\|,\|\widetilde G^{-1/2}\|\leq\kappa^{-1/2}\), and
apply (48.22).  Enlarging the two asymmetric coefficients to their
displayed sum gives (48.23).
\end{proof}

This turns entrance and action Gram errors into an explicit
representation-coordinate error, provided the lower frame floor is
certified.

\subsection{Perturbation of the commutant expectation}

Let \(\Pi,\widetilde\Pi\) be exact representations of
\({\mathfrak A}\) on one coordinate carrier.  Choose the finite
twirl group \({\cal G}_{\mathfrak A}\) of V47 and assume
\[
 \max_{g\in{\cal G}_{\mathfrak A}}
 \|\widetilde\Pi(g)-\Pi(g)\|_{\rm op}\leq\eta.                 \tag{48.24}
\]

\begin{theorem}[Finite-twirl expectation stability]
\label{thm:v13v48-expectation-stability}
For every \(Z\),
\[
 \|{\mathbb E}_{\widetilde\Pi({\mathfrak A})'}(Z)
   -{\mathbb E}_{\Pi({\mathfrak A})'}(Z)\|_{\rm op}
 \leq2\eta\|Z\|_{\rm op}.                                     \tag{48.25}
\]
If \(\|\widetilde X-X\|_{\rm op}\leq\varepsilon_X\), then
\[
 \|R_{\widetilde\Pi({\mathfrak A})}(\widetilde X)
   -R_{\Pi({\mathfrak A})}(X)\|_{\rm op}
 \leq
 2\varepsilon_X+2\eta\|X\|_{\rm op}.                           \tag{48.26}
\]
\end{theorem}

\begin{proof}
Use the common finite twirl:
\[
 {\mathbb E}_{\Pi({\mathfrak A})'}(Z)
 ={1\over|{\cal G}|}\sum_g\Pi(g)Z\Pi(g)^*,                     \tag{48.27}
\]
and similarly for \(\widetilde\Pi\).  Each conjugation difference
is at most \(2\eta\|Z\|\), proving (48.25).  For (48.26), split
the residual difference into the coefficient perturbation and the
expectation perturbation.  The map \(I-{\mathbb E}\) has norm at
most two, giving \(2\varepsilon_X\); use (48.25) on \(X\) for the
remaining term.
\end{proof}

The error \(\eta\) can be bounded from the normalized action tables
by (48.23), after evaluating the finite group words.  No
Kadison--Kastler hypothesis remains hidden.

\subsection{The Gram-only adjacent-card interface}

For an adjacent RPESM pair, a complete finite input now consists
of:
\[
\begin{array}{c|l}
G_S,G_T&\text{entrance source Grams},\\
M_S(a),M_T(a)&\text{complete action Grams on generators},\\
\Delta_T(a)&\text{supported-range leakage Grams},\\
{\rm rank}\,\sigma(z_\lambda)&\text{new-detail multiplicities},\\
X&\text{calibrated complete coefficient mismatch}.
\end{array}                                                    \tag{48.28}
\]
Zero entrance mismatch, zero leakage, integral multiplicities, and
the reconstructed action tables determine the unital represented
embedding up to a source-fixing unitary.  V47 then computes the
basis-free commutant residual, and V44 transports the channel germ.

\begin{assessment}[Status after thirteenth-cycle V48]
\label{ass:v13v48-status}
Complete source and action Grams reconstruct the operation
representation on a source-complete carrier and the source-fixing
isometry between equal-Gram carriers.  For a proper supported
corner, the leakage Grams decide whether it is a subrepresentation,
while the complement Wedderburn multiplicities are independent
data.  Rank loss has an explicit coefficient witness, and stable
frame floors turn approximate Grams into a quantitative
commutant-expectation bound.
\end{assessment}

\begin{firewall}[Claim boundary after thirteenth-cycle V48]
\label{fire:v13v48-claim}
The RPESM entrance Grams, action Grams, leakage panels, central
multiplicities, and calibrated coefficient mismatch have not been
numerically or symbolically extracted from the event grammar.
The source Gram alone does not determine the new-detail
representation.  No cofinal frame floor, commutant residual decay,
continuum fermion dynamics, or full SM--Einstein continuum is
proved.
\end{firewall}

\subsection{Next bottleneck}

The remaining finite population request is now minimal: acquire
the table (48.28) from one adjacent RPESM event grammar.  If direct
matrix entries are unavailable, the next note should show how the
existing held-out word and mixed-incidence writers determine each
entry of (48.28), and identify which single table entry is first
absent.  No further abstract representation theorem is needed
before that audit.

\subsection{Parent record}

This V48 note appends to the validated thirteenth-cycle V47 file with
record
\[
\begin{array}{c|c}
\text{lines}&15604\\
\text{bytes}&549849\\
\text{SHA--256}&
\texttt{CEEC4413B5D50DB7BD067201D520163B65617CCE9A06FC2B80E0E7940467995E}.
\end{array}
\]

% =============================================================================
% END APPENDED THIRTEENTH-CYCLE V48 MATERIAL
% =============================================================================

% =============================================================================
% APPENDED THIRTEENTH-CYCLE V49 MATERIAL
% =============================================================================

\section{Thirteenth-cycle V49: Gram tomography of the complete
adjacent operation card}
\label{sec:v13v49}

\subsection{What was found in the populated event grammar}

The attached RPESM retained-terminal theorem supplies, for every
frozen bounded grammar, one-history direct/staged word equality and
a summable endpoint debit.  Its complete mixed-short theorem
supplies the joint entrance Gram
\[
 \begin{pmatrix}S&T\\T^*&R\end{pmatrix}\succeq0                \tag{49.1}
\]
and the exact follower/innovation Schur complement.  The finite
Euclidean landing supplies a source-minimal OS carrier and a
represented operation algebra, unique up to a source-fixing
unitary.

These results prove occurrence of the finite ingredients.  They do
not display one adjacent, clock-calibrated, source-complete matrix
card.  This note proves that such a card can be computed entirely
from Gram and incidence writers; raw carrier coordinates are not
needed.

\subsection{Operator tomography from one complete source frame}

Let \(S:E\to{\cal H}\) be bijective onto a finite carrier and put
\[
 G=S^*S\succ0,\qquad \widehat S=SG^{-1/2}.                     \tag{49.2}
\]
For \(A\in B({\cal H})\), define its incidence Gram
\[
 K_S(A):=S^*AS.                                                \tag{49.3}
\]

\begin{theorem}[Complete-frame operator tomography]
\label{thm:v13v49-tomography}
The map \(\widehat S:E\to{\cal H}\) is unitary and the matrix of
\(A\) in the normalized source frame is
\[
 \boxed{\quad
 [A]_S:=\widehat S^*A\widehat S
 =G^{-1/2}K_S(A)G^{-1/2}.\quad}                                \tag{49.4}
\]
Consequently:
\[
 [A^*]_S=[A]_S^*,\qquad
 [AB]_S=[A]_S[B]_S,                                            \tag{49.5}
\]
and every spectrum, rank, singular value, unitarily invariant norm,
commutator, and polynomial relation of \(A\) is determined by
\(G\) and \(K_S(A)\).
\end{theorem}

\begin{proof}
Equation \(G=S^*S\) gives
\(\widehat S^*\widehat S=I\).  Bijectivity makes
\(\widehat S\) unitary.  Substitution proves (49.4), and unitary
conjugation proves all remaining statements.
\end{proof}

Equivalently, if incidence data for products are recorded before
normalization, the exact held-out multiplication residual is
\[
 K_S(AB)-K_S(A)G^{-1}K_S(B).                                  \tag{49.6}
\]
It vanishes for a complete frame.  A nonzero value proves either
an incomplete source carrier, inconsistent product data, or a
cross-history assembly.

\subsection{Supported frames}

If \(S\) is injective but not onto \({\cal H}\), formula (49.4)
reconstructs only the compression \(P_SAP_S\), where
\[
 P_S=SG^{-1}S^*.                                               \tag{49.7}
\]
The missing off-support part has positive leakage
\[
 L_S(A)
 :=
 S^*A^*(I-P_S)AS\succeq0.                                    \tag{49.8}
\]

\begin{proposition}[Tomography completeness gate]
\label{prop:v13v49-completeness}
For a star-generating set \({\cal G}\), the source range reduces
the represented algebra if and only if
\[
 L_S(a)=0\qquad(a\in{\cal G}).                                \tag{49.9}
\]
On the zero branch, (49.4) reconstructs the complete cyclic
subrepresentation.  It contains no information about a
source-orthogonal representation on \((I-P_S){\cal H}\).
\end{proposition}

\begin{proof}
This is Proposition~\ref{prop:v13v48-supported-gate} written only
in Gram coordinates.  The Gram in (49.8) vanishes exactly when
\((I-P_S)AS=0\); applying the same condition to adjoints makes the
range reducing.
\end{proof}

Thus the source-minimal landing can be reconstructed exactly, but
a claim about the full physical carrier requires an independent
source-completeness or complement-multiplicity row.

\subsection{Adjacent normalized inclusion}

Let
\[
 S_0:E_0\to{\cal H}_0,\qquad
 S_1:E_1\to{\cal H}_1                                         \tag{49.10}
\]
be complete frames with Grams \(G_0,G_1\succ0\).  Suppose the
declared old-source label injection is
\[
 J:E_0\to E_1                                                  \tag{49.11}
\]
and satisfies
\[
 J^*G_1J=G_0.                                                  \tag{49.12}
\]
Define
\[
 V:=G_1^{1/2}JG_0^{-1/2}.                                     \tag{49.13}
\]
Then \(V^*V=I\).

Let \(M_{r}(a)=S_r^*\pi_r(a)S_r\) be action Grams and put
\[
 \Pi_r(a)=G_r^{-1/2}M_r(a)G_r^{-1/2}.                         \tag{49.14}
\]

\begin{theorem}[Gram-only operation inclusion]
\label{thm:v13v49-inclusion}
The normalized source inclusion \(V\) intertwines the operation
representations exactly if and only if
\[
 \Pi_1(a)V=V\Pi_0(a)
 \qquad(a\in{\cal G})                                         \tag{49.15}
\]
for a star-generating family \({\cal G}\).  Its complete positive
intertwining residual is
\[
 {\cal I}(a)
 :=
 (\Pi_1(a)V-V\Pi_0(a))^*
 (\Pi_1(a)V-V\Pi_0(a))\succeq0.                               \tag{49.16}
\]
On the zero branch, \(\operatorname{Ran}V\) reduces
\(\Pi_1({\mathfrak A})\), and the complement representation and
all of its Wedderburn multiplicities are obtained from
\[
 (I-VV^*)\Pi_1(z_\lambda)(I-VV^*)                             \tag{49.17}
\]
by the rank formula (48.19).
\end{theorem}

\begin{proof}
Equation (49.12) gives \(V^*V=I\).  Vanishing of (49.16) is
equivalent to (49.15).  The same relation for adjoint generators
makes \(\operatorname{Ran}V\) reducing.  The central units then
restrict to projections on the complement; their ranks give the
multiplicities by Proposition
\ref{prop:v13v48-multiplicity}.
\end{proof}

Every term in (49.12)--(49.17) is a finite Gram or incidence
matrix.

\subsection{Coefficient tomography and the commutant card}

Let \(D_r=D_r^*\) be the finite landed coefficients and record
\[
 H_r:=K_{S_r}(D_r)=S_r^*D_rS_r.                               \tag{49.18}
\]
For physical clock calibrations \(c_r>0\), define
\[
 {\mathsf D}_r
 :=
 c_r^{-1}G_r^{-1/2}H_rG_r^{-1/2}.                             \tag{49.19}
\]
The complete calibrated adjacent mismatch in fine normalized
coordinates is
\[
 {\mathsf X}
 :=
 {\mathsf D}_1-V{\mathsf D}_0V^*.                             \tag{49.20}
\]
Let
\[
 {\mathfrak A}_1:=C^*(\Pi_1(a):a\in{\cal G}).                  \tag{49.21}
\]

\begin{theorem}[Pure-Gram adjacent germ card]
\label{thm:v13v49-pure-gram-card}
Assume (49.12) and the zero intertwining branch (49.16).  Then
\[
 \boxed{\quad
 {\mathsf R}
 =
 {\mathsf X}
 -
 {\mathbb E}_{{\mathfrak A}_1'}({\mathsf X})\quad}             \tag{49.22}
\]
is the basis-free calibrated adjacent germ residual of V47,
computed entirely from
\[
 (G_0,G_1,J,\{M_0(a),M_1(a)\}_{a\in{\cal G}},
   H_0,H_1,c_0,c_1).                                          \tag{49.23}
\]
It has the exact zero criterion
\[
 {\mathsf R}=0
 \quad\Longleftrightarrow\quad
 [-i{\mathsf D}_1,\,\cdot\,]V
 =
 V[-i{\mathsf D}_0,\,\cdot\,]                                 \tag{49.24}
\]
on the embedded operation algebra, with the natural interpretation
of \(V\) as the representation intertwiner.  Every nonzero branch
has the finite-twirl unitary witness of V47.
\end{theorem}

\begin{proof}
Theorem~\ref{thm:v13v49-tomography} gives the calibrated coefficient
matrices (49.19).  Theorem~\ref{thm:v13v49-inclusion} gives the
represented fine algebra and its old reducing subspace.
The conditional expectation in (49.22) is reconstructed from the
central decomposition of (49.21), or by the finite twirl (47.11).
The V47 zero criterion is precisely equality of the two
commutator derivations, which gives (49.24).
\end{proof}

This proves that no raw choice of carrier basis, Kraus frame, or
determinant-line scalar phase is needed for the finite adjacent
germ verdict.

\subsection{A pivoted source-basis compiler}

A finite populated grammar may provide a redundant ordered source
list \(v_1,\ldots,v_N\).  Let \(G=[\langle v_i,v_j\rangle]\).

\begin{proposition}[Canonical Gram pivot]
\label{prop:v13v49-pivot}
There is a deterministic source basis obtained by scanning the
ordered list and retaining \(v_j\) exactly when its Schur pivot
against the previously retained vectors is positive.  The number
of retained vectors is \({\rm rank}\,G\), their principal Gram is
positive definite, and they span the original source range.
A zero pivot returns the exact follower coefficients through the
Moore--Penrose formula.
\end{proposition}

\begin{proof}
At each step, the Schur pivot is the squared norm of the component
of \(v_j\) orthogonal to the retained span.  It is positive exactly
when the span grows.  Orthogonal Gram--Schmidt therefore proves
independence, spanning, and the rank count.  When the pivot is
zero, orthogonal projection onto the retained synthesis gives the
follower coefficients.
\end{proof}

This selects a finite source frame from the RPESM word grammar
without an external basis choice.  Cofinal stability still
requires its positive pivots to have lower bounds.

\subsection{Population verdict for the attached RPESM sequence}

The existing finite theorems supply:
\[
\begin{array}{c|c|l}
\text{datum}&\text{finite status}&\text{reason}\\ \hline
G_r&\text{occurs}&\text{reflected and mixed entrance Grams}\\
M_r(a)&\text{occurs for a frozen grammar}
 &\text{complete terminal mixed words}\\
J&\text{occurs on one terminal tower}
 &\text{direct/staged path inclusion}\\
H_r&\text{formally writable}
 &\text{coefficient and energy--Hamiltonian incidence}\\
c_r&\text{not fixed by the algebraic landing}
 &\text{physical clock calibration is independent}.
\end{array}                                                    \tag{49.25}
\]
The retained-terminal bound controls a countable bounded path
grammar.  The same theorem explicitly leaves common domains,
uniform integrability, and Mosco survival of unbounded forms as
continuum rows.  Therefore it does not by itself bound the
coefficient incidence \(H_r\) or the inverse source frames
uniformly.

\begin{proposition}[First absent quantitative row]
\label{prop:v13v49-first-absent}
At each fixed cutoff, the pivot compiler and the populated finite
writers suffice to define an uncalibrated source-minimal card.
The first absent input to the physical V49 card is the common clock
calibration \(c_r\) relating the marked channel duration to the
coefficient germ.  After a calibration is chosen, the next absent
cofinal inputs are:
\[
 \inf_r\lambda_{\min}(G_r|_{\rm pivot})>0,\qquad
 \text{uniform integrability/common-domain transport of }H_r,
                                                                    \tag{49.26}
\]
and decay or summability of the resulting \(\|{\mathsf R}_r\|\).
\end{proposition}

\begin{proof}
All algebraic formulas through (49.22) are finite once the listed
Grams and incidences are evaluated.  Equation (49.19) cannot be
formed physically without \(c_r\).  Even after it is formed,
inverse square roots require the frame floor and passage of the
unbounded coefficient incidence requires the declared analytic
rows.  V44 then requires residual decay or summability.
\end{proof}

\subsection{Finite acquisition packet}

The minimal next acquisition on one adjacent pair is therefore
\[
 {\cal P}_{r+1/r}^{\rm germ}
 =
 \bigl(
 G_r,G_{r+1},J_r,
 \{M_r(a),M_{r+1}(a)\}_{a\in{\cal G}},
 H_r,H_{r+1},c_r,c_{r+1}
 \bigr).                                                       \tag{49.27}
\]
The compiler is:
\[
 \boxed{
 \text{pivot}\ \longrightarrow\
 \text{normalize}\ \longrightarrow\
 \text{test (49.16)}\ \longrightarrow\
 \text{form (49.20)}\ \longrightarrow\
 \text{twirl (49.22)}.}                                       \tag{49.28}
\]
A positive entrance or intertwining residual stops before the
coefficient comparison.  A nonzero final residual returns the
intrinsic unitary witness of Proposition
\ref{prop:v13v47-obstruction}.

\begin{assessment}[Status after thirteenth-cycle V49]
\label{ass:v13v49-status}
The V48 acquisition table can now be compiled from physical Grams
alone.  Complete source frames reconstruct operators, normalized
source inclusions reconstruct the adjacent representation and
complement multiplicities, and coefficient incidence Grams give
the calibrated mismatch whose commutant twirl is the V47 germ
residual.  The RPESM grammar contains the finite algebraic writers,
but the physical clock calibration and the cofinal source-frame and
unbounded-form controls remain absent.
\end{assessment}

\begin{firewall}[Claim boundary after thirteenth-cycle V49]
\label{fire:v13v49-claim}
No numerical RPESM Gram card, clock calibration, uniform pivot
floor, unbounded coefficient-integrability bound, or adjacent
commutant residual has been populated.  The retained-terminal
bounded-path debit is not applied to the unbounded coefficient
without those hypotheses.  The source-minimal card does not prove
completeness of a source-orthogonal physical sector.  No continuum
fermion dynamics or full SM--Einstein continuum is proved.
\end{firewall}

\subsection{Next bottleneck}

The next note must attack the clock calibration rather than
representation theory.  Compare the marked channel parameter, the
RPESM Poisson duration, the symmetric OS Markov time, and the
coefficient normalization.  A calibration is valid only if a
same-history mixed timing Gram identifies their first derivatives;
equality of their marginal duration laws is insufficient.

\subsection{Parent record}

This V49 note appends to the validated thirteenth-cycle V48 file with
record
\[
\begin{array}{c|c}
\text{lines}&15995\\
\text{bytes}&563122\\
\text{SHA--256}&
\texttt{C32E3B83ABF1F394AA9B2F6EFA83042F0D504E09E03CEE82120D09E32F47DD9A}.
\end{array}
\]

% =============================================================================
% END APPENDED THIRTEENTH-CYCLE V49 MATERIAL
% =============================================================================

% =============================================================================
% BEGIN APPENDED THIRTEENTH-CYCLE V50 MATERIAL
% =============================================================================

\section{Thirteenth-cycle V50: compulsory companion audit and the
clock-calibration decision}
\label{sec:v13v50}

\subsection{Read-only record of the companion programmes}

The newest numerical Yang--Mills note and the newest
Navier--Stokes note in the shared directory were inspected
read-only:
\[
\begin{array}{c|c|r|r|c}
 \text{programme}&\text{version}&\text{lines}&\text{bytes}&
 \text{SHA--256}\\ \hline
 {\rm Yang\!-\!Mills}&V86&28920&986388&
 \texttt{F2C8D2B226997064E1CA581432E2AFB4FA8EF34C14BB97BB11A6F23D40694554}\\
 {\rm Navier\!-\!Stokes}&V26&5319&278283&
 \texttt{82C887F8C2C69E4F28FAD7C3DC8DE5B9099CB5A4BF431EBA1FFF3E91935978AD}.
\end{array}                                                     \tag{50.1}
\]
The auxiliary parallel Yang--Mills note remains at V5 and its
global charge obstruction has not changed.  No file in a companion
programme was modified.

Yang--Mills V86 proves a two-sided Banach-residual estimate.  If
\(H\succ0\), \(G=H^{-1}\), \(E_L=I-RH\), and \(E_R=I-HR\), then
\[
 \|H^{-1/2}YH^{-1/2}\|_F^2
 \leq
 \frac{\|RY\|_F\|YR\|_F}
 {(1-\|E_L\|_{\rm op})(1-\|E_R\|_{\rm op})}                    \tag{50.2}
\]
whenever both residual norms are below one.  Its exact
interval-product cards reduce the certified all-minus path lengths
to \(34.171961\) and \(54.878727\) at \(256\) cells.  That note also
identifies the remaining loss: even the zero-width limit of the
separately normed two-sided products is about three times the
dressed Hilbert--Schmidt diagnostic.  Its next move is therefore to
retain the complete trial trace before applying an ambient norm.

Navier--Stokes V26 is unchanged.  Its first rank-one Oseen-kernel
derivative norm equals the unrestricted symmetric-traceless norm,
while its second-derivative refinement yields only a localized
gain.  The transferable lesson remains that one should contract the
prescribed structured source before taking a larger ambient norm.

\subsection{What transfers to the SM--Einstein programme}

The V49 clock problem contains the same structural choice.  Let
\(\mathfrak K_n\) be the finite Hilbert space of source-table
superoperators at cutoff \(n\), with inner product
\[
 \langle A,B\rangle_{\mathfrak K_n}
 :=\sum_{\alpha}
 \operatorname{Tr}\!\left(A(F_\alpha)^*B(F_\alpha)\right),      \tag{50.3}
\]
where \((F_\alpha)\) is the fixed source-complete matrix-unit bank.
For a marked Hamiltonian direction \(X_n\) and a physical timing
direction \(Y_n\), the correct two-column card is the complete
mixed Gram matrix
\[
 \Gamma_n=
 \begin{pmatrix}
  \langle X_n,X_n\rangle&
  \langle X_n,Y_n\rangle\\
  \langle Y_n,X_n\rangle&
  \langle Y_n,Y_n\rangle
 \end{pmatrix}.                                                \tag{50.4}
\]
Taking separate bounds for the three entries before forming the
Schur complement can destroy precisely the cancellation which
decides whether the two directions are collinear.  Thus the
Yang--Mills trace-before-norm rule transfers directly: compute
\(\Gamma_n\), its real constrained Schur complement, and only then
bound perturbations.

\begin{proposition}[V50 audit decision]
\label{prop:smxiii-v50-decision}
No numerical constant and no continuum theorem is imported from
either companion programme.  The next SM--Einstein note shall:
\begin{enumerate}[label=\textup{(\roman*)}]
\item separate the symmetric Markov clock direction from the
antisymmetric Hamiltonian operation direction;
\item prove that a symmetric clock generator cannot calibrate a
nonzero Hamiltonian commutator through a positive real scale;
\item formulate the exact complete mixed-Gram criterion for a
separately observed antisymmetric physical operation germ; and
\item retain the mixed trace through the scalar minimization before
converting its residual into a completely bounded estimate.
\end{enumerate}
\end{proposition}

\begin{proof}
Equation (50.2) and the NS rank-one calculation are estimates in
different models, so neither supplies an SM coefficient.  Their
common algebraic lesson is order of operations: preserve the
structured contraction until the exact scalar or matrix quantity
has been formed.  In V49 the scalar calibration is determined by
the mixed entry of (50.4).  Hence the four stated tasks use the
transferable architecture without importing an unrelated bound.
\end{proof}

\begin{assessment}[Position after V50]
\label{ass:smxiii-v50}
The required cross-program audit is complete.  The finite
source-complete reconstruction of V49 remains valid, but it does
not supply a physical clock.  The next obstruction can now be
stated sharply: the reflected Markov clock and the Dirac--Yukawa
commutator have different adjoint parity, so their apparent timing
parameters cannot be identified without an additional
same-history antisymmetric operation measurement.
\end{assessment}

\begin{firewall}[Claim boundary after V50]
\label{fire:smxiii-v50}
This audit proves no continuum limit, physical clock calibration,
Standard Model parameter identification, Einstein coupling,
spectral gap, clustering theorem, or full SM--Einstein continuum.
The companion estimates are structural guides only.
\end{firewall}

\subsection{Parent record}

This V50 note appends the preceding material to the validated V49
source, whose record was
\[
\begin{array}{c|c}
 \text{lines}&16391\\
 \text{bytes}&577072\\
 \text{SHA--256}&
 \texttt{0A93ECDD7F5FE8DEC6296797CA5FA2ADBB6EAA154E4B8EBADFF9A8FE0C25A00E}.
\end{array}                                                     \tag{50.5}
\]

% =============================================================================
% END APPENDED THIRTEENTH-CYCLE V50 MATERIAL
% =============================================================================
% =============================================================================
% BEGIN APPENDED THIRTEENTH-CYCLE V51 MATERIAL
% =============================================================================

\section{Thirteenth-cycle V51: adjoint-parity obstruction and
identifiable physical clock calibration}
\label{sec:v13v51}

\subsection{The finite source metric and its adjoint}

Fix a cutoff \(n\), write \(\mathcal A_n=M_{d_n}(\mathbb C)\), and
give it the Hilbert--Schmidt inner product
\[
 \langle A,B\rangle_2=\operatorname{Tr}(A^*B).                 \tag{51.1}
\]
For a linear map \(T:\mathcal A_n\to\mathcal A_n\), let \(T^\dagger\)
denote its adjoint for (51.1).  On the Hilbert space
\(\mathfrak K_n=\mathcal L(\mathcal A_n,\mathcal A_n)\), use
\[
 \langle S,T\rangle_{\rm sup}
 =\sum_{i,j=1}^{d_n}
   \operatorname{Tr}\!\left(S(E_{ij})^*T(E_{ij})\right)
 =\operatorname{Tr}_{\rm sup}(S^\dagger T).                   \tag{51.2}
\]
Thus a source-complete matrix-unit table computes the
superoperator Hilbert--Schmidt inner product exactly.  This is the
complete mixed trace retained by the V50 decision.

For \(D=D^*\in\mathcal A_n\), put
\[
 X_D(A)=-i[D,A].                                               \tag{51.3}
\]
The unitary group \(A\mapsto e^{-itD}Ae^{itD}\) is unitary on
\((\mathcal A_n,\langle\cdot,\cdot\rangle_2)\), so
\[
 X_D^\dagger=-X_D.                                             \tag{51.4}
\]
A reversible or reflection-symmetric Markov generator in the same
finite source metric satisfies instead
\[
 L^\dagger=L.                                                  \tag{51.5}
\]
The equality (51.5) is a hypothesis on the common physical metric;
it must be checked after the OS quotient and is not inferred merely
from complete positivity.

\begin{theorem}[Adjoint-parity clock obstruction]
\label{thm:v13v51-parity}
If \(X^\dagger=-X\) and \(L^\dagger=L\) on a finite-dimensional
complex Hilbert space, then
\[
 \operatorname{Re}\langle X,L\rangle_{\rm sup}=0.              \tag{51.6}
\]
Consequently, if \(X\ne0\),
\[
 \|L-aX\|_{\rm sup}^2
 =\|L\|_{\rm sup}^2+a^2\|X\|_{\rm sup}^2
 \qquad(a\in\mathbb R),                                       \tag{51.7}
\]
and the unique least-squares real calibration is \(a=0\).  In
particular, no positive rescaling of a nonzero Hamiltonian
commutator can be calibrated from a self-adjoint Markov clock
generator.

\end{theorem}

\begin{proof}
Cyclicity of the finite-dimensional trace gives
\[
 \overline{\operatorname{Tr}_{\rm sup}(X^\dagger L)}
 =\operatorname{Tr}_{\rm sup}(L^\dagger X)
 =\operatorname{Tr}_{\rm sup}(LX)
 =\operatorname{Tr}_{\rm sup}(XL)
 =-\operatorname{Tr}_{\rm sup}(X^\dagger L).                  \tag{51.8}
\]
Hence the mixed trace is purely imaginary and (51.6) follows.
Expanding the square and using (51.6) gives (51.7).  Strict
convexity follows from \(\|X\|_{\rm sup}>0\).
\end{proof}

\begin{corollary}[The Poisson or OS Markov rate is not the missing
fermion clock]
\label{cor:v13v51-os-no-clock}
Suppose the reflected transfer generator is self-adjoint in the
source metric used to reconstruct \(X_{D_n}\).  Its Poisson
duration or OS Markov time cannot determine a positive physical
coefficient multiplying \(X_{D_n}\).  A separate antisymmetric
operation germ is necessary.
\end{corollary}

\begin{proof}
Apply Theorem~\ref{thm:v13v51-parity}.  Changing the positive
Markov rate rescales \(L\) but preserves (51.6).
\end{proof}

\subsection{The acquisition datum which removes the parity
obstruction}

Let \(\Psi_{n,\tau}^{+}\) and \(\Psi_{n,\tau}^{-}\) be two
differentiable same-history physical channel curves for
\(\tau\geq0\), with
\[
 \Psi_{n,0}^{+}=\Psi_{n,0}^{-}=I,\qquad
 L_n^\pm=\left.\frac{d}{d\tau}\Psi_{n,\tau}^{\pm}\right|_{\tau=0}.
                                                                    \tag{51.9}
\]
Their odd operation germ is
\[
 Z_n:=\frac12(L_n^+-L_n^-).                                   \tag{51.10}
\]
The required physical design condition is
\[
 Z_n^\dagger=-Z_n.                                             \tag{51.11}
\]
For example, (51.11) holds when the two channel germs have a common
self-adjoint dissipative part and opposite Hamiltonian parts.
Neither individual one-sided channel curve must extend to a
completely positive group.

Put
\[
 g_n=\|X_{D_n}\|_{\rm sup}^2,\qquad
 c_n=\langle X_{D_n},Z_n\rangle_{\rm sup},\qquad
 h_n=\|Z_n\|_{\rm sup}^2.                                     \tag{51.12}
\]
If both maps are skew-adjoint, the same trace calculation with two
minus signs shows that \(c_n\in\mathbb R\).  The complete timing
Gram is therefore
\[
 \Gamma_n=
 \begin{pmatrix}g_n&c_n\\c_n&h_n\end{pmatrix}\succeq0.          \tag{51.13}
\]

\begin{theorem}[Exact positive clock calibration from the mixed
timing Gram]
\label{thm:v13v51-gram-clock}
Assume \(X_{D_n}\ne0\) and (51.11).  Among all \(a\geq0\), the
unique minimizer of
\[
 q_n(a)=\|Z_n-aX_{D_n}\|_{\rm sup}^2                           \tag{51.14}
\]
is
\[
 a_n^+=\frac{\max\{c_n,0\}}{g_n},                              \tag{51.15}
\]
and the exact residual is
\[
 \varepsilon_n^2
 =h_n-\frac{\max\{c_n,0\}^2}{g_n}.                             \tag{51.16}
\]
There is an exact strictly positive physical calibration
\(Z_n=a_nX_{D_n}\), \(a_n>0\), if and only if
\[
 c_n>0,\qquad g_nh_n-c_n^2=0.                                 \tag{51.17}
\]
All quantities in (51.15)--(51.17) are entries or minors of the
complete two-column Gram card.

\end{theorem}

\begin{proof}
Expansion gives
\[
 q_n(a)=h_n-2ac_n+a^2g_n.                                     \tag{51.18}
\]
Since \(g_n>0\), the unconstrained minimizer is \(c_n/g_n\);
projection onto \([0,\infty)\) proves (51.15), and substitution
proves (51.16).  If \(c_n>0\), equality in (51.16) is equivalent
to \(g_nh_n-c_n^2=0\).  Equality in the Cauchy--Schwarz inequality
then gives \(Z_n=(c_n/g_n)X_{D_n}\).  Conversely such a positive
collinearity gives (51.17).
\end{proof}

\begin{corollary}[Source-complete residual upgrades to a cb
residual]
\label{cor:v13v51-cb-clock}
Let \(d=d_n\), let \(R_n=Z_n-a_n^+X_{D_n}\), and let
\(\varepsilon_n\) be (51.16).  Then
\[
 \|R_n\|_{\rm cb}\leq d\,\varepsilon_n.                        \tag{51.19}
\]
In particular, zero Gram residual proves equality of the two
linear maps on the whole cutoff algebra, not only on the recorded
sources.

\end{corollary}

\begin{proof}
The matrix units are an orthonormal Hilbert--Schmidt basis, so
\(\varepsilon_n\) is the Hilbert--Schmidt norm of the
superoperator \(R_n\).  Hence
\[
 \|(\operatorname{id}_{M_d}\otimes R_n)(A)\|_2
 \leq\varepsilon_n\|A\|_2
 \leq d\,\varepsilon_n\|A\|_{\rm op}                           \tag{51.20}
\]
for \(A\in M_d\otimes M_d\).  The operator norm is bounded by the
Hilbert--Schmidt norm.  Smith's finite-dimensional amplification
lemma says that the cb norm of a map with domain \(M_d\) is attained
at the \(d\)-fold amplification.  This proves (51.19).  If
\(\varepsilon_n=0\), every \(R_n(E_{ij})\) vanishes, and linearity
also proves the last statement directly.
\end{proof}

\subsection{Why the paired V43 construction does not itself set
physical units}

The paired unitary channels constructed in V43 have the form
\[
 \Phi^\pm_\tau=\operatorname{Ad}(e^{\mp i\tau D_n}).           \tag{51.21}
\]
Their odd germ is a fixed sign multiple of \(X_{D_n}\), so
Theorem~\ref{thm:v13v51-gram-clock} returns a zero residual.
However, replacing the curve parameter by \(b\tau\), \(b>0\),
replaces its odd germ by \(bZ_n\).

\begin{proposition}[Reparametrization non-identifiability]
\label{prop:v13v51-reparam}
Under \(\tau\mapsto b\tau\), the Gram entries transform as
\[
 (g_n,c_n,h_n)\longmapsto(g_n,bc_n,b^2h_n),                   \tag{51.22}
\]
and the calibrated coefficient transforms as
\[
 a_n^+\longmapsto b\,a_n^+.                                  \tag{51.23}
\]
Therefore a channel curve specified only up to positive
reparametrization cannot determine an absolute physical clock
coefficient, even when its Gram residual is exactly zero.

\end{proposition}

\begin{proof}
The chain rule gives \(Z_n\mapsto bZ_n\).  Equations
(51.12), (51.15), and (51.16) then give (51.22), (51.23), and
\(\varepsilon_n\mapsto b\varepsilon_n\).  Thus exact collinearity
is invariant but its scale is not.
\end{proof}

\begin{definition}[Clock-anchored same-history timing card]
\label{def:v13v51-clock-card}
A cutoff timing card is called clock-anchored if it contains:
\begin{enumerate}[label=\textup{(\roman*)}]
\item a physically specified duration observable \(\tau\), with
its unit and normalization fixed independently of \(D_n\);
\item the two same-history channel derivatives \(L_n^\pm\) with
respect to that \(\tau\);
\item source-complete tables for \(X_{D_n}\) and
\(Z_n=(L_n^+-L_n^-)/2\);
\item the adjoint-parity residual
\(\|Z_n+Z_n^\dagger\|_{\rm sup}\); and
\item the complete Gram card (51.13), evaluated before any
entrywise majorization.
\end{enumerate}
\end{definition}

\begin{decision}[Next acquisition after V51]
\label{dec:v13v51-next}
The first missing row is now exact: populate
Definition~\ref{def:v13v51-clock-card} from a duration observable
which is independent of the inserted Dirac--Yukawa mark.  After
that, prove adjacent-cutoff stability of (51.15)--(51.16) under
transported Gram errors and a uniform lower bound for \(g_n\).
Without the independent duration anchor, further manipulation of
the V43 paired channels would be circular.
\end{decision}

\begin{assessment}[Progress after V51]
\label{ass:v13v51}
The clock bottleneck has split into a proved impossibility and a
precise positive criterion.  The symmetric reflected Markov clock
cannot calibrate the Hamiltonian commutator.  A clock-anchored odd
same-history operation germ can do so, and its exact calibration,
orientation, residual, and cb error are computable from one
source-complete mixed Gram card.
\end{assessment}

\begin{firewall}[Claim boundary after V51]
\label{fire:v13v51}
No clock-anchored physical timing card has yet been populated.
Theorems~\ref{thm:v13v51-parity} and
\ref{thm:v13v51-gram-clock} do not identify an experimental clock,
a fermion mass, a Yukawa coupling, a continuum generator, or an
Einstein coupling.  They show exactly what data would identify the
finite-cutoff coefficient and why the currently available
symmetric clock does not.
\end{firewall}

\subsection{Parent record}

This V51 note appends the preceding material to the validated V50
source, whose record was
\[
\begin{array}{c|c}
 \text{lines}&16531\\
 \text{bytes}&582973\\
 \text{SHA--256}&
 \texttt{50957169750EE5829891559178BE51B5B286381AD725C7F00C9BE7911BABB237}.
\end{array}                                                     \tag{51.24}
\]

% =============================================================================
% END APPENDED THIRTEENTH-CYCLE V51 MATERIAL
% =============================================================================
% =============================================================================
% BEGIN APPENDED THIRTEENTH-CYCLE V52 MATERIAL
% =============================================================================

\section{Thirteenth-cycle V52: adjacent-cutoff stability of the
mixed-Gram clock}
\label{sec:v13v52}

\subsection{Transported timing directions}

For each cutoff \(n\), let \(\mathfrak K_n\) be the
source-superoperator Hilbert space of V51 and let
\[
 U_n:\mathfrak K_n\longrightarrow\mathfrak K_{n+1}             \tag{52.1}
\]
be an isometry induced by the chosen adjacent source-table
transport.  Write
\[
 X_n=X_{D_n},\qquad Z_n=\frac12(L_n^+-L_n^-),                  \tag{52.2}
\]
and define the transport defects
\[
 \alpha_n=\|X_{n+1}-U_nX_n\|_{\rm sup},\qquad
 \beta_n=\|Z_{n+1}-U_nZ_n\|_{\rm sup}.                         \tag{52.3}
\]
The clock at the two cutoffs is assumed to use the same transported
physical duration unit.  A change of clock unit is treated
separately below.

Set
\[
 g_n=\|X_n\|_{\rm sup}^2,\quad
 c_n=\operatorname{Re}\langle X_n,Z_n\rangle_{\rm sup},\quad
 h_n=\|Z_n\|_{\rm sup}^2.                                     \tag{52.4}
\]
This definition remains meaningful if a small adjoint-parity error
makes the mixed entry not exactly real.

\begin{lemma}[Adjacent Gram-entry perturbation]
\label{lem:v13v52-gram-pert}
The entries (52.4) satisfy
\[
\begin{split}
 |g_{n+1}-g_n|
 &\leq\alpha_n\bigl(\|X_{n+1}\|_{\rm sup}+\|X_n\|_{\rm sup}\bigr),\\
 |h_{n+1}-h_n|
 &\leq\beta_n\bigl(\|Z_{n+1}\|_{\rm sup}+\|Z_n\|_{\rm sup}\bigr),\\
 |c_{n+1}-c_n|
 &\leq\alpha_n\|Z_{n+1}\|_{\rm sup}
       +\beta_n\|X_n\|_{\rm sup}.                             \tag{52.5}
\end{split}
\]

\end{lemma}

\begin{proof}
Because \(U_n\) is an isometry,
\[
 \bigl|\|X_{n+1}\|^2-\|U_nX_n\|^2\bigr|
 \leq\|X_{n+1}-U_nX_n\|
      \bigl(\|X_{n+1}\|+\|U_nX_n\|\bigr),                     \tag{52.6}
\]
which is the first estimate; the second is identical.  For the
mixed entry, insert and subtract
\(\langle U_nX_n,Z_{n+1}\rangle\):
\[
\begin{split}
 |\langle X_{n+1},Z_{n+1}\rangle
   -\langle U_nX_n,U_nZ_n\rangle|
 &\leq\alpha_n\|Z_{n+1}\|
       +\|U_nX_n\|\beta_n.                                    \tag{52.7}
\end{split}
\]
Take real parts and use the isometry.
\end{proof}

\subsection{Stability away from rank and orientation loss}

The positive calibration branch is
\[
 a_n=\frac{c_n}{g_n},\qquad
 r_n=h_n-\frac{c_n^2}{g_n},                                  \tag{52.8}
\]
provided \(g_n>0\) and \(c_n>0\).  Here \(r_n\) is the squared
least-squares residual.

\begin{theorem}[Quantitative adjacent stability]
\label{thm:v13v52-stability}
Suppose for two adjacent cutoffs that
\[
 g_n,g_{n+1}\geq\gamma>0,\qquad
 0<c_n,c_{n+1}\leq C,\qquad
 |c_{n+1}-c_n|\leq\delta_c,\quad
 |g_{n+1}-g_n|\leq\delta_g,\quad
 |h_{n+1}-h_n|\leq\delta_h.                                  \tag{52.9}
\]
Then
\[
 |a_{n+1}-a_n|
 \leq \frac{\delta_c}{\gamma}
       +\frac{C\,\delta_g}{\gamma^2},                          \tag{52.10}
\]
and
\[
 |r_{n+1}-r_n|
 \leq\delta_h+\frac{2C\,\delta_c}{\gamma}
                  +\frac{C^2\delta_g}{\gamma^2}.               \tag{52.11}
\]
Consequently, with \(\varepsilon_k=\sqrt{r_k}\),
\[
 |\varepsilon_{n+1}-\varepsilon_n|
 \leq
 \left(
 \delta_h+\frac{2C\,\delta_c}{\gamma}
          +\frac{C^2\delta_g}{\gamma^2}
 \right)^{1/2}.                                                \tag{52.12}
\]

\end{theorem}

\begin{proof}
The quotient identity gives
\[
 \left|\frac{c_{n+1}}{g_{n+1}}-\frac{c_n}{g_n}\right|
 \leq
 \frac{|c_{n+1}-c_n|}{g_{n+1}}
 +\frac{|c_n|\,|g_{n+1}-g_n|}{g_ng_{n+1}},                    \tag{52.13}
\]
which proves (52.10).  Similarly,
\[
\begin{split}
 \left|\frac{c_{n+1}^2}{g_{n+1}}-\frac{c_n^2}{g_n}\right|
 &\leq
 \frac{|c_{n+1}-c_n|\,|c_{n+1}+c_n|}{g_{n+1}}\\
 &\quad+
 \frac{c_n^2|g_{n+1}-g_n|}{g_ng_{n+1}},                       \tag{52.14}
\end{split}
\]
and (52.11) follows.  Both residuals are nonnegative by
Cauchy--Schwarz.  For \(x,y\geq0\),
\(|\sqrt{x}-\sqrt{y}|\leq\sqrt{|x-y|}\), which proves (52.12).
\end{proof}

\begin{corollary}[A populated adjacent card]
\label{cor:v13v52-card}
If
\[
 \|X_n\|,\|X_{n+1}\|\leq B_X,\qquad
 \|Z_n\|,\|Z_{n+1}\|\leq B_Z,                                \tag{52.15}
\]
then Theorem~\ref{thm:v13v52-stability} applies with
\[
 \delta_g=2B_X\alpha_n,\qquad
 \delta_h=2B_Z\beta_n,\qquad
 \delta_c=B_Z\alpha_n+B_X\beta_n.                             \tag{52.16}
\]
Thus the entire clock-stability estimate is populated by the two
transport defects, two norm ceilings, one Gram floor, and one
mixed-entry ceiling.

\end{corollary}

\begin{proof}
Substitute Lemma~\ref{lem:v13v52-gram-pert} into
Theorem~\ref{thm:v13v52-stability}.
\end{proof}

The positive branch itself must also be protected.  If
\(c_n\geq\kappa>0\) and
\[
 B_Z\alpha_n+B_X\beta_n<\kappa,                               \tag{52.17}
\]
then (52.5) gives \(c_{n+1}>0\).  Thus loss of orientation is an
explicit card failure and cannot be concealed in an absolute-value
estimate.

\subsection{Cofinal convergence and the cb scaling gate}

\begin{theorem}[Cofinal clock coefficient in the Hilbert direct
limit]
\label{thm:v13v52-cofinal}
Assume coherent isometries \(U_n\), uniform bounds (52.15), and
\[
 \sum_{n=1}^{\infty}(\alpha_n+\beta_n)<\infty.                 \tag{52.18}
\]
Then \(X_n\) and \(Z_n\), embedded in the Hilbert direct limit,
converge to elements \(X_\infty\) and \(Z_\infty\).  If in addition
\[
 g_n\geq\gamma>0,\qquad c_n\geq\kappa>0,                      \tag{52.19}
\]
then \(a_n=c_n/g_n\) and \(r_n=h_n-c_n^2/g_n\) converge, with
\[
 a_\infty=
 \frac{\operatorname{Re}\langle X_\infty,Z_\infty\rangle}
      {\|X_\infty\|^2},\qquad
 r_\infty=
 \|Z_\infty-a_\infty X_\infty\|^2.                            \tag{52.20}
\]
If \(r_n\to0\), then
\[
 Z_\infty=a_\infty X_\infty,\qquad a_\infty>0.                \tag{52.21}
\]

\end{theorem}

\begin{proof}
For \(m>n\), telescope the transported differences in the common
direct-limit Hilbert space:
\[
 \|X_m-U_{m-1}\cdots U_nX_n\|
 \leq\sum_{k=n}^{m-1}\alpha_k,                                \tag{52.22}
\]
and similarly for \(Z\).  Equation (52.18) makes both nets Cauchy.
Continuity of norms and inner products gives convergence of
\(g_n,c_n,h_n\).  The floor in (52.19) makes division continuous,
which proves (52.20).  The orientation floor gives
\(a_\infty\geq\kappa/\sup_ng_n>0\).  Finally \(r_\infty=0\)
is precisely (52.21).
\end{proof}

The Hilbert direct-limit statement is weaker than the continuum
channel statement.  At cutoff \(n\), V51 gives
\[
 \|Z_n-a_nX_n\|_{\rm cb}\leq d_n\sqrt{r_n}.                   \tag{52.23}
\]
Hence the sufficient source-complete continuum requirement is
\[
 d_n\sqrt{r_n}\longrightarrow0,                               \tag{52.24}
\]
not merely \(r_n\to0\).  Growth of the matrix algebra can defeat a
vanishing unweighted Gram residual.

\begin{proposition}[Local continuum consequence of the cb gate]
\label{prop:v13v52-cb-local}
Let \((\mathcal A_n,\iota_n)\) be a unital completely isometric
inductive system.  Suppose \(a_n\to a_\infty\), and suppose that
for every \(A\in\mathcal A_N\) the transported values
\[
 X_n(\iota_{N,n}A)\longrightarrow X_\infty(A),\qquad
 Z_n(\iota_{N,n}A)\longrightarrow Z_\infty(A)                 \tag{52.25}
\]
exist in a common normed target.  If (52.24) holds, then
\[
 Z_\infty(A)=a_\infty X_\infty(A)                             \tag{52.26}
\]
for every local \(A\).

\end{proposition}

\begin{proof}
Put \(A_n=\iota_{N,n}A\).  Complete isometry gives
\(\|A_n\|=\|A\|\), while (52.23) gives
\[
 \|Z_n(A_n)-a_nX_n(A_n)\|
 \leq d_n\sqrt{r_n}\,\|A\|.                                  \tag{52.27}
\]
The right side tends to zero.  Passing to the three limits proves
(52.26).
\end{proof}
\subsection{Clock-unit cocycle}

Let the duration coordinate at cutoff \(n+1\) be
\(\tau_{n+1}=b_n\tau_n\), \(b_n>0\).  Differentiation then changes
the transported odd germ by a reciprocal factor, depending on
which coordinate is held fixed.  Fix the convention
\[
 Z_{n+1}\simeq b_n^{-1}U_nZ_n.                                \tag{52.28}
\]
The raw coefficients satisfy
\[
 a_{n+1}\simeq b_n^{-1}a_n.                                  \tag{52.29}
\]
Define the accumulated clock transport
\[
 B_1=1,\qquad B_n=\prod_{k=1}^{n-1}b_k,\qquad
 \widehat a_n=B_na_n.                                         \tag{52.30}
\]

\begin{proposition}[Clock-cocycle removal]
\label{prop:v13v52-clock-cocycle}
If the defects in (52.28) are summable after multiplication by
\(B_{n+1}\), if
\[
 \sum_n|\log b_n|<\infty,                                     \tag{52.31}
\]
and the hypotheses of Theorem~\ref{thm:v13v52-cofinal} hold for
the rescaled germs \(\widehat Z_n=B_nZ_n\), then
\(\widehat a_n\) converges to a positive finite coefficient.
Without a specified cocycle \((b_n)\), the raw sequence \((a_n)\)
has no invariant cross-cutoff meaning.

\end{proposition}

\begin{proof}
Condition (52.31) makes \(B_n\) converge to a finite positive
limit.  The rescaled transport relation has a summable additive
defect by hypothesis, so Theorem~\ref{thm:v13v52-cofinal} applies
to \((X_n,\widehat Z_n)\).  Its calibrated coefficient is
\(B_na_n\) by the homogeneity proved in
Proposition~\ref{prop:v13v51-reparam}.  The last assertion follows
from the freedom to replace \(b_n\) by any positive clock-unit
change.
\end{proof}

\begin{decision}[Acquisition card after V52]
\label{dec:v13v52-next}
A viable cofinal timing packet must now record, with rational
ceilings where appropriate,
\[
 (d_n,\alpha_n,\beta_n,B_X,B_Z,\gamma,\kappa,r_n,b_n).         \tag{52.32}
\]
The immediate upstream problem is to construct the physical
duration observable and its adjacent cocycle \(b_n\) from the
same-history regulator, independently of the Dirac--Yukawa mark.
Only after that construction is it meaningful to demand the
summability tests (52.18), (52.24), and (52.31).
\end{decision}

\begin{assessment}[Progress after V52]
\label{ass:v13v52}
Finite clock identification is now connected to a rigorous
cofinal criterion.  The coefficient is stable away from
source-rank and orientation loss; summable transported germ
defects produce a Hilbert direct-limit coefficient; and the
stronger dimension-weighted residual controls the cb mismatch.
The clock-unit cocycle has been isolated as independent data rather
than absorbed into the fermion coefficient.
\end{assessment}

\begin{firewall}[Claim boundary after V52]
\label{fire:v13v52}
The bounds above are conditional because no physical timing card,
Gram floor, orientation floor, clock cocycle, or
dimension-weighted residual has been populated.  Hilbert
direct-limit convergence alone does not construct a locally normal
continuum dynamics, an unbounded Dirac operator, the Standard Model
parameters, or the Einstein coupling.
\end{firewall}

\subsection{Parent record}

This V52 note appends the preceding material to the validated V51
source, whose record was
\[
\begin{array}{c|c}
 \text{lines}&16835\\
 \text{bytes}&593891\\
 \text{SHA--256}&
 \texttt{493CF85F86A6134D5E0C1B6E4B18AB852675A0BDDF5656D4B927A2DD2B24BC94}.
\end{array}                                                     \tag{52.33}
\]

% =============================================================================
% END APPENDED THIRTEENTH-CYCLE V52 MATERIAL
% =============================================================================
% =============================================================================
% BEGIN APPENDED THIRTEENTH-CYCLE V53 MATERIAL
% =============================================================================

\section{Thirteenth-cycle V53: pressure-duration anchoring and the
finite same-history timing card}
\label{sec:v13v53}

\subsection{Upstream result audit}

The Grand Tensor manuscript already contains two relevant finite
results which must not be repeated.

First, its pressure-adapted coherent ADM packet defines a completed
boundary duration by
\[
 \tau_h=h_{{\rm P},h}\mathcal N_h
       =c_{\rm time}\ell_{0,h},\qquad
 h_{{\rm P},h}=\frac{c_{\rm time}}{\beta_h},\qquad
 \mathcal N_h=\beta_h\ell_{0,h}.                              \tag{53.1}
\]
Here the notation \(c_{\rm time}\) is used in place of the
manuscript's overloaded \(c_*\): in the neighboring theorem the
same typography also denotes a source-frame floor, which has
different units and a different role.

Second, the manuscript's unique calibrated reversal-writer theorem
proves that, for a given physical half-step \(t>0\), reversal
oddness and
\[
 t(\pi_+-\pi_-)J=I                                            \tag{53.2}
\]
force
\[
 Jx=\frac1{2t}(x,-x).                                         \tag{53.3}
\]
V53 does not re-prove these facts.  It identifies the additional
incidence and metrological data needed to use (53.1) as the
physical \(t\) in (53.2), and then connects the resulting finite
difference to the mixed-Gram calibration of V51.

\subsection{Incidence and absolute normalization}

A completed boundary step and a centered forward/reverse panel
need not have the same duration convention.  Introduce an
incidence constant
\[
 \kappa_{\rm arm}>0,\qquad
 t_h=\kappa_{\rm arm}\tau_h.                                  \tag{53.4}
\]
For example, \(\kappa_{\rm arm}=1/2\) when the two endpoints of the
centered panel are separated by one completed boundary duration,
and \(\kappa_{\rm arm}=1\) when each arm itself is one completed
boundary step.  The value is a property of the represented
history incidence and cannot be selected after the channel has
been evaluated.

The conversion \(c_{\rm time}\) is fixed metrologically by one
reference duration:
\[
 \ell_{0,h_0}>0,\qquad
 \tau_{\rm ref}>0,\qquad
 c_{\rm time}:=\frac{\tau_{\rm ref}}{\ell_{0,h_0}}.            \tag{53.5}
\]
Without (53.5), (53.1) determines only duration ratios.

\begin{theorem}[Pressure-depth clock anchor]
\label{thm:v13v53-pressure-clock}
Assume:
\begin{enumerate}[label=\textup{(\roman*)}]
\item \(\beta_h>0\) and \(\ell_{0,h}>0\) are reconstructed from a
protected pressure/temporal marginal;
\item \(c_{\rm time}\) is fixed by (53.5);
\item \(\kappa_{\rm arm}\) is fixed by a same-history incidence
record; and
\item insertion of the Dirac--Yukawa operation mark is passive for
the marginal which determines \((\beta_h,\ell_{0,h})\).
\end{enumerate}
Then
\[
 t_h=\kappa_{\rm arm}c_{\rm time}\ell_{0,h}
    =\kappa_{\rm arm}h_{{\rm P},h}\mathcal N_h                \tag{53.6}
\]
is a positive physical half-step independent of the inserted
operator \(D_h\).  It is therefore an admissible duration anchor
for the reversal writer (53.3) and the timing card of V51.

\end{theorem}

\begin{proof}
Positivity follows from the first three hypotheses.  The
identities in (53.6) are (53.1) and (53.4).  The quantities on
their right side are reconstructed from the protected marginal
and the independent reference calibration.  Passive marking says
that this marginal is unchanged when the operation mark is
inserted.  Hence \(t_h\) is not defined through \(D_h\), which
removes the reparametrization circularity of
Proposition~\ref{prop:v13v51-reparam}.  The writer conclusion then
uses the already proved Grand Tensor result (53.2)--(53.3).
\end{proof}

The passivity hypothesis has a direct same-history test.  If
\(P_h\) is the joint law of the pressure-depth record and the
operation mark and \(P_h^{\rm clk}\) is its clock marginal, require
\[
 P_h^{\rm clk}=P_{h,0}^{\rm clk},                             \tag{53.7}
\]
where the right side is the unmarked clock law.  In a finite
cylinder, (53.7) is equality of finitely many probabilities, not
an independence assumption about unrecorded variables.

\subsection{Mesh ratio is not a unit cocycle}

From (53.6),
\[
 \frac{t_{h'}}{t_h}
 =\frac{\ell_{0,h'}}{\ell_{0,h}}                              \tag{53.8}
\]
when \(c_{\rm time}\) and \(\kappa_{\rm arm}\) are common.  This
ratio measures step refinement in one fixed physical unit.

\begin{proposition}[Separation of mesh scale from clock-unit
gauge]
\label{prop:v13v53-mesh-unit}
Suppose all \(t_h\) in (53.6) are expressed using the same
\(c_{\rm time}\).  Even if \(t_h\downarrow0\), the unit cocycle
\(b_n\) of Proposition~\ref{prop:v13v52-clock-cocycle} is
\[
 b_n=1.                                                       \tag{53.9}
\]
The ratio \(t_{n+1}/t_n\) in (53.8) must not be substituted for
\(b_n\).  Doing so would confuse a smaller physical step with a
change in the coordinate unit and would incorrectly renormalize
the recovered fermion coefficient.

\end{proposition}

\begin{proof}
A unit cocycle compares the numerical representations of the same
physical duration after a change of coordinate unit.  Equation
(53.5) fixes one common conversion from depth to that unit at all
cutoffs.  Therefore the coordinate conversion is the identity,
which is (53.9).  By contrast, (53.8) compares two different
physical durations in the same unit.  The two ratios have
different definitions even when their numerical values happen to
coincide.
\end{proof}

Thus the finite-positive-limit hypothesis (52.31) concerns
optional unit changes only.  It is not a condition that the
continuum mesh duration stay positive.  On the anchored branch it
is discharged identically by \(b_n=1\).

\subsection{Finite channels rather than formal derivatives}

Let \(\Psi_{h,s}^{\pm}\) be the clock-anchored forward and reverse
physical channels through the identity.  Suppose, for
\(0\leq s\leq t_h\),
\[
 \Psi_{h,s}^{\pm}
 =I+sL_h^\pm+R_{h,s}^{\pm},\qquad
 \|R_{h,s}^{\pm}\|_{\rm cb}\leq K_hs^2.                       \tag{53.10}
\]
The observable odd finite-difference germ is
\[
 \widetilde Z_h
 :=\frac{\Psi_{h,t_h}^{+}-\Psi_{h,t_h}^{-}}{2t_h},\qquad
 Z_h:=\frac12(L_h^+-L_h^-).                                  \tag{53.11}
\]

\begin{theorem}[Clock-anchored finite-difference error]
\label{thm:v13v53-finite-difference}
Under (53.10),
\[
 \|\widetilde Z_h-Z_h\|_{\rm cb}\leq K_ht_h.                  \tag{53.12}
\]
Let \(X_h=X_{D_h}\), \(g_h=\|X_h\|_{\rm sup}^2\geq\gamma>0\),
and assume \(\|X_h\|_{\rm sup}\leq B_X\).  Put
\[
 c_h=\operatorname{Re}\langle X_h,Z_h\rangle,\qquad
 \widetilde c_h=
 \operatorname{Re}\langle X_h,\widetilde Z_h\rangle.          \tag{53.13}
\]
Because the source-superoperator norm is bounded by the cb norm
times a finite source-bank constant \(C_{{\rm src},h}\),
\[
 |\widetilde c_h-c_h|
 \leq B_XC_{{\rm src},h}K_ht_h.                              \tag{53.14}
\]
On the positive branch, if both \(c_h\) and \(\widetilde c_h\)
are positive, their calibrated coefficients obey
\[
 |\widetilde a_h-a_h|
 \leq\frac{B_XC_{{\rm src},h}K_ht_h}{\gamma}.                 \tag{53.15}
\]

\end{theorem}

\begin{proof}
Subtract the two expansions in (53.10), divide by \(2t_h\), and
use
\[
 \frac{\|R_{h,t_h}^+\|_{\rm cb}
       +\|R_{h,t_h}^-\|_{\rm cb}}{2t_h}
 \leq K_ht_h,                                                 \tag{53.16}
\]
which proves (53.12).  Cauchy--Schwarz in the source-superoperator
Hilbert space and the definition of \(C_{{\rm src},h}\) prove
(53.14).  Since the \(X_h\) column is unchanged,
\(\widetilde a_h-a_h=(\widetilde c_h-c_h)/g_h\), proving
(53.15).
\end{proof}

For the matrix-unit bank of \(M_{d_h}\), one may take the explicit
but nonuniform ceiling
\[
 C_{{\rm src},h}\leq d_h^{3/2}.                                    \tag{53.17}
\]
Indeed, every \(E_{ij}\) has operator norm one and each output has
Hilbert--Schmidt norm at most \(d_h^{1/2}\|T\|_{\rm cb}\), so the
sum of the \(d_h^2\) squared output norms is at most
\(d_h^3\|T\|_{\rm cb}^2\).  A sharper populated source constant
should be retained whenever the operation acts on a smaller
supported carrier.

\begin{corollary}[Orientation margin for a finite timing card]
\label{cor:v13v53-orientation}
If
\[
 c_h\geq\kappa_h>0,\qquad
 B_XC_{{\rm src},h}K_ht_h<\kappa_h,                           \tag{53.18}
\]
then \(\widetilde c_h>0\), so the finite channel card selects the
same positive calibration branch as the tangent card.
\end{corollary}

\begin{proof}
This is (53.14) and the triangle inequality.
\end{proof}

\subsection{The first genuinely populatable timing packet}

Combining the existing Grand Tensor writer theorem with
Theorems~\ref{thm:v13v53-pressure-clock} and
\ref{thm:v13v53-finite-difference}, the required finite card is
\[
 \boxed{
 \left(
 \beta_h,\ell_{0,h},c_{\rm time},\kappa_{\rm arm},
 t_h,\Psi_{h,t_h}^+,\Psi_{h,t_h}^-,
 K_h,X_h,\widetilde\Gamma_h
 \right),}                                                    \tag{53.19}
\]
where
\[
 \widetilde\Gamma_h=
 \begin{pmatrix}
  \langle X_h,X_h\rangle&
  \langle X_h,\widetilde Z_h\rangle\\
  \langle\widetilde Z_h,X_h\rangle&
  \langle\widetilde Z_h,\widetilde Z_h\rangle
 \end{pmatrix}.                                               \tag{53.20}
\]
The card must also include the finite probability equalities
(53.7) and the history incidence which fixes
\(\kappa_{\rm arm}\).

\begin{decision}[Direction after V53]
\label{dec:v13v53-next}
The deterministic dense-step branch now has an explicit,
noncircular clock bridge.  The next upstream ambiguity is that the
renewal instrument also carries a random protected duration.
Replacing that random duration by its mean can bias a generator.
The next note shall derive the exact moment condition under which a
random holding-time channel recovers the same physical tangent and
shall keep the duration/channel joint law on one history.
\end{decision}

\begin{assessment}[Progress after V53]
\label{ass:v13v53}
The pressure clock can anchor the SM operation germ provided one
reference duration, one arm-incidence convention, and passive-mark
marginal equality are supplied.  A shrinking mesh is compatible
with a fixed physical time unit.  Actual finite forward/reverse
channels now approximate the clock-anchored odd germ with an
explicit \(K_ht_h\) error, so the timing card need not assume
access to a formal derivative.
\end{assessment}

\begin{firewall}[Claim boundary after V53]
\label{fire:v13v53}
Neither the reference duration nor the arm-incidence record has
been populated from a concrete regulator card in these notes.
The manuscript formula (53.1) does not by itself fix
\(c_{\rm time}\).  No uniform bound on \(K_h\), no orientation
margin, no cb-vanishing finite-difference tail, and no continuum
fermion generator is yet proved.
\end{firewall}

\subsection{Parent record}

This V53 note appends the preceding material to the validated V52
source, whose record was
\[
\begin{array}{c|c}
 \text{lines}&17190\\
 \text{bytes}&605334\\
 \text{SHA--256}&
 \texttt{0883344E43779C52DBBECFDA74CC9E5BCFD008C1DB1DFE628E52A44C388EF4EF}.
\end{array}                                                     \tag{53.21}
\]

% =============================================================================
% END APPENDED THIRTEENTH-CYCLE V53 MATERIAL
% =============================================================================
% =============================================================================
% BEGIN APPENDED THIRTEENTH-CYCLE V54 MATERIAL
% =============================================================================

\section{Thirteenth-cycle V54: random holding times and the
duration-weighted odd operation germ}
\label{sec:v13v54}

\subsection{The joint holding-time channel}

Let \((\Omega_h,\mathbb P_h)\) be the finite protected
same-history cylinder at cutoff \(h\).  Let
\[
 T_h:\Omega_h\longrightarrow(0,\infty)                        \tag{54.1}
\]
be its physical holding-time writer, expressed in the common unit
fixed in V53.  For each history \(\omega\), suppose the two
physical operation channels have expansions
\[
 \Psi_{h,s}^{\pm}(\omega)
 =I+sL_h^\pm(\omega)+R_{h,s}^{\pm}(\omega),\qquad
 \|R_{h,s}^{\pm}(\omega)\|_{\rm cb}
 \leq K_hs^2                                                   \tag{54.2}
\]
for \(0\leq s\leq T_h(\omega)\).  Define
\[
 Z_h(\omega)=\frac12\bigl(L_h^+(\omega)-L_h^-(\omega)\bigr),   \tag{54.3}
\]
and set
\[
 m_h=\mathbb E_hT_h,\qquad
 v_{2,h}=\mathbb E_hT_h^2.                                   \tag{54.4}
\]
The observable odd channel quotient and the duration-weighted
tangent are
\[
 \overline Z_h
 =\frac{\mathbb E_h[
  \Psi_{h,T_h}^{+}-\Psi_{h,T_h}^{-}]}{2m_h},\qquad
 Z_h^\sharp
 =\frac{\mathbb E_h[T_hZ_h]}{m_h}.                            \tag{54.5}
\]
All expectations in (54.5) are taken from one joint
duration/channel law.

\begin{theorem}[Random-duration finite-difference theorem]
\label{thm:v13v54-random}
Under (54.1)--(54.4),
\[
 \boxed{\;
 \|\overline Z_h-Z_h^\sharp\|_{\rm cb}
 \leq K_h\frac{v_{2,h}}{m_h}.
 \;}                                                          \tag{54.6}
\]
Consequently, a sufficient holding-time condition for recovery of
the duration-weighted tangent is
\[
 K_h\frac{\mathbb E_hT_h^2}{\mathbb E_hT_h}\longrightarrow0.  \tag{54.7}
\]
If \(T_h\leq t_h^{\max}\) almost surely, then the left side of
(54.7) is at most \(K_ht_h^{\max}\).

\end{theorem}

\begin{proof}
Subtract the two expansions in (54.2), evaluate at
\(s=T_h(\omega)\), and average:
\[
 \overline Z_h-Z_h^\sharp
 =\frac{\mathbb E_h[
 R_{h,T_h}^+-R_{h,T_h}^-]}{2m_h}.                             \tag{54.8}
\]
The Bochner integral triangle inequality and (54.2) give
\[
 \|\overline Z_h-Z_h^\sharp\|_{\rm cb}
 \leq\frac{2K_h\mathbb E_hT_h^2}{2m_h},
\]
which is (54.6).  The last assertion uses
\(T_h^2\leq t_h^{\max}T_h\).
\end{proof}

Condition (54.7) is stronger than \(m_h\to0\).  Rare long holding
times can make \(m_h\) small while keeping \(v_{2,h}/m_h\) away
from zero.  The correct card therefore retains the second moment
or an almost-sure duration ceiling.

\subsection{Size bias and the covariance row}

The ordinary mean tangent is
\[
 Z_h^0=\mathbb E_hZ_h.                                        \tag{54.9}
\]
In general it differs from the tangent in (54.5):
\[
 Z_h^\sharp-Z_h^0
 =\frac{\mathbb E_h[(T_h-m_h)(Z_h-Z_h^0)]}{m_h}.              \tag{54.10}
\]

\begin{proposition}[Exact holding-time bias and its bound]
\label{prop:v13v54-bias}
Let
\[
 \sigma_{T,h}^2=\mathbb E_h(T_h-m_h)^2,\qquad
 \sigma_{Z,h}^2=
 \mathbb E_h\|Z_h-Z_h^0\|_{\rm cb}^2.                         \tag{54.11}
\]
Then
\[
 \|Z_h^\sharp-Z_h^0\|_{\rm cb}
 \leq\frac{\sigma_{T,h}\sigma_{Z,h}}{m_h}.                    \tag{54.12}
\]
If \(T_h\) and the complete map-valued writer \(Z_h\) are
independent in their same-history joint law, then the bias in
(54.10) is exactly zero.

\end{proposition}

\begin{proof}
Equation (54.10) follows by expanding \(T_h=m_h+(T_h-m_h)\) and
using \(\mathbb E_h(T_h-m_h)=0\).  Apply scalar Cauchy--Schwarz to
the norm of the finite Bochner sum to obtain (54.12).
Independence factors the mixed expectation into
\(\mathbb E_h(T_h-m_h)\mathbb E_h(Z_h-Z_h^0)=0\).
\end{proof}

The independence clause is sufficient but is not required.  The
sharp datum is the map-valued covariance in (54.10), which may
vanish through a symmetry visible only in the joint cylinder.

\begin{countertheorem}[Separate marginals do not determine the
physical odd tangent]
\label{cth:v13v54-marginals}
Let \(X\ne0\) be skew-adjoint.  At every \(h>0\), consider two
equiprobable histories with duration marginal
\(\{h,3h\}\) and operation-germ marginal \(\{X,-X\}\).
There are two couplings with these same marginals:
\[
\begin{array}{c|cc}
 &\omega_1&\omega_2\\ \hline
 {\rm A} &(h,X)&(3h,-X)\\
 {\rm B} &(h,-X)&(3h,X).
\end{array}                                                    \tag{54.13}
\]
They have
\[
 Z_{h,{\rm A}}^\sharp=-\frac12X,\qquad
 Z_{h,{\rm B}}^\sharp=\frac12X.                               \tag{54.14}
\]
Thus the positive orientation of the clock calibration can reverse
while both marginal laws remain unchanged.

\end{countertheorem}

\begin{proof}
In both couplings \(m_h=2h\).  For A,
\[
 \mathbb E[T_hZ_h]=\frac12(hX-3hX)=-hX,
\]
while for B the value is \(hX\).  Divide by \(m_h\).
\end{proof}

This counterexample also has
\[
 \frac{\mathbb E T_h^2}{\mathbb E T_h}
 =\frac{5}{2}h\longrightarrow0.                              \tag{54.15}
\]
Hence perfect finite-difference control does not repair loss of
the joint duration/operation incidence.

\subsection{Calibration error from the random finite card}

Let \(X_h=X_{D_h}\), assume
\[
 \|X_h\|_{\rm sup}\leq B_X,\qquad
 g_h=\|X_h\|_{\rm sup}^2\geq\gamma>0,                         \tag{54.16}
\]
and use the finite source-bank constant \(C_{{\rm src},h}\) of
V53.  Define the positive-branch coefficients from
\(Z_h^\sharp\) and \(\overline Z_h\) by
\[
 a_h^\sharp=
 \frac{\operatorname{Re}\langle X_h,Z_h^\sharp\rangle}{g_h},
 \qquad
 \overline a_h=
 \frac{\operatorname{Re}\langle X_h,\overline Z_h\rangle}{g_h}.
                                                                    \tag{54.17}
\]

\begin{corollary}[Random-card calibration bound]
\label{cor:v13v54-calibration}
If the two numerators in (54.17) are positive, then
\[
 |\overline a_h-a_h^\sharp|
 \leq
 \frac{B_XC_{{\rm src},h}K_h}{\gamma}
 \frac{v_{2,h}}{m_h}.                                        \tag{54.18}
\]
If the exact orientation margin satisfies
\[
 \operatorname{Re}\langle X_h,Z_h^\sharp\rangle
 >
 B_XC_{{\rm src},h}K_h\frac{v_{2,h}}{m_h},                   \tag{54.19}
\]
then the observed finite card lies on the same positive branch.

\end{corollary}

\begin{proof}
Theorem~\ref{thm:v13v54-random} and the definition of
\(C_{{\rm src},h}\) bound the source-superoperator difference by
\(C_{{\rm src},h}K_hv_{2,h}/m_h\).  Pair with \(X_h\), divide by
\(g_h\), and use (54.16).  The same bound and the triangle
inequality prove (54.19).
\end{proof}

\subsection{Map-valued conditional short}

A nuisance record \(Q_h\) may influence both duration and the
operation germ.  The correct physical target is then retained
conditionally:
\[
 Z_h^\sharp(q)
 =
 \frac{\mathbb E_h[T_hZ_h\mid Q_h=q]}
      {\mathbb E_h[T_h\mid Q_h=q]}.                           \tag{54.20}
\]
It may be averaged only after the conditional calibration has been
formed.  If a single coefficient \(a_h\) is claimed across all
nuisance cells, the complete test is
\[
 \sum_q p_h(q)
 \left\|Z_h^\sharp(q)-a_hX_h(q)\right\|_{\rm sup}^2,           \tag{54.21}
\]
with \(a_h\geq0\) minimized after summation.  The numerator and
denominator in (54.20), the cross terms in (54.21), and the
duration moments in (54.4) are all finite same-history cylinder
moments.

\begin{proposition}[Conditional common-scale compiler]
\label{prop:v13v54-conditional}
Put
\[
 g_h^{\rm all}=\sum_qp_h(q)\|X_h(q)\|_{\rm sup}^2,\quad
 c_h^{\rm all}=\sum_qp_h(q)
 \operatorname{Re}\langle X_h(q),Z_h^\sharp(q)\rangle,\quad
 h_h^{\rm all}=\sum_qp_h(q)\|Z_h^\sharp(q)\|_{\rm sup}^2.     \tag{54.22}
\]
If \(g_h^{\rm all}>0\), the optimal common positive coefficient is
\[
 a_h^{\rm all}
 =\frac{\max\{c_h^{\rm all},0\}}{g_h^{\rm all}},              \tag{54.23}
\]
with squared residual
\[
 r_h^{\rm all}
 =h_h^{\rm all}
 -\frac{\max\{c_h^{\rm all},0\}^2}{g_h^{\rm all}}.            \tag{54.24}
\]
It vanishes on the strictly positive branch exactly when one
\(a_h>0\) satisfies
\(Z_h^\sharp(q)=a_hX_h(q)\) for every positive-mass nuisance cell.

\end{proposition}

\begin{proof}
The direct sum over \(q\), weighted by \(p_h(q)^{1/2}\), is a
finite Hilbert space.  Apply
Theorem~\ref{thm:v13v51-gram-clock} in that direct sum.  A sum of
nonnegative squared cell residuals vanishes exactly when every
positive-mass summand vanishes.
\end{proof}

\begin{decision}[Acquisition packet after V54]
\label{dec:v13v54-next}
The random timing packet must retain
\[
 \left(
 T_h,Q_h,\Psi_{h,T_h}^+,\Psi_{h,T_h}^-,
 \mathbb E[T_h\mid Q_h],
 \mathbb E[T_h^2\mid Q_h],
 \mathbb E[T_hZ_h\mid Q_h]
 \right)                                                      \tag{54.25}
\]
on one same-history cylinder.  The next task is to combine its
dimension-weighted finite-difference error with the V52 adjacent
transport defects and state one sufficient local continuum
criterion for the clock-calibrated bounded germs.
\end{decision}

\begin{assessment}[Progress after V54]
\label{ass:v13v54}
The random renewal clock no longer hides a normalization choice.
The observed channel quotient converges to the duration-size-biased
odd germ under the sharp moment ratio \(K_h\mathbb E
T_h^2/\mathbb E T_h\to0\).  A map-valued covariance measures the
bias from an ordinary mean, and an explicit two-history example
proves that the separate marginals cannot determine even the
orientation of the fermion clock.
\end{assessment}

\begin{firewall}[Claim boundary after V54]
\label{fire:v13v54}
No joint duration/operation cylinder has been numerically populated.
The moment, orientation, source-bank, and adjacent transport bounds
remain open.  The conditional compiler proves finite
identifiability only; it does not construct the unbounded continuum
Dirac operator, an interacting Lorentzian dynamics, numerical
Standard Model parameters, or Einstein coupling.
\end{firewall}

\subsection{Parent record}

This V54 note appends the preceding material to the validated V53
source, whose record was
\[
\begin{array}{c|c}
 \text{lines}&17505\\
 \text{bytes}&617183\\
 \text{SHA--256}&
 \texttt{3DE057277A442EA26A7292914AB38D2DADD04900E1BF5190DD6A557A7788341D}.
\end{array}                                                     \tag{54.26}
\]

% =============================================================================
% END APPENDED THIRTEENTH-CYCLE V54 MATERIAL
% =============================================================================
% =============================================================================
% BEGIN APPENDED THIRTEENTH-CYCLE V55 MATERIAL
% =============================================================================

\section{Thirteenth-cycle V55: compulsory companion audit and the
bounded-germ continuum route}
\label{sec:v13v55}

\subsection{Current companion record}

The newest companion notes in the shared directory were inspected
read-only:
\[
\begin{array}{c|c|r|r|c}
 \text{programme}&\text{version}&\text{lines}&\text{bytes}&
 \text{SHA--256}\\ \hline
 {\rm Yang\!-\!Mills}&V88&29307&998813&
 \texttt{E0C63D89385077672116BD14E5C3BE1A700A8BA32BDC7F8D435D30BD71406855}\\
 {\rm Navier\!-\!Stokes}&V26&5319&278283&
 \texttt{82C887F8C2C69E4F28FAD7C3DC8DE5B9099CB5A4BF431EBA1FFF3E91935978AD}.
\end{array}                                                     \tag{55.1}
\]
The auxiliary Yang--Mills note remains at V5 with SHA--256
\[
 \texttt{306F1BB84CFA8A8D47EA569EC17E9545F9867C8D32B548EC9D9C444B4F3C0512}.
                                                                    \tag{55.2}
\]
No companion file was modified.

Yang--Mills V87 retained a complete trial trace but perturbed its
inverse through an additive unwhitened error.  Its exact card was
worse than the preceding two-sided Banach bound.  V88 instead
perturbs the already energy-scaled products multiplicatively and
proves
\[
\begin{split}
 \|H^{-1/2}YH^{-1/2}\|_F^2
 \leq{}&\operatorname{Re}\operatorname{Tr}(A^*B)\\
 &+\|A\|_F\|B\|_F
 \left[
 \frac1{(1-\eta_L)(1-\eta_R)}-1
 \right].                                                     \tag{55.3}
\end{split}
\]
On the common \(256\)-cell card it wins every cell and reduces the
two certified path energies from \(34.171961,54.878727\) to
\(20.396984,45.963686\).  It still proves no Yang--Mills continuum
or mass gap.

Navier--Stokes V26 is unchanged.  Its structured-input lesson has
already been absorbed into the complete source-table and
conditional timing constructions of V49--V54.

\subsection{Transfer to the timing problem}

Let \(\overline Z_h\) be the observable random-duration channel
quotient of V54 and let \(X_h\) be the marked Hamiltonian
derivation.  The primary finite object is their complete Gram
matrix
\[
 \overline\Gamma_h=
 \begin{pmatrix}
  \langle X_h,X_h\rangle&
  \langle X_h,\overline Z_h\rangle\\
  \langle\overline Z_h,X_h\rangle&
  \langle\overline Z_h,\overline Z_h\rangle
 \end{pmatrix}.                                               \tag{55.4}
\]
Its Schur residual is formed before any separate norm bound.  The
physical duration remainder
\[
 e_h:=K_h\frac{\mathbb E_hT_h^2}{\mathbb E_hT_h}              \tag{55.5}
\]
is then added at map level using
Theorem~\ref{thm:v13v54-random}.  Replacing (55.4) by separate
ceilings before taking its Schur complement would repeat the
unsuccessful additive architecture diagnosed in YM V87.

\begin{proposition}[V55 audit decision]
\label{prop:v13v55-decision}
No companion numerical constant is imported.  The next note shall
prove one conditional bounded-germ continuum theorem from:
\begin{enumerate}[label=\textup{(\roman*)}]
\item a common pressure-duration unit and positive mixed-Gram
orientation;
\item summable adjacent compatibility of the Hamiltonian
derivations;
\item convergence of the Gram-defined coefficients;
\item the complete observed residual
\(d_h\sqrt{\overline r_h}\to0\); and
\item the random-duration remainder \(e_h\to0\).
\end{enumerate}
The proof shall pass the exact finite Schur residual to a cb bound
before adding (55.5).

\end{proposition}

\begin{proof}
The YM result is an ordering principle for a positive matrix
calculation, not an SM estimate.  The complete Gram in (55.4)
contains the cancellation which defines the best calibrated
coefficient.  V51 upgrades its residual to cb norm, while V54
controls the independent finite-difference remainder.  These are
the exact ingredients listed above.
\end{proof}

\begin{assessment}[Position after V55]
\label{ass:v13v55}
The required audit is complete.  The deterministic and random
clock branches now share one exact finite matrix card.  The next
logical target is a norm-continuous dynamics on the protected
quasi-local algebra.  This is a bounded-screen continuum result;
the unbounded Dirac/form-domain problem remains a later and
strictly stronger gate.
\end{assessment}

\begin{firewall}[Claim boundary after V55]
\label{fire:v13v55}
The companion audit supplies no physical timing data.  The
conditions listed in Proposition~\ref{prop:v13v55-decision} are
not yet populated.  No unbounded continuum Dirac operator,
interacting quantum field, numerical coupling, Einstein equation,
or full SM--Einstein continuum is claimed.
\end{firewall}

\subsection{Parent record}

This V55 note appends the preceding material to the validated V54
source, whose record was
\[
\begin{array}{c|c}
 \text{lines}&17831\\
 \text{bytes}&627742\\
 \text{SHA--256}&
 \texttt{729063A4CDC1732E494D7CEBDC5CB05BEDDCA86BD941AEC6491F76D77A958599}.
\end{array}                                                     \tag{55.6}
\]

% =============================================================================
% END APPENDED THIRTEENTH-CYCLE V55 MATERIAL
% =============================================================================
% =============================================================================
% BEGIN APPENDED THIRTEENTH-CYCLE V56 MATERIAL
% =============================================================================

\section{Thirteenth-cycle V56: a clock-calibrated bounded
operation continuum}
\label{sec:v13v56}

\subsection{Inductive system and observed cards}

Let
\[
 \mathcal A_1\xrightarrow{\iota_1}\mathcal A_2
 \xrightarrow{\iota_2}\cdots,\qquad
 \mathcal A_n=M_{d_n}(\mathbb C),                             \tag{56.1}
\]
be a unital completely isometric \(C^*\)-inductive system, and
let \(\mathcal A_{\rm ql}\) be the norm completion of its algebraic
direct limit.  Write \(\iota_{n,m}\) for the composite embedding.
At cutoff \(n\), let
\[
 X_n(A)=-i[D_n,A],\qquad D_n=D_n^*,                           \tag{56.2}
\]
let \(Z_n^\sharp\) be the duration-weighted physical odd tangent
of V54, and let \(\overline Z_n\) be the observed finite-channel
quotient.

Using the complete matrix-unit source table, put
\[
\begin{split}
 g_n&=\|X_n\|_{\rm sup}^2,\\
 c_n&=\operatorname{Re}\langle X_n,\overline Z_n\rangle_{\rm sup},\\
 h_n&=\|\overline Z_n\|_{\rm sup}^2,\\
 a_n&=\frac{c_n}{g_n},\qquad
 r_n=h_n-\frac{c_n^2}{g_n}.                                  \tag{56.3}
\end{split}
\]
The last line is used on the positive branch \(g_n,c_n>0\).
The number \(r_n\) is formed from the complete Gram before any
separate entrywise majorization.

Define the adjacent derivation defect
\[
 \rho_n=
 \|X_{n+1}\iota_n-\iota_nX_n\|_{\rm cb}.                      \tag{56.4}
\]
For the random holding-time card, put
\[
 e_n=K_n\frac{\mathbb E_nT_n^2}{\mathbb E_nT_n}.              \tag{56.5}
\]
Theorem~\ref{thm:v13v54-random} gives
\[
 \|\overline Z_n-Z_n^\sharp\|_{\rm cb}\leq e_n.               \tag{56.6}
\]

\subsection{The conditional bounded-continuum theorem}

\begin{theorem}[Clock-calibrated bounded operation continuum]
\label{thm:v13v56-bounded-continuum}
Assume that all duration variables are expressed in the common
physical unit fixed by Theorem~\ref{thm:v13v53-pressure-clock},
and suppose:
\[
 0<\gamma\leq g_n\leq G,\qquad
 0<\kappa\leq c_n\leq C;                                     \tag{56.7}
\]
\[
 \sum_{n=1}^{\infty}
 \bigl(|g_{n+1}-g_n|+|c_{n+1}-c_n|\bigr)<\infty;              \tag{56.8}
\]
\[
 \sup_n\|X_n\|_{\rm cb}\leq M<\infty,\qquad
 \sum_{n=1}^{\infty}\rho_n<\infty;                            \tag{56.9}
\]
and
\[
 d_n\sqrt{r_n}\longrightarrow0,\qquad
 e_n\longrightarrow0.                                        \tag{56.10}
\]
Then:
\begin{enumerate}[label=\textup{(\roman*)}]
\item the positive coefficients \(a_n\) converge to a finite
\(a_\infty\) satisfying
\[
 \frac{\kappa}{G}\leq a_\infty\leq\frac C\gamma;              \tag{56.11}
\]
\item the \(X_n\) converge on every local observable to a bounded
star derivation
\[
 X_\infty:\mathcal A_{\rm ql}\longrightarrow\mathcal A_{\rm ql},
 \qquad \|X_\infty\|\leq M;                                  \tag{56.12}
\]
\item both the observed and physical odd germs converge locally,
uniformly on the unit ball of each fixed local algebra, to
\[
 \delta_\infty=a_\infty X_\infty;                            \tag{56.13}
\]
\item \(\delta_\infty\) generates the norm-continuous
one-parameter group of star automorphisms
\[
 \alpha_t=\exp(t\delta_\infty),\qquad t\in\mathbb R.           \tag{56.14}
\]
For \(A\in\mathcal A_N\), embedded at cutoff \(n\geq N\), the
physical-germ error obeys the explicit tail bound
\[
\begin{split}
 \|Z_n^\sharp(\iota_{N,n}A)
      -\delta_\infty(A)\|
 \leq\Bigl[
 e_n+d_n\sqrt{r_n}
 +M|a_n-a_\infty|
 +a_\infty\sum_{k=n}^{\infty}\rho_k
 \Bigr]\|A\|.                                                 \tag{56.15}
\end{split}
\]

\end{enumerate}
\end{theorem}

\begin{proof}
The quotient estimate used in V52 gives
\[
 |a_{n+1}-a_n|
 \leq
 \frac{|c_{n+1}-c_n|}{\gamma}
 +\frac{C|g_{n+1}-g_n|}{\gamma^2}.                            \tag{56.16}
\]
Thus (56.8) makes \((a_n)\) Cauchy.  The bounds in (56.11)
follow directly from (56.7).

Fix \(A\in\mathcal A_N\).  In the algebraic direct limit,
telescoping (56.4) gives, for \(m>n\geq N\),
\[
\begin{split}
 &\|X_m(\iota_{N,m}A)
 -\iota_{n,m}X_n(\iota_{N,n}A)\|\\
 &\hspace{35mm}\leq
 \left(\sum_{k=n}^{m-1}\rho_k\right)\|A\|.                    \tag{56.17}
\end{split}
\]
Hence the values are Cauchy and define \(X_\infty(A)\).
The uniform bound in (56.9) extends this map uniquely to
\(\mathcal A_{\rm ql}\).  Each \(X_n\) satisfies
\[
 X_n(AB)=X_n(A)B+AX_n(B),\qquad
 X_n(A^*)=X_n(A)^*.                                           \tag{56.18}
\]
Pass to the limit on local \(A,B\), then use density and
boundedness.  This proves (56.12).

Let
\[
 R_n=\overline Z_n-a_nX_n.                                   \tag{56.19}
\]
Corollary~\ref{cor:v13v51-cb-clock}, applied to the complete
observed Gram, gives
\[
 \|R_n\|_{\rm cb}\leq d_n\sqrt{r_n}.                          \tag{56.20}
\]
For \(A_n=\iota_{N,n}A\), insert
\[
\begin{split}
 Z_n^\sharp(A_n)-a_\infty X_\infty(A)
={}&(Z_n^\sharp-\overline Z_n)(A_n)+R_n(A_n)\\
 &+(a_n-a_\infty)X_n(A_n)
 +a_\infty\bigl(X_n(A_n)-X_\infty(A)\bigr).                  \tag{56.21}
\end{split}
\]
Equations (56.6), (56.17), and (56.20) prove (56.15).
Its right side tends to zero by (56.8)--(56.10).  Omitting the
first term proves the same statement for \(\overline Z_n\).
All bounds are proportional to \(\|A\|\), hence convergence is
uniform on each fixed local unit ball.

Finally \(\delta_\infty\) is a bounded star derivation.  Its
exponential series converges in operator norm.  The differential
equation
\[
 \frac{d}{dt}\alpha_t(A)=\delta_\infty(\alpha_t(A))            \tag{56.22}
\]
and the Leibniz rule show that
\(\alpha_t(AB)=\alpha_t(A)\alpha_t(B)\) and
\(\alpha_t(A^*)=\alpha_t(A)^*\).  The inverse is
\(\alpha_{-t}\).  Thus (56.14) is a norm-continuous
one-parameter automorphism group.
\end{proof}

\subsection{Convergence of the calibrated finite-cutoff groups}

Put
\[
 \delta_n=a_nX_n,\qquad
 \alpha_t^{(n)}=\exp(t\delta_n).                              \tag{56.23}
\]
Every \(\alpha_t^{(n)}\) is an isometric star automorphism of
\(\mathcal A_n\).  Its adjacent generator defect satisfies
\[
\begin{split}
 \zeta_n
 &:=
 \|\delta_{n+1}\iota_n-\iota_n\delta_n\|_{\rm cb}\\
 &\leq
 \frac C\gamma\,\rho_n
 +M|a_{n+1}-a_n|.                                             \tag{56.24}
\end{split}
\]
Consequently (56.8)--(56.9) imply
\[
 \sum_n\zeta_n<\infty.                                        \tag{56.25}
\]

\begin{corollary}[Uniform compact-time convergence of the
calibrated groups]
\label{cor:v13v56-groups}
For every \(A\in\mathcal A_N\), every \(T<\infty\), and
\(m>n\geq N\),
\[
 \sup_{|t|\leq T}
 \|\alpha_t^{(m)}(\iota_{N,m}A)
 -\iota_{n,m}\alpha_t^{(n)}(\iota_{N,n}A)\|
 \leq T\left(\sum_{k=n}^{m-1}\zeta_k\right)\|A\|.             \tag{56.26}
\]
The finite-cutoff calibrated automorphism groups therefore
converge uniformly on compact time intervals to the group
\((\alpha_t)\) in (56.14).

\end{corollary}

\begin{proof}
For one adjacent pair, Duhamel's identity gives
\[
 \alpha_t^{(n+1)}\iota_n-\iota_n\alpha_t^{(n)}
 =
 \int_0^t
 \alpha_{t-s}^{(n+1)}
 (\delta_{n+1}\iota_n-\iota_n\delta_n)
 \alpha_s^{(n)}\,ds.                                         \tag{56.27}
\]
Both automorphism groups are isometric, so the norm is at most
\(|t|\zeta_n\).  Telescope from \(n\) to \(m\), use complete
isometry of the embeddings, and take \(|t|\leq T\).  The series
(56.25) gives the compact-time limit.  Its generator agrees with
\(\delta_\infty\) on the dense local algebra by
Theorem~\ref{thm:v13v56-bounded-continuum}, so the limit group is
(56.14).
\end{proof}

\subsection{What has and has not crossed the continuum boundary}

Theorem~\ref{thm:v13v56-bounded-continuum} is a genuine
conditional continuum statement: once its finite rows are
populated, the clock-calibrated protected operation germ and its
Hamiltonian approximants have one norm-continuous limit dynamics.
Its proof combines the exact Gram cancellation, the
dimension-weighted cb conversion, the random-duration remainder,
and the adjacent Duhamel transport without replacing the mixed
Gram by separate ambient norms.

The theorem deliberately lives on a bounded protected screen.
A spacetime Dirac operator in the continuum is expected to be
unbounded, and a relativistic quantum field dynamics is generally
not norm continuous on the full field algebra.  Passing beyond
V56 therefore requires graph/form convergence on an exhausting
family of compact energy screens.

\begin{decision}[Next bottleneck after V56]
\label{dec:v13v56-next}
Construct the unbounded upgrade in three layers:
\begin{enumerate}[label=\textup{(\roman*)}]
\item choose positive compact-resolvent anchors \(Q_n\) and
spectral screens \(P_{n,R}\);
\item require V56 uniformly on each transported screen, with a
tail estimate that vanishes as \(R\to\infty\); and
\item prove Mosco or strong-resolvent convergence of the
Dirac--Yukawa quadratic forms on one common local core.
\end{enumerate}
The first task is to state a screen-diagonal theorem which makes
the order of the two limits \(n\to\infty\) and \(R\to\infty\)
explicit.
\end{decision}

\begin{assessment}[Progress after V56]
\label{ass:v13v56}
The clock programme has reached a conditional bounded continuum:
the coefficient, derivation, physical odd germ, and calibrated
automorphism groups all converge, with the computable tail
(56.15).  The remaining operator-theoretic obstruction is no
longer finite clock identifiability; it is uniform control of
unbounded domains and high-energy tails.
\end{assessment}

\begin{firewall}[Claim boundary after V56]
\label{fire:v13v56}
The hypotheses (56.7)--(56.10) have not been populated on a
cofinal physical regulator.  The group in (56.14) is the limit of
the calibrated Hamiltonian approximants; convergence of the
original finite physical channels over noninfinitesimal time
requires their subdivision/Duhamel error estimates as well.
Nothing here proves an unbounded continuum Dirac operator,
Lorentzian quantum field theory, interacting Standard Model
measure, numerical running or matching, Einstein dynamics, or the
full SM--Einstein continuum.
\end{firewall}

\subsection{Parent record}

This V56 note appends the preceding material to the validated V55
source, whose record was
\[
\begin{array}{c|c}
 \text{lines}&17972\\
 \text{bytes}&633033\\
 \text{SHA--256}&
 \texttt{4BE6FE3D75047E3886D761CD2AA24DEDD9A3444C72DAA0C69D566BCD2E7C3737}.
\end{array}                                                     \tag{56.28}
\]

% =============================================================================
% END APPENDED THIRTEENTH-CYCLE V56 MATERIAL
% =============================================================================
% =============================================================================
% BEGIN APPENDED THIRTEENTH-CYCLE V57 MATERIAL
% =============================================================================

\section{Thirteenth-cycle V57: screen-diagonal calibration and
the unbounded scale bridge}
\label{sec:v13v57}

\subsection{Existing analytic input and the missing bridge}

The Grand Tensor manuscript already proves:
\begin{enumerate}[label=\textup{(\roman*)}]
\item Mosco convergence of lower-bounded closed forms is equivalent
to strong resolvent convergence of their associated self-adjoint
operators and to compact-time strong convergence of their heat
semigroups; and
\item a vanishing compact graph-screen tail upgrades that strong
resolvent convergence to collective compactness and norm-resolvent
convergence.
\end{enumerate}
Those theorems are not repeated here.  The missing SM-specific
step is that the physical clock coefficient is reconstructed on
finite operation screens.  One must first show that the
screenwise coefficients select one scale and then prove that this
moving scale is compatible with Mosco and resolvent passage.

Let \(P_{n,R}\) be the finite spectral or graph screen at cutoff
\(n\) and threshold \(R\).  Apply the complete timing-Gram
construction of V51--V56 on this screen and denote its positive
coefficient by
\[
 a_{n,R}>0.                                                    \tag{57.1}
\]
Let \(s_{n,R}\) be a proved adjacent variation ceiling:
\[
 |a_{n+1,R}-a_{n,R}|\leq s_{n,R}.                             \tag{57.2}
\]
The bound (52.10), populated with the screenwise Gram entries, is
one source of \(s_{n,R}\).

\subsection{Two-parameter diagonal selection}

\begin{theorem}[Screen-diagonal clock coefficient]
\label{thm:v13v57-diagonal}
Assume:
\[
 \sum_{n=1}^{\infty}s_{n,R}<\infty
 \quad\text{for every fixed }R,                               \tag{57.3}
\]
and suppose there is a decreasing modulus
\(\omega_R\downarrow0\) such that
\[
 |a_{n,S}-a_{n,R}|\leq\omega_R
 \qquad(n\geq1,\ S\geq R).                                   \tag{57.4}
\]
Then the fixed-screen limits
\[
 a_R=\lim_{n\to\infty}a_{n,R}                                \tag{57.5}
\]
exist and converge as \(R\to\infty\) to one coefficient
\(a_\infty\).  More precisely,
\[
 |a_S-a_R|\leq\omega_R\quad(S\geq R),\qquad
 |a_R-a_\infty|\leq\omega_R.                                 \tag{57.6}
\]
Put
\[
 u_{n,R}=\sum_{k=n}^{\infty}s_{k,R}.                          \tag{57.7}
\]
For every increasing choice \(R(n)\to\infty\) satisfying
\[
 u_{n,R(n)}+\omega_{R(n)}\longrightarrow0,                   \tag{57.8}
\]
the diagonal coefficients
\[
 \widehat a_n:=a_{n,R(n)}                                     \tag{57.9}
\]
satisfy
\[
 |\widehat a_n-a_\infty|
 \leq u_{n,R(n)}+\omega_{R(n)}
 \longrightarrow0.                                           \tag{57.10}
\]
Such a slowly increasing diagonal \(R(n)\) always exists.

\end{theorem}

\begin{proof}
For fixed \(R\), (57.2)--(57.3) make \((a_{n,R})_n\)
Cauchy, proving (57.5).  Let \(n\to\infty\) in (57.4) to get the
first estimate in (57.6).  Thus \((a_R)_R\) is Cauchy; its limit
and the second estimate follow by letting \(S\to\infty\).

Telescope (57.2) from \(n\) to infinity:
\[
 |a_{n,R}-a_R|\leq u_{n,R}.                                  \tag{57.11}
\]
Equations (57.6) and (57.11), with \(R=R(n)\), prove
(57.10).

For existence, choose increasing integers \(N_j\) such that
\(u_{n,j}\leq2^{-j}\) whenever \(n\geq N_j\), increasing the
\(N_j\) if necessary.  Set \(R(n)=j\) for
\(N_j\leq n<N_{j+1}\).  Then \(R(n)\to\infty\) and
\[
 u_{n,R(n)}+\omega_{R(n)}
 \leq2^{-R(n)}+\omega_{R(n)}\longrightarrow0.
\]
\end{proof}

The uniform cross-screen condition (57.4) is load bearing.  Fixed
\(R\) convergence alone permits the limits \(a_R\) to oscillate
or diverge with \(R\), so it does not define an unbounded physical
scale.

\begin{countertheorem}[Fixed-screen convergence is insufficient]
\label{cth:v13v57-no-diagonal}
The array
\[
 a_{n,R}=R+\frac1n                                             \tag{57.12}
\]
has a convergent cutoff sequence at every fixed \(R\), with
\(a_R=R\), but has no finite \(R\to\infty\) limit.  Thus neither
pointwise screen convergence nor arbitrarily small adjacent
cutoff defects can replace the screen-tail modulus (57.4).
\end{countertheorem}

\begin{proof}
For fixed \(R\), \(a_{n,R}\to R\) and
\(|a_{n+1,R}-a_{n,R}|=1/(n(n+1))\) is summable.  The limiting
sequence \(a_R=R\) diverges.
\end{proof}

\subsection{Mosco stability under the recovered scale}

\begin{theorem}[Positive scalar stability of Mosco convergence]
\label{thm:v13v57-scaled-mosco}
Let \(q_n\) be nonnegative closed forms on the changing Hilbert
spaces of the Grand Tensor Mosco theorem, and suppose
\[
 q_n\xrightarrow{\rm Mosco}q_\infty.                          \tag{57.13}
\]
If \(b_n>0\) and \(b_n\to b_\infty>0\), then
\[
 b_nq_n\xrightarrow{\rm Mosco}b_\infty q_\infty.              \tag{57.14}
\]
In particular, if \(\widehat a_n\) is obtained from
Theorem~\ref{thm:v13v57-diagonal}, with \(a_\infty>0\), then
\[
 \widehat a_n^2q_n
 \xrightarrow{\rm Mosco}a_\infty^2q_\infty.                   \tag{57.15}
\]

\end{theorem}

\begin{proof}
Let \(x_n\rightharpoonup x\).  Since \(q_n[x_n]\geq0\) and
\(b_n\to b_\infty>0\),
\[
 \liminf_n b_nq_n[x_n]
 \geq b_\infty\liminf_nq_n[x_n]
 \geq b_\infty q_\infty[x].                                  \tag{57.16}
\]
For a Mosco recovery sequence \(x_n\to x\) with
\(q_n[x_n]\to q_\infty[x]\), multiplication by \(b_n\to
b_\infty\) gives
\(b_nq_n[x_n]\to b_\infty q_\infty[x]\).  These are the lower and
recovery conditions.  Take \(b_n=\widehat a_n^2\) for (57.15).
\end{proof}

If \(q_n[u]=\|D_nu\|^2\), (57.15) controls the graph energies of
the clock-scaled Dirac family.  It does not determine the sign of
\(D_n\); direct strong-resolvent control of \(D_n\) supplies that
additional information.

\subsection{Strong-resolvent stability of the signed operator}

Let \(J_n:\mathcal H_n\to\mathcal H_\infty\) be the comparison
isometries.  Write
\[
 D_n\xrightarrow{\rm sr}D_\infty                             \tag{57.17}
\]
when, for one and hence every \(z\in\mathbb C\setminus\mathbb R\),
\[
 J_n(D_n-z)^{-1}J_n^*
 \longrightarrow(D_\infty-z)^{-1}                            \tag{57.18}
\]
strongly.

\begin{theorem}[Unbounded clock-scale resolvent bridge]
\label{thm:v13v57-scaled-resolvent}
Let \(D_n,D_\infty\) be self-adjoint and satisfy (57.17).  If
\[
 \widehat a_n\longrightarrow a_\infty>0,                     \tag{57.19}
\]
then
\[
 \widehat a_nD_n\xrightarrow{\rm sr}a_\infty D_\infty.        \tag{57.20}
\]
For \(z\notin\mathbb R\), the resolvents are related by
\[
 (\widehat a_nD_n-z)^{-1}
 =\widehat a_n^{-1}
   (D_n-z/\widehat a_n)^{-1}.                                 \tag{57.21}
\]

\end{theorem}

\begin{proof}
Put \(w_n=z/\widehat a_n\) and
\(w=z/a_\infty\).  For all sufficiently large \(n\), the imaginary
parts of \(w_n\) stay a positive distance from zero.  The
resolvent identity gives
\[
 \|(D_n-w_n)^{-1}-(D_n-w)^{-1}\|
 \leq
 \frac{|w_n-w|}{|\operatorname{Im}w_n|
                  |\operatorname{Im}w|}
 \longrightarrow0.                                           \tag{57.22}
\]
At the fixed point \(w\), (57.17) gives embedded strong
convergence to \((D_\infty-w)^{-1}\).  Multiply by
\(\widehat a_n^{-1}\to a_\infty^{-1}\) in (57.21).  The limit is
\[
 a_\infty^{-1}(D_\infty-z/a_\infty)^{-1}
 =(a_\infty D_\infty-z)^{-1},
\]
which proves (57.20).
\end{proof}

\begin{corollary}[Calibrated graph and heat handoff]
\label{cor:v13v57-graph-heat}
Under Theorems~\ref{thm:v13v57-scaled-mosco} and
\ref{thm:v13v57-scaled-resolvent}:
\begin{enumerate}[label=\textup{(\roman*)}]
\item the nonnegative operators associated with
\(\widehat a_n^2q_n\) converge in strong resolvent sense and their
heat semigroups converge strongly, uniformly on compact
nonnegative time intervals;
\item if the Grand Tensor compact graph-screen profile vanishes,
the corresponding nonnegative resolvents converge in norm and
their isolated eigenvalues and finite spectral projections
converge with multiplicity; and
\item the signed Dirac resolvents converge as in (57.20).
\end{enumerate}
\end{corollary}

\begin{proof}
Items \textup{(i)} and \textup{(ii)} are the existing Grand
Tensor Mosco and compact-screen theorems applied after
Theorem~\ref{thm:v13v57-scaled-mosco}.  Item \textup{(iii)} is
Theorem~\ref{thm:v13v57-scaled-resolvent}.
\end{proof}

\subsection{Populatable screen ledger}

For each screen, the new data required beyond the existing Mosco
packet are
\[
 \left(
 g_{n,R},c_{n,R},h_{n,R},r_{n,R},
 s_{n,R},\omega_R,e_{n,R},d_{n,R}
 \right),                                                     \tag{57.23}
\]
where \(d_{n,R}=\operatorname{rank}P_{n,R}\) and
\[
 d_{n,R}\sqrt{r_{n,R}}+e_{n,R}\longrightarrow0
 \quad\text{for every fixed }R.                              \tag{57.24}
\]
The adjacent Gram estimate of V52 populates \(s_{n,R}\); comparison
of nested complete timing Grams populates \(\omega_R\); and the
joint duration moments of V54 populate \(e_{n,R}\).

\begin{proposition}[Order of limits]
\label{prop:v13v57-order}
Under (57.3)--(57.4) and (57.24), the valid order is:
\[
 \text{first }n\to\infty\text{ at fixed }R,\qquad
 \text{then }R\to\infty,                                     \tag{57.25}
\]
or one diagonal \(R(n)\) satisfying (57.8).  An uncontrolled
choice \(R=R(n)\to\infty\) is not justified by fixed-screen
convergence.

\end{proposition}

\begin{proof}
Theorem~\ref{thm:v13v57-diagonal} proves both valid passages.
Countertheorem~\ref{cth:v13v57-no-diagonal} proves the last
assertion.
\end{proof}

\begin{decision}[Next bottleneck after V57]
\label{dec:v13v57-next}
The scalar clock can now cross an unbounded Mosco/resolvent limit
once (57.23)--(57.24) are populated.  The remaining SM operator
gate is stronger: show that the finite Dirac--Yukawa operators
themselves have a common sign-sensitive strong-resolvent limit,
rather than only convergence of their squared graph forms.
The next attack should formulate a bounded-transform criterion
using
\[
 F_n=D_n(1+D_n^2)^{-1/2},                                    \tag{57.26}
\]
because \(F_n\) is bounded, retains the sign, and separates domain
control from the high-energy inverse transform.
\end{decision}

\begin{assessment}[Progress after V57]
\label{ass:v13v57}
The clock coefficient no longer stops at bounded source screens.
A summable cutoff modulus plus a vanishing cross-screen modulus
selects one coefficient, and positive scalar multiplication
preserves both Mosco convergence and signed strong-resolvent
convergence.  The required diagonal order and its failure mode are
explicit.
\end{assessment}

\begin{firewall}[Claim boundary after V57]
\label{fire:v13v57}
The screen ledger (57.23) is not populated, and strong-resolvent
convergence of the finite Dirac--Yukawa family is still an
assumption.  Mosco convergence of \(\|D_n\cdot\|^2\) alone loses
the sign of \(D_n\).  No unbounded continuum Dirac operator,
spectral triple, chiral measure, interacting Standard Model
continuum, Einstein coupling, or Lorentzian reconstruction is
proved here.
\end{firewall}

\subsection{Parent record}

This V57 note appends the preceding material to the validated V56
source, whose record was
\[
\begin{array}{c|c}
 \text{lines}&18292\\
 \text{bytes}&643816\\
 \text{SHA--256}&
 \texttt{93189046C6FB9C71BBCE72485593CD19EE85869298927CA98E82BB05BC9E3DE7}.
\end{array}                                                     \tag{57.27}
\]

% =============================================================================
% END APPENDED THIRTEENTH-CYCLE V57 MATERIAL
% =============================================================================
% =============================================================================
% BEGIN APPENDED THIRTEENTH-CYCLE V58 MATERIAL
% =============================================================================

\section{Thirteenth-cycle V58: bounded-transform reconstruction of
the signed unbounded Dirac limit}
\label{sec:v13v58}

\subsection{Why the bounded transform is the correct carrier}

For a self-adjoint operator \(D_n\) on \(\mathcal H_n\), define
\[
 F_n=D_n(1+D_n^2)^{-1/2}.                                    \tag{58.1}
\]
Then \(F_n=F_n^*\), \(\|F_n\|\leq1\), and \(F_n\) retains the sign
of \(D_n\).  In contrast, the graph form
\(\|D_nu\|^2\) determines only \(D_n^2\).

Let \(J_n:\mathcal H_n\to\mathcal H_\infty\) be the isometries of
a strict Hilbert direct system, let \(P_n=J_nJ_n^*\), and assume
\(P_n\to I\) strongly.  Extend \(F_n\) to
\(\mathcal H_\infty\) by
\[
 \widetilde F_n=J_nF_nJ_n^*,                                 \tag{58.2}
\]
which is zero on \(P_n^\perp\).

Choose a countable dense family of compatible source vectors
\((\xi_j)_{j\geq1}\) in the algebraic direct limit.  For all
sufficiently large \(n\), \(\xi_j\in P_n\mathcal H_\infty\).
Suppose finite mixed Gram cards give ceilings
\[
 \|(\widetilde F_{n+1}-\widetilde F_n)\xi_j\|
 \leq\varepsilon_{n,j}.                                      \tag{58.3}
\]

\begin{lemma}[Dense source-table convergence]
\label{lem:v13v58-dense}
If
\[
 \sum_n\varepsilon_{n,j}<\infty
 \qquad\text{for every }j,                                   \tag{58.4}
\]
then \(\widetilde F_n\) converges strongly to a self-adjoint
contraction \(F_\infty\) on \(\mathcal H_\infty\).

\end{lemma}

\begin{proof}
For each \(j\), (58.3)--(58.4) make
\((\widetilde F_n\xi_j)_n\) Cauchy.  The operators are uniformly
bounded by one.  Approximation of an arbitrary vector by the dense
linear span of the \(\xi_j\) therefore extends convergence to
every vector and defines a contraction \(F_\infty\).  Finally,
for arbitrary \(x,y\),
\[
 \langle F_\infty x,y\rangle
 =\lim_n\langle\widetilde F_nx,y\rangle
 =\lim_n\langle x,\widetilde F_ny\rangle
 =\langle x,F_\infty y\rangle,                               \tag{58.5}
\]
so \(F_\infty\) is self-adjoint.
\end{proof}

\subsection{The endpoint-mass gate}

A self-adjoint contraction \(F\) is the bounded transform of a
densely defined self-adjoint operator on the whole Hilbert space
if and only if
\[
 \ker(I-F^2)=\{0\}.                                           \tag{58.6}
\]
The forbidden eigenspaces at \(+1\) and \(-1\) represent spectral
mass which has escaped to infinite positive or negative energy.

For \(0<\epsilon<1/4\), choose a continuous function
\(\chi_\epsilon:[-1,1]\to[0,1]\) satisfying
\[
 \chi_\epsilon(s)=0\quad(|s|\leq1-2\epsilon),\qquad
 \chi_\epsilon(s)=1\quad(|s|\geq1-\epsilon).                 \tag{58.7}
\]

\begin{theorem}[Source-visible endpoint tightness]
\label{thm:v13v58-endpoint}
Assume Lemma~\ref{lem:v13v58-dense} and suppose that, for every
dense source vector \(\xi_j\),
\[
 \lim_{\epsilon\downarrow0}\,
 \limsup_{n\to\infty}
 \|\chi_\epsilon(\widetilde F_n)\xi_j\|=0.                   \tag{58.8}
\]
Then
\[
 \ker(I-F_\infty^2)=\{0\}.                                   \tag{58.9}
\]

\end{theorem}

\begin{proof}
For fixed \(\epsilon\), strong convergence of uniformly bounded
self-adjoint operators implies
\[
 \chi_\epsilon(\widetilde F_n)
 \longrightarrow\chi_\epsilon(F_\infty)                      \tag{58.10}
\]
strongly.  This follows first for polynomials and then for
continuous functions by uniform polynomial approximation on
\([-1,1]\).  Hence (58.8) gives
\[
 \lim_{\epsilon\downarrow0}
 \|\chi_\epsilon(F_\infty)\xi_j\|=0.                         \tag{58.11}
\]
Let \(E_\infty\) be the spectral measure of \(F_\infty\).
Since \(\chi_\epsilon=1\) at both endpoints,
\[
 \|E_\infty(\{-1,1\})\xi_j\|
 \leq\|\chi_\epsilon(F_\infty)\xi_j\|.                        \tag{58.12}
\]
Let \(\epsilon\downarrow0\).  The endpoint projection vanishes on
a dense set, hence vanishes identically.  Its range is
\(\ker(I-F_\infty^2)\), proving (58.9).
\end{proof}

Condition (58.8) is a sign-sensitive high-energy tightness row.
It is weaker than a uniform spectral gap away from \(\pm1\),
which would incorrectly exclude an unbounded limit.

\subsection{Reconstruction and resolvent convergence}

Define
\[
 \varphi(s)=\frac{s}{\sqrt{1-s^2}},\qquad -1<s<1.             \tag{58.13}
\]
Under (58.9), spectral calculus defines the densely defined
self-adjoint operator
\[
 D_\infty=\varphi(F_\infty),\qquad
 \mathcal D(D_\infty)=
 \left\{u:
 \int_{(-1,1)}\frac{s^2}{1-s^2}\,
 d\langle u,E_\infty(s)u\rangle<\infty\right\}.               \tag{58.14}
\]
Its bounded transform is \(F_\infty\).

\begin{theorem}[Bounded-transform strong-resolvent criterion]
\label{thm:v13v58-transform-resolvent}
Suppose \(\widetilde F_n\to F_\infty\) strongly and
\(\ker(I-F_\infty^2)=\{0\}\).  Then the self-adjoint operators
\(D_n\) satisfy
\[
 D_n\xrightarrow{\rm sr}D_\infty.                            \tag{58.15}
\]

\end{theorem}

\begin{proof}
For \(z\in\mathbb C\setminus\mathbb R\), define on \([-1,1]\)
\[
 g_z(s)=
 \begin{cases}
 \displaystyle
 \frac{\sqrt{1-s^2}}{s-z\sqrt{1-s^2}},& |s|<1,\\
 0,&s=\pm1.
 \end{cases}                                                  \tag{58.16}
\]
The denominator cannot vanish for real \(s\) and nonreal \(z\),
and the endpoint limits are zero, so \(g_z\) is continuous.
Functional calculus gives
\[
 (D_n-z)^{-1}=g_z(F_n),\qquad
 (D_\infty-z)^{-1}=g_z(F_\infty).                             \tag{58.17}
\]
Because \(\widetilde F_n\) is zero on \(P_n^\perp\),
\[
 g_z(\widetilde F_n)
 =
 J_ng_z(F_n)J_n^*+g_z(0)(I-P_n).                             \tag{58.18}
\]
The left side converges strongly to \(g_z(F_\infty)\) by the same
continuous-functional-calculus argument as (58.10), while the
last term tends strongly to zero because \(P_n\to I\).  Equations
(58.17)--(58.18) prove embedded strong resolvent convergence.
\end{proof}

\begin{countertheorem}[Endpoint escape destroys the operator
limit]
\label{cth:v13v58-endpoint-escape}
On a fixed nonzero Hilbert space, let
\[
 F_n=(1-n^{-1})I,\qquad
 D_n=\frac{1-n^{-1}}
 {\sqrt{1-(1-n^{-1})^2}}\,I.                                 \tag{58.19}
\]
Then \(F_n\to I\) in norm, but
\[
 (D_n-z)^{-1}\longrightarrow0                                \tag{58.20}
\]
in norm for every nonreal \(z\).  The zero operator cannot be the
resolvent of a densely defined self-adjoint operator.  Thus strong
bounded-transform convergence without (58.8) is insufficient.

\end{countertheorem}

\begin{proof}
The scalar multiplying \(I\) in \(D_n\) tends to \(+\infty\), so
(58.20) is immediate.  A resolvent is injective with dense range,
whereas the zero map is neither.
\end{proof}

\subsection{Insertion of the screen-diagonal clock scale}

Let \(\widehat a_n\to a_\infty>0\) be the coefficient obtained in
Theorem~\ref{thm:v13v57-diagonal}.  The bounded transform of the
scaled operator \(\widehat a_nD_n\) is
\[
 \widehat F_n
 =
 \frac{\widehat a_nF_n}
 {\sqrt{\,1+(\widehat a_n^2-1)F_n^2\,}}.                     \tag{58.21}
\]
Define similarly
\[
 \widehat F_\infty
 =
 \frac{a_\infty F_\infty}
 {\sqrt{\,1+(a_\infty^2-1)F_\infty^2\,}}.                    \tag{58.22}
\]

\begin{theorem}[Clock-scaled bounded-transform handoff]
\label{thm:v13v58-scaled-transform}
Under the hypotheses of
Theorem~\ref{thm:v13v57-diagonal},
Lemma~\ref{lem:v13v58-dense}, and
Theorem~\ref{thm:v13v58-endpoint},
\[
 \widetilde{\widehat F}_n
 \longrightarrow\widehat F_\infty                           \tag{58.23}
\]
strongly, \(\ker(I-\widehat F_\infty^2)=\{0\}\), and
\[
 \widehat a_nD_n\xrightarrow{\rm sr}
 a_\infty D_\infty.                                          \tag{58.24}
\]
Thus the screen-reconstructed physical clock scale preserves the
sign-sensitive unbounded Dirac limit.

\end{theorem}

\begin{proof}
For \(a>0\), set
\[
 f_a(s)=\frac{as}{\sqrt{1+(a^2-1)s^2}}.                      \tag{58.25}
\]
On every compact interval \(a\in[a_-,a_+]\subset(0,\infty)\),
the functions \(f_a\) depend uniformly continuously on
\((a,s)\in[a_-,a_+]\times[-1,1]\).  Since
\(\widehat a_n\to a_\infty>0\) and
\(\widetilde F_n\to F_\infty\) strongly, polynomial
approximation gives
\(f_{\widehat a_n}(\widetilde F_n)\to
f_{a_\infty}(F_\infty)\) strongly.  The complement correction is
zero here because \(f_a(0)=0\), proving (58.23).

The map \(f_a\) sends \((-1,1)\) homeomorphically onto
\((-1,1)\) and sends only \(\pm1\) to \(\pm1\).  Hence the
endpoint spectral projection of \(\widehat F_\infty\) equals that
of \(F_\infty\) and vanishes.  Applying
Theorem~\ref{thm:v13v58-transform-resolvent} to
\(\widehat F_n\) proves (58.24).  The reconstructed inverse
transform is \(a_\infty D_\infty\) by direct substitution in
(58.13).
\end{proof}

\subsection{Finite acquisition form}

For a source vector \(\xi_j\), the squared defect in (58.3) is the
finite Gram quantity
\[
 \|(\widetilde F_{n+1}-\widetilde F_n)\xi_j\|^2
 =
 \|\widetilde F_{n+1}\xi_j\|^2
 +\|\widetilde F_n\xi_j\|^2
 -2\operatorname{Re}
 \langle\widetilde F_{n+1}\xi_j,
        \widetilde F_n\xi_j\rangle.                           \tag{58.26}
\]
The mixed term must be recorded on one transported source
carrier.  Separate marginal norms do not determine it.

For a polynomial ceiling \(p_{\epsilon,m}\) uniformly
approximating \(\chi_\epsilon\), the quantity
\[
 \|p_{\epsilon,m}(\widetilde F_n)\xi_j\|^2                   \tag{58.27}
\]
is a finite word Gram.  Uniform polynomial error converts it into
a rigorous ceiling for the endpoint tail in (58.8).  Therefore
both the convergence and no-escape hypotheses admit finite
source-table certificates.

\begin{decision}[Next bottleneck after V58]
\label{dec:v13v58-next}
The sign-sensitive unbounded operator gate now has a terminating
conditional certificate.  The next step is to connect
strong-resolvent convergence of the clock-scaled
Dirac--Yukawa operators to convergence of the chiral determinant
or Pfaffian line.  This requires a relative trace-class or
Schatten resolvent row; strong resolvent convergence alone does
not control determinant phases.
\end{decision}

\begin{assessment}[Progress after V58]
\label{ass:v13v58}
A bounded-transform source packet now reconstructs the signed
unbounded Dirac limit.  Summable dense-source defects give a
strong contraction limit, the endpoint-tail card excludes
spectral mass at infinite energy, and the screen-diagonal clock
coefficient passes through the inverse bounded transform.  This
removes the need to assume sign-sensitive strong-resolvent
convergence in V57 once the new finite rows are populated.
\end{assessment}

\begin{firewall}[Claim boundary after V58]
\label{fire:v13v58}
No bounded-transform defect series or endpoint-tail card has been
populated for the physical regulator.  Strong resolvent
convergence does not imply trace-norm resolvent convergence,
determinant convergence, phase selection, anomaly cancellation,
or reflection positivity of the interacting measure.  The
unbounded Dirac, chiral line, quantum Standard Model, Einstein
backreaction, and full SM--Einstein continuum remain open.
\end{firewall}

\subsection{Parent record}

This V58 note appends the preceding material to the validated V57
source, whose record was
\[
\begin{array}{c|c}
 \text{lines}&18637\\
 \text{bytes}&655596\\
 \text{SHA--256}&
 \texttt{5431AD0320F9AD0610F2E62CDA3CABE6D4099326CCD4ACA1801480789C94FAD8}.
\end{array}                                                     \tag{58.28}
\]

% =============================================================================
% END APPENDED THIRTEENTH-CYCLE V58 MATERIAL
% =============================================================================
% =============================================================================
% BEGIN APPENDED THIRTEENTH-CYCLE V59 MATERIAL
% =============================================================================

\section{Thirteenth-cycle V59: relative trace-class convergence
and the determinant-phase handoff}
\label{sec:v13v59}

\subsection{Scope relative to the Grand Tensor manuscript}

The manuscript already constructs finite determinant and Pfaffian
lines, keeps zero-mode divisors explicit, and proves that a complex
determinant does not orient a real Pfaffian line.  Those results are
not repeated.  V58 supplies only sign-sensitive strong-resolvent
convergence of the clock-scaled Dirac--Yukawa operators.  The
missing analytic implication is false without a stronger
Schatten-class row: strong resolvent convergence does not control
a Fredholm determinant or its phase.

Let \(D_n\) be the clock-scaled self-adjoint operator from V58 and
let \(D_n^0\) be a physically specified reference operator on the
same cutoff Hilbert space.  At a common
\(z\in\mathbb C\setminus\mathbb R\), assume the relative operator
\[
 K_n(z)
 :=
 (D_n-D_n^0)(D_n^0-z)^{-1}                                   \tag{59.1}
\]
extends from the reference domain to a trace-class operator.
Then
\[
 (D_n-z)(D_n^0-z)^{-1}=I+K_n(z),                              \tag{59.2}
\]
and the relative Fredholm determinant is
\[
 \Delta_n(z)=\det_{\!F}(I+K_n(z)).                            \tag{59.3}
\]
The reference family is part of the physical packet; changing it
changes the relative determinant by its own cocycle.

\subsection{Trace-norm convergence is the determinant gate}

\begin{theorem}[Relative determinant continuity]
\label{thm:v13v59-det-continuity}
Let \(K_n,K\) be trace-class operators on one Hilbert space and
suppose
\[
 \|K_n-K\|_1\longrightarrow0.                                \tag{59.4}
\]
Then
\[
 \det_{\!F}(I+K_n)\longrightarrow\det_{\!F}(I+K),             \tag{59.5}
\]
with the quantitative bound
\[
 \left|\det_{\!F}(I+K_n)-\det_{\!F}(I+K)\right|
 \leq
 \|K_n-K\|_1
 \exp\!\left(1+\|K_n\|_1+\|K\|_1\right).                     \tag{59.6}
\]
If the limiting determinant is nonzero, its normalized phase also
converges:
\[
 \frac{\det_{\!F}(I+K_n)}
      {|\det_{\!F}(I+K_n)|}
 \longrightarrow
 \frac{\det_{\!F}(I+K)}
      {|\det_{\!F}(I+K)|}.                                   \tag{59.7}
\]

\end{theorem}

\begin{proof}
The Fredholm expansion is
\[
 \det_{\!F}(I+K)
 =\sum_{m=0}^{\infty}\operatorname{Tr}(\wedge^mK).            \tag{59.8}
\]
The exterior-power estimate
\[
 \|\wedge^mA-\wedge^mB\|_1
 \leq
 \frac{\|A-B\|_1(\|A\|_1+\|B\|_1)^{m-1}}{(m-1)!}
 \qquad(m\geq1)                                               \tag{59.9}
\]
follows by telescoping the \(m\) tensor factors, restricting to
the alternating subspace, and using the ideal property of the
trace norm.  Summing (59.9), with a harmless factor \(e\) to cover
the zeroth and small-norm cases, gives (59.6).  Hence (59.5).
If the limit is nonzero, the determinants stay away from zero for
large \(n\), and continuity of \(w\mapsto w/|w|\) on
\(\mathbb C\setminus\{0\}\) proves (59.7).
\end{proof}

The constant in (59.6) is intentionally nonsharp; the decisive
feature is trace-norm continuity.

\begin{countertheorem}[Strong convergence does not control a
determinant]
\label{cth:v13v59-strong-no-det}
Let \((e_j)\) be the standard basis of \(\ell^2\), let
\(P_n=|e_n\rangle\langle e_n|\), and fix
\(\lambda\in\mathbb C\setminus\{-1,0\}\).  Then
\[
 K_n=\lambda P_n\longrightarrow0                             \tag{59.10}
\]
strongly, while
\[
 \det_{\!F}(I+K_n)=1+\lambda
 \not\longrightarrow1=\det_{\!F}(I).                         \tag{59.11}
\]
Thus even bounded strong convergence, and hence a fortiori bare
strong-resolvent convergence of an underlying operator family,
does not imply determinant convergence.

\end{countertheorem}

\begin{proof}
For every \(x\in\ell^2\),
\(\|P_nx\|=|\langle e_n,x\rangle|\to0\).  The only nontrivial
eigenvalue of \(K_n\) is \(\lambda\), giving (59.11).
\end{proof}

The obstruction is moving trace-class mass:
\[
 \|K_n\|_1=|\lambda|                                          \tag{59.12}
\]
at every \(n\), although no fixed vector detects it strongly.

\subsection{A summable adjacent acquisition criterion}

Embed the cutoff operators into the common comparison Hilbert
space and write them again as \(K_n\).  Suppose finite
source-complete singular-value cards give
\[
 \theta_n\geq\|K_{n+1}-K_n\|_1.                              \tag{59.13}
\]

\begin{theorem}[Summable nuclear transport]
\label{thm:v13v59-nuclear-transport}
If
\[
 \sum_{n=1}^{\infty}\theta_n<\infty,                          \tag{59.14}
\]
then \(K_n\) converges in trace norm to a trace-class operator
\(K_\infty\).  Moreover,
\[
 \|K_n-K_\infty\|_1
 \leq\Theta_n:=\sum_{k=n}^{\infty}\theta_k,                   \tag{59.15}
\]
and
\[
 \|K_n\|_1\leq
 M_1:=\|K_1\|_1+\sum_{k=1}^{\infty}\theta_k.                  \tag{59.16}
\]
The determinant tail satisfies
\[
 |\Delta_n-\Delta_\infty|
 \leq
 \Theta_n\exp(1+2M_1).                                       \tag{59.17}
\]

\end{theorem}

\begin{proof}
The trace-class ideal is a Banach space.  Telescoping (59.13) and
using (59.14) makes \((K_n)\) Cauchy in trace norm, proving the
limit and (59.15).  A second telescope gives (59.16).
Apply Theorem~\ref{thm:v13v59-det-continuity} with the uniform
bound (59.16) to obtain (59.17).
\end{proof}

At finite cutoff,
\[
 \|K_{n+1}-K_n\|_1
 =
 \operatorname{Tr}\sqrt{(K_{n+1}-K_n)^*
                         (K_{n+1}-K_n)}                      \tag{59.18}
\]
is computed from the singular values of one transported complete
mixed Gram.  The two marginal singular-value lists do not
determine (59.18).

\subsection{Compact parameter families and a local phase section}

Let \(U\) be a compact, connected, simply connected parameter
space, representing one anomaly-free chart of background fields.
Suppose \(K_n(q)\) and \(K_\infty(q)\) are trace-class and
\[
 \sup_{q\in U}\|K_n(q)-K_\infty(q)\|_1\longrightarrow0,       \tag{59.19}
\]
with
\[
 \sup_{n,q}\|K_n(q)\|_1\leq M_1.                              \tag{59.20}
\]

\begin{theorem}[Divisor-free local determinant phase]
\label{thm:v13v59-local-phase}
Under (59.19)--(59.20), the relative determinants converge
uniformly:
\[
 \sup_{q\in U}|\Delta_n(q)-\Delta_\infty(q)|
 \longrightarrow0.                                           \tag{59.21}
\]
If
\[
 \inf_{q\in U}|\Delta_\infty(q)|\geq m>0,                    \tag{59.22}
\]
then, for all sufficiently large \(n\), \(\Delta_n\) is nowhere
zero.  After fixing one phase value at a base point \(q_0\), there
are continuous logarithms on \(U\) for which
\[
 \log\Delta_n-\log\Delta_\infty\longrightarrow0              \tag{59.23}
\]
uniformly, modulo the fixed common \(2\pi i\) branch.

\end{theorem}

\begin{proof}
The determinant estimate (59.6), uniformly bounded by (59.20),
gives (59.21).  For large \(n\), its error is below \(m/2\), so
\(|\Delta_n|\geq m/2\).  A continuous nonvanishing complex
function on a simply connected parameter space admits a
continuous logarithm.  More explicitly, the ratio
\(\Delta_n/\Delta_\infty\) lies uniformly in a small disk around
one; its principal logarithm tends uniformly to zero.  Add this
to the base-point-anchored logarithm of \(\Delta_\infty\), proving
(59.23).
\end{proof}

The hypothesis (59.22) is the divisor-free branch.  When it fails,
zero multiplicity, kernel/cokernel lines, and crossing data must be
retained exactly as in the manuscript's zero-safe determinant and
Pfaffian packets.

\subsection{Relative Cayley phase}

For a self-adjoint \(D\), let
\[
 U_D=(D-i)(D+i)^{-1}=I-2i(D+i)^{-1}.                          \tag{59.24}
\]
If \(D,D^0\) are self-adjoint and
\[
 (D+i)^{-1}-(D^0+i)^{-1}\in\mathcal S_1,                     \tag{59.25}
\]
then \(U_D-U_{D^0}\in\mathcal S_1\), and
\[
 W(D,D^0):=U_DU_{D^0}^*=I+\mathcal S_1.                      \tag{59.26}
\]

\begin{proposition}[Trace-class Cayley phase]
\label{prop:v13v59-cayley}
For a convergent sequence, take \(D^0\) to be fixed on the
comparison carrier.  The Fredholm determinant
\[
 \zeta(D,D^0)=\det_{\!F}W(D,D^0)                              \tag{59.27}
\]
has unit modulus.  If the relative resolvent differences in
(59.25) converge in trace norm, then the phases (59.27) converge.
This supplies a cutoff-stable relative spectral phase, but not an
absolute chiral determinant-line trivialization.

\end{proposition}

\begin{proof}
The ideal property gives (59.26).  The operator \(W(D,D^0)\) is
unitary, so its Fredholm eigenvalues have unit modulus and their
absolutely convergent determinant product has unit modulus.
Trace-norm convergence of the resolvent differences gives
trace-norm convergence of the Cayley differences and then of
\(W-I\).  Apply Theorem~\ref{thm:v13v59-det-continuity}.
\end{proof}

\subsection{The required determinant packet}

The analytic rows which must accompany V58 are therefore
\[
 \boxed{
 \left(
 D_n^0,z,K_n(z),\theta_n,\Theta_n,
 \Delta_n(z),\text{divisor data},
 \text{base phase},\text{chart overlaps}
 \right).}                                                    \tag{59.28}
\]
For a parameter family, (59.14) must hold uniformly on each
protected compact chart, or be replaced by the direct uniform
trace-norm convergence (59.19).  The physical reference and its
change-of-reference cocycle must be transported on the same
history as the chiral operator.

\begin{decision}[Next bottleneck after V59]
\label{dec:v13v59-next}
The next compulsory companion audit occurs at V60.  After it, the
programme should separate:
\begin{enumerate}[label=\textup{(\roman*)}]
\item local determinant-section convergence, now controlled by
V59;
\item global line holonomy on chart overlaps;
\item the real Pfaffian orientation and crossing sign; and
\item scalar measure positivity and OS reflection.
\end{enumerate}
A convergent local Fredholm determinant cannot be used as a
substitute for the last three rows.
\end{decision}

\begin{assessment}[Progress after V59]
\label{ass:v13v59}
The unbounded clock-scaled Dirac limit now has an exact conditional
handoff to a local relative determinant and relative Cayley phase.
Summable adjacent nuclear defects give trace-norm convergence and
an explicit determinant tail.  Moving rank-one mass proves why
the strong-resolvent result of V58 cannot supply this conclusion
by itself.
\end{assessment}

\begin{firewall}[Claim boundary after V59]
\label{fire:v13v59}
No relative trace-class card, divisor-free margin, or phase anchor
has been populated for the physical SM regulator.  A local
relative determinant does not orient a real Pfaffian line, cancel
an anomaly, construct a scalarized fermion measure, prove
reflection positivity, or produce the interacting
SM--Einstein continuum.
\end{firewall}

\subsection{Parent record}

This V59 note appends the preceding material to the validated V58
source, whose record was
\[
\begin{array}{c|c}
 \text{lines}&18985\\
 \text{bytes}&667698\\
 \text{SHA--256}&
 \texttt{650DBA36E55164111B892A81C6FF45E1F70CF1049275B1587D51D44CB6C77999}.
\end{array}                                                     \tag{59.29}
\]

% =============================================================================
% END APPENDED THIRTEENTH-CYCLE V59 MATERIAL
% =============================================================================
% =============================================================================
% BEGIN APPENDED THIRTEENTH-CYCLE V60 MATERIAL
% =============================================================================

\section{Thirteenth-cycle V60: compulsory companion audit and the
global line decision}
\label{sec:v13v60}

\subsection{Current read-only companion record}

The newest companion notes were inspected without modification:
\[
\begin{array}{c|c|r|r|c}
 \text{programme}&\text{version}&\text{lines}&\text{bytes}&
 \text{SHA--256}\\ \hline
 {\rm Yang\!-\!Mills}&V89&29492&1005424&
 \texttt{A2F30094ED229D2708E7848C881DE56D1F45DD51DE04F8C777844DAE0413D6EF}\\
 {\rm Navier\!-\!Stokes}&V26&5319&278283&
 \texttt{82C887F8C2C69E4F28FAD7C3DC8DE5B9099CB5A4BF431EBA1FFF3E91935978AD}.
\end{array}                                                     \tag{60.1}
\]
The auxiliary Yang--Mills V5 record remains unchanged, with
SHA--256
\[
 \texttt{306F1BB84CFA8A8D47EA569EC17E9545F9867C8D32B548EC9D9C444B4F3C0512}.
                                                                    \tag{60.2}
\]

Yang--Mills V89 refines the multiplicative trace card at
action-dependent resolutions.  It certifies the Wilson path at
\(512\) cells with energy ceiling \(14.953051\), and the endpoint
path at \(1024\) cells with ceiling \(18.671264\).  Both lie within
the predeclared factor-two diagnostic scale.  It then freezes
those unequal resolutions and uses reflection and exact-zero
relations to reduce nine vertex families to five independent
batched products.  No torus cover or Yang--Mills mass gap is
claimed.

Navier--Stokes V26 remains unchanged and supplies no new constant
or theorem for the determinant-line problem.

\subsection{Transfer to V59}

A determinant line is assembled on a finite background cover
\((U_\alpha)\).  The relative trace-class complexity, divisor
distance, and required screen depth can differ between charts.
The YM V89 lesson therefore transfers as follows:
\begin{enumerate}[label=\textup{(\roman*)}]
\item allow a chart-dependent screen \(R_\alpha(n)\), provided
each satisfies the V57 diagonal condition and the V59 nuclear
tail;
\item compute all overlap transition functions on one common
reference carrier before taking phases;
\item batch conjugate, reversal-related, and symmetry-equivalent
overlaps, while retaining a direct card for every independent
cycle; and
\item freeze a resolution only after its own trace-norm and
divisor-margin gates pass.
\end{enumerate}
A common resolution is optional bookkeeping.  A common overlap
carrier and exact cocycle identities are mathematical
requirements.

\begin{proposition}[V60 audit decision]
\label{prop:v13v60-decision}
No numerical companion bound is imported.  The next note shall
prove that uniformly convergent finite transition functions on a
fixed finite cover define a limiting Hermitian line bundle, that
compatible local determinant sections glue in the limit, and that
overlap holonomies converge when their connection data converge.
It shall keep the real Pfaffian sign as a separate
\(\mathbb Z_2\) cocycle.
\end{proposition}

\begin{proof}
YM V89 concerns a different finite inverse problem, so its
numerical data are inapplicable.  Its valid structural transfer is
asymmetric resolution with common-product batching.  V59 already
provides local trace-norm convergence; the absent operation is
precisely the assembly of its local phases across overlaps.
\end{proof}

\begin{assessment}[Position after V60]
\label{ass:v13v60}
The required audit is complete.  The local analytic determinant
problem has a conditional trace-class solution, but the global
line can still carry nontrivial overlap holonomy or a real
Pfaffian sign.  These global data must be passed before any local
phase is promoted to a scalar fermionic measure.
\end{assessment}

\begin{firewall}[Claim boundary after V60]
\label{fire:v13v60}
No overlap transition card, global holonomy, Pfaffian orientation,
or scalarized measure is populated.  The audit proves no anomaly
cancellation, reflection positivity, quantum Standard Model
continuum, Einstein backreaction, or full SM--Einstein continuum.
\end{firewall}

\subsection{Parent record}

This V60 note appends the preceding material to the validated V59
source, whose record was
\[
\begin{array}{c|c}
 \text{lines}&19326\\
 \text{bytes}&679146\\
 \text{SHA--256}&
 \texttt{E379FF4AA306825443E21A0193970347E28B42B706A7BDB072B71D394D53D4AD}.
\end{array}                                                     \tag{60.3}
\]

% =============================================================================
% END APPENDED THIRTEENTH-CYCLE V60 MATERIAL
% =============================================================================
% =============================================================================
% BEGIN APPENDED THIRTEENTH-CYCLE V61 MATERIAL
% =============================================================================

\section{Thirteenth-cycle V61: cofinal determinant-line gluing,
holonomy, and the separate Pfaffian sign}
\label{sec:v13v61}

\subsection{Finite-cover transition data}

Let \(X\) be a compact background parameter space and let
\((U_\alpha)_{\alpha=1}^N\) be a fixed finite good cover.  At
cutoff \(n\), suppose the local relative determinants of V59 are
represented in unitary frames whose overlap functions are
\[
 g_{\alpha\beta}^{(n)}:
 U_\alpha\cap U_\beta\longrightarrow U(1).                   \tag{61.1}
\]
Assume the exact relations
\[
 g_{\alpha\alpha}^{(n)}=1,\qquad
 g_{\beta\alpha}^{(n)}=(g_{\alpha\beta}^{(n)})^{-1},\qquad
 g_{\alpha\beta}^{(n)}
 g_{\beta\gamma}^{(n)}
 g_{\gamma\alpha}^{(n)}=1                                    \tag{61.2}
\]
on every nonempty double or triple overlap.  These are line-frame
relations, not consequences of the separate determinant moduli.

\begin{theorem}[Uniform cofinal gluing of Hermitian lines]
\label{thm:v13v61-line-limit}
Suppose that, on every double overlap,
\[
 \|g_{\alpha\beta}^{(n)}-g_{\alpha\beta}\|_\infty
 \longrightarrow0.                                           \tag{61.3}
\]
Then the limits satisfy (61.2) and define a Hermitian line bundle
\(L_\infty\to X\).  If local sections
\(s_\alpha^{(n)}:U_\alpha\to\mathbb C\) obey
\[
 s_\alpha^{(n)}
 =g_{\alpha\beta}^{(n)}s_\beta^{(n)}                          \tag{61.4}
\]
and converge locally uniformly to \(s_\alpha\), then
\[
 s_\alpha=g_{\alpha\beta}s_\beta                             \tag{61.5}
\]
and the \(s_\alpha\) glue to a continuous section
\(s_\infty\) of \(L_\infty\).

\end{theorem}

\begin{proof}
Multiplication and inversion are continuous on \(U(1)\).
Taking the uniform limit in each of the finitely many identities
(61.2) preserves them.  The standard quotient of
\(\coprod_\alpha U_\alpha\times\mathbb C\) by
\[
 (x,z,\beta)\sim(x,g_{\alpha\beta}(x)z,\alpha)                \tag{61.6}
\]
is therefore a Hermitian line bundle.  On a compact subset of an
overlap,
\[
\begin{split}
 \|s_\alpha-g_{\alpha\beta}s_\beta\|
 \leq{}&
 \|s_\alpha-s_\alpha^{(n)}\|
 +\|g_{\alpha\beta}^{(n)}\|
   \|s_\beta^{(n)}-s_\beta\|\\
 &+\|g_{\alpha\beta}^{(n)}-g_{\alpha\beta}\|
   \|s_\beta\|,                                               \tag{61.7}
\end{split}
\]
and every term tends to zero.  This proves (61.5), which is
exactly the gluing condition.
\end{proof}

\begin{corollary}[Divisor-free trivialization criterion]
\label{cor:v13v61-trivial}
Under Theorem~\ref{thm:v13v61-line-limit}, if
\[
 \inf_{\alpha,x\in U_\alpha}|s_\alpha(x)|\geq m>0,            \tag{61.8}
\]
then \(s_\infty\) is nowhere zero and \(L_\infty\) is trivial.
Without a compatible global section, nonvanishing of each
separately computed local determinant does not prove triviality.

\end{corollary}

\begin{proof}
The compatible local sections define one global section.  The
lower bound makes it nowhere zero.  The map
\[
 X\times\mathbb C\longrightarrow L_\infty,\qquad
 (x,z)\longmapsto z\,s_\infty(x)                              \tag{61.9}
\]
is then a continuous fibrewise unitary isomorphism after
normalizing \(s_\infty\).  The last assertion follows because
arbitrary local sections need not obey (61.5).
\end{proof}

\subsection{Connection and holonomy passage}

Let \(\mathcal A_\alpha^{(n)}\) be the local imaginary-valued
connection one-form in the chosen unitary frame.  On overlaps,
assume
\[
 \mathcal A_\beta^{(n)}
 =
 \mathcal A_\alpha^{(n)}
 +(g_{\alpha\beta}^{(n)})^{-1}
  dg_{\alpha\beta}^{(n)}.                                    \tag{61.10}
\]
Let \(\gamma\) be a piecewise \(C^1\) loop, subdivided so that
each segment lies in one chart.  Its holonomy is the usual ordered
product of segment exponentials and transition values at the
subdivision points; because the structure group is \(U(1)\), no
path ordering is needed.

\begin{theorem}[Cofinal holonomy convergence]
\label{thm:v13v61-holonomy}
Assume (61.3), and for every segment \(\gamma_j\subset
U_{\alpha_j}\) suppose
\[
 \int_{\gamma_j}
 |\mathcal A_{\alpha_j}^{(n)}
  -\mathcal A_{\alpha_j}|
 \longrightarrow0.                                           \tag{61.11}
\]
Then
\[
 \operatorname{Hol}_{L_n}(\gamma)
 \longrightarrow
 \operatorname{Hol}_{L_\infty}(\gamma).                      \tag{61.12}
\]
The conclusion is invariant under compatible changes of local
unitary frame.

\end{theorem}

\begin{proof}
For the fixed finite subdivision, (61.11) gives convergence of
each exponential
\[
 \exp\!\left(-\int_{\gamma_j}
 \mathcal A_{\alpha_j}^{(n)}\right),                          \tag{61.13}
\]
and (61.3) gives convergence of every transition value at a
crossing point.  A finite product of convergent unit complex
numbers converges, proving (61.12).

Under a frame change \(u_\alpha:U_\alpha\to U(1)\), the transition
and connection data change by
\[
 g_{\alpha\beta}\mapsto
 u_\alpha^{-1}g_{\alpha\beta}u_\beta,\qquad
 \mathcal A_\alpha\mapsto
 \mathcal A_\alpha+u_\alpha^{-1}du_\alpha.                   \tag{61.14}
\]
The endpoint factors on successive segments telescope around the
closed loop, so the holonomy is unchanged.
\end{proof}

\begin{corollary}[Global phase obstruction survives the cutoff]
\label{cor:v13v61-phase-obstruction}
If a loop \(\gamma\) has
\[
 \operatorname{Hol}_{L_\infty}(\gamma)\ne1,                  \tag{61.15}
\]
then no choice of local phase frames can turn the limiting
connection into a globally exact phase on that component.
Local convergence of the Fredholm determinants in V59 does not
remove this obstruction.

\end{corollary}

\begin{proof}
A globally exact unitary connection has trivial holonomy around
every closed loop.  Holonomy is frame invariant by
Theorem~\ref{thm:v13v61-holonomy}.
\end{proof}

\subsection{Summable overlap acquisition}

Suppose the finite overlap panels give
\[
 \|g_{\alpha\beta}^{(n+1)}
   -g_{\alpha\beta}^{(n)}\|_\infty
 \leq\eta_{\alpha\beta,n},\qquad
 \sum_n\eta_{\alpha\beta,n}<\infty.                           \tag{61.16}
\]
Then the transition functions converge uniformly and
\[
 \|g_{\alpha\beta}^{(n)}-g_{\alpha\beta}\|_\infty
 \leq\sum_{k=n}^{\infty}\eta_{\alpha\beta,k}.                 \tag{61.17}
\]
Thus Theorem~\ref{thm:v13v61-line-limit} has a direct
finite-card input.  Each squared pointwise overlap defect is a
two-column source Gram; uniformity requires a certified chart
cover or a Lipschitz modulus in addition to sampled values.

The chart resolutions may differ as allowed in V60.  What must be
common is the transported overlap carrier on
\(U_\alpha\cap U_\beta\), so that the mixed term in
\[
 |g_{\alpha\beta}^{(n+1)}
   -g_{\alpha\beta}^{(n)}|^2
 =2-2\operatorname{Re}\!
 \left(\overline{g_{\alpha\beta}^{(n+1)}}
       g_{\alpha\beta}^{(n)}\right)                           \tag{61.18}
\]
is physically represented.

\subsection{The real Pfaffian cocycle is separate}

On a genuine real skew branch, let
\[
 \epsilon_{\alpha\beta}^{(n)}:
 U_\alpha\cap U_\beta\longrightarrow\{+1,-1\}                \tag{61.19}
\]
be the Pfaffian-frame transition sign.  It satisfies the real
analogue of (61.2).

\begin{theorem}[Cofinal Pfaffian orientation alternative]
\label{thm:v13v61-pfaffian}
Suppose the connected components of all overlaps are fixed and
\((\epsilon_{\alpha\beta}^{(n)})_n\) is uniformly Cauchy on each
component.  Then every sign is eventually constant in \(n\) on
that component and defines a limiting
\(\mathbb Z_2\)-valued cocycle \(\epsilon_{\alpha\beta}\).
The limiting real Pfaffian line is orientable if and only if
there are locally constant signs
\(\sigma_\alpha:U_\alpha\to\{\pm1\}\) such that
\[
 \epsilon_{\alpha\beta}
 =\sigma_\alpha\sigma_\beta^{-1}.                            \tag{61.20}
\]
The square of the Pfaffian cocycle is always trivial:
\[
 \epsilon_{\alpha\beta}^2=1.                                 \tag{61.21}
\]
Therefore convergence and triviality of the determinant-square
line do not determine Pfaffian orientation.

\end{theorem}

\begin{proof}
Distinct signs have distance two.  A uniformly Cauchy sequence
in \(\{\pm1\}\) is therefore eventually constant.  The finite
cocycle identities pass to that eventual value.  A choice of
\(\sigma_\alpha\) changes every local real frame to make all
transition signs \(+1\), so (61.20) is sufficient for an
orientation.  Conversely, an orientation supplies such
compatible positive frames and hence the signs \(\sigma_\alpha\).
Equation (61.21) is immediate and shows that squaring loses the
orientation class.
\end{proof}

\begin{countertheorem}[A determinant square cannot see the
Möbius sign]
\label{cth:v13v61-mobius}
The real Möbius line over \(S^1\) has a transition sign
\(-1\) around one overlap cycle and is nonorientable.  Its tensor
square has transition \(+1\) and is trivial.  Thus a positive
determinant square is compatible with a nontrivial Pfaffian
\(\mathbb Z_2\) holonomy.

\end{countertheorem}

\begin{proof}
Transport of a real frame once around the Möbius line returns its
negative, which is the nontrivial element of
\(H^1(S^1;\mathbb Z_2)\).  Tensoring the frame with itself squares
the sign to \(+1\).
\end{proof}

\subsection{Global line packet}

The finite-to-continuum line packet is
\[
 \boxed{
 \left(
 \Delta_{\alpha}^{(n)},
 g_{\alpha\beta}^{(n)},
 \eta_{\alpha\beta,n},
 \mathcal A_\alpha^{(n)},
 \operatorname{Hol}_{n}(\gamma_j),
 \text{divisors},
 \epsilon_{\alpha\beta}^{(n)}
 \right).}                                                    \tag{61.22}
\]
The loops \((\gamma_j)\) must generate the declared finite
topological audit class, and the cover must be fine enough to
resolve every retained divisor and zero-mode chart.

\begin{decision}[Next bottleneck after V61]
\label{dec:v13v61-next}
Once (61.22) is populated, a scalar fermionic amplitude exists
only on a branch where the complex line is trivialized, its
connection holonomy is physically accepted or cancelled, and the
real sign is oriented when a Pfaffian branch is used.  The next
analytic step is to combine that scalarized fermion weight with
the bosonic cylinder and prove preservation of
Osterwalder--Schrader positivity under cutoff passage with
uniform integrability.
\end{decision}

\begin{assessment}[Progress after V61]
\label{ass:v13v61}
The local relative determinants of V59 now have a rigorous global
handoff.  Uniform overlap convergence constructs the limiting
Hermitian line, compatible local sections glue, connection
holonomies converge, and the real Pfaffian sign persists as a
separate discrete cocycle.  Local nonvanishing, global
trivialization, flat phase, and real orientation are now distinct
finite-card decisions.
\end{assessment}

\begin{firewall}[Claim boundary after V61]
\label{fire:v13v61}
No physical overlap, connection, loop, divisor, or Pfaffian card
has been populated.  The limiting determinant line may be
nontrivial and may have nontrivial holonomy.  No scalar fermionic
measure, anomaly cancellation, reflection-positive interacting
continuum, Einstein backreaction, or full SM--Einstein continuum
is proved.
\end{firewall}

\subsection{Parent record}

This V61 note appends the preceding material to the validated V60
source, whose record was
\[
\begin{array}{c|c}
 \text{lines}&19441\\
 \text{bytes}&683878\\
 \text{SHA--256}&
 \texttt{A409F5F82A592368004EB12133802EF9A5603E12CC9AE9F2FF6F8174DC4AD912}.
\end{array}                                                     \tag{61.23}
\]

% =============================================================================
% END APPENDED THIRTEENTH-CYCLE V61 MATERIAL
% =============================================================================
% =============================================================================
% BEGIN APPENDED THIRTEENTH-CYCLE V62 MATERIAL
% =============================================================================

\section{Thirteenth-cycle V62: scalarized determinant weights,
normalization stability, and the OS projective handoff}
\label{sec:v13v62}

\subsection{Existing reflection-positive input}

The Grand Tensor manuscript already proves three facts used here:
\begin{enumerate}[label=\textup{(\roman*)}]
\item reflection positivity of a finite cylinder is equivalent to
positivity of its complete reflection Gram;
\item a positive reflected-square boundary factor preserves
reflection positivity by the Schur product theorem; and
\item reflection-positive finite probabilities with summable
projective total-variation defects have a reflection-positive
inverse-limit cylinder probability.
\end{enumerate}
V62 does not repeat those theorems.  It supplies the intervening
estimate when the bosonic and scalarized fermionic data first
produce unnormalized positive cylinders.

Assume the V61 line and phase gates have selected a scalar branch.
At cutoff \(n\), let \(\Lambda_n\) be the resulting finite positive
unnormalized measure on \(\Omega_n\), let
\[
 Z_n=\Lambda_n(\Omega_n)>0,\qquad
 \mu_n=Z_n^{-1}\Lambda_n,                                    \tag{62.1}
\]
and let
\(\pi_n:\Omega_{n+1}\to\Omega_n\) be the reflection-equivariant
coarse map.  The exact unnormalized mismatch is
\[
 \epsilon_n=
 \|(\pi_n)_*\Lambda_{n+1}-\Lambda_n\|_{\rm TV}.               \tag{62.2}
\]
Because the spaces are finite, (62.2) is a finite sum of absolute
weight differences and is directly computable from the complete
same-history cylinder.

\subsection{Normalization lemma}

\begin{lemma}[Positive normalization stability]
\label{lem:v13v62-normalization}
Let \(\Lambda,\Lambda'\) be finite positive measures with masses
\(Z,Z'\geq z_*>0\).  Then
\[
 \left\|\frac{\Lambda}{Z}-\frac{\Lambda'}{Z'}\right\|_{\rm TV}
 \leq\frac{2}{z_*}\|\Lambda-\Lambda'\|_{\rm TV}.              \tag{62.3}
\]

\end{lemma}

\begin{proof}
Put \(\epsilon=\|\Lambda-\Lambda'\|_{\rm TV}\).  Since the
constant function one has supremum norm one,
\[
 |Z-Z'|\leq\epsilon.                                          \tag{62.4}
\]
Add and subtract \(\Lambda'/Z\):
\[
\begin{split}
 \left\|\frac{\Lambda}{Z}-\frac{\Lambda'}{Z'}\right\|_{\rm TV}
 &\leq
 \frac{\|\Lambda-\Lambda'\|_{\rm TV}}{Z}
 +\left|\frac1Z-\frac1{Z'}\right|\|\Lambda'\|_{\rm TV}\\
 &=\frac{\epsilon+|Z-Z'|}{Z}
 \leq\frac{2\epsilon}{z_*}.                                  \tag{62.5}
\end{split}
\]
\end{proof}

The factor two is harmless; the mass floor is not.

\begin{countertheorem}[Vanishing unnormalized error without a
partition floor]
\label{cth:v13v62-no-floor}
On \(\Omega=\{0,1\}\), let
\[
 \Lambda_n=2^{-n}\delta_0,\qquad
 \Lambda_n'=2^{-n}\delta_1.                                  \tag{62.6}
\]
Then
\[
 \|\Lambda_n-\Lambda_n'\|_{\rm TV}=2^{1-n}\longrightarrow0,  \tag{62.7}
\]
while their normalized probabilities are \(\delta_0\) and
\(\delta_1\), whose total-variation distance is maximal.  Thus an
absolute determinant-weight error cannot be passed through
normalization when the partition mass collapses.

\end{countertheorem}

\begin{proof}
Both masses equal \(2^{-n}\).  Normalization cancels that factor,
giving the two disjoint point masses.
\end{proof}

\subsection{Unnormalized-to-projective theorem}

\begin{theorem}[Reflection-positive normalized projective
handoff]
\label{thm:v13v62-projective}
Suppose:
\[
 \inf_nZ_n\geq z_*>0,\qquad
 \sum_{n=1}^{\infty}\epsilon_n<\infty;                        \tag{62.8}
\]
every normalized \(\mu_n\) is reflection positive; and every
\(\pi_n\) is reflection equivariant and preserves the declared
positive-time algebra.  Then
\[
 \delta_n^{\rm TV}:=
 \|(\pi_n)_*\mu_{n+1}-\mu_n\|_{\rm TV}
 \leq\frac{2\epsilon_n}{z_*},                                \tag{62.9}
\]
so
\[
 \sum_n\delta_n^{\rm TV}<\infty.                              \tag{62.10}
\]
Consequently the existing Grand Tensor projective-limit theorem
produces an exactly compatible inverse-limit cylinder
probability whose restriction to every finite positive-time
algebra is reflection positive.  At level \(m\), its marginal
obeys
\[
 \|\overline\mu_m-\mu_m\|_{\rm TV}
 \leq\frac2{z_*}\sum_{j\geq m}\epsilon_j.                     \tag{62.11}
\]

\end{theorem}

\begin{proof}
Pushforward preserves total mass, so the two positive measures in
(62.2) have masses \(Z_{n+1}\) and \(Z_n\).  Apply
Lemma~\ref{lem:v13v62-normalization} to obtain (62.9).
Equations (62.8)--(62.10) then meet the hypotheses of the
manuscript's reflection-positive projective-limit theorem, whose
tail estimate gives (62.11).
\end{proof}

\subsection{How the determinant line enters the finite OS test}

On a boundary decomposition
\(\Omega_n=B_n^-\times B_n^+\), let \(K_n^{\rm bos}\) be the
bosonic reflection kernel.  After the V61 global scalarization,
let \(H_n^{\rm fer}\) be the scalar fermionic crossing kernel.
The unnormalized combined kernel is
\[
 K_n^{\rm tot}=K_n^{\rm bos}\circ H_n^{\rm fer}.              \tag{62.12}
\]
A sufficient exact finite gate is
\[
 K_n^{\rm bos}\succeq0,\qquad
 H_n^{\rm fer}\succeq0,\qquad
 K_n^{\rm tot}(a,b)\geq0
 \ \text{as a scalar measure weight}.                         \tag{62.13}
\]
The first two conditions give
\(K_n^{\rm tot}\succeq0\) by the already proved Schur-product
theorem.  The last condition makes the cylinder a positive
measure rather than only a positive sesquilinear kernel.

\begin{proposition}[Line triviality is earlier than OS
positivity]
\label{prop:v13v62-line-before-os}
The V61 existence of a global nowhere-zero determinant section
does not imply any condition in (62.13).  To enter
Theorem~\ref{thm:v13v62-projective}, one must additionally:
\begin{enumerate}[label=\textup{(\roman*)}]
\item choose or physically cancel its global phase so that the
unnormalized scalar weights are nonnegative;
\item test the complete fermionic crossing matrix
\(H_n^{\rm fer}\succeq0\); and
\item verify the partition floor and the exact pushforward error
in (62.8).
\end{enumerate}

\end{proposition}

\begin{proof}
A line trivialization identifies every fibre with \(\mathbb C\);
it does not restrict the resulting complex coordinate to
\([0,\infty)\), and it does not impose matrix positivity on the
values at two reflected boundary arguments.  These are precisely
the first two additional conditions.  Normalization and
projective passage require Lemma~\ref{lem:v13v62-normalization},
which supplies the third.
\end{proof}

This proposition uses the manuscript's existing counterexample
showing that pointwise positive boundary weights can still have a
negative reflection eigenvalue.  No new copy of that
counterexample is needed.

\subsection{Trace-norm determinant errors and cylinder weights}

Suppose a protected chart contains no divisor and V59 gives a
uniform determinant error
\[
 \sup_{\omega\in\Omega_n}
 |\Delta_n(\omega)-\Delta_n^{\rm tar}(\omega)|
 \leq q_n.                                                    \tag{62.14}
\]
Let \(B_n\) be the positive bosonic unnormalized measure on that
finite chart.  If the accepted scalar phase makes both determinant
weights nonnegative, then
\[
 \|\Delta_nB_n-\Delta_n^{\rm tar}B_n\|_{\rm TV}
 \leq q_n B_n(\Omega_n).                                     \tag{62.15}
\]
Thus the determinant estimate (59.17) enters the projective
ledger only after multiplication by the bosonic mass.

\begin{proof}
For a finite positive measure,
\[
 \int|\Delta_n-\Delta_n^{\rm tar}|\,dB_n
 \leq q_nB_n(\Omega_n),
\]
which is (62.15).
\end{proof}

If the bosonic masses are not uniformly bounded, a small uniform
determinant error need not be a small unnormalized cylinder error.
The correct acquisition is the weighted \(L^1(B_n)\) quantity in
(62.15), or a Hölder estimate using a populated bosonic moment
bound.

\subsection{Correlation passage}

Theorem~\ref{thm:v13v62-projective} constructs a cylinder
probability, not automatically a Radon field measure or
convergence of unbounded composite operators.  Those conclusions
use the manuscript's existing tightness and uniform-integrability
theorem.  The new determinant-specific input is a uniform moment
row for the scalar weight.  For a protected observable \(O_n\),
one must retain, for some \(\delta>0\),
\[
 \sup_n
 \frac{\int |O_n|^{1+\delta}\,d\Lambda_n}{Z_n}<\infty.        \tag{62.16}
\]
For variances or quadratic insertions the corresponding
\(2+\delta\) row is required.  The division by \(Z_n\) again shows
why the partition floor and normalized moment control cannot be
replaced by an unnormalized determinant ceiling.

\begin{decision}[Next bottleneck after V62]
\label{dec:v13v62-next}
The determinant-to-OS interface is now reduced to finite rows:
\[
 \boxed{
 \left(
 K_n^{\rm bos},H_n^{\rm fer},Z_n,z_*,
 \epsilon_n,q_nB_n(\Omega_n),
 \text{tightness},\text{uniform integrability}
 \right).}                                                    \tag{62.17}
\]
After these are populated, the existing Grand Tensor theorem
gives a protected Euclidean continuum.  The next genuinely
stronger step is Lorentzian reconstruction: identify the OS
Hamiltonian generated by Euclidean time, prove that the
clock-scaled Dirac/current sector is incident with that same
Hamiltonian, and establish locality and a spectrum condition.
\end{decision}

\begin{assessment}[Progress after V62]
\label{ass:v13v62}
A scalarized determinant weight now has a rigorous route into the
existing reflection-positive projective-limit theorem.
Summable unnormalized errors survive normalization precisely under
a partition-function floor, with the explicit tail (62.11).
Line triviality, scalar positivity, crossing-matrix positivity,
normalization, tightness, and uniform integrability remain
separate checkable rows.
\end{assessment}

\begin{firewall}[Claim boundary after V62]
\label{fire:v13v62}
None of the finite rows in (62.17) has been populated for the
physical SM regulator.  Even their positive branch yields only
the protected Euclidean continuum already described
conditionally in the Grand Tensor manuscript.  It does not prove
full Euclidean covariance, clustering, a mass gap, Lorentzian
reconstruction, numerical Standard Model parameters, Einstein
backreaction, or the full SM--Einstein continuum.
\end{firewall}

\subsection{Parent record}

This V62 note appends the preceding material to the validated V61
source, whose record was
\[
\begin{array}{c|c}
 \text{lines}&19785\\
 \text{bytes}&695699\\
 \text{SHA--256}&
 \texttt{A52D1DAF02689EB8092F02AD886E87F7810FD86D21CD86A3EFE820CBA6614828}.
\end{array}                                                     \tag{62.18}
\]

% =============================================================================
% END APPENDED THIRTEENTH-CYCLE V62 MATERIAL
% =============================================================================
% =============================================================================
% BEGIN APPENDED THIRTEENTH-CYCLE V63 MATERIAL
% =============================================================================

\section{Thirteenth-cycle V63: type-correct OS-Hamiltonian
incidence and the moving-core limit}
\label{sec:v13v63}

\subsection{Upstream audit and the type correction}

The Grand Tensor manuscript already proves the finite
energy-density/Hamiltonian incidence theorem: on a source-complete
common graph core, the Gram of \(E_\Sigma-H_{\rm OS}\) vanishes
exactly when the two operators agree, and equality on an
exhausting common core identifies their closures.  It also proves
the selected Euclidean local-net and source-cyclic translation
theorems.  Those statements are not repeated.

The V7 attachment proves a different, earlier obstruction: an
inserted gauge multiplication writer on a finite heat-bath
cylinder does not equal the actual derivative of that kinetic
Hamiltonian, even modulo a connection commutator and scalar
energy shift.  That result prevents reuse of the old inserted
column as an OS-energy source.

There is a further typing issue.  The clock-scaled
Dirac--Yukawa operator \(a_nD_n\) of V58 acts on a one-particle
spinor space.  The OS Hamiltonian acts on the reconstructed
many-body Hilbert space.  They cannot be compared before a
physical occurrence lift has been specified.  On the quasi-free
reference branch the canonical lift is fermionic second
quantization; the interacting candidate must also contain the
bosonic, gauge, interaction, counterterm, boundary, and vacuum
pieces.

\subsection{The quasi-free reference lift}

Let \(D_n=D_n^*\) be the signed one-particle operators of V58 and
let \(a_n\to a_\infty>0\).  For a fixed \(m\geq0\), define the
positive one-particle energy
\[
 h_n=\sqrt{a_n^2D_n^2+m^2},\qquad
 h_\infty=\sqrt{a_\infty^2D_\infty^2+m^2}.                   \tag{63.1}
\]
On the antisymmetric Fock space
\[
 \mathfrak F_-(\mathcal H)
 =\bigoplus_{k=0}^{\infty}\bigwedge^k\mathcal H,              \tag{63.2}
\]
write \(d\Gamma(h)\) for the self-adjoint second quantization,
with vacuum energy zero.

\begin{theorem}[Strong-resolvent one-particle-to-Fock handoff]
\label{thm:v13v63-fock}
Assume
\[
 D_n\xrightarrow{\rm sr}D_\infty,\qquad
 a_n\longrightarrow a_\infty>0.                              \tag{63.3}
\]
After the compatible Hilbert-space embeddings are applied,
\[
 h_n\xrightarrow{\rm sr}h_\infty                             \tag{63.4}
\]
and
\[
 d\Gamma(h_n)\xrightarrow{\rm sr}d\Gamma(h_\infty).           \tag{63.5}
\]
Equivalently, for \(t\geq0\),
\[
 e^{-t\,d\Gamma(h_n)}
 =\Gamma_-(e^{-th_n})
 \longrightarrow
 \Gamma_-(e^{-th_\infty})
 =e^{-t\,d\Gamma(h_\infty)}                                  \tag{63.6}
\]
strongly.  The convergence is uniform for \(t\) in compact
subsets of \([0,\infty)\) whenever the one-particle semigroups
converge with that uniformity.

\end{theorem}

\begin{proof}
For \(\lambda>0\), the scalar functions
\[
 f_n(x)=
 \left(\sqrt{a_n^2x^2+m^2}+\lambda\right)^{-1}                \tag{63.7}
\]
converge uniformly on \(\mathbb R\) to the analogous \(f_\infty\).
They are continuous and vanish at infinity.  Strong-resolvent
functional calculus for \(D_n\), together with the uniform change
from \(a_n\) to \(a_\infty\), gives
\[
 (h_n+\lambda)^{-1}\longrightarrow
 (h_\infty+\lambda)^{-1}                                     \tag{63.8}
\]
strongly, which is (63.4).

Hence \(e^{-th_n}\to e^{-th_\infty}\) strongly, with each
operator a contraction.  On a decomposable \(k\)-particle vector,
\[
 \bigwedge^ke^{-th_n}
 (u_1\wedge\cdots\wedge u_k)
 =
 e^{-th_n}u_1\wedge\cdots\wedge e^{-th_n}u_k,                \tag{63.9}
\]
so the \(k\)-particle semigroups converge strongly by telescoping
the \(k\) factors.  Finite-particle vectors are dense in Fock
space.  Truncating an arbitrary Fock vector at particle number
\(K\) and using the contraction bound controls the remaining
tail uniformly in \(n\).  This proves (63.6); the same argument
preserves compact-time uniformity.  The
semigroup--strong-resolvent equivalence gives (63.5).
\end{proof}

The theorem constructs the quasi-free reference energy.  It does
not manufacture the interacting terms or identify the result
with the OS Hamiltonian.

\subsection{One-particle agreement does not determine the OS
Hamiltonian}

On a finite-dimensional one-particle carrier, let
Let \(N=d\Gamma(I)\) be particle number and let
\(H_0=d\Gamma(h)\).

\begin{countertheorem}[Invisible two-particle interaction]
\label{cth:v13v63-two-particle}
For every real \(\lambda\), the self-adjoint operator
\[
 H_\lambda=H_0+\lambda N(N-1)                                \tag{63.10}
\]
agrees with \(H_0\) on the vacuum and one-particle sectors, but on
the \(k\)-particle sector
\[
 H_\lambda-H_0=\lambda k(k-1)I.                              \tag{63.11}
\]
Thus the complete one-particle Dirac operator, its clock
calibration, determinant, and all vacuum/one-particle matrix
elements cannot decide the interacting OS Hamiltonian.  A
two-particle source already separates the family when
\(\lambda\ne0\).

\end{countertheorem}

\begin{proof}
The number operator equals \(kI\) on
\(\bigwedge^k\mathcal H\).  Hence \(N(N-1)\) vanishes for
\(k=0,1\) and equals \(k(k-1)I\) on grade \(k\).
\end{proof}

This is the many-body form of the held-out-word requirement.  A
one-particle occurrence lift is necessary but not sufficient.

\subsection{The interacting candidate and its direct residual}

At cutoff \(n\), let the physically assembled energy candidate on
the OS Hilbert space be
\[
 E_n=
 H_{{\rm bos},n}
 +d\Gamma(h_n)
 +V_{{\rm int},n}
 +C_{{\rm ct},n}
 +B_{\partial,n}
 +c_{{\rm vac},n}I.                                          \tag{63.12}
\]
Equation (63.12) is a declaration of represented operators on a
common domain, not a formal sum: the packet must prove that the
sum is closable and that its chosen closure is self-adjoint or is
the closed integrated-energy operator required by the OS
incidence theorem.  Let \(H_n^{\rm OS}\geq0\) be reconstructed
from the same reflected cylinder.

For each graph screen \(R\), let
\[
 B_{n,R}:F_R\longrightarrow
 \operatorname{Dom}(E_n)\cap\operatorname{Dom}(H_n^{\rm OS}) \tag{63.13}
\]
be a finite source synthesis containing the vacuum, regional
one-field sources, and the declared held-out multi-field words.
Define
\[
 \mathbb C_{n,R}^{E\mid H}
 =
 B_{n,R}^*(E_n-H_n^{\rm OS})^*
 (E_n-H_n^{\rm OS})B_{n,R}\succeq0                           \tag{63.14}
\]
and
\[
 \varepsilon_{n,R}
 =\|\mathbb C_{n,R}^{E\mid H}\|^{1/2}
 =\|(E_n-H_n^{\rm OS})B_{n,R}\|.                             \tag{63.15}
\]

\subsection{Moving-core incidence theorem}

Let
\(J_n:\mathcal K_n^{\rm OS}\to\mathcal K_\infty^{\rm OS}\)
be the comparison isometries.

\begin{theorem}[Cofinal moving-core energy incidence]
\label{thm:v13v63-moving-core}
Suppose there are closed operators
\(E_\infty,H_\infty^{\rm OS}\) and finite maps
\[
 B_R:F_R\longrightarrow
 \operatorname{Dom}(E_\infty)
 \cap\operatorname{Dom}(H_\infty^{\rm OS})                   \tag{63.16}
\]
such that, for every fixed \(R\),
\[
\begin{split}
 J_nB_{n,R}&\longrightarrow B_R,\\
 J_nE_nB_{n,R}&\longrightarrow E_\infty B_R,\\
 J_nH_n^{\rm OS}B_{n,R}
 &\longrightarrow H_\infty^{\rm OS}B_R                      \tag{63.17}
\end{split}
\]
in operator norm on the finite coefficient space \(F_R\), and
\[
 \varepsilon_{n,R}\longrightarrow0.                          \tag{63.18}
\]
Then
\[
 E_\infty B_R=H_\infty^{\rm OS}B_R
 \qquad\text{for every }R.                                   \tag{63.19}
\]
If
\[
 \mathcal D_{\rm src}
 :=\bigcup_R\operatorname{Ran}B_R                             \tag{63.20}
\]
is a common core for both limit operators, then
\[
 \boxed{E_\infty=H_\infty^{\rm OS}.}                         \tag{63.21}
\]

\end{theorem}

\begin{proof}
For \(x\in F_R\),
\[
 \|J_n(E_n-H_n^{\rm OS})B_{n,R}x\|
 \leq\varepsilon_{n,R}\|x\|.                                 \tag{63.22}
\]
By (63.17), the left side converges to
\(\|(E_\infty-H_\infty^{\rm OS})B_Rx\|\).  Equation (63.18)
therefore proves (63.19).  Equality on (63.20) identifies the two
restrictions.  Since (63.20) is a core for both closed operators,
both are closures of that common restriction, proving (63.21).
\end{proof}

The graph convergence in (63.17) is stronger than convergence of
the vectors \(B_{n,R}\) alone.  It prevents a moving-domain defect
from vanishing on every finite table while surviving in the
operator closure.

\begin{countertheorem}[Vector convergence without graph
convergence is insufficient]
\label{cth:v13v63-graph}
On \(\ell^2(\mathbb N)\), let
\[
 A_ne_n=ne_n,\qquad A_ne_j=0\quad(j\ne n).                   \tag{63.23}
\]
For every fixed finitely supported vector \(x\),
\(A_nx=0\) eventually, while for
\[
 x=\sum_{j\geq1}j^{-1}e_j\in\ell^2                           \tag{63.24}
\]
one has \(\|A_nx\|=1\) for every \(n\).  Hence agreement on each
fixed finite source list does not control a moving graph tail or
define convergence on the completed domain.

\end{countertheorem}

\begin{proof}
The support of a fixed finite vector omits \(e_n\) for large
\(n\), proving the first statement.  Equation (63.24) gives
\(A_nx=e_n\).
\end{proof}

\subsection{A diagonal quantitative version}

Suppose the fixed-screen graph errors and incidence errors have
explicit tails
\[
\begin{split}
 \|J_nE_nB_{n,R}-E_\infty B_R\|
 &\leq u_{n,R}^{E},\\
 \|J_nH_n^{\rm OS}B_{n,R}-H_\infty^{\rm OS}B_R\|
 &\leq u_{n,R}^{H},\\
 \varepsilon_{n,R}&\leq u_{n,R}^{\rm inc},                  \tag{63.25}
\end{split}
\]
each tending to zero at fixed \(R\).  Let a cross-screen graph
tail \(\omega_R^{\rm gr}\downarrow0\) control the chosen common
core exhaustion.  A diagonal \(R(n)\uparrow\infty\) is admissible
when
\[
 u_{n,R(n)}^E+u_{n,R(n)}^H+
 u_{n,R(n)}^{\rm inc}+\omega_{R(n)}^{\rm gr}
 \longrightarrow0.                                           \tag{63.26}
\]
The slowly increasing diagonal construction of V57 applies
verbatim.  It yields incidence on the core without interchanging
two uncontrolled limits.

\subsection{The vacuum and multi-grade acquisition rows}

The source atlas in (63.13) must contain:
\begin{enumerate}[label=\textup{(\roman*)}]
\item the vacuum row, which fixes \(c_{{\rm vac},n}\) and removes
the scalar ambiguity invisible to commutators;
\item a source-complete one-particle row, which checks the
clock-scaled Dirac occurrence lift;
\item at least the held-out two-particle rows needed to detect
(63.10), followed by an exhausting finite-grade atlas if the
full Fock energy is claimed;
\item the direct interaction, counterterm, boundary, and
improvement sources in (63.12); and
\item graph-tail and adjacent-cutoff transport data sufficient
for (63.17)--(63.26).
\end{enumerate}

\begin{decision}[Next bottleneck after V63]
\label{dec:v13v63-next}
Once the moving-core residual (63.14) is populated and vanishes,
the integrated energy candidate and the reconstructed OS
Hamiltonian coincide.  The next Lorentzian gate is not another
energy comparison.  It is the continuation of the Euclidean
regional net and translation kernel to a real-time covariant net.
The next note shall isolate a source-complete strip-analytic
criterion, prove uniqueness of the continuation from a
determining Euclidean set, and state the additional locality
boundary-value condition needed for microcausality.
\end{decision}

\begin{assessment}[Progress after V63]
\label{ass:v13v63}
The Hamiltonian incidence problem is now type correct and
cofinal.  The clock-scaled Dirac operator has a rigorous
quasi-free Fock lift, a two-particle counterexample proves that
one-particle data cannot determine the interacting energy, and a
moving-core theorem passes the complete finite incidence Gram to
equality with the OS Hamiltonian under explicit graph convergence.
\end{assessment}

\begin{firewall}[Claim boundary after V63]
\label{fire:v13v63}
No physical interacting energy sum (63.12), common graph atlas, or
incidence tail has been populated.  The quasi-free second
quantization is a reference construction, not the interacting
Standard Model Hamiltonian.  No real-time local net, forward-cone
spectrum, Wightman theory, scattering theory, Einstein
backreaction, or full SM--Einstein continuum is proved.
\end{firewall}

\subsection{Parent record}

This V63 note appends the preceding material to the validated V62
source, whose record was
\[
\begin{array}{c|c}
 \text{lines}&20091\\
 \text{bytes}&706672\\
 \text{SHA--256}&
 \texttt{56AF2B53008B8A45C6765D4C60362A560CD90C1764409BA6C27A496FB03AA380}.
\end{array}                                                     \tag{63.27}
\]

% =============================================================================
% END APPENDED THIRTEENTH-CYCLE V63 MATERIAL
% =============================================================================
% =============================================================================
% BEGIN APPENDED THIRTEENTH-CYCLE V64 MATERIAL
% =============================================================================

\section{Thirteenth-cycle V64: spectral Wick rotation on the
source-cyclic sector}
\label{sec:v13v64}

\subsection{Purpose and separation of gates}

Section~\ref{sec:v13v63} isolated a type-correct route from a
declared physical energy operator to the OS Hamiltonian.  The next
question is what that identification actually supplies in real
time.  The answer is precise but narrower than a Wightman
reconstruction: a joint positive energy--momentum measure gives a
unique source-cyclic unitary translation representation and an
analytic Wick rotation of every measured source kernel.  The
forward-cone condition is an additional support test.  Local
commutativity remains a separate theorem about fields or local
algebras.

The selected-net and source-cyclic translation results already
present in the Grand Tensor manuscript are used as upstream input.
The time-only non-identifiability counterexample proved there is
also retained: momentum support cannot be reconstructed from
\(x=0\) data.  This note does not repeat either result.

\subsection{Finite source matrix measures}

Fix \(q<\infty\), write
\[
 X=[0,\infty)\times\mathbb R^d,
 \qquad k=(E,p),
\]
and let \(\Sigma\) be a finite positive
\(M_q(\mathbb C)\)-valued Borel measure on \(X\).  Thus
\(\Sigma(A)\geq0\) for every Borel set \(A\), and the scalar trace
measure
\[
 \tau(A):=\operatorname{tr}\Sigma(A)                         \tag{64.1}
\]
is finite.  For \(\Re z>0\) and \(x\in\mathbb R^d\), set
\[
 F_\Sigma(z,x)
   :=\int_X e^{-zE+ip\cdot x}\,d\Sigma(E,p).                 \tag{64.2}
\]
All matrix-valued integrals below are understood entrywise.  Since
every scalar entry of \(\Sigma\) has total variation dominated by
\(\tau\), this convention is independent of a basis.

\begin{lemma}[Half-plane analyticity and the real-time boundary]
\label{lem:v13v64-halfplane}
For every \(x\), the function \(z\mapsto F_\Sigma(z,x)\) is
holomorphic on \(\Re z>0\).  It extends continuously in operator
norm to \(\Re z=0\), where
\[
 W_\Sigma(s,x)
   :=F_\Sigma(is,x)
    =\int_X e^{-isE+ip\cdot x}\,d\Sigma(E,p).                \tag{64.3}
\]
More explicitly,
\[
 \lim_{\epsilon\downarrow0}
 \sup_{(s,x)\in K}
 \bigl\|F_\Sigma(\epsilon+is,x)-W_\Sigma(s,x)\bigr\|=0       \tag{64.4}
\]
for every compact \(K\subset\mathbb R\times\mathbb R^d\).
If, for some \(\eta>0\),
\[
 \int_X e^{\eta E}\,d\tau(E,p)<\infty,                       \tag{64.5}
\]
then (64.2) extends holomorphically to \(\Re z>-\eta\).

\begin{proof}
On a compact subset of \(\Re z>0\), choose
\(\delta>0\) with \(\Re z\geq\delta\).  For every integer
\(r\geq0\),
\[
 \bigl|(-E)^r e^{-zE+ip\cdot x}\bigr|
 \leq \sup_{u\geq0}u^r e^{-\delta u}<\infty.                \tag{64.6}
\]
Differentiation under each scalar measure entry is therefore
justified by dominated convergence.  This proves holomorphy.

At the boundary, the difference of the two integrands has modulus
at most \(2\) and tends pointwise to zero.  Entrywise dominated
convergence gives norm convergence because \(q\) is finite.  To
make it uniform on a compact \(K\), first choose a compact subset
\(L\subset X\) with \(\tau(X\setminus L)\) arbitrarily small.  On
\([0,\epsilon_0]\times K\times L\) the exponential is uniformly
continuous, while the complement contributes at most a fixed
multiple of \(\tau(X\setminus L)\).  This proves (64.4).

If \(\Re z>-\eta\), every compact subdomain lies in
\(\Re z\geq-\eta+\delta\) for some \(\delta>0\).  The derivative
integrands are then bounded by
\[
 E^r e^{(\eta-\delta)E}
 \leq C_{r,\delta}e^{\eta E},                                \tag{64.7}
\]
which is \(\tau\)-integrable by (64.5).  The same differentiation
argument proves the final assertion.
\end{proof}
\end{lemma}

The exponential moment in (64.5) is a genuine extra estimate.  A
finite measure always yields the closed right-half-plane boundary,
but it need not permit continuation to any negative real part.

\subsection{The minimal source-cyclic translation representation}

\begin{theorem}[Matrix-measure reconstruction]
\label{thm:v13v64-reconstruction}
Associated with \(\Sigma\) there are a Hilbert space
\(\mathcal H_\Sigma\), strongly commuting self-adjoint operators
\[
 H_\Sigma\geq0,\qquad
 P_{\Sigma,1},\ldots,P_{\Sigma,d},                            \tag{64.8}
\]
and a bounded source map \(J_\Sigma:\mathbb C^q\to\mathcal
H_\Sigma\) such that
\[
 F_\Sigma(t,x)
 =J_\Sigma^*e^{-tH_\Sigma}e^{ix\cdot P_\Sigma}J_\Sigma,
 \qquad t>0,                                                  \tag{64.9}
\]
and
\[
 W_\Sigma(s,x)
 =J_\Sigma^*U_\Sigma(s,x)J_\Sigma,\qquad
 U_\Sigma(s,x):=e^{-isH_\Sigma}e^{ix\cdot P_\Sigma}.          \tag{64.10}
\]
The closed span of
\[
 \{f(H_\Sigma,P_\Sigma)J_\Sigma c:
      f\in C_c(X),\ c\in\mathbb C^q\}                         \tag{64.11}
\]
is all of \(\mathcal H_\Sigma\).  Among representations with this
cyclicity property, the quadruple
\((\mathcal H_\Sigma,H_\Sigma,P_\Sigma,J_\Sigma)\) is unique up
to a unitary intertwiner.

\begin{proof}
Start with \(C_c(X;\mathbb C^q)\) and the positive semidefinite
form
\[
 \langle f,g\rangle_\Sigma
 :=\sum_{a,b=1}^q\int_X
       \overline{f_a(k)}g_b(k)\,d\Sigma_{ab}(k).              \tag{64.12}
\]
Quotient by its null space and complete.  Multiplication by \(E\)
and by each \(p_j\) gives strongly commuting self-adjoint
multiplication operators on their maximal domains.  Constant
vectors belong to the completion because \(\tau(X)<\infty\);
their equivalence classes define \(J_\Sigma\).  The joint
functional calculus gives (64.9) and (64.10), and the construction
gives (64.11).

For uniqueness, map
\(f(H_\Sigma,P_\Sigma)J_\Sigma c\) to the corresponding vector in
a second cyclic realization.  Equality of the matrix spectral
measure makes this map isometric on a dense subspace.  It
therefore extends to a unitary, and the functional calculus shows
that it intertwines \(H\), \(P\), and \(J\).
\end{proof}
\end{theorem}

Thus the boundary in Lemma~\ref{lem:v13v64-halfplane} is not a
formal substitution.  It is a matrix coefficient of a strongly
continuous unitary representation of real spacetime translations.

\subsection{What complete Euclidean source data determine}

\begin{theorem}[Uniqueness from an open Euclidean-time interval]
\label{thm:v13v64-uniqueness}
Let \(\Sigma_1,\Sigma_2\) be finite positive
\(M_q(\mathbb C)\)-valued measures on \(X\).  Suppose that for a
nonempty open interval \(I\subset(0,\infty)\),
\[
 F_{\Sigma_1}(t,x)=F_{\Sigma_2}(t,x)
 \quad\text{for every }t\in I,\ x\in\mathbb R^d.             \tag{64.13}
\]
Then \(\Sigma_1=\Sigma_2\).  Consequently their minimal
source-cyclic real-time representations are unitarily equivalent.

\begin{proof}
It suffices to prove equality of every scalar matrix entry.  Let
\(\mu\) be an entry of \(\Sigma_1-\Sigma_2\).  For fixed \(x\),
the Laplace--Fourier transform
\[
 z\longmapsto\int_Xe^{-zE+ip\cdot x}\,d\mu(E,p)              \tag{64.14}
\]
is holomorphic on the right half-plane and vanishes on \(I\).
The identity theorem makes it zero for every \(t>0\).

Fix such a \(t\).  Uniqueness of the Fourier transform of finite
complex Borel measures on \(\mathbb R^d\) shows that
\[
 \int_{[0,\infty)\times B}e^{-tE}\,d\mu(E,p)=0               \tag{64.15}
\]
for every Borel \(B\subset\mathbb R^d\).  For fixed \(B\), the
left side is the Laplace transform of the finite complex measure
\(A\mapsto\mu(A\times B)\).  Uniqueness of the Laplace transform
on \([0,\infty)\) makes that measure zero.  Hence \(\mu\)
vanishes on all measurable rectangles, and therefore vanishes
identically.  Applying this to every entry proves the claim.  The
last assertion follows from
Theorem~\ref{thm:v13v64-reconstruction}.
\end{proof}
\end{theorem}

Both the continuum in \(x\) and an open set of Euclidean times
matter.  Finite samples form a truncated moment problem and do
not meet the hypotheses of
Theorem~\ref{thm:v13v64-uniqueness}.  Likewise, the restriction
\(x=0\) integrates out the momentum variable and cannot certify
the joint spectral support.

\subsection{The forward-cone gate}

Let
\[
 \overline V_+
 :=\{(E,p)\in[0,\infty)\times\mathbb R^d:E\geq |p|\}.         \tag{64.16}
\]

\begin{proposition}[Support is exactly the spectral-condition test]
\label{prop:v13v64-cone}
For the minimal representation in
Theorem~\ref{thm:v13v64-reconstruction}, the following are
equivalent:
\[
 \Sigma(X\setminus\overline V_+)=0,                           \tag{64.17}
\]
\[
 {\bf1}_{X\setminus\overline V_+}
       (H_\Sigma,P_\Sigma)=0,                                \tag{64.18}
\]
and
\[
 \operatorname{sp}(H_\Sigma,P_\Sigma)
       \subset\overline V_+.                                 \tag{64.19}
\]
In particular, positivity \(H_\Sigma\geq0\) alone does not imply
the relativistic spectrum condition.

\begin{proof}
In the multiplication realization, the projection in (64.18) is
multiplication by the indicator of
\(X\setminus\overline V_+\).  If (64.17) holds, this
multiplication operator is zero.  Conversely, if it is zero,
then its compression to the constant source vectors is zero, so
\[
 c^*\Sigma(X\setminus\overline V_+)c=0
 \quad\text{for all }c\in\mathbb C^q,
\]
which gives (64.17) by polarization.  Equivalence with (64.19)
is the joint spectral theorem.

For the final statement, the scalar point mass
\(\Sigma=\delta_{(1,2e_1)}\) has nonnegative energy and satisfies
all conclusions of Lemma~\ref{lem:v13v64-halfplane}, but its
support lies outside \(\overline V_+\).
\end{proof}
\end{proposition}

The measurable acquisition row for this gate is therefore
\[
 \varepsilon^{\rm cone}_{n,R}
 :=\operatorname{tr}\Sigma_{n,R}
       (X\setminus\overline V_+),                             \tag{64.20}
\]
or, with a tolerance \(\delta_n\downarrow0\),
\[
 \varepsilon^{\rm cone,\delta}_{n,R}
 :=\operatorname{tr}\Sigma_{n,R}
   \{(E,p):E+\delta_n<|p|\}.                                  \tag{64.21}
\]
A continuum spectrum claim requires this row on a
source-complete exhaustion and a tightness argument preventing
mass from escaping the measured windows.

\subsection{Handoff from the V63 incidence theorem}

\begin{theorem}[Conditional OS-to-Lorentzian source-kernel handoff]
\label{thm:v13v64-handoff}
Assume the moving-core hypotheses of
Theorem~\ref{thm:v13v63-moving-core}, so that the declared
continuum physical energy satisfies
\[
 E_\infty=H_{\rm OS}.                                        \tag{64.22}
\]
Assume also that the reconstructed spatial translations have
strongly commuting generators \(P_1,\ldots,P_d\), strongly
commute with \(H_{\rm OS}\), and that a finite source map
\(J:\mathbb C^q\to\mathcal H_{\rm OS}\) is cyclic for their joint
functional calculus.  Let
\[
 \Sigma(A):=J^*{\bf1}_A(H_{\rm OS},P)J.                       \tag{64.23}
\]
Then:

\begin{enumerate}
\item the full Euclidean source kernel
\[
 F(t,x)=J^*e^{-tE_\infty}e^{ix\cdot P}J                      \tag{64.24}
\]
has the unique real-time boundary
\[
 W(s,x)=J^*e^{-isE_\infty}e^{ix\cdot P}J;                    \tag{64.25}
\]
\item complete knowledge of (64.24) for \(x\in\mathbb R^d\) and
\(t\) in any open interval determines the source-cyclic
translation representation up to unitary equivalence;
\item if and only if the joint source measure obeys (64.17), the
source-cyclic sector satisfies the forward-cone spectrum
condition.
\end{enumerate}

\begin{proof}
The joint spectral theorem gives the finite positive matrix
measure (64.23) and represents (64.24) by (64.2).  Item 1 is
Lemma~\ref{lem:v13v64-halfplane}; item 2 is
Theorems~\ref{thm:v13v64-uniqueness} and
\ref{thm:v13v64-reconstruction}; item 3 is
Proposition~\ref{prop:v13v64-cone}.  Cyclicity ensures that no
orthogonal spectral sector is invisible to the stated source
family.
\end{proof}
\end{theorem}

This theorem identifies the next noncircular order of work.  One
must first populate the V63 energy-incidence hypotheses, then
construct a jointly covariant spatial translation system, then
measure the full \((t,x)\) kernels and the cone leakage
(64.20).  Analytic continuation is automatic only after those
real-variable data and tightness conditions have been certified.

\subsection{Why this is not yet microcausality}

The boundary kernel (64.25) supplies translation matrix
coefficients.  It does not by itself define a Lorentzian local
algebra, an operator-valued field, or spacelike commutators.
Even a source-complete two-point spectral measure does not fix
interacting higher correlation functions.  The next locality
step must therefore use a family of Euclidean ordered products,
a common analytic domain, and a source-completeness condition
strong enough to upgrade scalar boundary identities to operator
commutators.

\begin{assessment}[V64 advance]
\label{assess:v13v64}
The Euclidean-to-Lorentzian step has been split into exact gates.
A finite joint positive energy--momentum measure yields
half-plane analyticity, a norm-continuous real-time boundary on
every finite source screen, and a unique minimal source-cyclic
translation representation.  An open Euclidean-time interval
together with all spatial separations uniquely determines that
measure.  The forward-cone spectrum condition is precisely a
support condition and is not a consequence of OS energy
positivity.
\end{assessment}

\begin{firewall}[Claim boundary after V64]
\label{fire:v13v64}
The V63 physical energy sum, moving common core, joint spatial
translations, continuum spectral measure, cone leakage, and
source completeness have not been populated.  No exponential
energy moment, Lorentzian local net, microcausality, higher
Wightman hierarchy, scattering theory, Einstein backreaction, or
full SM--Einstein continuum is proved.
\end{firewall}

\subsection{Parent record}

This V64 note appends the preceding material to the validated V63
source, whose record was
\[
\begin{array}{c|c}
 \text{lines}&20464\\
 \text{bytes}&719961\\
 \text{SHA--256}&
 \texttt{6DEE9B95E05AF40130B2B28FD5A53A62597985F18A731F024AF1B7483B86460B}.
\end{array}                                                     \tag{64.26}
\]

% =============================================================================
% END APPENDED THIRTEENTH-CYCLE V64 MATERIAL
% =============================================================================

% =============================================================================
% BEGIN APPENDED THIRTEENTH-CYCLE V65 MATERIAL
% =============================================================================

\section{Thirteenth-cycle V65: compulsory YM/NS audit before
Lorentzian locality}
\label{sec:v13v65}

\subsection{Files inspected}

The five-version companion audit was performed before continuing
the locality argument.  The newest complete companion sources
available in the shared folder were
\[
\begin{array}{c|c|r|r|c}
 \text{programme}&\text{version}&\text{lines}&\text{bytes}&
 \text{SHA--256}\\ \hline
 \text{Yang--Mills mass gap}&V91&30038&1025165&
 \texttt{E89F4BB2D4C115AC53D8DE15E095AF5B35593B7E1471D16BC55405BA77766E05}\\
 \text{Navier--Stokes stability}&V26&5319&278283&
 \texttt{82C887F8C2C69E4F28FAD7C3DC8DE5B9099CB5A4BF431EBA1FFF3E91935978AD}.
\end{array}                                                     \tag{65.1}
\]
The auxiliary YM parallel note remained at V5, with \(1351\)
lines, \(54174\) bytes, and SHA--256
\[
 \texttt{306F1BB84CFA8A8D47EA569EC17E9545F9867C8D32B548EC9D9C444B4F3C0512}.
                                                                  \tag{65.2}
\]
NS V26 and the auxiliary YM V5 are unchanged from the previous
audit.  They were nevertheless reread at their terminal result
blocks.

\subsection{Transferable conclusions}

YM V90--V91 adds a coefficient-exact quarter-turn grid of dressed
vertex screens.  Its most relevant logical correction is that
covariance of a final scalar response does not authorize an orbit
quotient of the raw Hessian and vertex matrices: points may be
omitted only after the required unitary raw-matrix intertwiner has
been proved.  It also stresses that a finite exact grid controls no
open neighborhood.

NS V26 calculates a rank-one restriction of Oseen-kernel
derivative norms.  At first derivative order the restriction loses
nothing, while at second order it can miss a mixed tensor channel.
The numerical constants do not transfer to the SM programme, but
the structural warning does: a restricted diagonal source class
may fail to detect mixed operator directions.

\begin{proposition}[Audit decision for the locality bridge]
\label{prop:v13v65-decision}
No numerical YM or NS estimate is imported.  The next SM
locality theorem shall obey all of the following rules.

\begin{enumerate}
\item A covariance or reflection reduction of raw ordered-product
matrices is permitted only when the corresponding unitary or
antiunitary intertwiner is part of the hypotheses.
\item Vanishing at finitely many Euclidean times is not treated as
vanishing on an analytic set.  An actual interval, or an
independently certified uniqueness set, is required.
\item The source screen contains mixed bra--ket pairs.  A diagonal
or rank-one screen may be compressed only after a polarization
identity and linear closure of the source class have been
verified.
\item The theorem separates analytic continuation of bounded
matrix elements from the construction and domain control of
unbounded fields.
\end{enumerate}

\begin{proof}
Rule 1 is the direct raw-versus-scalar distinction established by
the V91 YM audit.  Rule 2 follows both from the V91 finite-grid
firewall and from the truncated-moment limitation recorded after
Theorem~\ref{thm:v13v64-uniqueness}.  For rule 3, suppose that \(B\) is conjugate-linear in its first
slot and linear in its second slot on a complex linear source
space.  Its diagonal values obey the exact polarization identity
\[
 B(\psi,\phi)
 ={1\over4}\sum_{\ell=0}^3 i^{-\ell}
 B(\psi+i^\ell\phi,\psi+i^\ell\phi).                           \tag{65.3}
\]
Thus diagonal data recover mixed data only
when all four sums remain in the measured source class.  NS V26
shows why an analogous closure cannot be presumed for restricted
tensor inputs.  Rule 4 is required because analytic scalar
matrix elements alone do not supply invariant domains for
products of unbounded fields.
\end{proof}
\end{proposition}

\subsection{Adjusted next step}

The next note will therefore use bounded local observables
\(A,B\), a dense complex linear source domain, and a bounded
holomorphic commutator kernel on a connected complex-time
domain.  Source-complete vanishing on a genuine Euclidean
interval will be continued by the identity theorem and then
upgraded to an operator commutator.  The result will be
conditional: establishing the needed Euclidean ordered-product
identity and its common analytic domain remains an acquisition
problem.

\begin{assessment}[V65 audit outcome]
\label{assess:v13v65}
The audit changes the proposed locality bridge in two useful
ways.  It forbids an unjustified symmetry-orbit reduction of raw
source matrices, and it requires mixed source directions plus an
open analytic uniqueness set.  The exact YM grid and the NS
rank-one constants themselves supply no SM continuum constant.
\end{assessment}

\begin{firewall}[Claim boundary after V65]
\label{fire:v13v65}
This is an audit and proof-design step.  No Lorentzian
commutator, local net, Wightman hierarchy, interacting continuum
measure, Einstein equation, or full SM--Einstein continuum is
proved.
\end{firewall}

\subsection{Parent record}

This V65 note appends the preceding material to the validated V64
source, whose record was
\[
\begin{array}{c|c}
 \text{lines}&20854\\
 \text{bytes}&734412\\
 \text{SHA--256}&
 \texttt{E2ACD16E83F2C99BF21578BB522F36CB06E4634EF294DA48D52CF6A6EBBC5671}.
\end{array}                                                     \tag{65.4}
\]

% =============================================================================
% END APPENDED THIRTEENTH-CYCLE V65 MATERIAL
% =============================================================================

% =============================================================================
% BEGIN APPENDED THIRTEENTH-CYCLE V66 MATERIAL
% =============================================================================

\section{Thirteenth-cycle V66: a source-complete analytic
locality bridge}
\label{sec:v13v66}

\subsection{The missing input beyond spectral Wick rotation}

Theorem~\ref{thm:v13v64-handoff} reconstructs real-time
translations on a source-cyclic sector.  To speak about
microcausality one must additionally have bounded local
observables and identify the two boundary orderings of their
Euclidean correlation functions.  This section gives the exact
analytic statement needed for that upgrade.  It is deliberately
formulated as a certification theorem: it proves what follows
from the required ordered-product data without pretending that
the present finite models have already supplied those data.

Let \(U(s,x)=e^{-isH}e^{ix\cdot P}\) be the unitary
representation supplied by the V64 handoff, and suppose that it
acts covariantly on a represented algebra by
\[
 \alpha_{(s,x)}(A):=U(s,x)AU(s,x)^*.                          \tag{66.1}
\]
Fix bounded observables \(A,B\), a spatial displacement \(x\),
and a real-time interval \(S_x\) on which the translated
localization of \(A\) is spacelike to the localization of \(B\).

\subsection{Exact continuation of the ordering difference}

\begin{theorem}[Dense-source ordered-product continuation]
\label{thm:v13v66-continuation}
Let \(\mathcal D_0\) be a dense complex linear subspace of the
physical Hilbert space.  Suppose that there are

\begin{enumerate}
\item a connected open set \(\Omega_x\subset\mathbb C\);
\item a nonempty open interval \(I_x\subset\mathbb R\) with
\(iI_x\subset\Omega_x\);
\item a real interval \(S_x\subset\partial\Omega_x\);
\item for every \(\psi,\phi\in\mathcal D_0\), two scalar
functions
\[
 F^{AB}_{\psi,\phi}(\zeta;x),\qquad
 F^{BA}_{\psi,\phi}(\zeta;x),                                \tag{66.2}
\]
holomorphic on \(\Omega_x\) and continuous on
\(\Omega_x\cup S_x\);
\end{enumerate}
such that the Euclidean ordering identity
\[
 F^{AB}_{\psi,\phi}(it;x)
   =F^{BA}_{\psi,\phi}(it;x)
 \quad(t\in I_x)                                              \tag{66.3}
\]
holds for all \(\psi,\phi\in\mathcal D_0\), and the real boundary
difference is the commutator matrix element
\[
 F^{AB}_{\psi,\phi}(s;x)-F^{BA}_{\psi,\phi}(s;x)
 =\langle\psi,[\alpha_{(s,x)}(A),B]\phi\rangle
 \quad(s\in S_x).                                             \tag{66.4}
\]
Then
\[
 [\alpha_{(s,x)}(A),B]=0
 \qquad\text{for every }s\in S_x.                             \tag{66.5}
\]

\begin{proof}
For fixed \(\psi,\phi\in\mathcal D_0\), put
\[
 C_{\psi,\phi}(\zeta;x)
 :=F^{AB}_{\psi,\phi}(\zeta;x)
   -F^{BA}_{\psi,\phi}(\zeta;x).                              \tag{66.6}
\]
This function is holomorphic on the connected set \(\Omega_x\)
and vanishes on the line segment \(iI_x\).  Every point of that
segment is an accumulation point lying inside \(\Omega_x\).
The identity theorem therefore gives
\(C_{\psi,\phi}=0\) throughout \(\Omega_x\).  Continuity to
\(S_x\) and (66.4) imply
\[
 \langle\psi,[\alpha_{(s,x)}(A),B]\phi\rangle=0
 \quad(\psi,\phi\in\mathcal D_0).                             \tag{66.7}
\]
For real \(s\), the commutator is bounded and
\[
 \|[\alpha_{(s,x)}(A),B]\|
 \leq2\|A\|\,\|B\|.                                          \tag{66.8}
\]
Approximating arbitrary Hilbert-space vectors by vectors in
\(\mathcal D_0\) extends (66.7) to every pair of vectors.  Hence
the bounded commutator is zero.
\end{proof}
\end{theorem}

The theorem requires equality of the two ordered products, not
merely reflection symmetry of one scalar trace.  A reflection or
coordinate symmetry can establish (66.3) only when its
unitary/antiunitary action on \(A,B,\psi,\phi\) has been
certified.  This is the raw-object intertwiner discipline imported
at the V65 audit.

\subsection{Two sharp obstructions to weaker screens}

\begin{proposition}[A non-dense source screen cannot certify
locality]
\label{prop:v13v66-dark-sector}
Let \(\mathcal K\) be a proper closed subspace of a Hilbert space
\(\mathcal H\).  There exist bounded self-adjoint \(A,B\) such
that
\[
 \langle\psi,[A,B]\phi\rangle=0
 \quad(\psi,\phi\in\mathcal K),                               \tag{66.9}
\]
but \([A,B]\neq0\).

\begin{proof}
Take
\(\mathcal H=\mathcal K\oplus\mathbb C^2\).  Let both operators
vanish on \(\mathcal K\), and on \(\mathbb C^2\) take
\[
 A=\begin{pmatrix}0&1\\1&0\end{pmatrix},
 \qquad
 B=\begin{pmatrix}1&0\\0&-1\end{pmatrix}.                    \tag{66.10}
\]
Then (66.9) holds, whereas their commutator on the second summand
is
\[
 [A,B]=
 \begin{pmatrix}0&-2\\2&0\end{pmatrix}\neq0.                 \tag{66.11}
\]
\end{proof}
\end{proposition}

Thus cyclicity for the translation spectral algebra in V64 does
not automatically imply completeness for the enlarged algebra
generated by local observables.  The latter density must be
proved separately.

\begin{proposition}[Finite Euclidean samples are not an analytic
uniqueness set]
\label{prop:v13v66-finite-samples}
Given distinct points \(\zeta_1,\ldots,\zeta_N\in\mathbb C\), a
point \(s_0\) different from them, and a nonzero bounded operator
\(T\), there is an operator-valued entire function \(C\) such
that
\[
 C(\zeta_j)=0\quad(1\leq j\leq N),
 \qquad C(s_0)=T.                                             \tag{66.12}
\]

\begin{proof}
The polynomial
\[
 p(\zeta):=
 {\prod_{j=1}^N(\zeta-\zeta_j)\over
  \prod_{j=1}^N(s_0-\zeta_j)}                                \tag{66.13}
\]
satisfies \(p(\zeta_j)=0\) and \(p(s_0)=1\).  Set
\(C(\zeta)=p(\zeta)T\).
\end{proof}
\end{proposition}

Exact vanishing on a finite time grid therefore cannot replace
(66.3).  There is also no stable boundary estimate from small
values on one compact interior set alone.  On the unit disc,
\[
 f_N(z)=z^N,\qquad
 \sup_{|z|\leq r}|f_N(z)|\leq r^N\longrightarrow0
 \quad(0<r<1),                                                \tag{66.14}
\]
while \(f_N(1)=1\).  Quantitative continuation needs a uniform
analytic bound together with control on a set of positive
harmonic measure, or a comparably strong three-lines input.  This
is why the eventual continuum acquisition table must retain
analytic-domain bounds rather than only grid residuals.

\subsection{Upgrade to a local \(C^*\)-net}

\begin{corollary}[Locality from a dense analytic generating
family]
\label{cor:v13v66-local-net}
Suppose a covariant net
\(\mathcal O\mapsto\mathfrak A(\mathcal O)\subset B(\mathcal H)\)
has, for every region \(\mathcal O\), a norm-dense
\(*\)-subalgebra
\(\mathfrak A_0(\mathcal O)\).  Suppose that whenever
\(\mathcal O_1\) and \(\mathcal O_2\) are spacelike separated,
every
\[
 A\in\mathfrak A_0(\mathcal O_1),\qquad
 B\in\mathfrak A_0(\mathcal O_2)                              \tag{66.15}
\]
satisfies the hypotheses of
Theorem~\ref{thm:v13v66-continuation} for a dense source space
\(\mathcal D_0\), with the boundary point representing the given
pair of localizations.  Then
\[
 [\mathfrak A(\mathcal O_1),\mathfrak A(\mathcal O_2)]=0
 \quad\text{whenever }\mathcal O_1\perp\mathcal O_2.          \tag{66.16}
\]

\begin{proof}
The theorem first gives \([A,B]=0\) for the dense generating
subalgebras.  Choose norm-convergent sequences
\(A_n\to A\) and \(B_n\to B\) from those subalgebras.  Then
\[
\begin{split}
 \|[A,B]-[A_n,B_n]\|
 &\leq 2\|A-A_n\|\,\|B\|\\
 &\quad+2\|A_n\|\,\|B-B_n\|\longrightarrow0.                 \tag{66.17}
\end{split}
\]
Thus the commutator vanishes for the norm closures.
\end{proof}
\end{corollary}

This corollary closes the logical step from certified ordered
Euclidean products to Haag--Kastler locality for bounded
observables.  It does not create the covariant net appearing in
its hypotheses.  In particular, the V64 spectral representation
must still be shown to implement automorphisms of the selected
continuum algebra.

\subsection{A finite-to-continuum acquisition card}

For each localization pair, source screen \(R\), and lattice
index \(n\), the data must distinguish the following residuals:
\[
\begin{array}{c|l}
 \varepsilon^{\rm ord}_{n,R}
 &\displaystyle
  \sup_{t\in I_x}
  \|F^{AB}_{n,R}(it;x)-F^{BA}_{n,R}(it;x)\|\\
 \varepsilon^{\rm src}_{n,R}
 &\text{graph or Hilbert-space tail outside the mixed source span}\\
 \varepsilon^{\rm int}_{n,R}
 &\text{residual in the claimed covariance/reflection intertwiner}\\
 M^{\rm an}_{n,R}
 &\displaystyle
  \sup_{\zeta\in\Omega'_x}
  \bigl(\|F^{AB}_{n,R}(\zeta;x)\|
       +\|F^{BA}_{n,R}(\zeta;x)\|\bigr)
\end{array}                                                     \tag{66.18}
\]
for compact subdomains \(\Omega'_x\Subset\Omega_x\).  The first
row must be controlled on a continuum uniqueness set or by a
proved analytic interpolation rule; sampling it at finitely many
times is insufficient by
Proposition~\ref{prop:v13v66-finite-samples}.  The second row
must vanish on a cofinal mixed source exhaustion by
Proposition~\ref{prop:v13v66-dark-sector}.  The third row is
required before any symmetry orbit is removed.  The fourth row
provides normal-family compactness and is the starting point for
a quantitative boundary argument.

\begin{theorem}[Compact-normal-family version]
\label{thm:v13v66-normal-family}
Let \(F^{AB}_{n,\psi,\phi}\) and
\(F^{BA}_{n,\psi,\phi}\) be holomorphic on a connected
\(\Omega_x\), locally uniformly bounded there, and convergent
pointwise on a subset with an accumulation point in
\(\Omega_x\).  Suppose their limits agree on the Euclidean
interval \(iI_x\).  Then every locally uniform subsequential
limit has zero ordering difference throughout \(\Omega_x\).
If, in addition, its boundary values on \(S_x\) converge to the
bounded real-time commutator matrix elements and the limiting
source space is dense, then (66.5) holds.

\begin{proof}
Montel's theorem gives locally uniformly convergent subsequences
of each scalar matrix element.  The identity theorem makes the
difference of the two limiting holomorphic functions zero,
because it vanishes on \(iI_x\).  The stated boundary convergence
then yields (66.7), and the density argument in
Theorem~\ref{thm:v13v66-continuation} gives (66.5).
\end{proof}
\end{theorem}

The boundary-convergence assumption in the last theorem is not a
consequence of interior Montel convergence: example (66.14)
shows the gap.  A subsequent note must obtain a uniform
strip/tube estimate that reaches the desired spacelike boundary,
or weaken the conclusion to distributional boundary values and
prove the corresponding edge-of-the-wedge hypotheses.

\begin{assessment}[V66 advance]
\label{assess:v13v66}
A rigorous conditional microcausality bridge is now available.
Equality of the two Euclidean orderings on a genuine interval,
inside one connected analytic domain, propagates to the
spacelike real boundary.  Dense mixed sources upgrade scalar
identities to a bounded operator commutator, and norm-dense local
generators upgrade the result to a local \(C^*\)-net.  Explicit
counterexamples show why finite time samples and non-dense source
screens do not suffice.
\end{assessment}

\begin{firewall}[Claim boundary after V66]
\label{fire:v13v66}
The ordered Euclidean products, common analytic domains,
intertwiner residuals, normal-family bounds, boundary
convergence, covariant continuum net, and dense local generating
families in the theorems above have not been populated.  No
unbounded-field locality, Wightman hierarchy, interacting
continuum measure, scattering theory, Einstein backreaction, or
full SM--Einstein continuum is proved.
\end{firewall}

\subsection{Parent record}

This V66 note appends the preceding material to the validated V65
source, whose record was
\[
\begin{array}{c|c}
 \text{lines}&20992\\
 \text{bytes}&740165\\
 \text{SHA--256}&
 \texttt{2720AF3B74E08BDD0C0A63A57BA70A507D504B53C6FE702E193F68A8F0BFEBEB}.
\end{array}                                                     \tag{66.19}
\]

% =============================================================================
% END APPENDED THIRTEENTH-CYCLE V66 MATERIAL
% =============================================================================

% =============================================================================
% BEGIN APPENDED THIRTEENTH-CYCLE V67 MATERIAL
% =============================================================================

\section{Thirteenth-cycle V67: analytic-collar stability and
cofinal locality}
\label{sec:v13v67}

\subsection{Why an analytic collar is the next upstream gate}

The exact identity theorem in V66 is qualitative.  Finite
approximants instead provide nonzero ordering residuals.  Example
(66.14) shows that interior smallness need not control a distinct
boundary point.  A quantitative handoff becomes possible once
the desired spacelike points lie inside a common analytic collar.
This section proves the two elementary tools needed for that
handoff: gluing across a certified ordering seam and propagation
by harmonic measure.

\subsection{Weak operator gluing across a seam}

\begin{lemma}[One-variable weak edge gluing]
\label{lem:v13v67-gluing}
Let \(S=(a,b)\subset\mathbb R\).  Let \(\Omega_+\) and
\(\Omega_-\) be open subsets of the upper and lower half-planes,
respectively, each having \(S\) as a boundary interval, and
suppose
\[
 \Omega:=\Omega_+\cup S\cup\Omega_-                           \tag{67.1}
\]
is open.  Let
\[
 F_\pm:\Omega_\pm\longrightarrow B(\mathcal H)
\]
be weakly holomorphic, locally uniformly bounded in operator
norm up to compact subintervals of \(S\), and suppose all scalar
matrix elements have continuous boundary values satisfying
\[
 \langle\psi,F_+(s)\phi\rangle
 =\langle\psi,F_-(s)\phi\rangle
 \quad(s\in S,\ \psi,\phi\in\mathcal H).                     \tag{67.2}
\]
Then there is a unique weakly holomorphic
\(F:\Omega\to B(\mathcal H)\) whose restrictions are \(F_+\) and
\(F_-\).

\begin{proof}
For fixed \(\psi,\phi\), join the two scalar functions using their
common continuous boundary value.  The resulting function is
continuous on \(\Omega\) and holomorphic off \(S\).  For a
triangle compactly contained in \(\Omega\), cut it along the
lines \(\operatorname{Im}z=\pm\epsilon\).  Cauchy's theorem
annuls the integrals on the upper and lower pieces.  The two
integrals adjacent to \(S\) cancel in the limit
\(\epsilon\downarrow0\) by the common continuous boundary value.
Morera's theorem gives a holomorphic scalar function on
\(\Omega\).

The scalar extensions define a weak operator-valued map.  Local
uniform operator bounds and the uniform boundedness principle
make it a \(B(\mathcal H)\)-valued weakly holomorphic function.
Uniqueness follows from the scalar identity theorem.
\end{proof}
\end{lemma}

For ordered products, condition (67.2) must come from a proved
Euclidean permutation/reflection identity on the seam, including
the action on the raw observables and source vectors.  Equality of
a final scalar trace after unrecorded coordinate cancellation
does not meet the hypothesis.

\subsection{The harmonic-measure propagation bound}

Let \(D\subset\mathbb C\) be a bounded regular domain and
\(\Gamma\subset\partial D\) a Borel boundary set.  Write
\(\omega_D(z,\Gamma)\) for its harmonic measure viewed from
\(z\in D\).

\begin{theorem}[Operator two-constants estimate]
\label{thm:v13v67-two-constants}
Let \(C:D\to B(\mathcal H)\) be weakly holomorphic and satisfy
\[
 \sup_{z\in D}\|C(z)\|\leq M.                                \tag{67.3}
\]
Suppose \(C\) has norm-continuous boundary values on \(\Gamma\)
with
\[
 \sup_{\zeta\in\Gamma}\|C(\zeta)\|\leq\epsilon\leq M.         \tag{67.4}
\]
Then
\[
 \boxed{\quad
 \|C(z)\|
 \leq M^{\,1-\omega_D(z,\Gamma)}
       \epsilon^{\,\omega_D(z,\Gamma)}
 \quad(z\in D).
 \quad}                                                       \tag{67.5}
\]
If \(K\Subset D\) and
\[
 \theta_K:=\inf_{z\in K}\omega_D(z,\Gamma)>0,                 \tag{67.6}
\]
then
\[
 \sup_{z\in K}\|C(z)\|
 \leq M^{1-\theta_K}\epsilon^{\theta_K}.                      \tag{67.7}
\]

\begin{proof}
Fix unit vectors \(\psi,\phi\) and put
\(f(z)=\langle\psi,C(z)\phi\rangle\).  For \(\delta>0\), the
subharmonic function
\[
 u_\delta(z):=\log(|f(z)|+\delta)
\]
has boundary upper bound \(\log(M+\delta)\) off \(\Gamma\) and
\(\log(\epsilon+\delta)\) on \(\Gamma\).  The two-constants
theorem for subharmonic functions gives
\[
 u_\delta(z)\leq
 (1-\omega_D(z,\Gamma))\log(M+\delta)
 +\omega_D(z,\Gamma)\log(\epsilon+\delta).                   \tag{67.8}
\]
Letting \(\delta\downarrow0\) proves the scalar form of (67.5),
including the case \(\epsilon=0\).  Taking the supremum over unit
\(\psi,\phi\) proves the operator norm bound.  Equation (67.7)
follows from (67.5), (67.6), and \(\epsilon\leq M\).
\end{proof}
\end{theorem}

If a Euclidean interval lies in the interior of a larger
holomorphy domain, one may slit the domain along a compact
subinterval and use the two sides of the slit as \(\Gamma\).
The desired spacelike compact set \(K\) must remain in the same
regular component.  The positive number \(\theta_K\) is then a
geometric conditioning coefficient for analytic continuation.

\subsection{Cofinal source screens}

The harmonic estimate may first be available only after
compression to a finite source screen.  The following result
shows exactly how to pass from those compressions to an operator
identity.

\begin{theorem}[Cofinal compressed locality criterion]
\label{thm:v13v67-cofinal}
Let \(P_R\) be orthogonal projections on \(\mathcal H\) with
\[
 P_R\longrightarrow I\quad\text{strongly}.                   \tag{67.9}
\]
For each \(R\), let \(C_{n,R}:D_R\to P_RB(\mathcal H)P_R\) be
weakly holomorphic, where \(D_R\) is a bounded regular domain,
and let \(\Gamma_R\subset\partial D_R\) and
\(s\in K_R\Subset D_R\).  Assume, for each fixed \(R\),

\[
 \sup_n\sup_{z\in D_R}\|C_{n,R}(z)\|\leq M_R<\infty,          \tag{67.10}
\]
\[
 \epsilon_{n,R}:=
 \sup_{\zeta\in\Gamma_R}\|C_{n,R}(\zeta)\|
 \longrightarrow0,                                          \tag{67.11}
\]
and
\[
 \theta_R:=
 \inf_{z\in K_R}\omega_{D_R}(z,\Gamma_R)>0.                  \tag{67.12}
\]
Suppose also that a bounded operator \(C(s)\) satisfies
\[
 C_{n,R}(s)\longrightarrow P_RC(s)P_R
 \quad\text{in operator norm}                                \tag{67.13}
\]
for every fixed \(R\).  Then \(C(s)=0\).

\begin{proof}
Theorem~\ref{thm:v13v67-two-constants} gives
\[
 \|C_{n,R}(s)\|
 \leq M_R^{1-\theta_R}\epsilon_{n,R}^{\theta_R}
 \longrightarrow0                                           \tag{67.14}
\]
at fixed \(R\).  Passing to the limit in (67.13) yields
\(P_RC(s)P_R=0\).  For arbitrary \(\psi,\phi\in\mathcal H\),
strong convergence (67.9) and boundedness of \(C(s)\) give
\[
 \langle\psi,C(s)\phi\rangle
 =\lim_R\langle P_R\psi,C(s)P_R\phi\rangle=0.                \tag{67.15}
\]
Hence \(C(s)=0\).
\end{proof}
\end{theorem}

No uniform lower bound on \(\theta_R\) is needed for this exact
iterated limit: \(n\to\infty\) is taken at fixed \(R\).  A
uniform quantitative estimate before the cofinal limit would
need explicit control of \(M_R\), \(\theta_R\), and the
off-screen tails.

\begin{proposition}[Quantitative off-screen completion]
\label{prop:v13v67-tail}
For any orthogonal projection \(P\) and bounded operator \(C\),
\[
 \|C\|
 \leq\|PCP\|+\|(I-P)C\|+\|C(I-P)\|.                          \tag{67.16}
\]
Consequently, if the last two terms are at most \(\eta\), the
compressed harmonic estimate implies
\[
 \|C(s)\|
 \leq M^{1-\theta}\epsilon^\theta+2\eta.                     \tag{67.17}
\]

\begin{proof}
The identity
\[
 C=PCP+(I-P)C+PC(I-P)                                       \tag{67.18}
\]
and \(\|PC(I-P)\|\leq\|C(I-P)\|\) give (67.16).
Insert (67.7) for the first term.
\end{proof}
\end{proposition}

Strong convergence \(P_R\to I\) does not force either tail norm
in (67.16) to vanish.  For example, on \(\ell^2(\mathbb N)\)
with \(P_R\) the first-\(R\) coordinate projection, the unilateral
shift powers \(C_R=S^R\) have unit norm while
\[
 P_RC_RP_R=0.                                                \tag{67.19}
\]
Thus norm-level locality requires a separately measured tail,
whereas exact weak/operator locality of a fixed limit follows
from Theorem~\ref{thm:v13v67-cofinal}.

\subsection{Application to the ordering residual}

Let
\[
 C_{n,R}(\zeta;x)
 :=F^{AB}_{n,R}(\zeta;x)-F^{BA}_{n,R}(\zeta;x).              \tag{67.20}
\]
If Lemma~\ref{lem:v13v67-gluing} produces a common collar,
the Euclidean ordering segment can be made a boundary portion of
a slit domain \(D_R\).  The finite acquisition rows then have the
direct meanings
\[
 M_R=\sup_{n,z\in D_R}\|C_{n,R}(z;x)\|,\qquad
 \epsilon_{n,R}=\sup_{\zeta\in\Gamma_R}
                  \|C_{n,R}(\zeta;x)\|,                      \tag{67.21}
\]
\[
 \theta_R=\inf_{z\in K_R}\omega_{D_R}(z,\Gamma_R),\qquad
 \delta^{\rm bdry}_{n,R}
 =\|C_{n,R}(s;x)-P_R[\alpha_{(s,x)}(A),B]P_R\|.              \tag{67.22}
\]
Theorem~\ref{thm:v13v67-cofinal} applies when
\[
 \epsilon_{n,R}\to0,\qquad
 \delta^{\rm bdry}_{n,R}\to0
 \quad\text{at each fixed }R,                                \tag{67.23}
\]
and the screens are cofinal.  This converts the qualitative V66
requirements into a finite-to-continuum proof obligation with a
visible conditioning exponent.

\begin{assessment}[V67 advance]
\label{assess:v13v67}
The locality bridge now has a quantitative finite-level core.
A certified seam glues the two ordered continuations into one
analytic collar.  Harmonic measure propagates an ordering
residual with the explicit exponent \(\theta_R\), and cofinal
mixed source projections upgrade the vanishing compressed limit
to the full bounded commutator.  A separate tail estimate is
shown to be necessary only for uniform operator-norm error
bounds.
\end{assessment}

\begin{firewall}[Claim boundary after V67]
\label{fire:v13v67}
No SM ordered-product seam, analytic collar, harmonic-measure
coefficient, uniform analytic bound, boundary residual, or
cofinal local source projection has been populated.  The
one-variable gluing lemma is not a proof of the several-variable
Wightman edge-of-the-wedge theorem.  No covariant Lorentzian
local net, unbounded-field locality, interacting continuum
measure, Einstein backreaction, or full SM--Einstein continuum
is proved.
\end{firewall}

\subsection{Parent record}

This V67 note appends the preceding material to the validated V66
source, whose record was
\[
\begin{array}{c|c}
 \text{lines}&21321\\
 \text{bytes}&752362\\
 \text{SHA--256}&
 \texttt{1BDD15D4513FAE14060DCC2E09A3E5DB7DD2BC07CB9828F48806157DB46C391F}.
\end{array}                                                     \tag{67.24}
\]

% =============================================================================
% END APPENDED THIRTEENTH-CYCLE V67 MATERIAL
% =============================================================================

% =============================================================================
% BEGIN APPENDED THIRTEENTH-CYCLE V68 MATERIAL
% =============================================================================

\section{Thirteenth-cycle V68: bounded-observable Wightman
hierarchy after locality}
\label{sec:v13v68}

\subsection{Position relative to the existing stress compiler}

The Grand Tensor manuscript already gives a conditional cofinal
stress/current theorem, an energy-density--Hamiltonian incidence
test, and a selected Euclidean local-net landing theorem.  The
earlier cycles also give separate sufficient limit conditions for
mean Einstein equations, Schwinger--Dyson identities, and
unbounded stress/Einstein writers.  Those results are not
repeated here.

V64--V67 add the missing Lorentzian direction: joint spectral
Wick rotation, an ordered-product locality criterion, and a
quantitative cofinal continuation mechanism.  Once those
hypotheses have actually produced a positive-energy covariant
local net, one can extract a complete hierarchy of bounded local
correlation distributions.  The present note records that
implication and identifies why a finite correlation order is
insufficient.

\subsection{Net hypotheses}

Let
\(\mathcal O\mapsto\mathfrak A(\mathcal O)\subset B(\mathcal H)\)
be an isotone local net on Minkowski space.  Let \(U(a)\) be a
strongly continuous unitary translation representation with
\[
 U(a)\Omega=\Omega,\qquad
 \alpha_a(A)=U(a)AU(a)^*,                                    \tag{68.1}
\]
and suppose
\[
 \operatorname{sp}(H,P)\subset\overline V_+.                 \tag{68.2}
\]
For bounded local observables \(A_1,\ldots,A_n\), define
\[
 W_{A_1,\ldots,A_n}(x_1,\ldots,x_n)
 :=
 \langle\Omega,
 \alpha_{x_1}(A_1)\cdots\alpha_{x_n}(A_n)\Omega\rangle.       \tag{68.3}
\]

\begin{theorem}[Bounded local correlation hierarchy]
\label{thm:v13v68-hierarchy}
Under (68.1)--(68.2), the functions (68.3) have the following
properties.

\begin{enumerate}
\item They are bounded and jointly continuous, hence define
tempered distributions.
\item They are translation invariant:
\[
 W(x_1+a,\ldots,x_n+a)=W(x_1,\ldots,x_n).                    \tag{68.4}
\]
\item They obey Hermitian reversal:
\[
 \overline{W_{A_1,\ldots,A_n}(x_1,\ldots,x_n)}
 =
 W_{A_n^*,\ldots,A_1^*}(x_n,\ldots,x_1).                     \tag{68.5}
\]
\item The whole hierarchy is positive: for every finite family
of translated local words \(X_j\) and coefficients \(c_j\),
\[
 \sum_{i,j}\overline{c_i}c_j
 \langle\Omega,X_i^*X_j\Omega\rangle
 =
 \left\|\sum_jc_jX_j\Omega\right\|^2\geq0.                   \tag{68.6}
\]
\item Adjacent factors may be interchanged whenever their
translated localization regions are spacelike separated.
\item In consecutive difference variables
\[
 y_j=x_{j+1}-x_j,\qquad 1\leq j<n,                            \tag{68.7}
\]
the joint Fourier support is contained in
\[
 \overline V_+^{\,n-1}.                                      \tag{68.8}
\]
\end{enumerate}

If a strongly continuous unitary representation of the proper
orthochronous Poincar\'e group is also present, fixes \(\Omega\),
and acts covariantly on the net and on the chosen observable
multiplets, then the hierarchy has the corresponding Lorentz
covariance.

\begin{proof}
Strong continuity of \(U\) makes every vector-valued map
\(x\mapsto\alpha_x(A)\xi\) continuous.  Successive use of this
fact and the uniform bound \(\|\alpha_x(A)\|=\|A\|\) gives joint
continuity of (68.3), while
\[
 |W_{A_1,\ldots,A_n}|
 \leq\prod_{j=1}^n\|A_j\|.                                   \tag{68.9}
\]
Every bounded continuous function defines a tempered
distribution.  Vacuum invariance gives (68.4), taking adjoints
gives (68.5), and (68.6) is a Gram identity.  Locality of the net
gives item 5.

Using translation invariance to remove \(x_1\) and cancelling
neighboring unitaries gives
\[
 W(y_1,\ldots,y_{n-1})
 =
 \langle\Omega,
 A_1U(y_1)A_2U(y_2)\cdots
 A_{n-1}U(y_{n-1})A_n\Omega\rangle.                           \tag{68.10}
\]
Let \(\mathsf E\) be the joint spectral measure of \((H,P)\).
With the Fourier convention
\[
 \widehat f(E,p)
 :=\int_{\mathbb R^{1+d}}
 f(t,x)e^{-iEt+ip\cdot x}\,dt\,dx,                            \tag{68.11}
\]
one has
\[
 \int f(a)U(a)\,da
 =\int\widehat f(E,p)\,d\mathsf E(E,p).                       \tag{68.12}
\]
If \(\widehat f\) is supported off \(\overline V_+\), the last
operator is zero by (68.2).  Smearing any one variable of
(68.10) with such an \(f\) therefore gives zero.  Finite sums of
factorized Schwartz functions are dense in the Schwartz space
on the product, so the distributional support is contained in
(68.8).  The final covariance statement follows by inserting
the Poincar\'e representation and using invariance of \(\Omega\).
\end{proof}
\end{theorem}

This is a bounded-observable analogue of the corresponding
Wightman properties.  It avoids all domain questions because the
insertions lie in local \(C^*\)-algebras.

\subsection{When the hierarchy determines the state}

Let \(\mathfrak A_0(\mathcal O)\) be chosen norm-dense
\(*\)-subalgebras, closed under the recorded translation
operations whenever the translated region remains in the
regional ledger.  Let \(\mathfrak A_{\rm alg}\) be the unital
\(*\)-algebra generated by their translates.

\begin{theorem}[All orders on a dense local ledger determine the
vacuum state]
\label{thm:v13v68-determination}
Suppose \(\mathfrak A_{\rm alg}\) is norm dense in the
quasi-local algebra \(\mathfrak A\).  Knowledge of (68.3) at all
orders, for all ledger observables and all recorded translations,
determines the vacuum state
\[
 \omega(A)=\langle\Omega,A\Omega\rangle                      \tag{68.13}
\]
uniquely on \(\mathfrak A\).

\begin{proof}
Every element of \(\mathfrak A_{\rm alg}\) is a finite linear
combination of translated local words.  Linearity and the
correlation values therefore determine \(\omega\) on
\(\mathfrak A_{\rm alg}\).  States have norm one, so if
\(A_k\in\mathfrak A_{\rm alg}\) converges in norm to
\(A\in\mathfrak A\),
\[
 |\omega(A_k)-\omega(A_\ell)|
 \leq\|A_k-A_\ell\|.                                         \tag{68.14}
\]
The values extend uniquely and continuously to the norm closure.
\end{proof}
\end{theorem}

If the GNS vacuum is cyclic, the state then determines its
cyclic representation up to unitary equivalence.  This
determination statement uses all word lengths; no Gaussian Wick
closure is being assumed.

\subsection{A finite-order hierarchy is not complete}

\begin{proposition}[Matched finite moments with different states]
\label{prop:v13v68-finite-order}
For every \(N\geq1\), there are two distinct probability states
on \(C([-1,1])\) whose moments agree through order \(N\).

\begin{proof}
Let \(P_{N+1}\) be the Legendre polynomial normalized by
\(P_{N+1}(1)=1\).  It satisfies
\[
 |P_{N+1}(x)|\leq1\quad(-1\leq x\leq1),                       \tag{68.15}
\]
and is orthogonal in \(L^2([-1,1],dx)\) to every polynomial of
degree at most \(N\).  For \(0<\epsilon<1\), define
\[
 d\mu_\pm(x)
 :={1\over2}\bigl(1\pm\epsilon P_{N+1}(x)\bigr)\,dx.          \tag{68.16}
\]
These are positive probability measures because
\(\int_{-1}^1P_{N+1}(x)\,dx=0\).  For \(0\leq k\leq N\),
orthogonality gives
\[
 \int x^k\,d\mu_+
 =\int x^k\,d\mu_-.                                          \tag{68.17}
\]
They are distinct, since their difference has nonzero pairing
with \(P_{N+1}\); equivalently their \((N+1)\)-st moments differ
because \(x^{N+1}\) has a nonzero \(P_{N+1}\) coefficient.
\end{proof}
\end{proposition}

Embedding this commutative example as a central local observable
shows that agreement of all \(n\)-point functions only for
\(n\leq N\) cannot determine an interacting continuum state.
The nonzero stress--composite Wick excess in the Grand Tensor
programme is therefore a useful non-Gaussianity witness, but it
does not truncate the hierarchy.

\subsection{Finite-to-continuum hierarchy criterion}

\begin{theorem}[Diagonal landing of the bounded hierarchy]
\label{thm:v13v68-diagonal}
Let \(\{A_j\}_{j\geq1}\) be a countable norm-dense translated
local ledger, including adjoints and products through a
countable algebraic closure.  Suppose that for every word
\(\mathbf j=(j_1,\ldots,j_m)\) and every rational translation
tuple \(\mathbf x\), the regulated correlations
\[
 W_n(\mathbf j;\mathbf x)                                   \tag{68.18}
\]
converge.  Assume the common ledger is represented in one
candidate quasi-local \(C^*\)-algebra and every regulated
functional is contractive on every finite ledger polynomial:
\[
 |\omega_n(X)|\leq\|X\|.                                    \tag{68.19}
\]
In particular,
\(\lvert W_n(\mathbf j;\mathbf x)\rvert
 \leq\prod_{r=1}^m\|A_{j_r}\|\).
Assume the limiting values obey the compatibility, positivity,
locality, and translation-continuity moduli required to extend
from the rational ledger.  Then the limits define a unique state
on the norm closure of the ledger, and its bounded local
correlation hierarchy satisfies items 1--5 of
Theorem~\ref{thm:v13v68-hierarchy}.  If, in addition, the
limiting translations have joint spectrum in
\(\overline V_+\), item 6 also holds.

\begin{proof}
There are countably many pairs
\((\mathbf j,\mathbf x)\), so ordinary diagonal extraction gives
one subsequence on which all values converge if only
subsequential convergence is initially known.  Positivity is a
closed condition on each finite Gram matrix; algebraic
compatibility makes the limiting functional well defined on the
ledger algebra.  Bound (68.19) gives norm continuity, hence a
state on the closure.  The assumed translation modulus extends
the rational values continuously, and the limiting locality
identities extend by continuity.  The spectral-support argument
of Theorem~\ref{thm:v13v68-hierarchy} applies once the final
joint spectrum condition is available.
\end{proof}
\end{theorem}

The theorem deliberately lists compatibility and continuity as
hypotheses.  Separate convergence of scalar entries can otherwise
assign inconsistent values to two ledger words representing the
same algebra element.

\subsection{Acquisition consequence}

A continuum claim now needs a cofinal array indexed by
\[
 (\text{region},\text{word length},\text{source pair},
   \text{rational translation tuple}).                       \tag{68.20}
\]
At each fixed finite block one must retain positivity Grams,
adjoint/product consistency, the V67 ordering residual, the V64
cone leakage, and a translation-continuity modulus.  Word length
must grow cofinally; Proposition~\ref{prop:v13v68-finite-order}
forbids stopping at any fixed order.

\begin{assessment}[V68 advance]
\label{assess:v13v68}
Conditional on the Lorentzian net and forward-cone gate, all
bounded local vacuum correlations form a positive,
translation-invariant, local tempered hierarchy with the correct
difference-variable spectral support.  The all-order hierarchy on
a dense ledger determines the quasi-local vacuum state.  A
constructive Legendre perturbation proves that no fixed finite
correlation order can do so.
\end{assessment}

\begin{firewall}[Claim boundary after V68]
\label{fire:v13v68}
No covariant positive-energy Lorentzian SM net has yet been
populated, and the countable all-order ledger has not been shown
to converge.  Bounded observable distributions are not
operator-valued tempered fields.  Lorentz covariance, vacuum
uniqueness, clustering, field domains, stress-tensor
conservation, interacting measure construction, Einstein
backreaction, and the full SM--Einstein continuum remain open.
\end{firewall}

\subsection{Parent record}

This V68 note appends the preceding material to the validated V67
source, whose record was
\[
\begin{array}{c|c}
 \text{lines}&21627\\
 \text{bytes}&763074\\
 \text{SHA--256}&
 \texttt{7A2392775DF7DD99DBE56FED1858FE4B34F4DA47BD46DF35080CE0E2D0807D59}.
\end{array}                                                     \tag{68.21}
\]

% =============================================================================
% END APPENDED THIRTEENTH-CYCLE V68 MATERIAL
% =============================================================================

% =============================================================================
% BEGIN APPENDED THIRTEENTH-CYCLE V69 MATERIAL
% =============================================================================

\section{Thirteenth-cycle V69: the Bianchi--Einstein residual split}
\label{sec:v13v69}

\subsection{Why conservation is only one Einstein channel}

The previous cycles already separated mean Einstein equations,
Schwinger--Dyson identities, and almost-sure residual constraints,
and they supplied general uniform-integrability criteria for
passing unbounded Einstein writers to a limit.  The Grand Tensor
manuscript separately supplies a conditional represented
stress/current sector.  The remaining structural issue is to
couple the two without using stress conservation as a substitute
for the Einstein equation.

This note gives an orthogonal decomposition of the Einstein
residual into a diffeomorphism-tangent channel and a transverse
channel.  The contracted Bianchi identity annihilates the first
channel of the geometric tensor identically.  Consequently that
channel tests conservation of matter stress, while the second
channel carries the actual Einstein equation.

\subsection{Hilbert-screen geometry}

Let \(M\) be a compact smooth manifold without boundary and let
\(g\) be a \(C^2\) Riemannian metric.  The same calculation applies
on a relatively compact region when all test vector fields have
compact support.  Set
\[
 \mathscr H_g
 :=L^2(M,S^2T^*M;d{\rm vol}_g),\qquad
 \mathscr X_g:=L^2(M,T^*M;d{\rm vol}_g).                     \tag{69.1}
\]
On the Sobolev domain \(H^1(T^*M)\), define the infinitesimal
diffeomorphism map
\[
 K_gX:=\mathcal L_Xg=2\nabla_{(\mu}X_{\nu)}.                  \tag{69.2}
\]
For smooth tensors and vector fields,
\[
 \langle S,K_gX\rangle_g
 =-2\langle\operatorname{div}_gS,X\rangle_g,                 \tag{69.3}
\]
so \(K_g^*S=-2\operatorname{div}_gS\) on its natural domain.
Write
\[
 Q_g:\mathscr H_g\to\overline{\operatorname{Ran}K_g},
 \qquad P_g:=I-Q_g.
                                                                    \tag{69.4}
\]
The closed-range theorem gives
\[
 \mathscr H_g
 =\overline{\operatorname{Ran}K_g}
   \mathbin{\widehat\oplus}\ker K_g^*.                        \tag{69.5}
\]

Fix \(\kappa\in\mathbb R\setminus\{0\}\).  Let
\[
 \mathcal E_g:=G_g+\Lambda g,\qquad
 \mathcal R(g,T):=\mathcal E_g-\kappa T                       \tag{69.6}
\]
whenever the stress tensor \(T\) belongs to \(\mathscr H_g\).
The contracted Bianchi identity and metric compatibility give
\[
 K_g^*\mathcal E_g=0.                                        \tag{69.7}
\]

\begin{theorem}[Orthogonal Bianchi--Einstein split]
\label{thm:v13v69-split}
For every \(T\in\mathscr H_g\),
\[
 \boxed{
 \begin{aligned}
 Q_g\mathcal R(g,T)&=-\kappa Q_gT,\\
 P_g\mathcal R(g,T)&=\mathcal E_g-\kappa P_gT,
 \end{aligned}}                                              \tag{69.8}
\]
and
\[
 \boxed{\quad
 \|\mathcal R(g,T)\|_g^2
 =\kappa^2\|Q_gT\|_g^2
  +\|\mathcal E_g-\kappa P_gT\|_g^2.
 \quad}                                                       \tag{69.9}
\]
In particular,
\[
 Q_g\mathcal R(g,T)=0
 \iff K_g^*T=0,                                               \tag{69.10}
\]
whereas the full weak Einstein equation holds if and only if both
orthogonal terms in (69.9) vanish.

\begin{proof}
Equation (69.7) says
\(\mathcal E_g\in\ker K_g^*=\operatorname{Ran}P_g\).
Thus \(Q_g\mathcal E_g=0\) and
\(P_g\mathcal E_g=\mathcal E_g\).  Applying the two projections
to (69.6) proves (69.8).  Pythagoras in (69.5) gives (69.9).

Now \(Q_gT=0\) exactly when
\(T\in(\overline{\operatorname{Ran}K_g})^\perp
=\ker K_g^*\).  Combining this with the first line of (69.8)
proves (69.10).  Finally a vector in an orthogonal direct sum is
zero exactly when both components are zero.
\end{proof}
\end{theorem}

\begin{corollary}[Stress conservation does not imply Einstein]
\label{cor:v13v69-conservation-no-einstein}
The matter Ward identity
\[
 \operatorname{div}_gT=0                                     \tag{69.11}
\]
annuls only the first term in (69.9).  It does not imply
\(\mathcal E_g=\kappa T\).

\begin{proof}
The first claim is (69.10).  For a counterexample, choose any
metric for which
\(\mathcal E_g\neq0\) and take \(T=0\).  Both
\(\mathcal E_g\) and \(T\) are divergence free, but the Einstein
residual is \(\mathcal E_g\neq0\).  For example, a flat compact
torus with \(\Lambda\neq0\) has
\(\mathcal E_g=\Lambda g\).
\end{proof}
\end{corollary}

This rules out a circular argument in which a diffeomorphism Ward
identity is first used to prove stress conservation and that
conservation is then declared to be the gravitational field
equation.

\subsection{Mean, localized, and almost-sure equations}

Let \((\mathcal Z,\mathcal F,\mu)\) carry a random pair
\(z=(g,T)\).  Let \(\{h_j\}_{j\geq1}\) be a countable dense set of
smooth compactly supported symmetric test tensors and put
\[
 r_j(z):=\langle\mathcal R(g,T),h_j\rangle.                   \tag{69.12}
\]
Assume every \(r_j\) is measurable and integrable.

\begin{theorem}[Localized mean identities are equivalent to an
almost-sure weak equation]
\label{thm:v13v69-localized}
Suppose the residual is almost surely a continuous functional on
the test-tensor space and the countable family \(\{h_j\}\) is
dense in its chosen separable topology.  Then the following are
equivalent:

\begin{enumerate}
\item \(\mathcal R(g,T)=0\) as a weak tensor for
\(\mu\)-almost every \(z\);
\item for every \(j\) and every bounded real
\(\mathcal F\)-measurable function \(F\),
\[
 \int_{\mathcal Z}F(z)r_j(z)\,d\mu(z)=0.                     \tag{69.13}
\]
\end{enumerate}

The unlocalized mean equations
\[
 \int r_j\,d\mu=0                                             \tag{69.14}
\]
alone are strictly weaker.

\begin{proof}
Item 1 immediately implies item 2.  Conversely, fix \(j\) and
use in (69.13) the indicators of
\(\{r_j>0\}\) and \(\{r_j<0\}\).  The positive and negative
parts of \(r_j\) both have zero integral, so \(r_j=0\) almost
surely.  Intersecting the resulting probability-one sets over
the countable index \(j\) gives one set on which all \(r_j\)
vanish.  Density and continuity of the residual functional make
it vanish on every test tensor there.

To see strictness, fix a nonzero continuous residual functional
\(S\) and give \(+S\) and \(-S\) probability \(1/2\).  Equation
(69.14) holds for every test tensor, although the residual is
nonzero in every sample.
\end{proof}
\end{theorem}

When \(r_j\) is complex, the same proof is applied to its real and
imaginary parts.  It is enough in (69.13) to use a
measure-determining algebra of bounded cylinder functions and
then close it in the appropriate monotone-class argument.  This
is the operational meaning of coupling the Einstein writer to
the full joint state rather than checking only its vacuum mean.

\subsection{Positive residual Grams prevent cancellation}

For a finite test screen \(h_1,\ldots,h_q\), set
\[
 r(z):=(r_1(z),\ldots,r_q(z))^{\mathsf T},\qquad
 \mathbb G^{\rm Ein}_q
 :=\int_{\mathcal Z}r(z)r(z)^*\,d\mu(z).                     \tag{69.15}
\]

\begin{proposition}[Finite-screen almost-sure certificate]
\label{prop:v13v69-gram}
If \(r\in L^2(\mu;\mathbb C^q)\), then
\[
 \mathbb G^{\rm Ein}_q=0
 \iff r(z)=0\quad\text{for }\mu\text{-almost every }z.        \tag{69.16}
\]
By contrast, \(\int r\,d\mu=0\) permits cancellation between
nonzero residuals.

\begin{proof}
The matrix in (69.15) is positive and
\[
 \operatorname{tr}\mathbb G^{\rm Ein}_q
 =\int\|r(z)\|_{\mathbb C^q}^2\,d\mu(z).                     \tag{69.17}
\]
It is zero exactly when this nonnegative integrand vanishes
almost surely.  The final assertion is the two-point example in
Theorem~\ref{thm:v13v69-localized}.
\end{proof}
\end{proposition}

For a metric-dependent operator screen, let \(E_q\) be a fixed
finite coefficient space and let
\(R_q(z):E_q\to\mathscr H_{g(z)}\) synthesize the whitened
Einstein residuals on one transported physical tensor metric.
Define
\[
\begin{split}
 \mathbb G_q^{\rm Ein}
   &:=\int R_q(z)^*R_q(z)\,d\mu(z),\\
 \mathbb G_q^{\rm diff}
   &:=\int R_q(z)^*Q_{g(z)}R_q(z)\,d\mu(z),\\
 \mathbb G_q^{\rm trans}
   &:=\int R_q(z)^*P_{g(z)}R_q(z)\,d\mu(z).
\end{split}
\]
Then orthogonality gives the exact positive decomposition
\[
 \mathbb G_q^{\rm Ein}
 =
 \mathbb G_q^{\rm diff}
 +\mathbb G_q^{\rm trans}.                                   \tag{69.18}
\]
The first block measures the stress-conservation or
diffeomorphism-Ward defect.  The second measures the transverse
Einstein defect.  Their positivity forbids cancellation between
the two.

\subsection{Cofinal weak landing}

\begin{corollary}[Cofinal test-screen criterion]
\label{cor:v13v69-cofinal}
Let \(h_1,h_2,\ldots\) be as in
Theorem~\ref{thm:v13v69-localized}.  If
\[
 \mathbb G^{\rm Ein}_q=0
 \quad\text{for every }q,                                    \tag{69.19}
\]
then the weak Einstein equation holds almost surely.  More
generally, it is enough that for every fixed \(q\) a sequence of
regulated positive residual Grams converges to
\(\mathbb G^{\rm Ein}_q\) and tends to zero.

\begin{proof}
Proposition~\ref{prop:v13v69-gram} gives
\(r_j=0\) almost surely for every \(j\).  Countability, density,
and continuity finish the proof exactly as in
Theorem~\ref{thm:v13v69-localized}.  The convergence variant first
identifies each limiting Gram with zero.
\end{proof}
\end{corollary}

The same conclusion cannot be inferred by letting the test tensor
move with the cutoff.  A fixed cofinal ledger is needed to prevent
the Einstein residual from escaping every previously tested
screen, just as in the moving-core and dark-sector examples of
V63 and V66.

\subsection{Acquisition order for the coupled equation}

A noncircular SM--Einstein landing must now supply, on one joint
state and one transported tensor carrier,

\begin{enumerate}
\item \(C^2\) metric control and an inverse-metric floor sufficient
to define \(\mathcal E_g\), \(K_g\), and their weak pairings;
\item occurrence and uniform integrability of the same
renormalized stress tensor that entered the Ward and Hamiltonian
incidence compilers;
\item a vanishing contracted-Bianchi discretization defect for the
geometric writer, rather than an assumed exact continuum identity
at finite cutoff;
\item stable diffeomorphism-tangent projections, including
Killing-kernel rank changes and source transport;
\item separate positive Grams for the diffeomorphism and transverse
blocks in (69.18);
\item bounded cylinder multipliers or an equivalent conditional
test family, so that the desired claim is almost sure rather than
only annealed;
\item a countable cofinal test-tensor ledger with fixed-screen
limits before the screen is enlarged.
\end{enumerate}

The archived superlinear-moment theorem may pass the unbounded
writers once these occurrence and uniform-integrability rows are
populated.  It does not manufacture the rows or identify the
stress tensor with the one used in the Lorentzian local net.

\begin{assessment}[V69 advance]
\label{assess:v13v69}
The gravitational field equation has been decomposed into two
orthogonal, operationally distinct channels.  The
diffeomorphism-tangent channel is exactly stress conservation;
the transverse channel contains the remaining Einstein equation.
Positive residual Grams and localized cylinder tests prevent
random sign cancellation, and a countable cofinal tensor ledger
upgrades their zeros to an almost-sure weak equation.
\end{assessment}

\begin{firewall}[Claim boundary after V69]
\label{fire:v13v69}
No regulated SM stress tensor has been coupled to the same
random metric law and Lorentzian local state with the listed
uniform bounds.  The Bianchi projection, transverse residual
Gram, localized multipliers, and cofinal test ledger have not
been populated.  No almost-sure Einstein equation, quantum
metric, semiclassical solution, or full SM--Einstein continuum
is proved.
\end{firewall}

\subsection{Parent record}

This V69 note appends the preceding material to the validated V68
source, whose record was
\[
\begin{array}{c|c}
 \text{lines}&21951\\
 \text{bytes}&775255\\
 \text{SHA--256}&
 \texttt{0DE8BE381A199E3627D532B37F78E7556C754FC3D891609D490F78CE5CB52B2B}.
\end{array}                                                     \tag{69.20}
\]

% =============================================================================
% END APPENDED THIRTEENTH-CYCLE V69 MATERIAL
% =============================================================================

% =============================================================================
% BEGIN APPENDED THIRTEENTH-CYCLE V70 MATERIAL
% =============================================================================

\section{Thirteenth-cycle V70: compulsory companion audit before
metric-dependent Bianchi transport}
\label{sec:v13v70}

\subsection{Companion records}

The newest complete companion sources in the shared folder were
inspected before extending the V69 Einstein split:
\[
\begin{array}{c|c|r|r|c}
 \text{programme}&\text{version}&\text{lines}&\text{bytes}&
 \text{SHA--256}\\ \hline
 \text{Yang--Mills mass gap}&V94&30734&1048312&
 \texttt{EE5746919E3F8EDE5DA57FC1F2F1DC9A88C4C779222F3E20F4CA36EB908F0803}\\
 \text{Navier--Stokes stability}&V26&5319&278283&
 \texttt{82C887F8C2C69E4F28FAD7C3DC8DE5B9099CB5A4BF431EBA1FFF3E91935978AD}.
\end{array}                                                     \tag{70.1}
\]
The auxiliary YM parallel note remained at V5, with SHA--256
\[
 \texttt{306F1BB84CFA8A8D47EA569EC17E9545F9867C8D32B548EC9D9C444B4F3C0512}.
                                                                  \tag{70.2}
\]

\subsection{New YM information}

YM V92 constructs a separate exact rational inverse context at
every quarter-turn momentum.  V93 extends those anchors to a
rigorous full-torus cover by Laurent Lipschitz stocks and records
a residual-retaining local-cell lemma.  V94 propagates the global
stocks through the response word, but the resulting second-order
envelope is too large for a sign conclusion.  It identifies two
distinct sources of loss: vertex transport and normalized
covariance jets.

None of the numerical constants applies to the SM--Einstein
Bianchi operator.  Three proof rules do transfer:

\begin{enumerate}
\item exact anchor screens must be extended to neighborhoods by a
proved variation stock;
\item a finite global domination estimate remains useful even
when it is too coarse for a sign or contraction theorem;
\item when a global envelope is excessive, its separate factors
must be inspected upstream before refining cells.
\end{enumerate}

NS V26 and auxiliary YM V5 contain no new result since the V65
audit.  The NS rank-one computation continues to warn that a
restricted input geometry may miss a mixed channel.

\begin{proposition}[V70 direction decision]
\label{prop:v13v70-decision}
The next Einstein step is the stability of the
metric-dependent projections \(Q_g\) and \(P_g\) in (69.4).
It shall be treated upstream at the infinitesimal
diffeomorphism map \(K_g\), with:

\begin{enumerate}
\item a common transported tensor carrier;
\item a positive lower bound for the nonzero singular values of
\(K_g\);
\item explicit allowance for Killing vectors in \(\ker K_g\);
\item a perturbation bound for the range projection;
\item a counterexample when the nonzero singular-value floor
collapses;
\item propagation of the resulting projection error into both
positive Einstein residual Grams.
\end{enumerate}

\begin{proof}
The decomposition in Theorem~\ref{thm:v13v69-split} is meaningful
at every fixed metric, but a continuum limit requires its two
subspaces to converge.  Range projections are controlled by the
angle between the corresponding ranges, which in turn is
controlled by perturbation of \(K_g\) divided by its smallest
nonzero singular value.  This is the exact analogue of retaining
the inverse floor in the YM anchor transport.  Killing vectors
form the zero singular block and therefore must not be charged by
the positive floor.  If the next nonzero singular value approaches
zero, the projection can jump even while \(K_g\) converges in
norm, so no downstream Gram estimate can repair the missing gap.
\end{proof}
\end{proposition}

\begin{assessment}[V70 audit outcome]
\label{assess:v13v70}
The companion audit supplies no transferable constant, but it
selects the correct upstream bottleneck.  Before attempting a
cofinal Einstein residual, the programme must stabilize the
diffeomorphism-tangent and transverse tensor subspaces with a
protected nonzero singular-value floor and explicit Killing
kernel.
\end{assessment}

\begin{firewall}[Claim boundary after V70]
\label{fire:v13v70}
This audit proves no metric-projection estimate and populates no
SM--Einstein residual.  No stress conservation, Einstein
equation, Lorentzian interacting state, or full SM--Einstein
continuum is established.
\end{firewall}

\subsection{Parent record}

This V70 note appends the preceding material to the validated V69
source, whose record was
\[
\begin{array}{c|c}
 \text{lines}&22303\\
 \text{bytes}&788144\\
 \text{SHA--256}&
 \texttt{36EC907143AF45FCE0CC7F95AB1734DA0680D91BB63E8700A8A9F34D1EEAEED7}.
\end{array}                                                     \tag{70.3}
\]

% =============================================================================
% END APPENDED THIRTEENTH-CYCLE V70 MATERIAL
% =============================================================================

% =============================================================================
% BEGIN APPENDED THIRTEENTH-CYCLE V71 MATERIAL
% =============================================================================

\section{Thirteenth-cycle V71: Killing-safe stability of the
Bianchi projections}
\label{sec:v13v71}

\subsection{The correct positive floor}

For a bounded operator \(K:\mathcal X\to\mathcal H\) between
Hilbert spaces, define its reduced minimum modulus
\[
 \sigma(K):=
 \inf\{\|Kx\|:\ x\in(\ker K)^\perp,\ \|x\|=1\}.              \tag{71.1}
\]
The condition \(\sigma(K)>0\) is equivalent to closed range.  It
charges only the nonzero singular block and therefore permits a
nontrivial kernel.

For the geometric operator \(K_gX=\mathcal L_Xg\), the kernel is
the Killing-vector space.  After choosing the \(H^1\) norm on
vector fields and the \(L^2\) norm on symmetric tensors, a lower
bound for (71.1) is precisely a Korn-type estimate modulo Killing
fields:
\[
 \|K_gX\|_{L^2_g}\geq
 \sigma_0\|X\|_{H^1_g}
 \quad\bigl(X\perp\ker K_g\bigr).                            \tag{71.2}
\]
No positive lower bound is imposed on the Killing block itself.

\subsection{An output-range perturbation theorem}

\begin{theorem}[Reduced-gap range-projection bound]
\label{thm:v13v71-projection}
Let \(K,L:\mathcal X\to\mathcal H\) be bounded operators with
closed ranges, and let \(Q_K,Q_L\) be the orthogonal projections
onto their ranges.  Put
\[
 \delta:=\|K-L\|,\qquad
 \sigma_*:=\min\{\sigma(K),\sigma(L)\}>0.                    \tag{71.3}
\]
Then
\[
 \boxed{\quad
 \|Q_K-Q_L\|
 \leq \min\left\{1,{\delta\over\sigma_*}\right\}.
 \quad}                                                       \tag{71.4}
\]
If \(\mathcal X,\mathcal H\) are finite dimensional and
\(\delta<\sigma_*\), the two ranges have the same dimension and
the two kernels have the same dimension.

\begin{proof}
Let \(y\in\operatorname{Ran}K\) with \(\|y\|=1\).  The unique
minimal-norm preimage \(x=K^\dagger y\) lies in
\((\ker K)^\perp\) and satisfies
\[
 \|x\|\leq\sigma(K)^{-1}.                                    \tag{71.5}
\]
Since \(Lx\in\operatorname{Ran}L\),
\[
 \|(I-Q_L)y\|
 =\operatorname{dist}(Kx,\operatorname{Ran}L)
 \leq\|(K-L)x\|
 \leq{\delta\over\sigma(K)}.                                 \tag{71.6}
\]
Hence
\[
 \|(I-Q_L)Q_K\|\leq{\delta\over\sigma(K)}.                   \tag{71.7}
\]
Interchanging \(K,L\) gives
\[
 \|(I-Q_K)Q_L\|\leq{\delta\over\sigma(L)}.                   \tag{71.8}
\]
For orthogonal projections,
\[
 \|Q_K-Q_L\|
 =
 \max\{\|(I-Q_L)Q_K\|,\|(I-Q_K)Q_L\|\}.                      \tag{71.9}
\]
Equations (71.7)--(71.9), together with the universal bound
\(\|Q_K-Q_L\|\leq1\), prove (71.4).

In finite dimension, projections of different rank have norm
distance one.  Thus \(\delta<\sigma_*\) and (71.4) force equal
range ranks.  Rank-nullity gives equal kernel dimensions.
\end{proof}
\end{theorem}

The theorem needs no coordinate identification of individual
Killing vectors.  It controls the physical output subspace
\(\operatorname{Ran}K_g\) directly.  A basis-level Killing
transport is needed only if later acquisitions attach separate
labels or anchors to individual kernel elements.

\begin{proposition}[The reduced gap cannot be omitted]
\label{prop:v13v71-gap-no-go}
There are matrices \(K_\epsilon\to K_0\) in operator norm whose
range projections do not converge.

\begin{proof}
On \(\mathbb R^2\), let
\[
 K_\epsilon=
 \begin{pmatrix}1&0\\0&\epsilon\end{pmatrix}.
                                                                    \tag{71.10}
\]
For every \(\epsilon>0\),
\(\operatorname{Ran}K_\epsilon=\mathbb R^2\), so
\(Q_\epsilon=I\).  At \(\epsilon=0\),
\[
 Q_0=\begin{pmatrix}1&0\\0&0\end{pmatrix},\qquad
 \|Q_\epsilon-Q_0\|=1.                                      \tag{71.11}
\]
Yet \(\|K_\epsilon-K_0\|=\epsilon\to0\).  The disappearing
singular value is \(\sigma(K_\epsilon)=\epsilon\), and the kernel
dimension jumps at the limit.
\end{proof}
\end{proposition}

Thus a changing Killing dimension is not a harmless relabeling.
It is exactly a loss of the reduced floor required to transport
the Bianchi split.

\subsection{Metric-dependent carriers}

Let \(g,h\) lie in one uniformly elliptic metric chart.  Suppose
unitary whitening maps identify their vector and tensor Hilbert
spaces with fixed carriers \(\mathcal X,\mathcal H\), and denote
the transported symmetric-gradient operators by
\[
 \widetilde K_g,\widetilde K_h:\mathcal X\to\mathcal H.       \tag{71.12}
\]

\begin{corollary}[Korn-controlled metric variation]
\label{cor:v13v71-metric}
Assume
\[
 \sigma(\widetilde K_g),\sigma(\widetilde K_h)\geq\sigma_0>0,
 \qquad
 \|\widetilde K_g-\widetilde K_h\|
 \leq L_K\,d_{\rm met}(g,h).                                 \tag{71.13}
\]
Then the transported diffeomorphism-tangent projections obey
\[
 \|\widetilde Q_g-\widetilde Q_h\|
 \leq
 \min\left\{1,{L_K\over\sigma_0}d_{\rm met}(g,h)\right\},     \tag{71.14}
\]
and the transverse projections
\(\widetilde P_g=I-\widetilde Q_g\) obey the same bound.

\begin{proof}
Apply Theorem~\ref{thm:v13v71-projection} to the two transported
operators.  The equality
\(\widetilde P_g-\widetilde P_h
=-(\widetilde Q_g-\widetilde Q_h)\) proves the second claim.
\end{proof}
\end{corollary}

In a concrete metric core, (71.13) separates into two acquisitions:
a uniform Korn constant modulo Killing fields and a coefficient
variation estimate for \(K_g\) after density and index whitening.
Compactness of a \(C^1\) metric family does not by itself prove a
uniform Korn constant when Killing rank can change.

\subsection{Transport of the two Einstein channels}

Let \(E\) be a finite coefficient space and let
\[
 R_g,R_h:E\to\mathcal H                                      \tag{71.15}
\]
be transported whitened Einstein-residual syntheses.  Define
\[
 D_g:=\widetilde Q_gR_g,\qquad
 Z_g:=\widetilde P_gR_g,                                     \tag{71.16}
\]
and similarly at \(h\).  Here \(D\) is the
diffeomorphism-tangent channel and \(Z\) the transverse channel.

\begin{theorem}[Residual and Gram transport]
\label{thm:v13v71-gram-transport}
Suppose
\[
 \|R_g\|,\|R_h\|\leq M_R,\qquad
 \|R_h-R_g\|\leq\varepsilon_R,\qquad
 \|\widetilde Q_h-\widetilde Q_g\|\leq\varepsilon_Q.         \tag{71.17}
\]
Then
\[
 \boxed{
 \begin{aligned}
 \|D_h-D_g\|&\leq\varepsilon_R+M_R\varepsilon_Q,\\
 \|Z_h-Z_g\|&\leq\varepsilon_R+M_R\varepsilon_Q.
 \end{aligned}}                                              \tag{71.18}
\]
For the positive residual Grams
\[
 \mathbb G_g^{\rm diff}:=D_g^*D_g,\qquad
 \mathbb G_g^{\rm trans}:=Z_g^*Z_g,                           \tag{71.19}
\]
one has
\[
 \boxed{
 \begin{aligned}
 \|\mathbb G_h^{\rm diff}-\mathbb G_g^{\rm diff}\|
 &\leq2M_R(\varepsilon_R+M_R\varepsilon_Q),\\
 \|\mathbb G_h^{\rm trans}-\mathbb G_g^{\rm trans}\|
 &\leq2M_R(\varepsilon_R+M_R\varepsilon_Q).
 \end{aligned}}                                              \tag{71.20}
\]

\begin{proof}
Add and subtract \(\widetilde Q_hR_g\):
\[
 D_h-D_g
 =\widetilde Q_h(R_h-R_g)
  +(\widetilde Q_h-\widetilde Q_g)R_g.                       \tag{71.21}
\]
Projection contractivity gives the first line of (71.18).
The same calculation with
\(\widetilde P_h-\widetilde P_g
=-(\widetilde Q_h-\widetilde Q_g)\) gives the second.

For either channel \(Y\in\{D,Z\}\),
\[
 Y_h^*Y_h-Y_g^*Y_g
 =Y_h^*(Y_h-Y_g)+(Y_h^*-Y_g^*)Y_g.                           \tag{71.22}
\]
Both \(\|Y_h\|\) and \(\|Y_g\|\) are at most \(M_R\).
Using (71.18) in (71.22) proves (71.20).
\end{proof}
\end{theorem}

Combining Corollary~\ref{cor:v13v71-metric} with (71.20) gives
the explicit metric-chart estimate
\[
 \|\Delta\mathbb G^{\rm diff/trans}\|
 \leq
 2M_R\left(
 \varepsilon_R+
 {M_RL_K\over\sigma_0}d_{\rm met}(g,h)\right)               \tag{71.23}
\]
as long as the right side of (71.14) is on its
nontrivial branch.

\subsection{Finite Bianchi defects}

A discretized geometric writer need not satisfy the continuum
contracted Bianchi identity exactly.  Write its transported
source synthesis as
\[
 \mathcal E_g:E\to\mathcal H,\qquad
 T_g:E\to\mathcal H,\qquad
 R_g=\mathcal E_g-\kappa T_g,                                \tag{71.24}
\]
and define
\[
 \beta_g:=\|\widetilde Q_g\mathcal E_g\|.                    \tag{71.25}
\]

\begin{proposition}[Conservation residual with a Bianchi debit]
\label{prop:v13v71-bianchi-debit}
For \(\kappa\neq0\),
\[
 \left|
 \|\widetilde Q_gR_g\|
 -|\kappa|\,\|\widetilde Q_gT_g\|
 \right|
 \leq\beta_g.                                                \tag{71.26}
\]
Consequently,
\[
 \|\widetilde Q_gT_g\|
 \leq {\,\|\widetilde Q_gR_g\|+\beta_g\over|\kappa|}.         \tag{71.27}
\]

\begin{proof}
Projection of (71.24) gives
\[
 \widetilde Q_gR_g
 =\widetilde Q_g\mathcal E_g
  -\kappa\widetilde Q_gT_g.                                  \tag{71.28}
\]
The reverse triangle inequality proves (71.26), and the ordinary
triangle inequality gives (71.27).
\end{proof}
\end{proposition}

Therefore a small tangent Einstein residual is not yet a small
stress-conservation residual unless the finite Bianchi debit
\(\beta_g\) is also small.  In the transverse channel one must
separately control
\[
 \|\widetilde P_g\mathcal E_g
   -\kappa\widetilde P_gT_g\|,                               \tag{71.29}
\]
which no Ward identity can remove.

\subsection{Cofinal landing consequence}

\begin{corollary}[Stable Bianchi--Einstein split in the limit]
\label{cor:v13v71-landing}
Let transported metrics \(g_n\to g_\infty\) satisfy
\[
 \|\widetilde K_{g_n}-\widetilde K_{g_\infty}\|\to0,\qquad
 \inf_n\sigma(\widetilde K_{g_n})\wedge
 \sigma(\widetilde K_{g_\infty})>0.                          \tag{71.30}
\]
If \(R_{g_n}\to R_\infty\) in operator norm on every fixed
finite coefficient screen, then both projected residuals and
both positive Grams converge there.  If their cofinal limiting
Grams vanish, the almost-sure weak Einstein conclusion of
Corollary~\ref{cor:v13v69-cofinal} applies, provided the
localized random-state hypotheses of V69 also hold.

\begin{proof}
Theorem~\ref{thm:v13v71-projection} gives
\(\widetilde Q_{g_n}\to\widetilde Q_{g_\infty}\) and likewise for
\(P\).  Theorem~\ref{thm:v13v71-gram-transport} gives the
projected residual and Gram convergence.  The final assertion is
exactly the cofinal positive-Gram implication in V69.
\end{proof}
\end{corollary}

\begin{assessment}[V71 advance]
\label{assess:v13v71}
The metric-dependent Bianchi split is now stable under a precise
and Killing-safe hypothesis.  Norm variation of the symmetric
gradient divided by its reduced Korn gap controls both
diffeomorphism-tangent and transverse projections.  The estimate
propagates explicitly to the two positive Einstein residual
Grams, while a collapsing nonzero singular value is shown to
permit a unit jump.  A finite contracted-Bianchi error is retained
as a separate debit in the conservation channel.
\end{assessment}

\begin{firewall}[Claim boundary after V71]
\label{fire:v13v71}
No uniform Korn gap, metric-whitening chart, geometric coefficient
variation, finite Bianchi debit, or SM stress residual has been
populated on the coupled regulator.  Possible Killing-rank changes
have not been excluded on the gravitational background core.  No
almost-sure Einstein equation, Lorentzian quantum metric, or full
SM--Einstein continuum is proved.
\end{firewall}

\subsection{Parent record}

This V71 note appends the preceding material to the validated V70
source, whose record was
\[
\begin{array}{c|c}
 \text{lines}&22427\\
 \text{bytes}&793113\\
 \text{SHA--256}&
 \texttt{5E59CCB6693659719979563A1B68143DFE3EDD05980E771BF24EAF6CC69B9879}.
\end{array}                                                     \tag{71.31}
\]

% =============================================================================
% END APPENDED THIRTEENTH-CYCLE V71 MATERIAL
% =============================================================================

% =============================================================================
% BEGIN APPENDED THIRTEENTH-CYCLE V72 MATERIAL
% =============================================================================

\section{Thirteenth-cycle V72: covariant metric fibres and the
global positive-energy obstruction}
\label{sec:v13v72}

\subsection{The fixed-background assumption that must be removed}

The archived compactness theorem constructs conditional
classical--quantum states and trace-class disintegrations over a
metric background law.  V64--V68, by contrast, use one Hilbert
space carrying one positive-energy translation representation.
Passing from the first statement to the second is not automatic:
a spacetime symmetry may move the metric base as well as the
quantum fibre.

This note gives the exact direct-integral construction.  It also
exhibits a sharp obstruction: nonnegative Hamiltonians in every
fixed metric fibre do not imply a nonnegative generator after
time translation moves the metric variable.

\subsection{A unitary cocycle over the metric law}

Let \((X,\mathcal B,\mu)\) be a standard probability space of
admissible metric backgrounds, let
\(\{\mathcal H_x\}_{x\in X}\) be a measurable field of separable
Hilbert spaces, and set
\[
 \mathcal H
 :=\int_X^\oplus\mathcal H_x\,d\mu(x).                       \tag{72.1}
\]
Let a group \(G\) act measurably on \(X\) and preserve \(\mu\).
Suppose there is a measurable unitary cocycle
\[
 V(g,x):\mathcal H_x\longrightarrow\mathcal H_{gx}           \tag{72.2}
\]
satisfying
\[
 V(e,x)=I,\qquad
 V(gh,x)=V(g,hx)V(h,x)                                       \tag{72.3}
\]
for almost every \(x\).

\begin{theorem}[Covariant direct-integral representation]
\label{thm:v13v72-cocycle}
Define
\[
 (U_g\psi)_y
 :=V(g,g^{-1}y)\psi_{g^{-1}y}.                               \tag{72.4}
\]
Then \(U_g\) is unitary on \(\mathcal H\) and
\[
 U_gU_h=U_{gh}.                                               \tag{72.5}
\]
If \(g\mapsto U_g\psi\) is continuous for every \(\psi\) in one
dense fundamental family, then \(U\) is strongly continuous.

Let
\[
 A=\int_X^\oplus A_x\,d\mu(x)                                \tag{72.6}
\]
be a bounded decomposable operator.  Its transform is
decomposable and satisfies
\[
 (U_gAU_g^*)_y
 =
 V(g,g^{-1}y)A_{g^{-1}y}V(g,g^{-1}y)^*.                     \tag{72.7}
\]
If a measurable unit vector field obeys
\[
 V(g,x)\Omega_x=\Omega_{gx},                                 \tag{72.8}
\]
then
\[
 \Omega:=\int_X^\oplus\Omega_x\,d\mu(x)
\]
is \(U\)-invariant, and
\[
 \omega(A)
 :=\langle\Omega,A\Omega\rangle
 =\int_X\langle\Omega_x,A_x\Omega_x\rangle\,d\mu(x)           \tag{72.9}
\]
is a \(G\)-invariant state on every covariant decomposable
algebra.

\begin{proof}
Using invariance of \(\mu\),
\[
\begin{split}
 \|U_g\psi\|^2
 &=\int_X
 \|V(g,g^{-1}y)\psi_{g^{-1}y}\|^2\,d\mu(y)\\
 &=\int_X\|\psi_x\|^2\,d\mu(x).
\end{split}                                                   \tag{72.10}
\]
The inverse is \(U_{g^{-1}}\), by (72.3), so \(U_g\) is unitary.
For \(y=ghx\), the fibre map in \(U_gU_h\) is
\[
 V(g,hx)V(h,x)=V(gh,x),                                      \tag{72.11}
\]
which proves (72.5).  Continuity on a dense family extends to all
vectors because the \(U_g\) are isometries.

Direct substitution in (72.4) gives (72.7).  Equation (72.8)
gives \(U_g\Omega=\Omega\).  Positivity and normalization of
(72.9) are immediate, and (72.7)--(72.8), followed by the change
of variables \(y=gx\), prove invariance.
\end{proof}
\end{theorem}

For a quasi-invariant rather than invariant metric law, (72.4)
requires the square root of the Radon--Nikodym derivative.  That
density, its cocycle identity, and its integrability are then
additional acquisition rows; they cannot be hidden in the fibre
unitary.

\subsection{Locality and fibrewise covariance}

Suppose each fibre contains a represented net
\(\mathcal O\mapsto\mathfrak A_x(\mathcal O)\), with a countable
measurable generating ledger.  Assume
\[
 V(g,x)\mathfrak A_x(\mathcal O)V(g,x)^*
 =\mathfrak A_{gx}(g\mathcal O).                             \tag{72.12}
\]

\begin{proposition}[Fibre locality lands globally]
\label{prop:v13v72-locality}
If
\[
 [\mathfrak A_x(\mathcal O_1),
   \mathfrak A_x(\mathcal O_2)]=0                            \tag{72.13}
\]
for almost every \(x\) whenever
\(\mathcal O_1\perp\mathcal O_2\), then the corresponding
decomposable global algebras commute.  Equation (72.12) and
Theorem~\ref{thm:v13v72-cocycle} give their covariance.

\begin{proof}
For decomposable \(A=\int^\oplus A_x\) and
\(B=\int^\oplus B_x\),
\[
 [A,B]=\int_X^\oplus[A_x,B_x]\,d\mu(x).                      \tag{72.14}
\]
It is zero when its fibres vanish almost everywhere.  Covariance
is (72.7) applied to the measurable local generators and then to
their closures.
\end{proof}
\end{proposition}

The spacelike relation in (72.13) may depend on \(x\) when the
metric itself is dynamical.  A single background-independent
region label is legitimate only when the admitted metric core
has a common causal separation certificate for that pair of
regions, or when the local net is formulated relationally.

\subsection{When fibre positivity does pass to the global theory}

\begin{theorem}[Fixed-base positive-energy landing]
\label{thm:v13v72-positive}
Suppose the time-translation subgroup fixes the metric base:
\[
 t x=x\quad(t\in\mathbb R,\ \mu\text{-almost every }x),       \tag{72.15}
\]
and
\[
 V(t,x)=e^{-itH_x}                                           \tag{72.16}
\]
for a measurable field of self-adjoint operators satisfying
\(H_x\geq0\) almost everywhere.  Then
\[
 U_t=\int_X^\oplus e^{-itH_x}\,d\mu(x)                       \tag{72.17}
\]
has generator
\[
 H=\int_X^\oplus H_x\,d\mu(x)\geq0.                          \tag{72.18}
\]
More generally, if a fixed-base spacetime translation cocycle has
joint fibre spectrum in \(\overline V_+\) almost everywhere, the
global joint spectrum also lies in \(\overline V_+\).

\begin{proof}
Under (72.15), formula (72.4) reduces to the decomposable unitary
(72.17).  The direct-integral spectral calculus shows that its
Stone generator is (72.18); its quadratic form is
\[
 \langle\psi,H\psi\rangle
 =\int_X\langle\psi_x,H_x\psi_x\rangle\,d\mu(x)\geq0          \tag{72.19}
\]
on the natural form domain.  For spacetime translations, the
global joint spectral projection of a Borel set is the direct
integral of the fibre joint spectral projections.  It vanishes
outside \(\overline V_+\) when the fibre projections do.
\end{proof}
\end{theorem}

\begin{proposition}[Fibrewise zero energy can produce a
two-sided global spectrum]
\label{prop:v13v72-spectrum-no-go}
If time translation moves the base, positivity of every internal
fibre Hamiltonian does not imply positivity of the generator of
(72.4).

\begin{proof}
Take
\[
 X=S^1,\qquad d\mu={d\theta\over2\pi},\qquad
 \mathcal H_\theta=\mathbb C,
\]
and let time act by \(t\theta=\theta+t\pmod{2\pi}\).  Choose the
trivial fibre cocycle \(V(t,\theta)=1\).  Every internal fibre
Hamiltonian is then zero.  Formula (72.4) gives
\[
 (U_t\psi)(\theta)=\psi(\theta-t).                            \tag{72.20}
\]
On the Fourier basis \(e_n(\theta)=e^{in\theta}\),
\[
 U_te_n=e^{-int}e_n.                                         \tag{72.21}
\]
The global generator has spectrum \(\mathbb Z\), which is
unbounded above and below.
\end{proof}
\end{proposition}

This counterexample is not cured by averaging the fibre
Hamiltonians: the missing energy is the generator of motion in
the metric base.  A coupled Lorentzian construction must either
show that physical time acts within the selected metric fibres,
construct a positive global generator including the base motion,
or replace coordinate time by a certified relational clock.

\subsection{Global cyclicity and the metric centre}

Let \(\mathfrak A_{\rm dec}\) be a decomposable algebra containing
the scalar multiplication algebra \(L^\infty(X,\mu)I\).

\begin{theorem}[Fibre-cyclic fields plus the centre are globally
cyclic]
\label{thm:v13v72-cyclic}
Assume there is a countable measurable family of decomposable
operators \(A_1,A_2,\ldots\in\mathfrak A_{\rm dec}\) such that
\[
 \overline{\operatorname{span}}\{A_j(x)\Omega_x:j\geq1\}
 =\mathcal H_x                                                \tag{72.22}
\]
for almost every \(x\).  Then the vector field \(\Omega\) in
(72.9) is cyclic for \(\mathfrak A_{\rm dec}\).

\begin{proof}
Let \(\zeta\in\mathcal H\) be orthogonal to
\(\mathfrak A_{\rm dec}\Omega\).  For every
\(f\in L^\infty(X,\mu)\) and every \(j\),
\[
 0=\langle\zeta,fA_j\Omega\rangle
 =\int_X f(x)
   \langle\zeta_x,A_j(x)\Omega_x\rangle\,d\mu(x).             \tag{72.23}
\]
Since this holds for all bounded \(f\), the scalar integrand
vanishes almost everywhere.  Countability gives one
probability-one set on which it vanishes for all \(j\).  The
density (72.22) then gives \(\zeta_x=0\) almost everywhere.
Thus the orthogonal complement of
\(\overline{\mathfrak A_{\rm dec}\Omega}\) is zero.
\end{proof}
\end{theorem}

If the physical diffeomorphism quotient removes the full
multiplication centre, fibrewise cyclicity alone does not prove
cyclicity for the smaller invariant algebra.  The observable
quotient, central superselection variables, and vacuum-selection
step must therefore be stated before applying the V68
state-determination theorem.

\subsection{Coupled acquisition card}

A coupled metric--quantum Lorentzian landing now requires:

\begin{enumerate}
\item a standard measurable metric quotient and one measurable
Hilbert field;
\item exact or vanishing-defect unitary cocycles with the
composition law (72.3);
\item invariance or controlled quasi-invariance of the metric law;
\item measurable local generators and the raw covariance
intertwiner (72.12);
\item a common causal-separation screen over the admitted metric
core, or a relational localization scheme;
\item an invariant measurable vacuum field;
\item either fixed-base time evolution as in
Theorem~\ref{thm:v13v72-positive} or a direct proof of the global
forward-cone spectral projection;
\item the V69--V71 almost-sure Einstein residual in the same
metric fibres and the V67 locality residual on the same
decomposable observable ledger.
\end{enumerate}

\begin{assessment}[V72 advance]
\label{assess:v13v72}
The random-metric and Lorentzian constructions now meet in an
exact direct-integral theorem.  A measure-preserving unitary
cocycle yields the global covariant state and transports
fibrewise locality.  Positive energy passes directly only when
time fixes the metric base; a circle-shift model proves that
zero fibre energies can otherwise produce a two-sided global
spectrum.  The result identifies the base-motion or relational
clock row that a coupled continuum must add.
\end{assessment}

\begin{firewall}[Claim boundary after V72]
\label{fire:v13v72}
No invariant metric law, measurable SM Hilbert field, symmetry
cocycle, common causal screen, invariant vacuum field, or global
positive-energy generator has been populated.  Fibrewise
Einstein residuals and Lorentzian locality have not been placed
on one direct-integral ledger.  No dynamical quantum metric or
full SM--Einstein continuum is proved.
\end{firewall}

\subsection{Parent record}

This V72 note appends the preceding material to the validated V71
source, whose record was
\[
\begin{array}{c|c}
 \text{lines}&22793\\
 \text{bytes}&805037\\
 \text{SHA--256}&
 \texttt{6C5429039B9A76F10E12598AD222549C07D64DC7217D8CA8DC8CB7863302CA12}.
\end{array}                                                     \tag{72.24}
\]

% =============================================================================
% END APPENDED THIRTEENTH-CYCLE V72 MATERIAL
% =============================================================================

% =============================================================================
% BEGIN APPENDED THIRTEENTH-CYCLE V73 MATERIAL
% =============================================================================

\section{Thirteenth-cycle V73: base-motion compensation and the
vacuum ground anchor}
\label{sec:v13v73}

\subsection{Separating internal and base generators}

In a covariant metric-fibre representation, suppose a candidate
time generator decomposes at quadratic-form level into
\[
 q_{\rm tot}=q_{\rm fib}+q_{\rm base}.                        \tag{73.1}
\]
Here \(q_{\rm fib}\) is the nonnegative direct-integral form of
the fibre OS Hamiltonians, while \(q_{\rm base}\) contains metric
transport and any cocycle connection term.  Proposition
\ref{prop:v13v72-spectrum-no-go} shows that the second form cannot
be omitted even when every fibre energy is zero.

Let \(q_{\rm fib}\) be a densely defined closed nonnegative form
on \(\mathcal H\), with domain \(\mathcal Q\).  Let
\(q_{\rm base}\) be a symmetric form on the same domain.

\begin{theorem}[Relative-form compensation criterion]
\label{thm:v13v73-form-sum}
Assume that for some \(0\leq a<1\) and \(b\geq0\),
\[
 |q_{\rm base}[\psi]|
 \leq a\,q_{\rm fib}[\psi]+b\|\psi\|^2
 \qquad(\psi\in\mathcal Q).                                  \tag{73.2}
\]
Then \(q_{\rm tot}\) is closed and lower semibounded and defines a
unique self-adjoint operator \(H_{\rm tot}\) with
\[
 H_{\rm tot}\geq-bI.                                         \tag{73.3}
\]
If the sharper one-sided estimate
\[
 q_{\rm base}[\psi]\geq-a\,q_{\rm fib}[\psi]
 \qquad(\psi\in\mathcal Q)                                   \tag{73.4}
\]
holds, then
\[
 \boxed{\quad
 H_{\rm tot}\geq(1-a)H_{\rm fib}\geq0
 \quad}                                                       \tag{73.5}
\]
in quadratic-form order.

\begin{proof}
The KLMN theorem applied to (73.2) makes the form sum closed and
lower semibounded.  Its lower estimate is
\[
 q_{\rm tot}[\psi]
 \geq(1-a)q_{\rm fib}[\psi]-b\|\psi\|^2
 \geq-b\|\psi\|^2,                                           \tag{73.6}
\]
which proves (73.3).  Under (73.4),
\[
 q_{\rm tot}[\psi]
 \geq(1-a)q_{\rm fib}[\psi]\geq0.                            \tag{73.7}
\]
The representation theorem for closed forms turns (73.7) into
the form inequality (73.5).
\end{proof}
\end{theorem}

The two hypotheses play different roles.  The absolute relative
bound (73.2) constructs a self-adjoint generator.  The one-sided
bound (73.4) certifies its sign.  A finite lower bound \(b\) alone
does not establish the relativistic spectrum condition.

\subsection{The absolute vacuum-energy anchor}

Let
\[
 E_0:=\inf\operatorname{sp}(H_{\rm tot})>-\infty.             \tag{73.8}
\]
Suppose a normalized vector \(\Omega\) is an eigenvector,
\[
 H_{\rm tot}\Omega=E_\Omega\Omega.                            \tag{73.9}
\]
Rephasing the unitary group so that it fixes \(\Omega\) replaces
the generator by
\[
 H_\Omega:=H_{\rm tot}-E_\Omega I.                           \tag{73.10}
\]

\begin{theorem}[Vacuum invariance and positive energy are
compatible exactly at the bottom]
\label{thm:v13v73-ground}
Under (73.8)--(73.10),
\[
 H_\Omega\geq0
 \iff E_\Omega=E_0.                                          \tag{73.11}
\]
Thus a vacuum phase anchor fixes the additive scalar, while a
ground-state certificate is still required to prove positive
energy.

\begin{proof}
The bottom of the shifted spectrum is
\[
 \inf\operatorname{sp}(H_\Omega)=E_0-E_\Omega.               \tag{73.12}
\]
Since \(E_\Omega\in\operatorname{sp}(H_{\rm tot})\),
\(E_\Omega\geq E_0\).  Hence (73.12) is nonnegative exactly when
equality holds.
\end{proof}
\end{theorem}

\begin{proposition}[An invariant vector need not be a vacuum
ground state]
\label{prop:v13v73-invariant-no-ground}
There is a strongly continuous unitary group with a fixed unit
vector whose generator has negative spectrum.

\begin{proof}
On \(\mathbb C^2\), let
\[
 H=\begin{pmatrix}-1&0\\0&0\end{pmatrix},
 \qquad\Omega=\begin{pmatrix}0\\1\end{pmatrix}.              \tag{73.13}
\]
Then \(e^{-itH}\Omega=\Omega\), but
\(\inf\operatorname{sp}(H)=-1\).  The vector is invariant and is
not a ground state.
\end{proof}
\end{proposition}

This is the spectral version of the absolute energy/pressure
fibre retained in the Grand Tensor Hamiltonian-incidence
compiler.  Commutators determine a Hamiltonian only up to a
scalar; fixing the scalar by one vacuum expectation is valid only
after the chosen vacuum has also been shown to minimize the full
coupled generator.

\subsection{A source-complete ground test}

Let \(H\) be self-adjoint and lower semibounded, let
\(\Omega\in\operatorname{Dom}H\) satisfy
\(H\Omega=E_\Omega\Omega\), and let \(q_H\) be the closed
quadratic form of \(H\), with form domain \(\mathcal Q(H)\).
Let \(\mathcal D_0\) be a form core and put
\[
 q_\Omega[\psi]
 :=q_H[\psi]-E_\Omega\|\psi\|^2,
 \qquad\psi\in\mathcal Q(H).                                 \tag{73.14}
\]

\begin{theorem}[Form-core vacuum-ground criterion]
\label{thm:v13v73-form-core}
The following are equivalent:
\[
 E_\Omega=\inf\operatorname{sp}(H),                          \tag{73.15}
\]
and
\[
 q_\Omega[\psi]\geq0
 \quad\text{for every }\psi\in\mathcal D_0.                  \tag{73.16}
\]

\begin{proof}
If (73.15) holds, the spectral theorem gives
\(H-E_\Omega\geq0\), proving (73.16).  Conversely, the closed
quadratic form of \(H-E_\Omega\) is nonnegative on its form core.
Every form-domain vector is the form-norm limit of a sequence in
\(\mathcal D_0\), so closedness extends nonnegativity to the
whole form domain.  Hence \(H-E_\Omega\geq0\), and
Theorem~\ref{thm:v13v73-ground} gives (73.15).
\end{proof}
\end{theorem}

A dense Hilbert-space source set is insufficient here unless it
is a form core.  For an unbounded Hamiltonian, the negative
direction can escape in form norm while all fixed bounded-energy
vectors appear harmless.

\subsection{Finite form screens and a negative-energy witness}

Let \(B_R:E_R\to\mathcal Q(H)\) be finite source maps whose
ranges are cofinal in the form norm.  Using the polarized closed
form \(q_H(\,\cdot\,,\,\cdot\,)\), define the shifted energy Gram
by
\[
 \langle c,\mathbb G_R^{\rm grd}d\rangle
 :=q_H(B_Rc,B_Rd)-E_\Omega\langle B_Rc,B_Rd\rangle.           \tag{73.17}
\]
This Gram is Hermitian but need not be positive, and it is defined
without requiring \(\operatorname{Ran}B_R\subset\operatorname{Dom}H\).

\begin{proposition}[Cofinal ground certification]
\label{prop:v13v73-cofinal-ground}
If
\[
 \mathbb G_R^{\rm grd}\succeq0
 \quad\text{for every }R                                    \tag{73.18}
\]
and \(\bigcup_R\operatorname{Ran}B_R\) is a form core, then
\(\Omega\) is a ground vector.  Conversely, if \(\Omega\) is not
a ground vector, some sufficiently rich form-core screen has a
strictly negative Rayleigh direction.

\begin{proof}
Condition (73.18) says (73.16) on the union of the source ranges.
Form-core continuity extends it to the full form domain, and
Theorem~\ref{thm:v13v73-form-core} proves the first claim.

If \(\Omega\) is not a ground vector, the spectral theorem gives
a form-domain vector \(\psi\) with
\(q_\Omega[\psi]<0\).  Form-core approximation preserves this
strict inequality for all sufficiently close approximants.
One such approximant lies in a finite source range and supplies
a negative Rayleigh direction.
\end{proof}
\end{proposition}

For regulated approximants, a useful normalized debit is
\[
 \nu_{n,R}
 :=\max\left\{0,
 -\lambda_{\min}\!\left(
 G_{B,n,R}^{-1/2}
 \mathbb G_{n,R}^{\rm grd}
 G_{B,n,R}^{-1/2}\right)\right\},                            \tag{73.19}
\]
where \(G_{B,n,R}=B_{n,R}^*B_{n,R}\) has a protected source
floor.  The order of limits is fixed-screen \(n\to\infty\)
followed by \(R\to\infty\).  A moving lowest eigenvector without
form-core transport is not a continuum ground certificate.

\subsection{Cocycle connection form}

When the metric-base action is differentiable, its infinitesimal
contribution is not merely a scalar Koopman derivative.  In a
local trivialization it has the schematic form
\[
 K_{\rm base}
 =-i\,\mathcal L_v+\mathcal A_v,                             \tag{73.20}
\]
where \(v\) is the base flow and \(\mathcal A_v\) is the
self-adjoint connection term obtained by differentiating the
unitary cocycle.  A change of measurable fibre frame transforms
the two terms together.  Therefore the quantities to certify are
the invariant total form \(q_{\rm base}\), its relative bound,
and the full generator spectrum; separate coordinate bounds on
\(-i\mathcal L_v\) and \(\mathcal A_v\) have no frame-independent
meaning.

A finite coupled acquisition must consequently retain:

\begin{enumerate}
\item a common form domain for fibre energy and base motion;
\item an exact cocycle derivative including the connection term;
\item the two-sided relative form constants \(a,b\) in (73.2);
\item the one-sided compensation residual in (73.4);
\item an absolute vacuum-energy eigenvalue \(E_\Omega\);
\item a cofinal form-core ground Gram (73.17), rather than only
commutator or Hilbert-norm source data;
\item compatibility with the V64 forward-cone test for the full
spacetime generator.
\end{enumerate}

\begin{assessment}[V73 advance]
\label{assess:v13v73}
The base-motion obstruction of V72 now has a rigorous sufficient
resolution.  A relatively form-bounded cocycle/base generator
forms a self-adjoint coupled Hamiltonian; a one-sided relative
floor makes it nonnegative.  If only semiboundedness is known,
vacuum invariance and positive energy coincide exactly when the
anchored vacuum eigenvalue is the spectral bottom.  Cofinal
form-core Grams give a finite-screen test of that ground
property.
\end{assessment}

\begin{firewall}[Claim boundary after V73]
\label{fire:v13v73}
No differentiable metric-base flow, cocycle connection,
common coupled form domain, relative form bound, compensation
floor, or form-core ground Gram has been populated.  The
time-translation decomposition (73.1) is a target construction,
not a property already proved for the regulator.  No global
forward-cone spectrum, interacting vacuum, quantum Einstein
dynamics, or full SM--Einstein continuum is proved.
\end{firewall}

\subsection{Parent record}

This V73 note appends the preceding material to the validated V72
source, whose record was
\[
\begin{array}{c|c}
 \text{lines}&23130\\
 \text{bytes}&816850\\
 \text{SHA--256}&
 \texttt{DB4B9AC078AB39E73F5EC1D76E2D45EC309E8BF1FB54A5D9FF569B47E74B5EA4}.
\end{array}                                                     \tag{73.21}
\]

% =============================================================================
% END APPENDED THIRTEENTH-CYCLE V73 MATERIAL
% =============================================================================

% =============================================================================
% BEGIN APPENDED THIRTEENTH-CYCLE V74 MATERIAL
% =============================================================================

\section{Thirteenth-cycle V74: differentiating the metric-fibre
cocycle}
\label{sec:v13v74}

\subsection{A local trivialization of the metric fibres}

V72 identifies a unitary cocycle over the metric law as the
missing datum for global covariance.  V73 then treats its
infinitesimal base contribution as a quadratic form.  The present
note proves the link between those two descriptions.

Let \(X\) be a compact smooth manifold carrying a smooth
probability density \(\mu\), and let \(\varphi_t\) be a complete
smooth flow preserving \(\mu\), with vector field \(v\).  Work
first in a trivial measurable Hilbert bundle
\(X\times\mathcal K\), where \(\mathcal K\) is separable.  Let
\[
 A:X\longrightarrow B(\mathcal K)_{\rm sa}                  \tag{74.1}
\]
be bounded and strongly continuous.  Define \(V(t,x)\) by
\[
 {d\over dt}V(t,x)
 =-iA(\varphi_t x)V(t,x),\qquad V(0,x)=I.                    \tag{74.2}
\]

\begin{lemma}[Ordered exponential is a unitary cocycle]
\label{lem:v13v74-cocycle}
Equation (74.2) has a unique strongly differentiable solution,
and
\[
 V(t+s,x)=V(t,\varphi_sx)V(s,x),                             \tag{74.3}
\]
\[
 V(t,x)^*V(t,x)=I.                                           \tag{74.4}
\]

\begin{proof}
Bounded strong continuity gives the unique Dyson solution of
(74.2) on every compact time interval.  For fixed \(s,x\), both
\[
 t\longmapsto V(t+s,x)
 \quad\text{and}\quad
 t\longmapsto V(t,\varphi_sx)V(s,x)
\]
solve the same differential equation
\[
 \dot W(t)=-iA(\varphi_{t+s}x)W(t)
\]
and have the same value \(V(s,x)\) at \(t=0\).  Uniqueness proves
(74.3).  Since \(A=A^*\),
\[
 {d\over dt}(V^*V)
 =iV^*AV-iV^*AV=0,                                          \tag{74.5}
\]
which proves (74.4).
\end{proof}
\end{lemma}

\subsection{The global generator}

On
\[
 \mathcal H=L^2(X,\mu;\mathcal K),                            \tag{74.6}
\]
define the direct-integral transport
\[
 (U_t\psi)(y)
 :=V(t,\varphi_{-t}y)\psi(\varphi_{-t}y).                    \tag{74.7}
\]
Lemma~\ref{lem:v13v74-cocycle} and
Theorem~\ref{thm:v13v72-cocycle} make \(U_t\) a strongly
continuous unitary group.

Let \(K_0\) be the Stone generator of the Koopman group
\[
 (U_t^0\psi)(y)=\psi(\varphi_{-t}y).                          \tag{74.8}
\]
On smooth cylinder sections,
\[
 K_0=-i\mathcal L_v.                                         \tag{74.9}
\]

\begin{theorem}[Cocycle connection generator]
\label{thm:v13v74-generator}
Assume a smooth cylinder space \(\mathcal C\) is a core for
\(K_0\) and is preserved by multiplication by \(A\).  Then the
Stone generator of (74.7) is
\[
 \boxed{\quad
 K_{\rm base}=K_0+M_A=-i\mathcal L_v+M_A,
 \quad}                                                       \tag{74.10}
\]
self-adjoint on \(\operatorname{Dom}K_0\).

\begin{proof}
Invariance of \(\mu\) makes the Koopman group unitary, so \(K_0\)
is self-adjoint.  The multiplication operator \(M_A\) is bounded
and self-adjoint.  Therefore \(K_0+M_A\) is self-adjoint on
\(\operatorname{Dom}K_0\), and \(\mathcal C\) remains a core.

For \(\psi\in\mathcal C\), differentiate (74.7) at \(t=0\):
\[
 {d\over dt}U_t\psi\big|_{t=0}
 =-iM_A\psi-\mathcal L_v\psi
 =-i(K_0+M_A)\psi.                                           \tag{74.11}
\]
The Stone generator of \(U_t\) agrees with \(K_0+M_A\) on the
core \(\mathcal C\).  Self-adjointness and uniqueness of the
self-adjoint closure prove (74.10).
\end{proof}
\end{theorem}

Thus the schematic connection term in (73.20) is not optional:
it is the infinitesimal derivative of the same cocycle that
implements finite metric transport.

\subsection{Change of fibre frame}

Let
\[
 G:X\longrightarrow U(\mathcal K)                            \tag{74.12}
\]
be strongly differentiable along \(v\), with \(vG\) bounded.
The corresponding multiplication operator \(M_G\) is unitary on
\(\mathcal H\).  In the new frame the cocycle is
\[
 V'(t,x)
 :=G(\varphi_tx)V(t,x)G(x)^*.                                \tag{74.13}
\]

\begin{theorem}[Connection transformation and frame covariance]
\label{thm:v13v74-frame}
The cocycle \(V'\) satisfies (74.3), and its connection is
\[
 \boxed{\quad
 A'(x)=G(x)A(x)G(x)^*
       +i(vG)(x)G(x)^*.
 \quad}                                                       \tag{74.14}
\]
Moreover,
\[
 K'_{\rm base}
 =-i\mathcal L_v+M_{A'}
 =M_GK_{\rm base}M_G^*.                                     \tag{74.15}
\]
In particular, the spectrum and every relative quadratic-form
bound for the total base generator are invariant under a change
of fibre frame, when the fibre Hamiltonian is transformed by the
same \(M_G\).

\begin{proof}
Substitution of (74.13) into both sides of (74.3), followed by
cancellation of
\(G(\varphi_sx)^*G(\varphi_sx)\), proves the cocycle identity.
Differentiating (74.13) at zero gives
\[
 {d\over dt}V'(t,x)\big|_{t=0}
 =(vG)(x)G(x)^*-iG(x)A(x)G(x)^*
 =-iA'(x),                                                   \tag{74.16}
\]
which proves (74.14).  The derivative of \(GG^*=I\) shows that
\((vG)G^*\) is skew-adjoint, so \(A'\) is self-adjoint.

For a smooth section,
\[
\begin{split}
 M_G(-i\mathcal L_v+M_A)M_G^*
 &=-i\mathcal L_v
   -iM_{G(vG^*)}+M_{GAG^*}\\
 &=-i\mathcal L_v
   +M_{\,i(vG)G^*+GAG^*},
\end{split}                                                   \tag{74.17}
\]
where \(G(vG^*)=-(vG)G^*\).  This is (74.15).
Unitary conjugation preserves spectra and quadratic-form
inequalities.
\end{proof}
\end{theorem}

Separate estimates on \(-i\mathcal L_v\) and \(A\) are
frame-dependent.  The compensation inequality in V73 must be
applied to their sum after the same frame has also been used for
the fibre energy.

\subsection{Why an infinitesimal connection is insufficient}

\begin{proposition}[Equal first derivatives do not imply a
cocycle]
\label{prop:v13v74-infinitesimal-no-go}
A unitary family can have the correct value and derivative at
zero while failing the finite composition law.

\begin{proof}
On a one-point base and a nonzero bounded self-adjoint \(B\), let
\[
 V_1(t)=e^{-itA},\qquad
 V_2(t)=e^{-itA-it^2B},                                      \tag{74.18}
\]
where \(A\) and \(B\) commute.  Both are unitary and satisfy
\[
 V_j(0)=I,\qquad \dot V_j(0)=-iA.                            \tag{74.19}
\]
But for \(st\neq0\),
\[
 V_2(s)V_2(t)
 =e^{-i(s+t)A-i(s^2+t^2)B}
 \neq e^{-i(s+t)A-i(s+t)^2B}
 =V_2(s+t).                                                   \tag{74.20}
\]
Thus the infinitesimal row does not certify a group or cocycle.
\end{proof}
\end{proposition}

At finite cutoff the acquisition must retain both residuals
\[
 \varepsilon^{\rm coc}_n(s,t;x)
 :=
 \|V_n(s,\varphi_t x)V_n(t,x)-V_n(s+t,x)\|,                  \tag{74.21}
\]
\[
 \varepsilon^{\rm gen}_n(t;x)
 :=
 \left\|
 {V_n(t,x)-I\over t}+iA_n(x)
 \right\|.                                                    \tag{74.22}
\]
The first is a finite-time associativity row; the second is a
generator row.  Neither controls the other without a uniform
time-regularity estimate.

\subsection{Landing the global unitary group}

\begin{theorem}[Strong-resolvent cocycle landing]
\label{thm:v13v74-resolvent}
Let \(K_n,K\) be self-adjoint operators on one global Hilbert
space and suppose
\[
 K_n\longrightarrow K
 \quad\text{in the strong-resolvent sense}.                  \tag{74.23}
\]
Then for every \(\psi\in\mathcal H\) and \(T<\infty\),
\[
 \sup_{|t|\leq T}
 \|e^{-itK_n}\psi-e^{-itK}\psi\|
 \longrightarrow0.                                          \tag{74.24}
\]
Consequently, if the finite groups arise from exact cocycles and
the limiting representation retains the decomposable local
covariance identities on a dense ledger, those identities pass
to the limiting unitary group.

\begin{proof}
Strong-resolvent convergence gives strong convergence of the
bounded continuous functional calculus after a common spectral
cutoff.  For a fixed vector, the associated spectral measures
are tight because their vague limit has the same total mass.
On a compact spectral interval the functions
\(\lambda\mapsto e^{-it\lambda}\), \(|t|\leq T\), form an
equicontinuous bounded family and can be uniformly approximated
by finitely many functional-calculus test functions.  The
spectral tails are uniformly small, proving (74.24).  Covariance
on a dense ledger then passes by strong convergence, uniform
operator bounds, and closure.
\end{proof}
\end{theorem}

Mosco convergence of the closed form sums in V73 is one
sufficient route to (74.23).  Convergence of the connection
coefficients only pointwise on the metric base is not sufficient;
the common form domain and relative bounds remain necessary.

\subsection{Next finite and continuum rows}

A genuine coupled implementation now needs:

\begin{enumerate}
\item a metric-base flow preserving or quasi-preserving the
selected background law;
\item finite unitary transporters with a controlled cocycle
residual (74.21);
\item an adjacent-time derivative with residual (74.22);
\item a common fibre frame for the connection and internal
Hamiltonian forms;
\item a frame-covariant relative form bound for the total
operator (74.10);
\item Mosco or strong-resolvent convergence on the global direct
integral;
\item source-complete covariance checks for the local observable
ledger after the global group has landed.
\end{enumerate}

\begin{assessment}[V74 advance]
\label{assess:v13v74}
The finite metric-fibre cocycle and the V73 base-motion form are
now linked rigorously.  Differentiating the cocycle gives the
Koopman generator plus a connection multiplication operator.
The inhomogeneous connection transformation is exactly cancelled
by the differentiated fibre frame, so the total generator and
its form bounds are invariant.  Equal infinitesimal data are
proved insufficient without the finite cocycle law, and
strong-resolvent convergence supplies a precise landing theorem
for the global unitary groups.
\end{assessment}

\begin{firewall}[Claim boundary after V74]
\label{fire:v13v74}
No metric-base flow, finite SM fibre cocycle, cocycle derivative,
connection form, invariant background law, common form domain, or
strong-resolvent limit has been populated.  The compact
finite-dimensional connection model does not establish the
unbounded field-theory cocycle.  No global positive-energy
coupled group or full SM--Einstein continuum is proved.
\end{firewall}

\subsection{Parent record}

This V74 note appends the preceding material to the validated V73
source, whose record was
\[
\begin{array}{c|c}
 \text{lines}&23434\\
 \text{bytes}&827653\\
 \text{SHA--256}&
 \texttt{1FFDE23F83BF11E2BA6A84DA39A22FD0D12E311F265F33E80435B734B501C1F8}.
\end{array}                                                     \tag{74.25}
\]

% =============================================================================
% END APPENDED THIRTEENTH-CYCLE V74 MATERIAL
% =============================================================================

% =============================================================================
% BEGIN APPENDED THIRTEENTH-CYCLE V75 MATERIAL
% =============================================================================

\section{Thirteenth-cycle V75: compulsory companion audit before
a global metric-base form cover}
\label{sec:v13v75}

\subsection{Files inspected}

The newest complete companion notes in the shared directory were
\[
\begin{array}{c|c|r|r|c}
 \text{programme}&\text{version}&\text{lines}&\text{bytes}&
 \text{SHA--256}\\ \hline
 \text{Yang--Mills mass gap}&V95&30950&1056018&
 \texttt{2DB19DAF571B44B04B28A51FAE2708D87BB6234D9CE9A2FBC5530ADF081809C2}\\
 \text{Navier--Stokes stability}&V26&5319&278283&
 \texttt{82C887F8C2C69E4F28FAD7C3DC8DE5B9099CB5A4BF431EBA1FFF3E91935978AD}.
\end{array}                                                     \tag{75.1}
\]
The auxiliary YM parallel note remained V5 with SHA--256
\[
 \texttt{306F1BB84CFA8A8D47EA569EC17E9545F9867C8D32B548EC9D9C444B4F3C0512}.
                                                                  \tag{75.2}
\]

\subsection{New companion conclusion}

YM V95 computes exact normalized covariance jets at every
quarter-turn anchor.  Its second-order grid bounds improve the
previous global coefficient estimates by factors greater than
\(37\) and \(397\), but they remain discrete.  The note therefore
refuses to substitute them into a global response theorem until a
certified cell extension has been proved.  If absolute-radius
transport remains too lossy, it will test a relative Loewner
inequality directly.

This distinction transfers exactly to V73--V74.  A cocycle
generator and a compensation inequality at finitely many metric
backgrounds do not give a global direct-integral Hamiltonian.
The background cells need relative form variation stocks or
direct Loewner certificates.  Absolute operator bounds may be
far too coarse because the fibre Hamiltonian itself supplies the
natural scale.

NS V26 and auxiliary YM V5 are unchanged.  Their numerical
constants are not imported.

\begin{proposition}[V75 direction decision]
\label{prop:v13v75-decision}
The next note shall prove a finite-cover theorem transporting
anchor relative-form bounds over the metric base.  It must retain:

\begin{enumerate}
\item one common form domain and reference form;
\item a cellwise lower comparison of the fibre form with the
reference;
\item separate fibre-form and base-form variation radii;
\item both an absolute relative bound for self-adjointness and a
sharper one-sided bound for positivity;
\item the accumulated scalar debit;
\item a counterexample showing that exact anchor values without a
variation stock give no global form control.
\end{enumerate}

\begin{proof}
The KLMN and positivity conclusions of
Theorem~\ref{thm:v13v73-form-sum} require inequalities for almost
every point of the metric law.  Finite anchors cover that law only
after their estimates are transported throughout a finite
background cover.  A reference-form lower comparison converts
absolute cell variation into a relative debit measured in units
of the local fibre energy.  This is the form-theoretic counterpart
of the V95 proposed boxwise Loewner step.
\end{proof}
\end{proposition}

\begin{assessment}[V75 audit outcome]
\label{assess:v13v75}
No companion numerical bound transfers.  The audit identifies the
next noncircular task: extend the V73 compensation estimate over
the entire admitted metric base by relative form geometry, rather
than treating sharp anchor values as a global theorem.
\end{assessment}

\begin{firewall}[Claim boundary after V75]
\label{fire:v13v75}
This audit supplies no metric-cell variation stock, Loewner
certificate, global form sum, positive generator, or coupled
continuum state.  The full SM--Einstein continuum remains open.
\end{firewall}

\subsection{Parent record}

This V75 note appends the preceding material to the validated V74
source, whose record was
\[
\begin{array}{c|c}
 \text{lines}&23769\\
 \text{bytes}&839080\\
 \text{SHA--256}&
 \texttt{8931C564213C31C2287328DAF396C42DE3BCC9F8D4670531F6C410462FEAFD36}.
\end{array}                                                     \tag{75.3}
\]

% =============================================================================
% END APPENDED THIRTEENTH-CYCLE V75 MATERIAL
% =============================================================================

% =============================================================================
% BEGIN APPENDED THIRTEENTH-CYCLE V76 MATERIAL
% =============================================================================

\section{Thirteenth-cycle V76: a finite relative-form cover of
the metric base}
\label{sec:v13v76}

\subsection{Common form geometry}

Let \(X\) be a compact metric background space covered by finitely
many cells \(C_1,\ldots,C_N\), with anchor
\(x_j\in C_j\).  After the fibre whitening of V72--V74, suppose
all fibre and base forms have one dense common domain
\(\mathcal Q\subset\mathcal K\):
\[
 q_{{\rm fib},x},\qquad q_{{\rm base},x},
 \qquad x\in X.                                               \tag{76.1}
\]
Every fibre form is closed and nonnegative, every base form is
symmetric, and the measurable direct integral of the fibre forms is
assumed closed.
Let \(q_*\) be a fixed closed nonnegative reference form on
\(\mathcal Q\).  In cell \(C_j\), assume the fibre floor
\[
 q_{{\rm fib},x}[\psi]\geq
 m_jq_*[\psi],\qquad m_j>0,                                  \tag{76.2}
\]
and the two variation estimates
\[
 \left|q_{{\rm fib},x}[\psi]
       -q_{{\rm fib},x_j}[\psi]\right|
 \leq\delta_{f,j}q_*[\psi]+c_{f,j}\|\psi\|^2,                \tag{76.3}
\]
\[
 \left|q_{{\rm base},x}[\psi]
       -q_{{\rm base},x_j}[\psi]\right|
 \leq\delta_{b,j}q_*[\psi]+c_{b,j}\|\psi\|^2.                \tag{76.4}
\]
All constants are nonnegative.

At the anchor retain two different certificates:
\[
 |q_{{\rm base},x_j}[\psi]|
 \leq\alpha_jq_{{\rm fib},x_j}[\psi]
      +\beta_j\|\psi\|^2,                                    \tag{76.5}
\]
\[
 q_{{\rm base},x_j}[\psi]
 \geq-a_jq_{{\rm fib},x_j}[\psi]
      -b_j\|\psi\|^2.                                        \tag{76.6}
\]
The absolute estimate (76.5) is used for KLMN
self-adjointness.  The signed estimate (76.6) may be sharper and
is used for positivity.

\subsection{Cell transport}

\begin{theorem}[Relative-form anchor-cell theorem]
\label{thm:v13v76-cell}
For every \(x\in C_j\) and \(\psi\in\mathcal Q\),
\[
 |q_{{\rm base},x}[\psi]|
 \leq A_jq_{{\rm fib},x}[\psi]+B_j\|\psi\|^2,                \tag{76.7}
\]
where
\[
 \boxed{
 \begin{aligned}
 A_j&:=\alpha_j+
       {\delta_{b,j}+\alpha_j\delta_{f,j}\over m_j},\\
 B_j&:=\beta_j+c_{b,j}+\alpha_jc_{f,j},
 \end{aligned}}                                              \tag{76.8}
\]
and
\[
 q_{{\rm base},x}[\psi]
 \geq-C_jq_{{\rm fib},x}[\psi]-D_j\|\psi\|^2,                \tag{76.9}
\]
where
\[
 \boxed{
 \begin{aligned}
 C_j&:=a_j+
       {\delta_{b,j}+a_j\delta_{f,j}\over m_j},\\
 D_j&:=b_j+c_{b,j}+a_jc_{f,j}.
 \end{aligned}}                                              \tag{76.10}
\]

\begin{proof}
Equations (76.2)--(76.3) give
\[
\begin{split}
 q_{{\rm fib},x_j}[\psi]
 &\leq q_{{\rm fib},x}[\psi]
      +\delta_{f,j}q_*[\psi]+c_{f,j}\|\psi\|^2\\
 &\leq
 \left(1+{\delta_{f,j}\over m_j}\right)
 q_{{\rm fib},x}[\psi]+c_{f,j}\|\psi\|^2.                   \tag{76.11}
\end{split}
\]
Using (76.4), (76.5), and then (76.11),
\[
\begin{split}
 |q_{{\rm base},x}[\psi]|
 &\leq
 \alpha_jq_{{\rm fib},x_j}[\psi]
 +\delta_{b,j}q_*[\psi]
 +(\beta_j+c_{b,j})\|\psi\|^2\\
 &\leq
 \left(\alpha_j+
 {\,\delta_{b,j}+\alpha_j\delta_{f,j}\over m_j}\right)
 q_{{\rm fib},x}[\psi]\\
 &\quad+
 (\beta_j+c_{b,j}+\alpha_jc_{f,j})\|\psi\|^2,
\end{split}                                                   \tag{76.12}
\]
which proves (76.7)--(76.8).

For the signed bound,
\[
\begin{split}
 q_{{\rm base},x}[\psi]
 &\geq q_{{\rm base},x_j}[\psi]
       -\delta_{b,j}q_*[\psi]-c_{b,j}\|\psi\|^2\\
 &\geq-a_jq_{{\rm fib},x_j}[\psi]
       -\delta_{b,j}q_*[\psi]
       -(b_j+c_{b,j})\|\psi\|^2.
\end{split}                                                   \tag{76.13}
\]
Insert (76.11) and then (76.2).  The result is exactly
(76.9)--(76.10).
\end{proof}
\end{theorem}

The loss is now localized in four distinct places: the anchor
relative constants, the reference floor \(m_j\), the fibre
variation, and the base variation.  Refining a cell is useful
only when it improves the factor that dominates (76.8) or
(76.10).

\subsection{Global direct-integral conclusion}

Assume \(x\mapsto q_{{\rm fib},x}\) and
\(x\mapsto q_{{\rm base},x}\) are measurable form fields for a
probability law \(\mu\) on \(X\).  Define
\[
 Q_{\rm fib}[\Psi]
 :=\int_Xq_{{\rm fib},x}[\Psi_x]\,d\mu(x),\qquad
 Q_{\rm base}[\Psi]
 :=\int_Xq_{{\rm base},x}[\Psi_x]\,d\mu(x).                  \tag{76.14}
\]
Put
\[
 A:=\max_jA_j,\quad B:=\max_jB_j,\quad
 C:=\max_jC_j,\quad D:=\max_jD_j.                            \tag{76.15}
\]

\begin{corollary}[Finite-cover global form sum]
\label{cor:v13v76-global}
The pointwise cell estimates imply
\[
 |Q_{\rm base}[\Psi]|
 \leq A Q_{\rm fib}[\Psi]+B\|\Psi\|^2,                       \tag{76.16}
\]
\[
 Q_{\rm base}[\Psi]
 \geq-C Q_{\rm fib}[\Psi]-D\|\Psi\|^2.                       \tag{76.17}
\]
If \(A<1\), the sum
\[
 Q_{\rm tot}=Q_{\rm fib}+Q_{\rm base}                        \tag{76.18}
\]
is closed and defines a self-adjoint operator satisfying
\[
 H_{\rm tot}\geq-BI.                                         \tag{76.19}
\]
If also \(C\leq1\), the sharper signed estimate gives
\(H_{\rm tot}\geq-DI\); if in addition \(D=0\), then
\(H_{\rm tot}\geq0\).

\begin{proof}
For almost every \(x\), choose any cell containing \(x\).
The corresponding estimates are bounded by the maxima in
(76.15).  Integration gives (76.16)--(76.17).  The first and
\(A<1\) give the KLMN conclusion and
\[
 Q_{\rm tot}[\Psi]
 \geq(1-A)Q_{\rm fib}[\Psi]-B\|\Psi\|^2,
\]
which proves (76.19).  The signed estimate gives
\[
 Q_{\rm tot}[\Psi]
 \geq(1-C)Q_{\rm fib}[\Psi]-D\|\Psi\|^2.                     \tag{76.20}
\]
When \(C\leq1\), this proves the sharper lower bound by
\(-D\), and \(D=0\) proves positivity.
\end{proof}
\end{corollary}

When \(C\leq1\) and \(D>0\), the signed result is semibounded rather than
positive-energy.  The V73 ground anchor must then prove that the
chosen invariant vacuum realizes the bottom of the coupled
spectrum.  The scalar debit cannot be discarded as a
normal-ordering convention before that test.

\subsection{Direct Loewner alternative}

In a finite-dimensional screen, let \(F_x\succ0\) and
\(B_x=B_x^*\) represent the fibre and base forms.  Then
\[
 B_x\succeq-cF_x                                             \tag{76.21}
\]
is equivalent to
\[
 \lambda_{\min}
 \bigl(F_x^{-1/2}B_xF_x^{-1/2}\bigr)\geq-c.                  \tag{76.22}
\]
Thus a verified interval enclosure of the generalized minimum
eigenvalue on a cell can replace the separate variation estimates
(76.2)--(76.4).  It may be much sharper when the leading
variations of \(F_x\) and \(B_x\) cancel in relative coordinates.

\begin{proposition}[Loewner cell certificate]
\label{prop:v13v76-loewner}
If interval or exact arithmetic proves
\[
 F_x\succeq mI,\qquad B_x+cF_x\succeq0
 \quad\text{for every }x\in C_j,                             \tag{76.23}
\]
then the signed relative bound (76.9) holds throughout the cell
with \(C_j=c\) and \(D_j=0\), independently of the separate
absolute radii of \(F_x\) and \(B_x\).

\begin{proof}
The second matrix inequality in (76.23), tested on an arbitrary
coefficient vector, is exactly
\(q_{{\rm base},x}\geq-cq_{{\rm fib},x}\).  The first inequality
ensures that the whitening and generalized eigenvalue in
(76.22) are defined.
\end{proof}
\end{proposition}

This is the metric-base analogue of the boxwise relative
Loewner route proposed by YM V95.

\subsection{Anchor values alone give no global control}

\begin{proposition}[Finite-anchor no-go]
\label{prop:v13v76-anchor-no-go}
For any finite set of anchors in \([0,1]\) and any \(L>0\), there
are continuous scalar forms with
\[
 q_{{\rm fib},x}[\psi]=|\psi|^2,\qquad
 q_{{\rm base},x_j}[\psi]=0
\]
at every anchor, but
\[
 \inf_{x\in[0,1]}
 {q_{{\rm base},x}[\psi]\over
  q_{{\rm fib},x}[\psi]}
 =-L
 \quad(\psi\neq0).                                           \tag{76.24}
\]

\begin{proof}
For anchors \(x_1,\ldots,x_m\), let
\[
 p(x):=\prod_{j=1}^m(x-x_j)^2.                               \tag{76.25}
\]
This continuous nonnegative function is not identically zero.
Set
\[
 q_{{\rm base},x}[\psi]
 :=-L\,{p(x)\over\max_{[0,1]}p}\,|\psi|^2.                   \tag{76.26}
\]
It vanishes at every anchor and reaches relative value \(-L\)
at a maximizer of \(p\).
\end{proof}
\end{proposition}

No number of exact calculations at a fixed finite anchor set can
therefore establish the global compensation inequality without
a cell variation or direct interval certificate.

\subsection{Continuum preservation of positivity}

\begin{theorem}[Vanishing scalar debit under Mosco convergence]
\label{thm:v13v76-mosco}
Let \(Q_n\) be closed lower-semibounded forms on Hilbert spaces
identified with one carrier, and suppose
\[
 Q_n[\Psi]\geq-d_n\|\Psi\|^2,\qquad d_n\downarrow0.           \tag{76.27}
\]
If \(Q_n\) Mosco-converges to a closed form \(Q_\infty\), then
\[
 Q_\infty[\Psi]\geq0
 \quad\text{on }\operatorname{Dom}Q_\infty.                  \tag{76.28}
\]
The associated self-adjoint operators converge in the strong
resolvent sense, and the limiting generator is nonnegative.

\begin{proof}
For \(\Psi\in\operatorname{Dom}Q_\infty\), the Mosco recovery
condition supplies \(\Psi_n\to\Psi\) with
\[
 \limsup_n Q_n[\Psi_n]\leq Q_\infty[\Psi].                   \tag{76.29}
\]
But (76.27) gives
\[
 \liminf_n Q_n[\Psi_n]
 \geq-\lim_n d_n\|\Psi_n\|^2=0.                              \tag{76.30}
\]
Hence \(Q_\infty[\Psi]\geq0\).  The strong-resolvent conclusion
is the standard representation theorem for Mosco-convergent
closed forms; nonnegativity follows from (76.28).
\end{proof}
\end{theorem}

Combining this theorem with V74 gives uniform-on-compact-time
strong convergence of the coupled unitary groups.  It still does
not supply the forward spatial-momentum cone, which remains the
joint support test of V64.

\subsection{Acquisition consequence}

For every metric cell and cutoff, the ledger must store
\[
 (m_j,\delta_{f,j},c_{f,j},
   \delta_{b,j},c_{b,j},
   \alpha_j,\beta_j,a_j,b_j),                                \tag{76.31}
\]
or a direct relative Loewner certificate.  The finite cover is
accepted only if
\[
 \max_jA_j<1                                                 \tag{76.32}
\]
and either
\[
 \max_jC_j\leq1,\qquad\max_jD_j\to0,                         \tag{76.33}
\]
or the V73 vacuum-ground Grams close the retained scalar debit.
A failed cell is refined at the factor dominating (76.8) or
(76.10), rather than by globally enlarging every constant.

\begin{assessment}[V76 advance]
\label{assess:v13v76}
The pointwise V73 compensation condition now has a rigorous
finite-cover compiler.  Anchor absolute and signed form bounds,
a reference fibre floor, and separate cell variations produce
explicit global KLMN and positivity constants.  Direct Loewner
certification is available when relative cancellation beats
absolute radii.  Exact anchors without a variation stock are
proved insufficient, and a vanishing scalar debit survives a
Mosco continuum limit.
\end{assessment}

\begin{firewall}[Claim boundary after V76]
\label{fire:v13v76}
No metric-cell cover, common coupled form domain, reference-form
floor, fibre/base variation stock, Loewner interval, or Mosco
limit has been populated for the SM--Einstein regulator.  The
global forward cone and vacuum-ground rows remain open.  No full
SM--Einstein continuum is proved.
\end{firewall}

\subsection{Parent record}

This V76 note appends the preceding material to the validated V75
source, whose record was
\[
\begin{array}{c|c}
 \text{lines}&23879\\
 \text{bytes}&843474\\
 \text{SHA--256}&
 \texttt{8C4254B972D5473E8EE99A41EEF5EE746539468B7E79A721FFFB6D0D5DF11642}.
\end{array}                                                     \tag{76.34}
\]

% =============================================================================
% END APPENDED THIRTEENTH-CYCLE V76 MATERIAL
% =============================================================================

% =============================================================================
% BEGIN APPENDED THIRTEENTH-CYCLE V77 MATERIAL
% =============================================================================

\section{Thirteenth-cycle V77: directional form certification of
the forward cone}
\label{sec:v13v77}

\subsection{From energy positivity to the joint cone}

Let \(H,P_1,\ldots,P_d\) be strongly commuting self-adjoint
operators with joint spectral measure \(\mathsf E\).  Put
\[
 |P|:=\left(P_1^2+\cdots+P_d^2\right)^{1/2}                 \tag{77.1}
\]
and use the common form domain
\[
 \mathcal Q_{HP}
 :=\operatorname{Dom}|H|^{1/2}
   \cap\operatorname{Dom}|P|^{1/2}.                          \tag{77.2}
\]
For \(u\in S^{d-1}\), the directional energy form is
\[
 q_u[\psi]
 :=\int_{\mathbb R^{1+d}}(E-u\cdot p)\,
       d\langle\psi,\mathsf E(E,p)\psi\rangle.                \tag{77.3}
\]

\begin{theorem}[Directional characterization of the spectrum
condition]
\label{thm:v13v77-directional}
The following are equivalent:
\[
 \operatorname{sp}(H,P)\subset
 \overline V_+
 =\{(E,p):E\geq|p|\},                                        \tag{77.4}
\]
and
\[
 q_u[\psi]\geq0
 \quad\text{for every }u\in S^{d-1},\
 \psi\in\mathcal Q_{HP}.                                     \tag{77.5}
\]

\begin{proof}
If (77.4) holds, then
\[
 E-u\cdot p\geq E-|p|\geq0
\]
on the joint spectral support, proving (77.5).

Conversely, suppose the joint spectral projection of
\(\mathbb R^{1+d}\setminus\overline V_+\) is nonzero.  The open
sets
\[
 O_u:=\{(E,p):E-u\cdot p<0\},\qquad u\in S^{d-1},             \tag{77.6}
\]
cover that complement: for \(p\neq0\) choose
\(u=p/|p|\), while for \(p=0,E<0\) every direction works.
Second countability gives a countable directional subcover.
Hence some \(O_u\) has nonzero spectral projection.  Intersecting
with a bounded subset on which \(E-u\cdot p\leq-\eta<0\)
produces a nonzero vector \(\psi\in\mathcal Q_{HP}\) with
\(q_u[\psi]<0\), contradicting (77.5).
\end{proof}
\end{theorem}

Thus the V73--V76 conclusion \(H\geq0\) is only the zero-momentum
part of the relativistic spectrum condition.

\subsection{Finite angular screens with a reserve}

Let \(\mathcal N_\delta\subset S^{d-1}\) be a finite
\(\delta\)-net in Euclidean distance.

\begin{theorem}[Angular-reserve cone certificate]
\label{thm:v13v77-angular}
Suppose that for constants \(m,\epsilon\geq0\),
\[
 q_{u_k}[\psi]
 \geq m\,q_{|P|}[\psi]-\epsilon\|\psi\|^2
 \quad
 (u_k\in\mathcal N_\delta,\ \psi\in\mathcal Q_{HP}),          \tag{77.7}
\]
where
\[
 q_{|P|}[\psi]
 :=\int|p|\,d\langle\psi,\mathsf E(E,p)\psi\rangle.
\]
Then every \(u\in S^{d-1}\) satisfies
\[
 \boxed{\quad
 q_u[\psi]
 \geq(m-\delta)q_{|P|}[\psi]
      -\epsilon\|\psi\|^2.
 \quad}                                                       \tag{77.8}
\]
In particular, if \(\epsilon=0\) and \(m\geq\delta\), the full
forward-cone condition holds.

\begin{proof}
Choose \(u_k\in\mathcal N_\delta\) with
\(|u-u_k|\leq\delta\).  Pointwise in momentum,
\[
 |(u-u_k)\cdot p|\leq\delta|p|.                              \tag{77.9}
\]
Therefore
\[
 q_u[\psi]
 =q_{u_k}[\psi]
   -\int(u-u_k)\cdot p\,d\langle\psi,\mathsf E\psi\rangle
 \geq q_{u_k}[\psi]-\delta q_{|P|}[\psi].                    \tag{77.10}
\]
Insert (77.7) to obtain (77.8).  The final statement follows
from Theorem~\ref{thm:v13v77-directional}.
\end{proof}
\end{theorem}

The reserve \(m q_{|P|}\) is an angular interpolation budget.
Massless spectral points on the boundary of the cone generally
leave no fixed positive reserve, so an exact finite net need not
be enough.

\begin{proposition}[A finite zero-reserve screen can miss a
spacelike spectral point]
\label{prop:v13v77-finite-no-go}
For every finite
\(\mathcal N\subset S^{d-1}\) that does not contain every
direction, there is a one-dimensional joint spectral
representation satisfying
\[
 H-u_k\cdot P\geq0\quad(u_k\in\mathcal N)                    \tag{77.11}
\]
whose spectral point lies outside \(\overline V_+\).

\begin{proof}
Choose \(w\in S^{d-1}\setminus\mathcal N\) and set
\[
 p=w,\qquad E:=\max_{u_k\in\mathcal N}u_k\cdot w.             \tag{77.12}
\]
Because the net is finite and no \(u_k=w\), one has \(E<1=|p|\).
On the one-dimensional Hilbert space take \(H=E\) and \(P=p\).
Then \(E-u_k\cdot p\geq0\) for every sampled direction, proving
(77.11), but \(E<|p|\).
\end{proof}
\end{proposition}

\subsection{Cofinal angular landing without a fixed reserve}

A shrinking sequence of angular meshes can recover the exact
cone after the limiting operators have landed.

\begin{theorem}[Cofinal directional criterion]
\label{thm:v13v77-cofinal}
Let \(\delta_n\downarrow0\), let
\(\mathcal N_n\) be finite \(\delta_n\)-nets, and let
\(\epsilon_n\downarrow0\).  Suppose the fixed limiting strongly
commuting family \((H,P)\) satisfies
\[
 q_u[\psi]\geq-\epsilon_n\|\psi\|^2
 \quad
 (u\in\mathcal N_n,\ \psi\in\mathcal Q_{HP})                 \tag{77.13}
\]
for every \(n\).  Then
\[
 \operatorname{sp}(H,P)\subset\overline V_+.                 \tag{77.14}
\]

\begin{proof}
Fix \(u\in S^{d-1}\), \(\psi\in\mathcal Q_{HP}\), and choose
\(u_n\in\mathcal N_n\) with \(|u-u_n|\leq\delta_n\).
As in (77.10),
\[
 q_u[\psi]
 \geq-\epsilon_n\|\psi\|^2
      -\delta_n q_{|P|}[\psi].                               \tag{77.15}
\]
The right side tends to zero because
\(\psi\in\operatorname{Dom}|P|^{1/2}\).  Thus \(q_u[\psi]\geq0\)
for every \(u,\psi\), and
Theorem~\ref{thm:v13v77-directional} proves (77.14).
\end{proof}
\end{theorem}

The order is essential: first obtain one limiting strongly
commuting translation family, then refine its angular screen.
For changing cutoff operators, (77.13) requires convergence of
the directional forms on a common core or convergence of the
joint spectral measures.  Separate strong-resolvent convergence
of \(H_n\) and each \(P_{j,n}\) does not by itself guarantee
strong commutation of the limits.

\subsection{Finite relative-form implementation}

At cutoff \(n\), direction \(u\), and source screen \(R\), define
the Hermitian form Gram
\[
 \langle c,\mathbb G^{\rm cone}_{n,R,u}d\rangle
 :=
 q_{H_n-u\cdot P_n}(B_{n,R}c,B_{n,R}d).                     \tag{77.16}
\]
A rigorous cone ledger must retain:

\begin{enumerate}
\item a joint translation representation, not separately fitted
energy and momentum generators;
\item a common form-core source atlas for \(H_n\) and \(|P_n|\);
\item a finite angular net and either the reserve in (77.7) or a
cofinal sequence as in (77.13);
\item source-Gram floors before generalized eigenvalues are used;
\item convergence of the joint spectral calculus or all
directional forms on the fixed core;
\item the V64 spatial kernel data, which prevent momentum from
being integrated out before this test.
\end{enumerate}

The V76 metric-base cover can be applied separately to each
directional form \(H_x-u\cdot P_x\).  Its cell constants must then
be uniform over the chosen angular screen.  A background cover
for \(H_x\) alone does not control \(P_x\).

\begin{assessment}[V77 advance]
\label{assess:v13v77}
The global spectrum condition is now expressed as an exact family
of directional form inequalities.  A finite angular screen
certifies the whole cone when it carries an explicit
\(|P|\)-reserve; without that reserve a constructive
one-dimensional counterexample escapes between sampled
directions.  A cofinal shrinking angular mesh recovers the exact
cone after the joint limiting translation representation has
landed.
\end{assessment}

\begin{firewall}[Claim boundary after V77]
\label{fire:v13v77}
No coupled spatial-momentum generators, joint spectral measure,
common momentum form core, angular reserve, or cofinal cone Gram
has been populated.  Strong commutation through the continuum
limit remains unproved.  No forward-cone coupled spectrum or full
SM--Einstein continuum is established.
\end{firewall}

\subsection{Parent record}

This V77 note appends the preceding material to the validated V76
source, whose record was
\[
\begin{array}{c|c}
 \text{lines}&24262\\
 \text{bytes}&855900\\
 \text{SHA--256}&
 \texttt{41DA50522D160225AD2675C0A70A5FFD5E57EBAC5C6B82BEB560D9B89F48F90A}.
\end{array}                                                     \tag{77.17}
\]

% =============================================================================
% END APPENDED THIRTEENTH-CYCLE V77 MATERIAL
% =============================================================================

% =============================================================================
% BEGIN APPENDED THIRTEENTH-CYCLE V78 MATERIAL
% =============================================================================

\section{Thirteenth-cycle V78: landing the joint
energy--momentum calculus}
\label{sec:v13v78}

\subsection{The missing jointness condition}

V77 characterizes the forward cone by forms of
\(H-u\cdot P\).  That characterization presupposes a joint
spectral measure.  Componentwise limits must therefore retain
strong commutation.  This note gives two sufficient criteria:
exact jointness at every cutoff, or a vanishing commutator of
bounded resolvents.  It then passes the complete joint
\(C_0\)-functional calculus and closed-cone support.

Let
\[
 T_n=(T_{0,n},\ldots,T_{d,n}),\qquad
 T=(T_0,\ldots,T_d)                                          \tag{78.1}
\]
be self-adjoint tuples on one Hilbert space, with
\(T_{0,n}=H_n\) and \(T_{j,n}=P_{j,n}\).

\subsection{Exact cutoff commutation}

\begin{theorem}[Componentwise strong resolvent convergence
preserves an exact joint tuple]
\label{thm:v13v78-joint}
Assume that the operators in every tuple \(T_n\) strongly
commute and
\[
 T_{j,n}\longrightarrow T_j
 \quad\text{in the strong-resolvent sense}
 \qquad(0\leq j\leq d).                                      \tag{78.2}
\]
Then the limiting operators \(T_0,\ldots,T_d\) strongly commute.

\begin{proof}
Strong-resolvent convergence of self-adjoint operators implies
\[
 e^{itT_{j,n}}\longrightarrow e^{itT_j}
 \quad\text{strongly, locally uniformly in }t.               \tag{78.3}
\]
At every cutoff,
\[
 e^{isT_{j,n}}e^{itT_{k,n}}
 =e^{itT_{k,n}}e^{isT_{j,n}}.                                \tag{78.4}
\]
Products of uniformly bounded strongly convergent operators
converge strongly.  Passing to the limit in (78.4) gives
\[
 e^{isT_j}e^{itT_k}=e^{itT_k}e^{isT_j}
 \quad(s,t\in\mathbb R).                                     \tag{78.5}
\]
Commutation of all spectral unitary groups is equivalent to
commutation of their spectral measures.  Hence the limiting
tuple strongly commutes.
\end{proof}
\end{theorem}

Thus the caution at the end of V77 can be sharpened:
componentwise strong-resolvent convergence is sufficient when
exact strong commutation is retained at every cutoff.  It is
insufficient only when the cutoff construction has not supplied
that jointness.

\subsection{A bounded resolvent residual}

Put
\[
 R_{j,n}:=(T_{j,n}-i)^{-1},\qquad R_j:=(T_j-i)^{-1}.          \tag{78.6}
\]

\begin{theorem}[Vanishing resolvent-commutator criterion]
\label{thm:v13v78-resolvent-commutator}
Assume (78.2) and, for every pair \(j,k\),
\[
 [R_{j,n},R_{k,n}]
 \longrightarrow0\quad\text{strongly}.                      \tag{78.7}
\]
Then the limiting tuple strongly commutes.

\begin{proof}
Strong-resolvent convergence gives
\(R_{j,n}\to R_j\) strongly.  Since
\(\|R_{j,n}\|\leq1\), multiplication is jointly continuous on
these strongly convergent bounded sequences.  Therefore (78.7)
implies
\[
 R_jR_k=R_kR_j.                                               \tag{78.8}
\]
Each \(R_j\) is normal.  Fuglede's theorem upgrades (78.8) to
commutation with adjoints, so the abelian \(C^*\)-algebras
generated by \(R_j\) and \(R_k\) commute.  The spectral
projections of \(T_j\) are Borel functions of \(R_j\); hence the
spectral measures commute.  Applying this to every pair proves
strong commutation of the tuple.
\end{proof}
\end{theorem}

Because resolvents are bounded, (78.7) is a safer continuum row
than a formal commutator \([T_{j,n},T_{k,n}]\) on a moving
unbounded domain.  A finite source compression of (78.7) still
requires a cofinal dense source argument.

\subsection{Joint functional-calculus convergence}

Let \(\mathsf E_n,\mathsf E\) denote the joint spectral measures
given by Theorem~\ref{thm:v13v78-joint} or
\ref{thm:v13v78-resolvent-commutator}.

\begin{theorem}[Strong convergence of the joint
\(C_0\)-calculus]
\label{thm:v13v78-functional}
Under the hypotheses of
Theorem~\ref{thm:v13v78-joint}, for every
\(f\in C_0(\mathbb R^{d+1})\),
\[
 f(T_n):=\int f(\lambda)\,d\mathsf E_n(\lambda)
 \longrightarrow
 f(T):=\int f(\lambda)\,d\mathsf E(\lambda)
 \quad\text{strongly}.                                      \tag{78.9}
\]

\begin{proof}
First take a product function
\[
 f(\lambda_0,\ldots,\lambda_d)
 =\prod_{j=0}^df_j(\lambda_j),
 \qquad f_j\in C_0(\mathbb R).                               \tag{78.10}
\]
One-variable strong-resolvent functional calculus gives
\(f_j(T_{j,n})\to f_j(T_j)\) strongly.  All factors have a common
uniform norm bound, so a telescoping product proves (78.9) for
(78.10) and for finite sums of such products.

The algebraic tensor product
\(\bigotimes_{j=0}^dC_0(\mathbb R)\) is uniformly dense in
\(C_0(\mathbb R^{d+1})\).  Joint functional calculus is
contractive in the supremum norm.  Uniform approximation before
and after taking \(n\to\infty\) proves (78.9) for arbitrary
\(f\).
\end{proof}
\end{theorem}

\subsection{Closed spectral supports survive}

For a closed set \(C\subset\mathbb R^{d+1}\), let
\[
 C_\epsilon:=\{\lambda:
 \operatorname{dist}(\lambda,C)\leq\epsilon\}.               \tag{78.11}
\]

\begin{corollary}[Approximate closed-support landing]
\label{cor:v13v78-support}
Suppose, in addition to the hypotheses of
Theorem~\ref{thm:v13v78-functional}, that
\[
 \operatorname{sp}(T_n)\subset C_{\epsilon_n},
 \qquad\epsilon_n\downarrow0.                                \tag{78.12}
\]
Then
\[
 \operatorname{sp}(T)\subset C.                             \tag{78.13}
\]
In particular, taking \(C=\overline V_+\) passes the
forward-cone spectrum condition to the limit.

\begin{proof}
Let \(f\in C_c(\mathbb R^{d+1}\setminus C)\).  Compactness of its
support and closedness of \(C\) give a positive separation.
For all sufficiently large \(n\), (78.12) implies
\(f(T_n)=0\).  Theorem~\ref{thm:v13v78-functional} gives
\(f(T)=0\).  Such compactly supported functions exhaust the
complement of \(C\), so the joint spectral projection there is
zero, proving (78.13).
\end{proof}
\end{corollary}

This support theorem is stronger than checking only
\(H_n\geq0\).  It also avoids attempting to pass discontinuous
spectral projections directly at a boundary carrying spectral
mass.

\subsection{Cofinal source version}

Let \(\mathcal D_0\) be dense and suppose it is the union of a
directed family of finite source ranges.  Choose a countable
dense set
\[
 \mathcal F_C\subset
 C_c(\mathbb R^{d+1}\setminus C)                             \tag{78.14}
\]
in the uniform norm on each compact exhaustion.

\begin{proposition}[Source-complete spectral-support test]
\label{prop:v13v78-source}
Assume the limiting joint tuple exists and, for every
\(f\in\mathcal F_C\) and \(\psi\in\mathcal D_0\),
\[
 \|f(T)\psi\|=0.                                             \tag{78.15}
\]
Then \(\operatorname{sp}(T)\subset C\).  It is enough to obtain
(78.15) as the limit of finite-cutoff source rows
\[
 \|f(T_n)B_{n,R}c\|\longrightarrow0                          \tag{78.16}
\]
when the transported sources converge and
Theorem~\ref{thm:v13v78-functional} applies.

\begin{proof}
Each bounded operator \(f(T)\) vanishes on the dense set
\(\mathcal D_0\), hence vanishes identically.  Uniform density of
\(\mathcal F_C\) and contractivity of the functional calculus
extend this to all compactly supported functions off \(C\).
The proof of Corollary~\ref{cor:v13v78-support} then gives the
support claim.  Under the convergence assumptions, (78.16) and
(78.9) yield (78.15) on every fixed transported source.
\end{proof}
\end{proposition}

For positive source Grams, (78.16) is equivalently tested by
\[
 \mathbb G^{f}_{n,R}
 :=B_{n,R}^*f(T_n)^*f(T_n)B_{n,R}\longrightarrow0.           \tag{78.17}
\]
The function \(f\) must depend jointly on energy and momentum;
separate marginal Grams cannot detect correlations placing mass
outside the cone.

\subsection{Acquisition consequence}

The coupled translation ledger must now retain:

\begin{enumerate}
\item one common Hilbert carrier for all generators;
\item self-adjointness and componentwise strong-resolvent
convergence;
\item exact cutoff strong commutation or the bounded residual
(78.7);
\item joint \(C_0\)-functional-calculus source rows;
\item an approximate cone support (78.12), or a cofinal
measure-determining family of positive Grams (78.17);
\item compatibility with the V72 metric-fibre cocycle and V76
metric-base form cover.
\end{enumerate}

\begin{assessment}[V78 advance]
\label{assess:v13v78}
The energy and momentum limits can now be assembled into one
joint observable.  Exact cutoff commutation, or vanishing
commutators of bounded resolvents, preserves strong commutation.
Componentwise strong-resolvent convergence then yields the full
joint \(C_0\)-calculus, and approximate support in any closed set
lands in the limit.  This supplies a rigorous continuum route
from the V77 directional screens to the forward-cone spectral
measure required by V64.
\end{assessment}

\begin{firewall}[Claim boundary after V78]
\label{fire:v13v78}
No coupled momentum generators, cutoff resolvent-commutator
rows, componentwise strong-resolvent limits, or joint cone
functional-calculus Grams have been populated.  The theorem does
not construct translation covariance of the interacting local
net.  No global forward-cone SM--Einstein representation or full
continuum is proved.
\end{firewall}

\subsection{Parent record}

This V78 note appends the preceding material to the validated V77
source, whose record was
\[
\begin{array}{c|c}
 \text{lines}&24521\\
 \text{bytes}&864364\\
 \text{SHA--256}&
 \texttt{EC7E59842EF4A6F4713CEF162B41D566D59FFE768629BAFBD8EE630DB7106F0D}.
\end{array}                                                     \tag{78.18}
\]

% =============================================================================
% END APPENDED THIRTEENTH-CYCLE V78 MATERIAL
% =============================================================================

% =============================================================================
% BEGIN APPENDED THIRTEENTH-CYCLE V79 MATERIAL
% =============================================================================

\section{Thirteenth-cycle V79: extending covariance from a dense
symmetry ledger}
\label{sec:v13v79}

\subsection{The continuum-group gap}

The local ledgers in earlier cycles are countable, while the
Poincar\'e and diffeomorphism groups are not.  Exact covariance
on a countable dense subgroup determines a continuum
representation only when a uniform continuity modulus prevents
high-frequency symmetry oscillation.  This note proves that
extension theorem on the Hilbert carrier produced in V72--V78.

Let \(G\) be a second-countable metrizable topological group,
\(D\subset G\) a countable dense subgroup, and \(\mathcal H\) a
Hilbert space.  For each \(n\), let
\[
 U_n:D\longrightarrow U(\mathcal H)                          \tag{79.1}
\]
be a unitary representation.  Fix a dense linear subspace
\(\mathcal D_0\subset\mathcal H\).

\subsection{The extension theorem}

\begin{theorem}[Dense-subgroup strong landing]
\label{thm:v13v79-extension}
Assume:

\begin{enumerate}
\item for every \(g\in D\) and \(\psi\in\mathcal D_0\), the
sequence \(U_n(g)\psi\) converges;
\item for every \(\psi\in\mathcal D_0\),
\[
 \lim_{\substack{g\to e\\g\in D}}
 \sup_n\|(U_n(g)-I)\psi\|=0.                                 \tag{79.2}
\]
\end{enumerate}
Then there is a unique strongly continuous unitary
representation
\[
 U:G\longrightarrow U(\mathcal H)                            \tag{79.3}
\]
such that
\[
 U_n(g)\longrightarrow U(g)\quad\text{strongly}
 \qquad(g\in D).                                             \tag{79.4}
\]

If every \(U_n\) is already a strongly continuous representation
of \(G\) and (79.2) holds with \(g\to e\) in \(G\), then (79.4)
holds for every \(g\in G\).

\begin{proof}
For fixed \(g\in D\), convergence on \(\mathcal D_0\) extends to
strong convergence on all of \(\mathcal H\), because the
\(U_n(g)\) are isometries.  Denote the strong limit by
\(U_D(g)\).  Multiplication of uniformly bounded strongly
convergent sequences gives
\[
 U_D(g)U_D(h)=U_D(gh).                                       \tag{79.5}
\]
Applying the same argument to \(g^{-1}\) shows that
\(U_D(g^{-1})\) is the inverse of \(U_D(g)\); hence every
\(U_D(g)\) is unitary.

Passing \(n\to\infty\) in (79.2) gives strong continuity of
\(U_D\) at the identity on \(\mathcal D_0\).  Isometry and density
extend it to every vector.  If \(g_k\in D\) converges to
\(g\in G\), then for every \(\psi\),
\[
 \|U_D(g_k)\psi-U_D(g_\ell)\psi\|
 =
 \|(U_D(g_\ell^{-1}g_k)-I)\psi\|\longrightarrow0.            \tag{79.6}
\]
Thus \(U_D(g_k)\psi\) has a limit independent of the approximating
sequence.  Define that limit to be \(U(g)\psi\).  Equation
(79.5), continuity, and approximation from \(D\) prove the group
law, unitarity, and strong continuity on \(G\).  Density of \(D\)
gives uniqueness.

For the final assertion, fix \(g\in G\) and choose \(h\in D\)
near \(g\).  For \(\psi\in\mathcal D_0\),
\[
\begin{split}
 \|(U_n(g)-U(g))\psi\|
 &\leq
 \|(U_n(h^{-1}g)-I)\psi\|\\
 &\quad+\|(U_n(h)-U(h))\psi\|
 +\|(U(h^{-1}g)-I)\psi\|.                                   \tag{79.7}
\end{split}
\]
The first and third terms are uniformly small when \(h\) is
close to \(g\), and the middle term tends to zero by (79.4).
Density again extends the conclusion to all vectors.
\end{proof}
\end{theorem}

The modulus (79.2) must be measured after all cutoff Hilbert
spaces have been transported to one carrier.  Pointwise
continuity with an \(n\)-dependent modulus is insufficient.

\subsection{A sharp oscillation counterexample}

\begin{proposition}[Dense exact convergence without
equicontinuity]
\label{prop:v13v79-no-equicontinuity}
There are continuous one-dimensional unitary representations of
\(S^1\) that converge to the identity representation at every
finite-order point but have no strongly convergent continuum
limit.

\begin{proof}
Write \(S^1=\mathbb R/(2\pi\mathbb Z)\), let \(D\) be the
subgroup of rational angles, and define
\[
 U_n(\theta):=e^{in!\theta}.                                 \tag{79.8}
\]
If \(\theta=2\pi p/q\), then \(U_n(\theta)=1\) for every
\(n\geq q\).  Thus the sequence converges exactly to \(1\) at
every point of the dense subgroup \(D\).

However, with \(\theta_n=\pi/n!\),
\[
 \theta_n\to0,\qquad |U_n(\theta_n)-1|=2.                    \tag{79.9}
\]
Hence (79.2) fails.  If the sequence had a pointwise limit that
were a strongly continuous representation, its value \(1\) on
the dense subgroup \(D\) would force it to be the identity on
all of \(S^1\).  Dominated convergence with Haar measure would
then give
\[
 \int_{S^1}U_n(\theta)\,d\theta\longrightarrow1,
\]
whereas the integral of every nontrivial character in (79.8) is
zero.  This contradiction proves that no such continuum limit
exists.
\end{proof}
\end{proposition}

A growing list of exact lattice rotations can therefore coexist
with an unbounded symmetry frequency invisible at every fixed
ledger entry.

\subsection{Passing covariance of observables}

Let \(\mathfrak A_0\) be a norm-dense unital
\(*\)-subalgebra of a represented \(C^*\)-algebra
\(\mathfrak A\subset B(\mathcal H)\).  Suppose \(D\) acts on
\(\mathfrak A_0\) by \(*\)-automorphisms \(\alpha_g\).

\begin{proposition}[Dense-ledger covariance landing]
\label{prop:v13v79-covariance}
Under Theorem~\ref{thm:v13v79-extension}, suppose that for every
\(A\in\mathfrak A_0\) and \(g\in D\) there are uniformly bounded
operators \(A_n\to A\) and \(A_{n,g}\to\alpha_g(A)\) strongly,
with adjoints also converging strongly, and
\[
 \|U_n(g)A_nU_n(g)^*-A_{n,g}\|\longrightarrow0.              \tag{79.10}
\]
Then
\[
 U(g)AU(g)^*=\alpha_g(A)
 \qquad(A\in\mathfrak A_0,\ g\in D).                         \tag{79.11}
\]
If \(g\mapsto\alpha_g(A)\) is norm continuous, (79.11) extends
to \(g\in G\), and norm density extends it to all
\(A\in\mathfrak A\).

\begin{proof}
Uniform boundedness and strong convergence allow every product in
(79.10) to pass strongly to the limit, giving (79.11).  Approximate
an arbitrary \(g\in G\) by elements of \(D\); strong continuity
of \(U\) and norm continuity of \(\alpha_g(A)\) pass the identity.
For arbitrary \(A\in\mathfrak A\), use
\[
 \|U(g)(A-A_k)U(g)^*\|=\|A-A_k\|                             \tag{79.12}
\]
with \(A_k\in\mathfrak A_0\).
\end{proof}
\end{proposition}

For local nets, the action must also transport the localization
label:
\[
 U(g)\mathfrak A(\mathcal O)U(g)^*
 =\mathfrak A(g\mathcal O).                                 \tag{79.13}
\]
Continuity of the operator action does not by itself define or
continuously transport the regions; a common causal/localization
ledger remains necessary.

\subsection{Vacuum and spectral compatibility}

If normalized vectors \(\Omega_n\to\Omega\) and
\[
 \sup_{g\in K}\|U_n(g)\Omega_n-\Omega_n\|\longrightarrow0     \tag{79.14}
\]
on every compact \(K\subset G\), then
\(U(g)\Omega=\Omega\).  For the translation subgroup, V78 passes
the joint spectral support.  Lorentz covariance then preserves
the cone because the proper orthochronous Lorentz group preserves
\(\overline V_+\).  These are separate rows: group extension does
not supply spectral support, and the cone theorem does not supply
boost covariance.

\subsection{Acquisition consequence}

A continuum covariance compiler must retain:

\begin{enumerate}
\item a countable dense symmetry subgroup and its exact
multiplication table on the ledger;
\item strong convergence on a dense Hilbert source family;
\item the cutoff-uniform identity modulus (79.2);
\item raw observable covariance residuals (79.10), including
adjoints;
\item continuity of the limiting observable and region actions;
\item the invariant-vacuum residual (79.14);
\item the independent joint spectral rows of V78.
\end{enumerate}

\begin{assessment}[V79 advance]
\label{assess:v13v79}
Covariance on a countable ledger now has a rigorous route to the
full continuum group.  Strong pointwise convergence on the dense
subgroup plus a cutoff-uniform identity modulus produces a unique
strongly continuous unitary representation and transports
observable covariance.  The factorial-character counterexample
shows that exact convergence at every fixed dense-subgroup element
does not replace equicontinuity.
\end{assessment}

\begin{firewall}[Claim boundary after V79]
\label{fire:v13v79}
No Poincar\'e- or diffeomorphism-dense subgroup representation,
uniform identity modulus, raw local covariance residual, region
action, or invariant-vacuum limit has been populated for the
coupled regulator.  The theorem does not construct boost
generators or a causal net.  No covariant full SM--Einstein
continuum is proved.
\end{firewall}

\subsection{Parent record}

This V79 note appends the preceding material to the validated V78
source, whose record was
\[
\begin{array}{c|c}
 \text{lines}&24810\\
 \text{bytes}&874353\\
 \text{SHA--256}&
 \texttt{DB8D014A76D06D070B969594BCF8D253A42574C534C45F09D5FAB05289260944}.
\end{array}                                                     \tag{79.15}
\]

% =============================================================================
% END APPENDED THIRTEENTH-CYCLE V79 MATERIAL
% =============================================================================

% =============================================================================
% BEGIN APPENDED THIRTEENTH-CYCLE V80 MATERIAL
% =============================================================================

\section{Thirteenth-cycle V80: compulsory companion audit before
a common causal screen}
\label{sec:v13v80}

\subsection{Audit record}

The newest companion files remain
\[
\begin{array}{c|c|r|r|c}
 \text{programme}&\text{version}&\text{lines}&\text{bytes}&
 \text{SHA--256}\\ \hline
 \text{Yang--Mills mass gap}&V95&30950&1056018&
 \texttt{2DB19DAF571B44B04B28A51FAE2708D87BB6234D9CE9A2FBC5530ADF081809C2}\\
 \text{Navier--Stokes stability}&V26&5319&278283&
 \texttt{82C887F8C2C69E4F28FAD7C3DC8DE5B9099CB5A4BF431EBA1FFF3E91935978AD}.
\end{array}                                                     \tag{80.1}
\]
The auxiliary YM parallel V5 digest remains
\[
 \texttt{306F1BB84CFA8A8D47EA569EC17E9545F9867C8D32B548EC9D9C444B4F3C0512}.
                                                                  \tag{80.2}
\]
No companion file has changed since V75.  Their terminal result
and firewall blocks were checked again; no new numerical or
structural result is available to import.

YM V95 still reinforces the distinction between a sharp discrete
screen and a global neighborhood theorem.  NS V26 still
reinforces the need to test mixed directions rather than a
restricted rank-one family.  V76--V79 already incorporated those
lessons into metric-base form covers, joint spectral support, and
dense-group covariance.

\begin{proposition}[V80 direction decision]
\label{prop:v13v80-decision}
The next unresolved coupled row is a common causal-separation
certificate over the admitted metric background core.  V81 shall
derive such a certificate from uniform ADM lapse, shift, and
spatial-metric bounds, and shall retain a strict separation
margin.  This is required before fibrewise locality can be
promoted through Proposition~\ref{prop:v13v72-locality}.

\begin{proof}
The locality theorem in each fibre refers to that fibre's causal
cones.  A pair of coordinate-labelled regions can be spacelike
for one metric and causally related for another.  Neither the
V79 group modulus nor the V78 spectrum condition controls this
geometric change.  A uniform cone-speed bound converts the
metric uncertainty into one background-independent spatial
travel estimate and therefore supplies the missing common
screen.
\end{proof}
\end{proposition}

\begin{assessment}[V80 audit outcome]
\label{assess:v13v80}
The mandatory audit is complete.  With no changed companion
input, the programme continues at the first live coupled
bottleneck: robust causal separation across the random metric
support.
\end{assessment}

\begin{firewall}[Claim boundary after V80]
\label{fire:v13v80}
No common causal screen or Lorentzian metric-core estimate is
proved in this audit.  No local coupled continuum or full
SM--Einstein continuum is established.
\end{firewall}

\subsection{Parent record}

This V80 note appends the preceding material to the validated V79
source, whose record was
\[
\begin{array}{c|c}
 \text{lines}&25073\\
 \text{bytes}&883724\\
 \text{SHA--256}&
 \texttt{F52324AED8B8EC595749616E1F1E07DB51C50104FB2F97DEB648EDED425E6DEB}.
\end{array}                                                     \tag{80.3}
\]

% =============================================================================
% END APPENDED THIRTEENTH-CYCLE V80 MATERIAL
% =============================================================================

% =============================================================================
% BEGIN APPENDED THIRTEENTH-CYCLE V81 MATERIAL
% =============================================================================

\section{Thirteenth-cycle V81: a common causal screen over the
metric core}
\label{sec:v13v81}

\subsection{Uniform ADM cone control}

Let \(M=I\times\Sigma\), let \(h_0\) be a complete Riemannian
metric on \(\Sigma\), and consider a family \(\mathcal G\) of
time-oriented Lorentzian metrics with ADM form
\[
 g=-N_g^2dt^2+
 h_{g,t}\bigl(dx+\beta_gdt,dx+\beta_gdt\bigr).               \tag{81.1}
\]
Assume the uniform bounds
\[
 0<N_g\leq N_+,\qquad
 h_{g,t}(w,w)\geq\lambda |w|_{h_0}^2,\qquad
 |\beta_g|_{h_0}\leq B                                      \tag{81.2}
\]
for every admitted metric, time, point, and tangent vector, where
\(\lambda>0\).  Put
\[
 v_*:=B+{N_+\over\sqrt\lambda}.                              \tag{81.3}
\]

\begin{theorem}[Uniform causal-speed bound]
\label{thm:v13v81-speed}
Every future- or past-directed causal curve
\(\gamma=(t,x)\) for any \(g\in\mathcal G\) satisfies
\[
 d_{h_0}(x(t_1),x(t_2))
 \leq v_*|t_2-t_1|.                                         \tag{81.4}
\]

\begin{proof}
A nonconstant causal vector has nonzero \(dt\), since a vector
with \(dt=0\) has squared norm
\(h_{g,t}(dx,dx)>0\).  Thus \(t\) is strictly monotone along a
causal curve and may be used as parameter.  The causal inequality
from (81.1) is
\[
 h_{g,t}(\dot x+\beta_g,\dot x+\beta_g)\leq N_g^2.           \tag{81.5}
\]
Using (81.2),
\[
 |\dot x+\beta_g|_{h_0}
 \leq {N_+\over\sqrt\lambda},
 \qquad
 |\dot x|_{h_0}\leq v_*.                                    \tag{81.6}
\]
Integrating the \(h_0\)-length and using that distance is bounded
by length gives (81.4).  Reversing orientation proves the
past-directed case.
\end{proof}
\end{theorem}

\subsection{A background-independent spacelike pair}

Let \(J_1,J_2\subset I\) be compact time intervals and
\(K_1,K_2\subset\Sigma\) compact spatial sets.  Define
\[
 T_{12}:=\sup\{|s-t|:s\in J_1,\ t\in J_2\},\qquad
 D_{12}:=d_{h_0}(K_1,K_2).                                  \tag{81.7}
\]

\begin{corollary}[Common causal-separation certificate]
\label{cor:v13v81-common}
If
\[
 \boxed{\quad
 \mathfrak m_{12}:=D_{12}-v_*T_{12}>0,
 \quad}                                                       \tag{81.8}
\]
then
\[
 \mathcal O_1:=J_1\times K_1,\qquad
 \mathcal O_2:=J_2\times K_2                                 \tag{81.9}
\]
are causally disjoint for every \(g\in\mathcal G\).

\begin{proof}
If a causal curve joined points of the two sets, (81.4) would
give
\[
 D_{12}\leq d_{h_0}(x_1,x_2)
 \leq v_*|t_1-t_2|\leq v_*T_{12},                            \tag{81.10}
\]
contradicting (81.8).  The same argument covers either time
orientation.
\end{proof}
\end{corollary}

The positive number \(\mathfrak m_{12}\) is the causal reserve.
It is a geometric analogue of the angular reserve in V77 and
must remain explicit under cutoff and background transport.

\subsection{Why pointwise spacelikeness is insufficient}

\begin{proposition}[A zero causal margin can flip under an
arbitrarily small metric change]
\label{prop:v13v81-margin-no-go}
In \(1+1\) dimensions, arbitrarily close Lorentzian metrics can
classify the same pair of points on opposite sides of the null
cone.

\begin{proof}
Let
\[
 g_c=-c^2dt^2+dx^2.                                         \tag{81.11}
\]
For the displacement \((\Delta t,\Delta x)=(1,1)\), the squared
interval is \(1-c^2\).  It is null at \(c=1\), spacelike for
\(c<1\), and timelike for \(c>1\).  Metrics with \(c\) on either
side can be arbitrarily close in every coefficient norm.
Therefore a non-strict or single-background classification gives
no common causal screen.
\end{proof}
\end{proposition}

\subsection{Landing locality in the metric direct integral}

\begin{theorem}[Common-screen decomposable locality]
\label{thm:v13v81-locality}
Let the metric law \(\mu\) be supported on \(\mathcal G\), and
suppose the measurable fibre nets satisfy locality with respect
to their own metrics.  For every pair obeying (81.8), the
decomposable algebras satisfy
\[
 [\mathfrak A(\mathcal O_1),
   \mathfrak A(\mathcal O_2)]=0.                             \tag{81.12}
\]
If the cocycle covariance hypotheses of V72 and the continuum
group hypotheses of V79 hold, this locality identity is
covariant on the common-screen region ledger.

\begin{proof}
Corollary~\ref{cor:v13v81-common} makes the pair causally
disjoint in almost every metric fibre.  Fibre locality therefore
annuls the commutator almost everywhere.  Proposition
\ref{prop:v13v72-locality} makes the decomposable commutator zero.
The covariance statement follows from (72.7), (72.12), and the
V79 extension theorem.
\end{proof}
\end{theorem}

This theorem uses only a sufficient common screen.  Failure of
(81.8) does not prove causal contact; it returns the pair as
unclassified and calls for a sharper metric-dependent or
relational analysis.

\subsection{Finite acquisition rows}

For each background cell and region pair, retain
\[
 \lambda_n,\quad N_{+,n},\quad B_n,\quad
 v_{*,n}=B_n+N_{+,n}/\sqrt{\lambda_n},                       \tag{81.13}
\]
\[
 D_{12,n},\quad T_{12,n},\quad
 \mathfrak m_{12,n}=D_{12,n}-v_{*,n}T_{12,n}.                \tag{81.14}
\]
A positive continuum classification requires outward bounds with
\[
 \liminf_n\mathfrak m_{12,n}>0.                              \tag{81.15}
\]
The lapse lower bound, inverse-metric control, time orientation,
and measurability of the fibre local generators remain separate
rows even though only the upper cone speed appears in (81.3).

\begin{assessment}[V81 advance]
\label{assess:v13v81}
Uniform lapse, shift, and spatial-metric bounds now produce an
explicit causal speed valid over the whole admitted metric core.
A strict spatial-versus-temporal margin certifies one
background-independent spacelike region pair and promotes
fibrewise locality to the decomposable coupled algebra.  A null
pair is shown to change type under arbitrarily small metric
perturbations, proving that the strict reserve is necessary.
\end{assessment}

\begin{firewall}[Claim boundary after V81]
\label{fire:v13v81}
No Lorentzian metric law, uniform ADM bounds, common causal
region ledger, fibre local net, or cocycle covariance has been
populated.  The certificate is sufficient and does not classify
all spacelike pairs or construct relational observables.  No
coupled local SM--Einstein continuum is proved.
\end{firewall}

\subsection{Parent record}

This V81 note appends the preceding material to the validated V80
source, whose record was
\[
\begin{array}{c|c}
 \text{lines}&25163\\
 \text{bytes}&887241\\
 \text{SHA--256}&
 \texttt{79E04CFF9DEDD4A2542BA84BECAF514FAAD7D766F8D585678CB54DCA6FAB2629}.
\end{array}                                                     \tag{81.16}
\]

% =============================================================================
% END APPENDED THIRTEENTH-CYCLE V81 MATERIAL
% =============================================================================

% =============================================================================
% BEGIN APPENDED THIRTEENTH-CYCLE V82 MATERIAL
% =============================================================================

\section{Thirteenth-cycle V82: a countable cofinal common-causal
region ledger}
\label{sec:v13v82}

\subsection{Countable rational rectangles}

Retain the product spacetime and uniform speed \(v_*\) of V81.
Because \((\Sigma,h_0)\) is a complete finite-dimensional
Riemannian manifold, it is proper and separable.  Choose a
countable dense set of spatial centers.  Let \(\mathfrak R_0\)
consist of finite unions of rectangles
\[
 J\times B_{h_0}(z,r),                                      \tag{82.1}
\]
where the endpoints of \(J\), the radius \(r\), and the label of
\(z\) are drawn from countable dense sets.  Then
\(\mathfrak R_0\) is countable and is a basis, closed under finite
unions, for the product topology.

\begin{lemma}[Strict causal margins survive rational
enlargement]
\label{lem:v13v82-enlargement}
Let
\(\mathcal O_j=J_j\times K_j\) have compact factors and V81
margin
\[
 \mathfrak m_{12}
 =d_{h_0}(K_1,K_2)-v_*T_{12}>0.                             \tag{82.2}
\]
Suppose
\[
 K_j\subset K'_j\subset
 \{x:d_{h_0}(x,K_j)<\eta_x\},\qquad
 J_j\subset J'_j\subset
 \{t:d(t,J_j)<\eta_t\}.                                     \tag{82.3}
\]
Then
\[
 \mathfrak m'_{12}
 \geq\mathfrak m_{12}-2\eta_x-2v_*\eta_t.                   \tag{82.4}
\]
In particular, sufficiently small enlargements remain causally
disjoint for every metric in the admitted core.

\begin{proof}
For \(x'_j\in K'_j\), choose \(x_j\in K_j\) within
\(\eta_x\).  The triangle inequality gives
\[
 d_{h_0}(K'_1,K'_2)
 \geq d_{h_0}(K_1,K_2)-2\eta_x.                             \tag{82.5}
\]
The enlarged time intervals obey
\[
 T'_{12}\leq T_{12}+2\eta_t.                                \tag{82.6}
\]
Subtracting \(v_*T'_{12}\) from (82.5) proves (82.4).
Corollary~\ref{cor:v13v81-common} gives the final assertion.
\end{proof}
\end{lemma}

\begin{theorem}[Cofinality of the rational common screen]
\label{thm:v13v82-cofinal-regions}
Every compact product-region pair with positive margin (82.2)
admits enclosing regions
\[
 \mathcal O_j\subset R_j,\qquad R_j\in\mathfrak R_0,          \tag{82.7}
\]
whose closures still have a positive common causal margin.
Moreover, the \(R_j\) may be chosen inside any prescribed open
neighborhoods of \(\mathcal O_j\).

\begin{proof}
Compactness of each factor and the rational-ball basis give a
finite union \(R_j\in\mathfrak R_0\) contained in the prescribed
neighborhood and satisfying (82.3) for arbitrarily small
\(\eta_x,\eta_t\).  Choose them so that
\[
 2\eta_x+2v_*\eta_t<\mathfrak m_{12}.                       \tag{82.8}
\]
Lemma~\ref{lem:v13v82-enlargement} then gives a positive margin
for the enclosing pair.
\end{proof}
\end{theorem}

The result converts the uncountable family of strictly
common-spacelike compact pairs into a countable determining
ledger.  Pairs whose common margin is zero remain outside this
argument.

\subsection{Extending the ledger net}

Suppose a unital isotone net
\[
 R\longmapsto\mathfrak A_0(R),\qquad R\in\mathfrak R_0,       \tag{82.9}
\]
has been represented on one Hilbert space and satisfies
\[
 [\mathfrak A_0(R_1),\mathfrak A_0(R_2)]=0                  \tag{82.10}
\]
whenever the closures of \(R_1,R_2\) have positive common
margin.  For an arbitrary open product region \(\mathcal O\),
define
\[
 \mathfrak A(\mathcal O)
 :=
 C^*\!\left(
 \bigcup_{\substack{R\in\mathfrak R_0\\\overline R\subset
 \mathcal O}}\mathfrak A_0(R)\right).                        \tag{82.11}
\]

\begin{theorem}[Locality of the generated continuum net]
\label{thm:v13v82-net}
The assignment (82.11) is isotone.  If two open regions have
compact closures and a positive common causal margin, then
\[
 [\mathfrak A(\mathcal O_1),
   \mathfrak A(\mathcal O_2)]=0.                             \tag{82.12}
\]
If the V79 symmetry representation maps the ledger covariantly
and its region action preserves \(\mathfrak R_0\) or is
approximable from within it, the extended net is covariant.

\begin{proof}
Isotony follows immediately from inclusion of the indexing
families in (82.11).  Every pair of ledger regions with closures
inside \(\mathcal O_1,\mathcal O_2\) is still commonly causally
disjoint.  Hence their algebras commute by (82.10).  Finite
polynomials in the two unions commute, and norm continuity of
the commutator extends this to both \(C^*\)-closures, proving
(82.12).

For covariance, first apply the ledger covariance identity to
generators.  Inner approximation and isotony identify the
generated algebra of the transformed region; norm closure then
gives equality.
\end{proof}
\end{theorem}

\subsection{Finite residual version}

At cutoff \(n\), let \(A_{n,a}\) and \(B_{n,b}\) be norm-bounded
generators for two ledger regions.  The positive commutator Gram
\[
 \mathbb G^{\rm loc}_{n,R}
 :=
 \sum_{a,b}
 [A_{n,a},B_{n,b}]^*[A_{n,a},B_{n,b}]                       \tag{82.13}
\]
vanishes exactly when all generator commutators vanish.  If the
generators converge strongly together with their adjoints and
\(\mathbb G^{\rm loc}_{n,R}\to0\) strongly at every fixed ledger
screen, then the limiting generator algebras commute.  Norm-dense
generation and Theorem~\ref{thm:v13v82-net} finish the extension.

The geometric and operator rows must remain paired:
\[
 \liminf_n\mathfrak m_{12,n}>0,\qquad
 \mathbb G^{\rm loc}_{n,R}\longrightarrow0.                 \tag{82.14}
\]
A zero commutator for a pair whose causal margin was not
transported is not a certified locality statement.

\begin{assessment}[V82 advance]
\label{assess:v13v82}
Every compact region pair with strict common causal margin now
admits a countable rational enlargement retaining that margin.
Locality on this countable ledger extends by norm closure to the
generated continuum \(C^*\)-net, and covariance extends from the
V79 action under inner region approximation.  This supplies the
countable localization framework needed by the all-order
correlation ledger of V68.
\end{assessment}

\begin{firewall}[Claim boundary after V82]
\label{fire:v13v82}
No rational region algebra, finite commutator Gram, inner
regularity, causal-margin table, or covariant region action has
been populated for the coupled regulator.  Zero-margin and
metric-dependent relational regions are not classified.  No
full local SM--Einstein continuum is proved.
\end{firewall}

\subsection{Parent record}

This V82 note appends the preceding material to the validated V81
source, whose record was
\[
\begin{array}{c|c}
 \text{lines}&25373\\
 \text{bytes}&894340\\
 \text{SHA--256}&
 \texttt{50DD9DBEDA1FA13DE916453D471EE196B8CC99EAECCD6299820928F10D7B7685}.
\end{array}                                                     \tag{82.15}
\]

% =============================================================================
% END APPENDED THIRTEENTH-CYCLE V82 MATERIAL
% =============================================================================

% =============================================================================
% BEGIN APPENDED THIRTEENTH-CYCLE V83 MATERIAL
% =============================================================================

\section{Thirteenth-cycle V83: a norm-dense time-slice
reconstruction criterion}
\label{sec:v13v83}

\subsection{The algebraic gate}

Locality and covariance do not imply that observables in a
neighborhood of a Cauchy surface generate the algebra of its
causal development.  Let
\[
 \mathfrak A(\mathcal O)\subset
 \mathfrak A(\mathcal D(\mathcal O))                         \tag{83.1}
\]
be the isotonic inclusion for a region \(\mathcal O\) containing
a Cauchy neighborhood in its causal development.

\begin{theorem}[Dense-generator time-slice criterion]
\label{thm:v13v83-time-slice}
Let \(\{B_j\}_{j\geq1}\) generate
\(\mathfrak A(\mathcal D(\mathcal O))\) as a \(C^*\)-algebra.
Suppose that for every \(j\) there are
\(A_{j,k}\in\mathfrak A(\mathcal O)\) such that
\[
 \|A_{j,k}-B_j\|\longrightarrow0.                            \tag{83.2}
\]
Then
\[
 \boxed{\quad
 \mathfrak A(\mathcal O)
 =\mathfrak A(\mathcal D(\mathcal O)).
 \quad}                                                       \tag{83.3}
\]

\begin{proof}
A \(C^*\)-subalgebra is norm closed, so (83.2) places every
\(B_j\) in \(\mathfrak A(\mathcal O)\).  It therefore contains
the \(C^*\)-algebra generated by all \(B_j\), which is the
right-hand side of (83.3).  The reverse inclusion is isotony.
\end{proof}
\end{theorem}

This formulation allows the reconstruction operator to depend on
the generator and on the refinement index.  It requires norm
approximation; weak expectation agreement in one state does not
place an observable in the slice algebra.

\begin{proposition}[One-state slice data do not imply the
time-slice axiom]
\label{prop:v13v83-state-no-go}
There is a proper unital inclusion
\(\mathfrak A\subsetneq\mathfrak B\) and a faithful state on
\(\mathfrak B\) whose restriction supplies all expectation
values on \(\mathfrak A\), but \(\mathfrak A\neq\mathfrak B\).

\begin{proof}
Take
\[
 \mathfrak A=\mathbb CI\subset M_2(\mathbb C)=\mathfrak B
\]
and the normalized trace on \(\mathfrak B\).  Its restriction
completely determines the scalar-algebra expectations, yet the
three traceless Pauli directions are absent from
\(\mathfrak A\).  Thus state compatibility across the inclusion
does not prove algebraic generation.
\end{proof}
\end{proposition}

\subsection{Finite reconstruction residuals}

At cutoff \(n\), let \(B_{n,1},\ldots,B_{n,m}\) be a selected
generator screen in the causal-development algebra.  Let
\(\mathcal S_n(\mathcal O)\) be the finite synthesis range of
words declared to occur in the Cauchy neighborhood.  Define
\[
 \delta^{\rm ts}_{n,j}
 :=
 \inf_{A\in\mathcal S_n(\mathcal O)}
 \|B_{n,j}-A\|.                                              \tag{83.4}
\]

\begin{theorem}[Fixed-screen time-slice landing]
\label{thm:v13v83-landing}
Suppose all cutoff algebras are represented on one carrier,
\[
 B_{n,j}\longrightarrow B_j\quad\text{in norm}              \tag{83.5}
\]
for every fixed \(j\), and for each \(n,j\) choose
\(A_{n,j}\in\mathcal S_n(\mathcal O)\) with
\[
 \|B_{n,j}-A_{n,j}\|
 \leq\delta^{\rm ts}_{n,j}+n^{-1}.                           \tag{83.6}
\]
If
\[
 \delta^{\rm ts}_{n,j}\longrightarrow0                      \tag{83.7}
\]
and every norm limit of such slice words belongs to the limiting
slice algebra, then \(B_j\in\mathfrak A(\mathcal O)\).  If the
fixed screens grow cofinally through a countable dense generator
ledger, the limiting net satisfies (83.3).

\begin{proof}
Equations (83.5)--(83.7) give
\[
 \|A_{n,j}-B_j\|
 \leq\|A_{n,j}-B_{n,j}\|+\|B_{n,j}-B_j\|
 \longrightarrow0.                                         \tag{83.8}
\]
The closure hypothesis places \(B_j\) in the limiting slice
algebra.  Cofinality supplies every ledger generator, and
Theorem~\ref{thm:v13v83-time-slice} finishes the proof.
\end{proof}
\end{theorem}

The closure hypothesis in this theorem is the occurrence row.
Without it, an approximating matrix word may converge as an
operator while its support label escapes the Cauchy
neighborhood.

\subsection{Hyperbolic evolution as a sufficient mechanism}

Let \(\mathcal E_n\) be a finite Cauchy-evolution map from slice
data to the selected causal-development generator screen.
Suppose it produces \(A_{n,j}\) and that its equation-of-motion,
boundary, and support errors give
\[
 \|B_{n,j}-A_{n,j}\|
 \leq
 C_{n,j}\bigl(
 \varepsilon^{\rm EOM}_{n,j}
 +\varepsilon^{\rm bdry}_{n,j}
 +\varepsilon^{\rm supp}_{n,j}\bigr).                        \tag{83.9}
\]
If \(C_{n,j}\) is bounded at every fixed screen and the three
residuals vanish, then (83.7) follows.  The stability constant
must be retained: vanishing equation residuals multiplied by an
unbounded inverse evolution norm do not prove reconstruction.

For a random metric law, (83.9) must hold on the same
common-causal background cells as V81--V82, or hold almost surely
with a uniform-integrability estimate strong enough to pass its
operator norm.  A reconstruction valid at one reference metric
does not imply a decomposable time-slice identity.

\subsection{Compatibility with the all-order ledger}

Once (83.3) holds, every bounded correlation in V68 whose
insertions lie in \(\mathcal D(\mathcal O)\) is determined by
slice-algebra words.  The converse is false: equality of finitely
many correlations follows from neither norm generation nor all
operator relations.  Therefore the finite compiler retains
separately:

\begin{enumerate}
\item the common causal-development label;
\item the slice occurrence synthesis;
\item the norm reconstruction residual (83.4);
\item EOM, boundary, and support contributions with the stability
constant in (83.9);
\item adjoint and product compatibility of the reconstructed
words;
\item cofinality in both generator index and region index.
\end{enumerate}

\begin{assessment}[V83 advance]
\label{assess:v13v83}
The time-slice axiom now has a finite-to-continuum criterion.
Norm approximation of a countable dense causal-development
generator ledger by genuine slice words proves equality of the
two \(C^*\)-algebras.  A faithful-state counterexample shows that
expectation compatibility is insufficient, and the finite
residual separates equation, boundary, support, and inverse
stability costs.
\end{assessment}

\begin{firewall}[Claim boundary after V83]
\label{fire:v13v83}
No hyperbolic SM--Einstein evolution map, slice occurrence
synthesis, reconstruction stability bound, common causal
development, or cofinal generator residual has been populated.
No time-slice axiom or full SM--Einstein continuum is proved.
\end{firewall}

\subsection{Parent record}

This V83 note appends the preceding material to the validated V82
source, whose record was
\[
\begin{array}{c|c}
 \text{lines}&25576\\
 \text{bytes}&901399\\
 \text{SHA--256}&
 \texttt{D335A555618F3427E959B10ADE3D9949A4F92B61E5230E0534C62711C7790877}.
\end{array}                                                     \tag{83.10}
\]

% =============================================================================
% END APPENDED THIRTEENTH-CYCLE V83 MATERIAL
% =============================================================================

% =============================================================================
% BEGIN APPENDED THIRTEENTH-CYCLE V84 MATERIAL
% =============================================================================

\section{Thirteenth-cycle V84: local affiliation of unbounded
observables}
\label{sec:v13v84}

\subsection{Using the existing resolvent landing theorem}

The Grand Tensor charge compiler already proves the
self-adjoint resolvent landing statement: a norm- or
strong-resolvent limit with an injective dense-range resolvent
defines a self-adjoint operator.  V58 likewise treats bounded
transforms for the Dirac carrier.  Those results are not repeated.
The new issue is localization: the landed resolvent must belong
to the von Neumann algebra of the region, and resolvents in
spacelike commuting algebras must imply strong commutation of the
unbounded observables.

Let \(\mathcal M\subset B(\mathcal H)\) be a von Neumann algebra.
A self-adjoint operator \(\Phi\) is affiliated with \(\mathcal M\)
when its spectral projections belong to \(\mathcal M\).
Equivalently,
\[
 (\Phi-i)^{-1}\in\mathcal M.                                 \tag{84.1}
\]

\begin{theorem}[Local resolvent landing]
\label{thm:v13v84-affiliation}
Let \(\Phi_n\) be self-adjoint operators affiliated with one
von Neumann algebra \(\mathcal M\), and put
\[
 R_n=(\Phi_n-i)^{-1}.                                        \tag{84.2}
\]
Suppose
\[
 R_n\to R,\qquad R_n^*\to R^*
 \quad\text{strongly},                                      \tag{84.3}
\]
and suppose \(R\) is injective with dense range.  Then
\[
 \Phi:=R^{-1}+i                                              \tag{84.4}
\]
is self-adjoint, \(R=(\Phi-i)^{-1}\), and \(\Phi\) is affiliated
with \(\mathcal M\).

\begin{proof}
Every \(R_n\) and \(R_n^*\) lies in \(\mathcal M\) by
affiliation.  A von Neumann algebra is strongly closed on bounded
sets, and resolvents have norm at most one, so (84.3) gives
\(R,R^*\in\mathcal M\).  The resolvent identity
\[
 R_n-R_n^*=2iR_n^*R_n                                       \tag{84.5}
\]
passes strongly to the limit.  Together with injectivity and
dense range, the existing resolvent characterization gives the
self-adjoint operator (84.4).  Finally (84.1) gives affiliation
with \(\mathcal M\).
\end{proof}
\end{theorem}

The injective dense-range row cannot be dropped: a strong limit
of resolvents may lose a subspace and cease to be the resolvent
of any densely defined self-adjoint operator.

\subsection{Spacelike strong commutation}

\begin{theorem}[Affiliated observables in commuting local
algebras strongly commute]
\label{thm:v13v84-strong-locality}
Let \(\mathcal M_1,\mathcal M_2\subset B(\mathcal H)\) satisfy
\[
 [\mathcal M_1,\mathcal M_2]=0.                              \tag{84.6}
\]
If self-adjoint \(\Phi_j\) is affiliated with
\(\mathcal M_j\), then \(\Phi_1,\Phi_2\) strongly commute.
Equivalently, all their spectral projections commute, and
\[
 f(\Phi_1)g(\Phi_2)=g(\Phi_2)f(\Phi_1)                       \tag{84.7}
\]
for all bounded Borel functions \(f,g\).

\begin{proof}
Affiliation places every spectral projection
\(E_{\Phi_1}(A)\) in \(\mathcal M_1\) and every
\(E_{\Phi_2}(B)\) in \(\mathcal M_2\).  Equation (84.6) makes
these projections commute for all Borel sets \(A,B\), which is
the definition of strong commutation.  Bounded Borel functional
calculus then proves (84.7).
\end{proof}
\end{theorem}

This is stronger and more domain-safe than writing
\([\Phi_1,\Phi_2]\psi=0\) on a selected core.  Strong commutation
controls the joint spectral calculus and all bounded functions
of the observables.  Fermionic fields require the corresponding
graded field algebra; the theorem applies directly to
self-adjoint bosonic fields and even observables such as smeared
stress components.

\subsection{Application to the common-causal ledger}

Let \(R\mapsto\mathcal M(R)\) be the von Neumann closures of the
V82 local \(C^*\)-net.  For a real test tensor \(f\) supported in
\(R\), let \(T_n(f)\) be a self-adjoint regulated stress
insertion.

\begin{corollary}[Local stress landing]
\label{cor:v13v84-stress}
Assume
\[
 (T_n(f)-i)^{-1}\in\mathcal M(R),                            \tag{84.8}
\]
and that the hypotheses of
Theorem~\ref{thm:v13v84-affiliation} hold.  Then the limiting
stress insertion \(T(f)\) is affiliated with \(\mathcal M(R)\).
If \(f_1,f_2\) have supports in a pair of regions with positive
common causal margin, then \(T(f_1)\) and \(T(f_2)\) strongly
commute.

\begin{proof}
The first assertion is
Theorem~\ref{thm:v13v84-affiliation}.  V82 makes the two regional
von Neumann algebras commute, and
Theorem~\ref{thm:v13v84-strong-locality} gives the second.
\end{proof}
\end{corollary}

The stress tensor used here must be the same renormalized
insertion that appears in the Grand Tensor Ward/Hamiltonian
compiler and in the V69 Einstein residual.  Landing a different
affiliated operator with the same vacuum expectation does not
close that incidence requirement.

\subsection{Covariance and the time slice}

Suppose the V79 representation acts covariantly on the bounded
net and finite resolvents obey
\[
 U_n(g)(\Phi_n(f)-i)^{-1}U_n(g)^*
 -
 (\Phi_n(g\cdot f)-i)^{-1}\longrightarrow0                  \tag{84.9}
\]
strongly on the common carrier.  Strong convergence of the
unitaries and resolvents gives
\[
 U(g)(\Phi(f)-i)^{-1}U(g)^*
 =(\Phi(g\cdot f)-i)^{-1}.                                  \tag{84.10}
\]
Uniqueness of a self-adjoint operator from one nonreal resolvent
then yields
\[
 U(g)\Phi(f)U(g)^*=\Phi(g\cdot f).                           \tag{84.11}
\]

If a causal-development field resolvent is approximated in norm
by the V83 slice algebra, it belongs to that algebra's von
Neumann closure and the affiliated field is localized in the
slice algebra.  Equality of a finite collection of field moments
does not imply this resolvent membership.

\subsection{Finite acquisition rows}

For each real local smearing and fixed region screen retain:

\begin{enumerate}
\item self-adjoint cutoff realizations and one nonreal resolvent;
\item regional resolvent occurrence (84.8);
\item strong-star resolvent convergence;
\item injectivity and dense range of the limiting resolvent;
\item covariance residual (84.9);
\item identity with the stress/current and Einstein writers;
\item cofinal support smearings and common-causal region labels.
\end{enumerate}

A resolvent Gram such as
\[
 \mathbb G^{\rm aff}_{n,R}
 :=
 \sum_j
 \bigl((I-\mathbb E_R)R_{n,j}\bigr)^*
 \bigl((I-\mathbb E_R)R_{n,j}\bigr)                         \tag{84.12}
\]
may test leakage from a finite regional operator system
\(\mathbb E_R\), but exact von Neumann membership requires a
cofinal ultraweak/strong closure argument.

\begin{assessment}[V84 advance]
\label{assess:v13v84}
The bounded local net now has a precise bridge to unbounded
observables.  Strong-star limits of locally occurring resolvents
produce affiliated self-adjoint fields, and affiliation with
commuting spacelike algebras gives strong commutation without
multiplying unbounded operators on an uncontrolled domain.
Covariance and time-slice localization reduce to bounded
resolvent identities.
\end{assessment}

\begin{firewall}[Claim boundary after V84]
\label{fire:v13v84}
No regulated stress resolvent, regional affiliation row,
strong-star limit, dense-range certificate, or covariant
smearing family has been populated.  Fermionic graded fields and
operator-valued distribution domains remain separate.  No local
unbounded SM field system or full SM--Einstein continuum is
proved.
\end{firewall}

\subsection{Parent record}

This V84 note appends the preceding material to the validated V83
source, whose record was
\[
\begin{array}{c|c}
 \text{lines}&25780\\
 \text{bytes}&908684\\
 \text{SHA--256}&
 \texttt{084044C7D67D9B39C8E8F58D3C039DFE349950DC020CE7E529CBEAF630547AF4}.
\end{array}                                                     \tag{84.13}
\]

% =============================================================================
% END APPENDED THIRTEENTH-CYCLE V84 MATERIAL
% =============================================================================

% =============================================================================
% BEGIN APPENDED THIRTEENTH-CYCLE V85 MATERIAL
% =============================================================================

\section{Thirteenth-cycle V85: compulsory companion audit before
the operator-valued distribution step}
\label{sec:v13v85}

\subsection{Audit record}

The newest companion notes remain unchanged:
\[
\begin{array}{c|c|r|r|c}
 \text{programme}&\text{version}&\text{lines}&\text{bytes}&
 \text{SHA--256}\\ \hline
 \text{Yang--Mills mass gap}&V95&30950&1056018&
 \texttt{2DB19DAF571B44B04B28A51FAE2708D87BB6234D9CE9A2FBC5530ADF081809C2}\\
 \text{Navier--Stokes stability}&V26&5319&278283&
 \texttt{82C887F8C2C69E4F28FAD7C3DC8DE5B9099CB5A4BF431EBA1FFF3E91935978AD}.
\end{array}                                                     \tag{85.1}
\]
The auxiliary YM parallel note remains V5 with SHA--256
\[
 \texttt{306F1BB84CFA8A8D47EA569EC17E9545F9867C8D32B548EC9D9C444B4F3C0512}.
                                                                  \tag{85.2}
\]
Their terminal theorem and firewall blocks were reread.  No new
companion theorem or numerical constant is available.

YM V95 continues to require a cofinal all-order interacting
ledger beyond its finite response calculation.  V82--V84 now
supply the corresponding countable region basis, time-slice
criterion, and local affiliation step on the SM side.  The next
gap is that separately affiliated smearings do not yet form a
linear operator-valued distribution on one domain.

\begin{proposition}[V85 direction decision]
\label{prop:v13v85-decision}
V86 shall construct a closable operator-valued distribution from
a countable dense smearing ledger under common-domain linearity
and test-function seminorm bounds.  It shall also show how
covariance, locality, and distributional conservation extend
from the ledger.

\begin{proof}
Affiliation in V84 localizes each self-adjoint smearing but does
not relate different smearings or provide continuity in the test
function.  Stress conservation and the Wightman field axioms
require precisely that missing linear-continuous structure.  No
companion estimate addresses it, so the proof must be made at the
common-domain level.
\end{proof}
\end{proposition}

\begin{assessment}[V85 audit outcome]
\label{assess:v13v85}
The mandatory audit is complete.  The programme proceeds from
isolated affiliated observables to a common-domain
operator-valued distribution, with no imported numerical data.
\end{assessment}

\begin{firewall}[Claim boundary after V85]
\label{fire:v13v85}
No common field domain, test-function continuity estimate, or
operator-valued SM distribution is proved in this audit.  The
full SM--Einstein continuum remains open.
\end{firewall}

\subsection{Parent record}

This V85 note appends the preceding material to the validated V84
source, whose record was
\[
\begin{array}{c|c}
 \text{lines}&26006\\
 \text{bytes}&916839\\
 \text{SHA--256}&
 \texttt{CE5F840977C95A1BADB22BF03250990B8F6A5C9A527B3A0D713B27FD2C492E6A}.
\end{array}                                                     \tag{85.3}
\]

% =============================================================================
% END APPENDED THIRTEENTH-CYCLE V85 MATERIAL
% =============================================================================

% =============================================================================
% BEGIN APPENDED THIRTEENTH-CYCLE V86 MATERIAL
% =============================================================================

\section{Thirteenth-cycle V86: landing a common-domain
operator-valued distribution}
\label{sec:v13v86}

\subsection{Test and domain spaces}

Let \(\mathcal E\) be a separable complex Fr\'echet test-function
space with defining seminorms \(p_1,p_2,\ldots\), and let
\(\mathcal E_0\subset\mathcal E\) be a countable dense complex
linear subspace closed under conjugation.  For local statements,
assume
\[
 \mathcal E_0(\mathcal O)
 :=\mathcal E_0\cap\mathcal E(\mathcal O)
\]
is dense in the subspace of tests supported in every ledger
region \(\mathcal O\).

Let \(\mathcal D\subset\mathcal H\) be a dense complex linear
domain.  Suppose a ledger map
\[
 \Phi_0:\mathcal E_0\longrightarrow
 \operatorname{Hom}(\mathcal D,\mathcal H)                  \tag{86.1}
\]
is complex linear and satisfies
\[
 \langle\psi,\Phi_0(f)\phi\rangle
 =
 \langle\Phi_0(\overline f)\psi,\phi\rangle
 \quad(\psi,\phi\in\mathcal D).                              \tag{86.2}
\]

\subsection{Distributional extension}

\begin{theorem}[Common-domain field extension]
\label{thm:v13v86-extension}
Assume that for every \(\phi\in\mathcal D\) there are
\(m=m(\phi)\) and \(C_\phi<\infty\) such that
\[
 \|\Phi_0(f)\phi\|
 \leq C_\phi p_m(f)
 \qquad(f\in\mathcal E_0).                                  \tag{86.3}
\]
Then there is a unique complex-linear extension
\[
 \Phi:\mathcal E\longrightarrow
 \operatorname{Hom}(\mathcal D,\mathcal H)                  \tag{86.4}
\]
such that \(f\mapsto\Phi(f)\phi\) is continuous for every
\(\phi\in\mathcal D\).  It satisfies
\[
 \langle\psi,\Phi(f)\phi\rangle
 =
 \langle\Phi(\overline f)\psi,\phi\rangle.                  \tag{86.5}
\]
Consequently \(\Phi(f)|_{\mathcal D}\) is closable for every
test \(f\), and for fixed \(\psi,\phi\in\mathcal D\),
\[
 f\longmapsto
 \langle\psi,\Phi(f)\phi\rangle                              \tag{86.6}
\]
is a continuous distribution.  If \(\mathcal E\) is a Schwartz
space, (86.6) is tempered.

\begin{proof}
For \(f\in\mathcal E\), choose \(f_k\in\mathcal E_0\) with
\(f_k\to f\).  Estimate (86.3) makes
\(\Phi_0(f_k)\phi\) Cauchy in \(\mathcal H\), and the same
estimate shows that its limit is independent of the
approximating sequence.  Define that limit to be
\(\Phi(f)\phi\).  Linearity, uniqueness, and continuity follow
from the construction.

Apply (86.2) to approximants \(f_k\) and pass to the limit to
obtain (86.5).  It says
\[
 \Phi(\overline f)|_{\mathcal D}
 \subset\Phi(f)^*.                                           \tag{86.7}
\]
The adjoint has dense domain because it contains \(\mathcal D\);
hence \(\Phi(f)|_{\mathcal D}\) is closable.  Finally (86.3) and
Cauchy--Schwarz make (86.6) a continuous linear functional on
\(\mathcal E\).
\end{proof}
\end{theorem}

Separate self-adjoint closures of the \(\Phi_0(f)\) would not
prove this theorem: they do not impose linearity or continuity in
\(f\).  Conversely, closability from (86.7) does not prove
essential self-adjointness or local affiliation; those require
the V84 resolvent rows.

\subsection{Products on a graph domain}

For Ward identities and locality, equip \(\mathcal D\) with a
complete locally convex graph topology having seminorms
\(q_1,q_2,\ldots\), continuously embedded in \(\mathcal H\).
Assume \(\Phi_0(f)\mathcal D\subset\mathcal D\) and, for every
\(r\), there are \(s,m,C\) with
\[
 q_r(\Phi_0(f)\phi)
 \leq C\,p_m(f)q_s(\phi).                                   \tag{86.8}
\]

\begin{lemma}[Graph-continuous extension]
\label{lem:v13v86-graph}
Under (86.8), the extension in
Theorem~\ref{thm:v13v86-extension} maps
\(\mathcal D\) to itself and is jointly continuous on bounded
sets in the stated seminorm sense.  In particular,
\[
 (f,g,\phi)\longmapsto\Phi(f)\Phi(g)\phi                    \tag{86.9}
\]
is continuous after choosing the finitely many seminorms
provided by two applications of (86.8).

\begin{proof}
The proof of Theorem~\ref{thm:v13v86-extension} may be repeated in
the complete graph topology using (86.8).  Applying the estimate
first to \(\Phi(g)\phi\) and then to \(\Phi(f)\) gives continuity
of (86.9).
\end{proof}
\end{lemma}

\subsection{Covariance and locality from the ledger}

\begin{theorem}[Extension of structural identities]
\label{thm:v13v86-structure}
Assume the graph hypotheses above.

\begin{enumerate}
\item Let \(U(g)\mathcal D=\mathcal D\), with \(U(g)\) continuous
in the graph topology, and let \(g\) act continuously on
\(\mathcal E\).  If
\[
 U(g)\Phi_0(f)U(g)^*=\Phi_0(g\cdot f)
 \quad(f\in\mathcal E_0)                                    \tag{86.10}
\]
on \(\mathcal D\), then the same identity holds for every
\(f\in\mathcal E\).
\item If
\[
 [\Phi_0(f),\Phi_0(g)]\phi=0                                \tag{86.11}
\]
for all \(\phi\in\mathcal D\) and all ledger tests with supports
in a positively common-spacelike V82 region pair, then
\[
 [\Phi(f),\Phi(g)]\phi=0                                    \tag{86.12}
\]
for all tests supported in that pair.
\end{enumerate}

\begin{proof}
For item 1, approximate \(f\) in \(\mathcal E_0\), apply
(86.10), and use graph continuity of \(U(g)\), \(\Phi\), and the
test action.  For item 2, use the support-dense ledger to choose
\(f_k,g_k\) with the same regional supports and limits \(f,g\).
Equation (86.11) holds for every \(k\), while
Lemma~\ref{lem:v13v86-graph} passes both ordered products to the
limit.
\end{proof}
\end{theorem}

For fermionic fields, (86.11) is replaced by the graded
commutator with the same proof.  Strong commutation of
self-adjoint bosonic closures is obtained only after the V84
affiliation theorem is also applied.

\subsection{Distributional conservation and the Einstein writer}

Let \(\mathcal E_T\) be the test space of compactly supported
symmetric tensors and \(\mathcal E_X\) the test vector fields.
For a fixed metric \(g\), define
\[
 K_gX=\mathcal L_Xg.                                        \tag{86.13}
\]
A stress distribution \(T\) is weakly conserved when
\[
 T(K_gX)=0\quad(X\in\mathcal E_X).                           \tag{86.14}
\]

\begin{corollary}[Conservation from a dense vector-field ledger]
\label{cor:v13v86-conservation}
Suppose \(X\mapsto K_gX\) is continuous between the chosen test
topologies, \(\mathcal E_{X,0}\) is dense in
\(\mathcal E_X\), and
\[
 T_0(K_gX)\phi=0
 \quad(X\in\mathcal E_{X,0},\ \phi\in\mathcal D).            \tag{86.15}
\]
If \(T_0\) satisfies Theorem~\ref{thm:v13v86-extension}, then
the extended operator-valued stress distribution obeys
\[
 T(K_gX)\phi=0
 \quad(X\in\mathcal E_X,\ \phi\in\mathcal D).                \tag{86.16}
\]

\begin{proof}
Choose \(X_k\in\mathcal E_{X,0}\) converging to \(X\).
Continuity of \(K_g\) gives \(K_gX_k\to K_gX\), and
Theorem~\ref{thm:v13v86-extension} gives
\(T(K_gX_k)\phi\to T(K_gX)\phi\).  Every term on the left is zero
by (86.15).
\end{proof}
\end{corollary}

This supplies the operator-distribution version of the
diffeomorphism channel in V69.  The transverse Einstein residual
still requires occurrence and convergence of the geometric
Einstein distribution on the same test and graph domains.

\subsection{Finite-to-continuum landing}

Let \(\Phi_n\) be regulated maps on a transported common graph
domain.  A sufficient finite ledger consists of:

\begin{enumerate}
\item exact linearity on the countable smearing space;
\item uniform versions of (86.3) and (86.8);
\item pointwise graph convergence
\(\Phi_n(f)\phi\to\Phi_0(f)\phi\) on countable dense
\(f,\phi\) ledgers;
\item adjoint, covariance, graded-locality, and conservation
residuals tending to zero at every fixed screen;
\item V84 resolvent landing for those real smearings claimed to
be self-adjoint and affiliated;
\item cofinal support and graph-domain exhaustions.
\end{enumerate}

Uniform seminorm bounds extend convergence from the countable
ledger to every test function exactly as in the proof of
Theorem~\ref{thm:v13v86-extension}.  Without those bounds,
cutoff oscillations can vanish at every fixed ledger smearing
while diverging between them, just as the symmetry characters in
V79.

\begin{assessment}[V86 advance]
\label{assess:v13v86}
A countable smearing ledger now extends to a genuine
common-domain operator-valued distribution under explicit
test-seminorm estimates.  Every smearing is closable; covariance,
graded locality, and weak stress conservation pass from
support-dense ledgers by graph continuity.  Local self-adjoint
affiliation remains correctly separated as the bounded-resolvent
gate of V84.
\end{assessment}

\begin{firewall}[Claim boundary after V86]
\label{fire:v13v86}
No uniform SM field seminorm bound, common invariant graph
domain, smearing convergence, stress conservation ledger, or
local resolvent affiliation has been populated.  The geometric
Einstein distribution and transverse equation remain open.  No
operator-valued SM field theory or full SM--Einstein continuum is
proved.
\end{firewall}

\subsection{Parent record}

This V86 note appends the preceding material to the validated V85
source, whose record was
\[
\begin{array}{c|c}
 \text{lines}&26091\\
 \text{bytes}&920171\\
 \text{SHA--256}&
 \texttt{AD0198F491D5101E658653052FA7B3837D00366C46BF3703A213C6575153386B}.
\end{array}                                                     \tag{86.17}
\]

% =============================================================================
% END APPENDED THIRTEENTH-CYCLE V86 MATERIAL
% =============================================================================

% =============================================================================
% BEGIN APPENDED THIRTEENTH-CYCLE V87 MATERIAL
% =============================================================================

\section{Thirteenth-cycle V87: the unbounded-field Wightman
hierarchy}
\label{sec:v13v87}

\subsection{The additional all-order bound}

Let \(\mathcal E=\mathcal S(\mathbb R^{1+d};\mathbb C^r)\), let
\(\Phi\) be the operator-valued distribution obtained in V86,
and let the common graph domain \(\mathcal D\) contain a normalized
vacuum \(\Omega\).  Assume \(\Phi(f)\mathcal D\subset\mathcal D\).
For every \(n\geq1\), require a Schwartz seminorm \(p_{m_n}\) and
a constant \(C_n\) such that
\[
 \|\Phi(f_1)\cdots\Phi(f_n)\Omega\|
 \leq C_n\prod_{j=1}^np_{m_n}(f_j).                          \tag{87.1}
\]
This is an all-order product estimate; the one-field estimate
(86.3) does not imply it.

Define
\[
 W_n(f_1,\ldots,f_n)
 :=
 \langle\Omega,
 \Phi(f_1)\cdots\Phi(f_n)\Omega\rangle.                       \tag{87.2}
\]

\begin{theorem}[Existence of the tempered hierarchy]
\label{thm:v13v87-existence}
Under (87.1), \(W_n\) is a continuous \(n\)-linear functional on
\(\mathcal E^n\).  Hence there is a unique tempered distribution
on \((\mathbb R^{1+d})^n\), with the finite internal indices
included, whose smearing is (87.2).

\begin{proof}
Cauchy--Schwarz and (87.1) give
\[
 |W_n(f_1,\ldots,f_n)|
 \leq C_n\prod_{j=1}^np_{m_n}(f_j).                          \tag{87.3}
\]
Thus \(W_n\) is jointly continuous.  The Schwartz space is
nuclear, so its completed projective tensor power is canonically
the Schwartz space on the Cartesian product.  The Schwartz
kernel theorem turns the continuous multilinear form into the
claimed unique tempered distribution.
\end{proof}
\end{theorem}

\subsection{Wightman properties already forced by the landed
structure}

Assume the adjoint relation (86.5), the V79 Poincar\'e covariance
on \(\mathcal D\), the V78 forward-cone spectrum, and the V86
graded-locality identity.

\begin{theorem}[Conditional Wightman hierarchy]
\label{thm:v13v87-wightman}
The distributions of
Theorem~\ref{thm:v13v87-existence} satisfy:

\begin{enumerate}
\item Hermitian reversal;
\item Poincar\'e covariance for the declared finite-dimensional
field multiplet;
\item graded locality at spacelike-separated supports;
\item the Wightman positivity condition on the Borchers test
algebra;
\item translation spectral support in
\(\overline V_+^{\,n-1}\) in consecutive difference variables.
\end{enumerate}

\begin{proof}
The adjoint relation, applied successively on the invariant common
domain, gives
\[
 \overline{W_n(f_1,\ldots,f_n)}
 =
 W_n(\overline f_n,\ldots,\overline f_1)                    \tag{87.4}
\]
with the internal adjoint indices reversed.  Conjugating every
factor by the Poincar\'e representation and using
\(U(g)\Omega=\Omega\) proves covariance.  Graded commutation of
adjacent spacelike-supported factors proves locality; continuity
extends it from simple smearings to the distribution.

For a finite Borchers-algebra test polynomial \(F\), let
\(\Phi(F)\Omega\) be the corresponding finite sum of field
words.  Then the quadratic functional assembled from the
\(W_n\) satisfies
\[
 W(F^*\otimes F)
 =\|\Phi(F)\Omega\|^2\geq0.                                 \tag{87.5}
\]
All vectors exist by the common invariant domain and (87.1).

Finally, translation covariance and vacuum invariance rewrite
the unsmeared distribution in difference variables in the form
\[
 \langle\Omega,
 \Phi_1U(y_1)\Phi_2U(y_2)\cdots
 U(y_{n-1})\Phi_n\Omega\rangle,                              \tag{87.6}
\]
understood after smearing.  Each intermediate field word maps
the common domain into \(\mathcal H\).  Smearing one \(y_j\)
with a test whose Fourier transform is supported outside
\(\overline V_+\) inserts the zero operator by the V78 joint
spectral calculus.  Tensor-product density of Schwartz tests
then gives support in \(\overline V_+^{\,n-1}\).
\end{proof}
\end{theorem}

This is a conditional Wightman hierarchy for the selected
field multiplet.  It does not assert uniqueness of the vacuum,
irreducibility, asymptotic completeness, or that the multiplet
contains every physical field.

\subsection{Why one-field bounds do not control products}

\begin{proposition}[A one-point field may have no two-point
vacuum word]
\label{prop:v13v87-product-no-go}
There is a self-adjoint operator \(A\) and a normalized vector
\(\Omega\in\operatorname{Dom}A\) for which
\(A\Omega\notin\operatorname{Dom}A\).

\begin{proof}
On \(\ell^2(\mathbb N)\), let
\[
 Ae_k=ke_k,\qquad
 \Omega=c\sum_{k\geq1}k^{-2}e_k,                             \tag{87.7}
\]
where \(c\) normalizes the vector.  Then
\[
 \sum k^2|\Omega_k|^2=c^2\sum k^{-2}<\infty,
\]
so \(\Omega\in\operatorname{Dom}A\).  But
\[
 \sum k^4|\Omega_k|^2=c^2\sum1=\infty,
\]
so \(A\Omega\notin\operatorname{Dom}A\).

For any nonzero continuous scalar functional \(\ell(f)\), the
smearings \(\Phi(f)=\ell(f)A\) are self-adjoint and have
well-defined one-field vacuum vectors, while
\(\Phi(f)\Phi(g)\Omega\) is undefined when both coefficients are
nonzero.
\end{proof}
\end{proposition}

Thus local affiliation of every individual smearing in V84 and
distributional continuity in V86 still require the independent
all-order domain estimate (87.1).

\subsection{Cofinal finite hierarchy}

At cutoff \(N\), word order \(n\), region screen \(R\), and test
seminorm index \(m\), retain
\[
 C_{N,n,R,m}
 :=
 \sup_{\substack{p_m(f_j)\leq1\\
                  \operatorname{supp}f_j\subset R}}
 \|\Phi_N(f_1)\cdots\Phi_N(f_n)\Omega_N\|.                   \tag{87.8}
\]
For each fixed \(n,R\), a uniform finite value and convergence on
a countable dense tensor-product test ledger give the limiting
bound (87.1).  Positivity is tested on finite Borchers Grams,
locality on the V82 region pairs, covariance on the V79 dense
group, and spectral support with the V78 joint calculus.

The limits must be diagonalized simultaneously in
\[
 (n,R,m,\text{test tuple},\text{source pair}).               \tag{87.9}
\]
Proposition~\ref{prop:v13v68-finite-order} forbids stopping at
any fixed word order, and
Proposition~\ref{prop:v13v87-product-no-go} forbids inferring
higher orders from the one-field bound.

\begin{assessment}[V87 advance]
\label{assess:v13v87}
An explicit all-order product-seminorm condition now upgrades the
V86 field distribution to tempered Wightman distributions.
Adjoints, Poincar\'e covariance, graded locality, positivity, and
forward-cone support follow from the structures landed in
V78--V86.  A diagonal-operator counterexample proves that
individual self-adjoint smearings and one-field bounds do not
supply even the two-field vacuum word.
\end{assessment}

\begin{firewall}[Claim boundary after V87]
\label{fire:v13v87}
No all-order SM product bounds, common invariant vacuum domain,
Borchers positivity Grams, or cofinal field-word convergence has
been populated.  Vacuum uniqueness, clustering, LSZ limits,
field completeness, and the transverse Einstein equation remain
open.  No full SM--Einstein Wightman theory or continuum is
proved.
\end{firewall}

\subsection{Parent record}

This V87 note appends the preceding material to the validated V86
source, whose record was
\[
\begin{array}{c|c}
 \text{lines}&26369\\
 \text{bytes}&929821\\
 \text{SHA--256}&
 \texttt{430CC264AAB6AC64A5284AF18983D0DA0339FA2C787D5EEE74E201B155BA1569}.
\end{array}                                                     \tag{87.10}
\]

% =============================================================================
% END APPENDED THIRTEENTH-CYCLE V87 MATERIAL
% =============================================================================

% =============================================================================
% BEGIN APPENDED THIRTEENTH-CYCLE V88 MATERIAL
% =============================================================================

\section{Thirteenth-cycle V88: vacuum uniqueness and averaged
clustering}
\label{sec:v13v88}

\subsection{The invariant projection}

Let \(U(a)\), \(a\in\mathbb R^{1+d}\), be the strongly continuous
translation representation landed in V78--V79, and let
\(\Omega\) be invariant.  Denote by
\[
 P_{\rm inv}
 :=\mathsf E(\{(0,0)\}),\qquad
 P_\Omega:=|\Omega\rangle\langle\Omega|,                     \tag{88.1}
\]
the joint translation-invariant projection and the chosen vacuum
line.  Since \(P_\Omega\leq P_{\rm inv}\),
\[
 Q_{\rm vac}:=P_{\rm inv}-P_\Omega                           \tag{88.2}
\]
is a positive projection.  Vacuum uniqueness is exactly
\(Q_{\rm vac}=0\).

Let \(F_R\) be a tempered Følner sequence of bounded measurable
sets in \(\mathbb R^{1+d}\), and define
\[
 M_R:={1\over|F_R|}\int_{F_R}U(a)\,da.                       \tag{88.3}
\]

\begin{theorem}[Mean-ergodic vacuum projection]
\label{thm:v13v88-ergodic}
One has
\[
 M_R\longrightarrow P_{\rm inv}\quad\text{strongly}.         \tag{88.4}
\]
Consequently, for bounded observables \(A,B\),
\[
\begin{split}
 &\lim_R{1\over|F_R|}\int_{F_R}
 \bigl(\omega(A\alpha_a(B))-\omega(A)\omega(B)\bigr)\,da\\
 &\qquad=
 \langle A^*\Omega,
 Q_{\rm vac}(B-\omega(B)I)\Omega\rangle.                     \tag{88.5}
\end{split}
\]
If the vacuum is unique, every averaged connected correlation
vanishes.

\begin{proof}
The mean ergodic theorem for the unitary representation of the
amenable group \(\mathbb R^{1+d}\) gives strong convergence of
the Følner averages to the projection onto invariant vectors,
which is (88.4).  Vacuum invariance rewrites the averaged first
term as
\[
 \langle A^*\Omega,M_RB\Omega\rangle.                        \tag{88.6}
\]
Subtracting the product of expectations and passing to (88.4)
gives
\[
 \langle A^*\Omega,
 (P_{\rm inv}-P_\Omega)(B-\omega(B)I)\Omega\rangle,
\]
which is (88.5).  It vanishes when \(Q_{\rm vac}=0\).
\end{proof}
\end{theorem}

This is averaged clustering, not pointwise spacelike clustering.
It requires neither a mass gap nor exponential decay and
therefore remains compatible with photon and graviton sectors.

\subsection{A source-complete uniqueness Gram}

Let \(B_R:E_R\to\mathcal H\) be a directed family of bounded
source maps with dense union of ranges.  Define
\[
 \mathbb G_R^{\rm vac}
 :=B_R^*Q_{\rm vac}B_R.                                      \tag{88.7}
\]

\begin{theorem}[Cofinal vacuum-uniqueness criterion]
\label{thm:v13v88-gram}
The following are equivalent:
\[
 P_{\rm inv}=P_\Omega,                                       \tag{88.8}
\]
and
\[
 \mathbb G_R^{\rm vac}=0\quad\text{for every }R.             \tag{88.9}
\]

\begin{proof}
Equation (88.8) makes every Gram zero.  Conversely, positivity
gives
\[
 \mathbb G_R^{\rm vac}
 =(Q_{\rm vac}B_R)^*(Q_{\rm vac}B_R).                        \tag{88.10}
\]
Thus (88.9) implies \(Q_{\rm vac}B_R=0\) for all \(R\).
The union of the source ranges is dense and \(Q_{\rm vac}\) is
bounded, so \(Q_{\rm vac}=0\).
\end{proof}
\end{theorem}

If the joint spectrum lies in \(\overline V_+\), then
\(H=0\) already forces \(P=0\).  In that case
\(P_{\rm inv}={\bf1}_{\{0\}}(H)\), and the uniqueness Gram may
be acquired from the zero-energy spectral projection alone.
Before the cone gate has landed, energy-zero and
translation-invariant sectors must not be identified.

\subsection{Restricted clustering can miss a dark vacuum}

\begin{proposition}[A source-invisible invariant sector]
\label{prop:v13v88-dark-vacuum}
A chosen source algebra can have vanishing connected correlations
while the global invariant subspace is more than one
dimensional.

\begin{proof}
Let \(\mathcal H=\mathbb C^2\), \(U(a)=I\), and
\(\Omega=e_1\).  Take the tested algebra to be
\(\mathbb CI\).  Every connected correlation of tested
observables vanishes identically.  Nevertheless
\[
 P_{\rm inv}=I,\qquad
 Q_{\rm vac}=|e_2\rangle\langle e_2|\neq0.                   \tag{88.11}
\]
The missing invariant direction is invisible to the restricted
source algebra.
\end{proof}
\end{proposition}

Thus cluster tests must be paired with cyclic/source-complete
coverage.  A finite list of observables cannot establish global
vacuum uniqueness merely because its connected correlations are
small.

\subsection{Finite spectral acquisition}

At cutoff \(n\), exact zero-energy projections are unstable when
small eigenvalues approach zero.  Use a shrinking spectral
window
\[
 P_{n,\delta_n}
 :={\bf1}_{[0,\delta_n]}(H_n),\qquad\delta_n\downarrow0,      \tag{88.12}
\]
together with the vacuum line \(P_{\Omega_n}\), and record
\[
 \mathbb G^{\rm vac}_{n,R}
 :=
 B_{n,R}^*
 (P_{n,\delta_n}-P_{\Omega_n})
 B_{n,R}.                                                    \tag{88.13}
\]
This expression is positive only after
\(P_{\Omega_n}\leq P_{n,\delta_n}\) has been certified.  To land
(88.7), one needs strong convergence of the window projections
on fixed source screens and exclusion of spectral mass in the
shrinking annulus.  Strong-resolvent convergence alone does not
give spectral-projection convergence at an atom.

An alternative avoids discontinuous projectors.  Choose
continuous functions \(0\leq f_k\leq1\) with
\(f_k(0)=1\) and support shrinking to zero.  The V78 joint
functional calculus lands
\[
 B_R^*\bigl(f_k(H,P)-P_\Omega\bigr)^2B_R,                    \tag{88.14}
\]
and the limits are taken first in the cutoff, then the source
screen, then \(k\to\infty\).

\subsection{Relation to pointwise clustering and gaps}

Vacuum uniqueness gives (88.5), but not pointwise decay.  A
massive-channel gap can produce exponential Euclidean decay,
while massless sectors can prevent a global exponential bound.
Those channel statements already occur in the Grand Tensor
spectral ladder and are not repeated here.  The new Gram
(88.7) tests only degeneracy at joint momentum zero.

\begin{assessment}[V88 advance]
\label{assess:v13v88}
Vacuum uniqueness is now isolated as the vanishing of a positive
invariant-sector projection.  Følner averages of connected
correlations converge exactly to that dark-vacuum component, and
a cofinal source Gram detects it.  The construction makes no
mass-gap claim and remains compatible with the required massless
photon and graviton channels.
\end{assessment}

\begin{firewall}[Claim boundary after V88]
\label{fire:v13v88}
No invariant projection, vacuum line, shrinking spectral-window
control, source-complete vacuum Gram, or averaged clustering row
has been populated for the coupled theory.  Pointwise
clustering, phase selection, and mass-channel spectra remain
open.  No unique SM--Einstein vacuum or full continuum is proved.
\end{firewall}

\subsection{Parent record}

This V88 note appends the preceding material to the validated V87
source, whose record was
\[
\begin{array}{c|c}
 \text{lines}&26590\\
 \text{bytes}&937546\\
 \text{SHA--256}&
 \texttt{5C84F5526A54E65A0F8C11B235929F92C0F1B32AFDDF0485312F2E07BB98F257}.
\end{array}                                                     \tag{88.15}
\]

% =============================================================================
% END APPENDED THIRTEENTH-CYCLE V88 MATERIAL
% =============================================================================

% =============================================================================
% BEGIN APPENDED THIRTEENTH-CYCLE V89 MATERIAL
% =============================================================================

\section{Thirteenth-cycle V89: vacuum purity and the central
metric-mixture obstruction}
\label{sec:v13v89}

\subsection{Purity is separate from invariant-vector uniqueness}

Let \(\omega\) be a state on a unital \(C^*\)-algebra
\(\mathfrak A\), with GNS triple
\((\pi_\omega,\mathcal H_\omega,\Omega_\omega)\).

\begin{theorem}[GNS purity criterion]
\label{thm:v13v89-purity}
The state \(\omega\) is pure if and only if
\[
 \pi_\omega(\mathfrak A)'=\mathbb CI.                        \tag{89.1}
\]

\begin{proof}
If the commutant contains a nonscalar self-adjoint element, its
spectral calculus contains a nontrivial positive contraction
\(0<C<I\).  The positive functionals
\[
 \omega_1(A)=
 {\langle\Omega_\omega,C\pi_\omega(A)\Omega_\omega\rangle
  \over\langle\Omega_\omega,C\Omega_\omega\rangle},
\quad
 \omega_2(A)=
 {\langle\Omega_\omega,(I-C)\pi_\omega(A)\Omega_\omega\rangle
  \over\langle\Omega_\omega,(I-C)\Omega_\omega\rangle}       \tag{89.2}
\]
give a nontrivial convex decomposition of \(\omega\).

Conversely, a nontrivial convex decomposition
\(\omega=t\omega_1+(1-t)\omega_2\) yields, by the
Radon--Nikodym theorem for positive functionals dominated by
\(\omega\), a nonscalar positive contraction in the commutant.
Thus irreducibility is equivalent to extremality.
\end{proof}
\end{theorem}

Theorem~\ref{thm:v13v88-gram} tests the dimension of the
translation-invariant subspace.  It does not test the full
commutant (89.1).

\subsection{A unique invariant vector with a mixed state}

\begin{proposition}[Ergodicity does not imply purity]
\label{prop:v13v89-ergodic-no-pure}
There is a covariant representation with a unique invariant
vector and a non-pure invariant state.

\begin{proof}
Let
\[
 \mathfrak A=C(S^1),\qquad
 \mathcal H=L^2(S^1,d\theta/2\pi),\qquad
 \Omega=1,
\]
with \(\pi(f)\) acting by multiplication.  Let rotations act by
\[
 (U_t\psi)(\theta)=\psi(\theta-t).                            \tag{89.3}
\]
The invariant vectors are exactly the constant functions, so
\(\Omega\) spans the invariant subspace.  The state
\[
 \omega(f)=\int_{S^1}f(\theta)\,{d\theta\over2\pi}            \tag{89.4}
\]
is not pure: it is integration against a non-Dirac probability
measure, and \(\pi(\mathfrak A)'=L^\infty(S^1)\) is nonscalar.
\end{proof}
\end{proposition}

Thus the V88 vacuum Gram can vanish while a classical central
mixture survives.

\subsection{Central metric variance}

Let \(Z(\mathfrak A)\) be the center.  For self-adjoint
\(z\in Z(\mathfrak A)\), define
\[
 \operatorname{Var}_\omega(z)
 :=\omega(z^2)-\omega(z)^2.                                  \tag{89.5}
\]

\begin{proposition}[Pure states are characters on the center]
\label{prop:v13v89-center}
If \(\omega\) is pure, then
\[
 \operatorname{Var}_\omega(z)=0
 \quad(z=z^*\in Z(\mathfrak A)).                             \tag{89.6}
\]

\begin{proof}
By Theorem~\ref{thm:v13v89-purity}, the GNS representation is
irreducible.  Since
\[
 \pi_\omega(z)\in
 \pi_\omega(\mathfrak A)\cap\pi_\omega(\mathfrak A)',
\]
Schur's lemma gives \(\pi_\omega(z)=cI\).  Hence
\(\omega(z^2)-\omega(z)^2=c^2-c^2=0\).
\end{proof}
\end{proposition}

For the direct-integral construction of V72, central metric
observables act as multiplication by functions \(f(x)\).  Their
variance is
\[
 \operatorname{Var}_\mu(f)
 =\int f^2\,d\mu-\left(\int f\,d\mu\right)^2.                \tag{89.7}
\]

\begin{theorem}[A separating zero-variance ledger collapses the
central metric law]
\label{thm:v13v89-central-collapse}
Let \(X\) be Polish and let \(f_1,f_2,\ldots\) be bounded
continuous functions whose joint map
\[
 F:X\to\mathbb R^{\mathbb N},\qquad
 F(x)=(f_1(x),f_2(x),\ldots)                                 \tag{89.8}
\]
is injective.  If a probability measure \(\mu\) satisfies
\[
 \operatorname{Var}_\mu(f_j)=0
 \quad\text{for every }j,                                   \tag{89.9}
\]
then \(\mu\) is a Dirac measure.

\begin{proof}
Zero variance makes \(f_j=c_j\) almost surely.  Countability
gives one full-measure set on which all equalities hold.  Every
point in that set has the same image \((c_j)_j\) under the
injective map \(F\), so the set contains at most one point.
A probability measure supported there is Dirac.
\end{proof}
\end{theorem}

Therefore a nontrivial classical metric probability law in the
center cannot coexist with purity of the full represented
observable state.  The programme must choose among distinct
structures: a mixed state over classical backgrounds, a selected
ergodic background sector, or genuinely quantum metric
observables that do not remain a classical central algebra.

\subsection{Finite central covariance Gram}

For real central ledger observables \(z_1,\ldots,z_q\), define
\[
 \mathbb C_q^{\rm cen}
 :=
 \bigl(\omega(z_az_b)-\omega(z_a)\omega(z_b)\bigr)_{a,b=1}^q.
                                                                  \tag{89.10}
\]
This matrix is positive semidefinite because
\[
 c^*\mathbb C_q^{\rm cen}c
 =
 \omega\!\left(
 \left|\sum_ac_a(z_a-\omega(z_a)I)\right|^2\right)\geq0.     \tag{89.11}
\]
Its zero on a cofinal separating central ledger forces the
central law to collapse by
Theorem~\ref{thm:v13v89-central-collapse}.  A nonzero matrix is
a direct witness that the state is not pure on any algebra in
which those variables remain central.

The converse fails: zero central covariance removes only the
central obstruction and does not prove the full commutant is
scalar.  Noncentral superselection multiplicities require a
separate irreducibility or sector-selection argument.

\subsection{Compatibility with Einstein fluctuations}

The almost-sure Einstein equation of V69 does not require a
Dirac metric law; it can hold in every classical background
sample.  Purity, however, cannot then be claimed for the enlarged
algebra containing all classical metric multipliers.  Likewise,
nonzero metric variance in a semiclassical ensemble is not by
itself a quantum graviton fluctuation.  Quantum fluctuations must
be represented by noncentral metric operators in a pure or
factorial sector and tested through their local correlation
hierarchy.

\subsection{Acquisition consequence}

A final vacuum statement must separately record:

\begin{enumerate}
\item the V88 invariant-subspace Gram;
\item the central covariance Gram (89.10);
\item the chosen treatment of the metric center under the
diffeomorphism quotient;
\item factoriality or the stronger commutant-scalar purity test;
\item compatibility of sector selection with local covariance,
the Einstein residual, and the all-order Wightman hierarchy.
\end{enumerate}

\begin{assessment}[V89 advance]
\label{assess:v13v89}
Vacuum uniqueness and state purity are now rigorously separated.
A rotation-covariant example has a unique invariant vector but a
mixed Haar state.  Pure states have zero variance on the center,
and a countable separating zero-variance metric ledger collapses
a classical background law to one point.  A nontrivial random
metric continuum must therefore be described honestly as a
mixture, sector-selected, or represented by noncentral quantum
metric observables.
\end{assessment}

\begin{firewall}[Claim boundary after V89]
\label{fire:v13v89}
No central metric ledger, ergodic-sector selection,
factoriality, irreducibility, or quantum metric operator has been
populated.  The V88 vacuum-uniqueness Gram is also still
conditional.  No pure coupled vacuum, quantum-gravitational
sector, or full SM--Einstein continuum is proved.
\end{firewall}

\subsection{Parent record}

This V89 note appends the preceding material to the validated V88
source, whose record was
\[
\begin{array}{c|c}
 \text{lines}&26810\\
 \text{bytes}&944966\\
 \text{SHA--256}&
 \texttt{B13F33826F4187E608F8C5EBC6CFF3DF8795E52D5C9D1D0C3F4CE9B611305048}.
\end{array}                                                     \tag{89.12}
\]

% =============================================================================
% END APPENDED THIRTEENTH-CYCLE V89 MATERIAL
% =============================================================================

% =============================================================================
% BEGIN APPENDED THIRTEENTH-CYCLE V90 MATERIAL
% =============================================================================

\section{Thirteenth-cycle V90: compulsory companion audit before
covariant sector selection}
\label{sec:v13v90}

\subsection{Audit record}

The newest companion notes remain
\[
\begin{array}{c|c|r|r|c}
 \text{programme}&\text{version}&\text{lines}&\text{bytes}&
 \text{SHA--256}\\ \hline
 \text{Yang--Mills mass gap}&V95&30950&1056018&
 \texttt{2DB19DAF571B44B04B28A51FAE2708D87BB6234D9CE9A2FBC5530ADF081809C2}\\
 \text{Navier--Stokes stability}&V26&5319&278283&
 \texttt{82C887F8C2C69E4F28FAD7C3DC8DE5B9099CB5A4BF431EBA1FFF3E91935978AD}.
\end{array}                                                     \tag{90.1}
\]
Auxiliary YM parallel V5 remains unchanged with SHA--256
\[
 \texttt{306F1BB84CFA8A8D47EA569EC17E9545F9867C8D32B548EC9D9C444B4F3C0512}.
                                                                  \tag{90.2}
\]
Their terminal results and firewalls were reread.  There is no
new companion input.  YM V95 still requires an all-order
interacting state beyond its finite response screens; V87 now
states the corresponding SM product-bound gate.  NS V26 supplies
no estimate for the vacuum-sector problem.

\begin{proposition}[V90 direction decision]
\label{prop:v13v90-decision}
V91 shall combine the V72 covariant metric-base action with the
V89 purity theorem.  It will determine when a Dirac central
metric sector is invariant, when an ergodic metric law admits a
proper invariant central sector, and what finite separating
ledger certifies a fixed background class.

\begin{proof}
V89 shows that purity collapses a separating classical metric
center to a point.  Covariance then requires that point to be
fixed by the induced group action.  This additional condition is
not contained in the variance Gram and must be proved before a
pure covariant sector is claimed.
\end{proof}
\end{proposition}

\begin{assessment}[V90 audit outcome]
\label{assess:v13v90}
The audit is complete.  With no changed companion result, the
programme proceeds to the compatibility of purity, central
sector selection, and covariance.
\end{assessment}

\begin{firewall}[Claim boundary after V90]
\label{fire:v13v90}
No invariant pure metric sector or fixed quotient background is
proved in this audit.  The full SM--Einstein continuum remains
open.
\end{firewall}

\subsection{Parent record}

This V90 note appends the preceding material to the validated V89
source, whose record was
\[
\begin{array}{c|c}
 \text{lines}&27049\\
 \text{bytes}&953215\\
 \text{SHA--256}&
 \texttt{6B5981E0518180038315EC063D6FD936A3E49F45F9BE37FD07975E6ACDC88D7F}.
\end{array}                                                     \tag{90.3}
\]

% =============================================================================
% END APPENDED THIRTEENTH-CYCLE V90 MATERIAL
% =============================================================================

% =============================================================================
% BEGIN APPENDED THIRTEENTH-CYCLE V91 MATERIAL
% =============================================================================

\section{Thirteenth-cycle V91: covariant selection of a pure
metric sector}
\label{sec:v13v91}

\subsection{Central metric action}

Let \(X\) be a Polish space of metric classes and let a topological
group \(G\) act continuously on \(X\).  The induced action on
bounded continuous central observables is
\[
 (\alpha_gf)(x)=f(g^{-1}x).                                  \tag{91.1}
\]
Let \(f_1,f_2,\ldots\in C_b(X;\mathbb R)\) be a countable
separating family.

\begin{lemma}[Invariant Dirac laws are fixed points]
\label{lem:v13v91-dirac}
For \(x\in X\),
\[
 g_*\delta_x=\delta_x\quad\text{for every }g\in G            \tag{91.2}
\]
if and only if
\[
 gx=x\quad\text{for every }g\in G.                           \tag{91.3}
\]

\begin{proof}
One has \(g_*\delta_x=\delta_{gx}\).  Two Dirac measures are equal
exactly when their supporting points are equal.
\end{proof}
\end{lemma}

\begin{theorem}[Pure covariant central-metric obstruction]
\label{thm:v13v91-pure-covariant}
Suppose a \(G\)-invariant state \(\omega\) on
\(\mathfrak A\) contains \(C_b(X)\) as a separating central
metric algebra, and suppose its restriction is represented by a
Borel probability measure on \(X\).  If \(\omega\) is pure, then
its restriction to
that algebra is
\[
 \omega(f)=f(x_0)                                            \tag{91.4}
\]
for a point \(x_0\in X\) fixed by all of \(G\).

\begin{proof}
Purity and Proposition~\ref{prop:v13v89-center} make every
central variance vanish.  The separating-family theorem
\ref{thm:v13v89-central-collapse} makes the representing metric
law a Dirac measure \(\delta_{x_0}\).  Invariance of \(\omega\)
makes this Dirac law invariant, and
Lemma~\ref{lem:v13v91-dirac} gives \(gx_0=x_0\) for every \(g\).
\end{proof}
\end{theorem}

If the physical symmetry moves every point of the chosen
classical metric space, a pure invariant state cannot retain that
space as a separating center.  One must first quotient gauge
directions, select a fixed physical background class, accept a
mixed central law, or promote the metric variables to noncentral
quantum observables.

\subsection{Ergodic mixtures are different}

Let \(\mu\) be a \(G\)-invariant probability measure on \(X\).
It is ergodic when every invariant measurable set has measure
zero or one.

\begin{proposition}[No proper invariant central conditioning in
an ergodic law]
\label{prop:v13v91-ergodic}
If \(\mu\) is ergodic and \(A\subset X\) is invariant, then the
central projection \(1_A\) is either zero or the identity in
\(L^\infty(X,\mu)\).  Hence conditioning on an invariant central
event cannot select a proper positive-measure sector.

\begin{proof}
Ergodicity gives \(\mu(A)\in\{0,1\}\).  Therefore the
multiplication operator \(1_A\) equals \(0\) or \(I\) almost
everywhere.  If a measurable set has probability strictly
between zero and one, conditioning on it is nontrivial but the
set is not invariant, so the conditioned law is not
\(G\)-invariant.
\end{proof}
\end{proposition}

An ergodic metric law may still be non-Dirac, as the circle
rotation example in Proposition~\ref{prop:v13v89-ergodic-no-pure}
shows.  Ergodicity can give a unique invariant vector while the
state remains mixed on the central algebra.

\subsection{A finite fixed-sector residual}

For \(g\in G\), define the nonnegative separating residual
\[
 \mathcal R^{\rm fix}_{R,g}(x)
 :=
 \sum_{j=1}^R2^{-j}
 {\,|f_j(gx)-f_j(x)|^2\over
  1+|f_j(gx)-f_j(x)|^2}.                                    \tag{91.5}
\]

\begin{theorem}[Cofinal fixed-point certificate]
\label{thm:v13v91-fixed-ledger}
Let \(D\subset G\) be dense.  If
\[
 \mathcal R^{\rm fix}_{R,g}(x_0)=0
 \quad\text{for every }R\text{ and }g\in D,                 \tag{91.6}
\]
then \(x_0\) is fixed by all of \(G\).

\begin{proof}
For fixed \(g\in D\), vanishing for every \(R\) gives
\[
 f_j(gx_0)=f_j(x_0)\quad\text{for every }j.
\]
The functions separate points, so \(gx_0=x_0\).  If
\(g_n\in D\) converges to \(g\in G\), continuity of the action
gives
\[
 gx_0=\lim_ng_nx_0=x_0.                                     \tag{91.7}
\]
\end{proof}
\end{theorem}

For a metric law, the integrated positive row is
\[
 \varepsilon^{\rm fix}_{R,g}
 :=
 \int_X\mathcal R^{\rm fix}_{R,g}(x)\,d\mu(x).               \tag{91.8}
\]
Its zero on a cofinal separating ledger says that \(g\) fixes
\(\mu\)-almost every metric point, which is stronger than mere
invariance \(g_*\mu=\mu\).  For a pure central sector, V89 first
forces \(\mu=\delta_{x_0}\), and (91.8) then reduces exactly to
the fixed-point test (91.6).

\subsection{Gauge quotient versus physical symmetry}

The fixed-point condition must be imposed on the correct base.
If \(X_{\rm raw}\) is the space of metrics before
diffeomorphism quotient, gauge transformations generally move a
representative without changing the physical point.  Let
\[
 q:X_{\rm raw}\to X_{\rm phys}                              \tag{91.9}
\]
be the certified quotient map.  A pure physical central sector
requires a fixed point in \(X_{\rm phys}\), not a representative
fixed by every gauge diffeomorphism.  Conversely, declaring all
physical spacetime transformations to be gauge would erase the
covariance content rather than prove (91.3).

Thus the acquisition order is:

\begin{enumerate}
\item certify the diffeomorphism quotient and its separating
physical metric ledger;
\item determine the residual action of physical symmetries on
that quotient;
\item use the V89 central covariance Gram to decide whether the
sector is mixed or Dirac;
\item on a Dirac branch, apply (91.5)--(91.6);
\item on a noncentral quantum-metric branch, replace the classical
fixed-point test by covariance of the metric operator-valued
distribution.
\end{enumerate}

\subsection{Compatibility with Einstein dynamics}

A fixed background class in (91.4) must still satisfy the V69
transverse Einstein residual with the same stress state.  The
fixed-point theorem supplies neither the Einstein equation nor
metric fluctuations.  On the mixed branch, the almost-sure
Einstein Grams may vanish fibrewise while purity fails.  On the
noncentral branch, the metric field must satisfy the V84--V87
affiliation, common-domain, and Wightman product bounds.

\begin{assessment}[V91 advance]
\label{assess:v13v91}
Purity, covariance, and the classical metric center now have an
exact compatibility theorem.  A pure invariant state collapses
a separating central metric law to a Dirac sector, and covariance
then forces that physical metric class to be fixed.  Ergodic
non-Dirac laws admit no proper invariant central conditioning and
remain mixed.  A cofinal positive residual tests the fixed class
on a dense symmetry ledger.
\end{assessment}

\begin{firewall}[Claim boundary after V91]
\label{fire:v13v91}
No certified physical metric quotient, separating quotient
ledger, pure sector, residual symmetry action, fixed background
class, or noncentral quantum metric field has been populated.
The theorem does not choose between mixed, background-selected,
and quantum-metric branches.  No pure covariant SM--Einstein
continuum is proved.
\end{firewall}

\subsection{Parent record}

This V91 note appends the preceding material to the validated V90
source, whose record was
\[
\begin{array}{c|c}
 \text{lines}&27129\\
 \text{bytes}&956254\\
 \text{SHA--256}&
 \texttt{BB80D429BB680F19C68DE8394B52CD11094F8B2694E3C5647AE87E6D4322D6AD}.
\end{array}                                                     \tag{91.10}
\]

% =============================================================================
% END APPENDED THIRTEENTH-CYCLE V91 MATERIAL
% =============================================================================

% =============================================================================
% BEGIN APPENDED THIRTEENTH-CYCLE V92 MATERIAL
% =============================================================================

\section{Thirteenth-cycle V92: a noncentral coframe metric and
its invertibility gate}
\label{sec:v13v92}

\subsection{Algebraic coframe construction}

Let \(\mathfrak A\) be a unital \(C^*\)-algebra of local quantum
observables and let
\[
 \eta=\operatorname{diag}(-1,1,\ldots,1)
 \in M_{d+1}(\mathbb C)\subset M_{d+1}(\mathfrak A).         \tag{92.1}
\]
A bounded noncentral coframe is an element
\[
 E\in M_{d+1}(\mathfrak A).                                 \tag{92.2}
\]
Define its operator-valued metric matrix by
\[
 \boxed{\quad G:=E^*\eta E.\quad}                            \tag{92.3}
\]
Then \(G=G^*\), although its entries need not commute with one
another or with matter observables.

\begin{theorem}[Invertible quantum coframe metric]
\label{thm:v13v92-coframe}
If \(E\) is invertible in \(M_{d+1}(\mathfrak A)\), then \(G\)
is invertible and
\[
 G^{-1}=E^{-1}\eta(E^*)^{-1}.                                \tag{92.4}
\]
If \(L\in M_{d+1}(\mathfrak A)\) satisfies
\[
 L^*\eta L=\eta,                                             \tag{92.5}
\]
then the frame change \(E\mapsto LE\) leaves \(G\) invariant.
A right change \(E\mapsto EU\) gives
\[
 G\mapsto U^*GU.                                             \tag{92.6}
\]

\begin{proof}
Self-adjointness follows from \(\eta^*=\eta\).  Since
\(\eta^2=I\), direct multiplication gives
\[
 (E^*\eta E)
 \bigl(E^{-1}\eta(E^*)^{-1}\bigr)=I
\]
and the reverse product is also \(I\), proving (92.4).
Equation (92.5) gives
\[
 (LE)^*\eta(LE)=E^*L^*\eta LE=E^*\eta E.
\]
The right-change formula is immediate.
\end{proof}
\end{theorem}

This is an algebraic Lorentz-coframe class.  It does not by
itself define pointwise causal cones when the entries are
noncommuting; causal localization then requires a state-dependent
or relational construction compatible with V81--V82.

\subsection{Two-sided coframe floors}

\begin{theorem}[Two-sided \(C^*\)-invertibility certificate]
\label{thm:v13v92-two-sided}
For \(E\in M_{d+1}(\mathfrak A)\), suppose that for some
\(\gamma>0\),
\[
 E^*E\succeq\gamma I,\qquad
 EE^*\succeq\gamma I.                                       \tag{92.7}
\]
Then \(E\) is invertible and
\[
 \|E^{-1}\|\leq\gamma^{-1/2},\qquad
 \|G^{-1}\|\leq\gamma^{-1}.                                 \tag{92.8}
\]

\begin{proof}
The first inequality in (92.7) makes \(E^*E\) invertible and
gives the left inverse
\[
 L=(E^*E)^{-1}E^*,\qquad LE=I.                              \tag{92.9}
\]
The second makes \(EE^*\) invertible and gives the right inverse
\[
 R=E^*(EE^*)^{-1},\qquad ER=I.                              \tag{92.10}
\]
Then \(L=L(ER)=(LE)R=R\), so \(E\) is invertible.  Functional
calculus gives the first norm bound.  Formula (92.4) and
\(\|\eta\|=1\) give the second.
\end{proof}
\end{theorem}

Both floors are needed in a general infinite \(C^*\)-algebra.
For example, the unilateral shift \(S\in B(\ell^2)\) satisfies
\(S^*S=I\) but is not onto and hence is not invertible.

\subsection{Perturbation from an invertible anchor}

\begin{theorem}[Coframe Neumann cell]
\label{thm:v13v92-neumann}
Let \(E_0\) be invertible and suppose
\[
 \|E-E_0\|\leq\delta,\qquad
 \rho:=\|E_0^{-1}\|\delta<1.                                \tag{92.11}
\]
Then \(E\) and \(G=E^*\eta E\) are invertible, with
\[
 \|E^{-1}\|
 \leq{\|E_0^{-1}\|\over1-\rho},\qquad
 \|G^{-1}\|
 \leq{\|E_0^{-1}\|^2\over(1-\rho)^2}.                       \tag{92.12}
\]
Moreover,
\[
 \|G-G_0\|
 \leq\delta(\|E\|+\|E_0\|).                                 \tag{92.13}
\]

\begin{proof}
Factor
\[
 E=E_0\bigl(I+E_0^{-1}(E-E_0)\bigr).                        \tag{92.14}
\]
The second factor is invertible by the Neumann series, whose norm
bound gives the first inequality in (92.12).  Theorem
\ref{thm:v13v92-coframe} gives the second.  Finally,
\[
 G-G_0
 =(E-E_0)^*\eta E+E_0^*\eta(E-E_0),                         \tag{92.15}
\]
and the \(C^*\)-norm proves (92.13).
\end{proof}
\end{theorem}

This supplies a finite background-cell compiler analogous to
V76.  The coframe norm radius must be measured in the physical
local algebra, not only after taking a state expectation.

\subsection{Expectation nondegeneracy is not operator
nondegeneracy}

\begin{proposition}[A state can miss a coframe kernel]
\label{prop:v13v92-expectation-no-go}
There is a positive noninvertible operator whose expectation in
a faithful-looking selected vector is strictly positive.

\begin{proof}
On \(\mathbb C^2\), let
\[
 P=\begin{pmatrix}1&0\\0&0\end{pmatrix},\qquad\Omega=e_1.
                                                                    \tag{92.16}
\]
Then \(\langle\Omega,P\Omega\rangle=1\), but \(P\) has kernel
\(\mathbb Ce_2\) and is not invertible.  Embedding \(P\) as one
coframe block shows that nondegeneracy of a chosen expectation
does not imply either floor in (92.7).
\end{proof}
\end{proposition}

Source-complete positive Grams may detect the missing block, but
a single vacuum mean cannot.

\subsection{Covariance and noncentral fluctuations}

If \(\alpha_g\) is a \(C^*\)-automorphism acting entrywise and
the coframe transforms by
\[
 \alpha_g(E)=E_gU_g,                                        \tag{92.17}
\]
then
\[
 \alpha_g(G)=U_g^*G_gU_g                                   \tag{92.18}
\]
by Theorem~\ref{thm:v13v92-coframe}.  This is the raw
operator-valued covariance row required before applying the V79
dense-group extension.

Because the entries of \(G\) need not be central, a pure state
may have
\[
 \omega(G_{\mu\nu}^2)-\omega(G_{\mu\nu})^2>0                \tag{92.19}
\]
without contradicting V89.  Such variance is only a candidate
quantum metric fluctuation.  It becomes physical after quotient,
localization, covariance, positivity, and Einstein-incidence
rows have all landed.

\subsection{Differential and Einstein gates}

Curvature is not defined by (92.3) alone.  A noncentral Einstein
writer additionally needs:

\begin{enumerate}
\item a common differential calculus or derivations on the local
algebra;
\item domains stable under \(E,E^{-1}\) and the required
derivatives;
\item an ordering prescription for connection and curvature
products;
\item frame and diffeomorphism covariance of that prescription;
\item a Bianchi identity with an explicitly controlled
noncommutative defect;
\item incidence with the same stress distribution constructed in
V84--V87;
\item the transverse positive residual Grams of V69--V71.
\end{enumerate}

\begin{assessment}[V92 advance]
\label{assess:v13v92}
The noncentral quantum-metric branch now has a rigorous bounded
starting point.  An invertible operator-valued coframe produces
a self-adjoint invertible metric matrix, with exact frame
covariance.  Two-sided \(C^*\)-floors and a Neumann cell provide
finite invertibility certificates.  A two-dimensional
counterexample proves that expectation-level nondegeneracy cannot
replace operator invertibility.
\end{assessment}

\begin{firewall}[Claim boundary after V92]
\label{fire:v13v92}
No local quantum coframe, two-sided floor, differential calculus,
curvature ordering, noncommutative Bianchi identity, or Einstein
incidence has been populated.  The algebraic metric does not
define causal cones or a quantum Einstein dynamics.  No
noncentral quantum-gravity sector or full SM--Einstein continuum
is proved.
\end{firewall}

\subsection{Parent record}

This V92 note appends the preceding material to the validated V91
source, whose record was
\[
\begin{array}{c|c}
 \text{lines}&27348\\
 \text{bytes}&964326\\
 \text{SHA--256}&
 \texttt{818B1E0CAD2814BA7DB3AE9F9D255949AFBA0DF895C4ED0EEB6A0D2AF187AF7A}.
\end{array}                                                     \tag{92.20}
\]

% =============================================================================
% END APPENDED THIRTEENTH-CYCLE V92 MATERIAL
% =============================================================================

% =============================================================================
% BEGIN APPENDED THIRTEENTH-CYCLE V93 MATERIAL
% =============================================================================

\section{Thirteenth-cycle V93: differentiating the noncentral
coframe inverse}
\label{sec:v13v93}

\subsection{Smooth inverse-closed algebra}

Let \(\mathfrak A_\infty\subset\mathfrak A\) be a dense unital
\(*\)-subalgebra, equipped with \(*\)-derivations
\[
 \delta_a:\mathfrak A_\infty\to\mathfrak A_\infty,
 \qquad 1\leq a\leq r,                                      \tag{93.1}
\]
and assume \(\mathfrak A_\infty\) is inverse closed in
\(\mathfrak A\).  Extend the derivations entrywise to matrix
algebras.  Let
\[
 E\in M_m(\mathfrak A_\infty)                               \tag{93.2}
\]
be invertible in \(M_m(\mathfrak A)\), and put \(F=E^{-1}\).

\begin{theorem}[First and second inverse derivatives]
\label{thm:v13v93-inverse}
The inverse belongs to \(M_m(\mathfrak A_\infty)\) and
\[
 \boxed{\quad
 \delta_aF=-F(\delta_aE)F.
 \quad}                                                       \tag{93.3}
\]
Whenever \(\delta_a\delta_bE\) is defined,
\[
\boxed{
\begin{split}
 \delta_a\delta_bF
 ={}&F(\delta_aE)F(\delta_bE)F\\
 &+F(\delta_bE)F(\delta_aE)F
 -F(\delta_a\delta_bE)F.
\end{split}}                                                  \tag{93.4}
\]

\begin{proof}
Inverse closedness places \(F\) in the smooth algebra.
Differentiating \(EF=I\) gives
\[
 (\delta_aE)F+E(\delta_aF)=0,
\]
and left multiplication by \(F\) proves (93.3).  Apply
\(\delta_a\) to
\(\delta_bF=-F(\delta_bE)F\), use the Leibniz rule, and substitute
(93.3) in the two differentiated inverse factors.  The three
resulting terms are exactly (93.4).
\end{proof}
\end{theorem}

No commutativity of the entries is used.  The order of every
factor in (93.3)--(93.4) is load bearing.

\begin{corollary}[Quantitative inverse-jet bounds]
\label{cor:v13v93-bounds}
If
\[
 \|F\|\leq K,\qquad
 \|\delta_aE\|\leq L_a,\qquad
 \|\delta_a\delta_bE\|\leq L_{ab},                           \tag{93.5}
\]
then
\[
 \|\delta_aF\|\leq K^2L_a,                                  \tag{93.6}
\]
\[
 \|\delta_a\delta_bF\|
 \leq2K^3L_aL_b+K^2L_{ab}.                                  \tag{93.7}
\]

\begin{proof}
Apply the \(C^*\)-norm and submultiplicativity to
(93.3)--(93.4).
\end{proof}
\end{corollary}

\subsection{Metric jets}

Let \(G=E^*\eta E\) as in V92 and assume the derivations commute
with the involution.  Then
\[
 \delta_aG
 =(\delta_aE)^*\eta E+E^*\eta(\delta_aE),                    \tag{93.8}
\]
and
\[
\begin{split}
 \delta_a\delta_bG
 ={}&(\delta_a\delta_bE)^*\eta E
   +(\delta_bE)^*\eta(\delta_aE)\\
 &+(\delta_aE)^*\eta(\delta_bE)
   +E^*\eta(\delta_a\delta_bE).                              \tag{93.9}
\end{split}
\]
Consequently,
\[
 \|\delta_aG\|
 \leq2\|E\|L_a,\qquad
 \|\delta_a\delta_bG\|
 \leq2\|E\|L_{ab}+2L_aL_b.                                  \tag{93.10}
\]
The inverse metric is
\[
 G^{-1}=F\eta F^*,                                          \tag{93.11}
\]
so its first and second derivative bounds follow from
Corollary~\ref{cor:v13v93-bounds} and the same product rule.

These estimates are the bounded differential input for any
ordered connection or curvature polynomial.  They do not select
that ordering.

\subsection{Landing inverse jets}

\begin{theorem}[Norm landing of differentiated inverses]
\label{thm:v13v93-landing}
Let \(E_n,E_n^{-1}\in M_m(\mathfrak A_\infty)\) satisfy
\[
 E_n\to E,\qquad
 \sup_n\|E_n^{-1}\|\leq K,                                  \tag{93.12}
\]
\[
 \delta_aE_n\to X_a,\qquad
 \delta_a\delta_bE_n\to X_{ab}                              \tag{93.13}
\]
in norm.  Then \(E\) is invertible,
\[
 E_n^{-1}\to E^{-1},                                        \tag{93.14}
\]
and
\[
 \delta_a(E_n^{-1})
 \to-E^{-1}X_aE^{-1},                                       \tag{93.15}
\]
with the corresponding second derivative given by (93.4) using
\(X_a,X_b,X_{ab}\).  If the derivations are closed and
\(E_n\to E\) in their graph topology, these limits are the
actual derivatives of \(E^{-1}\).

\begin{proof}
The uniform inverse bound and norm convergence imply
invertibility of \(E\) for all sufficiently small perturbations,
and the resolvent identity
\[
 E_n^{-1}-E^{-1}
 =E_n^{-1}(E-E_n)E^{-1}                                     \tag{93.16}
\]
gives (93.14).  Substitute (93.12)--(93.14) into the ordered
products (93.3)--(93.4) to obtain the derivative limits.
Closedness of the derivations identifies graph limits with the
derivatives of the limiting element.
\end{proof}
\end{theorem}

The uniform inverse bound in (93.12) is essential.  Scalar
elements \(E_n=n^{-1}\) converge to zero while every \(E_n\) is
invertible and the inverse norms diverge.

\subsection{Why a polynomial differential ledger is not enough}

\begin{proposition}[Invertibility need not stay in a non-spectral
subalgebra]
\label{prop:v13v93-inverse-closed-no-go}
Let \(\mathfrak P\) be the polynomial subalgebra of
\(C([0,1])\).  The element
\[
 E(x)=2+x                                                     \tag{93.17}
\]
belongs to \(\mathfrak P\), is invertible in \(C([0,1])\), but
\(E^{-1}(x)=(2+x)^{-1}\notin\mathfrak P\).

\begin{proof}
The function is bounded below by two, hence invertible in the
ambient \(C^*\)-algebra.  Its reciprocal is not a polynomial,
since a polynomial identity
\((2+x)p(x)=1\) is impossible by degree.
\end{proof}
\end{proposition}

Thus a finite algebra closed under differentiation and products
can still lose the inverse needed for the connection.  The
continuum smooth algebra must be spectrally invariant or carry a
separate inverse-closure theorem.

\subsection{Finite acquisition rows}

For each local coframe screen retain:

\begin{enumerate}
\item both V92 invertibility floors;
\item membership of \(E^{-1}\) in the smooth algebra;
\item first and second coframe derivative stocks;
\item the ordered identities (93.3)--(93.4);
\item graph convergence for the derivations;
\item covariance of the derivations under the V79 symmetry
action;
\item compatibility of the differential algebra with local
affiliation and the common field domain of V84--V87.
\end{enumerate}

\begin{assessment}[V93 advance]
\label{assess:v13v93}
The noncentral coframe now has a rigorous differential inverse
calculus.  First and second inverse derivatives retain their
noncommutative order and have explicit norm bounds.  Uniform
inverse floors and graph convergence land those jets in the
continuum.  A polynomial example proves that ambient
invertibility does not keep an inverse inside an arbitrary
differential ledger.
\end{assessment}

\begin{firewall}[Claim boundary after V93]
\label{fire:v13v93}
No spectrally invariant SM differential algebra, coframe
derivative stock, inverse-jet graph limit, connection ordering,
curvature, or Bianchi identity has been populated.  No quantum
Einstein writer or full SM--Einstein continuum is proved.
\end{firewall}

\subsection{Parent record}

This V93 note appends the preceding material to the validated V92
source, whose record was
\[
\begin{array}{c|c}
 \text{lines}&27598\\
 \text{bytes}&972379\\
 \text{SHA--256}&
 \texttt{6326D42ADBD582048F8B6711063D599B82957C7F92B31AD38430307A1B4C8EB3}.
\end{array}                                                     \tag{93.18}
\]

% =============================================================================
% END APPENDED THIRTEENTH-CYCLE V93 MATERIAL
% =============================================================================

% =============================================================================
% BEGIN APPENDED THIRTEENTH-CYCLE V94 MATERIAL
% =============================================================================

\section{Thirteenth-cycle V94: noncommutative Cartan calculus and
exact Bianchi-defect identities}
\label{sec:v13v94}

\subsection{Why the connection layer must precede Einstein contraction}

Section~\ref{sec:v13v93} established an ordered differential calculus for
the inverse of a noncentral coframe.  The next operation is not yet the
Einstein contraction.  One must first identify a differential graded
algebra in which connection, curvature, and torsion are meaningful and
verify that the two covariant derivatives used below are not confused:
one acts on column-valued forms and the other on endomorphism-valued
forms.  This section supplies that layer.

The distinction is essential in a noncommutative coefficient algebra.
Moving a one-form through another factor, suppressing the right
connection term, or using cyclicity before a trace has been constructed
would invalidate the usual textbook cancellation.  The following
identities retain the order of every factor.

\begin{definition}[Differential graded coefficient algebra]
Let
\[
 \Omega^\bullet=\bigoplus_{q\geq0}\Omega^q
\]
be an associative unital graded algebra over \(\mathbb C\), with
\(\Omega^0=\mathfrak A_\infty\), and let
\(d:\Omega^q\to\Omega^{q+1}\) be linear.  In the exact regime we assume
\[
 d^2=0,\qquad
 d(\alpha\beta)=(d\alpha)\beta+(-1)^{|\alpha|}\alpha\,d\beta
                                                               \tag{94.1}
\]
for homogeneous \(\alpha\).  Matrix and column spaces over
\(\Omega^\bullet\) inherit the graded product and entrywise
differential.  No commutativity of \(\mathfrak A_\infty\) is assumed.
\end{definition}

\begin{definition}[Connection, curvature, and the two covariant
derivatives]
Fix \(m\geq1\) and a connection form
\(\omega\in M_m(\Omega^1)\).  Its curvature is
\[
 {\cal R}:=d\omega+\omega^2\in M_m(\Omega^2).                  \tag{94.2}
\]
For \(X\in M_{m,1}(\Omega^q)\), define the left-module covariant
derivative
\[
 \nabla^\mathrm{L}_\omega X:=dX+\omega X.                     \tag{94.3}
\]
For \(A\in M_m(\Omega^q)\), define the endomorphism covariant derivative
\[
 \nabla^\mathrm{ad}_\omega A
 :=dA+\omega A-(-1)^qA\omega.                                 \tag{94.4}
\]
If \(E\in M_{m,1}(\Omega^1)\) is a coframe column, its torsion is
\[
 {\cal T}:=\nabla^\mathrm{L}_\omega E=dE+\omega E
       \in M_{m,1}(\Omega^2).                                 \tag{94.5}
\]
\end{definition}

The notation \({\cal R}\) is chosen deliberately: the symbol \(F\) in
Section~\ref{sec:v13v93} denotes the inverse coframe and must not be
silently reassigned inside a calculation involving both objects.

\begin{lemma}[Curvature is the square of the left connection]
For every homogeneous column \(X\in M_{m,1}(\Omega^q)\),
\[
 (\nabla^\mathrm{L}_\omega)^2X={\cal R}X.                     \tag{94.6}
\]
\end{lemma}

\begin{proof}
The graded Leibniz rule and \(|\omega|=1\) give
\[
\begin{aligned}
 (\nabla^\mathrm{L}_\omega)^2X
 &=d(dX+\omega X)+\omega(dX+\omega X)\\
 &=d^2X+(d\omega)X-\omega dX+\omega dX+\omega^2X\\
 &=(d\omega+\omega^2)X.
\end{aligned}
\]
The two middle terms cancel without any permutation of factors.
\end{proof}

\begin{theorem}[Ordered Cartan--Bianchi identities]
In the exact differential graded algebra of
Definition~94.1,
\[
 \nabla^\mathrm{ad}_\omega{\cal R}=0,                         \tag{94.7}
\]
and
\[
 \nabla^\mathrm{L}_\omega{\cal T}={\cal R}E.                  \tag{94.8}
\]
Consequently, the torsion-free equation \({\cal T}=0\) implies the first
Bianchi relation
\[
 {\cal R}E=0.                                                  \tag{94.9}
\]
\end{theorem}

\begin{proof}
Since \(|\omega|=1\), (94.1) yields
\[
 d{\cal R}=d^2\omega+d(\omega^2)
           =(d\omega)\omega-\omega\,d\omega.
\]
Since \(|{\cal R}|=2\), formula (94.4) gives
\[
\begin{aligned}
 \nabla^\mathrm{ad}_\omega{\cal R}
 &= (d\omega)\omega-\omega d\omega
    +\omega(d\omega+\omega^2)
    -(d\omega+\omega^2)\omega\\
 &=0.
\end{aligned}
\]
Associativity cancels the two cubic terms, and the remaining terms
cancel in their displayed order.  Equation (94.8) follows either from
Lemma~94.3 with \(X=E\), or directly:
\[
 \nabla^\mathrm{L}_\omega{\cal T}
 =(\nabla^\mathrm{L}_\omega)^2E={\cal R}E.
\]
If \({\cal T}=0\), its covariant derivative vanishes and (94.9) follows.
\end{proof}

\begin{remark}[What has and has not been proved]
Equation (94.7) is the uncontracted second Bianchi identity, while
(94.9) is the torsion-free first Bianchi identity.  Neither statement
alone defines a Ricci tensor, scalar curvature, Einstein tensor, or
covariant divergence.  Those operations require an ordered contraction
built from the coframe and its inverse, together with a proof that the
contraction intertwines the relevant covariant derivatives.  This
firewall prevents a classical index argument from being imported
without checking its noncommutative hypotheses.
\end{remark}

\subsection{Gauge covariance}

\begin{proposition}[Change of frame]
Let \(u\in GL_m(\Omega^0)\), and set
\[
 E^u:=uE,\qquad
 \omega^u:=u\omega u^{-1}-(du)u^{-1}.                         \tag{94.10}
\]
Then for every column-valued form \(X\),
\[
 \nabla^\mathrm{L}_{\omega^u}(uX)
     =u\nabla^\mathrm{L}_\omega X.                            \tag{94.11}
\]
In particular,
\[
 {\cal T}^u=u{\cal T},\qquad
 {\cal R}^u=u{\cal R}u^{-1}.                                  \tag{94.12}
\]
Thus torsion-freeness and both Bianchi identities are frame covariant.
\end{proposition}

\begin{proof}
Because \(u\) has degree zero,
\[
\begin{aligned}
 \nabla^\mathrm{L}_{\omega^u}(uX)
 &=d(uX)+\bigl(u\omega u^{-1}-(du)u^{-1}\bigr)uX\\
 &=(du)X+u\,dX+u\omega X-(du)X\\
 &=u(dX+\omega X).
\end{aligned}
\]
Putting \(X=E\) proves the torsion formula.  Apply (94.11) twice:
\[
 {\cal R}^u uX
 =(\nabla^\mathrm{L}_{\omega^u})^2(uX)
 =u(\nabla^\mathrm{L}_\omega)^2X
 =u{\cal R}X.
\]
Right multiplication by \(u^{-1}\), entrywise on arbitrary columns,
gives \({\cal R}^u=u{\cal R}u^{-1}\).  Equations (94.7)--(94.9) are
therefore transported by \(u\).
\end{proof}

\begin{remark}
If the differential graded algebra has a compatible involution and the
geometric structure restricts changes of frame to a unitary or
pseudo-unitary subgroup, that restriction must be imposed in addition
to Proposition~94.6.  Algebraic invertibility alone proves covariance,
but it does not prove metric compatibility.
\end{remark}

\subsection{Exact identities for imperfect finite-stage calculi}

The cofinal programme does not begin with an exact continuum
differential.  It therefore needs an identity that exposes, rather than
hides, every finite-stage defect.

\begin{definition}[Nilpotence and Leibniz defects]
Let \(d_n\) be a degree-one linear map on an associative graded algebra.
For homogeneous \(\alpha,\beta\), define
\[
 {\cal L}_n(\alpha,\beta)
 :=d_n(\alpha\beta)-(d_n\alpha)\beta
       -(-1)^{|\alpha|}\alpha\,d_n\beta.                      \tag{94.13}
\]
For \(\omega_n\in M_m(\Omega_n^1)\) and
\(E_n\in M_{m,1}(\Omega_n^1)\), put
\[
 {\cal R}_n=d_n\omega_n+\omega_n^2,\qquad
 {\cal T}_n=d_nE_n+\omega_nE_n.                               \tag{94.14}
\]
The definitions of \(\nabla^{\mathrm L}_n\) and
\(\nabla^{\mathrm{ad}}_n\) are (94.3)--(94.4) with
\((d,\omega)\) replaced by \((d_n,\omega_n)\).
\end{definition}

\begin{theorem}[Bianchi-defect identities]
Without assuming \(d_n^2=0\) or the graded Leibniz rule, one has the
exact ordered formulas
\[
 \nabla^{\mathrm{ad}}_n{\cal R}_n
 =d_n^2\omega_n+{\cal L}_n(\omega_n,\omega_n),                \tag{94.15}
\]
and
\[
 \nabla^{\mathrm L}_n{\cal T}_n-{\cal R}_nE_n
 =d_n^2E_n+{\cal L}_n(\omega_n,E_n).                          \tag{94.16}
\]
No commutator, centrality, or cyclic-trace error is omitted.
\end{theorem}

\begin{proof}
Expanding the first left-hand side and using only associativity gives
\[
\begin{aligned}
 d_n{\cal R}_n+\omega_n{\cal R}_n-{\cal R}_n\omega_n
 &=d_n^2\omega_n+d_n(\omega_n^2)
   +\omega_n d_n\omega_n-d_n\omega_n\,\omega_n\\
 &\quad+\omega_n^3-\omega_n^3.
\end{aligned}
\]
By the definition (94.13),
\[
 d_n(\omega_n^2)
 =(d_n\omega_n)\omega_n-\omega_n d_n\omega_n
   +{\cal L}_n(\omega_n,\omega_n),
\]
which proves (94.15).  Similarly,
\[
\begin{aligned}
 \nabla^{\mathrm L}_n{\cal T}_n
 &=d_n^2E_n+d_n(\omega_nE_n)
      +\omega_n d_nE_n+\omega_n^2E_n\\
 &=d_n^2E_n+(d_n\omega_n)E_n-\omega_n d_nE_n
      +{\cal L}_n(\omega_n,E_n)\\
 &\quad+\omega_n d_nE_n+\omega_n^2E_n\\
 &= {\cal R}_nE_n+d_n^2E_n
      +{\cal L}_n(\omega_n,E_n),
\end{aligned}
\]
which is (94.16).
\end{proof}

\begin{corollary}[Quantitative cofinal Bianchi criterion]
Suppose the finite-stage form spaces carry submultiplicative norms and
all terms in (94.15)--(94.16) are defined.  Then
\[
 \|\nabla^{\mathrm{ad}}_n{\cal R}_n\|
 \leq \|d_n^2\omega_n\|
      +\|{\cal L}_n(\omega_n,\omega_n)\|,                     \tag{94.17}
\]
and
\[
 \|\nabla^{\mathrm L}_n{\cal T}_n-{\cal R}_nE_n\|
 \leq \|d_n^2E_n\|
      +\|{\cal L}_n(\omega_n,E_n)\|.                          \tag{94.18}
\]
Hence vanishing nilpotence and Leibniz defects on the actual connection
and coframe are sufficient for both Bianchi residuals to vanish.  A
global operator-norm estimate for \({\cal L}_n\) is sufficient but is
not required.
\end{corollary}

\begin{proof}
Take norms in the exact identities (94.15)--(94.16) and use the triangle
inequality.  The last assertion follows because only the two displayed
evaluations of \({\cal L}_n\) enter those identities.
\end{proof}

\subsection{Acquisition ledger and next obstruction}

\[
\begin{array}{p{0.27\linewidth}|p{0.25\linewidth}|p{0.38\linewidth}}
\textbf{object} & \textbf{status after V94} & \textbf{proof obligation}\\ \hline
\text{left connection square}
 & \text{acquired}
 & (\nabla^{\mathrm L}_\omega)^2X={\cal R}X
   with ordered factors\\
\text{uncontracted Bianchi identity}
 & \text{acquired}
 & \nabla^{\mathrm{ad}}_\omega{\cal R}=0\\
\text{torsion identity}
 & \text{acquired}
 & \nabla^{\mathrm L}_\omega{\cal T}={\cal R}E\\
\text{frame covariance}
 & \text{acquired}
 & {\cal T}^u=u{\cal T},\
   {\cal R}^u=u{\cal R}u^{-1}\\
\text{cofinal Bianchi control}
 & \text{reduced to explicit defects}
 & d_n^2\omega_n,\ d_n^2E_n,\
   {\cal L}_n(\omega_n,\omega_n),\
   {\cal L}_n(\omega_n,E_n)\to0\\
\text{Einstein contraction}
 & \text{open}
 & construct an ordered contraction and prove covariant
   intertwining
\end{array}                                                     \tag{94.19}
\]

The advance is therefore precise but limited.  Curvature conservation
has been proved before contraction, and its finite-stage failure has
been localized exactly.  The next mathematical bottleneck is the map
from \({\cal R}\) to a Ricci/scalar/Einstein object.  That map must
respect the inverse-coframe order from Section~\ref{sec:v13v93}; a
cyclic trace or a central metric cannot be inserted without its own
hypotheses.

\subsection{Parent-integrity record}

The installed parent was
\[
\begin{array}{c|c}
 \text{file}&
 \texttt{Renewal\_SM\_Einstein\_Continuum\_Research\_Notes\_V93.tex}\\
 \text{lines}&27841\\
 \text{bytes}&979930\\
 \text{SHA--256}&
 \texttt{082525CAC2C0D497BAB845CFCCA884D620A8A61ED5378B8B8635D6A16086CBAD}.
\end{array}                                                     \tag{94.20}
\]

% =============================================================================
% END APPENDED THIRTEENTH-CYCLE V94 MATERIAL
% =============================================================================
% =============================================================================
% BEGIN APPENDED THIRTEENTH-CYCLE V95 MATERIAL
% =============================================================================

\section{Thirteenth-cycle V95: compulsory companion audit and
contraction-direction decision}
\label{sec:v13v95}

\subsection{Files inspected at the five-version boundary}

The required audit was performed on 10 September 2026 after the
installation of V94.  The current active companion files were
\[
\begin{array}{p{0.39\linewidth}|r|r|p{0.27\linewidth}}
\textbf{file}&\textbf{lines}&\textbf{bytes}&\textbf{SHA--256}\\ \hline
\texttt{Renewal\_Yang\_Mills\_Mass\_Gap\_Research\_Notes\_V96.tex}
 &31165&1063317&
 \texttt{731EA8552A0E36269EC0407C9B776C673444BB5C527D3B498C78FC7AF475738A}\\
\texttt{Renewal\_NS\_Stability\_Research\_Notes\_V26.tex}
 &5319&278283&
 \texttt{82C887F8C2C69E4F28FAD7C3DC8DE5B9099CB5A4BF431EBA1FFF3E91935978AD}.
\end{array}                                                     \tag{95.1}
\]
The auxiliary shared-folder note
\[
 \texttt{Renewal\_YM\_Parallel\_Research\_Notes\_V5.tex}
\]
was also checked.  It has \(1351\) lines, \(54174\) bytes, and SHA--256
\[
 \texttt{306F1BB84CFA8A8D47EA569EC17E9545F9867C8D32B548EC9D9C444B4F3C0512}.
                                                                  \tag{95.2}
\]
Archived files were not mistaken for active companions and were not
modified.

\subsection{Mathematical findings of the audit}

The principal YM note has advanced since the preceding SM audit.  Its
V95 block obtains strong exact normalized-covariance bounds on a
quarter-turn grid, but explicitly refuses to promote a discrete card
to a global bound.  Its V96 block then populates a genuine Voronoi
cover and proves a no-go for the chosen absolute-radius,
fixed-anchor architecture: the transported residual is already too
large before response contraction.  It therefore goes upstream to the
relative Hermitian pencils
\[
 L_{j,\pm}(k;t)=tH(k)\pm H_j(k),                              \tag{95.3}
\]
whose positivity retains cancellation between the base operator and
its derivative.

Two parts of that result transfer to the present programme.

First, an ordered Einstein contraction must not be certified solely by
expanding all factors into separate absolute norms when a relative
order estimate is available.  The inverse-coframe bounds of
Section~\ref{sec:v13v93} provide a safe ambient estimate, but a sharper
continuum gate should also test the contraction on the actual
curvature/coframe input class.

Second, the YM V95 audit imports the all-order interacting-state
requirement obtained earlier in this SM programme.  This confirms that
no curvature identity at fixed correlation order can replace the
cofinal positive compatible observable ledger already required here.

The NS note remains at V26.  It computes rank-one Oseen-kernel norms
and records that the resulting improvement is only about two percent
in the relevant integral.  The transferable principle is precise:
restriction to the input class can sharpen a bound, but the measured
gain must be retained at its actual size.  No Oseen-kernel constant or
Navier--Stokes regularity assertion enters the present programme.

The auxiliary YM note proves an energy--entropy failure for a
volume-uniform tail of an extensive absolute topological charge.  Its
recommended replacement is a quasi-local algebra, phase
decomposition, or fixed-window/density control.  This agrees with the
local-net and all-order-ledger route already adopted in the SM notes;
it supplies no new Einstein contraction formula.

\begin{proposition}[Audit decision]
\label{prop:v13v95-decision}
The next contraction stage shall have three distinct outputs:
\begin{enumerate}
\item construct the universal graded cyclic trace and state exactly
      which Bianchi consequence survives modulo commutators;
\item define the covariant-intertwining defect of any proposed
      coframe contraction and derive an exact divergence identity;
\item require estimates both on the realized curvature/coframe input
      class and, where continuum landing uses it, on the complete
      contraction map.
\end{enumerate}
No cyclic trace will be identified with a Ricci or Einstein
contraction without an additional theorem.
\end{proposition}

\begin{proof}
The first item is forced by noncommutativity: an ordinary matrix trace
with values in \(\mathfrak A_\infty\) is not cyclic.  The second is
forced by Section~\ref{sec:v13v94}, which proves only the uncontracted
Bianchi identity.  The third follows from the two companion lessons:
YM V96 locates a decisive loss in absolute transport, while NS V26
shows that structured-input sharpening must be quantified rather than
presumed.  These are compatible requirements, so the stated direction
does not import any unproved companion conclusion.
\end{proof}

\subsection{Nonduplication and claim boundary}

V95 adds no new curvature algebra; that was completed in
Section~\ref{sec:v13v94}.  It also does not repeat the earlier
all-order moment obstruction, local-net landing, OS reconstruction,
or noncentral inverse-coframe calculus.  Its purpose is to record the
mandatory cross-programme read and to select the next unresolved
interface.

The following implications remain invalid without further hypotheses:
\[
 \nabla^{\mathrm{ad}}_\omega{\cal R}=0
 \ \centernot\Longrightarrow\
 \operatorname{Div}\operatorname{Ric}({\cal R})=0,
 \qquad
 \operatorname{tr}_{m}([A,B])
 \ \centernot=0
 \quad\hbox{in a noncommutative coefficient algebra}.        \tag{95.4}
\]
V96 will replace the second failed equality by a universal cyclic
quotient and will isolate the exact defect in the first implication.

\subsection{Parent-integrity record}

The installed parent was
\[
\begin{array}{c|c}
 \text{file}&
 \texttt{Renewal\_SM\_Einstein\_Continuum\_Research\_Notes\_V94.tex}\\
 \text{lines}&28190\\
 \text{bytes}&991626\\
 \text{SHA--256}&
 \texttt{272A20F0A5E77887E60DF07FC6CF4BF6A07335733746188B28C1E2F6AF556075}.
\end{array}                                                     \tag{95.5}
\]

% =============================================================================
% END APPENDED THIRTEENTH-CYCLE V95 MATERIAL
% =============================================================================
% =============================================================================
% BEGIN APPENDED THIRTEENTH-CYCLE V96 MATERIAL
% =============================================================================

\section{Thirteenth-cycle V96: universal cyclic trace and the exact
Einstein-contraction gate}
\label{sec:v13v96}

\subsection{The universal trace available without commutativity}

The V95 audit separated cyclic conservation from Einstein
contraction.  We first construct the former with no hidden cyclicity
assumption.

For homogeneous \(\alpha\in\Omega^p\) and
\(\beta\in\Omega^q\), write
\[
 [\alpha,\beta]_{\mathrm{gr}}
 :=\alpha\beta-(-1)^{pq}\beta\alpha.                          \tag{96.1}
\]
Let \(\mathfrak C^k\subseteq\Omega^k\) be the linear span of all
graded commutators of total degree \(k\), and define
\[
 \Omega^k_{\mathrm{cyc}}:=\Omega^k/\mathfrak C^k,\qquad
 q_k:\Omega^k\longrightarrow\Omega^k_{\mathrm{cyc}}.         \tag{96.2}
\]
If the form spaces are normed and a Hausdorff quotient is required,
replace \(\mathfrak C^k\) by its norm closure.  The algebraic
identities below are unchanged.

\begin{lemma}[The differential descends to the cyclic quotient]
\label{lem:v13v96-descend}
For the exact differential of Section~\ref{sec:v13v94},
\[
 d\mathfrak C^k\subseteq\mathfrak C^{k+1}.                   \tag{96.3}
\]
Consequently there is a well-defined differential
\(\bar d\,q_k(\alpha)=q_{k+1}(d\alpha)\) on
\(\Omega^\bullet_{\mathrm{cyc}}\), and \(\bar d^2=0\).
\end{lemma}

\begin{proof}
For \(|\alpha|=p\) and \(|\beta|=q\), the graded Leibniz rule gives
\[
 d[\alpha,\beta]_{\mathrm{gr}}
 =[d\alpha,\beta]_{\mathrm{gr}}
   +(-1)^p[\alpha,d\beta]_{\mathrm{gr}}.                     \tag{96.4}
\]
Both terms on the right are graded commutators of total degree
\(p+q+1\).  This proves (96.3), hence well-definedness.  Finally,
\(\bar d^2q_k(\alpha)=q_{k+2}(d^2\alpha)=0\).
\end{proof}

\begin{definition}[Universal graded matrix trace]
For \(A=(A_{ij})\in M_m(\Omega^k)\), set
\[
 \operatorname{Tr}_{\mathrm{cyc}}(A)
 :=q_k\!\left(\sum_{i=1}^m A_{ii}\right)
 \in\Omega^k_{\mathrm{cyc}}.                                 \tag{96.5}
\]
\end{definition}

\begin{proposition}[Graded cyclicity and compatibility with \(d\)]
\label{prop:v13v96-cyclic}
If \(A\in M_m(\Omega^p)\) and \(B\in M_m(\Omega^q)\) are
homogeneous, then
\[
 \operatorname{Tr}_{\mathrm{cyc}}(AB)
 =(-1)^{pq}\operatorname{Tr}_{\mathrm{cyc}}(BA),             \tag{96.6}
\]
and
\[
 \bar d\,\operatorname{Tr}_{\mathrm{cyc}}(A)
 =\operatorname{Tr}_{\mathrm{cyc}}(dA).                      \tag{96.7}
\]
\end{proposition}

\begin{proof}
By (96.2),
\[
 q_{p+q}(A_{ij}B_{ji})
 =(-1)^{pq}q_{p+q}(B_{ji}A_{ij}).
\]
Sum over \(i,j\) and relabel the two indices to obtain (96.6).
Equation (96.7) follows from Lemma~\ref{lem:v13v96-descend} and the
entrywise definition of \(d\).
\end{proof}

\begin{theorem}[Cyclic Chern forms are closed]
\label{thm:v13v96-chern}
Let \({\cal R}=d\omega+\omega^2\) be the curvature from
Section~\ref{sec:v13v94}.  For every integer \(r\geq1\),
\[
 \boxed{\quad
 \bar d\,\operatorname{Tr}_{\mathrm{cyc}}({\cal R}^r)=0.
 \quad}                                                       \tag{96.8}
\]
In particular, the uncontracted Bianchi identity has a canonical
scalar-valued consequence in the cyclic quotient.
\end{theorem}

\begin{proof}
The ordered Bianchi identity says
\[
 d{\cal R}={\cal R}\omega-\omega{\cal R}.
\]
Because \({\cal R}\) has even degree, the ordinary Leibniz expansion
and telescoping give
\[
 d({\cal R}^r)
 =\sum_{j=0}^{r-1}{\cal R}^j(d{\cal R}){\cal R}^{r-1-j}
 ={\cal R}^r\omega-\omega{\cal R}^r.                         \tag{96.9}
\]
Apply (96.7), then (96.6) with degrees \(2r\) and \(1\).
The two cyclic classes on the right of (96.9) coincide, so their
difference is zero.
\end{proof}

\begin{proposition}[Why the quotient is necessary]
\label{prop:v13v96-noordinarytrace}
An ordinary matrix trace with values in a noncommutative coefficient
algebra need not be cyclic.  Already for matrix size \(m=1\) and
\(\Omega^0=M_2(\mathbb C)\), let
\[
 a=\begin{pmatrix}0&1\\0&0\end{pmatrix},\qquad
 b=\begin{pmatrix}0&0\\1&0\end{pmatrix}.
\]
Then
\[
 \operatorname{tr}_1(ab)-\operatorname{tr}_1(ba)
 =ab-ba
 =\begin{pmatrix}1&0\\0&-1\end{pmatrix}\neq0.                \tag{96.10}
\]
The class of this difference vanishes under \(q_0\), exactly as
required by (96.6).
\end{proposition}

\begin{proof}
For \(m=1\), the matrix trace is the identity on the coefficient
algebra.  Direct multiplication gives (96.10), while every
commutator belongs to \(\mathfrak C^0\).
\end{proof}

\begin{remark}[Cyclic closure is not Einstein conservation]
The map \(q_k\) deliberately discards commutators and
\(\operatorname{Tr}_{\mathrm{cyc}}\) discards matrix-component
information.  Thus (96.8) is a Chern--Weil type statement.  It is not
a Ricci tensor, an Einstein tensor, or an operator-valued local
conservation law.
\end{remark}

\subsection{The trace-reversal ordering obstruction}

Even after an ordered Ricci contraction and scalar contraction have
been proposed, the classical trace reversal contains a new choice.
Let \(P\) denote a Ricci candidate, \(g\) a metric coefficient, and
\(s\) a scalar-curvature candidate in a unital algebra.  The two
literal orderings are
\[
 {\cal G}_{\mathrm L}:=P-\frac12\,gs,\qquad
 {\cal G}_{\mathrm R}:=P-\frac12\,sg.                         \tag{96.11}
\]

\begin{proposition}[Classical trace reversal is not order independent]
\label{prop:v13v96-ordering}
The candidates in (96.11) obey
\[
 {\cal G}_{\mathrm L}-{\cal G}_{\mathrm R}
 =\frac12[s,g].                                               \tag{96.12}
\]
Hence they coincide if and only if \([s,g]=0\).  The ambiguity is
genuine: for
\[
 g=\begin{pmatrix}1&0\\0&2\end{pmatrix},\qquad
 s=\begin{pmatrix}0&1\\0&0\end{pmatrix}
\]
in \(M_2(\mathbb C)\), one has
\[
 {\cal G}_{\mathrm L}-{\cal G}_{\mathrm R}
 =\frac12\begin{pmatrix}0&1\\0&0\end{pmatrix}\neq0.           \tag{96.13}
\]
\end{proposition}

\begin{proof}
Subtract the two formulas in (96.11).  In the displayed example,
\(gs=s\) and \(sg=2s\), which proves (96.13).
\end{proof}

If \(g=g^*\), \(s=s^*\), and \(P=P^*\), neither one-sided expression
in (96.11) need be self-adjoint.  The Jordan-ordered expression
\[
 {\cal G}_{\mathrm J}
 :=P-\frac14(gs+sg)                                          \tag{96.14}
\]
is self-adjoint, but self-adjointness alone does not prove the correct
classical limit, frame covariance, or covariant conservation.  Those
remain separate gates.

\subsection{The exact contraction-intertwining identity}

We now isolate the missing hypothesis without choosing a false
contraction.  Regrade the target tensor complex, if necessary, so that
the proposed contraction is degree preserving.

\begin{definition}[Contraction and its intertwining defect]
Let \(({\cal K}^\bullet,{\mathfrak D})\) be a graded vector space with
a degree-one linear operator \({\mathfrak D}\), interpreted as the
candidate covariant divergence after regrading.  Let
\[
 {\cal C}:M_m(\Omega^\bullet)\longrightarrow{\cal K}^\bullet
\]
be a degree-preserving linear contraction.  Its covariant
intertwining defect on a homogeneous \(A\) is
\[
 {\cal J}_{\cal C}(A)
 :={\mathfrak D}\,{\cal C}(A)
   -{\cal C}\bigl(\nabla^{\mathrm{ad}}_\omega A\bigr).        \tag{96.15}
\]
The contracted curvature candidate is
\[
 {\cal G}:={\cal C}({\cal R}).                                \tag{96.16}
\]
\end{definition}

\begin{theorem}[Exact Einstein-contraction gate]
\label{thm:v13v96-gate}
For an exact differential graded connection,
\[
 \boxed{\quad
 {\mathfrak D}{\cal G}={\cal J}_{\cal C}({\cal R}).
 \quad}                                                       \tag{96.17}
\]
For a finite-stage calculus with Bianchi residual
\[
 {\cal B}_n:=
 \nabla^{\mathrm{ad}}_n{\cal R}_n
 =d_n^2\omega_n+{\cal L}_n(\omega_n,\omega_n),
\]
one has the exact identity
\[
 {\mathfrak D}_n{\cal C}_n({\cal R}_n)
 ={\cal J}_{{\cal C}_n}({\cal R}_n)+{\cal C}_n({\cal B}_n).
                                                                  \tag{96.18}
\]
Thus uncontracted Bianchi conservation implies contracted
conservation precisely when the contraction intertwines the two
covariant differentials on the realized curvature.
\end{theorem}

\begin{proof}
Insert \(A={\cal R}\) into (96.15) and use
\(\nabla^{\mathrm{ad}}_\omega{\cal R}=0\).  The same rearrangement at
stage \(n\), without setting the Bianchi residual to zero, gives
(96.18).  No estimate or commutation is used.
\end{proof}

\begin{corollary}[Quantitative cofinal conservation criterion]
\label{cor:v13v96-cofinal}
Suppose the finite-stage target spaces are normed.  Then
\[
 \|{\mathfrak D}_n{\cal C}_n({\cal R}_n)\|
 \leq
 \|{\cal J}_{{\cal C}_n}({\cal R}_n)\|
 +\|{\cal C}_n({\cal B}_n)\|.                                \tag{96.19}
\]
Consequently,
\[
 \|{\cal J}_{{\cal C}_n}({\cal R}_n)\|\longrightarrow0,
 \qquad
 \|{\cal C}_n({\cal B}_n)\|\longrightarrow0                 \tag{96.20}
\]
are sufficient for asymptotic contracted conservation.  A uniform
complete bound
\(\|{\cal C}_n\|\leq M\) together with
\(\|{\cal B}_n\|\to0\) is sufficient for the second limit, but the
strictly weaker realized-input condition in (96.20) is enough.
\end{corollary}

\begin{proof}
Take norms in (96.18).  Under the optional complete bound,
\(\|{\cal C}_n({\cal B}_n)\|\leq M\|{\cal B}_n\|\).
\end{proof}

This corollary implements both lessons from the V95 audit.  It retains
a complete map bound when the landing proof needs one, but it does not
replace the actual Bianchi input by an unnecessarily large ambient
ball.

\subsection{Landing of finite ordered coframe contractions}

The algebraic part of an ordered contraction can already be closed
under the V93 inverse-coframe landing theorem.

\begin{lemma}[Ordered-word telescoping estimate]
\label{lem:v13v96-word}
Let \({\cal A}\) be a Banach algebra and, for
\(1\leq j\leq\ell\), let \(X_{j,n},X_j\in{\cal A}\).  Then
\[
\begin{split}
 \left\|\prod_{j=1}^{\ell}X_{j,n}
       -\prod_{j=1}^{\ell}X_j\right\|
 \leq\sum_{r=1}^{\ell}
 &\left(\prod_{j<r}\|X_{j,n}\|\right)
 \|X_{r,n}-X_r\|\\
 &\times\left(\prod_{j>r}\|X_j\|\right).                     \tag{96.21}
\end{split}
\]
In particular, componentwise convergence plus uniform boundedness
implies convergence of the ordered word, without commuting any
factors.
\end{lemma}

\begin{proof}
Use the exact telescoping identity
\[
 \prod_{j=1}^{\ell}X_{j,n}-\prod_{j=1}^{\ell}X_j
 =
 \sum_{r=1}^{\ell}
 \left(\prod_{j<r}X_{j,n}\right)
 (X_{r,n}-X_r)
 \left(\prod_{j>r}X_j\right),
\]
and apply submultiplicativity and the triangle inequality.
\end{proof}

\begin{theorem}[Norm landing of finite ordered contraction words]
\label{thm:v13v96-wordlanding}
Assume
\[
 E_n\to E,\qquad F_n=E_n^{-1}\to F=E^{-1},\qquad
 {\cal R}_n\to{\cal R}
\]
in the relevant matrix Banach algebra, with the adjoints converging as
well.  Let \({\cal P}\) be any fixed finite linear combination of
ordered words in
\[
 E_n,\ E_n^*,\ F_n,\ F_n^*,\ {\cal R}_n
\]
and fixed bounded coefficient tensors.  Then
\[
 {\cal P}(E_n,E_n^*,F_n,F_n^*,{\cal R}_n)
 \longrightarrow
 {\cal P}(E,E^*,F,F^*,{\cal R})                              \tag{96.22}
\]
in norm.  The same conclusion holds for words in the first and second
jets whenever the jet convergences proved in
Section~\ref{sec:v13v93} and the corresponding curvature-jet
convergences hold.
\end{theorem}

\begin{proof}
Apply Lemma~\ref{lem:v13v96-word} to each word.  Convergent factor
sequences are uniformly bounded, and fixed bounded coefficient maps
preserve convergence.  Finite summation gives (96.22).  The jet
statement is the same argument with each differentiated factor treated
as one of the convergent inputs; the ordered inverse-jet formulas of
Section~\ref{sec:v13v93} ensure that no unlisted inverse derivative is
needed.
\end{proof}

\begin{remark}
Theorem~\ref{thm:v13v96-wordlanding} proves analytic landing for a
chosen finite ordered formula.  It does not decide which formula is
geometrically correct.  Frame covariance, adjoint compatibility,
classical reduction, and the vanishing of
\({\cal J}_{\cal C}({\cal R})\) must still be proved for that formula.
\end{remark}

\subsection{Acquisition ledger and next target}

\[
\begin{array}{p{0.27\linewidth}|p{0.24\linewidth}|p{0.39\linewidth}}
\textbf{object}&\textbf{status after V96}&\textbf{content}\\ \hline
\text{universal graded trace}
 &\text{acquired}
 &\operatorname{Tr}_{\mathrm{cyc}}:
   M_m(\Omega^k)\to\Omega^k_{\mathrm{cyc}}\\
\text{cyclic curvature conservation}
 &\text{acquired}
 &\bar d\operatorname{Tr}_{\mathrm{cyc}}({\cal R}^r)=0\\
\text{ordinary trace cyclicity}
 &\text{refuted in general}
 &explicit \(M_2(\mathbb C)\) commutator\\
\text{trace-reversal uniqueness}
 &\text{refuted in general}
 &{\cal G}_{\mathrm L}-{\cal G}_{\mathrm R}
   =\frac12[s,g]\\
\text{contracted conservation}
 &\text{reduced to exact defects}
 &{\cal J}_{{\cal C}_n}({\cal R}_n)
   +{\cal C}_n({\cal B}_n)\\
\text{finite ordered-word landing}
 &\text{acquired}
 &noncommutative telescoping in norm\\
\text{physical Einstein contraction}
 &\text{open}
 &select and verify covariance, reality, classical limit,
   and covariant intertwining
\end{array}                                                     \tag{96.23}
\]

The next construction should start from adjoint compatibility rather
than choose a one-sided trace reversal.  A viable candidate must use
the coframe and inverse coframe in a declared order, reduce to the
classical Einstein tensor when the coefficient algebra is
commutative, transform covariantly under the frame change (94.10), and
make the defect (96.15) computable.  If these requirements conflict,
the programme must return upstream to the differential calculus or
metric-compatibility equation rather than assume the conflict away.

\subsection{Parent-integrity record}

The installed parent was
\[
\begin{array}{c|c}
 \text{file}&
 \texttt{Renewal\_SM\_Einstein\_Continuum\_Research\_Notes\_V95.tex}\\
 \text{lines}&28335\\
 \text{bytes}&997945\\
 \text{SHA--256}&
 \texttt{FE4ED461B20D9223901D9D7245B15D7BCB6C42A9044EE70AC14098CA368F04D2}.
\end{array}                                                     \tag{96.24}
\]

% =============================================================================
% END APPENDED THIRTEENTH-CYCLE V96 MATERIAL
% =============================================================================
% =============================================================================
% BEGIN APPENDED THIRTEENTH-CYCLE V97 MATERIAL
% =============================================================================

\section{Thirteenth-cycle V97: the symmetric trace reversal and its
factorized conservation defect}
\label{sec:v13v97}

\subsection{Selection inside the two-word ordering class}

Section~\ref{sec:v13v96} showed that the left and right versions of
trace reversal differ by a commutator.  We now impose the minimum
reality and reversal requirements on the part of the formula that
multiplies a metric coefficient \(g\) by a scalar-curvature
coefficient \(s\).

For real constants \(a,b\), consider
\[
 {\cal M}_{a,b}(g,s):=a\,gs+b\,sg.                            \tag{97.1}
\]
The classical normalization is
\[
 {\cal M}_{a,b}(g,s)=gs
 \quad\hbox{whenever }gs=sg,                                 \tag{97.2}
\]
and reversal symmetry is
\[
 {\cal M}_{a,b}(g,s)={\cal M}_{a,b}(s,g).                    \tag{97.3}
\]

\begin{proposition}[Unique symmetric two-word continuation]
\label{prop:v13v97-unique}
Within the real two-word class (97.1), conditions (97.2)--(97.3)
have the unique solution
\[
 a=b=\frac12,\qquad
 {\cal M}_{\mathrm J}(g,s):=\frac12(gs+sg).                  \tag{97.4}
\]
If \(g=g^*\) and \(s=s^*\), then
\({\cal M}_{\mathrm J}(g,s)^*={\cal M}_{\mathrm J}(g,s)\).
\end{proposition}

\begin{proof}
On commuting inputs, (97.1) equals \((a+b)gs\), so (97.2) gives
\(a+b=1\).  Reversal symmetry gives
\[
 (a-b)(gs-sg)=0
\]
for every pair \(g,s\).  Since noncommuting pairs exist by
Proposition~\ref{prop:v13v96-ordering}, \(a=b\).  Thus
\(a=b=1/2\).  The adjoint statement follows by reversing the order
of each product:
\[
 (gs+sg)^*=sg+gs.
\]
\end{proof}

\begin{remark}
The uniqueness in Proposition~\ref{prop:v13v97-unique} is deliberately
local to the declared two-word class.  It does not exclude longer
commutator corrections that vanish in the commutative limit.  Such a
correction would require an independent geometric or variational
principle.
\end{remark}

Let \({\cal P}\) be an ordered Ricci contraction, let
\({\cal S}\) be an ordered scalar contraction, and put
\[
 P(A):={\cal P}(A),\qquad s(A):={\cal S}(A).
\]
For a fixed metric coefficient \(g\), define the Jordan-ordered
Einstein contraction
\[
 {\cal C}_{\mathrm J}(A)
 :=P(A)-\frac14\bigl(g\,s(A)+s(A)\,g\bigr),                  \tag{97.5}
\]
and the curvature candidate
\[
 {\cal G}_{\mathrm J}:={\cal C}_{\mathrm J}({\cal R}).       \tag{97.6}
\]
When all coefficients commute, (97.5) reduces exactly to
\(P(A)-\frac12g\,s(A)\).

\subsection{Frame covariance and reality}

\begin{theorem}[Covariance of the symmetric trace reversal]
\label{thm:v13v97-covariance}
Suppose a frame change \(u\) acts by conjugation on the target
coefficients and that
\[
 g^u=ugu^{-1},\qquad
 P^u(uAu^{-1})=uP(A)u^{-1},\qquad
 s^u(uAu^{-1})=us(A)u^{-1}.                                 \tag{97.7}
\]
Then
\[
 {\cal C}_{\mathrm J}^u(uAu^{-1})
 =u{\cal C}_{\mathrm J}(A)u^{-1}.                            \tag{97.8}
\]
If, in addition, \(g=g^*\), \(P(A)=P(A)^*\), and
\(s(A)=s(A)^*\), then
\[
 {\cal C}_{\mathrm J}(A)^*={\cal C}_{\mathrm J}(A).          \tag{97.9}
\]
\end{theorem}

\begin{proof}
Insert (97.7) into (97.5).  Associativity and
\(u^{-1}u=1\) give
\[
 (ugu^{-1})(us(A)u^{-1})=u\,g\,s(A)\,u^{-1},
\]
and the same calculation gives
\((us(A)u^{-1})(ugu^{-1})=u\,s(A)\,g\,u^{-1}\).
This proves (97.8).  Equation (97.9) follows from
Proposition~\ref{prop:v13v97-unique}.
\end{proof}

The hypotheses (97.7) are gates, not conclusions.  In particular,
the scalar output is allowed to be noncentral; centrality is not
inserted to force the two orderings to agree.

\subsection{Factorization of the covariant-intertwining defect}

Let \({\mathfrak D}\) act on the target coefficient algebra as a
derivation on the degree-zero tensors occurring after regrading:
\[
 {\mathfrak D}(XY)=({\mathfrak D}X)Y+X({\mathfrak D}Y).       \tag{97.10}
\]
Define
\[
 Q:={\mathfrak D}g,                                         \tag{97.11}
\]
and, for a homogeneous curvature-complex input \(A\),
\[
\begin{aligned}
 {\cal J}_{P}(A)
 &:={\mathfrak D}P(A)
      -P(\nabla^{\mathrm{ad}}_\omega A),\\
 {\cal J}_{s}(A)
 &:={\mathfrak D}s(A)
      -s(\nabla^{\mathrm{ad}}_\omega A).                     \tag{97.12}
\end{aligned}
\]

\begin{theorem}[Exact defect decomposition]
\label{thm:v13v97-factor}
The intertwining defect of (97.5) is
\[
\boxed{
\begin{aligned}
 {\cal J}_{{\cal C}_{\mathrm J}}(A)
 ={}&{\cal J}_{P}(A)\\
 &-\frac14\Bigl(
   Q\,s(A)+g\,{\cal J}_{s}(A)
   +{\cal J}_{s}(A)\,g+s(A)\,Q
   \Bigr).
\end{aligned}}                                               \tag{97.13}
\]
For an exact curvature,
\[
\boxed{
\begin{aligned}
 {\mathfrak D}{\cal G}_{\mathrm J}
 ={}&{\cal J}_{P}({\cal R})\\
 &-\frac14\Bigl(
   Q\,s({\cal R})+g\,{\cal J}_{s}({\cal R})
   +{\cal J}_{s}({\cal R})\,g+s({\cal R})\,Q
   \Bigr).
\end{aligned}}                                               \tag{97.14}
\]
Thus the sufficient contracted-Bianchi condition is
\[
 Q=0,\qquad
 {\cal J}_{P}({\cal R})
 =\frac14\bigl(
   g\,{\cal J}_{s}({\cal R})
   +{\cal J}_{s}({\cal R})\,g
   \bigr).                                                    \tag{97.15}
\]
\end{theorem}

\begin{proof}
Apply \({\mathfrak D}\) to (97.5), use (97.10), and subtract
\[
 {\cal C}_{\mathrm J}
   (\nabla^{\mathrm{ad}}_\omega A)
 =
 P(\nabla^{\mathrm{ad}}_\omega A)
 -\frac14\Bigl(
  g\,s(\nabla^{\mathrm{ad}}_\omega A)
  +s(\nabla^{\mathrm{ad}}_\omega A)\,g
 \Bigr).
\]
Collecting the two differences in (97.12) gives (97.13).
Equation (97.14) follows from the exact contraction gate
(96.17).  Under (97.15), every term on the right of (97.14)
cancels in its displayed order.
\end{proof}

\begin{corollary}[Quantitative form]
\label{cor:v13v97-bound}
In a submultiplicative target norm,
\[
\begin{aligned}
 \|{\cal J}_{{\cal C}_{\mathrm J}}(A)\|
 \leq{}&
 \|{\cal J}_{P}(A)\|
 +\frac12\|Q\|\,\|s(A)\|
 +\frac12\|g\|\,\|{\cal J}_{s}(A)\|.                         \tag{97.16}
\end{aligned}
\]
At finite stage \(n\), with Bianchi residual \({\cal B}_n\),
\[
\begin{aligned}
 \|{\mathfrak D}_n{\cal G}_{{\mathrm J},n}\|
 \leq{}&
 \|{\cal J}_{P_n}({\cal R}_n)\|
 +\frac12\|Q_n\|\,\|s_n({\cal R}_n)\|\\
 &+\frac12\|g_n\|\,\|{\cal J}_{s_n}({\cal R}_n)\|
 +\|{\cal C}_{{\mathrm J},n}({\cal B}_n)\|.                 \tag{97.17}
\end{aligned}
\]
\end{corollary}

\begin{proof}
Apply the triangle inequality and submultiplicativity to (97.13),
pairing the two \(Q\)-terms and the two
\({\cal J}_s\)-terms.  Add the last term using the finite-stage
identity (96.18).
\end{proof}

\subsection{A cofinal landing statement}

\begin{proposition}[Landing of the symmetric trace reversal]
\label{prop:v13v97-landing}
Suppose, after embedding the finite stages in common Banach spaces,
\[
 g_n\to g,\qquad
 P_n({\cal R}_n)\to P({\cal R}),\qquad
 s_n({\cal R}_n)\to s({\cal R}).                             \tag{97.18}
\]
Then
\[
 {\cal G}_{{\mathrm J},n}
 \longrightarrow
 {\cal G}_{\mathrm J}                                       \tag{97.19}
\]
in norm.  If the four terms on the right of (97.17) tend to zero,
then the landed candidate is asymptotically covariantly conserved
in the target norm.
\end{proposition}

\begin{proof}
The product convergence in (97.19) follows from the two-factor case
of Lemma~\ref{lem:v13v96-word}, applied to
\(g_ns_n({\cal R}_n)\) and \(s_n({\cal R}_n)g_n\).
The conservation statement is immediate from (97.17).
\end{proof}

\subsection{Acquisition ledger}

\[
\begin{array}{p{0.28\linewidth}|p{0.23\linewidth}|p{0.39\linewidth}}
\textbf{gate}&\textbf{status after V97}&\textbf{remaining datum}\\ \hline
\text{trace-reversal order}
 &\text{selected in the symmetric two-word class}
 &{\cal C}_{\mathrm J}=P-\frac14(gs+sg)\\
\text{classical reduction}
 &\text{proved}
 &commuting coefficients give \(P-\frac12gs\)\\
\text{adjoint compatibility}
 &\text{proved conditionally}
 &g,P,s\text{ self-adjoint}\\
\text{frame covariance}
 &\text{proved conditionally}
 &the three transformation laws (97.7)\\
\text{contracted conservation}
 &\text{factorized}
 &Q,\ {\cal J}_{P},\ {\cal J}_{s},\
   {\cal C}_{\mathrm J}({\cal B})\\
\text{metric-compatible connection}
 &\text{open}
 &construct \(\omega\) with \(Q=0\) and controlled torsion\\
\text{Ricci/scalar contractions}
 &\text{open}
 &ordered maps \(P,s\) satisfying (97.7), reality,
   landing, and (97.15)
\end{array}                                                     \tag{97.20}
\]

The factor \(Q={\mathfrak D}g\) is now an upstream bottleneck.
Choosing \(Q=0\) by declaration would be circular because the
connection has not yet been constructed from the noncentral coframe.
The next version therefore returns to the connection space and
classifies the metric-compatible corrections to the coframe-induced
connection.

\subsection{Parent-integrity record}

The installed parent was
\[
\begin{array}{c|c}
 \text{file}&
 \texttt{Renewal\_SM\_Einstein\_Continuum\_Research\_Notes\_V96.tex}\\
 \text{lines}&28762\\
 \text{bytes}&1012428\\
 \text{SHA--256}&
 \texttt{87110AED2C9F7BF826D308075670D498B9101C4F307659DE0EEFA9ACFEE7E11F}.
\end{array}                                                     \tag{97.21}
\]

% =============================================================================
% END APPENDED THIRTEENTH-CYCLE V97 MATERIAL
% =============================================================================
% =============================================================================
% BEGIN APPENDED THIRTEENTH-CYCLE V98 MATERIAL
% =============================================================================

\section{Thirteenth-cycle V98: metric-compatible connection space and
the pure-gauge curvature obstruction}
\label{sec:v13v98}

\subsection{Coefficient coframe and metricity convention}

This section distinguishes the invertible degree-zero coefficient
matrix \(E\) from the column of coframe one-forms \(\theta\).  Let
\[
 E\in GL_m(\mathfrak A_\infty),\qquad F:=E^{-1},
\]
and let \(\eta=\eta^*\in GL_m(\mathbb C)\) be constant.  The
noncentral metric coefficient is
\[
 G:=E^*\eta E,\qquad
 G^{-1}=F\eta^{-1}F^*.                                      \tag{98.1}
\]
Assume the differential graded algebra has an involution satisfying
\[
 (\alpha\beta)^*=(-1)^{|\alpha||\beta|}\beta^*\alpha^*,
 \qquad d(\alpha^*)=(d\alpha)^*.                             \tag{98.2}
\]
For a connection one-form \(\Gamma\in M_m(\Omega^1)\), define its
metricity defect by
\[
 {\cal Q}_\Gamma
 :=dG-\Gamma^*G-G\Gamma.                                    \tag{98.3}
\]
Thus \({\cal Q}_\Gamma=0\) is the declared metric-compatibility
equation.

\begin{definition}[Coframe-induced connection]
Set
\[
 \Gamma_0:=F\,dE.                                            \tag{98.4}
\]
This is the left Maurer--Cartan, or pure-gauge, connection associated
with the invertible coefficient matrix \(E\).
\end{definition}

\begin{theorem}[Metric compatibility and flatness of \(\Gamma_0\)]
\label{thm:v13v98-flat}
The connection (98.4) satisfies
\[
 {\cal Q}_{\Gamma_0}=0,\qquad
 {\cal R}_{\Gamma_0}:=d\Gamma_0+\Gamma_0^2=0.                \tag{98.5}
\]
Hence the coframe-induced connection acquires the metricity gate in
V97 but cannot by itself supply nonzero Einstein curvature.
\end{theorem}

\begin{proof}
Differentiate (98.1), using \(d\eta=0\):
\[
 dG=(dE)^*\eta E+E^*\eta\,dE.                               \tag{98.6}
\]
Since \(F^*E^*=(EF)^*=1\) and \(EF=1\),
\[
 \Gamma_0^*G=(dE)^*F^*E^*\eta E=(dE)^*\eta E,
 \qquad
 G\Gamma_0=E^*\eta EF\,dE=E^*\eta\,dE.
\]
Equations (98.3) and (98.6) give
\({\cal Q}_{\Gamma_0}=0\).

Differentiating \(FE=1\) gives the ordered inverse identity
\[
 dF=-F(dE)F.                                                 \tag{98.7}
\]
Therefore, using \(d^2E=0\),
\[
 d\Gamma_0=d(F\,dE)=(dF)dE+F\,d^2E
 =-F(dE)F(dE)=-\Gamma_0^2.
\]
This proves the curvature statement.
\end{proof}

\begin{remark}[Upstream interpretation]
Theorem~\ref{thm:v13v98-flat} is not a contradiction.  It identifies
\(\Gamma_0\) as a flat metric-compatible, generally torsionful
connection of teleparallel type.  Equating it with the desired
torsion-free gravitational connection would erase curvature by
construction.
\end{remark}

\subsection{All metric-compatible corrections}

Write a general connection relative to \(\Gamma_0\) as
\[
 \Gamma=\Gamma_0+K,\qquad K\in M_m(\Omega^1).                \tag{98.8}
\]

\begin{theorem}[Affine classification of metric-compatible
connections]
\label{thm:v13v98-classification}
The metricity defect of (98.8) is
\[
 {\cal Q}_{\Gamma}=-(K^*G+GK).                               \tag{98.9}
\]
Consequently, \(\Gamma\) is metric compatible if and only if
\[
 K^*G+GK=0.                                                  \tag{98.10}
\]
Equivalently, there is a unique anti-self-adjoint one-form matrix
\(H\), \(H^*=-H\), such that
\[
 K=G^{-1}H,\qquad H=GK.                                     \tag{98.11}
\]
Thus the space of metric-compatible connections is the affine space
\[
 \Gamma_0+G^{-1}
 \{H\in M_m(\Omega^1):H^*=-H\}.                              \tag{98.12}
\]
\end{theorem}

\begin{proof}
Substitute (98.8) into (98.3) and use
\({\cal Q}_{\Gamma_0}=0\) to obtain (98.9).  If
\(H:=GK\), then, because \(G=G^*\),
\[
 H^*=K^*G.
\]
Hence (98.10) is equivalent to \(H^*+H=0\).  Conversely, for
anti-self-adjoint \(H\), define \(K=G^{-1}H\); then \(GK=H\) and
\(K^*G=H^*\), so (98.10) holds.  Invertibility of \(G\) makes the
correspondence unique.
\end{proof}

\begin{corollary}[Curvature carried by the correction]
\label{cor:v13v98-curvature}
Let
\[
 \nabla^{\mathrm{ad}}_{\Gamma_0}K
 :=dK+\Gamma_0K+K\Gamma_0
\]
for the degree-one form \(K\).  Then
\[
 {\cal R}_{\Gamma}
 =\nabla^{\mathrm{ad}}_{\Gamma_0}K+K^2.                     \tag{98.13}
\]
In the metric-compatible class (98.12), every nonzero curvature
contribution is therefore carried by the anti-self-adjoint datum
\(H\) through \(K=G^{-1}H\).
\end{corollary}

\begin{proof}
Expand \(d(\Gamma_0+K)+(\Gamma_0+K)^2\) and use
\({\cal R}_{\Gamma_0}=0\).  Since both \(\Gamma_0\) and \(K\) have
degree one, the two mixed products occur with plus signs, giving
(98.13).
\end{proof}

\subsection{Torsion is the selector, not an automatic consequence}

Let \(\theta\in M_{m,1}(\Omega^1)\) be the coframe-form column used in
Section~\ref{sec:v13v94}.  Define
\[
 {\cal T}_\Gamma:=d\theta+\Gamma\theta,\qquad
 {\cal T}_0:=d\theta+\Gamma_0\theta.                         \tag{98.14}
\]

\begin{theorem}[Exact torsion-selection equation]
\label{thm:v13v98-selector}
For the metric-compatible connection
\(\Gamma=\Gamma_0+G^{-1}H\), one has
\[
 {\cal T}_\Gamma
 ={\cal T}_0+G^{-1}H\theta.                                  \tag{98.15}
\]
Therefore simultaneous metric compatibility and torsion-freeness are
equivalent to the ordered feasibility system
\[
 \boxed{
 H^*=-H,\qquad
 G^{-1}H\theta=-{\cal T}_0.}                                 \tag{98.16}
\]
\end{theorem}

\begin{proof}
Insert (98.8) into (98.14):
\[
 {\cal T}_\Gamma
 =d\theta+\Gamma_0\theta+K\theta
 ={\cal T}_0+K\theta.
\]
Now use the classification \(K=G^{-1}H\).
\end{proof}

\begin{corollary}[Uniqueness criterion]
\label{cor:v13v98-unique}
Define the real-linear Cartan selector
\[
 {\cal L}_{G,\theta}:
 \{H:H^*=-H\}\longrightarrow M_{m,1}(\Omega^2),
 \qquad
 {\cal L}_{G,\theta}(H):=G^{-1}H\theta.                     \tag{98.17}
\]
If (98.16) is solvable, its solution is unique exactly when
\[
 \ker{\cal L}_{G,\theta}=\{0\}.                              \tag{98.18}
\]
It exists exactly when
\(-{\cal T}_0\in\operatorname{Ran}{\cal L}_{G,\theta}\).
\end{corollary}

\begin{proof}
Equation (98.16) is the linear equation
\({\cal L}_{G,\theta}(H)=-{\cal T}_0\) on the real vector space of
anti-self-adjoint one-form matrices.  The two assertions are the
elementary kernel and range criteria for a linear equation.
\end{proof}

This is the first noncircular connection gate in the present cycle.
Metricity is solved explicitly, and torsion-freeness has become a
declared linear selector.  Neither existence nor uniqueness of that
selector is assumed.

\subsection{Bounds and cofinal landing}

\begin{proposition}[Norm controls]
\label{prop:v13v98-bounds}
In submultiplicative matrix-form norms,
\[
 \|\Gamma_0\|\leq\|F\|\,\|dE\|,\qquad
 \|G^{-1}\|\leq\|F\|^2\|\eta^{-1}\|,\qquad
 \|K\|\leq\|G^{-1}\|\,\|H\|.                                \tag{98.19}
\]
Moreover,
\[
 \|{\cal R}_{\Gamma}\|
 \leq\|dK\|+2\|\Gamma_0\|\,\|K\|+\|K\|^2.                   \tag{98.20}
\]
\end{proposition}

\begin{proof}
The first three estimates follow from (98.1), (98.4), and (98.11).
Take norms in (98.13) for the last estimate.
\end{proof}

\begin{theorem}[Cofinal landing of the classified connection]
\label{thm:v13v98-landing}
Suppose, in common Banach form spaces,
\[
 E_n\to E,\quad F_n\to F,\quad dE_n\to dE,\quad
 H_n\to H,\quad dH_n\to dH,                                 \tag{98.21}
\]
with \(H_n^*=-H_n\), and suppose the differentiated inverse/metric
factors appearing in \(d(G_n^{-1}H_n)\) converge as supplied by the
V93 inverse-jet calculus.  Then
\[
 G_n\to G,\quad
 \Gamma_{0,n}\to\Gamma_0,\quad
 K_n=G_n^{-1}H_n\to K,\quad
 {\cal R}_{\Gamma_n}\to{\cal R}_{\Gamma}.                    \tag{98.22}
\]
If also \(\theta_n\to\theta\), \(d\theta_n\to d\theta\), then
\({\cal T}_{\Gamma_n}\to{\cal T}_{\Gamma}\).
\end{theorem}

\begin{proof}
Every expression in (98.22) is a finite sum of ordered words in the
listed factors and their first derivatives.  Apply
Theorem~\ref{thm:v13v96-wordlanding}.  The torsion statement follows
from (98.14) by the same argument.
\end{proof}

\subsection{Acquisition ledger and next local cell}

\[
\begin{array}{p{0.28\linewidth}|p{0.23\linewidth}|p{0.39\linewidth}}
\textbf{gate}&\textbf{status after V98}&\textbf{content}\\ \hline
\text{metricity}
 &\text{classified}
 &\Gamma=\Gamma_0+G^{-1}H,\ H^*=-H\\
\text{coframe-only connection}
 &\text{flat}
 &{\cal R}_{\Gamma_0}=0\\
\text{curvature carrier}
 &\text{identified}
 &\nabla^{\mathrm{ad}}_{\Gamma_0}K+K^2\\
\text{torsion-free metric connection}
 &\text{reduced to a linear gate}
 &{\cal L}_{G,\theta}(H)=-{\cal T}_0\\
\text{existence and uniqueness}
 &\text{open globally}
 &range and kernel of \({\cal L}_{G,\theta}\)\\
\text{cofinal connection landing}
 &\text{conditional}
 &convergence of \(E,F,H\) and their first jets
\end{array}                                                     \tag{98.23}
\]

The next version analyzes the selector near a central orthonormal
coframe, where the classical Cartan map has an explicit inverse.  A
Neumann cell around that anchor can then provide a genuine local
existence-and-uniqueness theorem.  Global coverage will remain a
separate problem; the failure of a local cell will trigger an upstream
change of calculus rather than an assertion of Levi--Civita
existence.

\subsection{Parent-integrity record}

The installed parent was
\[
\begin{array}{c|c}
 \text{file}&
 \texttt{Renewal\_SM\_Einstein\_Continuum\_Research\_Notes\_V97.tex}\\
 \text{lines}&29077\\
 \text{bytes}&1021935\\
 \text{SHA--256}&
 \texttt{CD57CBBD7157D7C43C083896D13F4E7D33C36E2469C21F72C8BAB13BC5F2300E}.
\end{array}                                                     \tag{98.24}
\]

% =============================================================================
% END APPENDED THIRTEENTH-CYCLE V98 MATERIAL
% =============================================================================
% =============================================================================
% BEGIN APPENDED THIRTEENTH-CYCLE V99 MATERIAL
% =============================================================================

\section{Thirteenth-cycle V99: the Cartan selector inverse and a
noncommutative Neumann cell}
\label{sec:v13v99}

\subsection{The central orthonormal anchor}

Let \(\theta_0^1,\ldots,\theta_0^m\) be a central coframe basis and
let \(\eta_{ab}\) be a constant nondegenerate symmetric matrix.
Write
\[
 d\theta_0^a
 =-\frac12 C^a{}_{bc}\,
   \theta_0^b\theta_0^c,\qquad
 C^a{}_{bc}=-C^a{}_{cb}.                                    \tag{99.1}
\]
For a connection
\[
 \omega^a{}_b=\Gamma^a{}_{bc}\theta_0^c,
\]
lower the first index with \(\eta\):
\(\Gamma_{abc}:=\eta_{ad}\Gamma^d{}_{bc}\) and
\(C_{abc}:=\eta_{ad}C^d{}_{bc}\).  Metric compatibility at this
anchor is
\[
 \Gamma_{abc}=-\Gamma_{bac}.                                 \tag{99.2}
\]

\begin{theorem}[Explicit Cartan coefficient formula]
\label{thm:v13v99-cartan}
With conventions (99.1)--(99.2), there is a unique
metric-compatible torsion-free connection, and its coefficients are
\[
 \boxed{
 \Gamma_{abc}
 =\frac12\bigl(-C_{abc}-C_{bca}+C_{cab}\bigr).}              \tag{99.3}
\]
\end{theorem}

\begin{proof}
The coefficient of
\(\frac12\theta_0^b\theta_0^c\) in
\(\omega^a{}_b\theta_0^b\) is
\(\Gamma^a{}_{cb}-\Gamma^a{}_{bc}\).  Thus
\(d\theta_0^a+\omega^a{}_b\theta_0^b=0\) is equivalent to
\[
 \Gamma_{abc}-\Gamma_{acb}=-C_{abc}.                         \tag{99.4}
\]
Using (99.2), the three cyclic instances of (99.4) are
\[
\begin{aligned}
 \Gamma_{abc}+\Gamma_{cab}&=-C_{abc},\\
 \Gamma_{cab}+\Gamma_{bca}&=-C_{cab},\\
 \Gamma_{bca}+\Gamma_{abc}&=-C_{bca}.
\end{aligned}                                                \tag{99.5}
\]
The first plus the third minus the second gives (99.3).
Conversely, antisymmetry of \(C_{abc}\) in \(b,c\) shows directly
that (99.3) satisfies (99.2), and substitution gives (99.4).
The same three-equation calculation shows uniqueness.
\end{proof}

\subsection{Invertibility of the anchor selector}

Let \({\cal X}\) be the real vector space of coefficient one-form
matrices \(K\) satisfying
\[
 K^*\eta+\eta K=0,
\]
and let \({\cal Y}\) be the real vector space of coframe-valued
two-forms.  Define
\[
 {\cal L}_0:{\cal X}\longrightarrow{\cal Y},\qquad
 {\cal L}_0(K):=K\theta_0.                                  \tag{99.6}
\]

\begin{proposition}[The central Cartan selector is an isomorphism]
\label{prop:v13v99-isomorphism}
The map \({\cal L}_0\) in (99.6) is a real-linear isomorphism.
Equivalently, for every prescribed torsion source
\(Y\in{\cal Y}\), there is a unique metric-skew \(K\) satisfying
\[
 K\theta_0=Y.                                                \tag{99.7}
\]
\end{proposition}

\begin{proof}
Expand \(Y\) in the central basis and absorb its coefficients into
a structure tensor \(C\) with the antisymmetry in (99.1).  Equation
(99.7), together with metric skewness, is the same linear system as
(99.2)--(99.4), with the right-hand side determined by \(Y\).
Formula (99.3) gives a solution linear in \(Y\), and the uniqueness
part of Theorem~\ref{thm:v13v99-cartan} gives a trivial kernel.
Hence \({\cal L}_0\) is both surjective and injective.
\end{proof}

To compare directly with the fixed-domain selector in V98, let
\[
 {\cal X}_{\mathrm{ah}}
 :=\{H:H^*=-H\},\qquad
 \widetilde{\cal L}_0(H):=\eta^{-1}H\theta_0.                \tag{99.8}
\]
The bijection \(H=\eta K\) identifies
\(\widetilde{\cal L}_0\) with \({\cal L}_0\), so it is also an
isomorphism.  Fix norms and put
\[
 \kappa_0:=\|\widetilde{\cal L}_0^{-1}\|<\infty.             \tag{99.9}
\]

\subsection{The noncommutative selector cell}

For a self-adjoint invertible metric coefficient \(G\) and a coframe
column \(\theta\), define on the same real domain
\({\cal X}_{\mathrm{ah}}\)
\[
 \widetilde{\cal L}_{G,\theta}(H):=G^{-1}H\theta.            \tag{99.10}
\]

\begin{lemma}[Explicit perturbation radius]
\label{lem:v13v99-radius}
In compatible submultiplicative norms,
\[
 \|\widetilde{\cal L}_{G,\theta}
      -\widetilde{\cal L}_0\|
 \leq
 \|G^{-1}-\eta^{-1}\|\,\|\theta\|
 +\|\eta^{-1}\|\,\|\theta-\theta_0\|
 =:\Delta(G,\theta).                                        \tag{99.11}
\]
\end{lemma}

\begin{proof}
For \(H\in{\cal X}_{\mathrm{ah}}\),
\[
\begin{aligned}
 G^{-1}H\theta-\eta^{-1}H\theta_0
 &=(G^{-1}-\eta^{-1})H\theta
   +\eta^{-1}H(\theta-\theta_0).
\end{aligned}
\]
Take norms, factor out \(\|H\|\), and take the supremum over
\(\|H\|=1\).
\end{proof}

\begin{theorem}[Quantitative noncommutative Levi--Civita cell]
\label{thm:v13v99-cell}
If
\[
 \boxed{\quad
 \kappa_0\Delta(G,\theta)<1,
 \quad}                                                       \tag{99.12}
\]
then \(\widetilde{\cal L}_{G,\theta}\) is invertible and
\[
 \|\widetilde{\cal L}_{G,\theta}^{-1}\|
 \leq
 \frac{\kappa_0}
 {1-\kappa_0\Delta(G,\theta)}.                               \tag{99.13}
\]
Consequently the simultaneous metricity and torsion equations of
V98 have the unique solution
\[
 H=-\widetilde{\cal L}_{G,\theta}^{-1}{\cal T}_0,\qquad
 \Gamma=\Gamma_0+G^{-1}H,                                   \tag{99.14}
\]
with
\[
 \|H\|
 \leq
 \frac{\kappa_0\|{\cal T}_0\|}
 {1-\kappa_0\Delta(G,\theta)}.                               \tag{99.15}
\]
\end{theorem}

\begin{proof}
Factor the perturbed selector as
\[
 \widetilde{\cal L}_{G,\theta}
 =\widetilde{\cal L}_0
 \left[
 I+\widetilde{\cal L}_0^{-1}
   (\widetilde{\cal L}_{G,\theta}-\widetilde{\cal L}_0)
 \right].
\]
By (99.9)--(99.12), the operator inside brackets is within distance
strictly less than one of the identity.  Its Neumann series converges,
so it is invertible and has inverse norm at most
\[
 \frac1{1-\kappa_0\Delta(G,\theta)}.
\]
This proves (99.13).  Equation (99.14) is exactly the selector
equation (98.16), and (99.15) follows from (99.13).
\end{proof}

\begin{remark}[Scope of the Levi--Civita terminology]
Theorem~\ref{thm:v13v99-cell} uses the term Levi--Civita cell for the
unique solution of the declared metricity and torsion equations
inside the specified differential calculus.  It does not assert
global existence, coordinate independence across undeclared
calculi, or equality with a spectral-triple connection.
\end{remark}

\subsection{Cofinal stability inside a uniform cell}

\begin{theorem}[Landing of selector solutions]
\label{thm:v13v99-landing}
Let
\[
 L_n:=\widetilde{\cal L}_{G_n,\theta_n},\qquad
 L:=\widetilde{\cal L}_{G,\theta}
\]
act between common finite-dimensional normed spaces.  Suppose
\[
 L_n\to L,\qquad
 {\cal T}_{0,n}\to{\cal T}_0,\qquad
 \sup_n\|L_n^{-1}\|\leq M.                                  \tag{99.16}
\]
If \(L\) is invertible, then the unique selector solutions
\(H_n=-L_n^{-1}{\cal T}_{0,n}\) satisfy
\[
 H_n\longrightarrow H=-L^{-1}{\cal T}_0.                    \tag{99.17}
\]
The associated metric-compatible torsion-free connections and
curvatures land under the hypotheses of
Theorem~\ref{thm:v13v98-landing}.
\end{theorem}

\begin{proof}
The resolvent identity gives
\[
 L_n^{-1}-L^{-1}=L_n^{-1}(L-L_n)L^{-1},
\]
so \(L_n^{-1}\to L^{-1}\) by (99.16).  Therefore
\[
\begin{aligned}
 \|H_n-H\|
 &\leq
 \|L_n^{-1}\|\,\|{\cal T}_{0,n}-{\cal T}_0\|
 +\|L_n^{-1}-L^{-1}\|\,\|{\cal T}_0\|
 \longrightarrow0.
\end{aligned}
\]
Apply Theorem~\ref{thm:v13v98-landing} for the final assertion.
\end{proof}

A uniform version of (99.12) supplies the inverse bound in (99.16).
This is the connection analogue of retaining a positive relative
margin in the companion YM programme.  Pointwise membership in the
cell is insufficient for continuum landing unless the margin is
uniform on the physical region and through the cofinal stages.

\subsection{Acquisition ledger and global obstruction}

\[
\begin{array}{p{0.28\linewidth}|p{0.23\linewidth}|p{0.39\linewidth}}
\textbf{gate}&\textbf{status after V99}&\textbf{content}\\ \hline
\text{central Cartan selector}
 &\text{acquired}
 &explicit inverse (99.3)\\
\text{local noncommutative selector}
 &\text{acquired}
 &Neumann condition \(\kappa_0\Delta<1\)\\
\text{metricity plus torsion}
 &\text{uniquely solvable in the cell}
 &formula and norm bound (99.14)--(99.15)\\
\text{cofinal connection landing}
 &\text{acquired inside a uniform cell}
 &uniform inverse margin and jet convergence\\
\text{global cell coverage}
 &\text{open}
 &cover the physical coframe spectrum by compatible anchors\\
\text{Einstein contraction}
 &\text{still open}
 &construct \(P,s\) for the landed connection and prove
   the V97 defect relation
\end{array}                                                     \tag{99.18}
\]

Failure of (99.12) would not prove nonexistence of a
metric-compatible torsion-free connection.  It would show only that
the chosen central anchor does not certify it.  The next alternatives
are additional algebraic anchors, relative positive-order cells for
\(G\), or an upstream change of differential calculus.  No global
Levi--Civita assertion is made.

\subsection{Parent-integrity record}

The installed parent was
\[
\begin{array}{c|c}
 \text{file}&
 \texttt{Renewal\_SM\_Einstein\_Continuum\_Research\_Notes\_V98.tex}\\
 \text{lines}&29398\\
 \text{bytes}&1031973\\
 \text{SHA--256}&
 \texttt{C7F52D1EBB2F8CE09064CE62B127F2CF5641D7F9A7942240A34B21FA80B042CF}.
\end{array}                                                     \tag{99.19}
\]

% =============================================================================
% END APPENDED THIRTEENTH-CYCLE V99 MATERIAL
% =============================================================================
% =============================================================================
% BEGIN APPENDED THIRTEENTH-CYCLE V100 MATERIAL
% =============================================================================

\section{Thirteenth-cycle V100: compulsory audit, cycle closure, and
uniform-selector direction}
\label{sec:v13v100}

\subsection{Current companion records}

The five-version companion audit was repeated on 10 September 2026
after V99.  The current active companion records are
\[
\begin{array}{p{0.39\linewidth}|r|r|p{0.27\linewidth}}
\textbf{file}&\textbf{lines}&\textbf{bytes}&\textbf{SHA--256}\\ \hline
\texttt{Renewal\_Yang\_Mills\_Mass\_Gap\_Research\_Notes\_V97.tex}
 &31165&1063317&
 \texttt{731EA8552A0E36269EC0407C9B776C673444BB5C527D3B498C78FC7AF475738A}\\
\texttt{Renewal\_NS\_Stability\_Research\_Notes\_V26.tex}
 &5319&278283&
 \texttt{82C887F8C2C69E4F28FAD7C3DC8DE5B9099CB5A4BF431EBA1FFF3E91935978AD}.
\end{array}                                                    \tag{100.1}
\]

The YM V97 file has the same line count, byte count, SHA--256 digest,
and terminal mathematical section as the YM V96 file inspected at
V95.  It is therefore byte-identical to that audited source despite
the changed filename.  No new theorem is inferred from the version
label.  Its operative conclusion remains the failure of the absolute
Voronoi-radius architecture and the upstream pivot to relative
Loewner pencils.

The NS note is unchanged.  Its operative conclusion remains the
rank-one derivative-norm calculation and the small, localized
improvement for the second derivative.  No new companion constant
enters the connection selector.

\begin{proposition}[V100 audit decision]
\label{prop:v13v100-decision}
The next SM cycle shall continue from the V99 selector theorem by
seeking a uniform cover of the physical coframe/metric range.  The
first cover variable should be relative metric order, with absolute
operator-norm distance retained as a fallback sufficient condition.
Any structured-input improvement must be propagated through the
actual selector inverse and curvature words before it is credited.
\end{proposition}

\begin{proof}
The YM result shows that an absolute radius can erase cancellation
already present at the anchor, while its relative Hermitian pencils
retain the relevant base-operator scale.  V99 has exactly the same
structural dependence: the useful quantity is the inverse margin of
\(\widetilde{\cal L}_{G,\theta}\), not a large independent bound on
each coefficient.  The NS result shows that a restricted norm must be
evaluated on the realized input and may yield only a small gain.
Together these facts justify the stated order of attack without
importing a Yang--Mills or Navier--Stokes conclusion.
\end{proof}

\subsection{Cycle result and unresolved gates}

The V91--V100 block has established:
\begin{enumerate}
\item the distinction between pure central metric sectors and genuine
      noncentral coframe metrics;
\item quantitative invertibility and first/second inverse-jet landing
      for \(G=E^*\eta E\);
\item an exact noncommutative Cartan calculus, both Bianchi identities,
      and their finite-stage nilpotence/Leibniz defect formulas;
\item the universal cyclic trace and the proof that it does not by
      itself yield an Einstein tensor;
\item the unique symmetric two-word trace reversal and its factorized
      covariant-conservation defect;
\item the affine classification of metric-compatible connections,
      the flatness of the pure-gauge connection, and the exact torsion
      selector;
\item the explicit classical Cartan inverse and a quantitative
      noncommutative Neumann cell with cofinal landing.
\end{enumerate}

The full SM--Einstein continuum is not yet proved.  The live gates are
\[
\begin{array}{p{0.31\linewidth}|p{0.57\linewidth}}
\textbf{gate}&\textbf{required next evidence}\\ \hline
\text{uniform connection coverage}
 &a cofinal physical-region cover with a positive selector-inverse
  margin\\
\text{ordered Ricci/scalar maps}
 &frame covariance, adjoint compatibility, classical reduction,
  and norm landing\\
\text{contracted Bianchi}
 &vanishing of the V97 Ricci/scalar/metricity defect combination\\
\text{source coupling}
 &identification with the already constructed all-order
  gauge-invariant stress/current ledger\\
\text{joint continuum law}
 &positivity, compatibility, locality, OS/Hamiltonian consistency,
  and Einstein support on the same limiting state
\end{array}                                                     \tag{100.2}
\]

The programme will now obey the requested size-control rule.  This
complete V100 source will be frozen with suffix
\(\texttt{\_f13}\), and the active note will restart at V1 with a
compact context and result summary before continuing the selector
cover.  The frozen source remains an immutable proof archive.

\subsection{Parent-integrity record}

The installed parent was
\[
\begin{array}{c|c}
 \text{file}&
 \texttt{Renewal\_SM\_Einstein\_Continuum\_Research\_Notes\_V99.tex}\\
 \text{lines}&29703\\
 \text{bytes}&1041592\\
 \text{SHA--256}&
 \texttt{7AD2655A4A3174CF15227F5123911771CF238AC4D130DBF164694FC1CAA8C635}.
\end{array}                                                     \tag{100.3}
\]

% =============================================================================
% END APPENDED THIRTEENTH-CYCLE V100 MATERIAL
% =============================================================================
\end{document}
